\documentclass[11pt,a4paper,openany]{book}
\usepackage[margin=25mm]{geometry}
\usepackage{fontspec}
\setmainfont{OLMarathiSerif-Regular.ttf}[Path=../../fonts/,Script=Devanagari,Language=Marathi,BoldFont=OLMarathiSerif-Bold.ttf,ItalicFont=OLMarathiSerif-Regular.ttf]
\setsansfont{OLMarathiSerif-Regular.ttf}[Path=../../fonts/,Script=Devanagari,Language=Marathi,BoldFont=OLMarathiSerif-Bold.ttf]
\usepackage{amsmath,amssymb,amsthm,nicefrac,graphicx,tikz}
\usepackage{stmaryrd}
\usepackage{../../upstream/sty/bussproofs-extra}
\usepackage[tableaux]{prooftrees}
\usepackage{flafter}
\usepackage[unicode,colorlinks=true,linkcolor=blue,urlcolor=blue]{hyperref}
\usepackage{microtype}
\setlength{\parindent}{0pt}
\setlength{\parskip}{0.45em}
\linespread{1.18}
\setlength{\emergencystretch}{2em}
\renewcommand{\contentsname}{अनुक्रमणिका}
\renewcommand{\chaptername}{प्रकरण}
\renewcommand{\figurename}{आकृती}
\renewcommand{\proofname}{सिद्धता}
\theoremstyle{definition}
\newtheorem{defn}{व्याख्या}[section]
\newtheorem{ex}[defn]{उदाहरण}
\newtheorem{prop}[defn]{प्रतिज्ञा}
\newtheorem{thm}[defn]{प्रमेय}
\newtheorem{prob}{सराव}[chapter]
\newenvironment{explain}{}{}
\newenvironment{digress}{\par\smallskip}{\par\smallskip}
\newcommand{\Setabs}[2]{\{#1:#2\}}
\newcommand{\Pow}[1]{\wp(#1)}
\newcommand{\tuple}[1]{\langle #1\rangle}
\newcommand{\Nat}{\mathbb{N}}
\newcommand{\Int}{\mathbb{Z}}
\newcommand{\Rat}{\mathbb{Q}}
\newcommand{\Real}{\mathbb{R}}
\newcommand{\PosInt}{\mathbb{Z}^{+}}
\newcommand{\Bin}{\mathbb{B}}
\newcommand{\len}[1]{\mathrm{len}(#1)}
\newcommand{\lif}{\mathbin{\to}}
\colorlet{oldiagcolorA}{black}
\colorlet{oldiagcolorB}{gray}
\definecolor{oldiagcolorC}{HTML}{a81c21}
\newcommand{\olasset}[2][0.52\textwidth]{\centering\resizebox{#1}{!}{\input{../../upstream/#2}}}
\hypersetup{pdftitle={मुक्त तर्कशास्त्र: संच, संबंध, फलने, संचांचे आकारमान, अंकगणितीकरण, अनंत संच, विधानीय तर्कशास्त्र, चार सिद्धता-पद्धती, संपूर्णता, प्रथम-क्रम विन्यासमीमांसा व चिन्हार्थमीमांसा, प्रतिमाने व उपपत्ती आणि प्रतिमान उपपत्ती},pdfauthor={Open Logic Project; Marathi machine translation by Codex}}
\newenvironment{intro}{}{}
\newcommand{\Id}[1]{\mathord{\mathrm{Id}_{#1}}}
\newcommand{\emptyseq}{\Lambda}
\newcommand{\liff}{\mathbin{\leftrightarrow}}
\newcommand{\equivrep}[2]{[#1]_{#2}}
\newcommand{\equivclass}[2]{#1/_{\!{#2}}}
\newcommand{\funrestrictionto}[2]{#1\mathord{\restriction}_{#2}}
\newcommand{\funimage}[2]{#1[#2]}
\newcommand{\citeyear}[1]{\hyperref[bib:#1]{1965}}
\newcommand{\dom}[1]{\mathrm{dom}(#1)}
\newtheorem{cor}[defn]{निष्कर्ष}
\newcommand{\ran}[1]{\mathrm{ran}(#1)}
\newcommand{\comp}[2]{#2\circ #1}
\newcommand{\pto}{\mathrel{\ooalign{\hfil$\mapstochar\mkern 5mu$\hfil\cr$\to$}}}
\newcommand{\fdefined}{\downarrow}
\newcommand{\fundefined}{\uparrow}
\newcommand{\cardle}[2]{#1 \preceq #2}
\newcommand{\cardless}[2]{#1 \prec #2}
\newcommand{\cardeq}[2]{#1 \approx #2}
\newcommand{\cardneq}[2]{#1 \not\approx #2}
\newcommand{\Intequiv}{\sim}
\newcommand{\Ratequiv}{\backsim}
\newcommand{\Realequiv}{\Bumpeq}
\newcommand{\defis}{=}
\newcommand{\closureofunder}[2]{\mathrm{clo}_{#1}(#2)}
\newcommand{\Closureofunder}[2]{\mathrm{Clo}_{#1}(#2)}
\newtheorem{lem}[defn]{सहायक प्रमेय}
\newenvironment{editorial}{\begin{quote}\small\itshape}{\end{quote}}
\definecolor{oldiagcolorD}{HTML}{1973ba}
\newlength{\olphotowidth}\setlength{\olphotowidth}{0.35\textwidth}

\newcommand{\applytofirst}[2]{{\expandafter#1#2}}
\newcommand{\Struct}[1]{\applytofirst{\mathfrak}{#1}}
\newcommand{\Lang}[1]{\applytofirst{\mathcal}{#1}}
\newcommand{\Obj}[1]{\mathsf{#1}}
\newcommand{\PVar}{\mathrm{At}_0}
\NewDocumentCommand{\Frm}{o}{\IfNoValueTF{#1}{\mathrm{Frm}}{\mathrm{Frm}(\Lang{#1})}}
\newcommand{\True}{\ensuremath{\mathbb{T}}}
\newcommand{\False}{\ensuremath{\mathbb{F}}}
\newcommand{\lfalse}{\bot}
\newcommand{\ltrue}{\top}
\newcommand{\ident}{\equiv}
\newcommand{\subst}[2]{#1/#2}
\newcommand{\SSubst}[2]{#1[#2]}
\newcommand{\Subst}[3]{#1[\subst{#2}{#3}]}
\NewDocumentCommand{\Entails}{t{/} o}{%
  \IfBooleanTF{#1}{\IfNoValueTF{#2}{\nvDash}{\nvDash_{#2}}}{\IfNoValueTF{#2}{\vDash}{\vDash_{#2}}}}
\newcommand{\pAssign}[1]{\applytofirst{\mathfrak}{#1}}
\NewDocumentCommand{\pValue}{m d()}{\overline{\pAssign{#1}}\IfNoValueF{#2}{(#2)}}
\NewDocumentCommand{\pSat}{t{/} m m}{\pAssign{#2}\IfBooleanTF{#1}{\nvDash}{\vDash}#3}
\newcommand{\indfrm}{}
\NewDocumentCommand{\indcase}{m m +m}{\renewcommand{\indfrm}{#1}$#1 \ident #2$: #3}
\def\startycommalist{\def\ycomma{\def\ycomma{, }}}

\newcommand{\Sequent}{\Rightarrow}
\renewcommand{\fCenter}{\ensuremath{\,\Sequent\,}}
\NewDocumentCommand{\LeftR}{m}{\ensuremath{{#1}\mathrm{L}}}
\NewDocumentCommand{\RightR}{m}{\ensuremath{{#1}\mathrm{R}}}
\newcommand{\Weakening}{\ensuremath{\mathrm{W}}}
\NewDocumentCommand{\Intro}{m o}{\ensuremath{{#1}\mathrm{Intro}\IfNoValueTF{#2}{}{_{#2}}}}
\NewDocumentCommand{\Elim}{m o}{\ensuremath{{#1}\mathrm{Elim}\IfNoValueTF{#2}{}{_{#2}}}}
\newcommand{\FalseInt}{\ensuremath{\lfalse_I}}
\NewDocumentCommand{\Discharge}{m m}{[#1]^{#2}}
\NewDocumentCommand{\DischargeRule}{m m}{\RightLabel{#1}\LeftLabel{\scriptsize $#2$}}
\NewDocumentCommand{\sFmla}{m m o}{\ensuremath{\IfNoValueTF{#3}{}{#3\,}\hbox to.8em{\ensuremath{#1}\hfil} #2}}
\NewDocumentCommand{\TRule}{m m o}{\ensuremath{{#2}{#1}\IfNoValueTF{#3}{}{\, #3}}}
\newcommand{\TAss}{\text{गृहीतक}}
\NewDocumentCommand{\Proves}{t{/} o}{%
  \IfBooleanTF{#1}{\IfNoValueTF{#2}{\nvdash}{\nvdash_{#2}}}{\IfNoValueTF{#2}{\vdash}{\vdash_{#2}}}}
\newcommand{\formula}[1]{#1}
\renewenvironment{prooftree}{\begin{center}\bottomAlignProof}{\DisplayProof\end{center}}
\newenvironment{oltableau}{\center\tableau{}}{\endtableau\endcenter}
\newenvironment{derivation}{~\begin{trivlist}\item\begin{tabular}[b]{@{}rll@{}}}{\end{tabular}\end{trivlist}}

\NewDocumentCommand{\lexists}{t{!} o o}{%
  \exists\IfBooleanT{#1}{\mathord{!}}%
  \IfNoValueF{#2}{#2}\IfNoValueF{#3}{\,#3}}
\NewDocumentCommand{\lforall}{o o}{%
  \IfNoValueTF{#1}{\forall}{\forall #1}\IfNoValueF{#2}{\,#2}}
\NewDocumentCommand{\eq}{t{/} o o}{%
  \IfNoValueTF{#3}{\IfBooleanTF{#1}{\neq}{=}}{%
    \IfBooleanTF{#1}{#2\neq#3}{#2=#3}}}
\NewDocumentCommand{\Atom}{m m}{\mathord{#1}(#2)}
\NewDocumentCommand{\Assign}{m m}{\mathord{#1^{\Struct{#2}}}}
\NewDocumentCommand{\Sat}{t{/} m m o}{%
  \IfBooleanTF{#1}{%
    \IfNoValueTF{#4}{\Struct{#2}\nvDash#3}{\Struct{#2},#4\nvDash#3}}{%
    \IfNoValueTF{#4}{\Struct{#2}\vDash#3}{\Struct{#2},#4\vDash#3}}}
\NewDocumentCommand{\varAssign}{m m m o}{%
  \IfNoValueTF{#4}{#1\sim_{#3}#2}{#1=#2[^{#4}/#3]}}
\NewDocumentCommand{\Value}{m m o}{%
  \IfNoValueTF{#3}{\mathrm{Val}^{\Struct{#2}}(#1)}{%
    \mathrm{Val}^{\Struct{#2}}_{#3}(#1)}}
\NewDocumentCommand{\Log}{m o}{%
  \ensuremath{\mathbf{#1}\IfNoValueF{#2}{_{#2}}}}
\newcommand{\Contraction}{\ensuremath{\mathrm{C}}}
\newcommand{\Exchange}{\ensuremath{\mathrm{X}}}
\newcommand{\Cut}{\ensuremath{\mathrm{Cut}}}
\newcommand{\FalseCl}{\ensuremath{\lfalse_C}}
\newenvironment{defish}{}{}
\newcommand{\readerexternalref}[1]{\textit{मूळ ग्रंथातील संबंधित विभाग}}
\makeatletter
\renewcommand{\@pnumwidth}{2em}
\renewcommand*{\l@section}{\@dottedtocline{1}{1.5em}{2.8em}}
\makeatother

\newtheorem{rem}[defn]{टीप}
\newenvironment{history}{\begin{quote}\small}{\end{quote}}
\newcommand{\MP}{\textsc{mp}}
\newcommand{\QR}{\textsc{qr}}
\newcommand{\Hyp}{\textsc{Hyp}}
\newcommand{\Nec}{\textsc{nec}}
\newcommand{\Var}{\mathrm{Var}}
\NewDocumentCommand{\Trm}{o}{\IfNoValueTF{#1}{\mathrm{Trm}}{\mathrm{Trm}(\Lang{#1})}}
\NewDocumentCommand{\Sent}{o}{\IfNoValueTF{#1}{\mathrm{Sent}}{\mathrm{Sent}(\Lang{#1})}}
\newcommand{\Th}[1]{\mathbf{#1}}
\newcommand{\Theory}[1]{\mathrm{Th}(\Struct{#1})}
\newcommand{\substruct}{\subseteq}
\NewDocumentCommand{\elemequiv}{t{/} o}{%
  \IfBooleanTF{#1}{\IfNoValueTF{#2}{\not\equiv}{\not\equiv_{#2}}}{%
  \IfNoValueTF{#2}{\equiv}{\equiv_{#2}}}}
\NewDocumentCommand{\iso}{t{/} o}{%
  \IfBooleanTF{#1}{\IfNoValueTF{#2}{\not\simeq}{\not\simeq_{#2}}}{%
  \IfNoValueTF{#2}{\simeq}{\simeq_{#2}}}}
\newcommand{\QuantRank}[1]{\mathrm{qr}(#1)}
\newcommand{\Domain}[1]{\left|\Struct{#1}\right|}
\newcommand{\Expan}[2]{(\Struct{#1},#2)}
\newcommand{\nszero}{\mathbf{z}}
\newcommand{\nssucc}{*}
\newcommand{\nsplus}{\oplus}
\newcommand{\nstimes}{\otimes}
\newcommand{\nsless}{\varolessthan}
\newcommand{\concat}{\frown}
\NewDocumentCommand{\lambd}{o o}{%
  \IfNoValueTF{#1}{\lambda}{\lambda #1}\IfNoValueF{#2}{.\,#2}}
\newcommand{\num}[1]{\overline{#1}}
\newcommand{\mModel}[1]{\applytofirst{\mathfrak}{#1}}
\NewDocumentCommand{\mSat}{t{/} m m o}{%
  \IfBooleanTF{#1}{\IfNoValueTF{#4}{\mModel{#2}\nVdash#3}{\mModel{#2},#4\nVdash#3}}{%
  \IfNoValueTF{#4}{\mModel{#2}\Vdash#3}{\mModel{#2},#4\Vdash#3}}}
\RenewDocumentCommand{\indcase}{s t{!} m m +m}{%
  \renewcommand{\indfrm}{#3}%
  \def\indfrmp{#3}%
  \def\indcomplex{#4}%
  \IfBooleanTF{#1}{$#3$ हे आण्विक सूत्र आहे: }{$#3 \ident #4$: }%
  \IfBooleanTF{#2}{सराव.}{#5}}
\NewDocumentCommand{\OPrf}{o}{\mathsf{Prf}\IfNoValueF{#1}{_{#1}}}
\NewDocumentCommand{\OCon}{o}{\mathsf{Con}\IfNoValueF{#1}{_{#1}}}
\newcommand{\PAx}{\mathrm{Ax}_0}
\newcommand{\PIso}[1]{\mathcal{#1}}
\newcommand{\Part}[2]{\Atom{\Obj P}{#1,#2}}
\newcommand{\gn}[1]{\ulcorner #1\urcorner}
\begin{document}
\begin{titlepage}
\vspace*{2cm}
{\Huge\bfseries मुक्त तर्कशास्त्र\par}
\vspace{1cm}
{\Large संच, संबंध, फलने, संचांचे आकारमान, अंकगणितीकरण, अनंत संच, विधानीय तर्कशास्त्र, चार सिद्धता-पद्धती, संपूर्णता, प्रथम-क्रम विन्यासमीमांसा व चिन्हार्थमीमांसा, प्रतिमाने व उपपत्ती आणि प्रतिमान उपपत्ती\par}
\vspace{1cm}
{\large मराठी आवृत्ती · वीस संपूर्ण प्रकरणे\par}
\vfill
मूळ ग्रंथ: Open Logic Project, \textit{The Open Logic Text}.\par
या आवृत्तीची तांत्रिक स्रोत-ओळख प्रकल्पाच्या नोंदींत दिली आहे.\par
हे प्रकाशन संच, संबंध, फलने, संचांचे आकारमान, अंकगणितीकरण, अनंत संच, विधानीय तर्कशास्त्र, चार सिद्धता-पद्धती, संपूर्णता, प्रथम-क्रम विन्यासमीमांसा व चिन्हार्थमीमांसा, प्रतिमाने व उपपत्ती आणि प्रतिमान उपपत्ती या प्रकरणांचे आहे: मूळ 722 स्रोत-एककांपैकी OLP-0004 ते OLP-0197, एकूण 194 स्रोत-एकके आणि 171 वाचक-विभाग. उर्वरित 528 विभागांचे काम अपूर्ण आहे.
संपूर्ण मराठी ग्रंथाचे काम सुरू आहे.\par
Codex कडून यंत्रानुवाद; स्रोताशी तुलना आणि यांत्रिक तपासण्या केल्या आहेत.
स्वतंत्र मानवी संपादनाचा दावा केलेला नाही. काही तांत्रिक संज्ञा तात्पुरत्या आहेत.\par
मूळ मजकूर आणि हे रूपांतर: Creative Commons Attribution 4.0.
मूळ घटकांचे स्वतंत्र परवाने लागू राहतात.\par
\url{https://github.com/OpenLogicProject/OpenLogic}\par
\url{https://github.com/KokunoYumeto/OpenLogic-translations}
\end{titlepage}
\tableofcontents
\chapter{संच}








\section{विस्तारात्मकता}\label{sfr:set:bas:sec}

वस्तूंचा समूह एकच वस्तू {म्हणून} विचारात घेतला, की त्याला \emph{संच} म्हणतात.
संच बनवणाऱ्या वस्तूंना त्या संचाचे \emph{घटक} किंवा \emph{सदस्य} म्हणतात.
$x$ हा $A$ या संचाचा घटक असेल, तर आपण $x \in A$ असे लिहितो;
नसेल, तर $x \notin A$ असे लिहितो. ज्या संचात एकही घटक नसतो,
त्याला \emph{रिक्त} संच म्हणतात आणि तो ``$\emptyset$'' या चिन्हाने दर्शवतात.

\begin{explain}
आपण संच \emph{कसा निर्दिष्ट करतो}, त्यातील घटक \emph{कोणत्या क्रमाने}
मांडतो किंवा तेच घटक \emph{किती वेळा} मोजतो, याला महत्त्व नसते.
त्यात कोणते घटक आहेत, एवढेच महत्त्वाचे असते. ही गोष्ट आपण पुढील
तत्त्वात नेमकेपणाने मांडतो.
\end{explain}

\begin{defn}[विस्तारात्मकता]
  $A$ आणि $B$ हे संच असतील, तर $A = B$ तेव्हा आणि केवळ तेव्हाच, जेव्हा
  $A$ मधील प्रत्येक घटक हा $B$ चाही घटक असतो, आणि उलटही.
\end{defn}

विस्तारात्मकतेमुळे काही चिन्हपद्धती वापरणे योग्य ठरते. सर्वसाधारणपणे,
$a_{1}$, \dots, $a_{n}$ या वस्तू दिलेल्या असतील, तर
$\{a_{1}, \dots, a_{n}\}$ म्हणजे ज्यातील घटक हे
$a_1, \ldots, a_n$ आहेत \emph{तो विशिष्ट} संच होय. येथे ``\emph{तो विशिष्ट}''
या शब्दांवर भर आहे, कारण विस्तारात्मकतेनुसार असा \emph{एकच} संच असू शकतो.
विस्तारात्मकतेमुळे पुढील समानताही योग्य ठरते:
  \[
    \{a, a, b\} = \{a, b\} = \{b,a\}.
  \]
म्हणजेच, संचांचा विचार करताना त्यातील घटक कोणत्या क्रमाने किंवा
किती वेळा निर्दिष्ट केले आहेत, याला महत्त्व नसते.


\begin{ex}
काही वस्तू दिलेल्या असतील, तर त्या एकत्र करून त्यांचा संच बनवता येतो.
उदाहरणार्थ, रिचर्डच्या भावंडांच्या संचात एक व्यक्ती आहे, आणि तो
$S=\{\textrm{Ruth}\}$ असा लिहिता येईल.
$4$ पेक्षा लहान धन पूर्णांकांचा संच $\{1, 2, 3\}$ आहे; पण तो
$\{3, 2, 1\}$ किंवा अगदी $\{1, 2, 1, 2, 3\}$ असाही लिहिता येतो.
विस्तारात्मकतेनुसार हे सर्व एकच संच आहेत. कारण $\{1, 2, 3\}$ मधील
प्रत्येक घटक हा $\{3, 2, 1\}$ चा (आणि $\{1,
2, 1, 2, 3\}$ चाही) घटक आहे, आणि उलटही.
\end{ex}


अनेकदा संचातील सर्व घटक यांना समान असलेला एखादा गुणधर्म सांगून आपण
संच निर्दिष्ट करू. त्यासाठी पुढील संक्षिप्त चिन्हपद्धती वापरू:
$\Setabs{x}{\phi(x)}$. येथे $\phi(x)$ हा असा गुणधर्म दर्शवतो की,
$x$ ची संचातील घटक मध्ये गणना होण्यासाठी त्यात तो गुणधर्म असला पाहिजे.


\begin{ex}
आपल्या उदाहरणातील $S$ हा संच पुढीलप्रमाणेही निर्दिष्ट करता आला असता:
\[
S = \Setabs{x}{x \text{ हे रिचर्डचे भावंड आहे}}.
\]
\end{ex}



\begin{ex}
एखादी संख्या तिच्या स्वतःखेरीज इतर सर्व धन विभाजकांच्या बेरजेइतकी असेल तेव्हा
आणि केवळ तेव्हाच तिला \emph{परिपूर्ण} म्हणतात (हे विभाजक म्हणजे त्या संख्येला
निःशेष भाग जाणाऱ्या, पण तिच्याशी समान नसलेल्या संख्या).
उदाहरणार्थ, $6$ ही परिपूर्ण संख्या आहे, कारण तिचे असे विभाजक
$1$, $2$ आणि~$3$ आहेत, आणि $6 = 1 + 2 + 3$.
खरे तर, $10$ पेक्षा लहान धन पूर्णांकांपैकी $6$ हीच एकमेव परिपूर्ण संख्या आहे.
म्हणून विस्तारात्मकतेचा उपयोग करून आपण असे म्हणू शकतो:
\[
	\{6\} = \Setabs{x}{x\text{ ही परिपूर्ण संख्या आहे आणि }0 \leq x \leq 10}
\]
उजवीकडील चिन्हलेखनाचे वाचन असे करतो: ``$x$ ही परिपूर्ण संख्या आहे आणि
$0 \leq x \leq 10$ अशा सर्व $x$ चा संच.'' संच कसा निर्दिष्ट केला आहे,
याला संचाचा विचार करताना महत्त्व नसते, ही गोष्ट वरील समानता स्पष्ट करते.
अधिक सर्वसाधारणपणे, $\phi(x)$ असणाऱ्या सर्व $x$ चा संच नेहमी एकच असतो,
याची विस्तारात्मकता खात्री देते.
त्यामुळे $\Setabs{x}{\phi(x)}$ याला $\phi(x)$ असणाऱ्या सर्व $x$ चा
\emph{तो विशिष्ट} संच म्हणणे विस्तारात्मकतेमुळे योग्य ठरते.
\end{ex}


संच समान असल्याचे दाखवण्याची पद्धत विस्तारात्मकतेमुळे मिळते:
$A = B$ हे दाखवण्यासाठी, $x \in A$ असेल तेव्हा $x \in B$ देखील असते,
आणि $y \in B$ असेल तेव्हा $y \in A$ देखील असते, हे दाखवा.

\begin{prob}
जास्तीत जास्त एकच रिक्त संच असतो हे सिद्ध करा. म्हणजे, $A$ आणि $B$ या
संचांत एकही घटक नसेल, तर $A = B$ हे दाखवा.
\end{prob}











\section{उपसंच आणि घातसंच}\label{sfr:set:sub:sec}

\begin{explain}
आपल्याला अनेकदा संचांची तुलना करावी लागेल. तुलना करण्याचा एक सरळ प्रकार
असा आहे: \emph{एका संचातील प्रत्येक वस्तू दुसऱ्या संचातही असते}.
ही स्थिती पुरेशी महत्त्वाची असल्यामुळे तिच्यासाठी आपण नवीन चिन्हपद्धतीची
ओळख करून घेऊ.
\end{explain}

\begin{defn}[उपसंच]
$A$ या संचातील प्रत्येक घटक हा $B$ चाही घटक असेल,
तर $A$ हा $B$ चा \emph{उपसंच} आहे असे म्हणतो, आणि $A \subseteq B$ असे लिहितो.
$A$ हा $B$ चा उपसंच नसेल, तर $A \not\subseteq B$ असे लिहितो.
$A \subseteq B$ पण $A \neq B$ असेल, तर $A \subsetneq B$ असे लिहितो
आणि $A$ हा $B$ चा \emph{उचित उपसंच} आहे असे म्हणतो.
\end{defn}

\begin{ex}
प्रत्येक संच स्वतःचाच उपसंच असतो, आणि $\emptyset$ हा प्रत्येक संचाचा उपसंच असतो.
सम नैसर्गिक संख्यांचा संच हा नैसर्गिक संख्यांच्या संचाचा उपसंच आहे.
तसेच, $\{ a, b \} \subseteq \{ a, b, c \}$.
पण $\{ a, b, e
\}$ हा $\{ a, b, c \}$ चा उपसंच नाही.
\end{ex}

\begin{ex}
$2$ ही संख्या पूर्णांकांच्या संचाचा घटक आहे, तर सम संख्यांचा संच
हा पूर्णांकांच्या संचाचा उपसंच आहे. मात्र, एखादा संच दुसऱ्या संचाचा
घटक आणि उपसंच \emph{दोन्हीही} असू शकतो. उदाहरणार्थ,
$\{0\} \in \{0, \{0\}\}$ आणि $\{0\} \subseteq \{0,
\{0\}\}$ देखील.
\end{ex}

विस्तारात्मकतेमुळे संचांच्या समानतेचा निकष मिळतो: $A = B$ तेव्हा आणि केवळ तेव्हाच,
जेव्हा $A$ मधील प्रत्येक घटक हा $B$ चाही घटक असतो, आणि उलटही.
``उपसंच'' या व्याख्येत $A \subseteq B$ याचा अर्थ या निकषाचा नेमका पहिला
भाग आहे: $A$ मधील प्रत्येक घटक हा $B$ चा घटक असतो.
अर्थात, $A$ आणि $B$ यांची अदलाबदल केली तरी ही व्याख्या लागू होते:
$B \subseteq A$ तेव्हा आणि केवळ तेव्हाच, जेव्हा $B$ मधील प्रत्येक घटक
हा $A$ चा घटक असतो. हाच विस्तारात्मकतेतील ``आणि उलटही'' हा भाग होय.
दुसऱ्या शब्दांत, संच एकमेकांचे उपसंच असतील तेव्हा आणि केवळ तेव्हाच ते समान
असतात, हे विस्तारात्मकतेवरून निष्पन्न होते.

\begin{prop}
$A = B$ तेव्हा आणि केवळ तेव्हाच, जेव्हा $A \subseteq B$ आणि $B \subseteq A$ दोन्ही असतात.
\end{prop}

आणखी काही उपयुक्त चिन्हपद्धतींची ओळख करून घेण्याची ही योग्य वेळ आहे.
$A$ हा $B$ चा उपसंच केव्हा असतो याची व्याख्या करताना आपण
``$A$ मधील प्रत्येक घटक हा \dots'' असे म्हटले, आणि त्या
``$\dots$'' च्या जागी ``$B$ चा घटक असतो'' असे घातले.
परंतु असा \emph{आकार} असणाऱ्या अभिव्यक्ती इतक्या नेहमी वापरल्या जातात की,
त्यांच्यासाठी आकारिक चिन्हपद्धती उपयोगी पडेल.

\begin{defn}\label{sfr:set:sub:forallxina}
$(\forall x \in A)\phi$ हे $\forall x(x \in A \lif
\phi)$ याचे संक्षिप्त रूप आहे. त्याचप्रमाणे, $(\exists x \in A)\phi$ हे
$\exists x(x
\in A \land \phi)$ याचे संक्षिप्त रूप आहे.
\end{defn}

या चिन्हपद्धतीचा उपयोग करून, $A \subseteq B$ तेव्हा आणि केवळ तेव्हाच, जेव्हा
$(\forall
x \in A)x \in B$, असे म्हणता येते.

आता आपण एका विशिष्ट प्रकारच्या संचाचा विचार करू: दिलेल्या संचाच्या सर्व
उपसंचांचा संच.

\begin{defn}[घातसंच]
$A$ या संचाच्या सर्व उपसंचांचा मिळून जो संच बनतो, त्याला
$A$ चा \emph{घातसंच} म्हणतात आणि तो $\Pow{A}$ असा लिहितात.
  \[
    \Pow{A} = \Setabs{B}{B \subseteq A}
  \]
\end{defn}

\begin{ex}
$\{ a, b, c \}$ चे सर्व शक्य उपसंच कोणते? ते असे आहेत:
$\emptyset$, $\{a \}$, $\{b\}$, $\{c\}$, $\{a, b\}$, $\{a, c\}$, $\{b,
c\}$, $\{a, b, c\}$. या सर्व उपसंचांचा संच म्हणजे
$\Pow{\{a,b,c\}}$:
\[
\Pow{\{ a, b, c \}} = \{\emptyset, \{a \}, \{b\}, \{c\}, \{a, b\},
\{b, c\}, \{a, c\}, \{a, b, c\}\}
\]
\end{ex}

\begin{prob}
$\{a, b, c, d\}$ चे सर्व उपसंच लिहा.
\end{prob}

\begin{prob}
$A$ मध्ये $n$ घटक असतील, तर $\Pow{A}$ मध्ये $2^n$
घटक असतात हे दाखवा.
\end{prob}











\section{काही महत्त्वाचे संच}\label{sfr:set:imp:sec}

\begin{ex}
आपण प्रामुख्याने ज्यातील घटक या गणिती वस्तू आहेत अशा संचांचा
विचार करणार आहोत. असे चार संच इतके महत्त्वाचे आहेत की त्यांना
विशिष्ट नावे दिलेली आहेत:
\begin{multline*}
    \Nat = \{0, 1, 2, 3, \ldots\} \\
    \shoveright{\text{नैसर्गिक संख्यांचा संच}}\\
    \shoveleft{\Int = \{\ldots, -2, -1, 0, 1, 2, \ldots\}} \\
    \shoveright{\text{पूर्णांकांचा संच}}\\
    \shoveleft{\Rat = \Setabs{\nicefrac{m}{n}}{m, n \in \Int\text{ आणि }n \neq 0}}\\
    \shoveright{\text{परिमेय संख्यांचा संच}}\\
    \shoveleft{\Real = (-\infty, \infty)}\\
    \text{वास्तव संख्यांचा संच (सांतत्य)}
\end{multline*}
हे सर्व \emph{अनंत} संच आहेत; म्हणजेच, प्रत्येकात अनंत घटक आहेत.

या संचांतून पुढे जाताना आपण आपल्या संग्रहात \emph{आणखी} संख्या समाविष्ट
करत असतो. खरे तर, $\Nat \subseteq \Int
\subseteq \Rat \subseteq \Real$ हे स्पष्ट असावे: कारण प्रत्येक नैसर्गिक संख्या
पूर्णांक असते; प्रत्येक पूर्णांक परिमेय असतो; आणि प्रत्येक परिमेय संख्या वास्तव असते.
त्याचप्रमाणे, $\Nat \subsetneq \Int \subsetneq
\Rat$ हेही स्पष्ट असावे, कारण $-1$ हा पूर्णांक आहे पण नैसर्गिक संख्या नाही,
आणि $\nicefrac{1}{2}$ ही परिमेय संख्या आहे पण पूर्णांक नाही.
$\Rat \subsetneq \Real$, म्हणजे परिमेय नसलेल्या काही वास्तव संख्या आहेत,
ही गोष्ट तितकी सहज स्पष्ट होत नाही.

कधीकधी आपण धन पूर्णांकांचा संच $\PosInt = \{1,
2, 3, \dots\}$ आणि फक्त पहिल्या दोन नैसर्गिक संख्या असणारा संच
$\Bin = \{0, 1\}$ यांचाही उपयोग करू.
\end{ex}


\begin{ex}[चिन्हमाला]
आणखी एक लक्षवेधी उदाहरण म्हणजे $A$ या वर्णमालेवरील \emph{सांत चिन्हमालांचा}
संच $A^{*}$: $A$ मधील घटकांचा कोणताही सांत अनुक्रम हा $A$ वरील चिन्हमाला
असतो. प्रत्येक $A$ वर्णमालेसाठी, $A$ वरील चिन्हमालांमध्ये आपण
\emph{रिक्त चिन्हमाला $\Lambda$} हिचाही समावेश करतो. उदाहरणार्थ,
\begin{multline*}
\Bin^*
=\{\Lambda,0,1,00,01,10,11,\\
000,001,010,011,100,101,110,111,0000,\ldots\}.
\end{multline*}
$x=x_{1}\ldots x_{n}\in A^{*}$ ही $A$ मधील $n$ ``अक्षरांनी'' बनलेली
चिन्हमाला असेल, तर तिची \emph{लांबी} $n$ आहे असे म्हणतो आणि $\len{x}=n$ असे लिहितो.
\end{ex}


\begin{ex}[अनंत अनुक्रम]
कोणत्याही $A$ संचासाठी आपण $A$ मधील घटक यांच्या अनंत अनुक्रमांचा
संच $A^\omega$ याचाही विचार करू शकतो. अनंत अनुक्रम
$a_1a_2a_3a_4\dots$ म्हणजे वस्तूंची एका दिशेने अनंत लांब जाणारी यादी;
यादीतील प्रत्येक वस्तू $A$ चा घटक असते.
\end{ex}











\section{संयोग आणि छेद}\label{sfr:set:uni:sec}

\begin{explain}
\ref{sfr:set:bas:sec} मध्ये आपण गुणधर्मानुसार संचांची व्याख्या करणे,
म्हणजे $\Setabs{x}{\phi(x)}$ या आकाराची व्याख्या करणे, पाहिले.
येथे आपण एखादा $\phi$ गुणधर्म वापरतो; या गुणधर्मात आधीच व्याख्या केलेल्या
संचांचा उल्लेख असू शकतो. उदाहरणार्थ, $A$ आणि $B$ हे संच असतील,
तर $\Setabs{x}{x \in A \lor x \in B}$ या संचात $A$ किंवा $B$ मधील
घटक असणाऱ्या सर्व वस्तू येतात. म्हणजेच, $A$ आणि $B$ मधील
घटक एकत्र करणारा हा संच आहे. \ref{sfr:set:uni:fig:union} मध्ये
दाखवल्याप्रमाणे त्याचे चित्र काढता येते; ठळक केलेला भाग $A$ आणि $B$
या दोन संचांतील घटक एकत्र दर्शवतो.

\begin{figure}
  \olasset{assets/diagrams/union.tikz}
  \caption{दोन संचांचा संयोग $A \cup B$ म्हणजे $A$ मधील आणि $B$ मधील
    घटक एकत्र घेऊन बनणारा संच.}
  \label{sfr:set:uni:fig:union}
\end{figure}

संच एकत्र करणारी ही क्रिया अतिशय उपयुक्त आहे आणि वारंवार वापरली जाते.
म्हणून तिला आपण आकारिक नाव आणि चिन्ह देतो.
\end{explain}

\begin{defn}[संयोग]
$A$ आणि $B$ या दोन संचांचा \emph{संयोग} $A \cup B$ असा लिहितात.
तो $A$, $B$ किंवा दोन्ही संचांतील घटक असणाऱ्या सर्व वस्तूंचा संच आहे.
\[
A \cup B = \Setabs{x}{x \in A \lor x \in B}
\]
\end{defn}

\begin{ex}
तेच घटक किती वेळा आले आहेत याला महत्त्व नसल्याने, दोन संचांत
सामाईक घटक असेल, तर त्यांच्या संयोगात तो घटक एकदाच येतो.
उदाहरणार्थ, $\{ a, b, c\} \cup \{ a, 0, 1\} = \{a, b, c, 0, 1\}$.

संच आणि त्याचा एखादा उपसंच यांचा संयोग म्हणजे मोठा संचच:
$\{a,
b, c \} \cup \{a \} = \{a, b, c\}$.

संचाचा रिक्त संचाशी संयोग हा मूळ संचाशी समान असतो:
$\{a,
b, c \} \cup \emptyset = \{a, b, c \}$.
\end{ex}

\begin{prob}
$A \subseteq B$ असेल, तर $A \cup B = B$ हे सिद्ध करा.
\end{prob}

\begin{explain}
संयोगाच्या ``द्वैत'' क्रियेचाही विचार करता येतो. या क्रियेत असे सर्व
घटक घेतले जातात की जे $A$ मधील घटक आणि
$B$ मधीलही घटक आहेत, आणि त्यांचा संच बनवला जातो.
या क्रियेला \emph{छेद} म्हणतात. तिचे चित्र \ref{sfr:set:uni:fig:intersection} प्रमाणे काढता येते.
\begin{figure}
  \olasset{assets/diagrams/intersection.tikz}
  \caption{दोन संचांचा छेद $A \cap B$ म्हणजे त्यांतील सामाईक
    घटक यांचा संच.}
  \label{sfr:set:uni:fig:intersection}
\end{figure}
\end{explain}

\begin{defn}[छेद]
$A$ आणि $B$ या दोन संचांचा \emph{छेद} $A \cap B$ असा लिहितात.
तो $A$ आणि $B$ या दोन्ही संचांतील घटक असणाऱ्या सर्व वस्तूंचा संच आहे.
\[
A \cap B = \Setabs{x}{x \in A \land x \in B}
\]
दोन संचांचा छेद रिक्त असेल, तर त्यांना \emph{विभक्त} म्हणतात.
याचा अर्थ, त्यांच्यात सामाईक घटक नसतात.
\end{defn}

\begin{ex}
दोन संचांत सामाईक घटक नसतील, तर त्यांचा छेद रिक्त असतो:
$\{ a, b, c\} \cap \{ 0, 1\} = \emptyset$.

दोन संचांत सामाईक घटक असतील, तर त्यांचा छेद म्हणजे त्या सर्वांचा संच:
$\{a, b, c \} \cap \{a, b, d \} = \{a, b\}$.

संच आणि त्याचा एखादा उपसंच यांचा छेद म्हणजे लहान संचच:
$\{a, b, c\} \cap \{a, b\} = \{a, b\}$.

कोणत्याही संचाचा रिक्त संचाशी छेद रिक्त असतो:
$\{a, b, c \}
\cap \emptyset = \emptyset$.
\end{ex}

\begin{prob}
$A \subseteq B$ असेल, तर $A \cap B = A$ हे काटेकोरपणे सिद्ध करा.
\end{prob}

\begin{explain}
दोनपेक्षा अधिक संचांचा संयोग किंवा छेदही बनवता येतो.
हे सर्वसाधारणपणे हाताळण्याची एक सुटसुटीत पद्धत पुढीलप्रमाणे आहे:
ज्या सर्व संचांचा संयोग (किंवा छेद) करायचा आहे, ते सर्व एका संचात एकत्र केले आहेत असे समजा.
मग या संचाच्या किमान एका घटक मध्ये असणाऱ्या सर्व वस्तूंचा संच
म्हणजे मूळ सर्व संचांचा संयोग अशी व्याख्या करता येते.
तसेच, या संचाच्या प्रत्येक घटक मध्ये असणाऱ्या सर्व वस्तूंचा संच
म्हणजे त्यांचा छेद अशी व्याख्या करता येते.
\end{explain}

\begin{defn}
$A$ हा संचांचा संच असेल, तर $\bigcup A$ म्हणजे $A$ मधील
घटक यांत असणाऱ्या घटक यांचा संच:
\begin{align*}
\bigcup A & = \Setabs{x}{x \text{ ही पुढील संचाच्या एका घटकात असते: } A},
\text{ म्हणजे,}\\
& = \Setabs{x}{\text{असा } B \in A
  \text{ असतो की } x \in B}
\end{align*}
\end{defn}

\begin{defn}
$A$ हा संचांचा संच असेल, तर $\bigcap A$ म्हणजे $A$ मधील सर्व घटकांत
सामाईक असणाऱ्या वस्तूंचा संच:
\begin{align*}
\bigcap A & = \Setabs{x}{x \text{ ही पुढील संचाच्या प्रत्येक घटकात असते: } A},
\text{ म्हणजे,}\\
 & = \Setabs{x}{\text{प्रत्येक } B \in A, x \in B}
\end{align*}
\end{defn}

\begin{ex}
$A = \{ \{ a, b \}, \{ a, d, e \}, \{ a, d \} \}$ असे समजा.
मग $\bigcup A = \{ a, b, d, e \}$ आणि $\bigcap A = \{ a \}$.
\end{ex}
\begin{prob}
	$A$ हा संच असेल आणि $A \in B$ असेल, तर $A \subseteq \bigcup B$ हे दाखवा.
\end{prob}

$A_1$, $A_2$, \dots या संचांच्या अनुक्रमासाठीही हेच करता येईल:
\begin{align*}
\bigcup_i A_i & = \Setabs{x}{x \text{ ही पुढीलपैकी एखाद्या संचात असते: } A_i}\\
\bigcap_i A_i & = \Setabs{x}{x \text{ ही पुढील प्रत्येक संचात असते: } A_i}.
\end{align*}

संचांसाठी \emph{निर्देशांक संच} दिलेला असेल, म्हणजे एखादा $I$ संच असा असेल की,
प्रत्येक $i \in I$ साठी आपण $A_i$ चा विचार करत असू, तर पुढील संक्षिप्त
रूपेही वापरता येतात:
\begin{align*}
	\bigcup_{i \in I} A_i & = \bigcup \Setabs{A_i }{i \in I}\\
	\bigcap_{i \in I} A_i & = \bigcap\Setabs{A_i}{i \in I}
\end{align*}

शेवटी, $A$ मध्ये असणारे पण $B$ मध्ये नसणारे सर्व घटक यांच्या
संचाचा विचार करायचा असू शकतो. त्याचे चित्र \ref{sfr:set:uni:difference} प्रमाणे काढता येते.

\begin{figure}
  \olasset{assets/diagrams/difference.tikz}
  \caption{दोन संचांचा फरक $A \setminus B$ म्हणजे $A$ मधील असे
    घटक की जे $B$ मधील घटक नाहीत, यांचा संच.}
  \label{sfr:set:uni:difference}
\end{figure}

\begin{defn}[फरक]
\emph{संचांचा फरक} $A \setminus B$ म्हणजे $A$ मधील असे सर्व
घटक की जे $B$ मधील घटक नाहीत, यांचा संच. म्हणजेच,
\[
A\setminus B = \Setabs{x}{x\in A \text{ आणि } x \notin B}.
\]
\end{defn}

\begin{prob}
	$A \subsetneq B$ असेल, तर $B \setminus A \neq \emptyset$ हे सिद्ध करा.
\end{prob}











\section{जोड्या, क्रमित घटकसमूह आणि कार्तीय गुणाकार}\label{sfr:set:pai:sec}

\begin{explain}
संचांतील घटकांना क्रम नसतो, हे विस्तारात्मकतेवरून निष्पन्न होते.
म्हणून क्रम दर्शवायचा असेल, तर आपण \emph{क्रमित जोड्या} $\tuple{x, y}$ वापरतो.
अक्रमित जोडी $\{x, y\}$ मध्ये क्रमाला महत्त्व नसते:
$\{x, y\} = \{y, x\}$. क्रमित जोडीमध्ये मात्र क्रम महत्त्वाचा असतो:
$x
\neq y$ असेल, तर $\tuple{x, y} \neq \tuple{y, x}$.

संच सिद्धांतात क्रमित जोड्यांचा विचार कसा करावा? दोन क्रमित जोड्यांचे
पहिले घटक समान असतील आणि दुसरे घटकही समान असतील तेव्हा आणि केवळ तेव्हाच
त्या जोड्या समान असतात, ही कल्पना जपणे आवश्यक आहे. म्हणजेच:
\[
  \tuple{a, b}= \tuple{c, d}\text{ तेव्हा आणि केवळ तेव्हाच, जेव्हा }a = c \text{ आणि }b=d.
\]
वीनर--कुराटोव्स्की यांच्या व्याख्येने संच सिद्धांतात क्रमित जोड्यांची व्याख्या करता येते.
\end{explain}

\begin{defn}[क्रमित जोडी]\label{sfr:set:pai:wienerkuratowski}
	$\tuple{a, b} = \{\{a\}, \{a, b\}\}$.
\end{defn}

\begin{prob}
	\ref{sfr:set:pai:wienerkuratowski} वापरून, $\tuple{a,
	b}= \tuple{c, d}$ तेव्हा आणि केवळ तेव्हाच, जेव्हा $a = c$ आणि $b=d$
दोन्ही असतात, हे सिद्ध करा.
\end{prob}

\begin{explain}
क्रमित जोडीची व्याख्या निश्चित केल्यावर तिच्या साहाय्याने आणखी संचांची
व्याख्या करता येते. उदाहरणार्थ, कधीकधी दोनपेक्षा अधिक वस्तूंचे क्रमित
अनुक्रम हवे असतात: \emph{त्रिके} $\tuple{x, y, z}$,
\emph{चतुष्के} $\tuple{x, y, z, u}$ इत्यादी.
त्रिक म्हणजे ज्याचा पहिला घटक स्वतःच क्रमित जोडी आहे अशी विशिष्ट क्रमित
जोडी असे मानता येते: $\tuple{x, y, z}$ म्हणजे $\tuple{\tuple{x, y},z}$.
चतुष्कांसाठीही हेच लागू होते: $\tuple{x,y,z,u}$ म्हणजे
$\tuple{\tuple{\tuple{x,y},z},u}$, आणि असेच पुढे.
सर्वसाधारणपणे आपण \emph{क्रमित $n$-घटकसमूह} $\tuple{x_1, \dots, x_n}$ यांचा उल्लेख करतो.

क्रमित जोड्यांचे, किंवा इतर क्रमित $n$-घटकसमूहांचे, काही संच उपयुक्त ठरतील.
\end{explain}

\begin{defn}[कार्तीय गुणाकार]
$A$ आणि $B$ हे संच दिलेले असतील, तर त्यांचा \emph{कार्तीय गुणाकार}
$A \times B$ याची व्याख्या अशी करतात:
\[
  A \times B = \Setabs{\tuple{x, y}}{x \in A \text{ आणि } y \in B}.
\]
\end{defn}

\begin{ex}
$A = \{0, 1\}$ आणि $B = \{1, a, b\}$ असतील, तर त्यांचा गुणाकार असा आहे:
\[
A \times B = \{ \tuple{0, 1}, \tuple{0, a}, \tuple{0, b},
    \tuple{1, 1}, \tuple{1, a}, \tuple{1, b} \}.
\]
\end{ex}

\begin{ex}
$A$ हा संच असेल, तर $A$ चा स्वतःशी गुणाकार $A \times A$ हा $A^2$
असाही लिहितात. तो $x, y \in A$ असणाऱ्या \emph{सर्व} $\tuple{x, y}$ जोड्यांचा
संच आहे. सर्व $\tuple{x, y, z}$ त्रिकांचा संच $A^3$ असतो, आणि असेच पुढे.
पुढील पुनरावर्ती व्याख्या देता येते:
\begin{align*}
  A^1 & = A\\
  A^{k+1} & = A^k \times A
\end{align*}
\end{ex}

\begin{prob}
$\{1, 2, 3\}^3$ मधील सर्व घटक लिहा.
\end{prob}

\begin{prop}\label{sfr:set:pai:cardnmprod}
$A$ मध्ये $n$ घटक आणि $B$ मध्ये $m$ घटक असतील,
तर $A
\times B$ मध्ये $n\cdot m$ घटक असतात.
\end{prop}

\begin{proof}
$A$ मधील प्रत्येक घटक $x$ साठी, $\tuple{x, y} \in A \times B$
या आकाराचे $m$ घटक असतात.
$B_x = \Setabs{\tuple{x, y}}{y
  \in B}$ असे मानू. $x_1 \neq x_2$ असेल तेव्हा $\tuple{x_1, y} \neq
\tuple{x_2, y}$ असल्यामुळे $B_{x_1} \cap B_{x_2} = \emptyset$.
पण $A = \{x_1,
\dots, x_n\}$ असेल, तर $A \times B = B_{x_1} \cup \dots \cup B_{x_n}$;
म्हणून त्यात $n\cdot m$ घटक असतात.

हे चित्ररूपाने पाहण्यासाठी $A \times B$ मधील घटक जाळीच्या स्वरूपात मांडा:
\[
\begin{array}{rcccc}
  B_{x_1} = & \{\tuple{x_1, y_1} & \tuple{x_1, y_2} & \dots & \tuple{x_1, y_m}\}\\
  B_{x_2} = & \{\tuple{x_2, y_1} & \tuple{x_2, y_2} & \dots & \tuple{x_2, y_m}\}\\
  \vdots & & \vdots\\
  B_{x_n} = & \{\tuple{x_n, y_1} & \tuple{x_n, y_2} & \dots & \tuple{x_n, y_m}\}
\end{array}
\]
सर्व $x_i$ वेगवेगळे आहेत आणि सर्व $y_j$ वेगवेगळे आहेत. त्यामुळे या जाळीतील
कोणत्याही दोन जोड्या समान नाहीत, आणि त्यांची संख्या $n\cdot m$ आहे.
\end{proof}

\begin{prob}
$k$ वरील गणिती विगमनाने असे दाखवा की, प्रत्येक $k \ge 1$ साठी,
$A$ मध्ये $n$ घटक असतील, तर $A^k$ मध्ये $n^k$ घटक असतात.
\end{prob}

\begin{ex}
$A$ हा संच असेल, तर $A$ वरील \emph{शब्द} म्हणजे $A$ मधील
घटक यांचा कोणताही अनुक्रम.
असा अनुक्रम म्हणजे $A$ मधील घटक यांचा $n$-घटकसमूह असे मानता येते.
उदाहरणार्थ, $A = \{a, b, c\}$ असेल, तर ``$bac$'' या अनुक्रमाचा
$\tuple{b, a, c}$ हे त्रिक म्हणून विचार करता येतो.
शब्दांना, म्हणजे चिन्हांच्या अनुक्रमांना, संगणकशास्त्रात अत्यंत महत्त्व आहे.
संकेतानुसार, $A$ मधील घटक यांना आपण $1$ लांबीचे अनुक्रम मानतो,
आणि $\emptyset$ याला $0$ लांबीचा अनुक्रम मानतो.
मग $A$ वरील \emph{सर्व} शब्दांचा संच असा आहे:
\[
A^* = \{\emptyset\} \cup A \cup A^2 \cup A^3 \cup \dots
\]
\end{ex}












\section{रसेलची विरोधापत्ती}\label{sfr:set:rus:sec}

विस्तारात्मकतेमुळे $\Setabs{x}{\phi(x)}$ ही चिन्हपद्धती $\phi(x)$ असणाऱ्या
सर्व $x$ च्या \emph{त्या विशिष्ट} संचासाठी वापरणे योग्य ठरते.
मात्र, विस्तारात्मकतेमुळे \emph{प्रत्यक्षात} फक्त पुढील विचार योग्य ठरतो.
ज्याचे सदस्य सर्व आणि केवळ $\phi$ असणाऱ्या वस्तू आहेत असा संच \emph{असेल},
\emph{तर} असा एकच संच असतो. दुसऱ्या शब्दांत: एखादा $\phi$ निश्चित केल्यावर,
$\Setabs{x}{\phi(x)}$ हा संच \emph{अस्तित्वात असल्यास} तो एकमेव असतो.

पण ही अट महत्त्वाची आहे!{} प्रत्येक गुणधर्मापासून \emph{गुणधर्माधारित संचरचना}
करता येतेच असे नाही. म्हणजेच, काही गुणधर्म संचांची व्याख्या करत \emph{नाहीत}.
सर्वच गुणधर्मांनी संचांची व्याख्या होत असती, तर आपल्याला थेट व्याघातांना
सामोरे जावे लागले असते. याचे सर्वांत प्रसिद्ध उदाहरण म्हणजे रसेलची विरोधापत्ती.

संच हे इतर संचांचे घटक असू शकतात; उदाहरणार्थ, $A$ या संचाचा घातसंच
हा संचांनी बनलेला असतो. म्हणून एखादा संच दुसऱ्या संचाचा घटक आहे का,
असा प्रश्न विचारणे किंवा त्याचा शोध घेणे अर्थपूर्ण आहे. संच स्वतःचाच सदस्य
असू शकतो का? संचाच्या कल्पनेत असे काही दिसत नाही की ज्यामुळे ही शक्यता
नाकारली जाईल. उदाहरणार्थ, \emph{सर्व} संच मिळून वस्तूंचा समूह बनत असेल,
तर त्या सर्वांना एका संचात एकत्र करता येईल असे वाटू शकते; तो म्हणजे सर्व
संचांचा संच. तो स्वतः एक संच असल्यामुळे सर्व संचांच्या संचाचा घटक असेल.

स्वतःला घटक म्हणून न सामावणाऱ्या, म्हणजे \emph{स्वतःचे सदस्य नसणाऱ्या}
संचांचा गुणधर्म विचारात घेतल्यावर रसेलची विरोधापत्ती उद्भवते.
स्वतःला घटक म्हणून न सामावणाऱ्या सर्व संचांचा एक संच आहे असे आपण
मानले, तर काय होईल? असा
\[
R = \Setabs{x}{x \notin x}
\]
संच अस्तित्वात आहे का? तो अस्तित्वात नसतो हे आपण सिद्ध करू शकतो.

\begin{thm}[रसेलची विरोधापत्ती]\label{sfr:set:rus:thm:russells-paradox}
	$R = \Setabs{x}{x \notin x}$ असा कोणताही संच नाही.
\end{thm}

\begin{proof}
$R = \Setabs{x}{x \notin x}$ अस्तित्वात असेल, तर
$R \in R$ तेव्हा आणि केवळ तेव्हाच, जेव्हा $R \notin R$; हा व्याघात आहे.
\end{proof}


\begin{explain}
ही सिद्धता अधिक सावकाश पाहू. $R$ अस्तित्वात असेल, तर $R \in
R$ आहे की नाही हा प्रश्न अर्थपूर्ण आहे. खरोखर $R \in R$ आहे असे समजा.
आता, स्वतःचे घटक नसणाऱ्या सर्व संचांचा संच अशी $R$ ची व्याख्या केली होती.
त्यामुळे $R \in R$ असेल, तर $R$ मध्ये स्वतःच $R$ चा व्याख्यक गुणधर्म नसतो.
पण हा गुणधर्म असणारेच संच $R$ मध्ये असतात. म्हणून $R$ हा $R$ चा
घटक असू शकत नाही; म्हणजेच $R \notin R$.
परंतु $R$ हा $R$ चा घटक आहे आणि नाही, असे दोन्ही होऊ शकत नाही;
म्हणून व्याघात मिळतो.

$R \in R$ या गृहीतकामुळे व्याघात येत असल्याने $R \notin R$ मिळते.
पण यामुळेही व्याघात येतो!{} कारण $R
\notin R$ असेल, तर $R$ मध्ये स्वतःच $R$ चा व्याख्यक गुणधर्म असतो;
म्हणून स्वतःचे सदस्य नसणाऱ्या इतर सर्व संचांप्रमाणे $R$ हा $R$ चा
घटक असेल. आणि पुन्हा, $R$ चा घटक नसणे आणि असणे
या दोन्ही गोष्टी एकत्र शक्य नाहीत.
\end{explain}


\begin{digress}
रसेलची विरोधापत्ती टाळणारा, म्हणजे $R = \Setabs{x}{x \notin x}$
अस्तित्वात आहे हा \emph{विसंगत} दावा न करणारा संच सिद्धांत कसा मांडायचा?
त्यासाठी संच केव्हा अस्तित्वात असतात (आणि केव्हा नसतात) हे सांगणाऱ्या
अगदी नेमक्या अटी देणारी स्वयंसिद्धके मांडावी लागतील.

या प्रकरणात आराखडा दिलेला संच सिद्धांत असे करत नाही. तो
\emph{खरोखरच अनौपचारिक} आहे. त्यातून एवढेच सांगितले जाते की संच
विस्तारात्मकतेचे पालन करतात आणि काही संच दिलेले असतील, तर त्यांचा संयोग,
छेद इत्यादी बनवता येतात. संच सिद्धांत यापेक्षा अधिक काटेकोरपणे विकसित करणे
शक्य आहे.
\end{digress}




\chapter{संबंध}\label{sfr:rel::chap}








\section{संच म्हणून संबंध}\label{sfr:rel:set:sec}

\begin{explain}
\ref{sfr:set:imp:sec} मध्ये आपण काही महत्त्वाच्या संचांचा उल्लेख केला:
$\Nat$, $\Int$, $\Rat$, $\Real$. यांपैकी काही संचांतील घटक
यांच्यातील काही लक्षवेधी संबंध तुम्हाला आठवत असतील. उदाहरणार्थ, या प्रत्येक
संचावर एक सर्वमान्य \emph{क्रमसंबंध} असतो. तसेच, प्रत्येक वस्तूचा फक्त
स्वतःशी आणि इतर कशाशीही नसणारा \emph{एकरूप असण्याचा} संबंध असतो. पुढे
आपल्याला आणखी अनेक लक्षवेधी संबंध भेटतील, आणि शक्य संबंध तर त्याहूनही
अधिक आहेत. त्यांचा आढावा घेण्यापूर्वी मात्र, संबंधांकडे एका विशिष्ट प्रकारचे
संच म्हणून पाहता येते, हे आपण दाखवू.

यासाठी \ref{sfr:set:pai:sec} मधील दोन गोष्टी आठवा. प्रथम,
\emph{क्रमित जोडी} ही संकल्पना: $a$ आणि $b$ दिलेले असतील, तर
$\tuple{a, b}$ बनवता येते. येथे घटकांचा क्रम महत्त्वाचा \emph{असतो}.
म्हणून $a \neq b$ असेल, तर $\tuple{a, b} \neq \tuple{b, a}$.
(याची अक्रमित जोड्यांशी, म्हणजे $2$-घटक संचांशी, तुलना करा: तेथे
$\{a, b\}=\{b, a\}$.) दुसरे म्हणजे, \emph{कार्तीय गुणाकार} ही संकल्पना
आठवा: $A$ आणि $B$ हे संच असतील, तर $A \times B$ बनवता येतो. तो
$x \in A$ आणि $y \in B$ असणाऱ्या सर्व $\tuple{x, y}$ जोड्यांचा संच
आहे. विशेषतः, $A^{2}= A \times A$ हा $A$ मधून घेतलेल्या सर्व क्रमित
जोड्यांचा संच आहे.

आता एका संचावरील विशिष्ट संबंध पाहू: नैसर्गिक संख्यांच्या $\Nat$ या
संचावरील $<$-संबंध. $n<m$ असणाऱ्या संख्यांच्या सर्व $\tuple{n, m}$
जोड्यांचा संच विचारात घ्या, म्हणजे,
\[
  R=\Setabs{\tuple{n, m}}{n, m \in \Nat \text{ आणि } n<m}.
\]
$n$ ही संख्या $m$ पेक्षा लहान असणे आणि $\tuple{n, m}$ ही जोडी
$R$ ची सदस्य असणे यांचा निकटचा संबंध आहे, तो असा:
\[
      n<m\text{ तेव्हा आणि केवळ तेव्हाच }\tuple{n, m} \in R.
\]
खरे तर, कोणतीही माहिती न गमावता $R$ हा संच \emph{म्हणजेच} $\Nat$
वरील $<$-संबंध आहे, असे आपण मानू शकतो.

त्याचप्रमाणे, संख्यांमधील कोणत्याही संबंधासाठी $\Nat^{2}$ चा एक उपसंच
तयार करता येतो. उलट, संख्यांच्या जोड्यांचा कोणताही $S \subseteq \Nat^{2}$
संच दिला असेल, तर त्याला अनुरूप संख्यांमधील एक संबंध असतो: $n$ चा $m$
शी तो संबंध तेव्हा आणि केवळ तेव्हाच असतो, जेव्हा $\tuple{n, m} \in S$.
यावरून पुढील व्याख्या समर्थित होते:
\end{explain}

\begin{defn}[द्विपदी संबंध]
$A$ या संचावरील \emph{द्विपदी संबंध} म्हणजे $A^{2}$ चा एक उपसंच.
$R \subseteq A^{2}$ हा $A$ वरील द्विपदी संबंध असेल आणि $x, y \in A$
असतील, तर $\tuple{x, y} \in R$ यासाठी आपण कधी कधी $Rxy$ (किंवा $xRy$)
असे लिहितो.
\end{defn}

\begin{ex}
  \label{sfr:rel:set:relations}
नैसर्गिक संख्यांच्या जोड्यांचा $\Nat^{2}$ हा संच अशा द्विमितीय सारणीत
मांडता येतो:
\[
  \begin{array}{ccccc}
  \mathbf{\tuple{ 0,0 }} & \tuple{ 0,1 } &
    \tuple{ 0,2 } & \tuple{ 0,3 } & \ldots\\
  \tuple{ 1,0 } & \mathbf{\tuple{ 1,1 }} &
    \tuple{ 1,2 } & \tuple{ 1,3 } & \ldots\\
  \tuple{ 2,0 } & \tuple{ 2,1 } &
    \mathbf{\tuple{ 2,2 }} & \tuple{ 2,3 } & \ldots\\
  \tuple{ 3,0 } & \tuple{ 3,1 } & \tuple{ 3,2 } &
    \mathbf{\tuple{ 3,3 }} & \ldots\\
  \vdots & \vdots & \vdots & \vdots & \mathbf{\ddots}
  \end{array}
\]
येथे कर्णावरील नोंदी ठळक केल्या आहेत, कारण त्या जोड्यांचा बनलेला
$\Nat^2$ चा उपसंच, म्हणजे,
\[
  \{\tuple{0,0 }, \tuple{ 1,1 }, \tuple{ 2,2 }, \dots\},
  \]
हा \emph{$\Nat$ वरील एकरूपता संबंध} आहे. (एकरूपता संबंध वारंवार वापरला
जातो, म्हणून कोणत्याही $A$ संचासाठी $\Id{A}=\Setabs{\tuple{ x,x }}{x \in
A}$ अशी व्याख्या करू.) कर्णाच्या वर असलेल्या सर्व जोड्यांचा उपसंच, म्हणजे,
\[
  L = \{\tuple{ 0,1 },\tuple{ 0,2 },\ldots,\tuple{ 1,2 },
  \tuple{ 1,3 }, \dots, \tuple{ 2,3 },\tuple{ 2,4 },\ldots\},
\]
हा \emph{लहान असण्याचा} संबंध आहे, म्हणजे $Lnm$ तेव्हा आणि केवळ तेव्हाच,
जेव्हा $n<m$. कर्णाच्या खाली असलेल्या जोड्यांचा उपसंच, म्हणजे,
\[
  G=\{ \tuple{ 1,0 },\tuple{ 2,0 },\tuple{
    2,1 }, \tuple{ 3,0 },\tuple{ 3,1 },\tuple{ 3,2 }, \dots\},
\]
हा \emph{मोठे असण्याचा} संबंध आहे, म्हणजे $Gnm$ तेव्हा आणि केवळ तेव्हाच,
जेव्हा $n>m$. $L$ आणि $I$ यांचा संयोग $K=L\cup I$ असा लिहू; तो
\emph{लहान किंवा समान असण्याचा} संबंध आहे: $Knm$ तेव्हा आणि केवळ तेव्हाच,
जेव्हा $n \le m$. त्याचप्रमाणे, $H=G \cup I$ हा \emph{मोठे किंवा समान
असण्याचा संबंध} आहे. $L$, $G$, $K$ आणि $H$ हे \emph{क्रम} नावाच्या
विशिष्ट प्रकारचे संबंध आहेत. $L$ आणि $G$ यांचा गुणधर्म असा: कोणत्याही संख्येचा स्वतःशी $L$ किंवा $G$
संबंध नसतो (म्हणजे, सर्व $n$ साठी $Lnn$ आणि $Gnn$ दोन्ही असत्य असतात).
हा गुणधर्म असणाऱ्या संबंधांना \emph{अपरावर्ती} म्हणतात; आणि ते क्रमही
असतील, तर त्यांना \emph{काटेकोर क्रम} म्हणतात.
\end{ex}

\begin{explain}
क्रम आणि एकरूपता हे महत्त्वाचे व स्वाभाविक संबंध असले, तरी आपल्या
व्याख्येनुसार $A^{2}$ चा \emph{कोणताही} उपसंच हा $A$ वरील संबंध असतो,
तो कितीही कृत्रिम किंवा ओढूनताणून बनवलेला वाटला तरीही. विशेषतः,
$\emptyset$ हा कोणत्याही संचावरील संबंध असतो (हा \emph{रिक्त संबंध};
कोणत्याही घटकजोडीमध्ये तो नसतो), आणि $A^{2}$ हाही $A$ वरील संबंध असतो
(प्रत्येक जोडीमध्ये असणारा); त्याला \emph{सार्वत्रिक संबंध} म्हणतात.
तसेच, $E=\Setabs{\tuple{n, m}}{n>5 \text{ किंवा } m \times n \ge 34}$
यासारखी गोष्टही संबंध ठरते.
\end{explain}

\begin{prob}
  $\Pow{\{a, b, c\}}$ या संचावरील $\subseteq$ संबंधाचे घटक लिहा.
\end{prob}











\section{तात्त्विक चिंतन}\label{sfr:rel:ref:sec}

\ref{sfr:rel:set:sec} मध्ये आपण संबंधांची व्याख्या काही विशिष्ट संच म्हणून
केली. येथे थांबून एक छोटासा तात्त्विक प्रश्न विचारू: अशी व्याख्या नेमके
काय \emph{करते}? काही सत्तामीमांसक एकरूपतेची तथ्ये आपण \emph{शोधून काढली},
असा दावा आपल्याला करायचा आहे का, याबद्दल खूपच शंका आहे. उदाहरणार्थ,
$\Nat$ वरील क्रमसंबंध हा \ref{sfr:rel:set:sec} मध्ये व्याख्यिलेला
$R=  \Setabs{\tuple{n,m}}{n, m \in \Nat\text{ आणि }
n < m}$ संचच \emph{निघाला}, असे म्हणायचे का? अशी शंका घेण्याची तीन
कारणे पुढीलप्रमाणे आहेत.

पहिले: \ref{sfr:set:pai:wienerkuratowski} मध्ये आपण
$\tuple{a, b} = \{\{a\}, \{a, b\}\}$ अशी व्याख्या केली. त्याऐवजी
$\lVert a,b\rVert = \{\{b\}, \{a, b\}\} = \tuple{b,a}$ ही व्याख्या
विचारात घ्या. $a \neq b$ असेल, तर $\tuple{a, b} \neq \lVert a,b\rVert$.
पण क्रमित जोडीची व्याख्या $\tuple{a,b}$ ऐवजी $\lVert a,b\rVert$ अशीही
तितक्याच योग्य रीतीने करता आली असती. दोन्ही व्याख्या तितक्याच उपयोगी
ठरल्या असत्या. त्यामुळे नैसर्गिक संख्यांवरील क्रमसंबंध ``असण्यासाठी''
आता दोन तितकेच योग्य उमेदवार आहेत:
\begin{align*}
		R &= \Setabs{\tuple{n,m}}{n, m \in \Nat \text{ आणि }n < m}\\
		S &= \Setabs{\lVert n,m\rVert}{n, m \in \Nat \text{ आणि }n < m}.
\end{align*}
विस्तारात्मकतेनुसार $R \neq S$ असल्याने ते \emph{दोन्ही} $\Nat$ वरील
क्रमसंबंधाशी एकरूप असू शकत नाहीत. पण प्रत्यक्षात $S$ ऐवजी $R$ (किंवा
उलट) \emph{हाच} क्रमसंबंध आहे, असा दावा करणे मनमानीचे आणि म्हणून काहीसे
अडचणीचे ठरेल. (संचसैद्धान्तिक न्यूनीकरणवादाविरुद्ध बेनासेराफ यांनी
\citeyear{Benacerraf1965} मध्ये प्रसिद्ध केलेल्या युक्तिवादाचे हे अतिशय
सोपे उदाहरण आहे. आपण त्याकडे आणखी काही वेळा परत येऊ.)

दुसरे: \emph{प्रत्येक} संबंध एखाद्या संचाशी एकरूप मानला पाहिजे असे आपल्याला
वाटत असेल, तर संचसदस्यत्वाचा $\in$ हा संबंधसुद्धा एक विशिष्ट संच असला
पाहिजे. तो $\Setabs{\tuple{x,y}}{x \in y}$ हा संच असावा लागेल. पण हा
संच अस्तित्वात आहे का? रसेलची विरोधापत्ती लक्षात घेता, असा संच अस्तित्वात
आहे हा दावा सहज स्वीकारता येत नाही. खरे तर,
संच सिद्धांत
काटेकोरपणे स्वयंसिद्धकांवर आधारित सिद्धांत म्हणून विकसित करता येतो,
आणि तो सिद्धांत या संचाचे अस्तित्व नाकारतो.
म्हणून काही संबंधांना संच म्हणून वागवता येत असले, तरी संचसदस्यत्वाचा
संबंध अपवादात्मक प्रकरण म्हणून घ्यावा लागेल.

तिसरे: संबंध संचांशी ``एकरूप'' मानताना, $\tuple{x,y} \in R$ यासाठी
$Rxy$ असे लिहिण्याची मुभा आपण घेतली. सदस्यत्वाचा ``$\in$'' हा संबंध
विधेय \emph{म्हणून} वापरला, तर हे योग्य आहे. पण ``$\in$'' हे एका
विशिष्ट प्रकारच्या संचाचे नाव आहे असे मानले, तर ``$\tuple{x,y} \in R$''
या व्यंजकात संचांचा निर्देश करणारी फक्त तीन एकवाची पदे राहतात:
``$\tuple{x,y}$'', ``$\in$'' आणि ``$R$''. अशा नावांच्या यादीतून विधान
व्यक्त होऊ शकत नाही; ``तो कप लेखणीधारक ते टेबल'' या निरर्थक शब्दमालेइतकीच
ती निष्फळ आहे. म्हणून पुन्हा, काही संबंधांना संच म्हणून वागवता येत असले,
तरी संचसदस्यत्वाचा संबंध अपवादात्मक प्रकरण असला पाहिजे. (यात फ्रेगे यांच्या
\emph{घोडा} संकल्पनेच्या विरोधापत्तीचे एक सोपे रूप आणि विट्गेन्स्टाइन यांनी
रसेल यांच्याविरुद्ध मांडलेला एक प्रसिद्ध आक्षेप एकत्र आले आहेत.)

मग यातून काय निष्पन्न होते? संख्यांवरील संबंध हे संच आहेत असे म्हणण्यात
काही \emph{चूक} नाही. ते म्हणण्यामागचा हेतू मात्र नीट समजून घ्यावा लागेल.
आपण सत्तामीमांसक एकरूपतेचे तथ्य सांगत नाही. फक्त एवढेच नोंदवत आहोत की,
काही संदर्भांत काही संबंधांना काही विशिष्ट संच \emph{म्हणून वागवता}
येते (आणि आपण तसे वागवणार आहोत).











\section{संबंधांचे विशेष गुणधर्म}\label{sfr:rel:prp:sec}

\begin{intro}
काही प्रकारचे संबंध इतके वारंवार आढळतात की त्यांना विशेष नावे दिली
आहेत. उदाहरणार्थ, $\le$ आणि $\subseteq$ हे आपापल्या प्रांतांतील घटकांना
सारख्याच प्रकारे जोडतात (उदा., $\le$ साठी $\Nat$ आणि $\subseteq$ साठी
$\Pow{A}$). त्यांच्यातील साम्य आणि भेद नेमके समजण्यासाठी संबंधांचे काही
विशेष गुणधर्मांनुसार वर्गीकरण करू. यांपैकी काही गुणधर्मांची संयोजने विशेष
महत्त्वाची ठरतात: क्रम आणि तुल्यता संबंध.
\end{intro}

\begin{defn}[परावर्तिता]
$R \subseteq A^2$ हा संबंध \emph{परावर्ती} तेव्हा आणि केवळ तेव्हाच असतो,
जेव्हा प्रत्येक $x \in A$ साठी $Rxx$ असते.
\end{defn}

\begin{defn}[संक्रमकता]
$R \subseteq A^2$ हा संबंध \emph{संक्रमक} तेव्हा आणि केवळ तेव्हाच असतो,
जेव्हा $Rxy$ आणि $Ryz$ असले, की $Rxz$ सुद्धा असते.
\end{defn}

\begin{defn}[सममिती]
$R \subseteq A^2$ हा संबंध \emph{सममित} तेव्हा आणि केवळ तेव्हाच असतो,
जेव्हा $Rxy$ असले, की $Ryx$ सुद्धा असते.
\end{defn}

\begin{defn}[प्रतिसममिती]
$R \subseteq A^2$ हा संबंध \emph{प्रति\-सम\-मित} तेव्हा आणि केवळ तेव्हाच
असतो, जेव्हा $Rxy$ आणि $Ryx$ दोन्ही असले, की $x=y$ असते (दुसऱ्या शब्दांत:
$x\neq y$ असेल, तर $\lnot Rxy$ किंवा $\lnot Ryx$ असते).
\end{defn}

\begin{explain}
सममित संबंधात $Rxy$ आणि $Ryx$ नेहमी एकत्र सत्य असतात किंवा दोन्ही असत्य
असतात. प्रतिसममित संबंधात $Rxy$ आणि $Ryx$ दोन्ही सत्य असण्यासाठी $x = y$
असणे आवश्यक असते. पण $x = y$ असताना $Rxy$ आणि $Ryx$ सत्य \emph{असलेच
पाहिजेत}, अशी अट यात नाही; फक्त त्यांना मनाई नाही. म्हणून प्रतिसममित
संबंध परावर्ती असू शकतो, पण प्रत्येक प्रतिसममित संबंध परावर्ती असतोच
असे नाही. तसेच, प्रतिसममित असणे आणि केवळ सममित नसणे या वेगळ्या अटी आहेत.
खरे तर, एकच संबंध एकाच वेळी सममित आणि प्रतिसममितही असू शकतो
(उदा., एकरूपता संबंध).
\end{explain}

\begin{defn}[संयोजितता]
$R \subseteq A^2$ या संबंधाला \emph{संयोजित} म्हणतात, जर सर्व $x,y\in A$
साठी $x \neq y$ असल्यास $Rxy$ किंवा $Ryx$ असते.
\end{defn}

\begin{prob}
पुढील गुणधर्म असणाऱ्या संबंधांची उदाहरणे द्या: (a) परावर्ती आणि सममित,
पण संक्रमक नाही; (b) परावर्ती आणि प्रतिसममित; (c) प्रतिसममित व संक्रमक,
पण परावर्ती नाही; आणि (d) परावर्ती, सममित व संक्रमक. संख्या किंवा संच
यांवरील संबंध वापरू नका.
\end{prob}

\begin{defn}[अपरावर्तिता]
$R \subseteq A^2$ या संबंधाला \emph{अपरावर्ती} म्हणतात, जर सर्व $x \in A$
साठी $Rxx$ असत्य असते.
\end{defn}

\begin{defn}[असममिती]
$R \subseteq A^2$ या संबंधाला \emph{असममित} म्हणतात, जर $x,y\in A$ अशा
कोणत्याही जोडीसाठी $Rxy$ आणि $Ryx$ दोन्ही सत्य नसतात.
\end{defn}

लक्षात घ्या: $A \neq \emptyset$ असेल, तर $A$ वरील कोणताही अपरावर्ती
संबंध परावर्ती नसतो, आणि $A$ वरील प्रत्येक असममित संबंध प्रतिसममितही
असतो. तरीही, काही $R \subseteq A^2$ संबंध परावर्तीही नसतात आणि
अपरावर्तीही नसतात; तसेच काही प्रतिसममित संबंध असममित नसतात.












\section{तुल्यता संबंध}\label{sfr:rel:eqv:sec}

एखाद्या संचावरील एकरूपता संबंध परावर्ती, सममित आणि संक्रमक असतो.
हे तिन्ही गुणधर्म असणारे $R$ संबंध वारंवार आढळतात.

\begin{defn}[तुल्यता संबंध]
परावर्ती, सममित आणि संक्रमक असणाऱ्या $R \subseteq A^2$ या संबंधाला
\emph{तुल्यता संबंध} म्हणतात. $A$ मधील घटक $x$ आणि $y$ यांना
\emph{$R$-तुल्य} म्हणतात, जर $Rxy$ असते.
\end{defn}

तुल्यता संबंधांमुळे \emph{तुल्यतावर्ग} ही संकल्पना मिळते. तुल्यता संबंध
प्रांताचे वेगवेगळ्या भागांत ``तुकडे'' करतो. प्रत्येक भागात सर्व वस्तू
एकमेकींशी संबंधित असतात; भिन्न भागांतील कोणत्याही वस्तू एकमेकींशी संबंधित
नसतात. काही वेळा या भागांबद्दल \emph{थेट} बोलणे उपयोगी ठरते. त्यासाठी
पुढील व्याख्या करू:

\begin{defn}\label{sfr:rel:eqv:def:equivalenceclass}
$R \subseteq A^2$ हा तुल्यता संबंध असू द्या. प्रत्येक $x \in A$ साठी,
$A$ मधील $x$ चा \emph{तुल्यतावर्ग} म्हणजे
$\equivrep{x}{R} = \Setabs{y \in A}{Rxy}$ हा संच. $R$ नुसार $A$ चा
\emph{भागाकारसंच} म्हणजे $\equivclass{A}{R} = \Setabs{\equivrep{x}{R}}{x \in A}$,
अर्थात या तुल्यतावर्गांचा संच.
\end{defn}

पुढील निष्कर्ष तुल्यतावर्गाच्या व्याख्येला आधार देतो: तुल्यतावर्ग हे खरोखर
$A$ चे विभाजन करणारे भाग असतात, हे तो सिद्ध करतो.

\begin{prop}
$R \subseteq A^2$ हा तुल्यता संबंध असेल, तर $Rxy$ तेव्हा आणि केवळ
तेव्हाच, जेव्हा $\equivrep{x}{R} = \equivrep{y}{R}$.
\end{prop}

\begin{proof}
डावीकडून उजवीकडे जाणाऱ्या दिशेसाठी, $Rxy$ असे समजा आणि
$z \in \equivrep{x}{R}$ असू द्या. व्याख्येनुसार $Rxz$. $R$ हा तुल्यता
संबंध असल्याने $Ryz$. (सविस्तर सांगायचे तर: $Rxy$ असल्याने आणि $R$
सममित असल्याने $Ryx$; तसेच $Rxz$ असल्याने आणि $R$ संक्रमक असल्याने
$Ryz$.) म्हणून $z \in \equivrep{y}{R}$. हे कोणत्याही अशा घटकासाठी
लागू असल्याने $\equivrep{x}{R} \subseteq \equivrep{y}{R}$. अगदी
त्याचप्रमाणे $\equivrep{y}{R} \subseteq \equivrep{x}{R}$. म्हणून
विस्तारात्मकतेनुसार $\equivrep{x}{R} = \equivrep{y}{R}$.

उजवीकडून डावीकडे जाणाऱ्या दिशेसाठी,
$\equivrep{x}{R} = \equivrep{y}{R}$ असे समजा. $R$ परावर्ती असल्याने
$Ryy$, म्हणून $y \in \equivrep{y}{R}$. मग
$\equivrep{x}{R} = \equivrep{y}{R}$ या गृहीतकामुळे
$y \in \equivrep{x}{R}$ सुद्धा. म्हणून $Rxy$.
\end{proof}

\begin{ex}
तुल्यता संबंधांचे एक चांगले उदाहरण शेषांक अंकगणितातून मिळते. कोणत्याही
$a$, $b$ आणि $n \in \PosInt$ साठी, $a$ ला $n$ ने भागल्यावर येणारी
बाकी आणि $b$ ला $n$ ने भागल्यावर येणारी बाकी सारखी असेल, तेव्हा आणि
केवळ तेव्हाच $a \equiv_n b$ असे म्हणू. (चिन्हांत मांडायचे तर:
$a \equiv_n b$ तेव्हा आणि केवळ तेव्हाच, जेव्हा एखाद्या $k \in \Int$
साठी $a - b = kn$.) मग कोणत्याही $n$ साठी $\equiv_n$ हा तुल्यता
संबंध असतो. $\equiv_n$ मुळे नेमके $n$ भिन्न तुल्यतावर्ग मिळतात; म्हणजे
$\equivclass{\Nat}{\equiv_n}$ मध्ये $n$ घटक असतात. ते असे:
$n$ ने निःशेष भाग जाणाऱ्या संख्यांचा संच, म्हणजे $\equivrep{0}{\equiv_n}$;
$n$ ने भागल्यावर बाकी $1$ येणाऱ्या संख्यांचा संच, म्हणजे
$\equivrep{1}{\equiv_n}$; \ldots; आणि $n$ ने भागल्यावर बाकी $n-1$
येणाऱ्या संख्यांचा संच, म्हणजे $\equivrep{n-1}{\equiv_n}$.
\end{ex}

\begin{prob}
कोणत्याही $n \in \PosInt$ साठी $\equiv_n$ हा तुल्यता संबंध आहे आणि
$\equivclass{\Nat}{\equiv_n}$ मध्ये नेमके $n$ सदस्य आहेत, हे दाखवा.
\end{prob}











\section{क्रम}\label{sfr:rel:ord:sec}

\begin{explain}
अनेक तुलनांमध्ये आपण काही वस्तू इतर वस्तूंपेक्षा एखाद्या दृष्टीने
``लहान'', ``समान'' किंवा ``मोठ्या'' आहेत असे म्हणतो. यात \emph{क्रम}
संबंध येतात. पण क्रमसंबंधांचे वेगवेगळे प्रकार आहेत. उदाहरणार्थ, काहींमध्ये
कोणत्याही दोन वस्तूंची तुलना करता आली पाहिजे, तर इतरांमध्ये तशी अट नसते.
काहींमध्ये एकरूपतेचा समावेश असतो (जसे $\le$), तर काहींमध्ये नसतो
(जसे $<$). येथे या प्रकारांचे वर्गीकरण उपयोगी पडेल.
\end{explain}

\begin{defn}[पूर्वक्रम]
परावर्ती आणि संक्रमक असणाऱ्या संबंधाला \emph{पूर्वक्रम} म्हणतात.
\end{defn}

\begin{defn}[अंशतः क्रम]
प्रतिसममितही असणाऱ्या पूर्वक्रमाला \emph{अंशतः क्रम} म्हणतात.
\end{defn}

\begin{defn}[रेषीय क्रम]\label{sfr:rel:ord:def:linearorder}
संयोजितही असणाऱ्या अंशतः क्रमाला \emph{पूर्ण क्रम} किंवा \emph{रेषीय क्रम}
म्हणतात.
\end{defn}

\begin{ex}
प्रत्येक रेषीय क्रम अंशतः क्रमही असतो आणि प्रत्येक अंशतः क्रम पूर्वक्रमही
असतो; पण यांची उलट विधाने सत्य नाहीत. $A$ वरील सार्वत्रिक संबंध परावर्ती
आणि संक्रमक असल्याने तो पूर्वक्रम असतो. पण $A$ मध्ये एकाहून अधिक
घटक असतील, तर सार्वत्रिक संबंध प्रतिसममित नसतो, म्हणून तो
अंशतः क्रम नसतो.
\end{ex}

\begin{ex}
$\Bin^*$ वरील \emph{अधिक लांब नसण्याचा} $\preccurlyeq$ हा संबंध पाहा:
$x \preccurlyeq y$ तेव्हा आणि केवळ तेव्हाच, जेव्हा $\len{x} \le \len{y}$.
हा पूर्वक्रम (परावर्ती आणि संक्रमक) आहे, शिवाय संयोजितही आहे; पण
प्रतिसममित नसल्याने अंशतः क्रम नाही. उदाहरणार्थ,
$01 \preccurlyeq 10$ आणि $10 \preccurlyeq 01$, पण $01 \neq 10$.
\end{ex}

\begin{ex}
संचांच्या संचावरील $\subseteq$ हा एक महत्त्वाचा अंशतः क्रम आहे.
सर्वसाधारणपणे तो रेषीय क्रम नसतो. कारण $a \neq b$ असेल आणि
$\Pow{\{a, b\}} = \{\emptyset, \{a\}, \{b\}, \{a,b\}\}$ विचारात घेतला,
तर $\{a\} \nsubseteq \{b\}$, $\{a\} \neq \{b\}$ आणि
$\{b\} \nsubseteq \{a\}$ असे दिसते.
\end{ex}

\begin{ex}
\emph{निःशेष विभाज्यतेचा} संबंध हा रेषीय नसलेला अंशतः क्रम देतो.
$n$ आणि $m$ या पूर्णांकांसाठी, $n \mid m$ म्हणजे $n$ ने $m$ ला निःशेष
भाग जातो, अर्थात एखादा पूर्णांक $k$ असा असतो की $m = kn$.
$\Nat$ वर हा अंशतः क्रम आहे, पण रेषीय क्रम नाही: उदाहरणार्थ,
$2 \nmid 3$ आणि $3 \nmid 2$. $\Int$ वरील संबंध म्हणून पाहिल्यास
विभाज्यता केवळ पूर्वक्रम आहे, कारण ती प्रतिसममित नाही:
$1 \mid -1$ आणि $-1 \mid 1$, पण $1 \neq -1$.
\end{ex}

\begin{ex}
अनुक्रमांच्या $A^*$ या संचावरील \emph{विस्तार} संबंध असा आहे:
$s \sqsubseteq s'$ तेव्हा आणि केवळ तेव्हाच, जेव्हा $s = \emptyseq$
(रिक्त अनुक्रम), $s = s'$, किंवा $s = \tuple{s_1, \dots, s_n}$ आणि
$s' = \tuple{s_1, \dots, s_n, s_{n+1}, \dots, s_m}$.
$s \sqsubseteq s'$ असल्यास $s$ हा $s'$ चा \emph{आरंभीचा खंड} आहे असेही
म्हणतो. $A^*$ वरील विस्तार संबंध अंशतः क्रम आहे, पण रेषीय क्रम नाही;
उदा., $a \neq b$ असेल, तर $ab \not\sqsubseteq ba$ आणि
$ba \not\sqsubseteq ab$.
\end{ex}

\begin{defn}[काटेकोर क्रम]
\emph{काटेकोर क्रम} म्हणजे अपरावर्ती, असममित आणि संक्रमक असणारा संबंध.
\end{defn}

\begin{defn}[काटेकोर रेषीय क्रम]\label{sfr:rel:ord:def:strictlinearorder}
संयोजितही असणाऱ्या काटेकोर क्रमाला \emph{काटेकोर पूर्ण क्रम} किंवा
\emph{काटेकोर रेषीय क्रम} म्हणतात.
\end{defn}

\begin{ex}
$\le$ हा काटेकोर रेषीय क्रम $<$ याला अनुरूप रेषीय क्रम आहे.
$\subseteq$ हा काटेकोर क्रम $\subsetneq$ याला अनुरूप अंशतः क्रम आहे.
\end{ex}

$A$ वरील कोणत्याही काटेकोर क्रम $R$ मध्ये कर्ण $\Id{A}$, म्हणजे सर्व
$\tuple{x, x}$ जोड्या, मिळवल्यास अंशतः क्रम तयार होतो. (याला $R$ चे
\emph{परावर्ती संवरण} म्हणतात.) उलट, अंशतः क्रमातून $\Id{A}$ काढून
टाकल्यास काटेकोर क्रम मिळतो. पुढील दोन निष्कर्ष हे नेमके स्पष्ट करतात.

\begin{prop}\label{sfr:rel:ord:prop:stricttopartial}
$R$ हा $A$ वरील काटेकोर क्रम असेल, तर $R^+ = R \cup \Id{A}$ हा अंशतः
क्रम असतो. शिवाय, $R$ हा काटेकोर रेषीय क्रम असेल, तर $R^+$ हा रेषीय
क्रम असतो.
\end{prop}

\begin{proof}
$R$ हा काटेकोर क्रम आहे असे समजा, म्हणजे $R \subseteq A^2$ आणि $R$
अपरावर्ती, असममित व संक्रमक आहे. $R^+ = R \cup \Id{A}$ असू द्या.
$R^+$ परावर्ती, प्रतिसममित आणि संक्रमक आहे हे दाखवायचे आहे.

$R^+$ स्पष्टपणे परावर्ती आहे, कारण सर्व $x \in A$ साठी
$\tuple{x, x} \in \Id{A} \subseteq R^+$.

$R^+$ प्रतिसममित आहे हे दाखवण्यासाठी व्याघात काढू. $R^+xy$ आणि $R^+yx$
असूनही $x \neq y$ आहे असे समजा. $\tuple{x,y} \in R \cup \Id{A}$,
पण $\tuple{x, y} \notin \Id{A}$ असल्याने $\tuple{x, y} \in R$, म्हणजे
$Rxy$ असले पाहिजे. त्याचप्रमाणे $Ryx$. पण $R$ असममित आहे या
गृहीतकाशी हे विसंगत आहे.

संक्रमकता दाखवण्यासाठी $R^+xy$ आणि $R^+yz$ असे समजा.
$\tuple{x, y} \in R$ आणि $\tuple{y,z} \in R$ दोन्ही असतील, तर $R$
संक्रमक असल्याने $\tuple{x, z} \in R$. अन्यथा $\tuple{x, y} \in \Id{A}$,
म्हणजे $x = y$, किंवा $\tuple{y, z} \in \Id{A}$, म्हणजे $y = z$.
पहिल्या प्रकरणात गृहीतकानुसार $R^+yz$ आणि $x = y$, म्हणून $R^+xz$.
दुसऱ्या प्रकरणातही त्याचप्रमाणे. प्रत्येक प्रकरणात $R^+xz$, म्हणून
$R^+$ संक्रमकही आहे.

``शिवाय'' या भागासाठी, $R$ संयोजितही आहे असे समजा. म्हणून सर्व
$x \neq y$ साठी $Rxy$ किंवा $Ryx$, म्हणजे $\tuple{x, y} \in R$ किंवा
$\tuple{y, x} \in R$. $R \subseteq R^+$ असल्याने $R^+$ बाबतही हे सत्य
राहते, म्हणून $R^+$ संयोजितही आहे.
\end{proof}

\begin{prop}\label{sfr:rel:ord:prop:partialtostrict}
$R$ हा $A$ वरील अंशतः क्रम असेल, तर $R^- = R \setminus \Id{A}$ हा
काटेकोर क्रम असतो. शिवाय, $R$ हा रेषीय क्रम असेल, तर $R^-$ हा काटेकोर
रेषीय क्रम असतो.
\end{prop}

\begin{proof}
हे सरावासाठी ठेवले आहे.
\end{proof}

\begin{prob}
\ref{sfr:rel:ord:prop:partialtostrict} याची सिद्धता द्या.
\end{prob}

काटेकोर रेषीय क्रम विस्तारात्मकतेसारखा एक गुणधर्म पाळतात, हे पुढील
सोप्या निष्कर्षातून दिसते:

\begin{prop}\label{sfr:rel:ord:prop:extensionality-strictlinearorders}
$<$ हा $A$ वरील काटेकोर रेषीय क्रम असेल, तर:
\[
  (\forall a, b \in A)((\forall x \in A)(x < a \liff x < b) \lif a = b).
\]
\end{prop}

\begin{proof}
$(\forall x \in A)(x < a \liff x < b)$ असे समजा. $a < b$ असेल, तर
$a < a$; पण $<$ अपरावर्ती आहे याच्याशी हे विसंगत आहे. म्हणून $a \nless b$.
अगदी त्याचप्रमाणे $b \nless a$. $<$ संयोजित असल्याने $a = b$.
\end{proof}











\section{आलेख}\label{sfr:rel:grp:sec}

\emph{आलेख} म्हणजे अशी आकृती की जिच्यात ``गाठी'' किंवा ``शिखरे''
(``शिखर'' या शब्दाचे अनेकवचन) म्हटले जाणारे बिंदू कडांनी जोडलेले असतात. विविक्त गणित आणि संगणकशास्त्रात
आलेख सर्वत्र वापरले जातात. विविध जाळ्यांसारख्या ठोस गोष्टींपासून ते
निर्णयांच्या शक्य परिणामांसारख्या अमूर्त संरचनांपर्यंत अनेक संबंध आणि
संरचना दर्शवण्यासाठी व त्यांचे दृश्य रूप देण्यासाठी ते अतिशय उपयोगी
असतात. ग्रंथसाहित्यात आलेखांचे अनेक प्रकार आढळतात. उदाहरणार्थ, कडांना
दिशा आहे की नाही, नावे आहेत की नाहीत, एखाद्या शिखरापासून त्याच शिखरापर्यंत
कड असू शकते की नाही, त्याच शिखरांमध्ये अनेक कडा असू शकतात की नाहीत,
इत्यादी बाबतींत हे प्रकार वेगळे असतात. \emph{दिशित आलेखांचा} संबंधांशी
विशेष संबंध आहे.

\begin{defn}[दिशित आलेख]
\emph{दिशित आलेख} $G = \tuple{V, E}$ म्हणजे \emph{शिखरांचा} $V$ हा संच
आणि \emph{कडांचा} $E \subseteq V^2$ हा संच.
\end{defn}

\begin{explain}
आपल्या व्याख्येनुसार आलेख म्हणजे एक संच आणि त्या संचावरील एक संबंध.
अर्थात, आलेखांबद्दल बोलताना त्यांचे चित्र काढले जावे ही अपेक्षा स्वाभाविक
आहे: $v_1$ आणि $v_2$ ही शिखरे बाणाने तेव्हा आणि केवळ तेव्हाच जोडून
आलेख काढता येतो, जेव्हा $\tuple{v_1, v_2} \in E$.
नुसता संबंध आणि आलेख यांतील फरक एवढाच की आलेखात शिखरांचा संचही स्पष्ट
केलेला असतो, म्हणजे आलेखात एकाकी शिखरे असू शकतात. महत्त्वाचा मुद्दा असा:
$X$ संचावरील प्रत्येक $R$ संबंधाकडे $\tuple{X, R}$ हा दिशित आलेख म्हणून
पाहता येते; आणि उलट, $\tuple{V, E}$ या दिशित आलेखाकडे, $V$ हा संच स्पष्ट
केलेला असताना, $E \subseteq V^2$ हा संबंध म्हणून पाहता येते.
\end{explain}

\begin{ex}
$V = \{1, 2, 3, 4\}$ आणि $E =
\{\tuple{1,1}, \allowbreak \tuple{1, 2}, \allowbreak \tuple{1, 3},
\allowbreak \tuple{2, 3}\}$ असलेला $\tuple{V, E}$ हा आलेख असा दिसतो:
\begin{align*}
& \begin{tikzpicture}[->,node distance=2cm]
  \node[draw,circle] (A) {$1$};
  \node[draw,circle] (B) [right of=A] {$2$};
  \node[draw,circle] (C) [below of=B] {$3$};
  \node[draw,circle] (D) [right of=B] {$4$};
  \draw (A) to [loop above]  (A);
  \draw (A) to  (B);
  \draw (A) to  (C);
  \draw (B) to  (C);
  \end{tikzpicture}
\intertext{हा आलेख $V' = \{1, 2, 3\}$ असणाऱ्या $\tuple{V', E}$ या
  आलेखापेक्षा वेगळा आहे; तो असा दिसतो:}
& \begin{tikzpicture}[->,node distance=2cm]
  \node[draw,circle] (A) {$1$};
  \node[draw,circle] (B) [right of=A] {$2$};
  \node[draw,circle] (C) [below of=B] {$3$};
  \draw (A) to [loop above]  (A);
  \draw (A) to  (B);
  \draw (A) to  (C);
  \draw (B) to  (C);
  \end{tikzpicture}
\end{align*}
\end{ex}

\begin{prob}
  $\{1, 2, 3, 4\}$ या संचावरील लहान-किंवा-समान असण्याचा $\le$ संबंध
  आलेख म्हणून पाहा आणि त्याची आकृती काढा.
\end{prob}











\section{वृक्ष}\label{sfr:rel:tre:sec}

तर्कशास्त्राच्या सर्व शाखांत महत्त्वाची भूमिका बजावणारा अंशतः क्रमाचा
एक विशिष्ट प्रकार म्हणजे \emph{वृक्ष}. तर्कशास्त्राच्या प्राथमिक भागांत
सांत वृक्ष येतात: उदाहरणार्थ, सूत्रे विन्यासवृक्षात विभागून
समजून घेता येतात, आणि अनेक निष्पत्ती पद्धतींतील
निष्पत्ती सुद्धा सांत वृक्षांच्या स्वरूपात असतात.

विधानीय आणि प्रथमस्तरीय तर्कशास्त्राच्या परिपूर्णता प्रमेयांच्या सिद्धतेतच
अनंत वृक्ष येऊ लागतात; पुढे संपूर्ण गणिती तर्कशास्त्रात त्यांचा वापर होतो.

वृक्षाची संचसैद्धान्तिक संकल्पना आणि आलेख सिद्धांतातील वृक्षाची संकल्पना
यांचा निकटचा संबंध आहे. एका सांत वृक्षाचे चित्र असे आहे:

\begin{center}
\begin{tikzpicture}[nodes={draw, circle}, -]
\node{$r$} [grow'=up]
    child { node {$a$}
        child { node {$c$} }
        child { node {$d$} }
        child { node {$e$} }
    }
    child { node {$b$} };
\end{tikzpicture}
\end{center}

सर्वांत खालचे $r$ हे शिखर मूळ आहे. $r$ व्यतिरिक्त प्रत्येक शिखराच्या
लगेच खाली नेमके एक पालक शिखर आहे. $x$ पासून कडांवरून वर जात $y$ पर्यंत
पोहोचता येत असेल, तर $x$ चा $y$ शी असणारा संबंध म्हणजे $x$ हा $y$ चा
\emph{पूर्वज} असणे, असे समजू शकतो.

वृक्षातील पूर्वज संबंध हा काटेकोर अंशतः क्रम असतो. यावरून संचसैद्धान्तिक
व्याख्या सुचते. ती मांडण्यासाठी दोन संकल्पना लागतात. $\le$ ने अंशतः
क्रमित केलेल्या $A$ या संचातील \emph{लघुतम घटक} म्हणजे असा
घटक $x \in A$ की सर्व $y \in A$ साठी $x \le y$ असते.
एखाद्या संचाच्या प्रत्येक अरिक्त उपसंचात लघुतम घटक असेल, तर तो संच
$\le$ ने \emph{सुक्रमित} आहे असे म्हणतात.

\begin{defn}[वृक्ष]
\emph{वृक्ष} म्हणजे $T = \tuple{A, \le}$ अशी जोडी की $A$ हा संच असतो,
$\le$ हा $A$ वरील अंशतः क्रम असतो, त्याचा $r \in A$ हा एकमेव लघुतम
घटक असतो (त्याला \emph{मूळ} म्हणतात), आणि सर्व $x \in A$ साठी
$\Setabs{y}{y \le x}$ हा संच $\le$ ने सुक्रमित असतो.
\end{defn}

\begin{defn}[उत्तरवर्ती]
$T = \tuple{A, \le}$ हा वृक्ष आहे असे समजा.
$x,y \in A$, $x < y$ असतील आणि $x < z < y$ असा कोणताही $z \in A$
नसेल, तर $y$ हा $x$ चा \emph{उत्तरवर्ती} आहे असे म्हणतो.
\end{defn}

$x \in A$ च्या उत्तरवर्तींना त्याची \emph{अपत्ये} असेही म्हणतात.
$y$ हा $x$ चा उत्तरवर्ती असेल, तर $x$ ला $y$ चा \emph{पूर्ववर्ती}
किंवा \emph{पालक} म्हणतो.

\begin{prop}
$\tuple{A,\le}$ हा वृक्ष असेल, तर मूळ वगळता प्रत्येक $x \in A$ ला
जास्तीत जास्त एक पूर्ववर्ती असतो.
\end{prop}

\begin{proof}
  $y_1 < x$, $y_2 < x$ आणि $y_1 \neq y_2$ असे समजा. मग
  $\{y_1, y_2\} \subseteq \Setabs{z}{z<x}$. $\Setabs{z}{z<x}$ हा संच
  $\le$ ने सुक्रमित असल्याने त्याच्या $\{y_1, y_2\}$ या उपसंचाला लघुतम
  घटक असतो. तो स्पष्टपणे $y_1$ किंवा $y_2$ असला पाहिजे. म्हणून
  $y_1 \le y_2$ किंवा $y_2 \le y_1$. आपण $y_1 \neq y_2$ गृहीत धरले,
  म्हणून प्रत्यक्षात $y_1 < y_2$ किंवा $y_2 < y_1$. शिवाय,
  $y_1 < x$ आणि $y_2 < x$ या गृहीतकांमुळे $y_1 < y_2 < x$ किंवा
  $y_2 < y_1 < x$. म्हणून $y_1$ आणि $y_2$ हे दोघेही $x$ चे पूर्ववर्ती
  असू शकत नाहीत.
\end{proof}

\begin{defn}
$T = \tuple{A, \le}$ या वृक्षातील $A$ हा अनंत संच असेल, तर वृक्षाला
\emph{अनंत}, आणि अन्यथा \emph{सांत} म्हणतात. $T$ मधील प्रत्येक
$x \in A$ ला फक्त सांत संख्येने उत्तरवर्ती असतील, तर $T$ ला
\emph{सांत शाखनाचा} वृक्ष म्हणतात.
\end{defn}

\begin{defn}[शाखा]
$T = \tuple{A, \le}$ हा वृक्ष दिला असता, $T$ ची \emph{शाखा} म्हणजे
$T$ मधील अधिकतम शृंखला; म्हणजे $B \subseteq A$ असा संच की कोणत्याही
$x, y \in B$ साठी $x \le y$ किंवा $y \le x$ असते, आणि प्रत्येक
$z \in X \setminus B$ साठी असा $u \in B$ असतो की $z \le u$ आणि
$u \le z$ दोन्ही असत्य असतात.

$T$ च्या सर्व शाखांचा संच दर्शवण्यासाठी आपण $[T]$ वापरतो.
\end{defn}

\begin{ex}
सांत शाखनाच्या वृक्षाचे एक प्रसिद्ध उदाहरण म्हणजे $0$ आणि $1$ यांच्या
सांत अनुक्रमांचा \emph{अनंत द्विमान वृक्ष}. त्याला कधी $\{0,1\}^*$ किंवा
$\Bin^*$ असे लिहितात आणि विस्तार संबंध $\sqsubseteq$ ने क्रमित करतात
(उदा., $101 \sqsubseteq 101101$). कोणत्याही द्विमान चिन्हमालेच्या शेवटी
$0$ किंवा $1$ जोडून तिचा विस्तार नेहमी करता येतो. म्हणून या वृक्षात
अनंत घटक असतात: प्रत्येक $s$ या घटकाला नेमके दोन उत्तरवर्ती, $s0$ आणि
$s1$, असतात. त्याचे मूळ म्हणजे $\emptyseq$ हा रिक्त अनुक्रम.
\end{ex}

\begin{ex}
थोडे अधिक सर्वसाधारण उदाहरण म्हणजे नैसर्गिक संख्यांच्या सांत अनुक्रमांचा
$\Nat^*$ हा संच आणि त्यावरील विस्तार संबंध $\sqsubseteq$; हाही वृक्ष
असतो. तो स्पष्टपणे सांत शाखनाचा नाही: प्रत्येक $s \in \Nat^*$ ला अनंत
उत्तरवर्ती $sn$ असतात, प्रत्येक $n \in \Nat$ साठी एक. $\sqsubseteq$
नुसार बंद असलेला प्रत्येक $A \subseteq \Nat^*$ हा $\Nat^*$ चा
\emph{उपवृक्ष} असतो. (म्हणजे $A$ असा आहे की $s \in A$ आणि
$s' \sqsubseteq s$ असतील, तर $s' \in A$ सुद्धा.) सर्व सांत वृक्ष
$\Nat^*$ चे सांत उपवृक्ष म्हणून दर्शवता येतात.
\end{ex}

\begin{prop}[K\H{o}nig यांचे उपप्रमेय]
$T = \tuple{A,\le}$ हा सांत शाखनाचा अनंत वृक्ष असेल, तर $T$ ला
एक अनंत शाखा असते.
\end{prop}

संगणनीयता सिद्धांतात वारंवार वापरले जाणारे K\H{o}nig यांच्या उपप्रमेयाचे
एक विशेष रूप, ज्याला \emph{दुर्बल K\H{o}nig उपप्रमेय} म्हणतात, असे आहे:
$\{0,1\}^*$ च्या प्रत्येक अनंत उपवृक्षाला एक अनंत शाखा असते.











\section{संबंधांवरील क्रिया}\label{sfr:rel:ops:sec}

संबंधांत बदल करणे किंवा त्यांना एकत्र करणे अनेकदा उपयोगी पडते.
\ref{sfr:rel:ord:prop:stricttopartial} मध्ये आपण संबंधांचा
\emph{संयोग} पाहिला; तो म्हणजे दोन संबंधांकडे जोड्यांचे संच म्हणून
पाहिल्यावर त्यांचा संयोग. त्याचप्रमाणे,
\ref{sfr:rel:ord:prop:partialtostrict} मध्ये संबंधांचा सापेक्ष फरक
पाहिला. संबंधांवर करता येणाऱ्या आणखी काही क्रिया पुढीलप्रमाणे आहेत.

\begin{defn}\label{sfr:rel:ops:relationoperations}
$R$, $S$ हे संबंध आणि $A$ हा कोणताही संच असू द्या.

$R$ चा \emph{व्युत्क्रम} म्हणजे $R^{-1} = \Setabs{\tuple{y, x}}{\tuple{x,
    y} \in R}$.

$R$ आणि $S$ यांचा \emph{सापेक्ष गुणाकार} म्हणजे $(R \mid S) =
\{\tuple{x, z} : \exists y(Rxy \land Syz)\}$.

$R$ चे $A$ वरील \emph{मर्यादन} म्हणजे $\funrestrictionto{R}{A}= R
\cap A^2$.

$R$ चा $A$ वर \emph{प्रयोग} म्हणजे $\funimage{R}{A} = \{y :
(\exists x \in A)Rxy\}$
\end{defn}

\begin{ex}
$S \subseteq \Int^2$ हा $\Int$ वरील उत्तरवर्ती संबंध असू द्या, म्हणजे
$S = \Setabs{\tuple{x, y} \in \Int^2}{x + 1 = y}$; त्यामुळे $Sxy$
तेव्हा आणि केवळ तेव्हाच, जेव्हा $x + 1 = y$.

$S^{-1}$ हा $\Int$ वरील पूर्ववर्ती संबंध आहे, म्हणजे
$\Setabs{\tuple{x,y}\in\Int^2}{x -1 =y}$.

$S\mid S$ म्हणजे
$ \Setabs{\tuple{x,y}\in\Int^2}{x + 2 =y}$

$\funrestrictionto{S}{\Nat}$ हा $\Nat$ वरील उत्तरवर्ती संबंध आहे.

$\funimage{S}{\{1,2,3\}}$ म्हणजे $\{2, 3, 4\}$.
\end{ex}

\begin{defn}[संक्रमक संवरण] $R \subseteq A^2$ हा द्विपदी संबंध असू द्या.

$R$ चे \emph{संक्रमक संवरण} म्हणजे $R^+ = \bigcup_{0 < n \in
\Nat} R^n$, जेथे $R^1 = R$ आणि $R^{n+1} = R^n \mid R$ अशी पुनरावर्ती
व्याख्या करतो.

$R$ चे \emph{परावर्ती संक्रमक संवरण} म्हणजे $R^* = R^+ \cup \Id{A}$.
\end{defn}

\begin{ex}
$S \subseteq \Int^2$ हा उत्तरवर्ती संबंध घ्या. $S^2xy$ तेव्हा आणि केवळ
तेव्हाच, जेव्हा $x + 2 = y$; $S^3xy$ तेव्हा आणि केवळ तेव्हाच, जेव्हा
$x + 3 = y$; इत्यादी. म्हणून $S^+xy$ तेव्हा आणि केवळ तेव्हाच, जेव्हा
एखाद्या $n \geq 1$ साठी $x + n = y$. दुसऱ्या शब्दांत, $S^+xy$ तेव्हा
आणि केवळ तेव्हाच, जेव्हा $x < y$; आणि $S^*xy$ तेव्हा आणि केवळ तेव्हाच,
जेव्हा $x \le y$.
\end{ex}

\begin{prob}
$R$ चे संक्रमक संवरण खरोखर संक्रमक आहे, हे दाखवा.
\end{prob}


\chapter{फलने}\label{sfr:fun::chap}





\section{मूलभूत संकल्पना}\label{sfr:fun:bas:sec}

\begin{explain}
दिलेल्या संचातील प्रत्येक घटक दुसऱ्या एखाद्या दिलेल्या संचातील
एका ठरावीक घटक शी जोडणारी संगती म्हणजे \emph{फलन}.
उदाहरणार्थ, $1$ मिळवण्याच्या क्रियेने एक फलन ठरते: प्रत्येक संख्या~$n$
एका एकमेव संख्या~$n+1$ शी जोडली जाते.

अधिक व्यापकपणे पाहता, फलनांना आदान म्हणून जोड्या, त्रिके इत्यादी देता
येतात आणि त्यांच्याकडून कोणत्या तरी प्रकारचे प्रदान मिळते. प्राथमिक
अंकगणितातून अनेक फलने आपल्याला परिचित आहेत. उदाहरणार्थ, बेरीज आणि
गुणाकार ही फलने आहेत. ती दोन संख्या आदान म्हणून घेतात आणि तिसरी संख्या देतात.

या गणिती, अमूर्त अर्थाने फलन म्हणजे एक \emph{काळी पेटी}: कोणत्या आदानाशी
कोणते प्रदान जोडलेले आहे एवढेच महत्त्वाचे असते; प्रदानाची गणना करण्याची
पद्धत महत्त्वाची नसते.
\end{explain}

\begin{defn}[फलन]
\emph{फलन} $f \colon A \to B$ ही~$A$ मधील प्रत्येक घटक ची~$B$
मधील एका घटक शी लावलेली संगती असते.

$A$ ला~$f$ चा \emph{प्रांत} आणि $B$ ला~$f$ चा \emph{सहप्रांत} म्हणतो.
~$A$ मधील घटक ना~$f$ ची आदाने किंवा \emph{फलसाधक} म्हणतात.
फलसाधक~$x$ शी~$f$ ने जोडलेल्या~$B$ मधील घटक ला फलसाधक~$x$
साठीचे \emph{$f$ चे मूल्य} म्हणतात आणि ते~$f(x)$ असे लिहितात.

~$f$ ची \emph{व्याप्ती} $\ran{f}$ म्हणजे सहप्रांताचा असा उपसंच की
त्याचे घटक कोणत्या तरी फलसाधकासाठीची~$f$ ची मूल्ये असतात;
$\ran{f} = \Setabs{f(x)}{x \in A}$.
\end{defn}

फलनांचा विचार करण्यासाठी \ref{sfr:fun:bas:fig:function} मधील आकृती उपयोगी पडू शकते.
डावीकडील लंबवर्तुळ फलनाचा \emph{प्रांत} दर्शवते; उजवीकडील लंबवर्तुळ
त्याचा \emph{सहप्रांत} दर्शवते; आणि बाण प्रांतातील \emph{फलसाधकापासून}
सहप्रांतातील त्याच्या संगत \emph{मूल्याकडे} जातो.

\begin{figure}
  \olasset{assets/diagrams/function.tikz}
  \caption{फलन एका संचातील प्रत्येक घटक ची दुसऱ्या संचातील
    घटक शी संगती लावते. बाण प्रांतातील फलसाधकापासून
    सहप्रांतातील त्याच्या संगत मूल्याकडे जातो.}
  \label{sfr:fun:bas:fig:function}
\end{figure}

\begin{ex}
गुणाकार नैसर्गिक संख्यांच्या जोड्या आदान म्हणून घेतो आणि त्यांना प्रदान
म्हणून नैसर्गिक संख्या जोडतो. म्हणून तो $\Nat \times \Nat$ या प्रांताकडून
$\Nat$ या सहप्रांताकडे जातो. त्याची व्याप्तीही $\Nat$ च असते, कारण
प्रत्येक $n \in \Nat$ हा $n \times 1$ असतो.
\end{ex}

\begin{ex}
गुणाकार हे फलन आहे, कारण तो प्रत्येक आदानाशी---नैसर्गिक संख्यांच्या
प्रत्येक जोडीशी---एकच प्रदान जोडतो: $\times \colon \Nat^2 \to \Nat$.
याउलट, नैसर्गिक संख्या~$n$ शी $x^2 = n$ असणारी वास्तव संख्या~$x$
जोडण्याची संगती फलन होत नाही, कारण प्रत्येक धन पूर्णांक~$n$ ची दोन
वर्गमुळे असतात: $\sqrt{n}$ आणि~$-\sqrt{n}$. फक्त धन वर्गमूळ देऊन
आपण तिचे फलन करू शकतो: फक्त ऋणेतर (प्रधान) वर्गमूळ परत द्यायचे,
$\sqrt{\phantom{X}} \colon \Nat \to \Real$.
\end{ex}
\begin{quote}\small\textbf{स्रोतदुरुस्ती.} OLFUN-002 — गोठवलेल्या इंग्रजी स्रोतामध्ये सर्व नैसर्गिक संख्यांसाठी positive square root म्हटले आहे. शून्याचे प्रधान वर्गमूळ शून्य असल्यामुळे येथे ऋणेतर (प्रधान) असे दुरुस्त केले आहे; धन पूर्णांकांना दोन वास्तव वर्गमुळे असतात हे आधीचे विधान योग्य आहे.\end{quote}

\begin{ex}
वर्गातील प्रत्येक विद्यार्थ्याशी त्याची अंतिम श्रेणी जोडणारा संबंध हे
फलन आहे---एकाच वर्गात कोणत्याही विद्यार्थ्याला दोन वेगवेगळ्या अंतिम
श्रेणी मिळू शकत नाहीत. वर्गातील प्रत्येक विद्यार्थ्याशी त्याचे पालक
जोडणारा संबंध फलन नाही: विद्यार्थ्यांचे पालक शून्य, दोन किंवा अधिक असू शकतात.
\end{ex}

\begin{explain}
प्रत्येक शक्य फलसाधकासाठी फलनाचे मूल्य काय असते हे नेमक्या पद्धतीने
सांगून फलनाची व्याख्या करता येते. सूत्र देणे, मूल्याची गणना करण्याची
पद्धत वर्णन करणे किंवा प्रत्येक फलसाधकासाठीची मूल्ये यादीत देणे या
त्याच्या वेगवेगळ्या पद्धती आहेत. फलनाची व्याख्या कोणत्याही पद्धतीने
केली, तरी प्रत्येक फलसाधकासाठी एक आणि केवळ एकच मूल्य दिलेले असेल
याची खात्री केली पाहिजे.
\end{explain}


\begin{ex}
$f \colon \Nat \to \Nat$ ची व्याख्या $f(x) = x+1$ अशी करू.
या व्याख्येत $f$ हे नैसर्गिक संख्या आदान म्हणून घेऊन नैसर्गिक संख्या
प्रदान करणारे फलन आहे असे सांगितले आहे. नैसर्गिक संख्या~$x$ दिली,
की $f$ तिची उत्तरवर्ती संख्या~$x+1$ देईल असे ती सांगते.
या उदाहरणात सहप्रांत $\Nat$ हा~$f$ ची व्याप्ती नाही, कारण नैसर्गिक
संख्या~$0$ ही कोणत्याही नैसर्गिक संख्येची उत्तरवर्ती संख्या नाही.
~$f$ ची व्याप्ती सर्व धन पूर्णांकांचा संच $\Int^{+}$ आहे.
\end{ex}

\begin{ex}\label{sfr:fun:bas:examplefunext}
$g \colon \Nat \to \Nat$ ची व्याख्या $g(x) = x+2-1$ अशी करू.
यावरून $g$ हे नैसर्गिक संख्या आदान म्हणून घेऊन नैसर्गिक संख्या प्रदान
करणारे फलन आहे हे कळते. नैसर्गिक संख्या x दिली, की $g$ हा~$x$ च्या
उत्तरवर्ती संख्येच्या उत्तरवर्ती संख्येची पूर्ववर्ती संख्या,
म्हणजेच $x+1$, देईल.
\end{ex}
\begin{quote}\small\textbf{स्रोतदुरुस्ती.} OLFUN-003 — गोठवलेल्या इंग्रजी स्रोताच्या गद्यामध्ये याच उदाहरणात आदानासाठी आधी $n$ आणि लगेच x वापरले आहेत. सूत्र न बदलता मराठी गद्याने x हे एकच नाव सातत्याने वापरले आहे.\end{quote}

\begin{explain}
आपण नुकतीच वेगवेगळ्या \emph{व्याख्या} असलेली $f$ आणि $g$ ही दोन
फलने पाहिली. पण ती \emph{एकच फलन} आहेत. कारण कोणत्याही नैसर्गिक
संख्या~$n$ साठी $f(n) = n+1 = n+2-1 = g(n)$ असते. दुसऱ्या शब्दांत,
$f$ आणि~$g$ यांच्या व्याख्या वेगवेगळ्या समीकरणांनी एकच संगती सांगतात.
म्हणून आपण अप्रत्यक्षपणे फलनांसाठीच्या विस्तारात्मकतेच्या तत्त्वावर
अवलंबून आहोत:
\[
  \text{जर }\forall x\, f(x) = g(x)\text{, तर }f = g
\]
मात्र $f$ आणि~$g$ यांचा प्रांत आणि सहप्रांत एकच असला पाहिजे.
\end{explain}

\begin{ex}
प्रकरणे पाडूनही फलनाची व्याख्या करता येते. उदाहरणार्थ,
$h \colon \Nat \to \Nat$ ची व्याख्या अशी करता येईल:
\[
h(x) =
\begin{cases}
  \frac{x}{2} & \text{जर $x$ सम असेल} \\
  \frac{x+1}{2} & \text{जर $x$ विषम असेल.}
\end{cases}
\]
प्रत्येक नैसर्गिक संख्या सम किंवा विषम असते, म्हणून या फलनाचे प्रदान
नेहमी नैसर्गिक संख्या असेल. प्रकरणे पाडून फलनाची व्याख्या करताना
प्रत्येक शक्य आदान नेमक्या एका प्रकरणात बसले पाहिजे, हे लक्षात ठेवा.
काही वेळा सर्व शक्यता त्या प्रकरणांत येतात आणि कोणतीही शक्यता दोन
प्रकरणांत येत नाही याची सिद्धता करावी लागेल.
\end{ex}








\section{फलनांचे प्रकार}\label{sfr:fun:kin:sec}

\begin{explain}
आपल्याला वारंवार भेटणाऱ्या फलनांच्या काही प्रकारांचे वर्गीकरण करून
घेणे उपयोगी ठरेल.

सुरुवातीला, सहप्रांतातील प्रत्येक सदस्य फलनाचे मूल्य असतो असा गुणधर्म
असलेल्या फलनांचा विचार करू. अशा फलनांना आच्छादक म्हणतात.
त्यांची आकृती \ref{sfr:fun:kin:fig:surjective} प्रमाणे काढता येते.

\begin{figure}
  \olasset{assets/diagrams/surjective.tikz}
  \caption{आच्छादक फलनात सहप्रांतातील प्रत्येक घटक
    हे फलनाचे मूल्य असते.}
  \label{sfr:fun:kin:fig:surjective}
\end{figure}
\end{explain}

\begin{defn}[आच्छादक फलन]
$f \colon A \rightarrow B$ हे फलन \emph{आच्छादक} असते तेव्हा
आणि केवळ तेव्हाच, जेव्हा $B$ हा~$f$ ची व्याप्तीही असतो. म्हणजे,
प्रत्येक $y \in B$ साठी~$f(x) = y$ असणारा किमान एक $x \in A$
असतो. चिन्हांत:
\[
  (\forall y \in B)(\exists x \in A)f(x) = y.
\]
अशा फलनाला $A$ पासून $B$ पर्यंतचे आच्छादन म्हणतो.
\end{defn}

\begin{explain}
$f$ हे आच्छादन आहे असे दाखवायचे असल्यास, $f$ च्या सहप्रांतातील
प्रत्येक वस्तू कोणत्या तरी आदान $x$ साठी $f(x)$ चे मूल्य आहे असे दाखवावे लागते.

कोणत्याही फलनापासून एक आच्छादन \emph{निर्माण होते}, हे लक्षात घ्या.
कारण $f \colon A \to B$ हे फलन दिले असल्यास, $f' \colon A \to \ran{f}$
ची व्याख्या $f'(x) = f(x)$ अशी करू. $\ran{f}$ ची \emph{व्याख्याच}
$\Setabs{f(x) \in B}{x \in A}$ अशी असल्यामुळे हे फलन $f'$ निश्चितच
आच्छादन असते.
\end{explain}

\begin{explain}
कोणतेही फलन प्रत्येक शक्य आदानाशी एकमेव प्रदान जोडते. पण वेगवेगळ्या
आदानांशी कधीच एकच प्रदान न जोडणारी फलनेही असतात. अशा फलनांना
एकास-एक म्हणतात. त्यांची आकृती \ref{sfr:fun:kin:fig:injective} प्रमाणे काढता येते.
\begin{figure}
  \olasset{assets/diagrams/injective.tikz}
  \caption{एकास-एक फलन कधीच दोन वेगवेगळ्या फलसाधकांशी
    एकच मूल्य जोडत नाही.}
  \label{sfr:fun:kin:fig:injective}
\end{figure}
\end{explain}

\begin{defn}[एकास-एक फलन]
$f \colon A \rightarrow B$ हे फलन \emph{एकास-एक} असते तेव्हा
आणि केवळ तेव्हाच, जेव्हा प्रत्येक $y \in B$ साठी~$f(x) = y$ असणारा
जास्तीत जास्त एक $x \in A$ असतो. अशा फलनाला $A$ पासून~$B$ पर्यंतचे
एकास-एक फलन म्हणतो.
\end{defn}

\begin{explain}
$f$ हे एकास-एक फलन आहे असे दाखवायचे असल्यास, $f$ च्या प्रांतातील
कोणत्याही घटक $x$ आणि $y$ साठी $f(x)=f(y)$ असल्यास $x=y$
असते असे दाखवावे लागते.
\end{explain}

\begin{ex}
$f(x) = 1$ ने दिलेले स्थिर फलन $f\colon \Nat \to \Nat$ हे
एकास-एक ही नाही आणि आच्छादक ही नाही.

$f(x) = x$ ने दिलेले एकरूपता फलन $f\colon \Nat \to \Nat$ हे
एकास-एक आणि आच्छादक दोन्ही आहे.

$f(x) = x+1$ ने दिलेले उत्तरवर्ती फलन $f \colon \Nat \to \Nat$
हे एकास-एक आहे, पण आच्छादक नाही.

पुढीलप्रमाणे व्याख्या केलेले फलन $f \colon \Nat \to \Nat$:
\[
  f(x) =
  \begin{cases}
    \frac{x}{2} & \text{जर $x$ सम असेल} \\
    \frac{x+1}{2} & \text{जर $x$ विषम असेल.}
  \end{cases}
\]
हे आच्छादक आहे, पण एकास-एक नाही.
\end{ex}

\begin{explain}
अनेकदा एकास-एक आणि आच्छादक हे दोन्ही गुणधर्म असलेल्या
फलनांचा विचार करायचा असतो. अशा फलनांना एकास-एक व आच्छादक म्हणतो.
ती \ref{sfr:fun:kin:fig:bijective} मधील फलनासारखी दिसतात. एकास-एक आच्छादने ना
कधी कधी \emph{एकास-एक संगती} असेही म्हणतात, कारण ती सहप्रांतातील
घटकांची प्रांतातील घटकांशी एकमेव जोडी लावतात.
\begin{figure}
  \olasset{assets/diagrams/bijective.tikz}
  \caption{एकास-एक व आच्छादक फलन सहप्रांतातील घटकांची प्रांतातील
    घटकांशी एकमेव जोडी लावते.}
  \label{sfr:fun:kin:fig:bijective}
\end{figure}
\end{explain}

\begin{defn}[एकास-एक आच्छादन]
$f \colon A \to B$ हे फलन \emph{एकास-एक व आच्छादक} असते तेव्हा आणि
केवळ तेव्हाच, जेव्हा ते आच्छादक आणि एकास-एक दोन्ही असते.
अशा फलनाला $A$ पासून~$B$ पर्यंतचे (किंवा $A$ आणि~$B$ यांमधील)
एकास-एक आच्छादन म्हणतो.
\end{defn}










\section{संबंध म्हणून फलने}\label{sfr:fun:rel:sec}

\begin{explain}
~$A$ मधील घटक ची~$B$ मधील घटक शी संगती लावणारे
फलन $A$ आणि~$B$ यांच्यामधील एक संबंध ठरवते: $f(x) = y$ असेल
तेव्हा आणि केवळ तेव्हाच $x$ आणि $y$ यांच्यात हा संबंध असतो.
खरे तर, उदाहरणार्थ गणिताच्या पायाभूत घटकांची संख्या कमी करायची असेल,
तर आपण फलन~$f$ आणि हा संबंध, म्हणजे जोड्यांचा एक संच, \emph{एकरूप}
मानू शकतो. मग असा प्रश्न उभा राहतो: कोणते संबंध अशा रीतीने फलने ठरवतात?
\end{explain}

\begin{defn}[फलनाचा आलेख] $f\colon A \to B$ हे फलन असू द्या.
~$f$ चा \emph{आलेख} म्हणजे पुढील व्याख्येने मिळणारा संबंध
$R_f \subseteq A \times B$:
\[
R_f = \Setabs{\tuple{x,y}}{f(x) = y}.
\]
\end{defn}

\begin{explain}
विस्तारात्मकतेमुळे फलनाचा आलेख एकमेव ठरतो. शिवाय, संचांसाठीच्या
विस्तारात्मकतेमुळे फलनांसाठीचे अप्रत्यक्ष विस्तारात्मकता तत्त्व लगेच
समर्थित होते: $f$ आणि~$g$ यांचा प्रांत व सहप्रांत एकच असेल आणि
सर्व मूल्यांवर ती जुळत असतील, तर ती एकच फलन असतात.

तसेच, एखादा संबंध ``फलनरूप'' असेल, तर तो एका फलनाचा आलेख असतो.
\end{explain}

\begin{prop}\label{sfr:fun:rel:prop:graph-function}
$R \subseteq A \times B$ हा संबंध पुढील अटी पूर्ण करतो असे समजा:
\begin{enumerate}
\item $Rxy$ आणि $Rxz$ असल्यास $y = z$ असते; आणि
\item प्रत्येक $x \in A$ साठी $\tuple{x,
y} \in R$ असणारा काही $y \in B$ असतो.
\end{enumerate}
मग $f(x) = y$ तेव्हा आणि केवळ तेव्हाच $Rxy$, अशा व्याख्येने मिळणाऱ्या
फलन $f\colon A \to B$ चा आलेख $R$ असतो.
\end{prop}

\begin{proof}
$Rxy$ असणारा एक $y$ आहे असे समजा. $Rxz$ असणारा दुसरा $z \neq y$
असता, तर~$R$ वरील अट मोडली गेली असती. म्हणून $Rxy$ असणारा एक $y$
असेल, तर हा $y$ एकमेव असतो. त्यामुळे $f$ सुस्पष्टपणे परिभाषित आहे.
उघडच, $R_f = R$.
\end{proof}

\begin{explain}
प्रत्येक फलन $f\colon A \to B$ ला एक आलेख असतो, म्हणजे $f(x) = y$
ने परिभाषित होणारा A आणि B यांच्यातील, $A \times B$ मध्ये
समाविष्ट असलेला एक संबंध असतो. उलट,
\ref{sfr:fun:rel:prop:graph-function} मधील गुणधर्म असलेला प्रत्येक संबंध~$R
\subseteq A \times B$ हा एका फलन~$f \colon A \to B$ चा आलेख असतो.
फलने आणि त्यांचे आलेख यांचा हा निकटचा संबंध असल्यामुळे फलन म्हणजे
त्याचा आलेखच असे आपण मानू शकतो. दुसऱ्या शब्दांत, फलने विशिष्ट
संबंधांशी, म्हणजे विशिष्ट क्रमित घटकसमूहांच्या संचांशी, एकरूप मानता
येतात. पण या ``एकरूप मानण्याचा'' आशय
\ref{sfr:rel:ref:sec} मधल्यासारखाच आहे, हे लक्षात घ्या: तो
फलनांच्या सत्तामीमांसेविषयीचा दावा नाही, तर फलनांना विशिष्ट संच
\emph{मानणे} सोयीचे आहे हे निरीक्षण आहे. ही सोय होण्याचे एक कारण असे की,
संबंधांवर जशा क्रिया केल्या तशाच क्रिया फलनांवर करण्याचा आता आपण
विचार करू शकतो (\ref{sfr:rel:ops:sec} पाहा). विशेषतः:
\end{explain}
\begin{quote}\small\textbf{स्रोतदुरुस्ती.} OLFUN-004 — इंग्रजी स्रोत येथे graph ला relation on A times B म्हणतो; पण त्याच ग्रंथाच्या व्याख्येनुसार त्याचा शब्दशः अर्थ (A times B) च्या वर्गाचा उपसंच होईल. मराठी मजकुरात अचूक अर्थ A आणि B यांच्यातील, A times B मध्ये समाविष्ट संबंध असा दिला आहे.\end{quote}

\begin{defn}\label{sfr:fun:rel:defn:funimage}
$f \colon A \to B$ हे फलन असू द्या आणि $C\subseteq A$ असू द्या.

~$f$ चे~$C$ पर्यंतचे \emph{मर्यादन} म्हणजे सर्व $x \in C$ साठी
$(\funrestrictionto{f}{C})(x) = f(x)$ या व्याख्येने मिळणारे
फलन~$\funrestrictionto{f}{C}\colon C \to B$. दुसऱ्या शब्दांत,
$\funrestrictionto{f}{C} = \Setabs{\tuple{x, y} \in R_f}{x \in C}$.

~$f$ चे~$C$ वरील \emph{उपयोजन} म्हणजे $\funimage{f}{C} =
\Setabs{f(x)}{x \in C}$. यालाच~$f$ अंतर्गत~$C$ ची \emph{प्रतिमा}
असेही म्हणतो.
\end{defn}

\begin{explain}
या व्याख्यांवरून कोणत्याही फलन~$f$ साठी $\ran{f} =
\funimage{f}{\dom{f}}$ असते.
\ref{sfr:rel:ops:sec} मधील
व्याख्या आणि फलनांना संबंधांशी एकरूप मानण्याची आपली पद्धत लक्षात
घेतल्यास या संकल्पना अनुरूप आहेत. मात्र फलनाचे C पर्यंतचे मर्यादन
केवळ आदान निर्देशांक C मध्ये ठेवते; संबंधाचे C वरील मर्यादन
दोन्ही निर्देशांक C मध्ये ठेवते. त्यामुळे त्या अगदी एकच क्रिया नाहीत.
व्युत्क्रम आणि सापेक्ष गुणाकार या इतर दोन क्रियांना थोडे अधिक
स्पष्टीकरण लागते. ते आपण
\ref{sfr:fun:inv:sec} आणि \ref{sfr:fun:cmp:sec} मध्ये देऊ.
\end{explain}
\begin{quote}\small\textbf{स्रोतदुरुस्ती.} OLFUN-005 — फलनाच्या मर्यादनाचे स्पष्ट सूत्र योग्य आहे: आलेखात पहिला निर्देशांक C मध्ये असतो. आधीचे संबंधावरील मर्यादन दोन्ही निर्देशांक C मध्ये ठेवते. इंग्रजी स्रोताचा exactly असा दावा येथे अनुरूप पण भिन्न असा स्पष्ट केला आहे.\end{quote}








\section{फलनांचे व्युत्क्रम}\label{sfr:fun:inv:sec}

\begin{explain}
फलनांचा विचार आपण संगती म्हणून करतो. मग ही संगती ``उलटवता'' येते का,
हा प्रश्न स्वाभाविकपणे पडतो. उदाहरणार्थ, उत्तरवर्ती फलन $f(x) = x + 1$
उलटवता येते, या अर्थाने की $g(y) = y - 1$ हे फलन $f$ ने केलेली
क्रिया ``पूर्ववत करते''.

पण येथे काळजी घेतली पाहिजे. ~$g$ च्या व्याख्येने $\Int \to \Int$
असे फलन ठरत असले, तरी $\Nat \to \Nat$ असे \emph{फलन} ठरत नाही,
कारण $g(0) \notin \Nat$. म्हणून साध्या उदाहरणांतसुद्धा फलन उलटवता
येते का हे तितकेसे उघड नसते; ते प्रांत आणि सहप्रांतावर अवलंबून असू शकते.

फलनाच्या व्युत्क्रमाच्या संकल्पनेने हा विचार अधिक नेमका करता येतो.
\end{explain}

\begin{defn}
सर्व $x \in A$ आणि $y \in B$ साठी $f(g(y)) = y$ आणि $g(f(x)) = x$
असेल, तर $g \colon B \to A$ हे फलन $f \colon A \to B$ या फलनाचा
\emph{व्युत्क्रम} असते.
\end{defn}

$f$ ला व्युत्क्रम~$g$ असेल, तर आपण अनेकदा~$g$ ऐवजी $f^{-1}$ असे लिहितो.

\begin{explain}
आता कोणत्या फलनांना व्युत्क्रम असतात हे ठरवू. $f\colon A \to B$
च्या व्युत्क्रमासाठी एक संभाव्य फलन म्हणजे पुढीलप्रमाणे ``व्याख्या केलेले''
$g\colon B \to A$:
\[
g(y) = \text{$f(x) = y$ असणारा ``तो'' $x$.}
\]
पण ``व्याख्या केलेले'' आणि ``तो'' यांभोवतीची अवतरणचिन्हे सुचवतात की
ही खरे तर व्याख्या नाही. किमान पूर्ण व्यापकतेने ती नेहमी लागू पडणार
नाही. कारण या व्याख्येने फलन ठरवायचे असेल, तर $f(x) = y$ असणारा
एक आणि केवळ एकच~$x$ असला पाहिजे---~$g$ चे प्रदान एकमेव ठरले पाहिजे.
शिवाय ते प्रत्येक $y \in B$ साठी ठरले पाहिजे. $x_1$ आणि $x_2 \in A$
हे $x_1 \neq x_2$ असूनही $f(x_1) = f(x_2)$ असतील, तर
$y = f(x_1) = f(x_2)$ साठी $g(y)$ एकमेव ठरणार नाही.
आणि $f(x) = y$ असणारा~$x$ मुळीच नसेल, तर $g(y)$ मुळीच ठरणार नाही.
दुसऱ्या शब्दांत, $g$ परिभाषित होण्यासाठी $f$ हे एकास-एक आणि
आच्छादक दोन्ही असले पाहिजे.

हे टप्प्याटप्प्याने पाहू. प्रश्नाचे दोन भाग करू: फलन~$f\colon A \to B$
दिले असता, $g(f(x)) = x$ असणारे फलन $g\colon B \to A$ कधी असते?
असे $g$ हे $f$ ने केलेली क्रिया ``पूर्ववत करते'', आणि त्याला~$f$ चा
\emph{डावा व्युत्क्रम} म्हणतात. दुसरे म्हणजे, $f(h(y)) = y$ असणारे
फलन $h\colon B \to A$ कधी असते? अशा $h$ ला~$f$ चा
\emph{उजवा व्युत्क्रम} म्हणतात---$f$ हे~$h$ ने केलेली क्रिया ``पूर्ववत करते''.
\end{explain}

\begin{prop}
A अरिक्त असेल आणि $f\colon A \to B$ हे एकास-एक असेल,
तर~$f$ चा एक
\emph{डावा व्युत्क्रम}~$g\colon B \to A$ असतो, आणि सर्व $x \in A$
साठी $g(f(x)) = x$ असते.
\end{prop}
\begin{quote}\small\textbf{स्रोतदुरुस्ती.} OLFUN-001 — गोठवलेल्या इंग्रजी प्रमेयात A अरिक्त असण्याची अट सुटली आहे. रिक्त A पासून अरिक्त B कडेचे रिक्त फलन एकास-एक असूनही B पासून A कडे फलन नसते. नेमकी सामान्य अट: एकास-एक फलनाला डावा व्युत्क्रम तेव्हा आणि केवळ तेव्हाच असतो, जेव्हा A अरिक्त असतो किंवा B रिक्त असतो. मराठी प्रमेयात साधी पुरेशी A अरिक्त ही अट जोडली आहे.\end{quote}

\begin{proof}
A अरिक्त आहे आणि $f\colon A \to B$ हे एकास-एक आहे असे समजा.
$y \in B$ विचारात घ्या.
$y \in \ran{f}$ असेल, तर $f(x) = y$ असणारा एक $x \in A$ असतो.
$f$ हे एकास-एक असल्यामुळे असा~$x \in A$ एकच असतो. म्हणून
$g(y) = x$ अशी व्याख्या करता येते; म्हणजे $g(y)$ हा $f(x) = y$
असणारा ``तो'' $x \in A$ असतो. $y \notin \ran{f}$ असेल, तर त्याची
संगती कोणत्याही~$a \in A$ शी लावता येते. म्हणून एक $a \in A$
निवडून $g\colon B \to A$ ची व्याख्या अशी करता येते:
\[
g(y) = \begin{cases}
    x & \text{जर $f(x) = y$ असेल}\\
    a & \text{जर $y \notin \ran{f}$ असेल.}
\end{cases}
\]
हे सर्व $y \in B$ साठी परिभाषित आहे, कारण अशा प्रत्येक $y \in \ran{f}$
साठी $f(x) = y$ असणारा नेमका एक $x \in A$ असतो. व्याख्येनुसार,
$y = f(x)$ असल्यास $g(y) = x$, म्हणजे $g(f(x)) = x$.
\end{proof}

\begin{prob}
$f\colon A \to B$ ला डावा व्युत्क्रम~$g$ असल्यास $f$ हे
एकास-एक असते, हे दाखवा.
\end{prob}

\begin{prop}
    $f\colon A \to B$ हे आच्छादक असेल, तर~$f$ चा एक
    \emph{उजवा व्युत्क्रम}~$h\colon B \to A$ असतो, आणि सर्व~$y \in B$
    साठी $f(h(y)) = y$ असते.
\end{prop}

\begin{proof}
$f\colon A \to B$ हे आच्छादक आहे असे समजा. $y \in B$ विचारात घ्या.
~$f$ हे आच्छादक असल्यामुळे $f(x_y) = y$ असणारा एक $x_y \in A$
असतो. म्हणून $h(y) = x_y$ अशी व्याख्या करता येते: प्रत्येक $y \in B$
साठी $f(x) = y$ असणारा काही $x \in A$ आपण निवडतो. ~$f$ हे
आच्छादक असल्यामुळे निवडण्यासाठी नेहमी किमान एक घटक असतो.\footnote{
$f$ हे आच्छादक असल्यामुळे प्रत्येक~$y \in B$ साठी
$\Setabs{x}{f(x) = y}$ हा संच अरिक्त असतो. ~$h$ ची आपली व्याख्या
या प्रत्येक संचातून एकच $x$ निवडण्याची मागणी करते. हे नेहमी करता येते
ही गोष्ट खरे तर उघड नाही---या निवडी करता येतात हे एक स्वयंसिद्धक म्हणून
गृहीत धरले जाते. दुसऱ्या शब्दांत, ही प्रतिज्ञा तथाकथित निवडीचे स्वयंसिद्धक
गृहीत धरते. या प्रश्नाकडे आपण सध्या दुर्लक्ष करू.
मात्र अनेक विशिष्ट उदाहरणांत, जसे $A = \Nat$ असेल किंवा सांत असेल,
किंवा $f$ हे एकास-एक व आच्छादक असेल, तेव्हा निवडीचे स्वयंसिद्धक लागत नाही.
($f$ हे एकास-एक व आच्छादक असण्याच्या विशिष्ट उदाहरणात प्रत्येक $y \in B$
साठी $\Setabs{x}{f(x) = y}$ या संचात नेमका एक घटक असतो,
म्हणून निवड करण्याचा प्रश्नच उरत नाही.)}
व्याख्येनुसार, $x = h(y)$ असल्यास $f(x) = y$ असते; म्हणजे कोणत्याही
$y \in B$ साठी $f(h(y)) = y$.
\end{proof}

\begin{prob}
$f\colon A \to B$ ला उजवा व्युत्क्रम~$h$ असल्यास $f$ हे
आच्छादक असते, हे दाखवा.
\end{prob}

\begin{explain}
  मागील सिद्धतेतील कल्पना एकत्र केल्यावर प्रत्येक एकास-एक आच्छादन ला
  व्युत्क्रम असतो हे मिळते. म्हणजे~$f$ चा डावा आणि उजवा व्युत्क्रम
  असे दोन्ही असणारे एकच फलन असते.
\end{explain}

\begin{prop}\label{sfr:fun:inv:prop:bijection-inverse}
$f\colon A \to B$ हे एकास-एक व आच्छादक असेल, तर एक
फलन~$f^{-1}\colon B \to A$ असते, आणि सर्व $x \in A$ साठी
$f^{-1}(f(x)) = x$ तसेच सर्व $y \in B$ साठी $f(f^{-1}(y)) = y$ असते.
\end{prop}

\begin{proof}
सरावासाठी.
\end{proof}

\begin{prob}
\ref{sfr:fun:inv:prop:bijection-inverse} सिद्ध करा. ~$f^{-1}$ ची
व्याख्या करा, ते फलन आहे हे दाखवा आणि ते~$f$ चा व्युत्क्रम आहे हे
दाखवा. म्हणजे, सर्व $x \in A$ आणि $y \in B$ साठी
$f^{-1}(f(x)) = x$ आणि $f(f^{-1}(y)) = y$ हे दाखवा.
\end{prob}

\begin{explain}
व्युत्क्रम मिळवण्याची थोडी अधिक व्यापक पद्धत आहे. \ref{sfr:fun:kin:sec}
मध्ये आपण पाहिले की प्रत्येक फलन $f$ पासून सर्व $x \in A$ साठी
$f'(x) = f(x)$ असे ठेवून आच्छादन $f' \colon A \to \ran{f}$
मिळते. उघडच, ~$f$ हे एकास-एक असेल, तर~$f'$ हे एकास-एक व आच्छादक
असते. म्हणून \ref{sfr:fun:inv:prop:bijection-inverse} नुसार त्याला एकमेव
व्युत्क्रम असतो. संकेतलेखनात अगदी थोडी सवलत घेऊन आपण कधी कधी
$f'$ च्या व्युत्क्रमाला फक्त ``~$f$ चा व्युत्क्रम'' म्हणतो.
\end{explain}

\begin{prop}\label{sfr:fun:inv:prop:left-right}%
  $f\colon A \to B$ ला डावा व्युत्क्रम~$g$ आणि उजवा व्युत्क्रम~$h$
  असल्यास $h = g$ असते, हे दाखवा.
\end{prop}

\begin{proof}
  सरावासाठी.
\end{proof}

\begin{prob}
  \ref{sfr:fun:inv:prop:left-right} सिद्ध करा.
\end{prob}

\begin{prop}\label{sfr:fun:inv:prop:inverse-unique}
प्रत्येक फलन~$f$ ला जास्तीत जास्त एक व्युत्क्रम असतो.
\end{prop}

\begin{proof}
  $g$ आणि $h$ हे दोन्ही~$f$ चे व्युत्क्रम आहेत असे समजा. मग विशेषतः
  ~$g$ हा~$f$ चा डावा व्युत्क्रम आणि~$h$ हा उजवा व्युत्क्रम असतो.
  \ref{sfr:fun:inv:prop:left-right} नुसार $g = h$.
\end{proof}








\section{फलनांचे संयोजन}\label{sfr:fun:cmp:sec}

\begin{explain}
\ref{sfr:fun:inv:sec} मध्ये आपण पाहिले की
एकास-एक आच्छादन~$f$ चा व्युत्क्रम~$f^{-1}$ हाही एक फलन असतो.
फलनांवरील आणखी एक क्रिया म्हणजे संयोजन: $f$ आणि~$g$ या दोन
फलनांचे संयोजन करून, म्हणजे आधी $f$ आणि मग~$g$ उपयोजून,
एक नवे फलन परिभाषित करता येते. अर्थात, व्याप्ती आणि प्रांत यांचा मेळ
असेल तेव्हाच हे शक्य असते; म्हणजे~$f$ ची व्याप्ती ही~$g$ च्या प्रांताचा
उपसंच असली पाहिजे. फलनांवरील ही क्रिया
\ref{sfr:rel:ops:sec} मधील संबंधांवरील सापेक्ष गुणाकाराच्या क्रियेशी
समांतर आहे.

संयोजनाची कल्पना स्पष्ट करण्यासाठी आकृती उपयोगी पडू शकते.
\ref{sfr:fun:cmp:fig:composition} मध्ये $f \colon A \to B$ आणि $g \colon B \to C$
ही दोन फलने आणि त्यांचे संयोजन~$(\comp{f}{g})$ दाखवले आहे.
$(\comp{f}{g}) \colon A \to C$ हे फलन~$A$ मधील प्रत्येक घटक
शी~$C$ मधील घटक जोडते. $A$ मधील घटक शी~$C$
मधील कोणता घटक जोडायचा हे असे ठरवतो: आदान $x \in A$
दिले असता, प्रथम~$x$ वर फलन $f$ उपयोजा. त्यामुळे काही
$f(x) = y \in B$ हे प्रदान मिळेल. मग~$y$ वर फलन $g$ उपयोजा.
त्यामुळे काही $g(f(x)) = g(y) = z \in C$ हे प्रदान मिळेल.
\begin{figure}
  \olasset[2\olphotowidth]{assets/diagrams/composition.tikz}
  \caption{$f$ आणि~$g$ या दोन फलनांचे संयोजन $g \circ f$.}
  \label{sfr:fun:cmp:fig:composition}
\end{figure}
\end{explain}

\begin{defn}[संयोजन]
$f\colon A \to B$ आणि $g\colon B \to C$ ही फलने असू द्या.
$f$ चे~$g$ बरोबरचे \emph{संयोजन} म्हणजे $\comp{f}{g} \colon A \to C$,
जेथे $(\comp{f}{g})(x) = g(f(x))$.
\end{defn}

\begin{ex}
$f(x) = x + 1$ आणि $g(x) = 2x$ ही फलने विचारात घ्या.
$(\comp{f}{g})(x) = g(f(x))$ असल्यामुळे प्रत्येक आदान~$x$ साठी
प्रथम त्याची उत्तरवर्ती संख्या घेऊन नंतर तिला दोनने गुणले पाहिजे.
म्हणून त्यांचे संयोजन $(\comp{f}{g})(x) = 2(x+1)$ असे मिळते.
\end{ex}

\begin{prob}
$f \colon A \to B$ आणि $g \colon B \to C$ ही दोन्ही फलने
एकास-एक असल्यास $\comp{f}{g}\colon A \to C$ हे
एकास-एक असते, हे दाखवा.
\end{prob}

\begin{prob}
$f \colon A \to B$ आणि $g \colon B \to C$ ही दोन्ही फलने
आच्छादक असल्यास $\comp{f}{g}\colon A \to C$ हे
आच्छादक असते, हे दाखवा.
\end{prob}

\begin{prob}
$f \colon A \to B$ आणि $g \colon B \to C$ असे समजा.
$\comp{f}{g}$ चा आलेख $R_f \mid R_g$ आहे हे दाखवा.
\end{prob}










\section{अंशतः फलने}\label{sfr:fun:par:sec}

\begin{explain}
कधी कधी फलनाच्या व्याख्येत थोडी सवलत देणे उपयोगी ठरते: सर्व शक्य
आदानांसाठी फलनाचे प्रदान परिभाषित असलेच पाहिजे ही अट काढून टाकता
येते. अशा संगतींना \emph{अंशतः फलने} म्हणतात.
\end{explain}

\begin{defn}
\emph{अंशतः फलन} $f \colon A \pto B$ ही अशी संगती असते की ती~$A$
मधील प्रत्येक घटक शी~$B$ मधील जास्तीत जास्त एक घटक
जोडते. $f$ ने $x \in A$ शी~$B$ मधील एखादा घटक जोडला असेल, तर
$f(x)$ \emph{परिभाषित} आहे असे म्हणतो; अन्यथा ते \emph{अपरिभाषित}
असते. $f(x)$ परिभाषित असल्यास $f(x) \fdefined$ असे लिहितो;
अन्यथा $f(x) \fundefined$ असे लिहितो. अंशतः फलन~$f$ चा \emph{प्रांत}
म्हणजे ते परिभाषित असलेल्या~$A$ मधील घटकांचा उपसंच:
$\dom{f} = \Setabs{x \in A}{f(x) \fdefined}$.
\end{defn}

\begin{ex}
प्रत्येक फलन $f\colon A \to B$ हे अंशतः फलनही असते. ~$A$ वरील
सर्व ठिकाणी परिभाषित असलेल्या अंशतः फलनांना---म्हणजे ज्यांना आपण
आतापर्यंत फक्त फलन म्हणत आलो, त्यांना---\emph{सर्वत्र परिभाषित}
फलने असेही म्हणतात.
\end{ex}

\begin{ex}
$f(x) = 1/x$ ने दिलेले अंशतः फलन $f \colon \Real \pto \Real$
हे $x = 0$ साठी अपरिभाषित असते आणि इतर सर्व ठिकाणी परिभाषित असते.
\end{ex}

\begin{prob}
$f\colon A \pto B$ दिले असता, अंशतः फलन $g\colon B \pto A$ ची
व्याख्या अशी करा: कोणत्याही $y \in B$ साठी $f(x) = y$ असणारा एकमेव
$x \in A$ असेल, तर $g(y) = x$; अन्यथा $g(y) \fundefined$.
$f$ हे एकास-एक असल्यास सर्व $x \in \dom{f}$ साठी $g(f(x)) = x$
आणि सर्व $y \in \ran{f}$ साठी $f(g(y)) = y$ असते, हे दाखवा.
\end{prob}

\begin{defn}[अंशतः फलनाचा आलेख]
$f\colon A \pto B$ हे अंशतः फलन असू द्या. ~$f$ चा \emph{आलेख}
म्हणजे पुढील व्याख्येने मिळणारा संबंध $R_f \subseteq A \times B$:
\[
R_f = \Setabs{\tuple{x,y}}{f(x) = y}.
\]
\end{defn}

\begin{prop}
$R \subseteq A \times B$ या संबंधाचा असा गुणधर्म आहे असे समजा की
$Rxy$ आणि $Rxy'$ असल्यास $y = y'$ असते. मग $R$ हा पुढीलप्रमाणे
व्याख्या केलेल्या अंशतः फलन $f\colon A \pto B$ चा आलेख असतो:
$Rxy$ असणारा एक $y$ असेल, तर $f(x) = y$; अन्यथा $f(x) \fundefined$.
शिवाय $R$ हा \emph{प्रत्येक आदानाशी किमान एक प्रदान जोडणारा} असेल,
म्हणजे प्रत्येक $x \in A$ साठी $Rxy$ असणारा एक $y \in B$ असेल,
तर $f$ हे सर्वत्र परिभाषित असते.
\end{prop}

\begin{proof}
$Rxy$ असणारा एक $y$ आहे असे समजा. $Rxy'$ असणारा दुसरा
$y' \neq y$ असता, तर $R$ वरील अट मोडली गेली असती. म्हणून
$Rxy$ असणारा एक $y$ असेल, तर तो $y$ एकमेव असतो. त्यामुळे $f$
सुस्पष्टपणे परिभाषित आहे. उघडच $R_f = R$, आणि~$R$ प्रत्येक आदानाशी
किमान एक प्रदान जोडत असेल, तर $f$ हे सर्वत्र परिभाषित असते.
\end{proof}



\chapter{संचांचे आकारमान}\label{sfr:siz::chap}






\section{परिचय}\label{sfr:siz:int:sec}

1870 च्या दशकात गेओर्क कँटर यांनी संच सिद्धांत
विकसित केला, तेव्हा अनंत समूहाची कल्पना---मध्ययुगीन विचारवंतांनी
म्हटल्याप्रमाणे प्रत्यक्ष अस्तित्वातील अनंतता---स्वीकारार्ह करणे हे
त्यांच्या उद्दिष्टांपैकी एक होते. वेगवेगळ्या संचांच्या \emph{आकारमानाचा}
त्यांनी केलेला विचार हा त्यातील महत्त्वाचा भाग होता. $a$, $b$ आणि $c$
हे सर्व भिन्न असतील, तर $\{a, b, c\}$ हा संच $\{a, b\}$ पेक्षा
\emph{मोठा} आहे असे सहज वाटते. पण अनंत संचांचे काय? ते सर्व
एकमेकांएवढेच मोठे असतात का? तसे नसते, हे दिसून येते.

येथील पहिली महत्त्वाची कल्पना म्हणजे प्रगणन. कोणत्याही सांत संचातील
सर्व घटक लिहून त्याची यादी देता येते. काही अनंत संचांतील
सर्व घटक चीही यादी देता येते, जर यादी स्वतः अनंत असण्याची
मुभा दिली, तर. अशा संचांना गणनीय म्हणतात. काही अनंत संच
गणनीय नसतात, हा कँटर यांचा आश्चर्यकारक निष्कर्ष होता.
या प्रकरणाच्या अखेरीस तो आपल्याला पूर्णपणे समजेल.









\section{प्रगणने आणि गणनीय संच}\label{sfr:siz:enm:sec}

\begin{editorial}
  या विभागात संचांच्या प्रगणनांची चर्चा आहे. त्यांची व्याख्या
  $\PosInt$ पासूनच्या आच्छादनांनी केली आहे. गणिताची फारशी पार्श्वभूमी
  नसलेल्या वाचकांसाठी विषय हळूहळू मांडला आहे. पर्यायी, अधिक संक्षिप्त
  आवृत्ती \ref{sfr:siz:enm-alt:sec} मध्ये दिली आहे. तेथे प्रगणनांची वेगळी
  व्याख्या केली आहे: $\Nat$ किंवा त्याच्या आरंभीच्या खंडाशी
  एकास-एक आच्छादन असणे.
\end{editorial}

\begin{explain}
संचांचे घटक यादीत देऊन आपण आधीच संचांची उदाहरणे दिली आहेत.
संच अनंत असला, तरी त्याच्या घटक ची यादी कशी आणि कधी
देता येते, याची आता अधिक व्यापकपणे चर्चा करू.
\end{explain}

\begin{defn}[प्रगणनाची अनौपचारिक व्याख्या]
अनौपचारिकपणे, संच~$A$ चे \emph{प्रगणन} म्हणजे~$A$ मधील
घटक ची अशी यादी (कदाचित अनंत) की $A$ मधील प्रत्येक
घटक त्या यादीत कोणत्या तरी सांत क्रमांकाच्या स्थानी येतो.
$A$ चे प्रगणन असेल, तर $A$ हा \emph{गणनीय} आहे असे म्हणतात.
\end{defn}

\begin{explain}
प्रगणनांविषयी काही मुद्दे:
\begin{enumerate}
\item ज्या यादीला सुरुवात असते आणि जिच्यात पहिल्या घटकाव्यतिरिक्त
  प्रत्येक घटक च्या लगेच आधी एकच घटक असतो, अशीच यादी
  आपण प्रगणन मानतो. दुसऱ्या शब्दांत, यादीच्या पहिल्या घटक
  आणि इतर कोणत्याही घटक यांच्यामध्ये केवळ सांत संख्येने घटक
  असतात. विशेषतः, प्रगणनातील प्रत्येक घटक चे स्थान सांत
  क्रमांकाचे असते: पहिल्या घटक चे स्थान~$1$, दुसऱ्याचे~$2$, इत्यादी.
\item एकाच संच~$A$ ची वेगवेगळी प्रगणने असू शकतात, ज्यांत
  घटक येण्याचा क्रम वेगळा असतो: $4$, $1$, $25$, $16$,~$9$
  हे पहिल्या पाच वर्गसंख्यांच्या संचाचे प्रगणन आहे, तसेच
  $1$, $4$, $9$, $16$,~$25$ हेही आहे.
\item पुनरुक्ती असलेली प्रगणनेही प्रगणनेच असतात: $1$, $1$, $2$,
  $2$, $3$, $3$,~\dots{} हे त्याच संचाचे प्रगणन आहे, ज्याचे
  प्रगणन $1$, $2$, $3$,~\dots{} हे आहे.
\item विशिष्ट प्रगणन सांगताना क्रम आणि पुनरुक्ती यांना महत्त्व
  \emph{असते}: धन पूर्णांकांचे प्रगणन $1$, $2$, $3$, $1$, \dots{}
  अशा सुरुवातीने करता येते; पण नेहमीच्या $1$, $2$, $3$, $4$,~\dots
  या पद्धतीने प्रगणन केले, तर त्यातील नमुना अधिक सहज दिसतो.
\item प्रगणनाला सुरुवात असली पाहिजे: \dots, $3$, $2$, $1$ हे
  धन पूर्णांकांचे प्रगणन नाही, कारण त्याला पहिला घटक नाही.
  अनौपचारिक व्याख्येवरून हे कसे येते ते पाहण्यासाठी स्वतःला विचारा:
  ``या यादीत 76 ही संख्या कोणत्या स्थानी येते?''
\item पुढील यादी धन पूर्णांकांचे प्रगणन नाही:
  $1$, $3$, $5$, \dots, $2$, $4$, $6$, \dots\@ अडचण अशी की
  सम संख्या सांत क्रमांकांच्या स्थानांवर येण्याऐवजी
  $\infty + 1$, $\infty + 2$, $\infty + 3$ या स्थानांवर येतात.
\item रिक्त संच गणनीय आहे: रिक्त यादी हे त्याचे प्रगणन आहे!{}
\end{enumerate}
\end{explain}

\begin{prop}
  $A$ चे प्रगणन असेल, तर त्याचे पुनरुक्ती नसलेले प्रगणनही असते.
\end{prop}

\begin{proof}
  $A$ चे $x_1$, $x_2$, \dots{} असे प्रगणन आहे, आणि प्रत्येक
  $x_i$ हा~$A$ चा घटक आहे असे समजा. प्रगणनातून पुन्हा
  आलेले घटक काढून टाकून त्यातील पुनरुक्ती काढता येते.
  उदाहरणार्थ, पुढीलप्रमाणे नवे प्रगणन मिळवता येते: $x_i$ हा~$A$
  मधील घटक असून $x_1$, \dots, $x_{i-1}$ यांमध्ये नसेल,
  तर त्याला यादीत लिहायचे; आणि $x_i$ हा आधीच $x_1$, \dots,~$x_{i-1}$
  यांमध्ये असेल, तर त्याला यादीतून काढून टाकायचे.
\end{proof}

मागील युक्तिवाद दाखवतो की प्रगणने आणि गणनीय संच नीट
समजून घेण्यासाठी, तसेच त्यांविषयी गोष्टी सिद्ध करण्यासाठी,
अधिक नेमकी व्याख्या हवी. ती पुढे दिली आहे.

\begin{defn}[प्रगणनाची औपचारिक व्याख्या]
$A \neq \emptyset$ या संचाचे \emph{प्रगणन} म्हणजे कोणतेही
आच्छादक फलन $f \colon \PosInt \to A$.
\end{defn}

\begin{explain}
औपचारिक व्याख्या आणि कदाचित अनंत असणाऱ्या यादीचा वापर करणारी
अनौपचारिक व्याख्या समतुल्य आहेत याची खात्री करू. प्रथम, $\PosInt$
पासून संच~$A$ पर्यंतचे कोणतेही आच्छादक फलन~$A$ चे प्रगणन करते.
अशा फलनाने वरील अनौपचारिक अर्थाचे प्रगणन ठरते: $f(1)$, $f(2)$,
$f(3)$, \dots{} ही यादी. $f$ हे आच्छादक असल्यामुळे~$A$
मधील प्रत्येक घटक हा कोणत्या तरी~$n \in \PosInt$ साठीचे
~$f(n)$ चे मूल्य असतो. म्हणून $A$ मधील प्रत्येक घटक यादीत
सांत क्रमांकाच्या स्थानी येतो. फलन एकास-एक असेलच असे नसल्याने
यादीत पुनरुक्ती असू शकते; पण वर म्हटल्याप्रमाणे ती चालते.

उलट,~$A$ मधील सर्व घटक चे प्रगणन करणारी यादी दिली असल्यास
आच्छादक फलन $f\colon \PosInt \to A$ असे परिभाषित करता येते:
$f(n)$ म्हणजे यादीतील $n$ व्या स्थानी असलेला घटक; आणि
$n$ व्या स्थानी घटक नसेल, तर यादीचा शेवटचा घटक.
$A$ रिक्त असतो आणि म्हणून यादीही रिक्त असते, फक्त या एकाच बाबतीत
या पद्धतीने आच्छादक फलन मिळत नाही. म्हणून प्रत्येक अरिक्त
यादी एक आच्छादक फलन $f\colon \PosInt \to A$ ठरवते.
\end{explain}

\begin{defn}
  \label{sfr:siz:enm:defn:enumerable}
  संच~$A$ हा गणनीय असतो तेव्हा आणि केवळ तेव्हाच,
  जेव्हा तो रिक्त असतो किंवा त्याचे प्रगणन असते.
\end{defn}

\begin{ex}
धन पूर्णांकांचे ($\PosInt$) प्रगणन करणारे फलन म्हणजे फक्त
$f(n) = n$ ने दिलेले एकरूपता फलन. नैसर्गिक संख्या $\Nat$ यांचे
प्रगणन करणारे फलन म्हणजे $g(n) = n - 1$.
\end{ex}

\begin{ex}
$f\colon \PosInt \to \PosInt$ आणि $g \colon \PosInt \to \PosInt$
ही फलने अशी दिली आहेत:
\begin{align*}
f(n) & = 2n \text{ आणि}\\
g(n) & = 2n - 1
\end{align*}
ती अनुक्रमे सम धन पूर्णांकांचे आणि विषम धन पूर्णांकांचे प्रगणन करतात.
मात्र यांपैकी कोणतेही फलन $\PosInt$ चे प्रगणन नाही, कारण कोणतेही
आच्छादक नाही.
\end{ex}

\begin{prob}
धन वर्गसंख्या $1$, $4$, $9$, $16$, \dots{} यांचे प्रगणन परिभाषित करा.
\end{prob}

\begin{ex}
$f(n) = (-1)^{n} \lceil \frac{(n-1)}{2}\rceil$ हे फलन पूर्णांकांचा
संच~$\Int$ प्रगणित करते. येथे $\lceil x \rceil$ हे
\emph{ऊर्ध्व पूर्णांक फलन} दर्शवते; ते $x$ ला त्याच्यापेक्षा कमी
नसलेल्या सर्वांत जवळच्या पूर्णांकाकडे नेते. धन आणि ऋण पूर्णांकांमध्ये
आलटूनपालटून ``उड्या मारून'' $f$ हे $\Int$ ची मूल्ये कशी निर्माण
करते ते पाहा:
\[
\begin{array}{c c c c c c c c}
f(1) & f(2) & f(3) & f(4) & f(5) & f(6) & f(7) & \dots \\ \\
- \lceil \tfrac{0}{2} \rceil & \lceil \tfrac{1}{2}\rceil & - \lceil \tfrac{2}{2} \rceil & \lceil \tfrac{3}{2} \rceil & - \lceil \tfrac{4}{2} \rceil  & \lceil \tfrac{5}{2}
\rceil & - \lceil \tfrac{6}{2} \rceil & \dots \\ \\
0 & 1 & -1 & 2 & -2 & 3 & -3 & \dots
\end{array}
\]
\begin{quote}\small\textbf{स्रोतदुरुस्ती.} OLSIZ-001: गोठवलेल्या इंग्रजी तक्त्याच्या शेवटच्या ओळीत f(7) साठी अपेक्षित −3 हे मूल्य सुटले आहे; सूत्र आणि वरील दोन ओळी ते मूल्य निश्चित करतात. मराठी तक्त्यात −3 समाविष्ट करून ही उणीव दुरुस्त केली आहे.\end{quote}
$f$ ची व्याख्या पुढीलप्रमाणे प्रकरणे पाडून केली आहे असेही मानता येते:
\[
f(n) = \begin{cases}
  0 & \text{जर $n = 1$ असेल}\\
  n/2 & \text{जर $n$ सम असेल}\\
  -(n-1)/2 & \text{जर $n$ विषम आणि $>1$ असेल}
  \end{cases}
\]
\end{ex}

\begin{prob}
  $A$ आणि $B$ हे गणनीय असतील, तर $A \cup B$ हाही
  गणनीय असतो, हे दाखवा. यासाठी आच्छादक फलने
  $f\colon \PosInt \to A$ आणि $g\colon \PosInt \to B$ आहेत असे
  समजा. एक आच्छादक फलन~$h\colon \PosInt \to A \cup B$
  परिभाषित करा आणि ते आच्छादक आहे हे सिद्ध करा. $A$ किंवा
  ~$B = \emptyset$ असण्याच्याही बाबतींचा विचार करा.
\end{prob}

\begin{prob}
  $B \subseteq A$ आणि $A$ हा गणनीय असेल, तर~$B$ हाही
  गणनीय असतो, हे दाखवा. यासाठी एक आच्छादक फलन
  $f\colon \PosInt \to A$ आहे असे समजा. एक आच्छादक
  फलन~$g\colon \PosInt \to B$ परिभाषित करा आणि ते
  आच्छादक आहे हे सिद्ध करा. $B = \emptyset$ असल्यास काय होते?
\end{prob}

\begin{prob}
  $A_1$, $A_2$, \dots, $A_n$ हे सर्व गणनीय असतील, तर
  $A_1 \cup \dots \cup A_n$ हाही गणनीय असतो, हे
  $n$ वर विगमन करून दाखवा. $A$ आणि~$B$ हे दोन संच
  गणनीय असल्यास~$A \cup B$ हाही गणनीय असतो
  हे तथ्य तुम्ही गृहीत धरू शकता.
\end{prob}

संचातील घटक ची यादी करताना $1$ व्या घटक पासून
मोजायला सुरुवात करणे कदाचित अधिक स्वाभाविक असले, तरी गणितज्ञांना
मोजणीसाठी नैसर्गिक संख्या~$\Nat$ वापरायला आवडतात. ते यादीतील
$0$ व्या, $1$ व्या, $2$ व्या, अशा घटक विषयी बोलतात.
त्याला अनुसरून प्रगणनाची व्याख्या $\Nat$ पासून~$A$ पर्यंतचे
आच्छादक फलन अशी करता येते. अर्थात, या दोन व्याख्या समतुल्य आहेत.

\begin{prop}\label{sfr:siz:enm:prop:enum-shift}
  आच्छादन $f\colon \PosInt \to A$ असते तेव्हा आणि
  केवळ तेव्हाच, जेव्हा आच्छादन $g\colon \Nat \to A$ असते.
\end{prop}

\begin{proof}
  आच्छादन $f\colon \PosInt \to A$ दिले असता, सर्व
  $n \in \Nat$ साठी $g(n) = f(n+1)$ अशी व्याख्या करता येते.
  $g\colon \Nat \to A$ हे आच्छादक आहे हे सहज दिसते.
  उलट, आच्छादन $g\colon \Nat \to A$ दिले असता,
  $f(n) = g(n-1)$ अशी व्याख्या करा.
\end{proof}

यावरून पुढील निष्कर्ष मिळतो:

\begin{cor}\label{sfr:siz:enm:cor:enum-nat}
संच $A$ हा गणनीय असतो तेव्हा आणि केवळ तेव्हाच, जेव्हा
तो रिक्त असतो किंवा आच्छादक फलन $f\colon \Nat \to A$ असते.
\end{cor}

संच~$A$ मधील घटक च्या यादीतून पुनरुक्ती नसलेली यादी
मिळवता येते याची वर चर्चा केली. प्रगणनांबाबतही हे खरे आहे,
पण ते काटेकोरपणे मांडणे आणि सिद्ध करणे थोडे कठीण आहे. कोणतेही
फलन $f\colon \PosInt \to A$ हे सर्व $n \in \PosInt$ साठी
परिभाषित असले पाहिजे. ~$A$ मध्ये फक्त सांत संख्येने घटक
असतील, तर~$\PosInt$ मधील अनंत संख्येने असलेल्या घटक
वर परिभाषित असणारे आणि~$A$ मधील सर्व घटक मूल्ये म्हणून
घेणारे, पण एकच मूल्य दोनदा कधीच न घेणारे फलन असू शकत नाही.
अशा बाबतीत, म्हणजे पुनरुक्ती नसलेली यादी सांत असते तेव्हा,~$f$
साठी वेगळा प्रांत निवडला पाहिजे; त्यात केवळ सांत संख्येने
घटक असले पाहिजेत. पुनरुक्ती नसणे म्हणजे $f$ हे
एकास-एक असले पाहिजे. ते आच्छादक ही असल्यामुळे
आपण कोणता तरी सांत संच $\{1, \dots, n\}$ किंवा $\PosInt$
आणि~$A$ यांच्यामधील एकास-एक आच्छादन शोधत आहोत.

\begin{prop}\label{sfr:siz:enm:prop:enum-bij}
$f\colon \PosInt \to A$ हे आच्छादक असेल (म्हणजे~$A$
चे प्रगणन असेल), तर एकास-एक आच्छादन $g\colon Z \to A$ असते,
जेथे $Z$ हा~$\PosInt$ किंवा कोणत्या तरी~$n \in \PosInt$
साठी $\{1, \dots, n\}$ असतो.
\end{prop}

\begin{proof}
  फलन $g$ ची व्याख्या आपण पुनरावर्ती रीतीने करतो: $g(1) = f(1)$
  असे ठेवा. $g(i)$ आधीच परिभाषित झाले असेल, तर $f(1)$, $f(2)$,
  \dots{} यांपैकी $g(1)$, \dots, $g(i)$ यांत आधीच न आलेले
  पहिले मूल्य, असे मूल्य असेल तर, $g(i+1)$ म्हणून घ्या.
  $A$ मध्ये नेमके $n$ घटक असतील, तर $g(1)$, \dots, $g(n)$
  ही सर्व मूल्ये परिभाषित असतात, आणि आपल्याला
  $g\colon \{1, \dots, n\} \to A$ असे फलन मिळते.
  $A$ मध्ये अनंत संख्येने घटक असतील, तर कोणत्याही $i$
  साठी $f(1)$, $f(2)$, \dots{} या प्रगणनात~$A$ मधील असा
  घटक असलाच पाहिजे, जो $g(1)$, \dots, $g(i)$ यांत
  आधीच आलेला नाही. अशा बाबतीत आपण $g\colon \PosInt \to A$
  हे फलन परिभाषित केले आहे.

  फलन $g$ हे आच्छादक आहे, कारण~$A$ मधील कोणताही घटक
  $f(1)$, $f(2)$, \dots{} यांमध्ये येतो ($f$ हे आच्छादक
  असल्यामुळे), आणि म्हणून कधी ना कधी तो कोणत्या तरी~$i$
  साठीचे~$g(i)$ चे मूल्य होईल. ते एकास-एक ही आहे:
  $g(j) = g(i)$ असणारे $j < i$ असते, तर $g(i)$ हे मूल्य
  आधीच $g(1)$, \dots, $g(i-1)$ यांमध्ये आलेले असते;
  पण हे~$g$ च्या आपल्या व्याख्येविरुद्ध आहे.
\end{proof}

\begin{cor}\label{sfr:siz:enm:cor:enum-nat-bij}
संच $A$ हा गणनीय असतो तेव्हा आणि केवळ तेव्हाच, जेव्हा
तो रिक्त असतो किंवा एकास-एक आच्छादन $f\colon N \to A$ असते,
जेथे $N = \Nat$ किंवा कोणत्या तरी $n \in \Nat$ साठी
$N = \{0, \dots, n\}$ असते.
\end{cor}

\begin{proof}
$A$ हा गणनीय असतो तेव्हा आणि केवळ तेव्हाच, जेव्हा $A$
रिक्त असतो किंवा आच्छादक $f\colon \PosInt \to A$ असते.
\ref{sfr:siz:enm:prop:enum-bij} नुसार दुसरी अट तेव्हा आणि केवळ तेव्हाच
खरी असते, जेव्हा एकास-एक व आच्छादक फलन~$f\colon Z \to A$ असते,
जेथे $Z = \PosInt$ किंवा कोणत्या तरी $n \in \PosInt$ साठी
$Z = \{1, \dots, n\}$ असते. \ref{sfr:siz:enm:prop:enum-shift} च्या
सिद्धतेतील युक्तिवादानुसार, हे तेव्हा आणि केवळ तेव्हाच शक्य होते,
जेव्हा एकास-एक आच्छादन $g\colon N \to A$ असते, जेथे
$N = \Nat$ किंवा $N = \{0, \dots, n-1\}$ असते.
\begin{quote}\small\textbf{स्रोतदुरुस्ती.} MRSIZ-001: विधानातील सांत प्रांत 0 ते n आणि सिद्धतेतील 0 ते n−1 हे वेगवेगळ्या चलनामांनी त्याच सांत आरंभीच्या खंडांच्या कुटुंबाचे वर्णन करतात; येथे स्रोताचे दोन्ही निर्देशांक जसेच्या तसे ठेवले आहेत.\end{quote}
\end{proof}

\begin{prob}
  \ref{sfr:siz:enm:defn:enumerable} नुसार संच $A$ हा गणनीय
  असतो तेव्हा आणि केवळ तेव्हाच, जेव्हा $A = \emptyset$ किंवा
  आच्छादक $f\colon \PosInt \to A$ असते.
  ``गणनीय संच'' याची नेमकी व्याख्या पुढीलप्रमाणेही करता येते:
  संच गणनीय असतो तेव्हा आणि केवळ तेव्हाच, जेव्हा एकास-एक
  फलन $g\colon A \to \PosInt$ असते. या व्याख्या समतुल्य आहेत
  हे दाखवा. म्हणजे, एकास-एक फलन $g\colon A \to \PosInt$
  असते तेव्हा आणि केवळ तेव्हाच, जेव्हा $A = \emptyset$ किंवा
  आच्छादक $f\colon \PosInt \to A$ असते, हे दाखवा.
\end{prob}









\section{कँटर यांची नागमोडी पद्धत}\label{sfr:siz:zigzag:sec}

\begin{explain}
काही ``सोप्या'' प्रगणनांचा आपण आधीच विचार केला आहे. आता जरा कठीण
उदाहरण पाहू. नैसर्गिक संख्यांच्या जोड्यांचा संच विचारात घ्या
; त्याची व्याख्या आपण
\ref{sfr:set:pai:sec} मध्ये पुढीलप्रमाणे केली:
\[
\Nat \times \Nat = \Setabs{\tuple{n,m}}{n,m \in \Nat}
\]
या क्रमित जोड्या आपण पुढीलप्रमाणे एका \emph{सरणीत} मांडू शकतो:
\[
\begin{array}{ c | c | c | c | c | c}
& \mathbf 0 & \mathbf 1 & \mathbf 2 & \mathbf 3 & \dots \\
\hline
\mathbf 0 & \tuple{0,0} & \tuple{0,1} & \tuple{0,2} & \tuple{0,3} & \dots \\
\hline
\mathbf 1 & \tuple{1,0} & \tuple{1,1} & \tuple{1,2} & \tuple{1,3} & \dots \\
\hline
\mathbf 2 & \tuple{2,0} & \tuple{2,1} & \tuple{2,2} & \tuple{2,3} & \dots \\
\hline
\mathbf 3 & \tuple{3,0} & \tuple{3,1} & \tuple{3,2} & \tuple{3,3} & \dots \\
\hline
\vdots & \vdots & \vdots & \vdots & \vdots & \ddots\\
\end{array}
\]
$\Nat \times \Nat$ मधील प्रत्येक क्रमित जोडी या सरणीत नेमकी एकदाच
येईल, हे स्पष्ट आहे. विशेषतः, $\tuple{n,m}$ ही जोडी $n$ व्या ओळीत
आणि $m$ व्या स्तंभात येईल. पण अशा सरणीतील घटकांची ``एकमितीय'' यादी
कशी करायची? पुढील सरणीतील नमुना ती करण्याचा एक मार्ग दाखवतो
(अर्थातच इतर अनेक मार्गही आहेत):
\[
\begin{array}{ c | c | c | c | c | c | c}
& \mathbf 0 & \mathbf 1 & \mathbf 2 & \mathbf 3 & \mathbf 4 &\dots \\
\hline
\mathbf 0 & 0  & 1& 3 & 6& 10 &\ldots \\
\hline
\mathbf 1 &2 & 4& 7 & 11 & \dots &\ldots \\
\hline
\mathbf 2 & 5 & 8 & 12 & \ldots & \dots&\ldots \\
\hline
\mathbf 3 & 9 & 13 & \ldots & \ldots & \dots & \ldots \\
\hline
\mathbf 4 & 14 & \ldots & \ldots & \ldots & \dots & \ldots \\
\hline
\vdots & \vdots & \vdots & \vdots & \vdots&\ldots & \ddots\\
\end{array}
\]\noindent
या नमुन्याला \emph{कँटर यांची नागमोडी पद्धत} म्हणतात. तो
$\Nat \times \Nat$ चे पुढीलप्रमाणे प्रगणन करतो:
\[
\tuple{0,0}, \tuple{0,1}, \tuple{1,0}, \tuple{0,2}, \tuple{1,1},
\tuple{2,0}, \tuple{0,3}, \tuple{1,2}, \tuple{2,1}, \tuple{3,0}, \dots
\]
यावरून पुढील विधान सिद्ध होते:
\end{explain}

\begin{prop}\label{sfr:siz:zigzag:natsquaredenumerable}
$\Nat \times \Nat$ हा गणनीय आहे.
\end{prop}

\begin{proof}
प्रत्येक $k \in \Nat$ ला कँटर यांच्या नागमोडी सरणीतील ज्या $n$ व्या
ओळीच्या आणि $m$ व्या स्तंभाच्या जागेचे मूल्य $k$ आहे, त्या
$\tuple{n,m} \in \Nat \times \Nat$ या क्रमित घटकाशी जोडणारे
$f \colon \Nat \to \Nat\times\Nat$ असे फलन घ्या.
\end{proof}

\begin{explain}
ही पद्धत अधिक व्यापक बाबतीतही सहज लागू होते. उदाहरणार्थ, नैसर्गिक
संख्यांच्या क्रमित तीन-घटकसमूहांचा पुढील संच प्रगणित करण्यासाठी ती
वापरता येते:
\[
\Nat \times \Nat \times \Nat = \Setabs{\tuple{n,m,k}}{n,m,k \in \Nat}
\]
$\Nat \times \Nat \times \Nat$ हा $\Nat \times \Nat$ आणि $\Nat$
यांचा कार्तीय गुणाकार आहे असे आपण मानतो; म्हणजे,
\[
\Nat^3 = (\Nat \times \Nat) \times \Nat =
\Setabs{\tuple{\tuple{n,m},k}}{n, m, k
  \in \Nat }
\]
म्हणून एका अक्षाला $\Nat$ चे प्रगणन आणि दुसऱ्या अक्षाला $\Nat^2$
चे प्रगणन अशी नावे देऊन, एका सरणीच्या साहाय्याने $\Nat^3$ चे
प्रगणन करता येते:
\[
\begin{array}{ c | c | c | c | c | c}
& \mathbf 0 & \mathbf 1 & \mathbf 2 & \mathbf 3 & \dots \\
\hline
\mathbf{\tuple{0,0}} & \tuple{0,0,0} & \tuple{0,0,1} & \tuple{0,0,2} & \tuple{0,0,3} & \dots \\
\hline
\mathbf{\tuple{0,1}} & \tuple{0,1,0} & \tuple{0,1,1} & \tuple{0,1,2} & \tuple{0,1,3} & \dots \\
\hline
\mathbf{\tuple{1,0}} & \tuple{1,0,0} & \tuple{1,0,1} & \tuple{1,0,2} & \tuple{1,0,3} & \dots \\
\hline
\mathbf{\tuple{0,2}} & \tuple{0,2,0} & \tuple{0,2,1} & \tuple{0,2,2} & \tuple{0,2,3} & \dots\\
\hline
\vdots & \vdots & \vdots & \vdots & \vdots & \ddots \\
\end{array}
\]
अशा रीतीने कँटर यांच्या नागमोडी पद्धतीसारखी पद्धत वापरून $\Nat^3$
चे प्रगणन मिळते. ही प्रक्रिया पुढे चालू ठेवून, प्रत्येक नैसर्गिक संख्या
$n$ साठी $\Nat^n$ चे प्रगणन मिळवता येते. म्हणून पुढील विधान मिळते:
\end{explain}

\begin{prop}
प्रत्येक $n \in \Nat$ साठी $\Nat^n$ हा गणनीय आहे.
\end{prop}

\begin{prob}\label{sfr:siz:zigzag:prob:posint-n}
प्रत्येक $n \in \Nat$ साठी $(\PosInt)^n$ हा गणनीय आहे, हे दाखवा.
\end{prob}

\begin{prob}\label{sfr:siz:zigzag:prob:posint-star}
  $(\PosInt)^*$ हा गणनीय आहे, हे दाखवा. तुम्ही \ref{sfr:siz:zigzag:prob:posint-n} गृहीत धरू शकता.
\end{prob}








\section{जोडीकरण फलने आणि संकेतांक}\label{sfr:siz:pai:sec}

\begin{explain}
कँटर यांच्या नागमोडी पद्धतीमुळे $\Nat^n$ ची गणनीयता चित्रातून
स्पष्ट दिसते. पण $\Nat^2$ दाखवणाऱ्या आपल्या सरणीवर लक्ष केंद्रित करू.
सरणीतील नागमोडी रेषेचा मागोवा घेऊन स्थाने मोजल्यास,
$\tuple{1,2}$ ही जोडी~$7$ या संख्येशी जोडली आहे हे तपासता येते.
पण ही संख्या अधिक थेट मोजता आली, तर सोयीचे होईल. म्हणजे, नागमोडी
प्रगणनाचे \emph{व्युत्क्रम} $g\colon \Nat^2 \to \Nat$ हाताशी असावे,
आणि त्यासाठी
\[
g(\tuple{0,0}) = 0, \;
g(\tuple{0,1}) = 1, \;
g(\tuple{1,0}) = 2, \; \dots,
g(\tuple{1,2}) = 7, \; \dots
\]
असे असावे. त्यामुळे $\tuple{n, m}$ ही जोडी आपल्या प्रगणनात नेमकी
कोणत्या स्थानी येईल ते मोजता येईल.

खरे तर दोन निरीक्षणांच्या साहाय्याने $g$ ची थेट व्याख्या करता येते.
पहिले निरीक्षण: $n$ व्या ओळीतील आणि $m$ व्या स्तंभातील मूल्य~$v$
असेल, तर $(n+1)$ व्या ओळीतील आणि $(m-1)$ व्या स्तंभातील मूल्य
$v + 1$ असते. दुसरे निरीक्षण: आपल्या प्रगणनाची पहिली ओळ $0$, $1$,
$3$, $6$, इत्यादींपासून सुरू होणाऱ्या त्रिकोणी संख्यांची आहे.
$k$ वी त्रिकोणी संख्या म्हणजे $< k$ ही अट पूर्ण करणाऱ्या नैसर्गिक संख्यांची
बेरीज; ती $k(k+1)/2$ अशी मोजता येते. ही दोन निरीक्षणे एकत्र करून
पुढील फलन विचारात घ्या:
\[
  g(n,m) = \frac{(n+m+1)(n+m)}{2} + n
\]
$g(\tuple{n, m})$ ऐवजी आपण बहुधा फक्त $g(n, m)$ असे लिहितो, कारण
ते पाहायला सोपे असते. या सूत्रानुसार आधी $(n+m)^\text{वी}$ त्रिकोणी
संख्या काढायची आणि नंतर तिच्यात $n$ मिळवायचा. त्यामुळे सरणी नेमकी
आपल्याला हवी तशी भरते. विशेषतः, $\tuple{1, 2}$ ही जोडी
$\frac{4 \times 3}{2} + 1 = 7$ कडे पाठवली जाते.

हे फलन $g$ म्हणजे जोड्यांच्या संचाच्या एका प्रगणनाचे \emph{व्युत्क्रम}
आहे. अशा फलनांना \emph{जोडीकरण फलने} म्हणतात.
\end{explain}

\begin{defn}[जोडीकरण फलन]
  $f\colon A \times B \to \Nat$ यातील $f$ हे एकास-एक असेल, तर ते
  \emph{अंकगणिती जोडीकरण फलन} असते. $f$ हे $A \times B$ चे
  \emph{सांकेतीकरण करते} आणि $f(x,y)$ हा $\tuple{x,y}$ चा
  \emph{संकेतांक} असतो असेही आपण म्हणतो.
\end{defn}

\begin{explain}
जोडीकरण फलनांनी, उदाहरणार्थ, नैसर्गिक संख्यांच्या जोड्यांचे
सांकेतीकरण करता येते; म्हणजेच घटकांची प्रत्येक \emph{जोडी} एका
\emph{एकाच} संख्येने दर्शवता येते. जोडीकरण फलनाचे व्युत्क्रम वापरून
संख्येचे \emph{विसांकेतीकरण} करता येते, म्हणजे ती संख्या कोणती जोडी
दर्शवते हे शोधता येते.
\end{explain}

\begin{prob}
सर्व ऋणेतर परिमेय संख्यांच्या संचाचे एक प्रगणन द्या.
\end{prob}

\begin{prob}
$\Rat$ हा गणनीय आहे, हे दाखवा. कोणतीही परिमेय संख्या
$z/m$ या अपूर्णांकाच्या रूपात लिहिता येते, जेथे $z \in \Int$ आणि
$m \in \Nat^+$ असतात, हे आठवा.
\end{prob}

\begin{prob}
$\Bin^*$ चे एक प्रगणन परिभाषित करा.
\end{prob}

\begin{prob}
प्रास्ताविक तर्कशास्त्राच्या अभ्यासक्रमातून आठवा की प्रत्येक संभाव्य
सत्यताकोष्टक एक सत्यता-फलन व्यक्त करतो. दुसऱ्या शब्दांत, सत्यता-फलने
म्हणजे कोणत्या तरी~$k$ साठी $\Bin^k \to \Bin$ अशी सर्व फलने आहेत.
सर्व सत्यता-फलनांचा संच गणनीय आहे, हे सिद्ध करा.
\end{prob}

\begin{prob}
कोणत्याही अनंत गणनीय संचाच्या सर्व सांत उपसंचांचा संच
गणनीय आहे, हे दाखवा.
\end{prob}

\begin{prob}
$\Nat$ च्या एखाद्या सांत उपसंचाचा पूरक असलेल्या $\Nat$ च्या उपसंचाला
\emph{सांत-पूरक} म्हणतात. म्हणजे, $A \subseteq \Nat$ हा सांत-पूरक
असतो तेव्हा आणि केवळ तेव्हाच, जेव्हा $\Nat\setminus A$ हा सांत असतो.
\begin{quote}\small\textbf{स्रोतदुरुस्ती.} OLSIZ-002: गोठवलेल्या इंग्रजी वाक्यात “complement of a finite set Nat” अशी संबंधदर्शक शब्दांशिवाय अपूर्ण रचना आहे. लगेचचे सूत्र अभिप्रेत अर्थ ठरवते: Nat मधील एखाद्या सांत उपसंचाचा Nat-सापेक्ष पूरक. मराठी मजकुरात हा संबंध स्पष्ट केला आहे.\end{quote}
$\Nat$ चे सांत आणि सांत-पूरक उपसंच हेच ज्याचे घटक आहेत,
असा संच $I$ घ्या. $I$ हा गणनीय आहे, हे दाखवा.
\end{prob}

\begin{prob}
गणनीय इतक्या गणनीय संचांचा संयोग गणनीय असतो,
हे दाखवा. म्हणजे, $A_1$, $A_2$, \dots{} हे संच असतील आणि प्रत्येक
$A_i$ हा गणनीय असेल, तर त्या सर्वांचा संयोग
$\bigcup_{i=1}^\infty A_i$ हाही गणनीय असतो. [टीप: हे कठीण आहे!]
\end{prob}

\begin{prob}
$f \colon A \times B \to \Nat$ हे कोणतेही जोडीकरण फलन असू द्या.
हे $f$ आच्छादकही असेल, तर त्याचे व्युत्क्रम $A \times B$ चे प्रगणन
आहे, हे दाखवा. सर्वसाधारण बाबतीत व्युत्क्रम केवळ फलनाच्या व्याप्तीवर
परिभाषित असते; यावरून जोड्यांचा हा संच गणनीय आहे, हे सिद्ध करा.
\begin{quote}\small\textbf{स्रोतदुरुस्ती.} MRSIZ-002: गोठवलेल्या इंग्रजी सरावात कोणत्याही एकास-एक जोडीकरण फलनाचे व्युत्क्रम थेट प्रगणन असल्याचा दावा आहे. फलन आच्छादक नसेल, तर त्याचे व्युत्क्रम सर्व नैसर्गिक संख्यांवर परिभाषित नसते. मराठी सरावात आच्छादक बाब आणि सामान्य अंशतः-व्युत्क्रम बाब वेगळ्या केल्या आहेत.\end{quote}
\end{prob}

\begin{prob}
$\Nat^3$ चे सांकेतीकरण करणारे फलन स्पष्टपणे द्या.
\end{prob}









\section{एक पर्यायी जोडीकरण फलन}\label{sfr:siz:pai-alt:sec}

\begin{explain}
$\Nat^2$ ची अशी इतर प्रगणनेही आहेत की त्यांची व्युत्क्रमे शोधणे अधिक
सोपे होते. त्यांपैकी एक पुढे दिले आहे. प्रगणन सरणीत दाखवण्याऐवजी,
सुरुवातीला रिकाम्या असलेल्या जागांशी जोडलेल्या धन पूर्णांकांच्या
यादीपासून सुरुवात करा. या जागा पुढीलप्रमाणे एकामागून एक
$\tuple{n,m}$ जोड्यांनी भरत आहोत, अशी कल्पना करा. पहिल्या स्थानी~$0$
असलेल्या जोड्यांपासून (म्हणजे $\tuple{0,m}$ जोड्यांपासून) सुरुवात करा.
पहिली जोडी (म्हणजे $\tuple{0,0}$) पहिल्या रिकाम्या जागेत ठेवा; मग एक
रिकामी जागा सोडून दुसरी जोडी पुढील रिकाम्या जागेत ठेवा आणि पुन्हा एक
जागा सोडा; ही प्रक्रिया अशीच पुढे चालू ठेवा. गोठवलेल्या इंग्रजी
मजकुरात दुसरी जोडी $\tuple{0,2}$ अशी दिली आहे; परंतु खालील सारणीनुसार
येथे $\langle 0,1\rangle$ अभिप्रेत आहे. आता आपल्या प्रगणनाची (अपूर्ण) सुरुवात अशी
दिसते:
\begin{quote}\small\textbf{स्रोतदुरुस्ती.} MRSIZ-003: गोठवलेल्या इंग्रजी मजकुरात एक रिकामी जागा सोडल्यानंतरची दुसरी जोडी $\langle 0,2\rangle$ अशी दिली आहे, पण लगेचची सारणी आणि वर्णनाची क्रमवार रचना ती $\langle 0,1\rangle$ असावी हे दाखवतात. सूत्र-संरचना जशीच्या तशी ठेवून मराठी मजकुरात अपेक्षित जोडी स्पष्ट केली आहे.\end{quote}
\[
\begin{array}{@{}c c c c c c c c c c c@{}}
\mathbf 1 & \mathbf 2 & \mathbf 3 & \mathbf 4 & \mathbf 5 & \mathbf 6 & \mathbf 7 & \mathbf 8 & \mathbf 9 & \mathbf{10} & \dots \\ \\
\tuple{0,0} &  & \tuple{0,1} &  & \tuple{0,2} &  & \tuple{0,3} & & \tuple{0,4} &  & \dots \\
\end{array}
\]
अजून रिकाम्या राहिलेल्या जागांसाठी $\tuple{1,m}$ जोड्यांबाबत हीच
प्रक्रिया पुन्हा करा; पुन्हा प्रत्येक दुसरी रिकामी जागा वगळा:
\[
\begin{array}{@{}c c c c c c c c c c c@{}}
\mathbf 1 & \mathbf 2 & \mathbf 3 & \mathbf 4 & \mathbf 5 & \mathbf 6 & \mathbf 7 & \mathbf 8 & \mathbf 9 & \mathbf{10} & \dots \\ \\
\tuple{0,0} & \tuple{1,0} & \tuple{0,1} &  & \tuple{0,2} & \tuple{1,1} &
\tuple{0,3} & & \tuple{0,4} &  \tuple{1,2} & \dots \\
\end{array}
\]
यानंतर $\tuple{2,m}$ जोड्या आणि मग $\tuple{3,m}$ जोड्या त्याच पद्धतीने
भरा; त्यापुढेही पहिला घटक अनुक्रमे 4, 5, इत्यादी करत पुढे जायचे आहे.
आपल्या पूर्ण प्रगणनाची सुरुवात म्हणून अशी दिसते:
\begin{quote}\small\textbf{स्रोतदुरुस्ती.} OLSIZ-003: गोठवलेल्या इंग्रजी मजकुरात “pairs $\langle 2,m\rangle$, $\langle 2,m\rangle$, etc.” अशी पुनरुक्ती आहे. पूर्ण सारणीतील $\langle 3,0\rangle$ आणि पुढील सर्व ओळींवरून दुसऱ्या ठिकाणी $\langle 3,m\rangle$ अभिप्रेत असल्याचे निश्चित होते; मराठी मजकुरात दुसरी जोडी त्यानुसार दुरुस्त केली आहे.\end{quote}
\[
\begin{array}{@{}cc c c c c c c c c c@{}}
\mathbf 1 & \mathbf 2 & \mathbf 3 & \mathbf 4 & \mathbf 5 & \mathbf 6 & \mathbf 7 & \mathbf 8 & \mathbf 9 & \mathbf{10} & \dots \\ \\
\tuple{0,0} & \tuple{1,0} & \tuple{0,1} & \tuple{2,0}  & \tuple{0,2} &
\tuple{1,1} & \tuple{0,3} & \tuple{3,0}  & \tuple{0,4} &  \tuple{1,2} & \dots \\
\end{array}
\]
वरील सरणीतील चौकटींना या प्रगणनानुसार क्रमांक दिले, तर आपल्याला
नीट नागमोडी रेषा दिसणार नाही; त्याऐवजी पुढील मांडणी दिसेल:
\[
\begin{array}{ c | c | c | c | c | c | c | c }
& \mathbf 0 & \mathbf 1 & \mathbf 2 & \mathbf 3 & \mathbf 4 & \mathbf 5 & \dots \\
\hline
\mathbf 0 & 1 & 3 & 5 & 7 & 9 & 11 & \dots \\
\hline
\mathbf 1 & 2 & 6 & 10 & 14 & 18 & \dots & \dots \\
\hline
\mathbf 2 & 4 & 12 & 20 & 28 & \dots & \dots & \dots \\
\hline
\mathbf 3 & 8 & 24 & 40 & \dots & \dots & \dots & \dots \\
\hline
\mathbf 4 & 16 & 48 & \dots & \dots & \dots & \dots & \dots \\
\hline
\mathbf 5 & 32 & \dots & \dots & \dots & \dots & \dots & \dots \\
\hline
\vdots & \vdots & \vdots & \vdots & \vdots & \vdots & \vdots & \ddots\\
\end{array}
\]
आपल्या प्रगणनात ओळ~$0$ मधील जोड्या विषम क्रमांकांच्या स्थानी आहेत,
हे दिसते; म्हणजे $\tuple{0,m}$ ही जोडी $2m+1$ या स्थानी आहे. दुसऱ्या
ओळीतील $\tuple{1,m}$ जोड्या विषम संख्येच्या दुप्पट क्रमांक असलेल्या
स्थानी आहेत; विशेषतः त्या $2 \cdot (2m+1)$ या स्थानी आहेत. तिसऱ्या
ओळीतील $\tuple{2,m}$ जोड्या विषम संख्येच्या चौपट क्रमांक असलेल्या
स्थानी, म्हणजे $4 \cdot (2m+1)$ या स्थानी आहेत; आणि ही रचना अशीच
पुढे चालते. प्रत्येक ओळीत $(2m+1)$ च्या गुणकांचे मूल्य $1$, $2$,
$4$, $8$, \dots{} असे आहे; हे नेमके~$2$ चे घात आहेत:
$1= 2^0$, $2 = 2^1$, $4 = 2^2$, $8 = 2^3$, \dots\@ विचाराधीन
जोडीचा पहिला घटक हाच नेहमी संबंधित घातांक असतो. म्हणून
$\tuple{n,m}$ या जोडीसाठी गुणक $2^n$ आहे. यातून आपल्याला
$2^n \cdot (2m+1)$ हे सर्वसाधारण सूत्र मिळते. पण या फलनाने जोड्या
\emph{धन} पूर्णांकांकडे पाठवल्या जातात; म्हणजे $\tuple{0,0}$ चे
स्थान~$1$ आहे. स्थानांची सुरुवात~$0$ पासून करायची असेल, तर उत्तरातून
$1$ वजा करावा लागतो. त्यामुळे पुढील फलन मिळते:
\end{explain}

\begin{ex}
$h\colon \Nat^2 \to \Nat$ हे पुढीलप्रमाणे दिलेले फलन घ्या:
\[
h(n,m) = 2^n (2m+1) - 1
\]
हे नैसर्गिक संख्यांच्या जोड्यांचा संच~$\Nat^2$ यासाठीचे जोडीकरण फलन
आहे.
\end{ex}

\begin{explain}
त्यानुसार, $\Nat^2$ च्या आपल्या दुसऱ्या प्रगणनात $\tuple{0,0}$ या
जोडीचा संकेतांक $h(0,0) = 2^0(2\cdot 0+1) - 1 = 0$ आहे;
$\tuple{1,2}$ हिचा संकेतांक $2^{1} \cdot (2 \cdot 2 + 1) - 1 = 2
\cdot 5 - 1 = 9$ आहे; आणि $\tuple{2,6}$ हिचा संकेतांक $2^{2} \cdot (2
\cdot 6 + 1) - 1 = 51$ आहे.
\end{explain}

कधी कधी नैसर्गिक संख्यांच्या जोड्यांचे~$\Nat^2$ सांकेतीकरण करताना ते
आच्छादक असण्याची आवश्यकता नसते. अशा सांकेतीकरणांची व्युत्क्रमे केवळ
अंशतः फलने असतात.

\begin{ex}
$j\colon \Nat^2 \to \Nat^+$ हे पुढीलप्रमाणे दिलेले फलन घ्या:
\[
j(n,m) = 2^n3^m
\]
हे $\Nat^2 \to \Nat$ असे एकास-एक फलन आहे.
\end{ex}









\section{अगणनीय संच}\label{sfr:siz:nen:sec}

\begin{editorial}
  या विभागात \ref{sfr:siz:enm:sec} मधील व्याख्या वापरून $\Bin^\omega$ आणि
  $\Pow{\PosInt}$ यांची अगणनीयता सिद्ध केली आहे. ही मांडणी
  \ref{sfr:siz:enm-alt:sec} मधील आवृत्तीपेक्षा थोडी अधिक प्राथमिक आणि
  थोडी अधिक तपशीलवार ठेवली आहे.
\end{editorial}

धन पूर्णांकांचा संच~$\PosInt$ यासारखे काही संच अनंत असतात. आतापर्यंत
पाहिलेल्या अनंत संचांची सर्व उदाहरणे गणनीय होती. पण हा गुणधर्म
नसलेले अनंत संचही असतात. अशा संचांना \emph{अगणनीय} म्हणतात.

सुरुवातीलाच, अगणनीय संच असतात ही गोष्ट कदाचित आश्चर्यकारक
वाटेल. कोणत्याही गणनीय संच~$A$ साठी $f \colon \PosInt \to A$
असे आच्छादक फलन असते. संच अगणनीय असेल, तर असे फलन
असू शकत नाही. म्हणजे, $\PosInt$ मधील अनंत घटक ना~$A$ कडे
पाठवणारे कोणतेही फलन~$A$ मधील सर्व घटक व्यापू शकत नाही. या अर्थाने,
अनंत असलेल्या धन पूर्णांकांपेक्षा~$A$ मध्ये ``अधिक'' घटक असतात.

एखादा संच अगणनीय आहे हे कसे सिद्ध करायचे? तसे कोणतेही
आच्छादक फलन अस्तित्वात असू शकत नाही, हे दाखवावे लागते. समतुल्यपणे,
$A$ मधील घटकांचे एकदिश अनंत यादीत प्रगणन करता येत नाही, हे दाखवावे
लागते. हे करण्याचा सर्वोत्तम मार्ग म्हणजे~$A$ मधील घटक च्या
प्रत्येक यादीतून किमान एक घटक सुटतो, किंवा $f\colon \PosInt \to A$
असे कोणतेही फलन आच्छादक असू शकत नाही, हे दाखवणे. यासाठी कँटर
यांची \emph{कर्णरेषीय पद्धत} वापरता येते. $A$ मधील घटक ची
एखादी यादी, म्हणजे $x_1$, $x_2$, \dots, दिली असता, आपण~$A$ मधील असा
दुसरा घटक रचतो की त्याच्या रचनेमुळे तो त्या यादीत असणे अशक्य ठरते.

आपले पहिले उदाहरण म्हणजे $0$ आणि $1$ यांच्या सर्व अनंत व मध्ये खंड
नसलेल्या अनुक्रमांचा संच~$\Bin^\omega$.

\begin{thm}
\label{sfr:siz:nen:thm:nonenum-bin-omega}
$\Bin^\omega$ हा अगणनीय आहे.
\end{thm}

\begin{proof}
व्याघाताने सिद्ध करण्यासाठी असे समजा की $\Bin^\omega$ हा गणनीय आहे; म्हणजे,
$\Bin^\omega$ मधील सर्व घटक ची $s_{1}$, $s_{2}$, $s_{3}$,
$s_{4}$, \dots{} अशी यादी आहे, असे समजा. यांपैकी प्रत्येक $s_i$ हा
स्वतः $0$ आणि~$1$ यांचा अनंत अनुक्रम आहे. या यादीतील $i$ व्या अनुक्रमाच्या
$j$ व्या घटकाला $s_i(j)$ म्हणू. मग $i$ वा अनुक्रम~$s_i$ असा आहे:
\[
s_i(1), s_i(2), s_i(3), \dots
\]

ही यादी आणि तिच्यातील प्रत्येक अनुक्रम~$s_i$ चे घटक आपण पुढील सरणीत
मांडू शकतो:
\[
\begin{array}{c|c|c|c|c|c}
& 1 & 2 & 3 & 4 & \dots \\\hline
1 & \mathbf{s_{1}(1)} & s_{1}(2) & s_{1}(3) & s_1(4) & \dots \\\hline
2 & s_{2}(1)& \mathbf{s_{2}(2)} & s_2(3) & s_2(4) & \dots \\\hline
3 & s_{3}(1)& s_{3}(2) & \mathbf{s_3(3)} & s_3(4) & \dots \\\hline
4 & s_{4}(1)& s_{4}(2) & s_4(3) & \mathbf{s_4(4)} & \dots \\\hline
\vdots & \vdots & \vdots & \vdots & \vdots & \mathbf{\ddots}
\end{array}
\]
डाव्या बाजूकडील खुणा यादीतील अनुक्रमांना $s_1$, $s_2$, \dots{} असे
क्रमांक देतात; वरच्या बाजूकडील संख्या प्रत्येक अनुक्रमातील घटक
ना क्रमांक देतात. उदाहरणार्थ, $s_{1}(1)$ हे अनुक्रम~$s_{1}$ मधील
पहिला घटक असलेल्या संख्येचे नाव आहे; ती संख्या $0$ किंवा~$1$
असू शकते. पुढेही याचप्रमाणे अर्थ घ्यायचा.

आता आपण $0$ आणि $1$ यांचा असा अनंत अनुक्रम~$\overline{s}$ रचू, जो
या यादीत असणे अशक्य आहे. $\overline{s}$ ची व्याख्या $s_1$, $s_2$,
\dots{} या यादीवर अवलंबून असेल. $0$ आणि $1$ यांच्या अनंत अनुक्रमांची
कोणतीही अनंत यादी असा अनंत अनुक्रम~$\overline{s}$ निर्माण करते, जो
त्या यादीत नसेल याची खात्री असते.

$\overline{s}$ ची व्याख्या करण्यासाठी आपण त्याचे सर्व घटक
ठरवतो; म्हणजे, प्रत्येक $n \in \PosInt$ साठी $\overline{s}(n)$
ठरवतो. वरील सरणीच्या कर्णरेषेवरून खाली वाचून (म्हणूनच ``कर्णरेषीय
पद्धत'' हे नाव) प्रत्येक $1$ चे $0$ मध्ये आणि प्रत्येक $0$ चे~$1$
मध्ये रूपांतर करून आपण हे करतो. अधिक अमूर्तपणे, कर्णरेषेवरील $n$ वा
घटक, म्हणजे $s_n(n)$, हा $1$ आहे की $0$ यानुसार
$\overline{s}(n)$ हे $0$ किंवा $1$ असे परिभाषित करतो:
\[
\overline{s}(n) =
\begin{cases}
1 & \text{जर $s_{n}(n) = 0$ असेल}\\
0 & \text{जर $s_{n}(n) = 1$ असेल}.
\end{cases}
\]
प्रकरणांनुसार व्याख्येपेक्षा सूत्रे अधिक आवडत असतील, तर
$\overline{s}(n) = 1 - s_n(n)$ अशीही व्याख्या करता येईल.

$\overline{s}$ हा स्पष्टपणे $0$ आणि $1$ यांचा अनंत अनुक्रम आहे, कारण
तो आपल्या सरणीच्या कर्णरेषेवर दिसणाऱ्या $0$ आणि $1$ यांच्या अनुक्रमाचा
पूरक अनुक्रम आहे. म्हणून $\overline{s}$ हा $\Bin^\omega$ चा
घटक आहे. पण तो $s_1$, $s_2$, \dots{} या यादीत असू शकत नाही.
असे का?

तो यादीतील पहिला अनुक्रम~$s_1$ असू शकत नाही, कारण त्याचा पहिला
घटक हा $s_1$ च्या पहिल्या घटकापेक्षा वेगळा आहे. $s_1(1)$
काहीही असो, $\overline{s}(1)$ त्याच्या विरुद्ध असेल अशी आपण व्याख्या
केली. तो यादीतील दुसरा अनुक्रमही असू शकत नाही, कारण $\overline{s}$
चा दुसरा घटक $s_2$ च्या दुसऱ्या घटकापेक्षा वेगळा आहे: $s_2(2)$ हा
$0$ असेल, तर $\overline{s}(2)$ हा $1$ आहे, आणि उलट. पुढेही असेच.

अधिक नेमकेपणाने सांगायचे तर, $\overline{s}$ यादीत असता, तर कोणत्या
तरी $k$ साठी $\overline{s} = s_{k}$ असते. दोन अनुक्रम प्रत्येक स्थानी
जुळतात तेव्हा आणि केवळ तेव्हाच ते एकरूप असतात; म्हणजे, कोणत्याही~$n$
साठी $\overline{s}(n) = s_{k}(n)$ असते. म्हणून, विशेष प्रकरण म्हणून
$n = k$ घेतल्यास $\overline{s}(k) = s_{k}(k)$ असणे आवश्यक ठरते.
$s_k(k)$ हे $0$ किंवा~$1$ आहे. ते $0$ असेल, तर $\overline{s}(k)$ हे~$1$
असले पाहिजे---कारण आपण $\overline{s}$ ची तशीच व्याख्या केली आहे. पण
$s_k(k) = 1$ असेल, तर पुन्हा $\overline{s}$ च्या व्याख्येमुळे
$\overline{s}(k) = 0$ होते. दोन्ही बाबतीत
$\overline{s}(k) \neq s_{k}(k)$.

सुरुवातीला आपण $\Bin^\omega$ मधील घटक ची $s_1$, $s_2$,
\dots{} अशी यादी आहे असे गृहीत धरले. या यादीपासून आपण असा
अनुक्रम~$\overline{s}$ रचला की तो यादीत असू शकत नाही, हे सिद्ध केले.
पण सर्व $s_i$ हे $0$ आणि $1$ यांचे अनुक्रम असतील, तर तो नक्कीच $0$
आणि $1$ यांचा अनुक्रम असतो; म्हणजे, $\overline{s} \in
\Bin^\omega$. यावरून विशेषतः $\Bin^\omega$ मधील \emph{सर्व}
घटक ची यादी असू शकत नाही. कारण अशी कोणतीही यादी दिली, तरी
तिच्यात नसेल असा अनुक्रम~$\overline{s}$ आपण रचू शकतो. म्हणून
$\Bin^\omega$ मधील सर्व अनुक्रमांची यादी आहे हे गृहीतक व्याघात निर्माण
करते.
\end{proof}

\begin{explain}
या सिद्धता-पद्धतीला ``कर्णीकरण'' म्हणतात, कारण तिच्यात सरणीच्या
कर्णरेषेचा वापर करून~$\overline{s}$ ची व्याख्या केली जाते. कर्णीकरणात
सरणी असलीच पाहिजे असे नाही: प्रत्यक्ष सरणी किंवा कर्णरेषा नसतानाही
हीच कल्पना वापरून संच गणनीय नाहीत, हे दाखवता येते.
\end{explain}

\begin{thm}
\label{sfr:siz:nen:thm:nonenum-pownat}
$\Pow{\PosInt}$ हा गणनीय नाही.
\end{thm}

\begin{proof}
आपण त्याच पद्धतीने पुढे जाऊ: $\PosInt$ च्या उपसंचांची प्रत्येक यादी
दिली असता त्या यादीत नसलेला $\PosInt$ चा एक उपसंच असतो, हे दाखवू.
$\PosInt$ च्या उपसंचांची पुढील यादी दिली आहे असे समजा:
\[
Z_{1}, Z_{2}, Z_{3}, \dots
\]
आता आपण $\overline{Z}$ हा असा संच परिभाषित करतो की प्रत्येक
$n \in \PosInt$ साठी $n \in \overline{Z}$ तेव्हा आणि केवळ तेव्हाच,
जेव्हा $n \notin Z_{n}$:
\[
\overline{Z} = \Setabs{n \in \PosInt}{n \notin Z_n}
\]
$\overline{Z}$ हा स्पष्टपणे धन पूर्णांकांचा संच आहे, कारण गृहीतकानुसार
प्रत्येक~$Z_n$ तसा आहे; म्हणून $\overline{Z} \in
\Pow{\PosInt}$. पण $\overline{Z}$ यादीत असू शकत नाही. हे दाखवण्यासाठी
प्रत्येक $k \in \PosInt$ साठी $\overline{Z} \neq Z_k$, हे सिद्ध करू.

आता $k \in \PosInt$ स्वैर असू द्या. प्रत्येक $n \in \PosInt$ साठी
$n \in \overline{Z}$ तेव्हा आणि केवळ तेव्हाच, जेव्हा $n \notin Z_n$,
अशी $\overline{Z}$ ची व्याख्या केली आहे. विशेषतः $n=k$ घेतल्यास,
$k \in \overline{Z}$ तेव्हा आणि केवळ तेव्हाच, जेव्हा $k \notin Z_k$.
पण यावरून $\overline{Z} \neq Z_k$ हे दिसते, कारण $k$ हा एकाचा
घटक आहे आणि दुसऱ्याचा नाही; म्हणून $\overline{Z}$ आणि $Z_k$
यांचे घटक वेगवेगळे आहेत. $k$ स्वैर असल्याने $\overline{Z}$ हा
$Z_1$, $Z_2$, \dots{} या यादीत नाही.
\end{proof}

\begin{explain}
मागील सिद्धतेत कर्णरेषेचा उल्लेख झाला नाही; पण पुढील चित्र मनात आणले,
तर तिच्यात कर्णरेषा आहे असे समजता येते. $Z_1$, $Z_2$, \dots{} हे संच
एका सरणीत लिहिले आहेत अशी कल्पना करा; प्रत्येक घटक~$j \in Z_i$
हा $j$ व्या स्तंभात लिहिला आहे. समजा, त्या यादीतील पहिले चार संच
$\{1,2,3,\dots\}$, $\{2, 4, 6, \dots\}$, $\{1,2,5\}$ आणि
$\{3,4,5,\dots\}$ आहेत. मग सरणीची सुरुवात अशी होईल:
\[
\begin{array}{r@{}rrrrrrr}
  Z_1 = \{ & \mathbf{1}, & 2, & 3, & 4, & 5, & 6, & \dots\}\\
  Z_2 = \{ &  & \mathbf{2}, &  & 4, &  & 6, & \dots\}\\
  Z_3 = \{ & 1, & 2, &  &  & 5\phantom{,} &  & \}\\
  Z_4 = \{ &  &  & 3, & \mathbf{4}, & 5, & 6, & \dots\}\\
  \vdots & & & & & \ddots
\end{array}
\]
यानंतर कर्णरेषेवरून खाली जाताना कर्णरेषेवर दिसणाऱ्या संख्या वगळून,
आणि $j$ व्या ओळीच्या त्याच क्रमांकाच्या स्तंभातील जागा रिकामी असेल अशा~$j$
चा समावेश करून, $\overline{Z}$ हा संच मिळतो. वरील उदाहरणात
आपण $1$ आणि $2$ वगळू, $3$ चा समावेश करू, $4$ वगळू, इत्यादी.
\end{explain}

\begin{prob}
कर्णरेषीय युक्तिवादाने $\Pow{\Nat}$ हा अगणनीय आहे, हे दाखवा.
\end{prob}

\begin{prob}\label{sfr:siz:nen:prob:f-posint}
$f \colon \PosInt \to \PosInt$ अशा फलनांचा संच अगणनीय
आहे, हे स्पष्ट कर्णरेषीय युक्तिवादाने दाखवा. म्हणजे, $f_1$, $f_2$,
\dots{} ही फलनांची यादी असेल आणि प्रत्येक $f_i\colon
\PosInt \to \PosInt$ असेल, तर या यादीत नसलेले कोणते तरी
$\overline{f}\colon \PosInt \to \PosInt$ असते, हे दाखवा.
\end{prob}









\section{न्यूनीकरण}\label{sfr:siz:red:sec}

\begin{editorial}
  या विभागात \ref{sfr:siz:nen:sec} मधील निष्कर्षांशी जुळणारी अगणनीयता
  न्यूनीकरणाने सिद्ध केली आहे. \ref{sfr:siz:nen-alt:sec} मधील निष्कर्षांशी
  जुळणारी, किंचित अधिक संक्षिप्त पर्यायी आवृत्ती \ref{sfr:siz:red-alt:sec}
  मध्ये दिली आहे.
\end{editorial}

कर्णीकरणाच्या युक्तिवादाने $\Pow{\PosInt}$ हा अगणनीय आहे,
हे आपण दाखवले. सर्व अनंत $0$--$1$ अनुक्रमांचा संच~$\Bin^\omega$ हा
अगणनीय आहे, याची सिद्धता आपल्याकडे आधीच होती.
$\Pow{\PosInt}$ हा अगणनीय आहे हे सिद्ध करण्याचा आणखी एक मार्ग
असा: \emph{$\Pow{\PosInt}$ हा गणनीय असेल, तर $\Bin^\omega$
  हाही गणनीय असतो} हे दाखवा. $\Bin^\omega$ हा गणनीय
नाही, हे माहीत असल्याने $\Pow{\PosInt}$ हाही तसा असू शकत नाही. यालाच
एक समस्या दुसऱ्या समस्येकडे \emph{न्यूनीत करणे} म्हणतात---या बाबतीत
$\Bin^\omega$ चे प्रगणन करण्याची समस्या आपण $\Pow{\PosInt}$ चे प्रगणन
करण्याच्या समस्येकडे न्यूनीत करतो. नंतरच्या समस्येचे उत्तर---म्हणजे
$\Pow{\PosInt}$ चे प्रगणन---पहिल्या समस्येचे उत्तर---म्हणजे
$\Bin^\omega$ चे प्रगणन---देईल.

संच~$B$ च्या प्रगणनाची समस्या संच~$A$ च्या प्रगणनाच्या समस्येकडे कशी
न्यूनीत करायची? $A$ चे प्रगणन $B$ च्या प्रगणनात रूपांतरित करण्याची
पद्धत आपण देतो. त्यासाठीचा सर्वांत सोपा मार्ग म्हणजे
$f\colon A \to B$ असे आच्छादक फलन परिभाषित करणे. $x_1$, $x_2$,
\dots{} हे~$A$ चे प्रगणन असेल, तर $f(x_1)$, $f(x_2)$, \dots{} हे~$B$
चे प्रगणन करील. आपल्या बाबतीत आपण
$f\colon \Pow{\PosInt} \to \Bin^\omega$ असे आच्छादक फलन शोधत आहोत.

\begin{prob}
$g\colon B \to A$ असे एकास-एक फलन असेल आणि $B$ हा
अगणनीय असेल, तर $A$ हाही अगणनीय असतो, हे दाखवा.
$g$ वापरून~$A$ चे प्रगणन~$B$ च्या प्रगणनात कसे रूपांतरित करता येते,
हे दाखवून सिद्धता द्या.
\end{prob}

\begin{proof}[न्यूनीकरणाने {\ref{sfr:siz:nen:thm:nonenum-pownat}} ची सिद्धता]
$\Pow{\PosInt}$ हा गणनीय आहे, आणि म्हणून त्याचे $Z_{1}$,
$Z_{2}$, $Z_{3}$, \dots{} असे प्रगणन आहे, असे समजा.

$f \colon \Pow{\PosInt} \to \Bin^\omega$ हे फलन पुढीलप्रमाणे परिभाषित
करा: $f(Z)$ हा असा अनुक्रम~$s$ असू द्या की $n \in Z$ असेल तेव्हा
आणि केवळ तेव्हाच $s(n) = 1$, आणि अन्यथा $s(n) = 0$. याने
फलन स्पष्टपणे परिभाषित होते, कारण $Z \subseteq \PosInt$ असताना कोणताही
$n \in \PosInt$ हा $Z$ चा घटक असतो किंवा नसतो. उदाहरणार्थ, धन सम संख्यांचा संच
$2\PosInt = \{2, 4, 6, \dots\}$ हा $010101\dots$ या अनुक्रमाकडे,
रिक्त संच $0000\dots$ कडे आणि $\PosInt$ हा संच स्वतः $1111\dots$ कडे
पाठवला जातो.
\begin{quote}\small\textbf{स्रोतदुरुस्ती.} OLSIZ-004: गोठवलेल्या इंग्रजी व्याख्येत f(Z) ला s\_k असे म्हटले आहे, पण k चा Z शी संबंध किंवा निवड दिलेली नाही. अभिप्रेत अनुक्रम Z नेच ठरतो; मराठी व्याख्येत एकच बद्ध निर्गत नाव s आणि त्याचे घटक s(n) सातत्याने वापरले आहेत.\end{quote}

हे फलन आच्छादक सुद्धा आहे: $0$ आणि $1$ यांचा प्रत्येक अनुक्रम धन
पूर्णांकांच्या कोणत्या तरी संचाशी जुळतो; तो संच म्हणजे अनुक्रमात ज्या
स्थानी~$1$ आहे त्या स्थानांना अनुरूप पूर्णांकांचा संच. अधिक नेमकेपणाने,
$s \in \Bin^\omega$ असे समजा. पुढीलप्रमाणे $Z \subseteq \PosInt$
परिभाषित करा:
\[
Z = \Setabs{n \in \PosInt}{s(n) = 1}
\]
मग $f$ च्या व्याख्येवरून तपासता येते की $f(Z) = s$.

आता पुढील यादी विचारात घ्या:
\[
f(Z_1), f(Z_2), f(Z_3), \dots
\]
$f$ हे आच्छादक असल्याने $\Bin^\omega$ मधील प्रत्येक सदस्य हा
$f$ च्या कोणत्या तरी फलसाधकावरील मूल्य असला पाहिजे, आणि म्हणून तो
यादीत आला पाहिजे. त्यामुळे ही यादी $\Bin^\omega$ मधील सर्व घटकांचे
प्रगणन करते.

म्हणून $\Pow{\PosInt}$ हा गणनीय असता, तर $\Bin^\omega$ हाही
गणनीय असता. पण $\Bin^\omega$ हा अगणनीय आहे
(\ref{sfr:siz:nen:thm:nonenum-bin-omega}). त्यामुळे $\Pow{\PosInt}$ हा
अगणनीय आहे.
\end{proof}

\begin{explain}
न्यूनीकरण कोणत्या दिशेने केले जाते याबद्दल गोंधळ होणे सोपे आहे.
उदाहरणार्थ, $g \colon \Bin^\omega \to B$ असे आच्छादक फलन
$B$ हा अगणनीय आहे, हे \emph{सिद्ध करत नाही}. ($g
\colon \Bin^\omega \to \Bin$ हे $g(s) = s(1)$ ने परिभाषित केलेले,
$0$ आणि $1$ यांच्या अनुक्रमाला त्याच्या पहिल्या घटक कडे
पाठवणारे फलन विचारात घ्या. काही अनुक्रम $0$ ने आणि काही $1$ ने सुरू
होत असल्याने ते आच्छादक आहे. पण $\Bin$ हा सांत संच आहे.)
तसेच फलन~$f$ हे आच्छादक असलेच पाहिजे; अन्यथा युक्तिवाद लागू
होत नाही: $f(x_1)$, $f(x_2)$, \dots{} यात~$B$ मधील सर्व घटक
असतील याची खात्री उरत नाही. उदाहरणार्थ,
\[
h(n) = \underbrace{000\dots0}_{\text{$n$ वेळा $0$}}111\dots
\]
याने $h\colon \PosInt \to
\Bin^\omega$ असे फलन परिभाषित होते: प्रत्येक मूल्याची सुरुवात n शून्यांनी
होते आणि त्यानंतर कायम एक येतात. या फलनाच्या व्याप्तीत अखेरीस कायम एक
न होणारे अनुक्रम येत नाहीत, म्हणून ते आच्छादक नाही; पण $\PosInt$ हा
गणनीय आहे.
\begin{quote}\small\textbf{स्रोतदुरुस्ती.} OLSIZ-005: गोठवलेल्या इंग्रजी प्रदर्शनात h(n) चे मूल्य फक्त n शून्यांची सांत चिन्हमाला आहे, जी अनंत अनुक्रमांच्या Bin-omega संचात येत नाही. मराठी प्रदर्शनात त्यानंतर अनंत एकांची शेपटी जोडून वैध, अनाच्छादक फलन दिले आहे.\end{quote}
\end{explain}

\begin{prob}
धन पूर्णांकांच्या जोड्यांच्या \emph{सर्व संचांचा} संच न्यूनीकरणाच्या
युक्तिवादाने अगणनीय आहे, हे दाखवा.
\end{prob}

\begin{prob}\label{sfr:siz:red:prob:nat-nat}
  $f\colon \Nat \to \Nat$ अशा सर्व फलनांचा संच~$X$ हा न्यूनीकरणाच्या
  युक्तिवादाने अगणनीय आहे, हे दाखवा. (सूचना: $X$ पासून
  $\Bin^\omega$ कडे जाणारे आच्छादक फलन द्या.)
\end{prob}

\begin{prob}
नैसर्गिक संख्यांच्या सर्व अनंत अनुक्रमांचा संच~$\Nat^\omega$ हा
न्यूनीकरणाच्या युक्तिवादाने अगणनीय आहे, हे दाखवा.
\end{prob}

\begin{prob}
$P$ हा धन पूर्णांकांच्या संचापासून $\{0\}$ या संचाकडे जाणाऱ्या फलनांचा
संच असू द्या; आणि $Q$ हा धन पूर्णांकांच्या संचापासून $\{0\}$ या
संचाकडे जाणाऱ्या \emph{अंशतः} फलनांचा संच असू द्या. $P$ हा
गणनीय आहे आणि $Q$ नाही, हे दाखवा. (सूचना: $\Bin^\omega$ चे
प्रगणन करण्याची समस्या~$Q$ चे प्रगणन करण्याच्या समस्येकडे न्यूनीत करा.)
\end{prob}

\begin{prob}
$S$ हा धन पूर्णांकांच्या संचापासून \{0,1\} या संचाकडे जाणाऱ्या सर्व
आच्छादक फलनांचा संच असू द्या; म्हणजे, $S$ मध्ये सर्व
आच्छादक~$f\colon \PosInt \to \Bin$ आहेत. $S$ हा
अगणनीय आहे, हे दाखवा.
\end{prob}

\begin{prob}
सर्व वास्तव संख्यांचा संच~$\Real$ हा अगणनीय आहे, हे दाखवा.
\end{prob}









\section{तुल्यबलता}\label{sfr:siz:equ:sec}

संचांच्या ``आकारमाना''ची एक अंतःप्रज्ञ कल्पना आपल्याकडे आहे आणि ती
सांत संचांसाठी व्यवस्थित चालते. पण अनंत संचांचे काय? \emph{कोणत्याही}
आकारमानाच्या दोन संचांच्या आकारमानांची औपचारिक रीतीने तुलना करायची
असेल, तर दोन संच एकाच आकारमानाचे कधी असतात याची व्याख्या प्रथम करणे
योग्य ठरते. फ्रेगे यांचे पुढील उदाहरण पाहा:
\begin{quote}
  उपाहारगृहातील कर्मचाऱ्याला मेजावर ताटांइतक्याच सुऱ्या मांडल्या आहेत
  याची खात्री करायची असेल, तर प्रत्येक ताटाच्या उजवीकडे एक सुरी ठेवून
  आणि मेजावरील प्रत्येक सुरी कोणत्या तरी ताटाच्या उजवीकडे येईल असे
  केल्यास त्याला ताटे किंवा सुऱ्या मोजण्याची गरज नसते. अशा प्रकारे
  ताटे आणि सुऱ्या एकास-एक संगतीने परस्परांशी जोडली जातात, आणि तीही
  त्याच अवकाशीय संबंधाने. (Frege, 1884, \S70)
\end{quote}
या उताऱ्यातील अंतर्दृष्टी पुढील औपचारिक व्याख्येत मांडता येते:

\begin{defn}\label{sfr:siz:equ:comparisondef}
  $A$ हा $B$ शी \emph{तुल्यबल} असतो, आणि ते $\cardeq{A}{B}$ असे लिहितो,
  तेव्हा आणि केवळ तेव्हाच, जेव्हा $f \colon A \to B$ असे
  एकास-एक आच्छादन असते.
\end{defn}

\begin{prop}\label{sfr:siz:equ:equinumerosityisequi}
तुल्यबलता हा तुल्यता संबंध आहे.
\end{prop}

\begin{proof}
तुल्यबलता परावर्ती, सममित आणि संक्रमक आहे, हे दाखवावे लागेल. $A, B$
आणि $C$ हे संच असू द्या.

\emph{परावर्तिता.} प्रत्येक $x \in A$ साठी $\Id{A} (x) = x$ असे
असलेले एकरूपता फलन $\Id{A} \colon A \to A$ हे एकास-एक आच्छादन आहे.
म्हणून $\cardeq{A}{A}$.

\emph{सममितता.} $\cardeq{A}{B}$ आहे, म्हणजे $f\colon A \to B$ असे
एकास-एक आच्छादन आहे, असे समजा. $f$ हे एकास-एक व आच्छादक असल्याने त्याचे
व्युत्क्रम $f^{-1}$ अस्तित्वात आहे आणि तेही एकास-एक व आच्छादक आहे. म्हणून
$f^{-1}\colon B \to A$ हे एकास-एक आच्छादन आहे; त्यामुळे
$\cardeq{B}{A}$.

\emph{संक्रमकता.} $\cardeq{A}{B}$ आणि $\cardeq{B}{C}$ आहे, म्हणजे
$f\colon A \to B$ आणि $g\colon B \to
C$ अशी एकास-एक आच्छादने आहेत, असे समजा. मग संयोजन
$\comp{f}{g}\colon A \to C$ हे एकास-एक व आच्छादक आहे; म्हणून
$\cardeq{A}{C}$.
\end{proof}

\begin{prop}
$\cardeq{A}{B}$ असेल, तर $A$ हा गणनीय असतो तेव्हा आणि केवळ
तेव्हाच $B$ तसा असतो.
\end{prop}

\begin{editorial}
पुढील सिद्धतेत \ref{sfr:siz:enm:sec} समाविष्ट असेल, तर
\ref{sfr:siz:enm:defn:enumerable} वापरली जाते; अन्यथा
\ref{sfr:siz:enm-alt:defn:enumerable} वापरली जाते.
\end{editorial}

\begin{proof}
$\cardeq{A}{B}$ आहे, म्हणून $f \colon A
\to B$ असे काही एकास-एक आच्छादन आहे, आणि $A$ हा गणनीय आहे,
असे समजा.

  मग एकतर $A = \emptyset$ असतो, किंवा $g\colon \PosInt \to A$ असे
  आच्छादक फलन असते. $A = \emptyset$ असेल, तर $B = \emptyset$
  सुद्धा असतो (अन्यथा एखादा घटक~$y \in B$ असेल,
  पण $x \in A$ असा कोणताही घटक नसल्याने $f(x) = y$ अशी पूर्वप्रतिमा
  मिळणार नाही). दुसरीकडे, $g\colon \PosInt \to A$ हे आच्छादक
  असेल, तर $\comp{g}{f} \colon \PosInt \to B$ हे आच्छादक आहे.
  हे पाहण्यासाठी $y \in B$ असू द्या. $f$ हे आच्छादक असल्याने
  $f(x) = y$ असा एखादा $x \in A$ घटक असतो. $g$ हे आच्छादक
  असल्याने $g(n) =
  x$ अशी एखादी $n \in \PosInt$ संख्या असते. त्यामुळे
\[
(\comp{g}{f})(n) = f(g(n)) = f(x) = y
\]
  आणि म्हणून $\comp{g}{f}$ हे आच्छादक आहे. $\comp{g}{f}$ हे~$B$
  चे प्रगणन आहे; त्यामुळे $B$ हा गणनीय आहे.
\begin{quote}\small\textbf{स्रोतदुरुस्ती.} OLSIZ-006: गोठवलेल्या इंग्रजी सिद्धतेच्या दोन्ही पर्यायी शाखांतील रिक्त-संच प्रकरणात g(x)=y असे लिहिले आहे. त्या शाखेत g अजून अस्तित्वात नसते; A पासून B पर्यंतचे आधी दिलेले द्विएक फलन f वापरूनच B रिक्त असल्याचा निष्कर्ष मिळतो. मराठी मजकुरात दोन्ही ठिकाणी f(x)=y केले आहे.\end{quote}
\begin{quote}\small\textbf{स्रोतदुरुस्ती.} MRSIZ-004: गोठवलेल्या पर्यायी शाखेत एकाच प्रांताला आधी “initial sequence of natural numbers” आणि नंतर “initial segment” म्हटले आहे. आधीच्या औपचारिक व्याख्येशी सुसंगत म्हणून मराठीत दोन्हीकडे “आरंभीचा खंड” वापरला आहे.\end{quote}

$B$ हा गणनीय असेल, तर~$f$ ऐवजी
एकास-एक आच्छादन $f^{-1}\colon B \to A$ घेऊन तोच युक्तिवाद पुन्हा केल्यास
$A$ हा गणनीय आहे, हे मिळते.
\end{proof}

\begin{prob}
$\cardeq{A}{C}$ आणि $\cardeq{B}{D}$ असेल, आणि $A \cap B =
C \cap D = \emptyset$ असेल, तर $\cardeq{A \cup B}{C \cup D}$ आहे,
हे दाखवा.
\end{prob}

\begin{prob}
$A$ हा अनंत आणि गणनीय असेल, तर $\cardeq{A}{\Nat}$ आहे, हे
दाखवा.
\end{prob}









\section{वेगवेगळ्या आकारमानांचे संच आणि कँटर यांचे प्रमेय}\label{sfr:siz:car:sec}

\begin{explain}
दोन संचांचे आकारमान समान असते या कल्पनेचे नेमके विधान आपण दिले आहे.
एक संच दुसऱ्यापेक्षा लहान असतो या कल्पनेचेही नेमके विधान करता येते.
``लहान (किंवा तुल्यबल) असणे'' याच्या व्याख्येत संचांदरम्यानचे
एकास-एक आच्छादन न घेता पहिल्या संचापासून दुसऱ्या संचाकडे जाणारे
एकास-एक फलन घ्यावे लागेल. असे फलन अस्तित्वात असेल, तर पहिल्या
संचाचे आकारमान दुसऱ्या संचाच्या आकारमानापेक्षा लहान किंवा त्याच्याइतके
असते. अंतःप्रज्ञेने, एका संचापासून दुसऱ्या संचाकडे जाणारे
एकास-एक फलन हे फलनाच्या व्याप्तीत प्रांताइतके तरी घटक आहेत
याची खात्री देते, कारण प्रांतातील कोणतेही दोन घटक व्याप्तीतील
एकाच घटक कडे पाठवले जात नाहीत.
\end{explain}

\begin{defn}
$A$ हा $B$ पेक्षा \emph{आकारमानाने मोठा नाही}, हे $\cardle{A}{B}$ असे
लिहितो, तेव्हा आणि केवळ तेव्हाच, जेव्हा $f \colon A \to B$ असे
एकास-एक फलन असते.
\end{defn}

हा संबंध परावर्ती आणि संक्रमक आहे, पण सममित नाही, हे स्पष्ट आहे (हे
सरावासाठी सोडले आहे). एक संच दुसऱ्या संचापेक्षा (काटेकोरपणे) लहान आहे
असे सांगणारी कल्पनाही मांडता येते.

\begin{defn}
$A$ हा $B$ पेक्षा \emph{आकारमानाने लहान} आहे, हे $\cardless{A}{B}$ असे
लिहितो, तेव्हा आणि केवळ तेव्हाच, जेव्हा $f\colon A \to B$ असे
एकास-एक फलन असते, पण $g\colon A
\to B$ असे कोणतेही एकास-एक आच्छादन नसते; म्हणजे, $\cardle{A}{B}$ आणि
$\cardneq{A}{B}$.
\end{defn}

हा संबंध अपरावर्ती आणि संक्रमक आहे, हे स्पष्ट आहे. (हे सरावासाठी
सोडले आहे.) या संकेतनाने संच $A$ हा गणनीय असतो तेव्हा आणि
केवळ तेव्हाच $\cardle{A}{\Nat}$, आणि $A$ हा अगणनीय असतो
तेव्हा आणि केवळ तेव्हाच $\cardless{\Nat}{A}$, असे म्हणता येते. त्यामुळे
%
\ref{sfr:siz:nen-alt:thm:nonenum-pownat}
हे प्रमेय $\cardless{\Nat}{\Pow{\Nat}}$ असे पुन्हा मांडता येते.
खरे तर, हा शेवटचा मुद्दा \emph{पूर्णपणे व्यापक} आहे, हे
Cantor (1892) यांनी सिद्ध केले:

\begin{thm}[कँटर]\label{sfr:siz:car:thm:cantor}
कोणत्याही संच $A$ साठी $\cardless{A}{\Pow{A}}$.
\end{thm}

\begin{proof}
$f(x) = \{x\}$ हे $f \colon A \to \Pow{A}$ असे एकास-एक फलन आहे;
कारण $x \neq y$ असेल, तर विस्तारात्मकतेमुळे $\{x\} \neq \{y\}$ सुद्धा
असते, आणि म्हणून $f(x) \neq f(y)$. त्यामुळे $\cardle{A}{\Pow{A}}$.

\begin{editorial}
\ref{sfr:siz:nen:sec} समाविष्ट असेल, तर आपण संथ सिद्धता देतो; अन्यथा
\ref{sfr:siz:nen-alt:sec} शी जुळणारी अधिक जलद सिद्धता देतो.
\end{editorial}

%
  आता $g\colon A \to
  \Pow{A}$ असे आच्छादक फलन, आणि त्यामुळे एकास-एक व आच्छादक फलनही,
  असू शकत नाही हे दाखवू; म्हणून $\cardneq{A}{\Pow{A}}$. त्यासाठी
  $g\colon A \to \Pow{A}$ असे समजा. $g$ हे सर्वत्र परिभाषित असल्याने
  प्रत्येक $x \in A$ हा $g(x)
  \subseteq A$ अशा उपसंचाकडे पाठवला जातो. $g$ हे आच्छादक असू शकत नाही,
  हे दाखवता येते. त्यासाठी व्याख्येमुळेच~$g$ च्या व्याप्तीत येऊ न
  शकणारा $\overline{A} \subseteq A$ असा उपसंच परिभाषित करू. पुढील संच घ्या:
  \[
  \overline{A} = \Setabs{x \in A}{x \notin g(x)}.
  \]
  प्रत्येक $x \in A$ साठी $g(x)$ परिभाषित असल्याने $\overline{A}$ हा
  स्पष्टपणे~$A$ चा सुयोग्य रीतीने परिभाषित उपसंच आहे. पण तो~$g$ च्या
  व्याप्तीत असू शकत नाही. $x \in A$ स्वैर असू द्या; आपण
  $\overline{A} \neq g(x)$ दाखवू. $x \in g(x)$ असेल, तर तो
  $x \notin g(x)$ हा गुणधर्म पूर्ण करत नाही, आणि म्हणून~$\overline{A}$
  च्या व्याख्येवरून $x \notin
  \overline{A}$. $x \in \overline{A}$ असेल, तर त्याने~$\overline{A}$
  चा व्याख्यात्मक गुणधर्म पूर्ण केला पाहिजे; म्हणजे $x \in A$ आणि
  $x \notin
  g(x)$. $x$ स्वैर असल्याने प्रत्येक $x \in A$ साठी $x \in g(x)$ तेव्हा
  आणि केवळ तेव्हाच $x \notin \overline{A}$; म्हणून
  $g(x) \neq \overline{A}$.
  दुसऱ्या शब्दांत, $\overline{A}$ हा~$g$ च्या व्याप्तीत येऊ शकत नाही;
  त्यामुळे~$g$ आच्छादक आहे या गृहीतकाला व्याघात होतो.{}
\begin{quote}\small\textbf{स्रोतदुरुस्ती.} OLSIZ-007: गोठवलेल्या संथ सिद्धतेत x स्वैरपणे A मधून घेतल्यानंतर निष्कर्षाचा आवाका चुकून “प्रत्येक x जो A-bar मध्ये आहे” असा दिला आहे. g(x) आणि A-bar वेगळे आहेत हे प्रत्येक $x\in A$ साठी दाखवणे आवश्यक आहे; मराठी सिद्धतेत तो आवाका दुरुस्त केला आहे.\end{quote}
\end{proof}

\begin{explain}

  \ref{sfr:siz:car:thm:cantor} ची सिद्धता आणि
  \ref{sfr:siz:nen:thm:nonenum-pownat} ची सिद्धता यांची तुलना करणे उद्बोधक
  आहे. तेथे~$\PosInt$ च्या उपसंचांची कोणतीही $Z_1$, $Z_2$, \dots{}
  अशी यादी दिली असता, यादीत नसेल असा संख्यांचा संच~$\overline{Z}$
  रचता येतो, हे दाखवले. तो यादीत नसतो याची खात्री यामुळे मिळते की
  प्रत्येक $n \in
  \PosInt$ साठी $n \in Z_n$ तेव्हा आणि केवळ तेव्हाच
  $n \notin \overline{Z}$. त्यामुळे $Z_n$ आणि $\overline{Z}$ यांपैकी
  एकाचा घटक असलेली, पण दुसऱ्याचा नसलेली कोणती तरी संख्या
  नेहमी असते. येथेही तीच कल्पना वापरतो; फरक एवढाच की निर्देशांक~$n$
  आता~$\PosInt$ चे नसून~$A$ चे घटक आहेत. $\overline{A}$ ची
  व्याख्या अशी केली आहे की प्रत्येक $x \in A$ साठी तो~$g(x)$ पेक्षा
  वेगळा असतो, कारण $x \in g(x)$ तेव्हा आणि केवळ तेव्हाच
  $x \notin \overline{A}$. पुन्हा, $g(x)$ आणि $\overline{A}$ यांपैकी
  एकाचा घटक असलेला, पण दुसऱ्याचा नसलेला~$A$ चा कोणता तरी
  घटक नेहमी असतो. म्हणूनच जसा $\overline{Z}$ हा $Z_1$,
  $Z_2$, \dots{} या यादीत येऊ शकत नाही, तसा $\overline{A}$ हा~$g$ च्या
  व्याप्तीत येऊ शकत नाही.


  \ref{sfr:siz:car:thm:cantor} ची सिद्धता आणि
  \ref{sfr:siz:nen-alt:thm:nonenum-pownat} ची सिद्धता यांची तुलना करणे
  उद्बोधक आहे. तेथे~$\Nat$ च्या उपसंचांची कोणतीही $N_0$, $N_1$,
  $N_2$, \dots{} अशी यादी दिली असता, यादीत नसेल असा संख्यांचा संच~$D$
  रचता येतो, हे दाखवले. प्रत्येक $n \in \Nat$ साठी $n \in N_n$ तेव्हा
  आणि केवळ तेव्हाच $n \notin
  D$ असल्यामुळे तो यादीत नसतो याची खात्री मिळते. येथेही तीच कल्पना
  वापरतो; फरक एवढाच की निर्देशांक~$n$ आता~$\Nat$ चे नसून~$A$ चे
  घटक आहेत. $D$ ची व्याख्या अशी केली आहे की प्रत्येक $x
  \in A$ साठी तो~$g(x)$ पेक्षा वेगळा असतो, कारण $x \in g(x)$ तेव्हा
  आणि केवळ तेव्हाच $x \notin D$.

या सिद्धतेची रसेल यांच्या विरोधापत्तीच्या सिद्धतेशीही तुलना करणे उपयुक्त
आहे: \ref{sfr:set:rus:thm:russells-paradox}. खरे तर, कँटर यांच्या
प्रमेयातूनच रसेल यांना स्वतःची विरोधापत्ती सुचली.
\end{explain}

\begin{prob}
  कोणत्याही संच~$A$ साठी $g\colon \Pow{A} \to
  A$ असे एकास-एक फलन असू शकत नाही, हे दाखवा. सूचना: $g\colon
  \Pow{A} \to A$ हे एकास-एक आहे, असे समजा. पुढील संच विचारात घ्या:
  $D = \Setabs{g(B)}{B \subseteq A \text{ आणि
  } g(B) \notin B}$. $x = g(D)$ असे घ्या. $g$ हे एकास-एक आहे हा
  हे तथ्य वापरून व्याघात मिळवा.
\end{prob}









\section{आकारमानाची कल्पना आणि श्रेडर--बर्नस्टाइन प्रमेय}\label{sfr:siz:sb:sec}

\begin{explain}
पुढील विचार अंतर्ज्ञानाला पटणारा आहे: $A$ हा $B$ पेक्षा आकारमानाने मोठा नाही
आणि $B$ हा $A$ पेक्षा आकारमानाने मोठा नाही, तर $A$ आणि $B$ तुल्यबल
आहेत. खरे सांगायचे तर, हा विचार \emph{चुकीचा} असता, तर तुल्यबलतेच्या
आपल्या व्याख्येचा संचांच्या ``आकारमानां''च्या तुलनेशी काही संबंध आहे,
असे म्हणण्याला आपल्याजवळ जवळजवळ आधारच उरला नसता!{} सुदैवाने हा
अंतःप्रज्ञ विचार बरोबर आहे. श्रेडर--बर्नस्टाइन प्रमेय त्याचे समर्थन
करते.
\end{explain}

\begin{thm}[श्रेडर--बर्नस्टाइन]
	\label{sfr:siz:sb:thm:schroder-bernstein}
	$\cardle{A}{B}$ आणि $\cardle{B}{A}$ असेल,
	तर $\cardeq{A}{B}$.
\end{thm}

\begin{explain}
दुसऱ्या शब्दांत, $A$ पासून~$B$ कडे जाणारे एकास-एक फलन आणि $B$
पासून~$A$ कडे जाणारे एकास-एक फलन असेल, तर $A$ पासून~$B$ कडे
जाणारे एकास-एक आच्छादन असते.

पण या निष्कर्षाची सिद्धता खरोखरच बरीच \emph{कठीण} आहे. कँटर यांनी
निष्कर्ष मांडला असला, तरी इतरांनी तो सिद्ध केला.
\footnote{इतिहासाच्या अधिक माहितीसाठी, उदाहरणार्थ,
Potter (2004, pp.~165--6) पाहा.}
%
आत्तासाठी हा निष्कर्ष विश्वासाने स्वीकारावा लागेल.

सुदैवाने श्रेडर--बर्नस्टाइन प्रमेय \emph{बरोबर} आहे आणि म्हणून आपण
परिभाषित केलेले $\cardeq{A}{B}$ आणि $\cardle{A}{B}$ हे संबंध
``आकारमाना''शी निगडित आहेत, असे मानणे समर्थित होते. शिवाय
श्रेडर--बर्नस्टाइन प्रमेय फार \emph{उपयुक्त} आहे. दोन तुल्यबल संचांमध्ये
एकास-एक आच्छादन शोधणे कठीण असू शकते. या प्रमेयामुळे तुलना दोन दिशांत
विभागता येते: दोन्ही दिशांतील एकास-एक फलन स्वतंत्रपणे शोधावी
लागतात.
\end{explain}








\section{प्रगणने आणि गणनीय संच}\label{sfr:siz:enm-alt:sec}

\begin{editorial}
  या विभागात संचसिद्धान्तात रूढ असलेल्या पद्धतीने, $\Nat$ च्या
  (आरंभीच्या खंडां)शी असलेली एकास-एक आच्छादने म्हणून प्रगणनांची
  व्याख्या केली आहे. त्यामुळे \ref{sfr:siz:enm:sec} मधील व्याख्यांशी तिचा
  किंचित विरोध होतो, आणि तेथील सर्व उदाहरणे येथे पुन्हा येतात. हा
  विभाग त्या विभागापेक्षा थोडा अधिक संक्षिप्तही आहे.
\end{editorial}

सांत संच त्यातील घटक केवळ क्रमाने मांडून सांगता येतो. पुढीलप्रमाणे
संचाची व्याख्या करताना आपण हेच करतो:
\[
  A = \{a_1, a_2, \ldots, a_n\}.
\]
घटक $a_1$, \dots, $a_n$ हे सर्व भिन्न आहेत असे गृहीत धरल्यास,
यातून $A$ आणि पहिल्या $n$ नैसर्गिक संख्या $0$, \dots, $n-1$ यांच्यात
एकास-एक आच्छादन मिळते. उलट, प्रत्येक सांत संचात सांत संख्येनेच
घटक असल्यामुळे प्रत्येक सांत संच अशी संगती लावून मांडता येतो.
दुसऱ्या शब्दांत, $A$ सांत असेल, तर $A$ आणि $\{0, \dots, n-1\}$
यांच्यात एकास-एक आच्छादन असते; येथे~$A$ मध्ये $n$ इतके घटक
असतात.

विशिष्ट प्रकारचे अनंत संच मान्य केल्यास, काही अनंत संचांचेही प्रगणन
करता येईल. अनंत संच आणि संपूर्ण~$\Nat$ यांच्यातील एकास-एक आच्छादन त्या
संचाचे प्रगणन करते, असे म्हणून हे नेमकेपणाने मांडता येते.

\begin{defn}[प्रगणनाची संचसैद्धान्तिक व्याख्या]
संच $A$ चे \emph{प्रगणन} म्हणजे असे एकास-एक आच्छादन की ज्याची व्याप्ती
$A$ असते आणि ज्याचा प्रांत नैसर्गिक संख्यांचा एखादा आरंभीचा खंड $\{0,
1, \ldots, n\}$ {किंवा} नैसर्गिक संख्यांचा संपूर्ण संच~$\Nat$ असतो.
\end{defn}

\begin{explain}
\emph{प्रगणन} या शब्दाच्या अशा वापरामागे सहज समजणारा आधार आहे.
संच $A$ चे प्रगणन केले आहे असे म्हणणे म्हणजे संच~$A$ मधील घटक मोजत
जाण्यास अनुमती देणारे एकास-एक आच्छादन $f$ आहे असे म्हणणे. $0$ वा घटक
$f(0)$, पहिला $f(1)$, \ldots आणि $n$ वा $f(n)$,~\ldots असतो.
\footnote{होय, आपण $0$ पासून मोजतो. अर्थात~$1$ पासूनही सुरुवात करता
येईल. त्यामुळे मोठा फरक पडणार नाही; फक्त~$\Nat$ ऐवजी~$\PosInt$
घ्यावा लागेल.} पुढील व्याख्या जोडल्यावर या मागचे कारण आणखी स्पष्ट होते:
\end{explain}

\begin{defn}
  \label{sfr:siz:enm-alt:defn:enumerable}
  संच~$A$ हा गणनीय असतो तेव्हा आणि केवळ तेव्हाच, जेव्हा
  $A = \emptyset$ असते किंवा~$A$ चे प्रगणन असते. $A$ हा
  अगणनीय आहे असे म्हणतो तेव्हा आणि केवळ तेव्हाच, जेव्हा
  $A$ हा गणनीय नसतो.
\end{defn}

\begin{explain}
म्हणून संच रिक्त असेल किंवा प्रगणन वापरून त्यातील घटक मोजता
येत असतील तेव्हा आणि केवळ तेव्हाच तो गणनीय असतो.
\end{explain}

\begin{ex}
नैसर्गिक संख्यांचे प्रगणन करणारे फलन म्हणजे फक्त $\Id{\Nat} \colon
\Nat \to \Nat$ हे $\Id{\Nat}(n) = n$ ने दिलेले एकरूपता फलन.
\emph{धन} नैसर्गिक संख्या $\Nat^+ = \Nat \setminus \{0\}$ यांचे
प्रगणन करणारे फलन $g(n) = n + 1$ आहे; म्हणजेच ते उत्तरवर्ती फलन आहे.
\end{ex}

\begin{prob}
संच $A$ हा गणनीय असतो तेव्हा आणि केवळ तेव्हाच $A =
\emptyset$ असते किंवा $f\colon \Nat \to A$ असे आच्छादन
असते, हे दाखवा. $A$ हा गणनीय असतो तेव्हा आणि केवळ तेव्हाच
$g\colon A \to \Nat$ असे एकास-एक फलन असते, हेही दाखवा.
\end{prob}

\begin{ex}
पुढीलप्रमाणे दिलेली $f\colon \Nat \to \Nat$ आणि $g \colon \Nat
\to \Nat$ ही फलने
\begin{align*}
  f(n) & = 2n \text{ आणि}\\
  g(n) & = 2n+1
\end{align*}
अनुक्रमे सम नैसर्गिक संख्यांचे आणि विषम नैसर्गिक संख्यांचे प्रगणन
करतात. पण यांपैकी कोणतेही फलन आच्छादक नाही; म्हणून कोणतेही
$\Nat$ चे प्रगणन नाही.
\end{ex}

\begin{prob}
वर्गसंख्या $1$, $4$, $9$, $16$, \dots{} यांचे प्रगणन परिभाषित करा.
\end{prob}

\begin{ex}
$\lceil x \rceil$ हे $x$ ला त्याच्यापेक्षा कमी नसलेल्या सर्वांत जवळच्या
पूर्णांकाकडे नेणारे \emph{ऊर्ध्व पूर्णांक} फलन असू द्या. मग पुढीलप्रमाणे
दिलेले फलन $f \colon \Nat \to \Int$:
\[
  f(n) = (-1)^{n} \left\lceil\tfrac{n}{2}\right\rceil
\]
पूर्णांकांचा संच~$\Int$ पुढीलप्रमाणे प्रगणित करते:
\[
\begin{array}{c c c c c c c c}
f(0) & f(1) & f(2) & f(3) & f(4) & f(5) & f(6) & \dots \\ \\
\big\lceil \tfrac{0}{2} \big\rceil & -\big\lceil \tfrac{1}{2}\big\rceil &  \big\lceil \tfrac{2}{2} \big\rceil & -\big\lceil \tfrac{3}{2} \big\rceil & \big\lceil \tfrac{4}{2} \big\rceil  & -\big\lceil \tfrac{5}{2}\big\rceil & \big\lceil \tfrac{6}{2} \big\rceil & \dots \\ \\
0 & -1 & 1 & -2 & 2 & -3 & 3& \dots
\end{array}
\]
धन आणि ऋण पूर्णांकांमध्ये आलटूनपालटून ``उड्या मारून'' $f$ हे $\Int$
ची मूल्ये कशी निर्माण करते ते पाहा. $f$ ची व्याख्या पुढीलप्रमाणे
प्रकरणे पाडून केली आहे असेही मानता येते:
\[
f(n) = \begin{cases}
  \frac{n}{2} & \text{जर $n$ सम असेल}\\
  -\frac{n+1}{2} & \text{जर $n$ विषम असेल}
  \end{cases}
\]
\end{ex}

\begin{prob}
$A$ आणि $B$ हे गणनीय असतील, तर $A \cup B$ सुद्धा गणनीय
आहे, हे दाखवा.
\end{prob}

\begin{prob}
$A_1$, $A_2$, \dots, $A_n$ हे सर्व गणनीय असतील, तर
$A_1 \cup \dots \cup A_n$ सुद्धा गणनीय आहे, हे $n$ वरील गणिती विगमनाने
दाखवा.
\end{prob}









\section{अगणनीय संच}\label{sfr:siz:nen-alt:sec}

\begin{editorial}
  या विभागात \ref{sfr:siz:enm-alt:sec} मधील व्याख्या वापरून
  $\Bin^\omega$ आणि $\Pow{\Nat}$ यांची अगणनीयता सिद्ध केली आहे;
  म्हणजे $\PosInt$ पासून आच्छादन मागण्याऐवजी~$\Nat$ शी एकास-एक
  आच्छादन आवश्यक मानले आहे.
\end{editorial}

\begin{explain}
नैसर्गिक संख्यांचा संच $\Nat$ हा अनंत आहे. तो सहजपणे
गणनीय देखील आहे. परंतु
\emph{अगणनीय} संच, म्हणजे गणनीय नसलेले संच,
अस्तित्वात आहेत ही लक्षणीय गोष्ट आहे
(पाहा \ref{sfr:siz:enm-alt:defn:enumerable}).

हे आश्चर्यकारक वाटू शकते. कारण $A$ हा अगणनीय आहे असे
म्हणणे म्हणजे \emph{कोणतेही} एकास-एक आच्छादन $f \colon \Nat \to A$
नाही असे म्हणणे; म्हणजे~$\Nat$ मधील अनंतपणे अनेक घटक~$A$
मध्ये पाठवणारे कोणतेही फलन संपूर्ण~$A$ व्यापत नाही. म्हणून $A$ हा
अगणनीय असेल, तर नैसर्गिक संख्यांपेक्षा~$A$ मध्ये ``अधिक''
घटक असतात.

संच अगणनीय आहे हे सिद्ध करण्यासाठी अनुरूप एकास-एक आच्छादन
अस्तित्वात असू शकत नाही हे दाखवावे लागते. याचा सर्वोत्तम मार्ग म्हणजे
~$A$ मधील घटक प्रगणित करण्याचा प्रत्येक प्रयत्न किमान एक
घटक वगळतो हे दाखवणे; यावरून कोणतेही फलन $f\colon \Nat \to A$
हे आच्छादक नाही असे दिसते. हे प्रस्थापित करण्याची एक सर्वसाधारण
युक्ती म्हणजे कँटर यांची \emph{कर्णरेषीय पद्धत} वापरणे. $A$ मधील
घटक $x_1$, $x_2$, \dots{} असे मांडले असता, रचनेमुळे त्या
यादीत असूच शकणार नाही असा~$A$ चा आणखी एक घटक आपण बनवतो.

हे सर्व उदाहरणाने उत्तम समजते. म्हणून आपले पहिले उदाहरण म्हणजे $0$
आणि $1$ यांच्या सर्व अनंत चिन्हमालांचा संच~$\Bin^\omega$. (`$\Bin$'
हे द्विमान दर्शवते आणि त्याचा अर्थ फक्त दोन घटकांचा संच $\{0,1\}$
असा घेता येतो.)
तरी सध्या या किंचित शिथिल मांडणीमुळे गोंधळ होऊ नये.
\end{explain}

\begin{thm}
\label{sfr:siz:nen-alt:thm:nonenum-bin-omega}
$\Bin^\omega$~हा अगणनीय आहे.
\end{thm}

\begin{proof}
$\Bin^\omega$ च्या एखाद्या उपसंचाचे कोणतेही प्रगणन घ्या. म्हणजे
$s_{0}$, $s_{1}$, $s_{2}$, \dots{} अशी यादी आपल्याकडे आहे, जेथे
प्रत्येक $s_n$ ही $0$ आणि~$1$ यांची अनंत चिन्हमाला आहे. या यादीतील
$n$ व्या चिन्हमालेतील $m$ वा अंक $s_n(m)$ असू द्या.
\begin{quote}\small\textbf{स्रोतदुरुस्ती.} OLSIZ-008: गोठवलेल्या इंग्रजी वर्णनात s\_n(m) ला m व्या चिन्हमालेतील n वा अंक म्हटले आहे; लगेचचे वाक्य आणि सारणी मात्र तो n व्या ओळीतील, म्हणजे n व्या चिन्हमालेतील, m वा अंक असल्याचे निश्चित करतात. मराठी मजकुरात ओळ-स्तंभाशी सुसंगत निर्देशांकक्रम वापरला आहे.\end{quote}
मग यादीकडे सरणी म्हणून पाहता येते, ज्यात $s_n(m)$ हा $n$ व्या ओळीत
आणि $m$ व्या स्तंभात ठेवलेला आहे:
\[
\begin{array}{c|c|c|c|c|c}
& 0 & 1 & 2 & 3 & \dots \\\hline
0 & \mathbf{s_{0}(0)} & s_{0}(1) & s_{0}(2) & s_0(3) & \dots \\\hline
1 & s_{1}(0)& \mathbf{s_{1}(1)} & s_1(2) & s_1(3) & \dots \\\hline
2 & s_{2}(0)& s_{2}(1) & \mathbf{s_2(2)} & s_2(3) & \dots \\\hline
3 & s_{3}(0)& s_{3}(1) & s_3(2) & \mathbf{s_3(3)} & \dots \\\hline
\vdots & \vdots & \vdots & \vdots & \vdots & \mathbf{\ddots}
\end{array}
\]
आता या यादीत नसलेली $0$ आणि $1$ यांची $d$ ही अनंत चिन्हमाला आपण
बनवू. तिची प्रत्येक नोंद सांगून, म्हणजे सर्व $n \in \Nat$ साठी
$d(n)$ सांगून, हे करू. सहज कल्पना अशी: वरील सरणीचा कर्ण खाली वाचायचा
(त्यामुळे ``कर्णरेषीय पद्धत'' हे नाव) आणि प्रत्येक $1$ चा $0$, तर
प्रत्येक $0$ चा~$1$ करायचा.
\begin{quote}\small\textbf{स्रोतदुरुस्ती.} OLSIZ-009: गोठवलेल्या इंग्रजी सूचनेत “प्रत्येक 1 चा 0” हेच परिवर्तन दोनदा आले आहे. लगेचचे प्रकरणनिहाय सूत्र परस्परपूरक दोन्ही बदल निश्चित करते: 1 चा 0 आणि 0 चा 1. मराठी मजकुरात हे दोन्ही बदल स्पष्ट केले आहेत.\end{quote}
अधिक अमूर्तपणे, कर्णावरील $n$ वा घटक, $s_n(n)$, हा $1$ किंवा
$0$ आहे त्यानुसार $d(n)$ हा $0$ किंवा $1$ असावा अशी व्याख्या करतो;
म्हणजे:
\[
d(n) =
\begin{cases}
1 & \text{जर $s_{n}(n) = 0$ असेल}\\
0 & \text{जर $s_{n}(n) = 1$ असेल}
\end{cases}
\]
स्पष्टच, $d \in \Bin^\omega$, कारण ती $0$ आणि $1$ यांची अनंत
चिन्हमाला आहे. पण आपण $d$ अशी बनवली की कोणत्याही $n \in \Nat$
साठी $d(n) \neq s_n(n)$. म्हणजे $d$ ही $s_n$ पेक्षा तिच्या $n$ व्या
नोंदीत वेगळी आहे. म्हणून कोणत्याही $n\in \Nat$ साठी $d \neq s_n$.
त्यामुळे $d$ ही $s_0$, $s_1$, $s_2$, \dots{} या यादीत असू शकत नाही.

$\Bin^\omega$ च्या एखाद्या उपसंचाचे स्वैर प्रगणन दिले असता ते
$\Bin^\omega$ चा कोणता तरी घटक वगळेल हे आपण दाखवले आहे.
म्हणून $\Bin^\omega$ या संचाचे प्रगणन नाही; म्हणजे $\Bin^\omega$ हा
अगणनीय आहे.
\end{proof}

\begin{explain}
या सिद्धतापद्धतीला ``कर्णीकरण'' म्हणतात, कारण ती~$d$ ची व्याख्या
करण्यासाठी सरणीचा कर्ण वापरते. मात्र कर्णीकरणात सरणी असणे आवश्यक नाही.
प्रत्यक्षात, सरणी किंवा प्रत्यक्ष कर्ण नसतानाही हिच्याशी मिळतीजुळती कल्पना
वापरून एखादा संच अगणनीय आहे हे दाखवता येते. पुढील फलित हे
कसे करता येते ते दाखवते.
\end{explain}

\begin{thm}
\label{sfr:siz:nen-alt:thm:nonenum-pownat}
$\Pow{\Nat}$ हा गणनीय नाही.
\end{thm}

\begin{proof}
~$\Nat$ च्या उपसंचांची प्रत्येक यादी $\Nat$ चा कोणता तरी उपसंच वगळते
हे दाखवून आपण त्याच पद्धतीने पुढे जाऊ. म्हणून $\Nat$ च्या उपसंचांची
$N_0, N_1, N_2, \ldots$ अशी एखादी यादी आहे असे समजा. $D$ या संचाची
व्याख्या पुढीलप्रमाणे करतो: $n \in D$ तेव्हा आणि केवळ तेव्हाच, जेव्हा
$n \notin N_{n}$:
\[
D = \Setabs{n \in \Nat}{n \notin N_n}
\]
स्पष्टच $D\subseteq \Nat$. पण $D$ हा यादीत असू शकत नाही. कारण
रचनेनुसार $n \in D$ तेव्हा आणि केवळ तेव्हाच, जेव्हा $n\notin N_n$;
म्हणून कोणत्याही $n \in \Nat$ साठी $D \neq N_n$.
\end{proof}

\begin{explain}
आधीच्या सिद्धतेत कर्णाचा उल्लेख नव्हता. तरी पुढीलप्रमाणे चित्र डोळ्यांसमोर
आणल्यास त्यात कर्ण आहे असे मानता येते: $N_0$, $N_1$, \dots{} हे संच
एका सरणीत लिहिल्याची कल्पना करा; $m \in N_n$ असेल तेव्हा आणि केवळ
तेव्हाच $m$ व्या स्तंभात $m$ लिहून $N_n$ हा $n$ व्या ओळीत लिहायचा.
उदाहरणार्थ, त्या यादीतील पहिले चार संच $\{0,1,2,\dots\}$,
$\{1, 3, 5, \dots\}$, $\{0,1,4\}$ आणि $\{2,3,4,\dots\}$ आहेत असे
समजा; मग आपल्या सरणीची सुरुवात अशी असेल:
\[
\begin{array}{r@{}rrrrrrr}
  N_0 = \{ & \mathbf{0}, & 1, & 2, & & & & \dots\}\\
  N_1 = \{ &  & \mathbf{1}, &  & 3, &  & 5, & \dots\}\\
  N_2 = \{ & 0, & 1, &  &  & 4\phantom{,} &  & \}\\
  N_3 = \{ &  &  & 2, & \mathbf{3}, & 4, & & \dots\}\\
  &\vdots & & & & & & \ddots\phantom{\}}
  \end{array}
\]
मग कर्णाच्या खाली जात, $n$ कर्णावर \emph{नसेल} तेव्हा आणि केवळ
तेव्हाच $n \in D$ असे ठेवून $D$ हा संच मिळतो. म्हणून वरील उदाहरणात
आपण $0$ आणि $1$ वगळू,~$2$ समाविष्ट करू,~$3$ वगळू, इत्यादी.
\end{explain}

\begin{prob}
सर्व फलने $f \colon \Nat \to \Nat$ यांचा संच स्पष्ट कर्णरेषीय
युक्तिवादाने अगणनीय आहे हे दाखवा. म्हणजे $f_1$, $f_2$,
\dots{} ही फलनांची यादी असून प्रत्येक $f_i\colon \Nat \to \Nat$
असेल, तर या यादीत नसलेले कोणते तरी $g \colon \Nat \to \Nat$ असते,
हे दाखवा.
\end{prob}









\section{न्यूनीकरण}\label{sfr:siz:red-alt:sec}

\begin{editorial}
  या विभागात \ref{sfr:siz:nen-alt:sec} मधील फलितांशी जुळणारी अगणनीयता
  न्यूनीकरणाने सिद्ध केली आहे. \ref{sfr:siz:nen:sec} मधील फलितांशी जुळणारी
  किंचित अधिक विस्तृत पर्यायी आवृत्ती \ref{sfr:siz:red:sec} मध्ये दिली आहे.
\end{editorial}

कर्णीकरणाच्या युक्तिवादाने $\Bin^\omega$ हा अगणनीय आहे,
हे आपण सिद्ध केले. तसाच कर्णीकरणाचा युक्तिवाद वापरून $\Pow{\Nat}$ हा
अगणनीय आहे हेही दाखवले. पण $\Pow{\Nat}$ हा
अगणनीय असल्याचे सिद्ध करण्याचा आणखी एक मार्ग आहे:
\emph{$\Pow{\Nat}$ हा गणनीय असेल, तर $\Bin^\omega$ हाही
गणनीय असतो} हे दाखवा. $\Bin^\omega$ हा अगणनीय
आहे हे आपल्याला माहीत असल्यामुळे $\Pow{\Nat}$ हाही तसाच आहे असा
निष्कर्ष निघेल.

याला एक समस्या दुसऱ्या समस्येकडे \emph{न्यूनीत करणे} म्हणतात. या
बाबतीत $\Bin^\omega$ चे प्रगणन करण्याची समस्या आपण $\Pow{\Nat}$ चे
प्रगणन करण्याच्या समस्येकडे न्यूनीत करतो. नंतरच्या समस्येचे उत्तर---
$\Pow{\Nat}$ चे प्रगणन---पहिल्या समस्येचे उत्तर---$\Bin^\omega$ चे
प्रगणन---देईल.

संच~$B$ च्या प्रगणनाची समस्या संच~$A$ च्या प्रगणनाच्या समस्येकडे
न्यूनीत करण्यासाठी~$A$ चे प्रगणन~$B$ च्या प्रगणनात रूपांतरित करण्याची
पद्धत आपण देतो. हे करण्याचा सर्वांत सोपा मार्ग म्हणजे
$f\colon A \to B$ असे आच्छादन परिभाषित करणे. $x_1$, $x_2$,
\dots{} हे~$A$ चे प्रगणन असेल, तर $f(x_1)$, $f(x_2)$, \dots{} हे~$B$
चे प्रगणन करील. आपल्या बाबतीत आपण
$f\colon \Pow{\Nat} \to \Bin^\omega$ असे आच्छादन शोधत आहोत.

\begin{prob}
$g\colon B \to A$ असे एकास-एक फलन असेल आणि $B$ हा
अगणनीय असेल, तर $A$ हाही तसाच असतो, हे दाखवा. $g$ वापरून
~$A$ चे प्रगणन~$B$ च्या प्रगणनात कसे रूपांतरित करता येते हे दाखवून
सिद्धता द्या.
\end{prob}

\begin{proof}[न्यूनीकरणाने {\ref{sfr:siz:nen-alt:thm:nonenum-pownat}} ची सिद्धता]
न्यूनीकरणासाठी, $\Pow{\Nat}$ हा गणनीय आहे आणि त्यामुळे त्याचे
$N_{1}$, $N_{2}$, $N_{3}$, \dots{} असे प्रगणन आहे, असे समजा.

$f \colon \Pow{\Nat} \to \Bin^\omega$ हे फलन पुढीलप्रमाणे परिभाषित
करा: $f(N)$ ही अशी चिन्हमाला~$s$ असू द्या की $s(n) = 1$ तेव्हा आणि
केवळ तेव्हाच, जेव्हा $n \in N$; अन्यथा $s(n) = 0$.
\begin{quote}\small\textbf{स्रोतदुरुस्ती.} OLSIZ-010: गोठवलेल्या इंग्रजी व्याख्येत f(N) ला s\_k असे म्हटले आहे, पण k बांधलेला किंवा N वरून ठरवलेला नाही. घटकनिहाय समीकरण N ची अभिलक्षण-चिन्हमाला एकमेव ठरवते; मराठी व्याख्येत एकच बद्ध निर्गत नाव s आणि त्याचे घटक s(n) सातत्याने वापरले आहेत.\end{quote}

याने फलन स्पष्टपणे परिभाषित होते, कारण $N \subseteq \Nat$ असताना
कोणताही $n \in \Nat$ हा $N$ चा घटक असतो किंवा नसतो.
उदाहरणार्थ, सम नैसर्गिक संख्यांचा संच
$2\Nat = \Setabs{2n}{n \in \Nat} = \{0,2, 4, 6, \dots\}$ हा
$1010101\dots$ या चिन्हमालेकडे; $\emptyset$ हा $0000\dots$ कडे; आणि
$\Nat$ हा $1111\dots$ कडे पाठवला जातो.

हे फलन आच्छादक देखील आहे: $0$ आणि $1$ यांची प्रत्येक चिन्हमाला
नैसर्गिक संख्यांच्या कोणत्या तरी संचाशी अनुरूप असते; तो संच म्हणजे
चिन्हमालेत ज्या स्थानांवर~$1$ आहे त्यांना अनुरूप नैसर्गिक संख्या
सदस्य म्हणून असलेला संच. अधिक नेमकेपणाने, $s \in \Bin^\omega$ असेल,
तर पुढीलप्रमाणे $N \subseteq \Nat$ परिभाषित करा:
\[
N = \Setabs{n \in \Nat}{s(n) = 1}
\]
मग~$f$ ची व्याख्या पाहून $f(N) = s$ हे तपासता येते.

आता पुढील यादी विचारात घ्या:
\[
f(N_1), f(N_2), f(N_3), \dots
\]
$f$ हे आच्छादक असल्यामुळे $\Bin^\omega$ चा प्रत्येक सदस्य
हा~$f$ चे कोणत्या तरी फलसाधकावरील मूल्य असला पाहिजे आणि म्हणून यादीत
आला पाहिजे. त्यामुळे ही यादी संपूर्ण~$\Bin^\omega$ चे प्रगणन करते.

म्हणून $\Pow{\Nat}$ हा गणनीय असता, तर $\Bin^\omega$ हाही
गणनीय असता. पण $\Bin^\omega$ हा अगणनीय आहे
(\ref{sfr:siz:nen-alt:thm:nonenum-bin-omega}). त्यामुळे $\Pow{\Nat}$ हा
अगणनीय आहे.
\end{proof}


\begin{prob}\label{sfr:siz:red:prob:nat-nat}
  $f\colon \Nat \to \Nat$ अशा सर्व फलनांचा संच~$X$ हा
  न्यूनीकरणाच्या युक्तिवादाने अगणनीय आहे, हे दाखवा.
  (सूचना: $X$ पासून~$\Bin^\omega$ कडे जाणारे आच्छादक फलन द्या.)
\end{prob}

\begin{prob}
नैसर्गिक संख्यांच्या जोड्यांच्या \emph{सर्व संचांचा} संच, म्हणजे
$\Pow{\Nat \times \Nat}$, न्यूनीकरणाच्या युक्तिवादाने
अगणनीय आहे, हे दाखवा.
\end{prob}

\begin{prob}
नैसर्गिक संख्यांच्या अनंत अनुक्रमांचा संच $\Nat^\omega$ हा
न्यूनीकरणाच्या युक्तिवादाने अगणनीय आहे, हे दाखवा.
\end{prob}


\begin{prob}
$S$ हा $\Nat$ पासून $\{0,1\}$ या संचाकडे जाणारी सर्व
आच्छादक फलने असलेला संच असू द्या; म्हणजे $S$ मध्ये
$f \colon \Nat \to \Bin$ अशी सर्व आच्छादक फलने आहेत. $S$ हा
अगणनीय आहे, हे दाखवा.
\end{prob}

\begin{prob}
सर्व वास्तव संख्यांचा संच~$\Real$ हा अगणनीय आहे, हे दाखवा.
\end{prob}



\chapter{अंकगणितीकरण}\label{sfr:arith::chap}
\begin{editorial}
या प्रकरणातील सामग्री अनौपचारिक संचसिद्धान्तात संख्याप्रणालींची रचना
मांडते. ती टिम बटन यांच्या Open Set Theory या ग्रंथातून घेतली आहे.
\end{editorial}






\section{$\Nat$ पासून $\Int$ कडे}\label{sfr:arith:int:sec}

येथे दोन मूलभूत निरीक्षणे आहेत:
	\begin{enumerate}
		\item प्रत्येक पूर्णांक $n - m$ या रूपात लिहिता येतो, जेथे $n, m \in\Nat$.
		\item $n - m$ या व्यंजकात सांकेतिक केलेली माहिती $\tuple{n, m}$ या क्रमित जोडीद्वारेही तितकीच सांकेतिक करता येते.
	\end{enumerate}
नैसर्गिक संख्यांच्या क्रमित जोड्या $\Nat^2$ चे घटक आहेत हे आपल्याला आधीच माहीत आहे. आणि आपण $\Nat$ समजतो असे गृहीत धरत आहोत. त्यामुळे या दोन निरीक्षणांवर आधारलेली एक भोळी सूचना पुढे करता येईल: \emph{नैसर्गिक संख्यांच्या क्रमित जोड्यांना पूर्णांक म्हणून हाताळू या}.

प्रत्यक्षात ही सूचना अतिशय भोळी आहे. अर्थात, $0- 2 = 4 - 6$ असले पाहिजे असे आपल्याला हवे आहे. पण स्पष्टपणे $\tuple{0, 2 }\neq \tuple{4, 6}$. त्यामुळे $\Nat^2$ हाच पूर्णांकांचा संच आहे असे सरळ म्हणता येत नाही.

मागील अडचणीचे सामान्यीकरण केल्यास आपल्याला पुढील अट हवी आहे:
	$$a - b = c - d \text{ तेव्हा आणि केवळ तेव्हाच }a + d = c + b$$
(पूर्णांकांचे वर्तन \emph{असेच अभिप्रेत} आहे हे स्पष्ट असावे: दोन्ही बाजूंना फक्त $b$ आणि $d$ मिळवा.) हे वर्तन निश्चित करण्याचा सोपा मार्ग म्हणजे क्रमित जोड्यांमधील $\Intequiv$ हा तुल्यता संबंध पुढीलप्रमाणे व्याख्यायित करणे:
$$\tuple{ a, b } \Intequiv \tuple{c, d}\text{ तेव्हा आणि केवळ तेव्हाच }a + d = c + b$$
आता हा तुल्यता संबंध आहे हे दाखवावे लागेल.
\begin{prop} $\Intequiv$ हा तुल्यता संबंध आहे.
\end{prop}
	\begin{proof}
	$\Intequiv$ परावर्ती, सममित आणि संक्रमक आहे हे दाखवले पाहिजे.

	\emph{परावर्तिता:} $a + b = b + a$ असल्यामुळे स्पष्टपणे $\tuple{a, b} \Intequiv \tuple{a, b}$.

	\emph{सममितता:} $\tuple{a, b} \Intequiv \tuple{c, d}$ असे समजा; म्हणजे $a + d = c + b$. मग $c + b = a + d$, त्यामुळे $\tuple{c, d} \Intequiv \tuple{a, b}$.

	\emph{संक्रमकता:} $\tuple{a, b} \Intequiv \tuple{c, d}\Intequiv \tuple{m, n}$ असे समजा. त्यामुळे $a + d = c + b$ आणि $c + n = m + d$. म्हणून $a + d + c + n = c + b + m + d$, व म्हणून $a + n = m + b$. परिणामी $\tuple{a, b } \Intequiv \tuple{m, n}$.
\end{proof}

आता या तुल्यता संबंधाचा वापर करून तुल्यतावर्ग घेता येतात:

\begin{defn}
		नैसर्गिक संख्यांच्या क्रमित जोड्यांचे $\Intequiv$ अंतर्गत तुल्यतावर्ग म्हणजे पूर्णांक; म्हणजेच $\Int = \equivclass{\Nat^2}{\Intequiv}$.
\end{defn}

या संकेतजनक व्याख्येबद्दल अनेक निरनिराळ्या \emph{तात्त्विक} प्रतिक्रिया असू शकतात. त्या प्रतिक्रियांचा विचार करण्यापूर्वी मात्र काही तांत्रिक तपशील पुढे नेणे उपयुक्त ठरेल.

पूर्णांक म्हणजे काय हे ठरवल्यावर त्यांच्यावरील मूलभूत फलने आणि संबंध व्याख्यायित करावे लागतील. ज्या $\Intequiv$ अंतर्गत तुल्यतावर्गात $\tuple{m, n}$ हा घटक आहे, त्या वर्गासाठी $\equivrep{m, n}{\Intequiv}$ असे लिहू या.\footnote{टीप: \ref{sfr:rel:eqv:def:equivalenceclass} मध्ये दिलेले संकेतन वापरल्यास त्याच वर्गासाठी आपण $\equivrep{\tuple{m,n}}{\Intequiv}$ असे लिहिले असते. पण ते वाचायला जरा कठीण आहे.} म्हणजे:
	$$\equivrep{m, n}{\Intequiv} = \Setabs{\tuple{a, b} \in \Nat^2}{\tuple{a, b}\Intequiv \tuple{m, n}}$$
आता काही व्याख्या देऊ:
	\begin{align*}
		\equivrep{a, b}{\Intequiv} + \equivrep{c, d}{\Intequiv} &= \equivrep{a + c, b + d}{\Intequiv}\\
		\equivrep{a, b}{\Intequiv} \times \equivrep{c, d}{\Intequiv} &= \equivrep{a c + b  d, a  d + b c}{\Intequiv}\\
		\equivrep{a, b}{\Intequiv} \leq \equivrep{c, d}{\Intequiv} &\text{ तेव्हा आणि केवळ तेव्हाच }a + d \leq b + c
	\end{align*}
(प्रथेप्रमाणे, स्वयंसिद्धे वाचायला सोपी व्हावीत म्हणून मी `$(a \times b)$' या अर्थी `$ab$' लिहित आहे.) आता या व्याख्या \emph{अपेक्षेप्रमाणे} वागतात याची खात्री केली पाहिजे. याचा नेमका अर्थ उलगडून संपूर्ण पडताळणी करणे बरेच कष्टाचे आहे; तपशील आपण \ref{sfr:arith:check:sec} कडे सोपवतो. पण थोडक्यात: सर्व काही व्यवस्थित चालते!

आता फक्त एक गोष्ट उरते. आपण नैसर्गिक संख्यांच्या साहाय्याने
पूर्णांकांची रचना केली आहे. पण त्यामुळे नैसर्गिक संख्या \emph{स्वतः
पूर्णांक नसतील}. याचे तात्त्विक महत्त्व आपण \ref{sfr:arith:ref:sec} मध्ये
पुन्हा पाहू. मात्र पूर्णपणे तांत्रिक दृष्टीने, नैसर्गिक संख्यांना
पूर्णांक \emph{म्हणून} हाताळण्याची काही पद्धत आपल्याला हवी. कल्पना अगदी
सोपी आहे: प्रत्येक $n \in \Nat$ साठी $n_\Int = \equivrep{n, 0}{\Intequiv}$
असे आपण ठरवतो. ही व्याख्या सुयोग्य रीतीने वागते याची, म्हणजे कोणत्याही
$m, n \in \Nat$ साठी पुढील अटी पूर्ण होतात याची खात्री केली पाहिजे:
\begin{align*}
	(m + n)_\Int &= m_\Int + n_\Int\\
	(m \times n)_\Int &= m_\Int \times n_\Int\\
	m \leq n &\liff m_\Int \leq n_\Int
\end{align*}
ही सर्व पडताळणी अगदी सरळ आहे. उदाहरणार्थ, दुसरी अट खरी आहे हे
दाखवताना नैसर्गिक संख्यांच्या परिचित वर्तनाचा थेट वापर करून पुढीलप्रमाणे
युक्तिवाद करता येतो:
	\begin{align*}
		(m \times n)_\Int  &= \equivrep{m \times n, 0}{\Intequiv} \\
		&= \equivrep{m \times n + 0 \times 0, m \times 0 + 0 \times n}{\Intequiv} \\
		&= \equivrep{m, 0}{\Intequiv} \times \equivrep{n, 0}{\Intequiv} \\
		&= m_\Int \times n_\Int
	\end{align*}
इतर दोन अटी पडताळण्याचे काम सरावासाठी सोडतो.
\begin{prob}
	कोणत्याही $m, n \in \Nat$ साठी $(m + n)_\Int = m_\Int + n_\Int$ आणि $m \leq n \liff m_\Int \leq n_\Int$ हे दाखवा.
\end{prob}







\section{$\Int$ पासून $\Rat$ कडे}\label{sfr:arith:rat:sec}

अनौपचारिक संचसिद्धान्त वापरून नैसर्गिक संख्यांपासून पूर्णांकांची रचना कशी
करायची हे आपण नुकतेच पाहिले. आता अगदी त्याच प्रकारे पूर्णांकांपासून
परिमेय संख्यांची रचना कशी करायची ते पाहू. सुरुवातीची निरीक्षणे अशी:
\begin{enumerate}
	\item प्रत्येक परिमेय संख्या $\nicefrac{i}{j}$ या रूपात लिहिता येते,
	जेथे $i$ आणि $j$ दोन्ही पूर्णांक असतात, पण $j$ शून्येतर असतो.
	\item $\nicefrac{i}{j}$ या व्यंजकात सांकेतिक केलेली माहिती
	$\tuple{ i, j}$ या क्रमित जोडीद्वारेही तितकीच सांकेतिक करता येते.
\end{enumerate}
यावरून परिमेय संख्यांकडे $\Int \times (\Int \setminus \{0_\Int\})$
मधून घेतलेल्या क्रमित जोड्या \emph{म्हणून} पाहणे हा स्वाभाविक मार्ग
वाटेल. मात्र आधीप्रमाणे ही कल्पना जरा अतिभोळी ठरेल, कारण आपल्याला
$\nicefrac{3}{2} = \nicefrac{6}{4}$ हवे आहे, पण $\tuple{ 3, 2}\neq \tuple{ 6,
4}$. अधिक सर्वसाधारणपणे आपल्याला पुढील अट हवी:
\[
	\nicefrac{a}{b} = \nicefrac{c}{d} \text{ तेव्हा आणि केवळ तेव्हाच } a \times d = b \times c
\]
ही अट मिळवण्यासाठी $\Int \times (\Int \setminus \{0_\Int\})$ वरील
{तुल्यता संबंध} आपण पुढीलप्रमाणे व्याख्यायित करतो:
\[
	\tuple{ a, b }\Ratequiv \tuple{ c, d} \text{ तेव्हा आणि केवळ तेव्हाच }a \times d = b \times c
\]
हा तुल्यता संबंध आहे हे तपासले पाहिजे. ही तपासणी बरीचशी $\Intequiv$
च्या प्रकरणासारखीच आहे आणि ती आपण सरावासाठी सोडू.
\begin{prob}
$\Ratequiv$ हा तुल्यता संबंध आहे हे दाखवा.
\end{prob}
त्यामुळे पुढील व्याख्या देता येते:
\begin{defn}
दुसरा घटक शून्येतर असलेल्या पूर्णांकांच्या जोड्यांचे $\Ratequiv$ अंतर्गत
तुल्यतावर्ग म्हणजे परिमेय संख्या. म्हणजेच $\Rat =
\equivclass{(\Int \times (\Int\setminus \{0_\Int\}))}{\Ratequiv}$.
\end{defn}

पूर्णांकांप्रमाणे येथेही काही मूलभूत क्रिया व्याख्यायित करायच्या आहेत.
$\tuple{i, j}$ हा घटक असलेला $\Ratequiv$ अंतर्गत तुल्यतावर्ग
$\equivrep{i,j}{\Ratequiv}$ असेल, तर आपण पुढील व्याख्या करतो:
\begin{align*}
	\equivrep{a, b}{\Ratequiv} + \equivrep{c, d}{\Ratequiv} &= \equivrep{ad + bc,  bd}{\Ratequiv}\\
	\equivrep{a, b}{\Ratequiv} \times \equivrep{c, d}{\Ratequiv} &= \equivrep{a  c, b d}{\Ratequiv}.
\intertext{या परिमेय संख्यांवर $r \leq s$ ची व्याख्या करण्यासाठी आपण पुढील तथ्य वापरतो: $r \le s$ तेव्हा आणि केवळ तेव्हाच $s - r$ हा फरक ऋण नसतो; म्हणजे $s - r$ हा $\nicefrac{i}{j}$ या रूपात लिहिता येतो, जेथे $i$ ऋणेतर आणि $j$~धन असतो:}
	\equivrep{a, b}{\Ratequiv} \leq \equivrep{c, d}{\Ratequiv} &\text{ तेव्हा आणि केवळ तेव्हाच }
	\equivrep{c, d}{\Ratequiv} - \equivrep{a, b}{\Ratequiv} =
	\equivrep{i_\Int, j_\Int}{\Ratequiv}
\end{align*}
\begin{quote}\small\textbf{स्रोतदुरुस्ती.} MRARITH-001: गोठवलेल्या इंग्रजी स्पष्टीकरणात “s-r is not negative” नंतर चुकून r-s ला ऋणेतर अपूर्णांकाच्या रूपात लिहिले आहे. लगेचची प्रदर्शित व्याख्या पुन्हा s-r वापरते. मराठी मजकुरात या दोन्ही नियंत्रक ठिकाणांशी सुसंगत s-r वापरले आहे.\end{quote}
असे काही $i \in \Nat$ आणि $0 \neq j \in \Nat$ असले पाहिजेत.

त्यानंतर या व्याख्या \emph{अपेक्षेप्रमाणे} वागतात हे तपासावे लागेल;
ही तपासणी आपण \ref{sfr:arith:check:sec} कडे सोपवतो. आणि त्या खरोखर तशाच
वागतात!{} शेवटी, पूर्णांकांना परिमेय संख्या \emph{म्हणून} हाताळण्याची
काही पद्धत हवी; म्हणून प्रत्येक $i \in \Int$ साठी $i_\Rat =
\equivrep{i, 1_\Int}{\Ratequiv}$ असे आपण ठरवतो. हे सर्व सुयोग्य रीतीने
वागते याची तपासणीही आपण \ref{sfr:arith:check:sec} मध्ये करतो.

\begin{prob}
कोणत्याही $i, j \in \Int$ साठी $(i + j)_\Rat = i_\Rat+ j_\Rat$ आणि $(i \times j)_\Rat =
i_\Rat \times j_\Rat$ आणि $i \leq j \liff i_\Rat \leq j_\Rat$ हे दाखवा.
\end{prob}








\section{वास्तव संख्या-रेषा}\label{sfr:arith:real:sec}

पुढील पायरी म्हणजे परिमेय संख्यांपासून वास्तव संख्यांची रचना कशी करायची
हे दाखवणे. त्याआधी वास्तव संख्यांचे \emph{वैशिष्ट्य} नेमके काय आहे हे
समजून घ्यावे लागेल.

वास्तव संख्या परिमेय संख्यांसारख्याच बऱ्याच अंशी वागतात. (तांत्रिक
भाषेत, दोन्ही \emph{क्रमित क्षेत्रांची} उदाहरणे आहेत; त्याची व्याख्या
\ref{sfr:arith:check:orderedfield} मध्ये पाहा.) आता, तुम्ही \ref{sfr:siz::chap}
मधील सराव सोडवले असतील, तर वास्तव संख्या परिमेय संख्यांपेक्षा काटेकोरपणे
अधिक आहेत, म्हणजे $\cardless{\Rat}{\Real}$, हे तुम्हाला माहीत असेल.
हे कँटर यांनी प्रथम सिद्ध केले. परंतु अपरिमेय संख्या, म्हणजे परिमेय
नसलेल्या वास्तव संख्या, अस्तित्वात आहेत हे सुमारे अडीच सहस्रकांपासून
माहीत आहे. खरे तर:
\begin{quote}\small\textbf{स्रोतदुरुस्ती.} MRARITH-002: गोठवलेल्या इंग्रजीतील निष्क्रिय केलेली तळटीप $\sqrt{2}$ हे धन आणि ऋण अशी दोन्ही मुळे दर्शवते असे म्हणते. पारंपरिक मूलचिन्ह मात्र मुख्य, ऋणेतर वर्गमूळ दर्शवते; म्हणून मराठीतील निष्क्रिय तळटीप दुरुस्त केली आहे. सक्रिय वाचनीय मजकुरावर याचा परिणाम नाही.\end{quote}
\begin{thm}\label{sfr:arith:real:root2irrational}
$\sqrt{2}$ ही परिमेय संख्या नाही; म्हणजे $\sqrt{2} \notin \Rat$.
\end{thm}

\begin{proof}
व्याघातासाठी $\sqrt{2}$ ही परिमेय संख्या आहे असे समजा. मग काही नैसर्गिक
संख्या $m$ आणि $n$ यांसाठी $\sqrt{2} = \nicefrac{m}{n}$. हा अपूर्णांक
यापुढे संक्षिप्त करता येणार नाही अशा रीतीने आपण $m$ आणि $n$ निवडू शकतो.
पुनर्मांडणी केल्यास $m^{2} = 2n^{2}$. येथून सिद्धता दोन प्रकारे पूर्ण
करता येते:

\emph{प्रथम, भूमितीय रीतीने} (टेनेनबाउम यांच्या पद्धतीने).\footnote{ही
सिद्धता Conway (2006) यांनी नोंदवली आहे.} पुढील चौरस पाहा:
\begin{center}
	\begin{tikzpicture}
		\draw[thick] (0,0) rectangle (3,3);
		\draw[thick, fill=red!50] (0,0) rectangle (2.3,2.3);
		\draw[thick, fill=yellow!50] (0.7,0.7) rectangle (3, 3);
		\draw[thick, fill=orange!50] (0.7,0.7) rectangle (2.3, 2.3);
		\draw[<->] (4, 0.7)--(4, 3);
		\node at (4.25, 1.85) (n) {$n$};
		\draw[<->] (5, 0)--(5, 3);
		\node at (5.25, 1.5) (m) {$m$};
	\end{tikzpicture}
\end{center}
$m^2 = 2n^2$ असल्यामुळे $n$ बाजू असलेले दोन चौरस ज्या प्रदेशात
एकमेकांवर येतात त्याचे क्षेत्रफळ, या दोन्ही चौरसांपैकी एकानेही न झाकलेल्या
प्रदेशाच्या क्षेत्रफळाइतके आहे; म्हणजे नारिंगी चौरसाचे क्षेत्रफळ हे दोन
न रंगवलेल्या चौरसांच्या क्षेत्रफळांच्या बेरजेइतके आहे. नारिंगी चौरसाची
बाजू $p$ आणि प्रत्येक न रंगवलेल्या चौरसाची बाजू $q$ असेल, तर $p^2 =
2q^2$. पण आता $p < m$, $q < n$ आणि $p, q \in \Nat$ असून
$\sqrt{2} = \nicefrac{p}{q}$. हे $m$ आणि $n$ शक्य तितके लहान निवडले
होते या तथ्याशी विसंगत आहे.

\emph{दुसरी, आकारिक रीतीने.} $m^{2} = 2n^{2}$ असल्यामुळे $m$ सम आहे.
($x$ विषम असेल, तर $x^2$ विषम असतो हे सहज दाखवता येते.) म्हणून काही
$r \in \Nat$ साठी $m = 2r$. पुनर्मांडणी केल्यास $2r^2 = n^2$,
त्यामुळे $n$ देखील सम आहे. म्हणून $m$ आणि $n$ दोन्ही सम आहेत, आणि
त्यामुळे $\nicefrac{m}{n}$ हा अपूर्णांक पुढे \emph{संक्षिप्त करता येतो}.
व्याघात!
\end{proof}

जाताजाता, या आकृतीवरील सिद्धतेमुळे पर्यायी विभाग
मधील सामग्रीचा पुन्हा विचार करता येतो. टेनेनबाउम (1927--2005) हे पूर्णपणे
आधुनिक गणितज्ञ होते; तरी ही सिद्धता निर्विवादपणे सुंदर, पूर्णतः तर्कशुद्ध
आहे आणि भूमितीय अंतःप्रज्ञेला आवाहन करते!
\begin{quote}\small\textbf{स्रोतदुरुस्ती.} MRARITH-003: गोठवलेल्या इंग्रजी मजकुरात निधनवर्ष 2006 दिले आहे. टेनेनबाउम यांच्या स्मृतिस्थळावरील चरित्र नोंद निधनाची तारीख 4 मे 2005 देते; मराठी मजकुरात 1927--2005 वापरले आहे.\end{quote}

कसेही असो, वास्तव संख्या परिमेय संख्यांपेक्षा ``अधिक विस्तृत'' आहेत.
एका अर्थाने परिमेय संख्यांमध्ये ``पोकळ्या'' आहेत आणि वास्तव संख्या त्या
भरून काढतात. वायश्ट्रास यांच्या लक्षात आले की या वर्णनातून वास्तव
संख्यांचा एकच गुणधर्म व्यक्त होतो, जो त्यांना परिमेय संख्यांपासून वेगळे
करतो; तो म्हणजे संपूर्णता गुणधर्म: \emph{ऊर्ध्वबंध असलेल्या वास्तव
संख्यांच्या प्रत्येक अरिक्त संचाला लघुतम ऊर्ध्वबंध असतो.}

परिमेय संख्यांना संपूर्णता गुणधर्म नाही हे सहज दिसते. उदाहरणार्थ,
$\sqrt{2}$ पेक्षा लहान परिमेय संख्यांचा पुढील संच पाहा:
\[
	\Setabs{p \in \Rat}{p^2 < 2 \text{ किंवा }p < 0}
\]
या संचाला परिमेय संख्यांमध्ये ऊर्ध्वबंध आहे; उदाहरणार्थ, त्याचे सर्व
घटक< $3$ पेक्षा लहान आहेत. पण त्याचा लघुतम ऊर्ध्वबंध कोणता?
आपल्याला ` $\sqrt{2}$ ' असे म्हणायचे आहे; पण $\sqrt{2}$ ही \emph{परिमेय
संख्या नाही} हे आपण नुकतेच पाहिले. आणि $\sqrt{2}$ पेक्षा मोठी अशी कोणतीही
\emph{लघुतम} परिमेय संख्या नाही. म्हणून या संचाला ऊर्ध्वबंध आहे, पण लघुतम
ऊर्ध्वबंध नाही. त्यामुळे परिमेय संख्यांमध्ये संपूर्णता गुणधर्माचा अभाव आहे.

याउलट, सांतत्याला संपूर्णता गुणधर्म स्वभावतः असायला हवा. $\sqrt{2}$ ही
केवळ वास्तव संख्या असावी एवढेच आपल्याला नको; परिमेय संख्या-रेषेतील सर्व
``पोकळ्या'' भरायच्या आहेत. खरे तर सांतत्यात स्वतःमध्येच कोणतीही ``पोकळी''
नसावी असे आपल्याला हवे आहे. संपूर्णतेद्वारे नेमके हेच मिळेल.







\section{$\Rat$ पासून $\Real$ कडे}\label{sfr:arith:cuts:sec}

मूलतः, संपूर्णता गुणधर्म असे दाखवतो की वास्तव संख्या-रेषेवरील कोणताही
$\alpha$ बिंदू त्या रेषेचे अचूक दोन भाग करतो: खालच्या भागाचा $\alpha$
हा लघुतम ऊर्ध्वबंध असतो आणि वरच्या भागाचा $\alpha$ हा महत्तम निम्नबंध
असतो. परिमेय संख्यांपासून वास्तव संख्यांची \emph{रचना} करण्यासाठी
डेडेकिंड यांनी वास्तव संख्यांना परिमेय संख्यांचे विभाजन करणारे
\emph{छेद} मानण्याची सूचना केली. म्हणजेच, $< \sqrt{2}$ असलेल्या परिमेय
संख्यांना $> \sqrt{2}$ असलेल्या परिमेय संख्यांपासून वेगळे करणाऱ्या
\emph{छेदाशी} आपण $\sqrt{2}$ ची एकरूपता मानतो.

ही मांडणी नीट नेमकी करू या. परिमेय संख्यांचे दोन भाग केले, तर आपण केलेले
विभाजन केवळ त्याचा \emph{खालचा} भाग विचारात घेऊन अनन्यपणे ओळखता येते.
म्हणून नेमकी पुढील व्याख्या देऊ:

\begin{defn}[छेद]
\emph{छेद} $\alpha$ म्हणजे परिमेय संख्यांचा असा कोणताही अरिक्त उचित
आरंभीचा खंड ज्याला महत्तम घटक नाही. म्हणजे $\alpha$ हा छेद तेव्हा आणि
केवळ तेव्हाच असतो, जेव्हा:
\begin{enumerate}
	\item \emph{अरिक्त, उचित}: $\emptyset \neq \alpha \subsetneq \Rat$
	\item \emph{आरंभीचा}: सर्व $p,q \in \Rat$ साठी: $p < q \in \alpha$ असेल, तर $p \in \alpha$
	\item \emph{महत्तम घटक नाही}: प्रत्येक $p \in \alpha$ साठी $p < q$ असे काही $q \in \alpha$ असते
\end{enumerate}
मग $\Real$ हा छेदांचा संच आहे.
\end{defn}

आता $\sqrt{2} = \Setabs{p \in \Rat}{p^2 < 2\text{
किंवा }p < 0}$ असे म्हणता येते. अर्थात, हा खरोखर \emph{छेद} आहे हे
तपासले पाहिजे; ती तपासणी आपण \ref{sfr:arith:check:sec} कडे सोपवतो.

आधीप्रमाणेच, काही वस्तूंची व्याख्या केल्यानंतर त्यांच्यावरील मूलभूत
फलने आणि संबंध व्याख्यायित करावे लागतात. एका सोप्या व्याख्येपासून सुरुवात करू:
\begin{align*}
	\alpha \leq \beta \text{ तेव्हा आणि केवळ तेव्हाच }\alpha \subseteq \beta
\end{align*}
क्रमाची ही व्याख्या छेदांच्या संचाला संपूर्णता गुणधर्म आहे हा मध्यवर्ती
निष्कर्ष \emph{मांडू} देते. पूर्णपणे उलगडल्यास विधानाचे रूप असे आहे.
$S$ हा ऊर्ध्वबंध असलेला छेदांचा अरिक्त संच असेल, तर $S$ ला लघुतम
ऊर्ध्वबंध असतो. अधिक तपशीलाने: $S$ चा ऊर्ध्वबंध असणारा $\lambda$ हा छेद
असतो, म्हणजे $(\forall \alpha \in S)\alpha \subseteq \lambda$, आणि
$\lambda$ हा अशा छेदांपैकी लघुतम असतो, म्हणजे $(\forall \beta \in \Real)((\forall \alpha \in S)\alpha \subseteq \beta \lif \lambda \subseteq \beta)$. आता या निष्कर्षाची सिद्धता पाहा:

\begin{thm}\label{sfr:arith:cuts:realcompleteness}
छेदांच्या संचाला संपूर्णता गुणधर्म आहे.
\end{thm}

\begin{proof}
$S$ हा ऊर्ध्वबंध असलेला छेदांचा कोणताही अरिक्त संच असू द्या. $\lambda
= \bigcup S$ असे ठरवू.
प्रथम $\lambda$ हा छेद आहे असा दावा करू:
\begin{enumerate}
\item $S$ अरिक्त असल्यामुळे $S$ मध्ये किमान एक छेद आहे, म्हणून
$\emptyset \neq \lambda$. $S$ हा छेदांचा संच असल्यामुळे $\lambda
\subseteq \Rat$. $S$ ला ऊर्ध्वबंध असल्यामुळे प्रत्येक $\alpha \in S$
मधून अनुपस्थित असलेली अशी काही $p \in \Rat$ आहे. म्हणून $p\notin \lambda$,
आणि परिणामी $\lambda \subsetneq \Rat$.
\begin{quote}\small\textbf{स्रोतदुरुस्ती.} MRARITH-004: गोठवलेल्या इंग्रजी सिद्धतेत $\lambda$ अरिक्त असल्याचे पहिले कारण चुकून “S has an upper bound” असे दिले आहे. आवश्यक आणि आधीच गृहीत असलेले कारण S अरिक्त असणे हे आहे; मराठी सिद्धतेत तेच वापरले आहे. पुढील उचित-उपसंच युक्तिवादासाठी मात्र ऊर्ध्वबंधाचे गृहीतक योग्य आहे.\end{quote}
\item $p < q \in \lambda$ असे समजा. मग $q \in \alpha$ असे काही
$\alpha \in S$ आहे. $\alpha$ हा छेद असल्यामुळे $p \in \alpha$.
म्हणून $p \in \lambda$.
\item $p \in \lambda$ असे समजा. मग $p \in \alpha$ असे काही
$\alpha \in S$ आहे. $\alpha$ हा छेद असल्यामुळे $p < q$ असे काही
$q \in \alpha$ आहे. म्हणून $q \in \lambda$.
\end{enumerate}
यावरून दावा सिद्ध होतो. शिवाय, स्पष्टपणे $(\forall \alpha \in S)\alpha
\subseteq \bigcup S = \lambda$, म्हणजे $\lambda$ हा $S$ चा ऊर्ध्वबंध
आहे. आता $\beta \in \mathbb{R}$ हाही ऊर्ध्वबंध आहे असे समजा, म्हणजे
$(\forall \alpha \in S)\alpha \subseteq \beta$. कोणत्याही $p \in \Rat$
साठी, $p \in \lambda$ असेल तर $p \in \alpha$ असे काही $\alpha \in S$
आहे, आणि त्यामुळे $p \in \beta$. सार्वत्रिकीकरण केल्यास $\lambda
\subseteq \beta$. म्हणून $\lambda$ हा $S$ चा \emph{लघुतम} ऊर्ध्वबंध आहे.
\end{proof}

म्हणून आपल्याकडे संपूर्णता गुणधर्म पूर्ण करणाऱ्या अनेक वस्तू आहेत.
हे सांगण्याचा एक मार्ग असा: आपल्या छेदांमध्ये ``पोकळ्या'' नाहीत.
(त्यामुळे परिमेय संख्यांऐवजी वास्तव संख्यांचे आणखी ``छेद'' घेतल्यास
काही रोचक नवी वस्तू मिळणार नाही.)

पुढे वास्तव संख्यांवरील काही क्रिया व्याख्यायित केल्या पाहिजेत. प्रत्येक
$p \in \Rat$ साठी $p_\Real = \Setabs{q \in \Rat}{q < p}$ असे ठरवून
परिमेय संख्यांना वास्तव संख्यांत अंतःस्थापित करण्यापासून सुरुवात करू.
मग पुढील व्याख्या करतो:
\begin{align*}
	\alpha + \beta &= \Setabs{p + q}{p \in \alpha \land q \in \beta}\\
	\alpha \times \beta &=
	\Setabs{p \times q}{0 \leq p \in \alpha \land 0 \leq q \in \beta} \cup 0_\Real & \text{जर }\alpha, \beta \geq 0_\Real
\end{align*}
\begin{quote}\small\textbf{स्रोतदुरुस्ती.} MRARITH-005: गोठवलेल्या इंग्रजी गुणाकार-व्याख्येत शून्याचा छेद एकदाच $0^R$ असा विसंगत ऊर्ध्वनिर्देशांक वापरून लिहिला आहे; त्याच परिच्छेदातील सर्व नियंत्रक संकेतन 0\_R वापरते. मराठी सूत्रात 0\_R वापरले आहे.\end{quote}
गुणाकाराची इतर प्रकरणे हाताळण्यासाठी प्रथम $0_\Real \times \alpha = 0_\Real = \alpha \times 0_\Real$ असे ठरवू आणि पुढील व्याख्या करू:
\begin{quote}\small\textbf{स्रोतदुरुस्ती.} MRARITH-006: गोठवलेल्या इंग्रजी स्रोताने शून्यासोबतच्या गुणाकाराची ही आवश्यक व्याख्या त्याच ओळीवर निष्क्रिय केली आहे; त्यामुळे एक घटक ऋण व दुसरा शून्य असलेली प्रकरणे पुढील प्रकरणनिहाय व्याख्येत येत नाहीत. स्रोताच्याच निष्क्रिय सूत्राला मराठीत सक्रिय केले आहे.\end{quote}
\begin{align*}
	-\alpha &= \Setabs{p - q}{p < 0 \land q \notin \alpha}
\end{align*}
आणि मग पुढीलप्रमाणे ठरवू:
\begin{align*}
	\alpha \times \beta &\defis
	\begin{cases}
		\mathord{-}\alpha \times \mathord{-}\beta &\text{जर }\alpha < 0_\Real\text{ आणि }\beta < 0_\Real\\
		\mathord{-}(\mathord{-}\alpha \times \beta) &\text{जर }\alpha < 0_\Real \text{ आणि }\beta > 0_\Real\\
		\mathord{-}(\alpha \times \mathord{-}\beta) &\text{जर }\alpha > 0_\Real \text{ आणि }\beta < 0_\Real
	\end{cases}
\end{align*}
यानंतर यांपैकी प्रत्येक व्याख्येतून नेहमी छेदच मिळतो हे तपासावे लागेल.
आणि शेवटी, अशा प्रकारे व्याख्यायित केलेले छेद अपेक्षेप्रमाणेच वागतात
हे दाखवणारी सोपी (पण लांबलचक) सिद्धता पूर्ण केली पाहिजे. मात्र हे सर्व
आपण \ref{sfr:arith:check:sec} कडे सोपवतो.







\section{काही तात्त्विक चिंतन}\label{sfr:arith:ref:sec}

तांत्रिक तपशील एवढेच. पण त्यातून साधले तरी काय?

एक तर, अगदी निर्विवादपणे, त्यातून आपल्याला काही {सुंदर} शुद्ध गणित
मिळाले. शिवाय काही खोल संकल्पनात्मक कामगिरीही झाली. वास्तव संख्या आणि
परिमेय संख्या यांतील निर्णायक फरक संपूर्णता गुणधर्म व्यक्त करतो, हे
ओळखणे ही एक सखोल अंतर्दृष्टी होती. तसेच, डेडेकिंड छेद म्हणून वास्तव
संख्यांची स्पष्ट रचना केल्यामुळे विश्लेषणाच्या विषयाला भक्कम आधार मिळतो.
छेदांपासून नेमके असेच क्षेत्र बनत असल्यामुळे \emph{संपूर्ण क्रमित क्षेत्र}
ही संकल्पना सुसंगत आहे, हे आपल्याला माहीत आहे.

एवढे असूनही, या कामगिरीबद्दलच्या काही शंका उघडपणे मांडल्या पाहिजेत.

पहिली शंका अशी की, वास्तव संख्यांचा परिचित (कदाचित अनंत) दशांश विस्तार
विचारात घेण्यापेक्षा छेदांच्या रूपात त्यांचा विचार करणे \emph{अधिक}
काटेकोर आहे, हे स्पष्ट नाही. दशांश विस्तारांनी वास्तव संख्यांची केलेली
ही ``रचना'' कोशी क्रमिकांद्वारे केलेल्या रचनेस काहीशी मिळतीजुळती आहे;
परंतु प्रत्यक्षात ती सतराव्या शतकाच्या प्रारंभापासूनच गणितज्ञांना मूलतः
माहीत होती (पाहा \ref{sfr:arith:cauchy:sec}). वास्तव संख्यांना संपूर्णता गुणधर्म
आहे याची जाणीव झाल्यामुळे काटेकोरपणात खरी वाढ झाली; वास्तव संख्यांची काही
विशिष्ट संच म्हणून रचना करता येणे हे कदाचित स्वतःपुरते फारसे रोचक नाही.

पूर्णांकांचे किंवा परिमेय संख्यांचे (त्याहून खूप सोपे) अंकगणितीकरण त्या
क्षेत्रांतील काटेकोरपणा वाढवते, हे तर आणखी कमी स्पष्ट आहे. येथे एक साधे
निरीक्षण उपयुक्त ठरेल. नैसर्गिक संख्यांच्या क्रमित जोड्यांचे तुल्यतावर्ग
म्हणून पूर्णांकांची \emph{रचना} केल्यानंतर, मग पूर्णांकांच्या क्रमित
जोड्यांचे तुल्यतावर्ग म्हणून परिमेय संख्यांची रचना केल्यानंतर, आणि मग
परिमेय संख्यांचे संच म्हणून वास्तव संख्यांची रचना केल्यानंतर आपण त्या
रचना लगेचच \emph{विसरून जातो}. विशेषतः, कोणालाही गणिती सिद्धतेत या
रचनांचा \emph{उपयोग} करावासा वाटणार नाही (त्या रचना अपेक्षेप्रमाणेच
वागतात हे सिद्ध करताना मात्र अर्थातच त्यांचा उपयोग करावा लागेल).
नैसर्गिक संख्यांच्या संचांच्या संचांच्या संचांच्या संचांच्या संचांच्या
संचांच्या एखाद्या संचाबद्दल बोलण्यापेक्षा वास्तव संख्येबद्दल थेट बोलणे
खूपच सोपे आहे.

या व्याख्या आपल्याला पूर्णांक, परिमेय संख्या किंवा वास्तव संख्या
\emph{नेमक्या काय आहेत}, \emph{सत्तामीमांसक दृष्ट्या सांगायचे तर}, हे
सांगतात, याबद्दल तर
सर्वाधिक शंका आहे. म्हणजे, वास्तव संख्या (उदाहरणार्थ) काही विशिष्ट संच
(संचांच्या संचांचे संच\ldots) \emph{आहेत}, हे संशयास्पद आहे. अशा
दृष्टिकोनासमोरील मुख्य अडथळा हा की ही रचना अनेक वेगवेगळ्या पद्धतींनी करता
आली असती. वास्तव संख्यांच्या बाबतीत खरोखर रोचक रीतीने भिन्न असलेल्या काही
रचना आहेत (पाहा \ref{sfr:arith:cauchy:sec}). पण काही वेगळ्या रचना मिळवण्याची
एक अगदी क्षुल्लक पद्धत अशी: \ref{sfr:rel:ref:sec} प्रमाणे आपण
क्रमित जोडीची व्याख्या किंचित वेगळी करू शकलो असतो; क्रमित जोडीची ही
पर्यायी संकल्पना वापरली असती, तर आपल्या रचना तितक्याच यशस्वी झाल्या
असत्या, पण शेवटी मिळालेल्या वस्तू वेगळ्या असत्या. त्यामुळे पूर्णांक,
परिमेय संख्या आणि वास्तव संख्या यांच्या अनेक प्रतिस्पर्धी संचसैद्धान्तिक
रचना आहेत. आता पूर्णांक (उदाहरणार्थ) \emph{त्या} संचांऐवजी \emph{हेच}
संच आहेत, असा दावा करणे केवळ मनमानीचे (आणि अडचणीचे) ठरेल.
(\ref{sfr:rel:ref:sec} प्रमाणे, Benacerraf (1965) यांनी
प्रसिद्ध केलेल्या युक्तिवादाचे हे एक उदाहरण आहे.)

आणखी एक मुद्दा उपस्थित करायला हवा: आपल्या रचनांमध्ये काहीतरी फारच
\emph{विचित्र} आहे. आपण नैसर्गिक संख्यांपासून सुरुवात केली. मग पूर्णांकांची
रचना केली आणि ``पूर्णांकांचा $0$'', म्हणजे $ \equivrep{0,0}{\Intequiv}$,
याची रचना केली. पण $0 \neq \equivrep{0,0}{\Intequiv}$. खरे तर, आपल्या
रचनांनुसार \emph{एकही} नैसर्गिक संख्या पूर्णांक नाही. पण हे सहजज्ञानाच्या
अगदी विरुद्ध वाटते. प्रत्यक्षात \ref{sfr:set:imp:sec} मध्ये आपण
फारसा युक्तिवाद न करता $\Nat \subseteq \Rat$ असा दावा केला होता. संख्या
नेमक्या \emph{काय आहेत} हे या रचना सांगत असतील, तर तो दावा उघडपणे असत्य
होता.

मग थोडे मागे हटून पाहिल्यास आपण कुठवर पोहोचलो? अनौपचारिक संचसिद्धान्तात
काम करताना आणि नैसर्गिक संख्या गृहीत धरताना, पूर्णांक, परिमेय संख्या आणि
वास्तव संख्या यांना काही विशिष्ट संच \emph{म्हणून वागवणे} आपल्याला शक्य
होते. त्या अर्थाने, या वस्तूंचे सिद्धान्त संचसिद्धान्तात
\emph{अंतःस्थापित} करता येतात. पण या अंतःस्थापनाचा तात्त्विक आशय अजिबात
सरळसोट नाही.

अर्थात, याने या विषयावरील शेवटचा शब्द सांगितलेला नाही!{} मुद्दा फक्त
एवढाच आहे. वास्तव संख्यांच्या अंकगणितीकरणाला खोल तात्त्विक महत्त्व
\emph{आहे} हे दाखवण्यासाठी आणखी काही \emph{तात्त्विक} युक्तिवाद आवश्यक असेल.









\section{क्रमित वलये आणि क्षेत्रे}\label{sfr:arith:check:sec}

या प्रकरणभर काही व्याख्या ``जशा वागायला हव्यात'' तशा वागतात, असे
आपण म्हटले. या तांत्रिक परिशिष्टात त्याचा अर्थ स्पष्ट करू आणि त्या
व्याख्या ``योग्य रीतीने'' वागतात हे कसे दाखवायचे याचा आराखडा देऊ.

\ref{sfr:arith:int:sec} मध्ये आपण $\Int$ वरील बेरीज आणि गुणाकार व्याख्यायित
केले. व्याख्येप्रमाणे या क्रिया $\Int$ ला आपण ``इच्छित'' असलेली संरचना
देतात, हे दाखवायचे आहे. विशेषतः, येथे अभिप्रेत संरचना क्रमनिरपेक्ष
वलयाची आहे.

\begin{defn}
	\emph{क्रमनिरपेक्ष वलय} म्हणजे विशिष्ट $0$ व $1$ घटक आणि $+$ व $\times$
	क्रिया असलेला असा संच $S$, जो पुढील आठ सूत्रे पूर्ण करतो:
	\begin{align*}
		\emph{सहचारिता}&&a + (b+ c) & = (a + b) + c \\
		&& (a \times b) \times c & = a \times (b\times c)\\
		\emph{क्रमनिरपेक्षता}&&a + b &= b+ a  \\
		&&  a \times b&= b\times a\\
		\emph{अविकारक}&&a + 0 &= a \\
		&& a \times 1 &= a\\
		\emph{बेरीज व्यस्त}&&(\exists b\in S)0&=a + b\\
		\emph{वितरणता}&&a \times (b+ c ) &= (a \times b) + (a \times c)
	\end{align*}
	येथे ही सर्व सूत्रे $S$ पुरते मर्यादित सार्वत्रिक परिमाणकांनी बद्ध आहेत,
	हे अध्याहृत आहे. तसेच येथील $0$ आणि~$1$ हे घटक याच नावाच्या नैसर्गिक
	संख्या असतीलच असे नाही.
\end{defn}

म्हणून पूर्णांक क्रमनिरपेक्ष वलय बनवतात हे तपासण्यासाठी या आठ अटी पूर्ण
होतात एवढेच तपासायचे आहे. यांपैकी कोणतीही अट सिद्ध करणे {अवघड} नाही,
पण ते थोडे कष्टाचे आहे. उदाहरणार्थ, बेरजेच्या बाबतीत
\emph{सहचारिता} अशी सिद्ध करता येते:

\begin{proof}
$i, j, k \in \Int$ निश्चित करा. मग $i = \equivrep{a_1, b_1}{}$ आणि $j =
\equivrep{a_2,b_2}{}$ आणि $k = \equivrep{a_3, b_3}{}$ असे होतील अशा
$a_1, b_1, a_2, b_2, a_3, b_3 \in \Nat$ संख्या आहेत. (वाचनीयतेसाठी
``$\equivrep{x, y}{\Intequiv}$'' ऐवजी ``$\equivrep{x, y}{}$'' असे लिहितो;
या संपूर्ण विभागात आपण हेच करू.) आता:
\begin{align*}
	i + (j + k) &= \equivrep{a_1, b_1}{}+(\equivrep{a_2, b_2}{} + \equivrep{a_3, b_3}{}) \\
	&= \equivrep{a_1,  b_1}{} + \equivrep{a_2+a_3, b_2+b_3}{}\\
	&= \equivrep{a_1 + (a_2 + a_3), b_1 + (b_2 + b_3)}{}\\
	&= \equivrep{(a_1 + a_2) + a_3, (b_1 + b_2) + b_3}{}\\
	&= \equivrep{a_1 + a_2, b_1 + b_2}{} + \equivrep{a_3, b_3}{}\\
	&= (\equivrep{a_1, b_1}{} + \equivrep{a_2, b_2}{}) + \equivrep{a_3, b_3}{}\\
	&= (i+j) + k
\end{align*}
येथे $\Nat$ वरील बेरजेच्या वर्तनाचा आपण मुक्तपणे उपयोग केला आहे.
\end{proof}

त्याचप्रमाणे, \emph{बेरीज व्यस्त} असा सिद्ध करता येतो:

\begin{proof}
$i \in \Int$ निश्चित करा; म्हणजे काही $a,b \in \Nat$ साठी $i =
\equivrep{a,b}{}$. $j = \equivrep{b,a}{} \in \Int$ असे ठरवू. नैसर्गिक
संख्यांच्या वर्तनाचा उपयोग केल्यास $(a+b) + 0 = 0 + (a+b)$; म्हणून
व्याख्येनुसार $\tuple{a+b, b+a} \sim_\Int \tuple{0,0}$, आणि त्यामुळे
$\equivrep{a+b, b+a}{} = \equivrep{0, 0}{} = 0_\Int$. आता $i + j =
\equivrep{a,b}{}+\equivrep{b,a}{}=\equivrep{a+b, b+a}{}= \equivrep{0,
0}{} = 0_\Int$.
\end{proof}

आणि \emph{वितरणतेची} सिद्धता अशी आहे:

\begin{proof}
वरीलप्रमाणे, $i = \equivrep{a_1, b_1}{}$ आणि $j =
\equivrep{a_2,b_2}{}$ आणि $k = \equivrep{a_3, b_3}{}$ निश्चित करा. आता:
\begin{align*}
	i \times (j + k)
	&= \equivrep{a_1, b_1}{} \times (\equivrep{a_2,b_2}{} + \equivrep{a_3, b_3}{})\\
	&= \equivrep{a_1, b_1}{} \times \equivrep{a_2 + a_3,b_2+b_3}{}\\
	&= \equivrep{a_1  (a_2 + a_3) + b_1  (b_2+b_3), a_1  (b_2 + b_3) + b_1 (a_2 + a_3)}{}\\
	&= \equivrep{a_1 a_2 + a_1a_3 + b_1 b_2+b_1b_3, a_1 b_2 + a_1b_3 + a_2b_1 + a_3b_1}{}\\
	&= \equivrep{a_1a_2 + b_1b_2, a_1b_2 + a_2b_1}{} + \equivrep{a_1a_3 + b_1b_3, a_1b_3 + a_3b_1}{}\\
	&= (\equivrep{a_1, b_1}{} \times \equivrep{a_2,b_2}{}) + (\equivrep{a_1, b_1}{} \times  \equivrep{a_3, b_3}{})\\
	&= (i \times j) + (i \times k)
\end{align*}
\end{proof}

उरलेल्या पाच अटी सिद्ध करण्याचे काम सरावासाठी सोडतो. ते पूर्ण केल्यावर
$\Int$ हे क्रमनिरपेक्ष वलय आहे, म्हणजे व्याख्यायित केलेली बेरीज आणि
गुणाकार अपेक्षेप्रमाणे वागतात, हे आपण दाखवलेले असेल.

\begin{prob}
$\Int$ हे क्रमनिरपेक्ष वलय आहे हे सिद्ध करा.
\end{prob}

पण आपले काम अजून संपलेले नाही. $\Int$ वरील बेरीज आणि गुणाकार यांच्याबरोबरच
आपण $\leq$ हा क्रमसंबंध व्याख्यायित केला होता; तो अपेक्षेप्रमाणे वागतो हेही
तपासले पाहिजे. अधिक नेमकेपणाने, $\Int$ हे \emph{क्रमित} वलय आहे हे दाखवले
पाहिजे.\footnote{\ref{sfr:rel:ord:def:linearorder} मधील व्याख्या आठवा:
	पूर्ण क्रम म्हणजे परावर्ती, संक्रमक, प्रतिसममित आणि संयोजित संबंध.
	क्रमसंबंधांच्या संदर्भात संयोजिततेला कधीकधी \emph{त्रिपर्याय} म्हणतात,
	कारण कोणत्याही $a$ आणि $b$ साठी $a \leq b \lor a
	= b \lor a \geq b$.}

\begin{defn}
\emph{क्रमित वलय} म्हणजे $\leq$ या पूर्ण क्रमसंबंधानेही सज्ज असलेले
असे क्रमनिरपेक्ष वलय की:
\begin{align*}
	a \leq b &\lif a + c \leq b + c\\
	(a \leq b \land 0 \leq c) &\lif a \times c \leq b \times c
\end{align*}
\end{defn}

\begin{prob}
$\Int$ हे क्रमित वलय आहे हे सिद्ध करा.
\end{prob}

पूर्वीप्रमाणेच, रचलेला $\Int$ हा क्रमित वलय आहे हे दाखवणे कष्टाचे पण
नित्याचे काम आहे. ते आम्ही तुमच्यासाठी सोडतो.

यामुळे पूर्णांकांचे काम संपले. आता परिमेय संख्यांबद्दल अगदी तशाच गोष्टी
दाखवायच्या आहेत. विशेषतः, $+$, $\times$ आणि $\leq$ यांच्या आपल्या
व्याख्यांनुसार परिमेय संख्या क्रमित \emph{क्षेत्र} बनवतात हे दाखवायचे आहे:
\begin{defn}\label{sfr:arith:check:orderedfield}
\emph{क्रमित क्षेत्र} म्हणजे पुढील अटही पूर्ण करणारे क्रमित वलय:
\begin{align*}
	\emph{गुणाकार व्यस्त}& & (\forall a \in S \setminus \{0\})(\exists b \in S) a\times b& = 1
\end{align*}
\end{defn}

$\Int$ हे क्रमित वलय आहे हे दाखवल्यानंतर $\Rat$ हे क्रमित क्षेत्र आहे
हे दाखवणे सोपे पण कष्टाचे आहे.

\begin{prob}
$\Rat$ हे क्रमित क्षेत्र आहे हे सिद्ध करा.
\end{prob}

पूर्णांक आणि परिमेय संख्या हाताळल्यानंतर केवळ वास्तव संख्या उरतात.
विशेषतः, $\Real$ हे \emph{संपूर्ण} क्रमित क्षेत्र, म्हणजे संपूर्णता गुणधर्म
असलेले क्रमित क्षेत्र आहे, हे दाखवले पाहिजे. \ref{sfr:arith:cuts:realcompleteness}
मध्ये $\Real$ ला संपूर्णता गुणधर्म आहे हे सिद्ध केले. मात्र $\Real$ हे
क्रमित क्षेत्र आहे याची (कंटाळवाणी) तपासणी अजून करायची आहे.
\begin{quote}\small\textbf{स्रोतदुरुस्ती.} MRARITH-007: गोठवलेल्या इंग्रजी वाक्यात “run through the (tedious) of checking” येथे आवश्यक नाम गाळले आहे. मराठीत अभिप्रेत “कंटाळवाणी तपासणी” स्पष्टपणे दिली आहे.\end{quote}

\emph{त्या} कष्टाच्या सरावाकडे वळण्यापूर्वी आणखी काही ``तात्काळ'' गोष्टी
तपासल्या पाहिजेत. उदाहरणार्थ, कोणत्याही $\alpha$ आणि $\beta$ छेदांसाठी
व्याख्यायित केलेला $\alpha + \beta$ खरोखरच \emph{छेद} आहे याची खात्री
हवी. त्याची सिद्धता अशी:

\begin{proof}
$\alpha$ आणि $\beta$ हे दोन्ही छेद असल्यामुळे $\alpha + \beta = \Setabs{p
+ q}{p \in \alpha \land q \in \beta}$ हा $\Rat$ चा अरिक्त उचित उपसंच
आहे. आता काही $p \in \alpha$ आणि $q \in \beta$ साठी $x < p + q$ असे
समजा. मग $x - p < q$, म्हणून $x - p \in \beta$, आणि $x = p + (x - p)
\in \alpha + \beta$. म्हणून $\alpha + \beta$ हा $\Rat$ चा आरंभीचा खंड
आहे. शेवटी, कोणत्याही $p + q \in \alpha + \beta$ साठी, $\alpha$ आणि
$\beta$ हे दोन्ही छेद असल्यामुळे $p < p_1$ आणि $q < q_1$ असे काही
$p_1 \in \alpha$ आणि $q_1 \in \beta$ असतात; म्हणून $p + q < p_1 + q_1
\in \alpha + \beta$; त्यामुळे $\alpha + \beta$ ला महत्तम घटक नाही.
\end{proof}

अशाच प्रयत्नांनी $\alpha - \beta$, $\alpha \times \beta$ आणि $\alpha
\div \beta$ हे छेद आहेत हे तुम्हाला तपासता येईल (शेवटच्या प्रकरणात
$\beta$ हा शून्य-छेद असलेले प्रकरण वगळून). पुन्हा, हे काम आम्ही
तुमच्यासाठीच सोडतो.
\begin{quote}\small\textbf{स्रोतदुरुस्ती.} MRARITH-008: गोठवलेल्या इंग्रजी स्रोताच्या आधीच्या cuts विभागात वजाबाकी आणि भागाकाराची सूत्रे निष्क्रिय टिप्पण्यांतच आहेत; त्यामुळे हे वाक्य सक्रियपणे पूर्ण व्याख्या न दिलेल्या क्रियांना गृहीत धरते. मराठी मजकूर दावा जतन करतो आणि ही स्रोत-अपूर्णता स्पष्ट करतो.\end{quote}

\begin{prob}
$\Real$ हे क्रमित क्षेत्र आहे हे सिद्ध करा.
\end{prob}

पण एक लहानसा अपूर्ण मुद्दा आवरायचा आहे. \ref{sfr:arith:cuts:sec} मध्ये आपण
$\sqrt{2} = \Setabs{p \in \Rat}{p < 0 \text{ किंवा }p^2 < 2}$ असे घेता
येईल असे सुचवले. मात्र हा संच \emph{छेद} आहे हे दाखवणे आवश्यक आहे.
त्याची सिद्धता अशी:

\begin{proof}
हा परिमेय संख्यांचा अरिक्त उचित आरंभीचा खंड आहे हे स्पष्ट आहे; म्हणून
त्याला महत्तम घटक नाही एवढे दाखवणे पुरेसे आहे. विशेषतः, $p$ ही $p^2 < 2$
असलेली धन परिमेय संख्या आणि $q = \frac{2p+2}{p+2}$ असेल, तर $p < q$
आणि $q^2 < 2$ हे दोन्ही दाखवणे पुरेसे आहे. $p < q$ हे पाहण्यासाठी फक्त
पुढील नोंद घ्या:
\begin{align*}
	p^2 &< 2\\
	p^2 + 2p &< 2 + 2p\\
	p(p + 2) &< 2 + 2p\\
	p &< \tfrac{2+2p}{p+2} = q
\end{align*}
$q^2 < 2$ हे पाहण्यासाठी फक्त पुढील नोंद घ्या:
\begin{align*}
	p^2 &< 2\\
	2p^2 + 4p + 2 &< p^2 + 4p+ 4\\
	4p^2 + 8p + 4 &< 2(p^2 + 4p + 4)\\
	(2p+2)^2 & <2(p+2)^2\\
	\tfrac{(2p+2)^2}{(p+2)^2} &< 2\\
	q^2 &< 2
\end{align*}
\end{proof}








\section{परिशिष्ट: कोशी क्रमिकांच्या रूपातील वास्तव संख्या}\label{sfr:arith:cauchy:sec}

\ref{sfr:arith:cuts:sec} मध्ये आपण डेडेकिंड छेद म्हणून वास्तव संख्यांची रचना
केली. या विभागात आपण एक पर्यायी रचना स्पष्ट करू. आज आपण जिला कोशी क्रमिका
म्हणतो तिची कोशी यांनी दिलेली व्याख्या या रचनेचा आधार आहे; पण त्या
व्याख्येचा उपयोग करून वास्तव संख्यांची \emph{रचना} करण्याचे श्रेय
एकोणिसाव्या शतकातील इतर लेखकांना, विशेषतः वायश्ट्रास, हाइने, मेरे आणि
कँटर यांना जाते. (या इतिहासाच्या छानशा आढाव्यासाठी
O'Connor आणि Robertson (2005) यांचा लेख पाहा.)

एकोणिसाव्या शतकाकडे वळण्यापूर्वी सायमन स्टेव्हिन (1548--1620) यांचा
विचार करणे उपयुक्त ठरेल. थोडक्यात, प्रत्येक वास्तव संख्येचा विचार तिच्या
दशांश विस्ताराच्या रूपात करता येतो हे स्टेव्हिन यांच्या लक्षात आले.
त्यामुळे $\sqrt{2}$ सारख्या अपरिमेय संख्येलाही एक व्यवस्थित दशांश विस्तार
असतो; त्याची सुरुवात अशी:
\[
	1.41421356237\ldots
\]
संचसिद्धान्तात दशांश विस्तारांचे प्रतिमानीकरण करणे अगदी सोपे आहे: त्यांना
$d \colon \Nat \to \Nat$ अशी फलने माना, जिथे $d(n)$ हे आपल्याला हवे
असलेले $n$ वे दशांश स्थान आहे. मग दशांश चिन्हाच्या आधी येणारा वास्तव
संख्येचा भाग (येथे फक्त $1$) हाताळण्यासाठी या मांडणीत थोडा बदल करावा
लागेल. तसेच, उदाहरणार्थ $0.999\ldots = 1$ याची खात्री करण्यासाठी आणखी
एक बदल (तुल्यता संबंध) करावा लागेल. पण संचसिद्धान्ताच्या चौकटीत
स्टेव्हिन यांच्या पद्धतीने वास्तव संख्यांची पूर्णतः काटेकोर रचना देणे
कठीण नाही.

स्टेव्हिन हा आपला मुख्य विषय नाही. (स्टेव्हिन यांच्याविषयी अधिक माहितीसाठी
Katz आणि Katz (2012) पाहा.) पण त्याच्याशी जवळून संबंधित असा एक विचार
करू. $\sqrt{2}$ चा दशांश विस्तार थेट हाताळण्याऐवजी अधिकाधिक अचूक होत
जाणारे विस्तार घेऊन आपण $\sqrt{2}$ च्या अधिकाधिक अचूक परिमेय निकटनांची
\emph{क्रमिका} विचारात घेऊ शकतो:
\[
1, 1.4, 1.414, 1.4142, 1.41421,\ldots
\]
वास्तव संख्यांचा विचार “अधिकाधिक चांगल्या निकटनांद्वारे” करता येतो या
कल्पनेतून अंतर्दृष्टींची आणखी एक मालिका मिळते (जशी आपण $\Int$ पासून
$\Rat$ किंवा $\Nat$ पासून $\Int$ रचताना वापरली होती). त्या मूलभूत
अंतर्दृष्टी अशा:
\begin{enumerate}
	\item प्रत्येक वास्तव संख्या एका (कदाचित अनंत) दशांश विस्ताराने
	लिहिता येते.
	\item एका (कदाचित अनंत) दशांश विस्तारात सांकेतिक केलेली माहिती
	परिमेय संख्यांच्या क्रमिकेनेही सांकेतिक करता येते.
	\item परिमेय संख्यांची क्रमिका ही $\Nat$ पासून $\Rat$ कडे जाणारे
	फलन मानता येते; क्रमिकेतील $n$ वी परिमेय संख्या $f(n)$ माना.
\end{enumerate}
अर्थात, $\Nat$ पासून $\Rat$ कडे जाणारे \emph{कोणतेही} फलन आपल्याला वास्तव
संख्या देईल असे नाही. उदाहरणार्थ, हे फलन पाहा:
\[
	f(n) =\begin{cases}
			1 & \text{जर }n\text{ विषम असेल}\\
			0 &\text{जर }n\text{ सम असेल}
		\end{cases}
\]
येथील मुख्य अडचण अशी की $0,1,0,1,0,1,0,\ldots$ ही क्रमिका कोणत्याही
वास्तव संख्येकडे जवळ जाताना दिसत नाही. म्हणून एखाद्या वास्तव संख्येकडे
जवळ जाणाऱ्या क्रमिकांचाच विचार करण्यासाठी आपले लक्ष काही \emph{सीमा}
असलेल्या क्रमिकांपुरते मर्यादित केले पाहिजे.

सीमेची कल्पना पर्यायी विभाग मध्ये आपण यापूर्वी पाहिली
आहे. पण तेथे वापरलेली \emph{अगदी तीच} व्याख्या येथे वापरता येणार नाही.
तेथील “$(\forall \epsilon>0)$” या पदावलीत वास्तव संख्यांवरील
परिमाणन गृहीत धरले होते; आणि आपण वास्तव संख्यांवरील फलनांच्या सीमा
विचारात घेत होतो. म्हणून ती व्याख्या येथे वापरणे म्हणजे वास्तव संख्या
आधीच उपलब्ध मानणे होईल; पण नेमकी त्यांचीच आपल्याला \emph{रचना} करायची
आहे. सुदैवाने, सीमेच्या अगदी जवळच्या कल्पनेने आपण काम करू शकतो.
\begin{defn}\label{sfr:arith:cauchy:def:CauchySequence}
  $f: \Nat \to \Rat$ हे फलन \emph{कोशी क्रमिका} आहे तेव्हाच आणि
  तेव्हाच, जेव्हा प्रत्येक धन $\epsilon \in \Rat$ साठी $(\exists \ell
  \in \Nat)(\forall m, n > \ell)|f(m) - f(n)| < \epsilon$.
\end{defn}

सीमेची व्यापक कल्पना पूर्वीसारखीच आहे: यथार्थतेची एखादी पातळी
(जिचे मापन~$\epsilon$ ने केले आहे) हवी असेल, तर पाहण्यासाठी एक
“प्रदेश” असतो ($\ell$ पेक्षा मोठे कोणतेही आदान). आपल्या $1$, $1.4$,
$1.414$, $1.4142$, $1.41421$\ldots या क्रमिकेला सीमा आहे हेही सहज
दिसते: $\sqrt{2}$ चे निकटन $\nicefrac{1}{10^{n}}$ इतक्या त्रुटीच्या
आत हवे असेल, तर $n$ व्या पदानंतरचे कोणतेही पद पाहा.

मग वास्तव संख्या ही कोणतीही कोशी क्रमिका \emph{आहेच} असे म्हणण्याचा
विचार साहजिक वाटेल. पण $\Int$ आणि $\Rat$ यांच्या रचनांप्रमाणेच तो फार
भोळा ठरेल: कोणत्याही एका वास्तव संख्येचा निर्देश अनेक भिन्न कोशी क्रमिका
करतात. हे पाहण्याची एक सोपी पद्धत अशी. $f$ ही कोशी क्रमिका दिली असता,
$g(0)\neq f(0)$ एवढा अपवाद सोडून $g$ हे $f$ सारखेच फलन ठरवा. पहिल्या
संख्येनंतर दोन्ही क्रमिका सर्वत्र जुळत असल्यामुळे
\ref{sfr:arith:cauchy:def:CauchySequence} मध्ये वापरलेल्या अर्थाने त्यांची सीमा समान
आहे, आणि म्हणून त्या एकाच वास्तव संख्येची “व्याख्या” करतात, असे
आपल्याला शेवटी म्हणायचे असेल. त्यामुळे या कोशी क्रमिका एकच वास्तव
संख्या आहेत असे आपण खरे तर मानले पाहिजे.

म्हणून कोशी क्रमिकांवर पुन्हा एक तुल्यता संबंध परिभाषित करून वास्तव
संख्यांची ओळख त्या संबंधाखालील तुल्यतावर्गांशी करावी लागेल.
\begin{quote}\small\textbf{स्रोतदुरुस्ती.} MRARITH-009: गोठवलेल्या इंग्रजी मजकुरात वास्तव संख्यांची ओळख equivalence relations शी करावी असे म्हटले आहे; पुढील वाक्ये आणि रचना स्पष्टपणे equivalence classes वापरतात. मराठी मजकुरात अपेक्षित तुल्यतावर्ग वापरले आहेत.\end{quote}
सर्वप्रथम ज्याची सीमा $0$ आहे अशा फलनाची कल्पना हवी. कोणत्याही $h :
\Nat \to \Rat$ या फलनाविषयी असे म्हणा की \emph{$h$ ची सीमा $0$ आहे}
तेव्हाच आणि तेव्हाच, जेव्हा प्रत्येक धन $\epsilon \in \Rat$ साठी
$(\exists \ell \in \Nat)(\forall n > \ell)|h(n)| < \epsilon$.
\footnote{याची तुलना पर्यायी विभाग मधील $\lim_{x
\mathord{\rightarrow}\infty}f(x) = 0$ या व्याख्येशी करा.}
शिवाय $f$ आणि $g$ ही $\Nat \to \Rat$ अशी फलने असतील, तर
$(f-g)(n) = f(n) - g(n)$ असे माना. आता अशी व्याख्या करा:
\[
	f \Realequiv g \text{ तेव्हाच आणि तेव्हाच जेव्हा $(f-g)$ ची सीमा $0$ आहे}.
\]
$\Realequiv$ हा तुल्यता संबंध आहे हे तपासावे लागेल; आणि तो तसा आहे.
मग, हवे असल्यास, $\Nat \to \Rat$ अशा सर्व कोशी क्रमिकांचे
$\Realequiv$ खालील तुल्यतावर्ग म्हणजे वास्तव संख्या, अशी व्याख्या करता
येईल.

\begin{prob}
प्रत्येक $n$ साठी $f(n) = 0$ असे माना. $g(n) =
\frac{1}{(n+1)^2}$ असे माना. ही दोन्ही फलने कोशी क्रमिका आहेत, आणि
दोन्ही फलनांची सीमा $0$ आहे, हे दाखवा; त्यामुळे $f \sim_\Real g$
हेही दाखवा.
\end{prob}

हे केल्यानंतर, नेहमीप्रमाणे, $f$ हा घटक असलेल्या तुल्यतावर्गासाठी
आपण $\equivrep{f}{\Realequiv}$ असे लिहू. मात्र मजकूर वाचनीय ठेवण्यासाठी
पुढे अधोलिखित भाग वगळून फक्त $\equivrep{f}{}$ असे लिहू. प्रत्येक $q
\in \Rat$ साठी $q_{\Real} = \equivrep{c_{q}}{}$ असेही आपण ठरवतो,
जिथे प्रत्येक $n \in \Nat$ साठी $c_q(n) = q$ असे स्थिर फलन $c_{q}$
आहे. मग वास्तव संख्यांवरील मूलभूत संबंध आणि क्रिया परिभाषित करतो;
उदाहरणार्थ:
\begin{align*}
	\equivrep{f}{} + \equivrep{g}{} &= 	\equivrep{(f + g)}{} \\
	\equivrep{f}{} \times 	\equivrep{g}{} &= \equivrep{(f \times g)}{}
\end{align*}
जिथे $(f + g)(n) = f(n) + g(n)$ आणि $(f \times g)(n) = f(n) \times
g(n)$. $f$ आणि $g$ कोशी क्रमिका असतील तेव्हा $(f + g)$, $(f-g)$ आणि
$(f\times g)$ यांपैकी प्रत्येकही कोशी क्रमिका आहे हे अर्थात तपासावे
लागेल; तसे ते आहेत, आणि ही सिद्धता तुमच्यासाठी सोडली आहे.
\begin{quote}\small\textbf{स्रोतदुरुस्ती.} MRARITH-012: गोठवलेला मजकूर येथे बेरीज, वजाबाकी आणि गुणाकाराने कोशी क्रमिका मिळतात एवढीच संवृति तपासायला सांगतो. तुल्यतावर्गांवरील क्रिया सुव्याख्यायित होण्यासाठी प्रतिनिधी बदलल्यावर वर्ग बदलत नाही हेही तपासणे आवश्यक आहे; तो स्वतंत्र तपशील स्रोतामध्ये दिलेला नाही.\end{quote}

शेवटी क्रमाची संकल्पना परिभाषित करू. $\equivrep{f}{}$ हा वर्ग
\emph{धन} आहे तेव्हाच आणि तेव्हाच, जेव्हा $\equivrep{f}{}\neq 0_\Real$
आणि $(\exists \ell \in \Nat)(\forall n > \ell)0 < f(n)$
या दोन्ही अटी पूर्ण होतात. मग $\equivrep{f}{} <
\equivrep{g}{}$ तेव्हाच आणि तेव्हाच, जेव्हा $\equivrep{(g - f)}{}$
धन आहे, असे म्हणा.
\begin{quote}\small\textbf{स्रोतदुरुस्ती.} MRARITH-010: गोठवलेल्या इंग्रजीतील धनतेच्या व्याख्येत वास्तव तुल्यतावर्गाची तुलना 0\_Rat शी केली आहे. येथे त्याच रचनेतील वास्तव शून्य 0\_Real वापरले आहे.\end{quote}
ही व्याख्या सुव्याख्यायित आहे, म्हणजे ती तुल्यतावर्गातून निवडलेल्या
“प्रतिनिधी” फलनावर अवलंबून नाही, हे तपासावे लागेल.
हे केल्यानंतर या व्याख्यांमुळे योग्य बैजिक गुणधर्म मिळतात हे दाखवणे
बरेच सोपे आहे; म्हणजे:
\begin{thm}\label{sfr:arith:cauchy:thm:cauchyorderedfield}
कोशी क्रमिकांचे वरील तुल्यतावर्ग एक क्रमित क्षेत्र बनवतात.
\end{thm}
\begin{quote}\small\textbf{स्रोतदुरुस्ती.} MRARITH-011: गोठवलेल्या इंग्रजीतील प्रमेय, पुढील प्रश्न आणि संपूर्णतेचे विधान Cauchy sequences हेच क्षेत्र बनवतात अशी संक्षिप्त भाषा वापरतात. प्रत्यक्ष क्षेत्राचे घटक Realequiv संबंधाखालील तुल्यतावर्ग आहेत; मराठी मजकुरात ते स्पष्ट केले आहे. पुढील सिद्धतेत जेव्हा क्रमिकेवर वास्तव-क्रम लावला जातो तेव्हा तिचा तुल्यतावर्ग अभिप्रेत आहे.\end{quote}

\begin{proof}
सराव.
\end{proof}

\begin{prob}
कोशी क्रमिकांचे तुल्यतावर्ग एक क्रमित क्षेत्र बनवतात हे सिद्ध करा.
\end{prob}

अशा रीतीने रचलेल्या वास्तव संख्यांना संपूर्णता गुणधर्म आहे हे सिद्ध
करणे अधिक कठीण आहे; म्हणून त्याची सिद्धता आपण देऊ.

\begin{thm}
ऊर्ध्वबंध असलेल्या कोशी क्रमिकांच्या तुल्यतावर्गांच्या प्रत्येक अरिक्त
संचाला लघुतम ऊर्ध्वबंध असतो.
\end{thm}

\begin{proof}[सिद्धतेचा आराखडा]
$S$ हा ऊर्ध्वबंध असलेल्या कोशी क्रमिकांचा कोणताही अरिक्त संच असू द्या.
म्हणून काही $p \in \Rat$ असे आहे की $p_{\Real}$ हा $S$ चा ऊर्ध्वबंध
आहे. $r \in S$ असू द्या; मग काही $q \in \Rat$ असे आहे की
$q_{\Real} < r$. त्यामुळे $S$ ला लघुतम ऊर्ध्वबंध असेल, तर तो
$q_\Real$ आणि $p_\Real$ यांच्या दरम्यान (दोन्ही टोकांसह) असेल.

लघुतम ऊर्ध्वबंधाकडे खालून आणि वरून एकाच वेळी जवळ जात आपण तो शोधू.
विशेषतः, $f, g \colon \Nat \to \Rat$ अशी दोन फलने परिभाषित करू:
$f$ ने वरून आणि $g$ ने खालून लघुतम ऊर्ध्वबंधाकडे जवळ जावे हा उद्देश
आहे. सुरुवात अशा व्याख्यांनी करू:
\begin{align*}
	f(0) &= p \\
	g(0) &= q
\end{align*}
मग $a_n = \frac{f(n) + g(n)}{2}$ असेल, तर पुढील व्याख्या
करा:\footnote{ही पुनरावर्ती व्याख्या आहे. पण पुनरावर्ती व्याख्या ग्राह्य
आहेत असे मानण्याचे कोणतेही कारण आपण \emph{अजून} दिलेले नाही.}
\begin{align*}
	f(n+1) &=
	\begin{cases}
		a_n &\text{जर }(\forall h \in S)\equivrep{h}{} \leq (a_n)_\Real\\
		f(n)&\text{अन्यथा}
	\end{cases}\\
	g(n+1) &=
	\begin{cases}
		a_n &\text{जर }(\exists h \in S)\equivrep{h}{} \geq (a_n)_\Real\\
	 	g(n) &\text{अन्यथा}
	\end{cases}
\end{align*}
$f$ आणि $g$ या दोन्ही कोशी क्रमिका आहेत. (हे बऱ्यापैकी सहज तपासता
येते, पण तो सराव तुमच्यासाठी सोडला आहे.) प्रत्येक पायरीवर $f$ आणि
$g$ मधील फरक निम्मा होत असल्यामुळे $(f-g)$ या फलनाची सीमा $0$ आहे.
त्यामुळे $\equivrep{f}{} = \equivrep{g}{}$.

सर्वप्रथम $\equivrep{f}{}$ हा $S$ चा ऊर्ध्वबंध आहे, म्हणजे
$(\forall h \in S)\equivrep{h}{} \leq \equivrep{f}{}$, हे दाखवू.
(वाटेत आपण \ref{sfr:arith:cauchy:thm:cauchyorderedfield} वापरू.) $h \in S$ असू द्या
आणि व्याघातासाठी $\equivrep{f}{} < \equivrep{h}{}$ असे समजा; म्हणजे
$0_\Real < \equivrep{(h-f)}{}$. $f$ ही एकस्वनिक घटती कोशी क्रमिका
असल्यामुळे काही $n \in \Nat$ असे आहे की
$\equivrep{(c_{f(n)} - f)}{} < \equivrep{(h-f)}{}$. म्हणून:
\[
	(f(n))_\Real = \equivrep{c_{f(n)}}{} < \equivrep{f}{} + \equivrep{(h-f)}{} = \equivrep{h}{},
\]
पण रचनेनुसार $\equivrep{h}{} \leq (f(n))_\Real$ असल्यामुळे हा
व्याघात आहे.

आता $\equivrep{f}{} = \equivrep{g}{}$ हाच $S$ चा \emph{लघुतम}
ऊर्ध्वबंध आहे हे दाखवू. $j$ ही कोणतीही कोशी क्रमिका असू द्या आणि
$\equivrep{j}{} < \equivrep{g}{}$ असे समजा. वरच्याप्रमाणेच युक्तिवाद
केल्यास ($g$ हे \emph{वाढते} आहे या तथ्याचा उपयोग करून) काही $n \in
\Nat$ असे आहे की $\equivrep{j}{} < (g(n))_\Real$. पण रचनेनुसार काही
$h \in S$ असे आहे की $(g(n))_\Real \leq \equivrep{h}{}$; म्हणून
$\equivrep{j}{} < \equivrep{h}{}$. त्यामुळे $\equivrep{j}{}$ हा
$S$ चा ऊर्ध्वबंध नाही.
\end{proof}



\chapter{अनंत संच}\label{sfr:infinite::chap}
\begin{editorial}
अनंत संचांवरील हे प्रकरण टिम बटन यांच्या \emph{Open Set Theory}
या ग्रंथातून घेतले आहे.
\end{editorial}




\section{हिल्बर्ट यांचे हॉटेल}\label{sfr:infinite:hilbert:sec}

नैसर्गिक संख्यांचा संच स्पष्टपणे अनंत आहे. म्हणून, नैसर्गिक संख्या
स्वतःच \emph{आयत्या गृहीत धरायच्या} नसतील, तर नैसर्गिक संख्यांचा उल्लेख
न करता अनंत संचाचे वैशिष्ट्य सांगणे ही आपली पहिली पायरी असली पाहिजे.
सन 1924 मधील एका व्याख्यानात हिल्बर्ट यांनी मांडलेला पुढील दृष्टिकोन
सुबक आहे. ते आपल्याला अशी कल्पना करायला सांगतात:
\begin{quote}
[\ldots] सांत संख्येने खोल्या असलेले एक हॉटेल. या सर्व खोल्यांपैकी
प्रत्येक खोलीत नेमका एक पाहुणा असावा. आता पाहुण्यांनी कशाही प्रकारे
आपापल्या खोल्या बदलल्या, [पण] प्रत्येक खोलीत अजूनही एकापेक्षा अधिक
व्यक्ती राहणार नाही अशा रीतीने, तर एकही खोली मोकळी होणार नाही आणि
हॉटेलमालकाला या मार्गाने नव्याने येणाऱ्या पाहुण्यासाठी नवी जागा निर्माण
करता येणार नाही [\ldots
\textparagraph \ldots]

आता हॉटेलमध्ये $1$, $2$, $3$, $4$, $5$, \dots अशा क्रमांकांच्या अनंत
खोल्या असतील आणि प्रत्येक खोलीत नेमका एक पाहुणा असेल, असे आपण ठरवू.
नवा पाहुणा येताच मालकाने प्रत्येक जुन्या पाहुण्याला त्याच्या खोलीच्या
क्रमांकापेक्षा एकाने मोठा क्रमांक असलेल्या खोलीत हलवायचे असते; मग नव्याने
आलेल्या पाहुण्यासाठी खोली~$1$ मोकळी होईल.
	\begin{center}
		\begin{tikzpicture}[scale = .75]
		\foreach \x in {1, 2, 3, 4, 5, 6, 7, 8, 9}
		{
			\node (\x a) at (\x, 1) {\small{\x}};
			\node (\x b) at (\x, 2) {\small{\x}};
		}
		\node (dotsa) at (10, 1) {\small{\ldots}};
		\node (dotsb) at (10, 2) {\small{\ldots}};
		\draw[->] (1b)--(2a);
		\draw[->] (2b)--(3a);
		\draw[->] (3b)--(4a);
		\draw[->] (4b)--(5a);
		\draw[->] (5b)--(6a);
		\draw[->] (6b)--(7a);
		\draw[->] (7b)--(8a);
		\draw[->] (8b)--(9a);
		\draw[->] (9b)--(dotsa);
		\draw (1,1) circle (.4);
		\end{tikzpicture}
	\end{center}
	(Ewald आणि Sieg (2013, p.~730) मध्ये प्रकाशित; अनुवाद आमचा)
\end{quote}
निर्णायक मुद्दा असा की हिल्बर्ट यांच्या हॉटेलमध्ये अनंत खोल्या आहेत;
आणि असे म्हणण्याचा अर्थ काय, याची व्याख्या म्हणून आपण त्यांचे स्पष्टीकरण
घेऊ शकतो. खरे तर हाच डेडेकिंड यांचा दृष्टिकोन होता (अर्थातच येथे फार
मोठा कालविपर्यास करून तो मांडला आहे; डेडेकिंड यांची व्याख्या
1888 मधील आहे):

\begin{defn}\label{sfr:infinite:hilbert:defn:DedekindInfinite}
संच $A$ हा \emph{डेडेकिंड-अनंत} असतो जर आणि तरच $A$ पासून~$A$ च्या
एखाद्या उचित उपसंचाकडे एकास-एक फलन असते. म्हणजे, असा काही $o \in A$
आणि एकास-एक फलन $f \colon A \to A$ असतात की $o \notin \ran{f}$.
\end{defn}







\section{डेडेकिंड बीजसंरचना}\label{sfr:infinite:dedekind:sec}

नैसर्गिक संख्या केवळ अनंत असाव्यात असेच आपल्याला वाटत नाही; त्यांना काही
(बीजगणितीय) गुणधर्म असावेत असेही आपल्याला वाटते: बेरीज, गुणाकार इत्यादी
क्रियांखाली त्यांनी योग्य रीतीने वागले पाहिजे.

\emph{उत्तरवर्ती फलनाची} कल्पना मूलभूत मानून, पुढील गुणधर्म असलेल्या
वस्तू म्हणजे संख्या असे त्यांचे वैशिष्ट्य सांगावे, ही डेडेकिंड यांची
कल्पना होती:
\begin{enumerate}
	\item अशी $0$ ही संख्या आहे जी कोणत्याही संख्येची उत्तरवर्ती नाही
	\\म्हणजे, $0 \notin \ran{s}$
	\\म्हणजे, $\forall x\ s(x) \neq 0$
	\item भिन्न संख्यांच्या उत्तरवर्ती संख्या भिन्न असतात
	\\म्हणजे, $s$ हे एकास-एक फलन आहे
	\\म्हणजे, $\forall x \forall y (s(x) = s(y) \lif x = y)$
	\item\label{sfr:infinite:dedekind:repeatedapplication} प्रत्येक संख्या $0$ पासून
	उत्तरवर्ती फलन वारंवार लावून मिळते.
\end{enumerate}
पहिल्या दोन अटी प्रथम-क्रम तर्कशास्त्र वापरून सहज हाताळता येतात (वर
पाहा). पण केवळ प्रथम-क्रम तर्कशास्त्र वापरून \ref{sfr:infinite:dedekind:repeatedapplication}
ही अट हाताळता येत नाही. \ref{sfr:infinite:dedekind:repeatedapplication} ही अट पुढीलप्रमाणे
संचसैद्धान्तिक रूपात पुन्हा मांडणे ही डेडेकिंड यांची निर्णायक कामगिरी होती:
\begin{enumerate}
	\item[3$'$.] नैसर्गिक संख्यांचा संच हा \emph{उत्तरवर्ती फलनाखाली
	बंद} असलेला लघुतम संच आहे: म्हणजे, संचातील कोणत्याही घटक वर
	$s$ लावल्यास संचातीलच दुसरा घटक मिळतो.
\end{enumerate}
पण हे आपल्याला सावकाश आणि स्पष्टपणे मांडावे लागेल.

\begin{defn}\label{sfr:infinite:dedekind:Closure}
	कोणत्याही फलन $f$ साठी, संच $X$ हा $f$-\emph{बंद} असतो जर आणि तरच
	$(\forall x \in X)f(x) \in X$. आता कोणत्याही $o$ साठी अशी व्याख्या करू:
	$$\closureofunder{f}{o} = \bigcap\Setabs{X}{o \in X\text{ आणि }X
\text{ हा $f$-बंद आहे}}$$
\end{defn}

म्हणून $\closureofunder{f}{o}$ हा $o$ हा घटक असलेल्या सर्व
$f$-बंद संचांचा छेद आहे. त्यामुळे सहजज्ञानानुसार,
$\closureofunder{f}{o}$ हा $o$ हा घटक असलेला \emph{लघुतम}
$f$-बंद संच आहे. पुढील निष्कर्ष हे सहजज्ञान नेमकेपणाने मांडतो:
\begin{lem}\label{sfr:infinite:dedekind:closureproperties}
	कोणतेही फलन $f$ आणि कोणतेही $o \in A$ यांसाठी:
	\begin{enumerate}
		\item\label{sfr:infinite:dedekind:closurehaselem} $o \in \closureofunder{f}{o}$; आणि
		\item\label{sfr:infinite:dedekind:closureclosed} $\closureofunder{f}{o}$ हा $f$-बंद आहे; आणि
		\item\label{sfr:infinite:dedekind:closuresmallest} $X$ हा $f$-बंद असून $o \in
		X$ असेल, तर $\closureofunder{f}{o} \subseteq X$
	\end{enumerate}
\end{lem}

\begin{proof}
$o$ हा घटक असलेला किमान एक $f$-बंद संच आहे, हे लक्षात घ्या;
तो म्हणजे $\ran{f}\cup \{o\}$. म्हणून अशा \emph{सर्व} संचांचा छेद
$\closureofunder{f}{o}$ अस्तित्वात आहे. आता आपल्याला
\ref{sfr:infinite:dedekind:closurehaselem}--\ref{sfr:infinite:dedekind:closuresmallest} तपासावे लागतील.

\ref{sfr:infinite:dedekind:closurehaselem} बाबत: छेदातील सर्व संचांत $o$ हा घटक
असल्यामुळे $o \in \closureofunder{f}{o}$.

\ref{sfr:infinite:dedekind:closureclosed} बाबत: $x \in \closureofunder{f}{o}$ असे समजा.
म्हणून $o \in X$ आणि $X$ हा $f$-बंद असेल, तर $x \in X$; आणि $X$ हा
$f$-बंद असल्यामुळे $f(x) \in X$. म्हणून
$f(x) \in \closureofunder{f}{o}$.

\ref{sfr:infinite:dedekind:closuresmallest} बाबत: सर्वसाधारणपणे, $X \in C$ असेल, तर
$\bigcap C \subseteq X$.
\end{proof}

याचा उपयोग करून आपण असे म्हणू शकतो:

\begin{defn}
संच $A$, फलन $f \colon A \to A$ आणि असा काही $o \in A$ यांनी मिळून
\emph{डेडेकिंड बीजसंरचना} बनते, जर:
	\begin{enumerate}
		\item \label{sfr:infinite:dedekind:ded:proper} $o \notin \ran{f}$
		\item \label{sfr:infinite:dedekind:ded:injection} $f$ हे एकास-एक फलन आहे
		\item \label{sfr:infinite:dedekind:ded:closure} $A = \closureofunder{f}{o}$
	\end{enumerate}
\end{defn}

$A = \closureofunder{f}{o}$ असल्यामुळे, आधीच्या निष्कर्षावरून $A$ हा
$o$ हा घटक असलेला लघुतम $f$-बंद संच आहे. डेडेकिंड बीजसंरचना
स्पष्टपणे डेडेकिंड-अनंत असते; व्याख्येतील \ref{sfr:infinite:dedekind:ded:proper} आणि
\ref{sfr:infinite:dedekind:ded:injection} ही कलमे पाहा. पण कोणत्याही डेडेकिंड-अनंत संचाचे
डेडेकिंड बीजसंरचनेत रूपांतर करता येते, ही अधिक रोचक गोष्ट आहे.

\begin{thm}\label{sfr:infinite:dedekind:thm:DedekindInfiniteAlgebra}
डेडेकिंड-अनंत संच अस्तित्वात असेल, तर डेडेकिंड बीजसंरचना अस्तित्वात असते.
\end{thm}

\begin{proof}
$D$ हा डेडेकिंड-अनंत असू द्या. मग एकास-एक फलन $g \colon D \to D$
आणि घटक $o  \in D \setminus \ran{g}$ अस्तित्वात आहेत. आता $A =
\closureofunder{g}{o}$ असे मानू; \ref{sfr:infinite:dedekind:closureproperties} नुसार $A$
अस्तित्वात आहे आणि $o \in A$. $f = \funrestrictionto{g}{A}$ असे मानू.
$A, f, o$ हे मिळून डेडेकिंड बीजसंरचना बनतात, हे आपण दाखवू.

\ref{sfr:infinite:dedekind:ded:proper} बाबत: $o \notin \ran{g}$ आणि $\ran{f}
\subseteq \ran{g}$, म्हणून $o\notin \ran{f}$.

\ref{sfr:infinite:dedekind:ded:injection} बाबत: $g$ हे $D$ वरील एकास-एक फलन आहे; म्हणून
$f \subseteq g$ हेही एकास-एक फलन असले पाहिजे.

\ref{sfr:infinite:dedekind:ded:closure} बाबत: \ref{sfr:infinite:dedekind:closureproperties} नुसार $A$ हा
$g$-बंद आहे; त्यामुळे विशेषतः $A$ हा $f$-बंद आहे. म्हणून
\ref{sfr:infinite:dedekind:closureproperties} नुसार $\closureofunder{f}{o} \subseteq A$.
तसेच $\closureofunder{f}{o}$ हा $f$-बंद आहे आणि
$f = \funrestrictionto{g}{A}$ असल्यामुळे $\closureofunder{f}{o}$ हा
$g$-बंद आहे. म्हणून \ref{sfr:infinite:dedekind:closureproperties} नुसार
$A \subseteq \closureofunder{f}{o}$.
\end{proof}








\section{डेडेकिंड बीजसंरचना आणि अंकगणितीय विगमन}\label{sfr:infinite:induction:sec}

आता निर्णायक बाब अशी की एखादी डेडेकिंड बीजसंरचना---खरे तर,
\emph{कोणतीही} डेडेकिंड बीजसंरचना---नैसर्गिक संख्यांचे पर्यायी प्रतिरूप
म्हणून उपयोगी पडेल. हे पुढील सहज निष्पत्तीमुळे घडते:

\begin{thm}[अंकगणितीय विगमन]\label{sfr:infinite:induction:thm:dedinfiniteinduction}
$N, s, o$ यांनी मिळून डेडेकिंड बीजसंरचना बनते असे समजा. मग कोणत्याही
संच $X$ साठी:
	\begin{center}
		जर $o \in X$ आणि $(\forall {n} \in N \cap X){s}({n}) \in X$, {तर} $N \subseteq X$.
	\end{center}
\end{thm}

\begin{proof}
डेडेकिंड बीजसंरचनेच्या व्याख्येनुसार $N = \closureofunder{s}{o}$.
आता ${o} \in X$ आणि $(\forall {n} \in N)(n \in X \lif
{s}({n}) \in X)$ या दोन्ही अटी पूर्ण होत असतील, तर
$N = \closureofunder{s}{o} \subseteq X$.
\end{proof}

विगमन हा नैसर्गिक संख्यांचा वैशिष्ट्यपूर्ण गुणधर्म असल्यामुळे मुद्दा असा
आहे: कोणताही डेडेकिंड-अनंत संच दिला असता, आपण डेडेकिंड बीजसंरचना बनवू
शकतो आणि ती बीजसंरचना नैसर्गिक संख्यांचे पर्यायी प्रतिरूप म्हणून वापरू
शकतो.

हे मान्यच की \ref{sfr:infinite:induction:thm:dedinfiniteinduction} मध्ये विगमन
\emph{संचसैद्धान्तिक} भाषेत मांडले आहे. पण हे तत्त्व अधिक परिचित वाटेल
अशा भाषेत आपण सहज मांडू शकतो:

\begin{cor}\label{sfr:infinite:induction:natinductionschema}
$N, s, o$ यांनी मिळून डेडेकिंड बीजसंरचना बनते असे समजा. मग प्राचले असू
शकतील अशा कोणत्याही सूत्र $\phi(x)$ साठी:
\begin{center}
	जर $\phi(o)$ आणि $(\forall {n} \in N)(\phi(n)\lif
	\phi({s}({n})))$, {तर} $(\forall n \in N)\phi(n)$
\end{center}
\end{cor}

\begin{proof}
$X = \Setabs{n \in N}{\phi(n)}$ असे मानून आता
\ref{sfr:infinite:induction:thm:dedinfiniteinduction} वापरा.
\end{proof}

या निष्कर्षात आपण सूत्राला ``प्राचले असण्याबद्दल'' बोललो. याचा ढोबळ
अर्थ असा की कोणत्याही वस्तू $c_1, \ldots, c_k$ साठी आपण
$\phi(x, c_1, \ldots, c_k)$ बरोबर काम करू शकतो. अधिक नेमकेपणाने,
``प्राचले'' असा उल्लेख न करता निष्कर्ष पुढीलप्रमाणे मांडता येतो.
ज्याची सर्व मुक्त चरे दाखवली आहेत अशा कोणत्याही सूत्र
$\phi(x, v_1, \ldots, v_k)$ साठी:
	\begin{align*}
			\forall v_1 \ldots \forall v_k((&\phi(o, v_1,\ldots, v_k) \land {}\\
			&	(\forall x \in N)(\phi(x,v_1, \ldots, v_k) \lif \phi(s(x), v_1,\ldots, v_k))) \lif {}\\
			&\hspace{3em} (\forall x \in N)\phi(x, v_1,\ldots, v_k))
	\end{align*}
``प्राचले असणे'' असे म्हणाल्याने मजकूर वाचणे फारच सोपे होते, हे उघड आहे.
(पर्यायी विभाग मध्ये आपण ही युक्ती वारंवार वापरू.)

पुन्हा डेडेकिंड बीजसंरचनांकडे वळू: कोणतीही डेडेकिंड बीजसंरचना दिली
असता, बेरीज, गुणाकार आणि घातांक क्रिया यांची नेहमीची अंकगणितीय फलनेही
आपण परिभाषित करू शकतो. मात्र हे क्षुल्लक नाही आणि त्यात
\emph{पुनरावर्ती व्याख्येचे} तंत्र वापरावे लागते. हे तंत्र आपण बरेच नंतर,
त्यापेक्षा अधिक व्यापक संदर्भात मांडू आणि समर्थित करू.
(उत्सुक वाचकांनी पर्यायी विभाग नंतर याकडे पुन्हा
वळावे किंवा Potter (2004, pp.~95--8) यांसारखा पर्यायी विवेचनग्रंथ
वाचावा.) पण $N, s, o$ यांनी डेडेकिंड बीजसंरचना बनत असेल, तर शेवटी
आपल्याला पुढीलप्रमाणे अर्थ ठरवता येईल:
\begin{align*}
	{a} + {o} &= {a} & & & {a} \times {o} &= {o} & & & {a}^{o} &= s(o)\\
{a} + {s}({b}) &= {s}({a}+{b}) &&& {a} \times {s}({b}) &= ({a}\times {b}) + {a}  & & & {a}^{{s}({b})} &= {a}^{b} \times {a}
\end{align*}
आणि ही फलने आपल्या अपेक्षेप्रमाणे वागतात, हे दाखवता येईल.









\section{अनंत संचाच्या अस्तित्वाची
डेडेकिंड यांची ``सिद्धता''}\label{sfr:infinite:dedekindsproof:sec}

या प्रकरणात आपण डेडेकिंड बीजसंरचनांच्या साहाय्याने नैसर्गिक संख्यांचे
संचसैद्धान्तिक विवेचन दिले. \ref{sfr:arith:ref:sec} मध्ये आपण इतर
गोष्टींबरोबर विश्लेषणाच्या अंकगणितीकरणाच्या तात्त्विक महत्त्वाचा विचार
केला. आता येथे आपण जे साध्य केले आहे त्याच्या महत्त्वाचा विचार करायला हवा.

संपूर्ण \ref{sfr:arith::chap} मध्ये आपण नैसर्गिक संख्या आयत्या
मानल्या आणि त्यांचा वापर करून पूर्णांक, परिमेय संख्या आणि वास्तव संख्या
यांची स्पष्ट रचना केली. या प्रकरणात आपण नैसर्गिक संख्यांची स्पष्ट रचना
दिलेली नाही. आपण एवढेच दाखवले आहे की \emph{कोणताही डेडेकिंड-अनंत संच
दिला असता}, आपण असा संच परिभाषित करू शकतो जो आपल्याला~$\Nat$ ने जसे
वागावेसे वाटते अगदी तसाच वागेल.

त्यामुळे \emph{कोणत्या वस्तू नैसर्गिक संख्या आहेत} अशा सत्तामीमांसक
प्रश्नाचे उत्तर आपण दिले आहे, असा दावा अर्थातच करता येणार नाही. पण ही
चांगलीच गोष्ट आहे. शेवटी, \ref{sfr:arith:ref:sec} मध्ये आपण यावर
भर दिला होता की $\Real$ म्हणजे डेडेकिंड छेदांचा संच ही व्याख्या सोयीची
संकेतजनक व्याख्या न मानता \emph{शोध} मानणे चुकीचे ठरेल. डेडेकिंड छेद
वास्तव संख्यांच्या प्रमुख गणितीय गुणधर्मांचे उदाहरण देतात, हे निर्णायक
निरीक्षण आहे. येथेही तसेच: \emph{कोणतीही} डेडेकिंड बीजसंरचना नैसर्गिक
संख्यांच्या प्रमुख गणितीय गुणधर्मांचे उदाहरण देते, हे निर्णायक निरीक्षण
आहे. (खरे तर सर्व डेडेकिंड बीजसंरचना परस्पर \emph{समरूपी} आहेत, हे
सिद्ध करून डेडेकिंड यांनी हा मुद्दा ठामपणे दाखवला
(1888, प्रमेये 132--3). त्यामुळे आजचे अनेक
``संरचनावादी'' डेडेकिंड यांना पूर्वसूरी मानतात, यात आश्चर्य नाही.)

 शिवाय आपण नैसर्गिक संख्यांचा सिद्धान्त अशा अनौपचारिक साध्या
 संचसिद्धान्तात कसा अंतःस्थापित करायचा हे दाखवले आहे, जो स्वतः अजूनही
 बराच अनौपचारिक आहे, पण नैसर्गिक संख्या आयत्या गृहीत धरत नाही असे
 (वरवर तरी) दिसते. त्यामुळे डेडेकिंड यांचा स्वतःचा महत्त्वाकांक्षी
 प्रकल्प प्रत्यक्षात आणण्याच्या मार्गावर आपण असू शकतो. त्यांनी तो असा
 स्पष्ट केला:
\begin{quote}
	विज्ञानात सिद्ध करण्याजोगी कोणतीही गोष्ट सिद्धतेशिवाय मान्य करू नये.
	ही मागणी रास्त वाटत असली, तरी अगदी साध्या विज्ञानाचा पाया घालण्याच्या
	सर्वांत अलीकडील पद्धतींमध्येही ती पूर्ण झाली आहे असे मला मानता येत
	नाही; म्हणजे, संख्यासिद्धान्ताशी संबंधित तर्कशास्त्राच्या त्या भागातही.
	अंकगणित (बीजगणित, विश्लेषण) हे तर्कशास्त्राचाच केवळ एक भाग आहे असे
	म्हणताना माझा आशय असा की संख्या ही संकल्पना अवकाश आणि काळ यांच्या
	कल्पना किंवा अंतःप्रज्ञा यांपासून पूर्णपणे स्वतंत्र आहे---ती विचारांच्या
	शुद्ध नियमांचे थेट फलित आहे असेच मी मानतो.
	(Dedekind, 1888, प्रस्तावना)
\end{quote}
डेडेकिंड यांची धाडसी कल्पना अशी आहे. आपण नुकतेच पाहिले की केवळ
(अनौपचारिक) संचसिद्धान्त वापरून नैसर्गिक संख्या कशा घडवायच्या.
\ref{sfr:arith::chap} मध्ये आपण पाहिले की नैसर्गिक संख्या आणि काही
संचसिद्धान्त दिला असता वास्तव संख्यांची रचना कशी करायची. त्यामुळे कदाचित
``अंकगणित (बीजगणित, विश्लेषण)'' हे ``तर्कशास्त्राचाच केवळ एक भाग'' ठरते
(डेडेकिंड यांच्या ``तर्कशास्त्र'' या शब्दाच्या विस्तारित अर्थाने).

ही झाली कल्पना. पण क्षणभर थांबा. डेडेकिंड बीजसंरचनेची आपली रचना
(म्हणजे नैसर्गिक संख्यांचे आपले पर्यायी प्रतिरूप) डेडेकिंड-अनंत संचाच्या
अस्तित्वावर अवलंबून आहे. (मागे \ref{sfr:infinite:dedekind:thm:DedekindInfiniteAlgebra}
पाहा.) डेडेकिंड-अनंत संचाचे अस्तित्व ``तर्कशास्त्र'' किंवा ``विचारांचे
शुद्ध नियम'' यांच्या साहाय्याने प्रस्थापित करता आले नाही, तर हा प्रकल्प
अडून पडतो.

मग डेडेकिंड-अनंत संचाचे अस्तित्व ``विचारांच्या शुद्ध नियमांनी''
प्रस्थापित \emph{करता येते का}? डेडेकिंड यांचा प्रयत्न असा होता:
\begin{quote}
  माझे स्वतःचे विचारविश्व, म्हणजे माझ्या विचारांचे विषय होऊ शकणाऱ्या सर्व
  वस्तूंची समग्रता $S$, अनंत आहे. कारण $s$ हा~$S$ चा घटक दर्शवत असेल,
  तर $s$ हा माझ्या विचाराचा विषय होऊ शकतो हा विचार $s'$ स्वतः~$S$ चा
  घटक आहे. याला घटक~$s$ ची प्रतिमा $\phi(s)$ मानले, तर
  \dots~$S$ हा [डेडेकिंडच्या अर्थाने] अनंत आहे, हेच सिद्ध करायचे होते.
	(Dedekind, 1888, \S66)
\end{quote}
मुख्यतः अत्यंत काटेकोर गणितीय सिद्धतांनी भरलेल्या ग्रंथाच्या मध्यभागी
अशी गोष्ट सापडणे खरोखरच चकित करणारे आहे. येथे दोन टिप्पण्या कराव्याशा
वाटतात.

पहिली: या ``सिद्धतेला'' आज आपण ``गणितीय'' स्वरूप म्हणून ओळखू असे स्वरूप
जवळजवळ नाहीच. ती मानसशास्त्रीय वस्तूंबद्दल (विचारांबद्दल) बोलते आणि
त्यातही केवळ \emph{संभाव्य} विचारांबद्दल बोलते.

दुसरी: निदान आपण डेडेकिंड बीजसंरचना ज्या रीतीने मांडल्या आहेत त्या
रीतीनुसार, या ``सिद्धतेत'' एक सरळ तांत्रिक उणीव आहे. डेडेकिंड यांचा
युक्तिवाद यशस्वी झाला, तरी तो केवळ अनंत वस्तू (विशेषतः अनंत विचार)
अस्तित्वात आहेत एवढेच प्रस्थापित करतो. पण~$S$ ला अनेक घटक असलेला
एकच \emph{संच} का मानावे,~$S$ ला केवळ \emph{काही वस्तू} (अनेकवचनात) का मानू
नये, याचे कारणही डेडेकिंड यांनी द्यायला हवे.

डेडेकिंड यांना येथे उणीव दिसली नाही, यावरून कदाचित त्यांनी वापरलेला
``समग्रता'' हा शब्द \emph{आपण} वापरत असलेल्या ``संच'' या शब्दाशी
तंतोतंत जुळत नाही असे सूचित होते.\footnote{तसे तो जुळत नव्हता असे
मानण्याची इतर कारणेही आपल्याकडे आहेत; पाहा Potter (2004, p.~23).}
पण हे फार आश्चर्यकारक ठरणार नाही. मागील दोन प्रकरणांत आपण राबवलेला
प्रकल्प---नैसर्गिक संख्यांची ``रचना'' आणि त्यांच्यापासून पूर्णांक,
वास्तव संख्या व परिमेय संख्या यांची ``रचना''---संपूर्णपणे अनौपचारिक
रीतीने राबवला आहे. नेमके कोणते संच कधी अस्तित्वात असतात याविषयी अनेक
सामान्य तत्त्वे न मांडता, गरज पडेल तेव्हा आपण हा संच किंवा तो संच आयता
मानला. पण काही सामान्य तत्त्वांची आपल्याला \emph{गरज} आहे हे माहीत आहे;
नाहीतर आपल्याला रसेलच्या विरोधापत्तीला सामोरे जावे लागेल.

 आता अनौपचारिकतेच्या या मर्यादा ओलांडण्याची वेळ आली आहे.








\section{परिशिष्ट: श्रेडर--बर्नस्टाइनची सिद्धता}\label{sfr:infinite:card-sb:sec}

अनौपचारिक संचसिद्धान्ताचा निरोप घेण्यापूर्वी, आणखी एका अनौपचारिक
(पण प्रगत!) सिद्धतेचा विचार करायचा आहे. ही श्रेडर--बर्नस्टाइनची सिद्धता
आहे (\ref{sfr:siz:sb:thm:schroder-bernstein}): जर
$\cardle{A}{B}$ आणि $\cardle{B}{A}$, तर $\cardeq{A}{B}$; म्हणजे,
एकास-एक फलने $f \colon A \to B$ आणि $g \colon B \to A$ दिली असता
एकास-एक आच्छादन $h \colon A \to B$ असते.

या प्रकरणात आपण डेडेकिंड यांच्या \emph{संवरणांच्या} संकल्पनेचा वापर
केला. किंबहुना याच संकल्पनेचा उपयोग करून डेडेकिंड यांनी
श्रेडर--बर्नस्टाइनची एक सुंदर सिद्धता दिली; ती येथे मांडू. ही सिद्धता
Potter (2004, pp.~157--8) यांच्या विवेचनाला अगदी जवळून अनुसरते;
तेथील मांडणी किंचित वेगळी पण मूलतः समान आहे. इंटरनेटवर थोडा शोध
घेतल्यास असेही दिसेल की---$\sqrt{2}$ च्या अपरिमेयतेप्रमाणेच---या
प्रमेयाच्या \emph{अनेक} रोचक आणि वेगवेगळ्या सिद्धता आहेत.

\ref{sfr:infinite:dedekind:Closure} मधील संकेतनासारखेच संकेतन
वापरून, प्रत्येक संच $B$ आणि फलन~$f$ यांसाठी
\[
\Closureofunder{f}{B} = \bigcap \Setabs{X}{B \subseteq X
\text{ आणि $X$ हा $f$-बंद आहे}}
\]
असे मानू. अशा रीतीने परिभाषित केलेला $\Closureofunder{f}{B}$ हा~$B$
चा समावेश असलेला लघुतम $f$-बंद संच आहे; म्हणजे:

\begin{lem}\label{sfr:infinite:card-sb:Closureprops}
	कोणतेही फलन $f$ आणि कोणताही $B$ यांसाठी:
	\begin{enumerate}
		\item\label{sfr:infinite:card-sb:Closurehaselem} $B \subseteq \Closureofunder{f}{B}$; आणि
		\item\label{sfr:infinite:card-sb:Closureclosed} $\Closureofunder{f}{B}$ हा $f$-बंद आहे; आणि
		\item\label{sfr:infinite:card-sb:Closuresmallest} $X$ हा $f$-बंद असून $B
		\subseteq X$ असेल, तर $\Closureofunder{f}{B} \subseteq X$.
	\end{enumerate}
\end{lem}

\begin{proof}
\ref{sfr:infinite:dedekind:closureproperties} मधील सिद्धतेप्रमाणेच.
\end{proof}

श्रेडर--बर्नस्टाइनपर्यंत पोहोचण्यासाठी आपल्याला आणखी एका तथ्याची गरज आहे:

\begin{prop}\label{sfr:infinite:card-sb:sbhelper}
जर $A \subseteq B \subseteq C$ आणि $A \approx C$, तर $\cardeq{B}{C}$.
\end{prop}
\begin{quote}\small\textbf{स्रोतदुरुस्ती.} MRINF-003 पात्रता: गोठवलेल्या इंग्रजीतील अंतःस्थ cardeq रचना $A\approx B\approx C$ अशी वैध रीतीने विस्तारते; ती स्रोतदोष नाही. मराठीतील स्पष्ट $B\approx C$ निष्कर्ष त्याच तुल्यबलता-साखळीतील, पुढील उपयोगासाठी केंद्रित केलेली सममूल्य मांडणी आहे.\end{quote}

\begin{proof}
एकास-एक आच्छादन  $f \colon C \to A$ दिले असता, $F =
\Closureofunder{f}{C \setminus B}$ असे मानू आणि प्रांत $C$ असलेले फलन
$g$ पुढीलप्रमाणे परिभाषित करू:
\[
	g(x) =
	\begin{cases}
			f(x) &\text{जर $x \in F$}\\
			x & \text{अन्यथा}
	\end{cases}
\]
$g$ हे $C \to B$ असे एकास-एक आच्छादन आहे, हे आपण दाखवू. त्यावरून
$\comp{f^{-1}}{g} \colon A \to B$ हे एकास-एक आच्छादन आहे आणि सिद्धता
पूर्ण होते.

प्रथम आपला दावा असा आहे की $x \in F$ पण $y\notin F$ असेल, तर $g(x) \neq
g(y)$. विरोधाभासाने सिद्ध करण्यासाठी तसे नाही असे समजा; मग $y = g(y) =
g(x) = f(x)$. $x \in F$ आहे आणि \ref{sfr:infinite:card-sb:Closureprops} नुसार $F$ हा
$f$-बंद आहे, म्हणून $y = f(x) \in F$; हा विरोधाभास आहे.

आता $g(x) = g(y)$ असे समजा. वरील दाव्यानुसार, $x \in F$ तेव्हा आणि
केवळ तेव्हाच, जेव्हा $y \in F$. जर $x, y \in F$, तर $f(x) = g (x) =
g(y) = f(y)$; आणि $f$ हे एकास-एक आच्छादन असल्यामुळे $x = y$. जर
$x, y \notin F$, तर $x = g(x) = g(y) = y$. म्हणून $g$ हे
एकास-एक फलन आहे.

$\ran{g} = B$ हे दाखवणे उरते. म्हणून $x \in B \subseteq C$ स्थिर करू.
जर $x \notin F$, तर $g(x) = x$. जर $x \in F$, तर कोणत्यातरी
$y \in F$ साठी $x = f(y)$; कारण तसे नसते, तर $F \setminus \{x\}$ हा
$f$-बंद संच असता आणि त्यात $C\setminus B$ चा समावेश असता; पण
\ref{sfr:infinite:card-sb:Closureprops} नुसार हे अशक्य आहे. आता $g(y) = f(y) = x$.
\end{proof}

शेवटी, मुख्य निष्कर्षाची सिद्धता पुढीलप्रमाणे. फलन $h$ आणि संच $D$ दिले
असता आपण $\funimage{h}{D} = \Setabs{h(x)}{x \in D}$ अशी व्याख्या करतो,
हे आठवा.

\begin{proof}[श्रेडर--बर्नस्टाइनची सिद्धता] $f \colon A \to B$
आणि $g \colon B \to A$ ही एकास-एक फलने असू द्या.
$\funimage{f}{A} \subseteq B$ असल्यामुळे
$\funimage{g}{\funimage{f}{A}} \subseteq g[B] \subseteq A$.
तसेच, $\comp{f}{g} \colon A \to
\funimage{g}{\funimage{f}{A}}$ हे एकास-एक फलन आहे, कारण $g$ आणि
$f$ दोन्ही तशी आहेत; आणि त्याचा सहप्रांत ज्या रीतीने परिभाषित केला आहे,
त्यावरून $\comp{f}{g}$ हे खरे तर एकास-एक आच्छादन आहे. म्हणून
$\cardeq{\funimage{g}{\funimage{f}{A}}}{A}$; आणि त्यामुळे
\ref{sfr:infinite:card-sb:sbhelper} नुसार $h \colon A \to \funimage{g}{B}$ असे
एकास-एक आच्छादन आहे. शिवाय, $g^{-1}$ हे
$\funimage{g}{B} \to B$ असे एकास-एक आच्छादन आहे. म्हणून
$\comp{h}{g^{-1}} \colon A \to B$ हे एकास-एक आच्छादन आहे.
\end{proof}



\chapter{विधानीय तर्कशास्त्र: विन्यासमीमांसा आणि चिन्हार्थमीमांसा}\label{pl:syn::chap}
\begin{editorial}
  या भागात अभिजात विधानीय तर्कशास्त्रावरील सामग्री आहे. पहिले
  प्रकरण तुलनेने प्राथमिक स्वरूपाचे आहे आणि त्यात केवळ व्याख्या
  व निष्पत्त्या दिल्या आहेत; अनेक सिद्धता पूर्ण केलेल्या नसून त्या
  सरावासाठी सोडल्या आहेत. सिद्धता-पद्धतींवरील आणि संपूर्णता
  प्रमेयावरील सामग्री प्रथम-क्रम तर्कशास्त्राच्या भागातून, ``FOL''
  हा टॅग false असा ठेवून, समाविष्ट केली आहे. त्यामुळे विधेये,
  पदे आणि संख्यापक यांच्याशी संबंधित सर्व मजकूर वगळला जातो;
  तसेच संरचना~$\Struct{M}$ या संरचनांविषयीच्या चर्चेऐवजी
  सत्य-मूल्यांकने~$\pAssign{v}$ या सत्य-मूल्यांकनांविषयी चर्चा येते.

  या भागात अधिक तपशील समाविष्ट करण्यासाठी त्याचा विस्तार करणे,
  आणि सत्यता-फलनात्मक संपूर्णतेसारखे पुढील विषय व निष्पत्त्या जोडणे,
  असे नियोजन आहे.
\end{editorial}
\begin{editorial}
  हा केवळ व्याख्यांचा अतिशय संक्षिप्त आढावा आहे. सूत्रांवरील
  विगमनाने सिद्धता करण्याचा सुलभ परिचय आणि आणखी बरीच उदाहरणे
  देऊन त्याचा विस्तार करायला हवा.
\end{editorial}






\section{परिचय}\label{pl:syn:int:sec}

विधानीय तर्कशास्त्रात विधानीय संयोजके
$\lnot$, $\land$, $\lor$, $\lif$ आणि $\liff$ वापरून
विधानीय चलांपासून बनवलेल्या सूत्रांचा विचार
केला जातो. अंतःप्रेरणेने पाहता, विधानीय चल~$p$
हे सत्य किंवा असत्य असलेल्या एखाद्या वाक्याच्या किंवा विधानाच्या
जागी येते. सूत्र मधील विधानीय चलाचे
``सत्यतामूल्य'' ठरले की, विधानीय संयोजके वापरून त्यांच्यापासून
बनवलेल्या कोणत्याही सूत्रांचे सत्यतामूल्यही ठरते. विधानीय
तर्कशास्त्राला आपण \emph{सत्यता-फलनात्मक} म्हणतो, कारण त्याची
चिन्हार्थमीमांसा सत्यतामूल्यांच्या फलनांनी दिली जाते. विशेषतः,
विधानीय तर्कशास्त्रात सत्य आणि असत्य यांचे पुढील प्रकारचे निर्धारण
आपण विचारात घेत नाही: उदाहरणार्थ, एखादी गोष्ट केवळ परिस्थितीवश
सत्य असण्याऐवजी अनिवार्यपणे सत्य आहे का, ती सत्य आहे हे ज्ञात आहे
का, किंवा ती आत्ता सत्य आहे की पूर्वी सत्य होती अथवा पुढे सत्य
असेल. आपण फक्त सत्य ($\True$) आणि असत्य ($\False$) ही दोन
सत्यतामूल्ये विचारात घेतो. म्हणून एखादे विधान सत्यही नसते व
असत्यही नसते, किंवा ते केवळ अर्धसत्य असते, ही शक्यता चर्चेतून
वगळतो. तसेच, ज्यांच्यापासून बनवलेल्या सूत्राचे सत्यतामूल्य
त्याच्या भागांच्या सत्यतामूल्यांनी पूर्णपणे ठरते (त्याच्या अर्थाने,
उदाहरणार्थ, ठरत नाही), अशाच संयोजकांवर आपण लक्ष केंद्रित करतो.
विशेषतः, इंग्रजीतील सशर्त विधानांचे सत्यतामूल्य या अर्थाने
सत्यता-फलनात्मक असते की नाही, हा वादग्रस्त प्रश्न आहे. वास्तविक
अभिव्यंजन~$\lif$ सत्यता-फलनात्मक आहे; सत्यता-फलनात्मक नसलेल्या सशर्त
विधानांचा विचार इतर तर्कपद्धती करतात.

सत्यता-फलनात्मक विधानीय तर्कशास्त्राची उपपत्ती आणि अधिउपपत्ती
विकसित करण्यासाठी, आपण प्रथम त्याच्या व्यंजकांची विन्यासमीमांसा
आणि चिन्हार्थमीमांसा परिभाषित केली पाहिजे. संयोजके वापरून
विधानीय चलांपासून सूत्रे तयार करण्याची एक
पद्धत आपण वर्णन करू. पर्यायी व्याख्या शक्य आहेत. इतर पद्धती वेगळी
चिन्हे निवडतील, संयोजकांचे वेगळे संच मूलभूत म्हणून निवडतील आणि
कंसांचा वेगळ्या रीतीने वापर करतील (किंवा तथाकथित पोलिश
संकेतपद्धतीप्रमाणे कंस अजिबात वापरणार नाहीत). तरी सर्व दृष्टिकोनांत
सामाईक बाब ही की रचनानियम सूत्रांचा संच \emph{विगमनाने}
परिभाषित करतात. हे योग्य रीतीने केल्यास, रचनानियमांनुसार प्रत्येक
व्यंजक मूलतः एकाच रीतीने तयार होऊ शकते. त्यामुळे तयार होणारी
व्यंजके \emph{एकमेव रीतीने वाचनीय} असतात. परिणामी, त्यांची ही
विगमनाधारित व्याख्या आपल्याला त्याच पद्धतीने---विगमनाधारित
व्याख्येने---त्या व्यंजकांना अर्थ देता येतो, याची खात्री देते.

व्यंजकांना अर्थ देणे हा चिन्हार्थमीमांसेचा विषय आहे. विधानीय
तर्कशास्त्राच्या चिन्हार्थमीमांसेतील मध्यवर्ती संकल्पना म्हणजे
सत्य-मूल्यांकन मधील पूर्ति. सत्य-मूल्यांकन~$\pAssign{v}$ हे
विधानीय चलांना $\True$, $\False$ ही सत्यतामूल्ये
नेमून देते. कोणतेही सत्य-मूल्यांकन कोणत्याही सूत्र~$\varphi$
साठी $\pValue{v}(\varphi)$ हे सत्यतामूल्य ठरवते. सूत्राची
सत्य-मूल्यांकन~$\pAssign{v}$ मध्ये पूर्ति तेव्हा आणि केवळ तेव्हाच
होते जेव्हा $\pValue{v}(\varphi) = \True$; हे आपण $\pSat{v}{\varphi}$ असे
लिहितो. तार्किक संयोजकांची सत्यता-फलने वापरून, उदाहरणार्थ
$\varphi \land \psi$ ची पूर्ति $\varphi$ आणि~$\psi$ यांची पूर्ति होते (किंवा
होत नाही) यावरून परिभाषित करून, हा संबंधही~$\varphi$ च्या रचनेवरील
विगमनाने परिभाषित करता येतो.

वाक्यांसाठी असलेल्या $\pSat{v}{\varphi}$ या पूर्ति-संबंधाच्या आधारे
आपण सर्वतः सत्यता, तार्किक निष्पन्नता आणि पूर्ततायोग्यता या मूलभूत
चिन्हार्थविषयक संकल्पना परिभाषित करू शकतो. सूत्र हे सर्वतः
सत्य सूत्र, $\Entails \varphi$, असते, जर प्रत्येक सत्य-मूल्यांकन त्याची
पूर्ति करत असेल, म्हणजे कोणत्याही~$\pAssign{v}$ साठी
$\pValue{v}(\varphi) = \True$ असेल. सूत्रांच्या एखाद्या संचापासून,
$\Gamma \Entails \varphi$, ते तार्किकरीत्या निष्पन्न होते, जर~$\Gamma$
मधील सर्व सूत्रांची पूर्ति करणारे प्रत्येक सत्य-मूल्यांकन
~$\varphi$ चीही पूर्ति करत असेल. तसेच, सूत्रांचा एखादा संच
पूर्ततायोग्य असतो, जर एखादे सत्य-मूल्यांकन त्यातील सर्व सूत्रांची एकाच वेळी पूर्ति करत असेल. सूत्रे विगमनाने
परिभाषित केलेली असल्यामुळे, आणि पूर्तिही सूत्रांच्या रचनेवरील
विगमनाने परिभाषित केलेली असल्यामुळे, आपल्या चिन्हार्थमीमांसेचे
गुणधर्म सिद्ध करण्यासाठी आणि परिभाषित केलेल्या चिन्हार्थविषयक
संकल्पनांचा परस्परसंबंध दाखवण्यासाठी आपण विगमन वापरू शकतो.









\section{विधानीय सूत्र}\label{pl:syn:fml:sec}

विधानीय तर्कशास्त्रातील सूत्रे ही
\emph{विधानीय चले}, विधानीय
  स्थिर~$\lfalse$
    आणि विधानीय स्थिर~$\ltrue$ यांपासून \emph{तार्किक
  संयोजके} वापरून तयार होतात.

\begin{enumerate}
\item प्रगणनीय संच~$\PVar$: त्यातील विधानीय चले
  $\Obj p_0$, $\Obj p_1$, \dots
\item असत्यता साठीचे विधानीय स्थिर~$\lfalse$.
\item सत्यता साठीचे विधानीय स्थिर~$\ltrue$.
\item तार्किक संयोजके:
  \startycommalist
  \ycomma $\lnot$ (नकरण)%
  \ycomma $\land$ (संधियोग)%
  \ycomma $\lor$ (विकल्पयोग)%
  \ycomma $\lif$ (शर्त)%
  \ycomma $\liff$ (द्विशर्त)%
\item विरामचिन्हे: (, ), आणि स्वल्पविराम.
\end{enumerate}

विधानीय तर्कशास्त्राची ही भाषा आपण $\Lang L_0$ ने दर्शवतो.

.




\begin{intro}
वर आपण वापरलेल्या संज्ञांपेक्षा आणि चिन्हांपेक्षा वेगळ्या संज्ञा व
चिन्हे तुम्हाला परिचित असू शकतात. तर्कशास्त्राच्या ग्रंथांत (आणि
अध्यापनात) ``नकरण'' यासाठी सामान्यतः $\sim$, $\neg$, आणि~!\ यांपैकी
एखादे, तर ``संधियोग'' यासाठी $\wedge$, $\cdot$, आणि $\&$ यांपैकी
एखादे चिन्ह वापरतात. ``सशर्त (संकेतार्थक)'' किंवा ``अभिव्यंजन'' यांसाठी
$\rightarrow$, $\Rightarrow$, आणि $\supset$ ही चिन्हे सामान्यतः
वापरली जातात.
``द्विसशर्त'', ``द्वि-अभिव्यंजन'' किंवा
  ``(वास्तविक) सममूल्यता'' यांसाठी $\leftrightarrow$,
  $\Leftrightarrow$, आणि $\equiv$ ही चिन्हे वापरतात.
$\lfalse$ या चिन्हाला वेगवेगळ्या ठिकाणी
  ``असत्यता'', ``फाल्सम'', ``विसंगती'' किंवा ``तळ'' म्हणतात.
$\ltrue$ या चिन्हाला वेगवेगळ्या ठिकाणी
  ``सत्यता'', ``वेरम'' किंवा ``शिखर'' म्हणतात.
\end{intro}

\begin{defn}[सूत्र]
\label{pl:syn:fml:defn:formulas}
विधानीय तर्कशास्त्रातील \emph{सूत्रे} चा संच~$\Frm[L_0]$
पुढीलप्रमाणे विगमनाने परिभाषित केला जातो:
\begin{enumerate}
\item $\lfalse$ हे आण्विक सूत्र आहे.

\item $\ltrue$ हे आण्विक सूत्र आहे.

\item प्रत्येक विधानीय चल~$\Obj p_i$ हे आण्विक
  सूत्र आहे.

\item $\varphi$ हे सूत्र असेल, तर $\lnot \varphi$ हे
  सूत्र आहे.

\item $\varphi$ आणि $\psi$ ही सूत्रे असतील, तर $(\varphi \land
  \psi)$ हे सूत्र आहे.

\item $\varphi$ आणि $\psi$ ही सूत्रे असतील, तर $(\varphi \lor \psi)$
  हे सूत्र आहे.

\item $\varphi$ आणि $\psi$ ही सूत्रे असतील, तर $(\varphi \lif \psi)$
  हे सूत्र आहे.

\item $\varphi$ आणि $\psi$ ही सूत्रे असतील, तर $(\varphi \liff \psi)$
  हे सूत्र आहे.

\item इतर काहीही सूत्र नाही.
\end{enumerate}
\end{defn}

\begin{explain}
सूत्रांची ही व्याख्या \emph{विगमनाधारित व्याख्या} आहे. मूलतः,
आपण सूत्रांचा संच अनंतपणे अनेक टप्प्यांत तयार करतो. आरंभीच्या
टप्प्यात आपण सर्व आण्विक सूत्रांना सूत्रे घोषित करतो; हे व्याख्येतील
पहिल्या काही बाबींशी, म्हणजे
$\ltrue$, %
$\lfalse$, %
$\Obj p_i$ यांच्या बाबींशी जुळते. म्हणून ``आण्विक सूत्र''
म्हणजे या रूपातील कोणतेही सूत्र.

व्याख्येतील इतर बाबी आधीच तयार केलेल्या सूत्रांपासून नवी
सूत्रे तयार करण्याचे नियम देतात. दुसऱ्या टप्प्यात या नियमांनी
आण्विक सूत्रांपासून सूत्रे तयार करता येतात. तिसऱ्या
टप्प्यात आण्विक सूत्रे आणि दुसऱ्या टप्प्यात मिळालेली सूत्रे यांपासून
नवी सूत्रे तयार करतो, आणि ही प्रक्रिया पुढे चालू राहते. अशा एखाद्या
टप्प्यात शेवटी तयार होणारी प्रत्येक गोष्ट सूत्र असते; इतर
काहीही सूत्र नसते.
\end{explain}

द्विस्थानी संयोजक~$\ast$ वापरून $\psi$, $\chi$ पासून तयार केलेले
सूत्र $(\psi \ast \chi)$ लिहिताना आपण अनेकदा सर्वांत बाहेरची कंसांची
जोडी वगळून केवळ~$\psi \ast \chi$ असे लिहू.



\begin{defn}[विन्यासात्मक अनन्यता]
$\ident$ हे चिन्ह चिन्हमालांमधील विन्यासात्मक अनन्यता व्यक्त करते;
म्हणजे $\varphi \ident \psi$ तेव्हा आणि केवळ तेव्हाच, जेव्हा $\varphi$ आणि $\psi$
या समान लांबीच्या चिन्हमाला असतात आणि प्रत्येक स्थानी त्यांच्यात तेच
चिन्ह असते.
\end{defn}

$\ident$ या चिन्हाच्या दोन्ही बाजूंना जोडणीने मिळालेल्या चिन्हमाला
येऊ शकतात. उदाहरणार्थ, $\varphi \ident (\psi \lor \chi)$ याचा अर्थ असा:
आरंभीचा कंस, चिन्हमाला~$\psi$, $\lor$ हे चिन्ह, चिन्हमाला~$\chi$ आणि
शेवटचा कंस या क्रमाने जोडून मिळणारी चिन्हमाला आणि चिन्हमाला~$\varphi$
या एकच आहेत. असे असेल, तर $\varphi$ चे पहिले चिन्ह आरंभीचा कंस आहे;
$\varphi$ मध्ये दुसऱ्या चिन्हापासून सुरू होणारी $\psi$ ही उपचिन्हमाला आहे;
तिच्यानंतर $\lor$ येते; इत्यादी बाबी आपल्याला माहीत होतात.









\section{पूर्वतयारी}\label{pl:syn:pre:sec}

\begin{thm}[\emph{सूत्रांवरील विगमनाचे तत्त्व}]
  \label{pl:syn:pre:thm:induction}
  एखादा गुणधर्म~$P$ सर्व आण्विक सूत्रांसाठी लागू होत असेल आणि
  पुढील अटी पूर्ण करत असेल:
  \begin{enumerate}
    \item तो~$\varphi$ साठी लागू असताना $\lnot \varphi$ साठीही
      लागू होतो;
    \item तो~$\varphi$ आणि~$\psi$ यांच्यासाठी लागू असताना
      $(\varphi \land \psi)$ साठीही लागू होतो;
    \item तो~$\varphi$ आणि~$\psi$ यांच्यासाठी लागू असताना
      $(\varphi \lor \psi)$ साठीही लागू होतो;
    \item तो~$\varphi$ आणि~$\psi$ यांच्यासाठी लागू असताना
      $(\varphi \lif \psi)$ साठीही लागू होतो;
    \item तो~$\varphi$ आणि~$\psi$ यांच्यासाठी लागू असताना
      $(\varphi \liff \psi)$ साठीही लागू होतो;
  \end{enumerate}
  तर $P$ सर्व सूत्रांसाठी लागू होतो.
\end{thm}

\begin{proof}
  गुणधर्म~$P$ असलेल्या सर्व सूत्रांचा संग्रह~$S$ असू द्या.
  स्पष्टपणे $S \subseteq \Frm[L_0]$. \ref{pl:syn:fml:defn:formulas} मधील
  सर्व अटी $S$~पूर्ण करतो: त्यात सर्व आण्विक सूत्रे आहेत आणि
  तो संकारकांखाली बंद आहे.  $\Frm[L_0]$ हा असा सर्वांत लहान
  वर्ग आहे; म्हणून $\Frm[L_0] \subseteq S$. त्यामुळे $\Frm[L_0] = S$,
  आणि प्रत्येक सूत्राला गुणधर्म~$P$ आहे.
\end{proof}

\begin{prop}\label{pl:syn:pre:prop:balanced}
   $\Frm[L_0]$ मधील प्रत्येक सूत्र \emph{संतुलित} आहे; म्हणजे
   त्यात जेवढे डावे कंस आहेत तेवढेच उजवे कंस आहेत.
\end{prop}

\begin{prob}
\ref{pl:syn:pre:prop:balanced} सिद्ध करा.
\end{prob}

\begin{prop} \label{pl:syn:pre:prop:noinit}
कोणत्याही सूत्राचा कोणताही उचित आरंभीचा खंड सूत्र नसतो.
\end{prop}

\begin{prob}
  \ref{pl:syn:pre:prop:noinit} सिद्ध करा.
\end{prob}

\begin{prop}[एकमेव वाचनीयता]
$ \Frm[L_0]$ मधील कोणत्याही सूत्र~$\varphi$ चे पुढीलपैकी नेमक्या एका
रूपात रचनाविश्लेषण करता येते:
\begin{enumerate}
\item $\lfalse$.

\item $\ltrue$.

\item एखाद्या $\Obj p_n \in  \PVar$ साठी $\Obj p_n$.

\item एखाद्या सूत्र~$\psi$ साठी $\lnot \psi$.

\item काही सूत्रे $\psi$ आणि~$\chi$ यांच्यासाठी
  $(\psi \land \chi)$.

\item काही सूत्रे $\psi$ आणि~$\chi$ यांच्यासाठी
  $(\psi \lor \chi)$.

\item काही सूत्रे $\psi$ आणि~$\chi$ यांच्यासाठी
  $(\psi \lif \chi)$.

\item काही सूत्रे $\psi$ आणि~$\chi$ यांच्यासाठी
  $(\psi \liff \chi)$.
\end{enumerate}
शिवाय, हे रचनाविश्लेषण \emph{एकमेव} असते.
\end{prop}

\begin{proof}
$\varphi$ वरील विगमनाने सिद्ध करू. उदाहरणार्थ, $\varphi$ चे $(\psi \lif \chi)$ आणि
$(\psi' \lif \chi')$ अशी दोन भिन्न रचनाविश्लेषणे आहेत असे समजा. मग $\psi$
आणि $\psi'$ एकच असले पाहिजेत (अन्यथा त्यांपैकी एक हा दुसऱ्याचा उचित
आरंभीचा खंड असेल); म्हणून $\varphi$ ची ही दोन रचनाविश्लेषणे भिन्न असतील,
तर त्याचे कारण असेच असले पाहिजे की $\chi$ आणि $\chi'$ ही एकाच
चिन्हमालेची भिन्न रचनाविश्लेषणे आहेत; पण विगमन गृहीतकानुसार ते अशक्य
आहे.
\end{proof}


\begin{defn}[एकरूप आदेशन]
$\varphi$ आणि $\psi$ ही सूत्रे असतील आणि $\Obj p_i$ हे विधानीय चल असेल, तर $\varphi$ मधील $\Obj p_i$ ची प्रत्येक
आस्थिति $\psi$ च्या एका आस्थितीने बदलल्यावर मिळणारा परिणाम
$\Subst{\varphi}{\psi}{\Obj p_i}$ ने दर्शवला जातो; त्याचप्रमाणे,
$\Obj p_1$, \dots,~$\Obj p_n$ यांची सूत्रे $\psi_1$,
\dots,~$\psi_n$ यांनी एकाच वेळी केलेले आदेशन
$\SSubst{\varphi}{\subst{\psi_1}{\Obj p_1},\dots,\subst{\psi_n}{\Obj p_n}}$
ने दर्शवले जाते.
\end{defn}

\begin{prob} खालील पाच सूत्रांपैकी प्रत्येकासाठी, ते
 सूत्र \( \Subst{\varphi}{\psi}{\Obj p_i} \) या आदेशनाच्या रूपात व्यक्त
 करता येते का ते ठरवा; येथे \( \varphi \) हे अनुक्रमे (i) \( \Obj p_0 \);
 (ii) \( ( \lnot \Obj p_0 \land \Obj p_1) \); आणि (iii)
 \( ( ( \lnot \Obj p_0 \lif \Obj p_1 ) \land \Obj p_2 ) \) आहे.
 प्रत्येक बाबतीत संबद्ध आदेशन स्पष्ट करा.
  \begin{enumerate}
    \item \( \Obj p_1 \)
    \item \( ( \lnot \Obj p_0 \land \Obj p_0 ) \)
    \item \( ( ( \Obj p_0 \lor \Obj p_1 ) \land \Obj p_2 )  \)
    \item \( \lnot ( ( \Obj p_0 \lif \Obj p_1 ) \land \Obj p_2 ) \)
    \item \( (( \lnot ( \Obj p_0 \lif \Obj p_1 ) \lif ( \Obj p_0 \lor \Obj p_1 )) \land \lnot ( \Obj p_0 \land \Obj p_1 )) \)
  \end{enumerate}
\end{prob}

\begin{prob}
$\Subst{\varphi}{\psi}{p}$ ची विगमनाधारित, गणितीयदृष्ट्या काटेकोर व्याख्या
द्या.
\end{prob}









\section{रचनाक्रमिका}\label{pl:syn:fseq:sec}

सूत्रांची विगमनाधारित व्याख्या देणे आणि त्याला पूरक तंत्र म्हणून
सूत्रांचे गुणधर्म विगमनाने सिद्ध करणे हा सुबक व कार्यक्षम
दृष्टिकोन आहे. तथापि, सूत्रांच्या रचनेचा तळापासून वर जाणारा,
टप्प्याटप्प्याचा दृष्टिकोनही उपयुक्त ठरू शकतो; येथे आपण तो
\emph{रचनाक्रमिका} या संकल्पनेच्या साहाय्याने मांडतो.

\begin{defn}[सूत्रांसाठी रचनाक्रमिका]
\label{pl:syn:fseq:defn:fseq-frm}
भाषा~$\Lang L_0$ मधील चिन्हांपासून बनलेल्या चिन्हमालांची सांत क्रमिका
$\tuple{\varphi_0,\dotsc,\varphi_n}$ ही $\varphi$ साठी \emph{रचनाक्रमिका} असते, जर
$\varphi \ident \varphi_n$ आणि प्रत्येक $i \leq n$ साठी एकतर $\varphi_i$ हे आण्विक
सूत्र असेल किंवा पुढीलपैकी एखादी अट पूर्ण करणारे $j,k < i$ अस्तित्वात
असतील:
\begin{enumerate}
    \item $\varphi_i \ident \lnot \varphi_j$.%
    \item $\varphi_i \ident (\varphi_j \land \varphi_k)$.%
    \item $\varphi_i \ident (\varphi_j \lor \varphi_k)$.%
    \item $\varphi_i \ident (\varphi_j \lif \varphi_k)$.%
    \item $\varphi_i \ident (\varphi_j \liff \varphi_k)$.%
\end{enumerate}
\end{defn}

\begin{ex}
\[
\tuple{
    \Obj p_0,
    \Obj p_1,
    (\Obj p_1 \land \Obj p_0),
    \lnot (\Obj p_1 \land \Obj p_0)
}
\]
ही $\lnot (\Obj p_1 \land \Obj p_0)$ ची एक रचनाक्रमिका आहे; तसेच
पुढील क्रमिकाही तिची रचनाक्रमिका आहे:
\[
\tuple{
    \Obj p_0,
    \Obj p_1,
    \Obj p_0,
    (\Obj p_1 \land \Obj p_0),
    (\Obj p_0 \lif \Obj p_1),
    \lnot (\Obj p_1 \land \Obj p_0)
}.
\]
दुसऱ्या उदाहरणावरून दिसते की रचनाक्रमिकेत `अडगळ' असू शकते: अनावश्यक
असलेली किंवा रचनेत काहीही हातभार न लावणारी सूत्रे.
\end{ex}

\begin{prop}
\label{pl:syn:fseq:prop:formed}
$\Frm[L_0]$ मधील प्रत्येक सूत्र~$\varphi$ ला एक रचनाक्रमिका असते.
\end{prop}

\begin{proof}
$\varphi$ आण्विक आहे असे समजा. मग क्रमिका~$\tuple{\varphi}$ ही $\varphi$ साठी
रचनाक्रमिका आहे.
आता $\psi$ आणि~$\chi$ यांना अनुक्रमे $\tuple{\psi_0,\dotsc,\psi_n}$ आणि
$\tuple{\chi_0,\dotsc,\chi_m}$ या रचनाक्रमिका आहेत असे समजा.
\begin{enumerate}
  \item जर $\varphi \ident \lnot \psi$, तर
      $\tuple{\psi_0,\dotsc,\psi_n,\lnot \psi_n}$ ही $\varphi$ साठी
      रचनाक्रमिका आहे.
  \item जर $\varphi \ident (\psi \land \chi)$, तर
      $\tuple{\psi_0,\dotsc,\psi_n,\chi_0,\dotsc,\chi_m,(\psi_n \land \chi_m)}$
      ही $\varphi$ साठी रचनाक्रमिका आहे.
  \item जर $\varphi \ident (\psi \lor \chi)$, तर
      $\tuple{\psi_0,\dotsc,\psi_n,\chi_0,\dotsc,\chi_m,(\psi_n \lor \chi_m)}$
      ही $\varphi$ साठी रचनाक्रमिका आहे.
  \item जर $\varphi \ident (\psi \lif \chi)$, तर
      $\tuple{\psi_0,\dotsc,\psi_n,\chi_0,\dotsc,\chi_m,(\psi_n \lif \chi_m)}$
      ही $\varphi$ साठी रचनाक्रमिका आहे.
  \item जर $\varphi \ident (\psi \liff \chi)$, तर
      $\tuple{\psi_0,\dotsc,\psi_n,\chi_0,\dotsc,\chi_m,(\psi_n \liff \chi_m)}$
      ही $\varphi$ साठी रचनाक्रमिका आहे.
\end{enumerate}
सूत्रांवरील विगमनाच्या तत्त्वानुसार प्रत्येक सूत्राला एक
रचनाक्रमिका असते.
\end{proof}

याचे उलट विधानही आपण सिद्ध करू शकतो. ते महत्त्वाचे आहे, कारण त्यामुळे
सूत्रे परिभाषित करण्याच्या आपल्या दोन्ही पद्धती सममूल्य आहेत, म्हणजे
त्या समान परिणाम देतात, हे दिसते. तसेच रचनाक्रमिकांच्या लांबीवर साधे
विगमन वापरून सूत्रांविषयीची प्रमेये सिद्ध करता येतात.

\begin{lem}
\label{pl:syn:fseq:lem:fseq-init}
$\tuple{\varphi_0,\dotsc,\varphi_n}$ ही $\varphi_n$ साठी रचनाक्रमिका आहे आणि
$k \leq n$ असे समजा. मग $\tuple{\varphi_0,\dotsc,\varphi_k}$ ही $\varphi_k$ साठी
रचनाक्रमिका आहे.
\end{lem}

\begin{proof}
अभ्यासप्रश्न.
\end{proof}

\begin{thm}
\label{pl:syn:fseq:thm:fseq-frm-equiv}
$\Frm[L_0]$ हा भाषा~$\Lang L_0$ मधील ज्या सर्व चिन्हमालांना
रचनाक्रमिका आहे, त्यांचा संच आहे.
\end{thm}

\begin{proof}
भाषा~$\Lang L_0$ मधील ज्या सर्व चिन्हमालांना रचनाक्रमिका आहे, त्यांचा
संच~$F$ असू द्या. $\Frm[L_0] \subseteq F$ हे
\ref{pl:syn:fseq:prop:formed} मध्ये आपण पाहिले आहे; म्हणून आता उलट
समावेश सिद्ध करू.

$\varphi$ ला $\tuple{\varphi_0,\dotsc,\varphi_n}$ ही रचनाक्रमिका आहे असे समजा.
$n$ वरील प्रबल विगमनाने $\varphi \in \Frm[L_0]$ हे सिद्ध करू. लांबी
$m < n$ असलेली रचनाक्रमिका ज्या प्रत्येक चिन्हमालेला आहे, ती
$\Frm[L_0]$ मध्ये आहे, हे आपले विगमन गृहीतक आहे. रचनाक्रमिकेच्या
व्याख्येनुसार एकतर $\varphi_n$ आण्विक आहे किंवा पुढीलपैकी एखादी अट पूर्ण
करणारे $j,k < n$ अस्तित्वात असले पाहिजेत:
\begin{enumerate}
    \item $\varphi_n \ident \lnot \varphi_j$.%
    \item $\varphi_n \ident (\varphi_j \land \varphi_k)$.%
    \item $\varphi_n \ident (\varphi_j \lor \varphi_k)$.%
    \item $\varphi_n \ident (\varphi_j \lif \varphi_k)$.%
    \item $\varphi_n \ident (\varphi_j \liff \varphi_k)$.%
\end{enumerate}
आता प्रत्येक शक्यतेचा स्वतंत्र विचार करू. $\varphi_n$ आण्विक असेल, तर
$\varphi_n \in \Frm[L_0]$. त्याऐवजी $\varphi \ident
(\varphi_j \land \varphi_k)$ असे समजा.
\begin{quote}\small\textbf{स्रोतदुरुस्ती.} MRPL-002: गोठवलेल्या इंग्रजीत या एकाच शक्यतेत विन्यासात्मक अनन्यता-चिन्हाऐवजी पर्यायी सममूल्यता-चिन्ह आले आहे. आधीची स्पष्ट व्याख्या आणि इतर सर्व रचना-बाबी यांच्याशी जुळण्यासाठी येथे विन्यासात्मक अनन्यता वापरली आहे.\end{quote}
\ref{pl:syn:fseq:lem:fseq-init} नुसार
$\tuple{\varphi_0,\dotsc,\varphi_j}$ आणि $\tuple{\varphi_0,\dotsc,\varphi_k}$ या अनुक्रमे
$\varphi_j$ आणि~$\varphi_k$ साठी रचनाक्रमिका आहेत. या $\varphi$ च्या
रचनाक्रमिकेच्या उचित आरंभीच्या उपक्रमिका असल्यामुळे दोन्हींची लांबी
$n$ पेक्षा कमी आहे. म्हणून विगमन गृहीतकानुसार $\varphi_j$ आणि~$\varphi_k$ ही
$\Frm[L_0]$ मध्ये आहेत; त्यामुळे सूत्राच्या व्याख्येनुसार
$(\varphi_j \land \varphi_k)$ हेही त्यात आहे. इतर शक्यता समांतर युक्तिवादाने
सिद्ध होतात.
\end{proof}









\section{सत्य-मूल्यांकन आणि पूर्ति}\label{pl:syn:val:sec}

\begin{defn}[सत्य-मूल्यांकने]
``सत्य'' आणि ``असत्य'' ही दोन सत्यतामूल्ये असलेला संच
$\{\True, \False\}$ असू द्या. $\Lang{L_0}$ साठीचे
\emph{सत्य-मूल्यांकन} म्हणजे भाषेतील विधानीय चलांना
$\True$ किंवा $\False$ यांपैकी एक मूल्य देणारे फलन~$\pAssign{v}$;
म्हणजेच $\pAssign{v} \colon \PVar \to \{\True, \False \}$.
\end{defn}

\begin{defn}
  दिलेले सत्य-मूल्यांकन~$\pAssign{v}$ असताना, मूल्यन फलन
  $\pValue{v} \colon \Frm[L_0] \to \{\True, \False \}$ पुढीलप्रमाणे
  विगमनाने परिभाषित करू:
  \begin{align*}
    \pValue{v}(\lfalse) & = \False; \\
    \pValue{v}(\ltrue) & = \True; \\
    \pValue{v}(\Obj p_n) & = \pAssign{v}(\Obj p_n); \\
    \pValue{v}(\lnot \varphi) & = \begin{cases}
        \True & \text{जर } \pValue{v}(\varphi) = \False;\\
        \False & \text{अन्यथा.}
      \end{cases} \\
    \pValue{v}(\varphi \land \psi) & = \begin{cases}
        \True &
        \text{जर $\pValue{v}(\varphi) = \True$ आणि $\pValue{v}(\psi) = \True$;}\\
        \False &
        \text{जर $\pValue{v}(\varphi) = \False$ किंवा $\pValue{v}(\psi) = \False$}.
      \end{cases}\\
    \pValue{v}(\varphi \lor \psi) & = \begin{cases}
        \True &
        \text{जर $\pValue{v}(\varphi) = \True$ किंवा $\pValue{v}(\psi) = \True$;}\\
        \False &
        \text{जर $\pValue{v}(\varphi) = \False$ आणि $\pValue{v}(\psi) = \False$}.
      \end{cases}\\
    \pValue{v}(\varphi \lif \psi) & = \begin{cases}
        \True &
        \text{जर $\pValue{v}(\varphi) = \False$ किंवा $\pValue{v}(\psi) = \True$;}\\
        \False &
        \text{जर $\pValue{v}(\varphi) = \True$ आणि $\pValue{v}(\psi) = \False$}.
      \end{cases}\\
    \pValue{v}(\varphi \liff \psi) & = \begin{cases}
        \True &
        \text{जर $\pValue{v}(\varphi) = \pValue{v}(\psi)$;}\\
        \False &
        \text{जर $\pValue{v}(\varphi) \neq \pValue{v}(\psi)$}.
      \end{cases}
\end{align*}
\end{defn}

\begin{explain}
या बाबींना पुढील सत्यताकोष्टके अनुरूप आहेत:
\begin{center}

  \begin{tabular}{|c||c|} \hline
    $\varphi$ & $ \lnot \varphi$ \\
    \hline \hline
    $\True$ & $\False$ \\
    $\False$ & $\True$ \\
    \hline
  \end{tabular}


  \begin{tabular}{|cc||c|} \hline
    $\varphi$ & $\psi$ & $\varphi \land \psi$ \\
    \hline \hline
    $\True$ & $\True$ & $\True$ \\
    $\True$ & $\False$ & $\False$ \\
    $\False$ & $\True$ & $\False$ \\
    $\False$ & $\False$ & $\False$ \\
    \hline
  \end{tabular}


  \begin{tabular}{|cc||c|} \hline
    $\varphi$ & $\psi$ & $\varphi \lor \psi$ \\
    \hline \hline
    $\True$ & $\True$ & $\True$ \\
    $\True$ & $\False$ & $\True$ \\
    $\False$ & $\True$ & $\True$ \\
    $\False$ & $\False$ & $\False$ \\
    \hline
  \end{tabular}
%
\\[1em]

  \begin{tabular}{|cc||c|} \hline
    $\varphi$ & $\psi$ & $\varphi \lif \psi$ \\
    \hline \hline
    $\True$ & $\True$ & $\True$ \\
    $\True$ & $\False$ & $\False$ \\
    $\False$ & $\True$ & $\True$ \\
    $\False$ & $\False$ & $\True$ \\
    \hline
  \end{tabular}


  \begin{tabular}{|cc||c|} \hline
    $\varphi$ & $\psi$ & $\varphi \liff \psi$ \\
    \hline \hline
    $\True$ & $\True$ & $\True$ \\
    $\True$ & $\False$ & $\False$ \\
    $\False$ & $\True$ & $\False$ \\
    $\False$ & $\False$ & $\True$ \\
    \hline
  \end{tabular}

\end{center}
\end{explain}

\begin{prob}
$\Lang{L_0}$ मध्ये $\diamondsuit$ हा त्रिस्थानी संयोजक वाढवण्याचा विचार
करा; त्याचे मूल्यन पुढीलप्रमाणे दिले आहे:
\begin{gather*}
  \pValue{v}(\diamondsuit( \varphi , \psi , \chi ) ) = \begin{cases}
    \pValue{v}( \psi ) &
    \text{जर $\pValue{v}(\varphi) = \True$;}\\
    \pValue{v}( \chi ) &
    \text{जर $\pValue{v}(\varphi) = \False $}.
  \end{cases}
\end{gather*}
या संयोजकाचे सत्यताकोष्टक लिहा.
\end{prob}

\begin{thm}[स्थानिक निर्धारण]
\label{pl:syn:val:thm:LocalDetermination} $\pAssign{v_1}$ आणि
$\pAssign{v_2}$ ही अशी सत्य-मूल्यांकने आहेत की $\varphi$ मध्ये आस्थित असलेल्या
विधानीय चलांना ती समान मूल्ये देतात; म्हणजे, एखादे $\Obj p_n$ हे
सूत्र~$\varphi$ मध्ये आस्थित असेल तेव्हा $\pAssign{v_1}(\Obj p_n) =
\pAssign{v_2}(\Obj p_n)$. मग $\pValue{v_1}$ आणि $\pValue{v_2}$ हीही
$\varphi$ वर एकमत असतात; म्हणजे $\pValue{v_1}(\varphi) = \pValue{v_2}(\varphi)$.
\end{thm}

\begin{proof}
$\varphi$ वरील विगमनाने.
\end{proof}

\begin{defn}[पूर्ति]
\label{pl:syn:val:defn:satisfaction} आपण
  \emph{सत्य-मूल्यांकन~$\pAssign{v}$ कडून सूत्र~$\varphi$ ची पूर्ति},
  $\pSat{v}{\varphi}$, ही संकल्पना पुढीलप्रमाणे विगमनाने परिभाषित करू शकतो.
  (``$\pSat{v}{\varphi}$ नाही'' हे दर्शवण्यासाठी आपण $\pSat/{v}{\varphi}$ लिहितो.)
\begin{enumerate}
\item %
  \indcase{\varphi}{\lfalse}{$\pSat/{v}{\indfrm}$.}

\item %
  \indcase{\varphi}{\ltrue}{$\pSat{v}{\indfrm}$.}

\item \indcase{\varphi}{\Obj p_i}{$\pSat{v}{\indfrm}$ तेव्हा आणि केवळ
  तेव्हाच, जेव्हा $\pAssign{v}(\Obj p_i) = \True$.}

\item %
  \indcase{\varphi}{\lnot \psi}{$\pSat{v}{\indfrm}$ तेव्हा आणि केवळ
    तेव्हाच, जेव्हा $\pSat/{v}{\psi}$.}

\item %
  \indcase{\varphi}{(\psi \land \chi)}{$\pSat{v}{\indfrm}$ तेव्हा आणि केवळ
    तेव्हाच, जेव्हा $\pSat{v}{\psi}$ आणि $\pSat{v}{\chi}$.}

\item %
  \indcase{\varphi}{(\psi \lor \chi)}{$\pSat{v}{\indfrm}$ तेव्हा आणि केवळ
    तेव्हाच, जेव्हा $\pSat{v}{\psi}$ किंवा $\pSat{v}{\chi}$ (किंवा दोन्ही).}

\item %
  \indcase{\varphi}{(\psi \lif \chi)}{$\pSat{v}{\indfrm}$ तेव्हा आणि केवळ
    तेव्हाच, जेव्हा $\pSat/{v}{\psi}$ किंवा $\pSat{v}{\chi}$ (किंवा दोन्ही).}

\item %
  \indcase{\varphi}{(\psi \liff \chi)}{$\pSat{v}{\indfrm}$ तेव्हा आणि केवळ
    तेव्हाच, जेव्हा $\pSat{v}{\psi}$ आणि $\pSat{v}{\chi}$ ही दोन्ही, किंवा
    $\pSat{v}{\psi}$ आणि $\pSat{v}{\chi}$ यांपैकी कोणतेही नाही.}
\end{enumerate}
$\Gamma$ हा सूत्रांचा संच असेल, तर $\pSat{v}{\Gamma}$ तेव्हा आणि
केवळ तेव्हाच, जेव्हा प्रत्येक~$\varphi \in \Gamma$ साठी $\pSat{v}{\varphi}$.
\end{defn}

\begin{prop}\label{pl:syn:val:prop:sat-value}
  $\pSat{v}{\varphi}$ तेव्हा आणि केवळ तेव्हाच, जेव्हा
  $\pValue{v}(\varphi) = \True$.
\end{prop}

\begin{proof}
  $\varphi$ वरील विगमनाने.
\end{proof}

\begin{prob}
  \ref{pl:syn:val:prop:sat-value} सिद्ध करा.
\end{prob}









\section{चिन्हार्थविषयक संकल्पना}\label{pl:syn:sem:sec}

पुढील चिन्हार्थविषयक संकल्पना परिभाषित करू:

\begin{defn}
\begin{enumerate}
\item सूत्र~$\varphi$ हे \emph{पूर्ततायोग्य} असते, जर एखाद्या
  $\pAssign{v}$ साठी $\pSat{v}{\varphi}$; आणि कोणत्याही $\pAssign{v}$ साठी
  $\pSat{v}{\varphi}$ नसल्यास ते \emph{अपूर्ततायोग्य} असते;
\item सर्व सत्य-मूल्यांकने~$\pAssign{v}$ साठी $\pSat{v}{\varphi}$ असेल, तर
  सूत्र~$\varphi$ हे \emph{सर्वतः सत्य सूत्र} असते;
\item सूत्र~$\varphi$ हे पूर्ततायोग्य असेल, पण सर्वतः सत्य सूत्र नसेल,
  तर ते \emph{परिस्थितीवश} असते;
\item $\Gamma$ हा सूत्रांचा संच असेल, तर $\Gamma \Entails \varphi$
  (``$\Gamma$ पासून $\varphi$ तार्किकरीत्या निष्पन्न होते'') तेव्हा आणि केवळ
  तेव्हाच, जेव्हा $\pSat{v}{\Gamma}$ असलेल्या प्रत्येक
  सत्य-मूल्यांकन~$\pAssign{v}$ साठी $\pSat{v}{\varphi}$.
\item $\Gamma$ हा सूत्रांचा संच असेल, तर एखाद्या
  सत्य-मूल्यांकन~$\pAssign{v}$ साठी $\pSat{v}{\Gamma}$ असल्यास $\Gamma$
  हा \emph{पूर्ततायोग्य} असतो; अन्यथा $\Gamma$ हा
  \emph{अपूर्ततायोग्य} असतो.
\end{enumerate}
\end{defn}

\begin{prob}
 पुढील चार सूत्रांपैकी प्रत्येक सूत्र (a)~पूर्ततायोग्य,
 (b)~सर्वतः सत्य सूत्र आणि (c)~परिस्थितीवश आहे की नाही ते ठरवा.
\begin{enumerate}
  \item \( ( \Obj p_0 \lif ( \lnot \Obj p_1 \lif \lnot \Obj p_0 ) ) \).
  \item \( ( ( \Obj p_0 \land \lnot \Obj p_1 ) \lif ( \lnot \Obj p_0 \land \Obj p_2 )) \liff ( ( \Obj p_2 \lif \Obj p_0 ) \lif ( \Obj p_0 \lif \Obj p_1 )) \).
  \item \( ( \Obj p_0 \liff \Obj p_1 ) \lif ( \Obj p_2 \liff \lnot \Obj p_1 ) \).
  \item \( (( \Obj p_0 \liff ( \lnot \Obj p_1 \land \Obj p_2 )) \lor ( \Obj p_2 \lif ( \Obj p_0 \liff \Obj p_1 ))) \).
\end{enumerate}
\end{prob}

\begin{prop}
\label{pl:syn:sem:prop:semanticalfacts}
\begin{enumerate}
\item $\varphi$ हे सर्वतः सत्य सूत्र तेव्हा आणि केवळ तेव्हाच असते, जेव्हा
  $\emptyset \Entails \varphi$;
\item $\Gamma \Entails \varphi$ आणि $\Gamma \Entails \varphi \lif \psi$ असल्यास
  $\Gamma \Entails \psi$;
\item $\Gamma$ पूर्ततायोग्य असेल, तर $\Gamma$ चा प्रत्येक सांत उपसंचही
  पूर्ततायोग्य असतो;
\item \label{pl:syn:sem:def:monotonicity} एकस्वनिकता: $\Gamma \subseteq \Delta$
  आणि $\Gamma \Entails \varphi$ असल्यास $\Delta \Entails \varphi$ हेही असते;
\item \label{pl:syn:sem:def:Cut} संक्रमकता: $\Gamma \Entails \varphi$ आणि
  $\Delta \cup \{ \varphi\} \Entails \psi$ असल्यास $\Gamma \cup \Delta \Entails
  \psi$.
\end{enumerate}
\end{prop}

\begin{proof}
सरावासाठी.
\end{proof}

\begin{prob}
\ref{pl:syn:sem:prop:semanticalfacts} सिद्ध करा.
\end{prob}

\begin{prop}\label{pl:syn:sem:prop:entails-unsat}
  $\Gamma \Entails \varphi$ तेव्हा आणि केवळ तेव्हाच, जेव्हा
  $\Gamma \cup \{\lnot \varphi\}$ अपूर्ततायोग्य असतो.
\end{prop}

\begin{proof}
सरावासाठी.
\end{proof}

\begin{prob}
\ref{pl:syn:sem:prop:entails-unsat} सिद्ध करा.
\end{prob}

\begin{thm}[चिन्हार्थविषयक निगमन प्रमेय]
  \label{pl:syn:sem:thm:sem-deduction} $\Gamma \Entails \varphi \lif \psi$ तेव्हा आणि केवळ
  तेव्हाच, जेव्हा $\Gamma \cup \{\varphi\} \Entails \psi$.
\end{thm}

\begin{proof}
सरावासाठी.
\end{proof}

\begin{prob}
\ref{pl:syn:sem:thm:sem-deduction} सिद्ध करा.
\end{prob}



\chapter{सिद्धता-पद्धती}\label{fol:prf::chap}
\begin{editorial}
या प्रकरणात निष्पत्ती-पद्धतींविषयीची सर्वसाधारण सामग्री
एकत्रित केली आहे. एखादी विशिष्ट पद्धत वापरणारे पाठ्यपुस्तक,
ज्या प्रकरणात त्या पद्धतीचा परिचय करून दिला आहे त्या प्रकरणाच्या
सुरुवातीस प्रस्तावना-विभाग आणि संबंधित आढावा-विभाग समाविष्ट करू शकते.
\end{editorial}






\section{परिचय}\label{fol:prf:int:sec}

तर्कप्रणालींमध्ये सामान्यतः चिन्हार्थमीमांसा आणि निष्पत्ती-पद्धती हे दोन्ही घटक असतात. चिन्हार्थमीमांसेत सत्यता, पूर्ततायोग्यता,
वैधता आणि तार्किक निष्पन्नता अशा संकल्पनांचा विचार केला जातो.
तार्किक निष्पन्नता आणि वैधता प्रस्थापित करण्यासाठी निव्वळ विन्यासात्मक
रीत उपलब्ध करून देणे हा निष्पत्ती-पद्धतींचा उद्देश असतो. त्या
निव्वळ विन्यासात्मक असतात याचा अर्थ असा की अशा पद्धतीतील
निष्पत्ती ही एक सांत विन्यासात्मक वस्तू असते; सहसा ती
वाक्ये किंवा सूत्रे यांचा अनुक्रम (किंवा इतर एखादी सांत
मांडणी) असते. वाक्ये किंवा सूत्रे यांचा कोणताही दिलेला
अनुक्रम किंवा मांडणी ``योग्य'' आहे, हे यांत्रिक रीतीने पडताळता येते,
हा चांगल्या निष्पत्ती-पद्धतींचा गुणधर्म असतो.

प्रथम-क्रम तर्कशास्त्रासाठीच्या सर्वांत सोप्या (आणि इतिहासात प्रथम
आलेल्या) निष्पत्ती-पद्धती \emph{स्वयंसिद्धकीय} होत्या. अशा
पद्धतीत सूत्रांचा एखादा अनुक्रम तेव्हाच निष्पत्ती मानला
जातो, जेव्हा त्यातील प्रत्येक स्वतंत्र सूत्र एकतर
``स्वयंसिद्धकांच्या'' एका निश्चित संचात असते किंवा त्या अनुक्रमात
त्याच्या आधी आलेल्या सूत्रांपासून ठरावीक संख्येतील ``निगमन
नियमां''पैकी एखाद्या नियमाने निष्पन्न होते. शिवाय, सूत्र
स्वयंसिद्धक आहे का आणि ते इतर सूत्रांपासून एखाद्या निगमन
नियमानुसार योग्य रीतीने निष्पन्न होते का, हे यांत्रिक रीतीने पडताळता
येते. स्वयंसिद्धकीय निष्पत्ती-पद्धतींचे वर्णन करणे सोपे असते आणि
अधिउपपत्तीच्या पातळीवरही त्या हाताळणे सोपे असते; पण त्यांतील
निष्पत्त्या वाचणे आणि समजून घेणे कठीण असते, तसेच त्या तयार करणेही
कठीण असते.

निष्पत्त्या तयार करणे सोपे व्हावे किंवा पूर्ण झालेल्या
निष्पत्त्या समजून घेणे सोपे व्हावे, या उद्देशाने इतर
निष्पत्ती-पद्धती विकसित करण्यात आल्या आहेत. नैसर्गिक निगमन,
सत्यता-वृक्ष---ज्यांना टॅब्लो सिद्धता असेही म्हणतात---आणि क्रमवर्ती कलन
ही त्यांची उदाहरणे आहेत. काही निष्पत्ती-पद्धती विशेषतः
यांत्रिकीकरण डोळ्यांसमोर ठेवून रचल्या जातात. उदाहरणार्थ, निराकरण पद्धत
संगणकीय प्रणालीत कार्यान्वित करणे सोपे असते (पण तिच्यातील
निष्पत्त्या समजून घेणे जवळजवळ अशक्य असते). यांपैकी बहुतेक इतर
निष्पत्ती-पद्धती निष्पत्त्यांना अनुक्रमांऐवजी सूत्रांचे
वृक्ष म्हणून दर्शवतात. त्यामुळे निष्पत्तीचा कोणता भाग इतर
कोणत्या भागांवर अवलंबून आहे, हे पाहणे सोपे होते.

म्हणून प्रथम-क्रम तर्कशास्त्रासारखी एखादी तर्कप्रणाली घेतल्यास,
वाक्य हे \emph{प्रमेय} असणे म्हणजे काय आणि वाक्य हे
इतर काही वाक्यांपासून निष्पन्न करता येण्याजोगे असणे म्हणजे काय, यांची
वेगवेगळ्या निष्पत्ती-पद्धती वेगवेगळी नेमकी स्पष्टीकरणे देतील. हे
कसेही केले (स्वयंसिद्धकीय निष्पत्त्या, नैसर्गिक निगमने, क्रमवर्ती
निष्पत्त्या, सत्यता-वृक्ष किंवा निराकरण-खंडने यांच्या साहाय्याने),
तरी हे संबंध वैधता आणि तार्किक निष्पन्नता या चिन्हार्थविषयक संकल्पनांशी
जुळावेत, अशी आपली अपेक्षा असते. ``$\varphi$~हे प्रमेय आहे'' यासाठी
$\Proves \varphi$ आणि ``$\varphi$~हे~$\Gamma$ पासून निष्पन्न करता येण्याजोगे आहे'' यासाठी
``$\Gamma \Proves \varphi$'' लिहू या. $\Proves$~ची व्याख्या कोणतीही असो,
ती $\Entails$~शी जुळावी, म्हणजे:
\begin{enumerate}
\item $\Proves \varphi$ जर आणि फक्त जर $\Entails \varphi$
\item $\Gamma \Proves \varphi$ जर आणि फक्त जर $\Gamma \Entails \varphi$
\end{enumerate}
वरील विधानांतील ``फक्त जर'' दिशेला \emph{निर्दोषता} म्हणतात.
निष्पत्ती-पद्धत तेव्हा निर्दोष असते, जेव्हा निष्पन्नतेमुळे
तार्किक निष्पन्नतेची (किंवा वैधतेची) हमी मिळते. प्रत्येक समाधानकारक
निष्पत्ती-पद्धत निर्दोष असलीच पाहिजे; सदोष निष्पत्ती-पद्धती
मुळीच उपयोगी नसतात. कारण निष्पत्तीचा संपूर्ण उद्देशच वैधतेची
किंवा तार्किक निष्पन्नतेची विन्यासात्मक हमी देणे हा असतो. आपण मांडत
असलेल्या निष्पत्ती-पद्धतींची निर्दोषता पुढे सिद्ध करू.

उलट ``जर'' दिशेलाही महत्त्व आहे: तिला \emph{संपूर्णता} म्हणतात.
एखादी संपूर्ण निष्पत्ती-पद्धत इतकी समर्थ असते की $\varphi$~वैध असेल
तेव्हा $\varphi$~हे प्रमेय आहे, तसेच $\Gamma \Entails \varphi$ असेल तेव्हा
$\Gamma \Proves \varphi$, हे ती दाखवू शकते. संपूर्णता प्रस्थापित करणे अधिक
कठीण असते आणि काही तर्कप्रणालींसाठी संपूर्ण निष्पत्ती-पद्धतीच
नसतात. प्रथम-क्रम तर्कशास्त्रासाठी मात्र त्या आहेत. कुर्ट गोडेल यांनी
आपल्या १९२९ सालच्या प्रबंधात प्रथम-क्रम तर्कशास्त्राच्या
निष्पत्ती-पद्धतीची संपूर्णता सर्वप्रथम सिद्ध केली.

निष्पत्ती-पद्धतींशी जोडलेली आणखी एक संकल्पना म्हणजे
\emph{सुसंगतता}. वाक्यांच्या एखाद्या संचापासून कोणतीही गोष्ट
निष्पन्न करता येत असेल तर तो संच विसंगत म्हणतात; अन्यथा तो सुसंगत
असतो. विसंगतता ही अपूर्ततायोग्यतेची विन्यासात्मक समकक्ष संकल्पना आहे:
जसे अपूर्ततायोग्य संच चांगल्या उपपत्ती ठरत नाहीत, तसेच वाक्यांचे
विसंगत संचही चांगल्या उपपत्ती ठरत नाहीत; त्यांच्यात मूलभूत स्वरूपाचा
दोष असतो. वाक्यांचे सुसंगत संच सत्य किंवा उपयोगी असतीलच असे
नाही, पण ते निदान तार्किक उपयोगितेचा किमान उंबरठा ओलांडतात. वेगवेगळ्या
निष्पत्ती-पद्धतींसाठी वाक्यांच्या संचांच्या सुसंगततेची नेमकी
व्याख्या वेगळी असू शकते; पण~$\Proves$ प्रमाणेच सुसंगतताही तिच्या
चिन्हार्थविषयक समकक्ष संकल्पनेशी, म्हणजे पूर्ततायोग्यतेशी, जुळली पाहिजे.
$\Gamma$ सुसंगत आहे जर आणि फक्त जर तो पूर्ततायोग्य आहे, हे नेहमी खरे
असावे अशी आपली अपेक्षा आहे. येथे ``फक्त जर'' दिशा संपूर्णतेइतकीच ठरते
(सुसंगततेमुळे पूर्ततायोग्यतेची हमी मिळते), आणि ``जर'' दिशा
निर्दोषतेइतकीच ठरते (पूर्ततायोग्यतेमुळे सुसंगततेची हमी मिळते).
प्रत्यक्षात, अभिजात प्रथम-क्रम तर्कशास्त्रासाठी निर्दोषता आणि संपूर्णता
यांची ही दोन रूपे सममूल्य आहेत.









\section{क्रमवर्ती कलन}\label{fol:prf:seq:sec}

अनेक निष्पत्ती-पद्धती वाक्यांच्या मांडण्यांवर कार्य करतात,
तर क्रमवर्ती कलन \emph{क्रमवर्तींवर} कार्य करते. क्रमवर्ती ही पुढील रूपाची
अभिव्यक्ती असते:
\[
\varphi_1, \dots, \varphi_m \Sequent \psi_1, \dots, \psi_m,
\]
म्हणजे वाक्यांच्या अशा दोन अनुक्रमांची जोडी, ज्यांना क्रमवर्ती
चिन्ह~$\Sequent$ वेगळे करते. यांपैकी कोणताही अनुक्रम रिकामा असू शकतो.
क्रमवर्ती कलनातील निष्पत्ती म्हणजे क्रमवर्तींचा वृक्ष होय. त्यातील
सर्वांत वरच्या क्रमवर्ती विशेष रूपाच्या असतात (त्यांना ``आरंभीच्या
क्रमवर्ती'' किंवा ``स्वयंसिद्धके'' म्हणतात), आणि इतर प्रत्येक क्रमवर्ती
तिच्या लगत वर असलेल्या क्रमवर्तींपासून निगमन नियमांपैकी एखाद्या नियमाने
निष्पन्न होते. निगमन नियम एकतर क्रमवर्तींतील वाक्ये हाताळतात
(डाव्या किंवा उजव्या बाजूला त्यांची भर घालतात, त्यांना काढतात किंवा त्यांचा क्रम
बदलतात), किंवा नियमाच्या निष्कर्षात एखादे जटिल सूत्र आणतात.
उदाहरणार्थ, $\LeftR{\land}$ नियमामुळे $\varphi, \Gamma \Sequent \Delta$ पासून
$\varphi \land \psi, \Gamma \Sequent \Delta$ हा निष्कर्ष काढता येतो; आणि
$\RightR{\lif}$ नियमामुळे $\varphi, \Gamma \Sequent \Delta, \psi$ पासून
$\Gamma \Sequent \Delta, \varphi \lif \psi$ हा निष्कर्ष काढता येतो. हे कोणत्याही
$\Gamma$, $\Delta$, $\varphi$ आणि~$\psi$ साठी लागू आहे. (विशेषतः, $\Gamma$
आणि~$\Delta$ रिकामे असू शकतात.)

क्रमवर्ती कलनावर आधारित $\Proves$ संबंधाची व्याख्या पुढीलप्रमाणे केली
जाते: $\Gamma \Proves \varphi$ तेव्हा आणि केवळ तेव्हाच, जेव्हा असा एखादा
अनुक्रम~$\Gamma_0$ अस्तित्वात असतो की $\Gamma_0$ मधील प्रत्येक $\varphi$
हा~$\Gamma$ मध्ये असतो आणि ज्याच्या मुळाशी $\Gamma_0 \Sequent \varphi$ ही
क्रमवर्ती आहे अशी निष्पत्ती अस्तित्वात असते. $\Sequent \varphi$ या क्रमवर्तीची
निष्पत्ती असेल तर क्रमवर्ती कलनात $\varphi$ हे प्रमेय असते. उदाहरणार्थ,
$\Proves (\varphi \land \psi) \lif \varphi$ हे दाखवणारी निष्पत्ती पुढे दिली आहे:
\begin{prooftree}
  \Axiom$\varphi \fCenter \varphi$
  \RightLabel{\LeftR{\land}}
  \UnaryInf$\varphi \land \psi \fCenter \varphi$
  \RightLabel{\RightR{\lif}}
  \UnaryInf$\fCenter (\varphi \land \psi) \lif \varphi$
\end{prooftree}

प्रत्येक $\varphi \in \Gamma_0$ हे~$\Gamma$ मध्ये असेल अशा एखाद्या
$\Gamma_0 \Sequent$ ची निष्पत्ती अस्तित्वात असेल (म्हणजे
क्रमवर्तीची उजवी बाजू रिकामी असेल), तर क्रमवर्ती कलनात संच~$\Gamma$
विसंगत असतो. \RightR{\Weakening} नियम वापरून विसंगत संचापासून कोणतेही
वाक्य निष्पन्न करता येते.

क्रमवर्ती कलनाची निर्मिती १९३० च्या दशकात गरहार्ड गेंटझेन यांनी केली.
त्याची रचना पद्धतशीर आणि सममित असल्यामुळे निष्पत्त्यांची उपपत्ती
विकसित करण्यासाठी ते अतिशय उपयुक्त आकारिक साधन आहे. क्रमवर्ती कलनात
निष्पत्त्या शोधणे तुलनेने सोपे असते; पण या निष्पत्त्या अनेकदा
वाचायला कठीण असतात आणि त्यांचा सिद्धतांशी असलेला संबंधही काही वेळा
सहज दिसत नाही. तरी निष्पत्ती-पद्धतींकडे पाहण्याची ही अतिशय सुबक
रीत ठरली आहे, आणि अनेक तर्कप्रणालींसाठी क्रमवर्ती-कलन पद्धती उपलब्ध आहेत.









\section{नैसर्गिक निगमन}\label{fol:prf:ntd:sec}

नैसर्गिक निगमन ही प्रत्यक्ष तर्कविचाराचे अनुकरण करण्यासाठी घडवलेली
निष्पत्ती-पद्धत आहे (विशेषतः गणितज्ञ वापरतात त्या नियमबद्ध
तर्कविचाराचे). प्रत्यक्ष तर्कविचार अनेक ``नैसर्गिक'' पद्धतींनी पुढे जातो.
उदाहरणार्थ, प्रकरणांनुसार सिद्धतेमुळे विकल्पयोगी आधारविधानाच्या आधारे
निष्कर्ष प्रस्थापित करता येतो: विकल्पयोगाच्या प्रत्येक घटकापासून
तो निष्कर्ष निष्पन्न होतो, हे स्वतंत्रपणे दाखवले जाते. अप्रत्यक्ष सिद्धतेत
निष्कर्षाचे नकरण व्याघाताकडे नेते हे दाखवून निष्कर्ष प्रस्थापित करता येतो.
सोपाधिक सिद्धतेत पूर्वांगापासून उत्तरांग निष्पन्न होते हे दाखवून
``जर \dots तर \dots'' असा सोपाधिक दावा प्रस्थापित केला जातो. नैसर्गिक
निगमन म्हणजे अशा नैसर्गिक अनुमानांपैकी काहींचे आकारिकीकरण होय. प्रत्येक
तार्किक संयोजक आणि संख्यापक यांना दोन नियम असतात—एक प्रवेशन नियम आणि एक
विलोपन नियम—आणि हे नियम प्रत्येकी अशाच एका नैसर्गिक अनुमान-पद्धतीशी
जुळतात. उदाहरणार्थ, $\Intro{\lif}$ सोपाधिक सिद्धतेशी, तर $\Elim{\lor}$
प्रकरणांनुसार सिद्धतेशी जुळतो. $\Elim{\land}$ हा विशेषतः सोपा नियम आहे;
त्यामुळे $\varphi \land \psi$ पासून~$\varphi$ (किंवा $\psi$) असा निष्कर्ष काढता येतो.

नैसर्गिक निगमनाला इतर निष्पत्ती-पद्धतींपासून वेगळे करणारे एक
वैशिष्ट्य म्हणजे त्यात गृहीतकांचा होणारा वापर. नैसर्गिक निगमनातील निष्पत्ती म्हणजे सूत्रांचा वृक्ष होय. या सूत्रांच्या
वृक्षाच्या मुळाशी एकच सूत्र असते, आणि वृक्षाची ``पाने'' ही अशी
सूत्रे असतात की ज्यांच्यापासून निष्कर्ष काढला जातो. नैसर्गिक
निगमनात पानांवरील काही सूत्रे ही निष्पत्तीमध्ये भूमिका बजावतात;
पण निष्पत्ती निष्कर्षापर्यंत पोहोचेपर्यंत ती ``वापरून संपवलेली'' असतात.
प्रत्यक्ष तर्कविचारात फक्त थोड्या काळापुरती लागू असणारी गृहीतके मांडण्याच्या
प्रथेशी हे जुळते. उदाहरणार्थ, प्रकरणांनुसार सिद्धतेत आपण विकल्पयोगाच्या
प्रत्येक घटकाची सत्यता गृहीत धरतो; सोपाधिक सिद्धतेत पूर्वांगाची सत्यता
गृहीत धरतो; आणि अप्रत्यक्ष सिद्धतेत निष्कर्षाच्या नकरणाची सत्यता गृहीत
धरतो. अशी तात्पुरती गृहीतके मांडणे आणि एखादी मधली पायरी प्रस्थापित करण्यासाठी
ती नंतर बाजूला काढणे, हे नैसर्गिक निगमनाचे खास वैशिष्ट्य आहे. नैसर्गिक
निगमनातील निष्पत्तीच्या पानांवरील सूत्रांना गृहीतके म्हणतात,
आणि काही निगमन नियम त्यांना ``मुक्त'' करू शकतात. उदाहरणार्थ, काही
गृहीतकांपासून~$\psi$ ची अशी निष्पत्ती असेल आणि त्या गृहीतकांत~$\varphi$
असेल, तर $\Intro{\lif}$ नियमामुळे~$\varphi \lif \psi$ असा निष्कर्ष काढता येतो
आणि~$\varphi$ या रूपातील कोणतेही गृहीतक मुक्त करता येते. (कोणती गृहीतके
कोणत्या अनुमानाच्या वेळी मुक्त केली जातात याची नोंद ठेवण्यासाठी, त्या
अनुमानाला आणि त्याने मुक्त केलेल्या गृहीतकांना
एक क्रमांक देतो.) निष्पत्तीच्या शेवटी जी गृहीतके मुक्त न केलेली
राहतात, ती निष्कर्षाच्या सत्यतेसाठी एकत्रितपणे पुरेशी असतात; म्हणून
निष्पत्ती ही प्रस्थापित करते की तिची मुक्त न केलेली गृहीतके तिचा निष्कर्ष तार्किकरीत्या निष्पन्न करतात.

नैसर्गिक निगमनावर आधारित $\Gamma \Proves \varphi$ हा संबंध तेव्हा आणि केवळ
तेव्हाच लागू होतो, जेव्हा अशी निष्पत्ती असते की त्यात वृक्षातील
शेवटचे वाक्य हे~$\varphi$ असते आणि मुक्त न केलेले प्रत्येक पान
हे~$\Gamma$ मध्ये असते. नैसर्गिक निगमनात $\varphi$ हे तेव्हा आणि केवळ तेव्हाच
प्रमेय असते, जेव्हा अशी निष्पत्ती असते की त्यातील शेवटचे
वाक्य हे~$\varphi$ असते आणि सर्व गृहीतके मुक्त झालेली असतात.
उदाहरणार्थ, $\Proves (\varphi \land \psi) \lif \varphi$ हे दाखवणारी निष्पत्ती
पुढे दिली आहे:
\begin{prooftree}
  \AxiomC{$\Discharge{\varphi \land \psi}{1}$}
  \RightLabel{\Elim{\land}}
  \UnaryInfC{$\varphi$}
  \DischargeRule{\Intro{\lif}}{1}
  \UnaryInfC{$(\varphi \land \psi) \lif \varphi$}
\end{prooftree}
$1$ ही खूण दाखवते की $\varphi \land \psi$ हे गृहीतक \Intro{\lif} अनुमानाच्या
वेळी मुक्त केले जाते.

नैसर्गिक निगमनात संच~$\Gamma$ तेव्हा आणि केवळ तेव्हाच विसंगत असतो,
जेव्हा $\Gamma \Proves \lfalse$. \FalseInt{} नियमामुळे विसंगत संचापासून
कोणतेही वाक्य निष्पन्न करता येते.

नैसर्गिक निगमन पद्धती १९३० च्या दशकात गरहार्ड गेंटझेन आणि स्टॅनिस्लॉ
यास्कोव्हस्की (Stanis\l{}aw Ja\'skowski) यांनी विकसित केल्या; पुढे डॅग
प्रावित्झ आणि फ्रेडरिक फिच यांनी त्यांचा आणखी विकास केला. त्यांतील अनुमाने
सिद्धतेच्या नैसर्गिक पद्धतींचे अनुकरण करत असल्यामुळे तत्त्वज्ञ नैसर्गिक
निगमनाला पसंती देतात.
फिच यांनी विकसित केलेल्या आवृत्त्या तर्कशास्त्राच्या प्रास्ताविक
पाठ्यपुस्तकांमध्ये अनेकदा वापरल्या जातात. तर्कशास्त्राच्या तत्त्वज्ञानात
नैसर्गिक निगमनाचे नियम तार्किक कारकांचे अर्थ देतात, असे काही वेळा मानले
गेले आहे (``सिद्धता-उपपत्तीय चिन्हार्थमीमांसा'').









\section{टॅब्लो}\label{fol:prf:tab:sec}

अनेक निष्पत्ती-पद्धती वाक्यांच्या मांडणीवर काम करतात, तर
टॅब्लो चिन्हांकित सूत्रांवर काम करतात. चिन्हांकित सूत्र म्हणजे
सत्यतामूल्याचे चिन्ह ($\True$ किंवा $\False$) आणि वाक्य यांची जोडी:
\[
\sFmla{\True}{\varphi} \text{ or } \sFmla{\False}{\varphi}.
\]
टॅब्लोमध्ये चिन्हांकित सूत्रे खालीच्या दिशेने शाखा फुटणाऱ्या
वृक्षात मांडलेली असतात. त्याची सुरुवात काही \emph{गृहीतकांपासून} होते आणि
त्यांच्या वर असलेल्या चिन्हांकित सूत्रांपैकी एखाद्याला निगमन नियमांपैकी एक
नियम लावल्यामुळे मिळणाऱ्या चिन्हांकित सूत्रांनी तो पुढे चालू राहतो.
प्रत्येक नियमामुळे एखाद्या शाखेच्या शेवटी एक किंवा अधिक चिन्हांकित सूत्रे
जोडता येतात, किंवा दोन चिन्हांकित सूत्रे शेजारी शेजारी जोडता येतात---दुसऱ्या
स्थितीत शाखा दोन शाखांत फुटते आणि नव्याने जोडलेली दोन चिन्हांकित सूत्रे त्या दोन शाखांची टोके बनतात.

संमिश्र चिन्हांकित सूत्राला नियम लावल्यावर तात्काळ उप-सूत्रे असलेली
चिन्हांकित सूत्रे जोडली जातात. हे नियम जोड्यांनी येतात—दोन्ही चिन्हांसाठी
प्रत्येकी एक नियम. उदाहरणार्थ, $\TRule{\True}{\land}$ हा नियम
$\sFmla{\True}{\varphi \land \psi}$ ला लागू होतो आणि
$\sFmla{\True}{\varphi \land \psi}$ असलेल्या कोणत्याही शाखेच्या शेवटी
$\sFmla{\True}{\varphi}$ आणि~$\sFmla{\True}{\psi}$ ही दोन्ही चिन्हांकित सूत्रे
जोडू देतो; तसेच $\TRule{\False}{\varphi \land \psi}$ हा नियम
$\sFmla{\False}{\varphi}$ आणि $\sFmla{\False}{\psi}$ शेजारी शेजारी जोडून शाखा
फोडू देतो. टॅब्लोच्या प्रत्येक शाखेत $\sFmla{\True}{\varphi}$ आणि
$\sFmla{\False}{\varphi}$ ही चिन्हांकित सूत्रांची जुळणारी जोडी असेल, तर तो
बंद असतो.

टॅब्लोवर आधारित $\Proves$ संबंधाची व्याख्या पुढीलप्रमाणे केली जाते:
$\Gamma \Proves \varphi$ तेव्हा आणि केवळ तेव्हाच, जेव्हा असा काही सांत संच
~$\Gamma_0 = \{\psi_1, \dots, \psi_n\} \subseteq \Gamma$ असतो की पुढील
गृहीतकांसाठी बंद टॅब्लो असतो:
\[
\{\sFmla{\False}{\varphi}, \sFmla{\True}{\psi_1}, \dots, \sFmla{\True}{\psi_n}\}
\]
उदाहरणार्थ, $\Proves (\varphi \land \psi) \lif \varphi$ हे दाखवणारा बंद टॅब्लो
पुढे दिला आहे:
\begin{oltableau}
  [\sFmla{\False}{(\varphi \land \psi) \lif \varphi}, just = \TAss
    [\sFmla{\True}{\varphi \land \psi}, just = {\TRule{\False}{\lif}[1]}
      [\sFmla{\False}{\varphi}, just = {\TRule{\False}{\lif}[1]}
        [\sFmla{\True}{\varphi}, just = {\TRule{\True}{\land}[2]}
          [\sFmla{\True}{\psi}, just = {\TRule{\True}{\land}[2]}, close]
        ]
      ]
    ]
  ]
\end{oltableau}

काही $\psi_i \in \Gamma$ साठी पुढील गृहीतकांचा बंद टॅब्लो असेल, तर
संच $\Gamma$ हा टॅब्लो कलनात विसंगत असतो:
\[
\{\sFmla{\True}{\psi_1}, \dots, \sFmla{\True}{\psi_n}\}
\]

टॅब्लोचा शोध १९५० च्या दशकात एव्हर्ट बेथ आणि याक्को हिंटिक्का यांनी
स्वतंत्रपणे लावला; रेमंड स्मलियन यांनी त्यांचे सुलभीकरण करून ते लोकप्रिय
केले. टॅब्लो तयार करणे ही अत्यंत पद्धतशीर प्रक्रिया असल्यामुळे त्यांचा
वापर अगदी सोपा आहे. टॅब्लोच्या पद्धतशीर स्वरूपामुळे त्यांची
संगणकीय अंमलबजावणीही सुलभ होते. तथापि, टॅब्लो वाचणे अनेकदा कठीण
असते आणि त्यांचा सिद्धतांशी असलेला संबंध काही वेळा सहज दिसत नाही. ही
दृष्टीही बरीच व्यापक आहे आणि अनेक भिन्न तर्कशास्त्रांसाठी टॅब्लो
पद्धती उपलब्ध आहेत. दिलेली वाक्ये (किंवा त्यांचे संच) पूर्ण करणाऱ्या
संरचना शोधण्यासही टॅब्लो मदत करतात: संच पूर्ततायोग्य असेल,
तर त्याचा बंद टॅब्लो असणार नाही; म्हणजेच कोणत्याही टॅब्लोमध्ये
एक खुली शाखा असेल. त्या शाखेवर लागू करता येणारा प्रत्येक नियम लागू केलेला
असेल, तर खुल्या शाखेवरून पूर्ति करणारी संरचना ``वाचून काढता'' येते.
क्रमवर्ती कलनाशीही याचा अत्यंत निकटचा संबंध आहे: मूलतः, बंद टॅब्लो
म्हणजे उलटी लिहिलेली क्रमवर्ती कलनातील संक्षेपित निष्पत्ती होय.









\section{स्वयंसिद्धकीय निष्पत्ती}\label{fol:prf:axd:sec}

स्वयंसिद्धकीय निष्पत्त्या या सर्वांत जुन्या आणि सोप्या तार्किक
निष्पत्ती-पद्धती आहेत. त्यांतील निष्पत्त्या म्हणजे केवळ
वाक्यांचे अनुक्रम असतात. वाक्यांचा अनुक्रम तेव्हा योग्य
निष्पत्ती मानला जातो, जेव्हा त्यातील प्रत्येक वाक्य~$\varphi$ पुढील
अटींपैकी एक अट पूर्ण करते:
\begin{enumerate}
\item $\varphi$ हे स्वयंसिद्धक असते, किंवा
\item $\varphi$ हे दिलेल्या वाक्यांच्या संच~$\Gamma$ चा घटक असते, किंवा
\item $\varphi$ ला निगमन नियमामुळे समर्थन मिळते.
\end{enumerate}
$\varphi$ हे स्वयंसिद्धक असण्यासाठी त्याचे रूप काही ठरावीक वाक्य
रूपबंधांपैकी एक असावे लागते. प्रथम-क्रम तर्कशास्त्रासाठी समाधानकारक
(निर्दोष आणि संपूर्ण) निष्पत्ती-पद्धत देणारे स्वयंसिद्धक-रूपबंधांचे
अनेक संच आहेत. त्यांपैकी काही संच, ते ज्या संयोजकांना नियंत्रित करतात
त्यांनुसार मांडलेले असतात; उदाहरणार्थ, पुढील रूपबंध
\[
\varphi \lif (\psi \lif \varphi) \qquad \psi \lif (\psi \lor \chi) \qquad (\psi \land \chi) \lif \psi
\]
हे $\lif$, $\lor$ आणि~$\land$ यांना नियंत्रित करणारे सर्वसाधारण स्वयंसिद्धक
आहेत. काही स्वयंसिद्धकीय पद्धतींमध्ये स्वयंसिद्धकांची संख्या कमीत कमी
ठेवण्याचे उद्दिष्ट असते. कोणते संयोजक आदिम म्हणून स्वीकारले आहेत यानुसार,
एकाच स्वयंसिद्धकाची पद्धत मिळणेही शक्य असते.

निगमन नियम म्हणजे असे सोपाधिक विधान की जे निष्पत्तीमधील
वाक्य समर्थित असण्यासाठी पुरेशी अट देते. विध्यात्मक अनुमान हा असा
एक अतिशय प्रचलित नियम आहे: $\varphi$ आणि $\varphi \lif \psi$ यांना आधीच समर्थन मिळाले
असेल, तर $\psi$ लाही समर्थन मिळते, असे तो सांगतो. म्हणजे, निष्पत्तीमध्ये वाक्य~$\psi$ असलेल्या ओळीला समर्थन मिळते, मात्र $\varphi$ आणि
$\varphi \lif \psi$ ही दोन्ही सूत्रे (काही वाक्य~$\varphi$ साठी) त्या
निष्पत्तीमध्ये~$\psi$ च्या आधी आलेली असली पाहिजेत.

स्वयंसिद्धकीय निष्पत्त्यांवर आधारित $\Proves$ संबंधाची व्याख्या
पुढीलप्रमाणे केली जाते: $\Gamma \Proves \varphi$ तेव्हा आणि केवळ तेव्हाच, जेव्हा
अशी निष्पत्ती असते की तिच्या शेवटच्या सूत्राच्या जागी
वाक्य~$\varphi$ असते (आणि वरच्या~(2) मुळे समर्थन मिळालेल्या त्या
निष्पत्तीमधील वाक्यांचा संच म्हणजे $\Gamma$ असे घेतले जाते).
~$\varphi$ ची अशी निष्पत्ती असेल की तिच्यासाठी~$\Gamma$ रिक्त असेल, म्हणजेच
त्या निष्पत्तीमधील प्रत्येक वाक्याला (1) किंवा~(3) मुळे समर्थन
मिळाले असेल, तर~$\varphi$ हे प्रमेय असते. उदाहरणार्थ, $\Proves
\varphi  \lif (\psi \lif (\psi \lor \varphi))$ हे दाखवणारी निष्पत्ती पुढे दिली आहे:
\begin{derivation}
  1. & $\psi \lif (\psi \lor \varphi)$ \\
  2. & $(\psi \lif (\psi \lor \varphi)) \lif (\varphi  \lif (\psi \lif (\psi \lor \varphi)))$\\
  3. & $\varphi  \lif (\psi \lif (\psi \lor \varphi))$
\end{derivation}
ओळ~1 वरील वाक्य हे $\varphi \lif (\varphi \lor \psi)$ या स्वयंसिद्धकाच्या
रूपाचे आहे ($\varphi$ आणि $\psi$ यांच्या भूमिका उलट आहेत). ओळ~2 वरील वाक्य
$\varphi \lif (\psi \lif \varphi)$ या स्वयंसिद्धकाच्या रूपाचे आहे. त्यामुळे दोन्ही
ओळींना समर्थन मिळते. ओळ~3 ला विध्यात्मक अनुमानामुळे समर्थन मिळते: त्याचा
संक्षेप $\theta$ असा केल्यास, ओळ~2 चे रूप $\chi \lif \theta$ असे आहे, आणि त्यातील
$\chi$ म्हणजे $\psi \lif (\psi \lor \varphi)$, म्हणजेच ओळ~1 होय.

$\Gamma \Proves \lfalse$ असेल, तर संच $\Gamma$ विसंगत असतो. संपूर्ण
स्वयंसिद्धकीय पद्धत कोणत्याही~$\varphi$ साठी $\lfalse \lif \varphi$ हेही सिद्ध करेल;
त्यामुळे $\Gamma$ विसंगत असेल, तर कोणत्याही~$\varphi$ साठी $\Gamma \Proves \varphi$.

तर्कशास्त्रासाठी स्वयंसिद्धकीय निष्पत्त्यांच्या पद्धती सर्वप्रथम गोटलोप
फ्रेग यांनी त्यांच्या १८७९ मधील \emph{Begriffsschrift} या ग्रंथात दिल्या;
म्हणून हा ग्रंथ आधुनिक तर्कशास्त्रातील पहिला ग्रंथ मानला जातो. ॲल्फ्रेड
नॉर्थ व्हाइटहेड आणि बर्ट्रंड रसेल यांच्या \emph{Principia Mathematica} मध्ये,
तसेच १९२० च्या दशकात डाव्हीट हिल्बर्ट आणि त्यांच्या विद्यार्थ्यांकडून, या
पद्धतींना परिपूर्ण रूप मिळाले. म्हणून त्यांना अनेकदा ``फ्रेग पद्धती'' किंवा
``हिल्बर्ट पद्धती'' म्हणतात. एखाद्या तर्कशास्त्रासाठी स्वयंसिद्धकीय पद्धत
शोधणे बहुधा सोपे असल्यामुळे या पद्धती अतिशय बहुपयोगी आहेत. निष्पत्त्यांची रचना अत्यंत सोपी असून त्यांत फक्त एक किंवा दोन निगमन नियम असतात;
त्यामुळे \emph{त्यांच्याविषयीच्या} गोष्टी सिद्ध करणेही तुलनेने सोपे असते. तथापि,
प्रत्यक्ष व्यवहारात त्यांचा वापर फार कठीण असतो; म्हणजेच सिद्धता शोधणे आणि
लिहिणे कठीण असते.



\chapter{क्रमवर्ती कलन}\label{fol:seq::chap}
\begin{editorial}
  या प्रकरणात अभिजात प्रथम-क्रम तर्कशास्त्रासाठी गेंटझेन यांचे मानक
  क्रमवर्ती कलन LK मांडले आहे. यात आणखी उदाहरणे आणि सराव समाविष्ट करणे
  उपयुक्त ठरेल. सिद्धता-पद्धती म्हणून क्रमवर्ती कलनाशी संबंधित मजकूर
  समाविष्ट करण्यासाठी किंवा वगळण्यासाठी ``prfLK'' हा टॅग वापरा.
\end{editorial}






\section{नियम आणि निष्पत्ती}\label{fol:seq:rul:sec}

पुढील मजकुरात, $\Gamma, \Delta, \Pi, \Lambda$ या वाक्यांच्या सांत
क्रमिका दर्शवतील, असे मानू.

\begin{defn}[क्रमवर्ती]
\emph{क्रमवर्ती} म्हणजे पुढील रूपाची अभिव्यक्ती होय:
\[
\Gamma \Sequent \Delta
\]
येथे $\Gamma$ आणि $\Delta$ या $\Lang L$ भाषेतील वाक्यांच्या
सांत (संभवतः रिक्त) क्रमिका आहेत. $\Gamma$ ला \emph{पूर्वांग}, तर
$\Delta$ ला \emph{उत्तरांग} म्हणतात.
\end{defn}

\begin{explain}
क्रमवर्तीमागील अंतःप्रेरित कल्पना अशी आहे: पूर्वांगातील सर्व वाक्ये
सत्य असतील, तर उत्तरांगातील वाक्यांपैकी किमान एक सत्य असते. म्हणजे,
$\Gamma = \tuple{\varphi_1, \dots, \varphi_m}$ आणि
$\Delta = \tuple{\psi_1, \dots, \psi_n}$ असतील, तर $\Gamma \Sequent \Delta$
तेव्हा आणि केवळ तेव्हाच सत्य असते, जेव्हा
\[
(\varphi_1 \land \cdots \land \varphi_m) \lif (\psi_1 \lor \cdots \lor
\psi_n)
\]
सत्य असते. याची दोन विशेष प्रकरणे आहेत: $\Gamma$ रिक्त असण्याचे आणि
$\Delta$~रिक्त असण्याचे. $\Gamma$ रिक्त असेल, म्हणजेच $m = 0$ असेल, तर
$\quad \Sequent \Delta$ तेव्हा आणि केवळ तेव्हाच सत्य असते, जेव्हा
$\psi_1 \lor \dots \lor \psi_n$ सत्य असते. $\Delta$ रिक्त असेल, म्हणजेच
$n = 0$ असेल, तर $\Gamma \Sequent \quad$ तेव्हा आणि केवळ तेव्हाच सत्य
असते, जेव्हा $\lnot(\varphi_1 \land \dots \land \varphi_m)$ सत्य असते. एखाद्या
क्रमवर्तीला तेव्हा आणि केवळ तेव्हाच वैध म्हणतो, जेव्हा तिच्याशी संबंधित
वाक्य वैध असते.
\end{explain}

$\Gamma$ ही वाक्यांची क्रमिका असेल, तर $\Gamma$ च्या उजव्या टोकाला
$\varphi$ जोडून मिळणाऱ्या क्रमिकेसाठी आपण $\Gamma, \varphi$ लिहितो (आणि $\Gamma$
च्या डाव्या टोकाला $\varphi$ जोडून मिळणाऱ्या क्रमिकेसाठी $\varphi, \Gamma$
लिहितो). $\Delta$ हीदेखील वाक्यांची क्रमिका असेल, तर $\Gamma,
\Delta$ ही त्या दोन क्रमिकांची जोडणी असते.

\begin{defn}[आरंभीची क्रमवर्ती]
\emph{आरंभीची क्रमवर्ती} म्हणजे, भाषेतील कोणत्याही वाक्य~$\varphi$ साठी,
पुढीलपैकी एका रूपाची क्रमवर्ती होय:
  \begin{enumerate}
    \item $\varphi \Sequent \varphi$
    \item $\quad \Sequent \ltrue$
    \item $\lfalse \Sequent \quad$
  \end{enumerate}.
\end{defn}

क्रमवर्ती कलनातील निष्पत्त्या हे क्रमवर्तींचे विशिष्ट वृक्ष असतात.
त्यांतील सर्वांत वरच्या क्रमवर्ती आरंभीच्या क्रमवर्ती असतात; आणि एखादी
क्रमवर्ती इतर एक किंवा दोन क्रमवर्तींच्या खाली असेल, तर ती निगमनाच्या
एखाद्या नियमानुसार त्यांच्यापासून योग्य रीतीने निष्पन्न झाली पाहिजे.
$\Log{LK}$ चे नियम दोन मुख्य प्रकारांत विभागले आहेत: \emph{तार्किक} नियम
आणि \emph{संरचनात्मक} नियम. खालच्या क्रमवर्तीतील $\varphi$ आणि/किंवा $\psi$
असलेल्या वाक्याच्या मुख्य संकारक वरून तार्किक नियमांना नावे दिली
जातात. प्रत्येक नियमाच्या दोन आवृत्त्या असतात: संकारक असलेले
वाक्य डावीकडे असणारी क्रमवर्ती निष्पन्न करण्यासाठी एक, आणि ते
वाक्य उजवीकडे असणारी क्रमवर्ती निष्पन्न करण्यासाठी दुसरी.









\section{विधानीय नियम}\label{fol:seq:prl:sec}

\subsection{$\lnot$ साठीचे नियम}

\begin{defish}
\Axiom$ \Gamma \fCenter \Delta, \varphi $
\RightLabel{\LeftR{\lnot}}
\UnaryInf$ \lnot \varphi, \Gamma \fCenter \Delta$
\DisplayProof
\hfill
\Axiom$\varphi, \Gamma \fCenter \Delta$
\RightLabel{\RightR{\lnot}}
\UnaryInf$ \Gamma \fCenter \Delta, \lnot \varphi $
\DisplayProof
\end{defish}

\subsection{$\land$ साठीचे नियम}

\begin{defish}\noindent
\begin{tabular}{l}
\Axiom$ \varphi, \Gamma \fCenter \Delta$
\RightLabel{\LeftR{\land}}
\UnaryInf$ \varphi \land \psi, \Gamma \fCenter \Delta$
\DisplayProof
\\[3ex]
\Axiom$\psi, \Gamma \fCenter \Delta$
\RightLabel{\LeftR{\land}}
\UnaryInf$\varphi \land \psi, \Gamma \fCenter \Delta$
\DisplayProof
\end{tabular}
\hfill
\Axiom$\Gamma \fCenter \Delta, \varphi$
\Axiom$ \Gamma \fCenter \Delta, \psi$
\RightLabel{\RightR{\land}}
\BinaryInf$ \Gamma \fCenter \Delta, \varphi \land \psi $
\DisplayProof
\end{defish}

\subsection{$\lor$ साठीचे नियम}

\begin{defish}
\Axiom$\varphi, \Gamma \fCenter \Delta$
\Axiom$\psi, \Gamma \fCenter \Delta$
\RightLabel{\LeftR{\lor}}
\BinaryInf$\varphi \lor \psi, \Gamma \fCenter \Delta$
\DisplayProof
\hfill
\begin{tabular}{r}
\Axiom$\Gamma \fCenter \Delta, \varphi$
\RightLabel{\RightR{\lor}}
\UnaryInf$ \Gamma \fCenter \Delta, \varphi \lor \psi$
\DisplayProof
\\[3ex]
\Axiom$ \Gamma \fCenter \Delta, \psi$
\RightLabel{\RightR{\lor}}
\UnaryInf$ \Gamma \fCenter \Delta, \varphi \lor \psi$
\DisplayProof
\end{tabular}
\end{defish}

\subsection{$\lif$ साठीचे नियम}

\begin{defish}
\Axiom$ \Gamma \fCenter \Delta, \varphi$
\Axiom$ \psi, \Pi \fCenter \Lambda$
\RightLabel{\LeftR{\lif}}
\BinaryInf$ \varphi \lif \psi, \Gamma, \Pi \fCenter \Delta, \Lambda$
\DisplayProof
\hfill
\Axiom$ \varphi, \Gamma \fCenter \Delta, \psi$
\RightLabel{\RightR{\lif}}
\UnaryInf$ \Gamma \fCenter \Delta, \varphi \lif \psi $
\DisplayProof
\end{defish}









\section{संख्यापकांचे नियम}\label{fol:seq:qrl:sec}

\subsection{$\lforall$ साठीचे नियम}

\begin{defish}
\Axiom$ \varphi(t), \Gamma \fCenter \Delta$
\RightLabel{\LeftR{\lforall}}
\UnaryInf$ \lforall[x][\varphi(x)],\Gamma \fCenter \Delta$
\DisplayProof
\hfill
\Axiom$ \Gamma \fCenter \Delta, \varphi(a) $
\RightLabel{\RightR{\lforall}}
\UnaryInf$ \Gamma \fCenter \Delta, \lforall[x][\varphi(x)]$
\DisplayProof
\end{defish}

\LeftR{\lforall} मध्ये, $t$ हे बंद पद आहे (म्हणजे, चरे नसलेले पद).
\RightR{\lforall} मध्ये, $a$~हा असा स्थिरांक आहे की तो
\RightR{\lforall} नियमाच्या खालच्या क्रमवर्तीमध्ये कोठेही येता कामा नये.
\RightR{\forall} अनुमानाचा $a$~हा \emph{आयगेन चर} आहे, असे आपण
म्हणतो.\footnote{वरील नियमात $a$ हा स्थिरांक असला तरी आपण
``आयगेन चर'' ही संज्ञा वापरतो. यामागे ऐतिहासिक कारणे आहेत.}

\enlargethispage{1pt}
\subsection{$\lexists$ साठीचे नियम}

\begin{defish}
\Axiom$ \varphi(a), \Gamma \fCenter \Delta $
\RightLabel{\LeftR{\lexists}}
\UnaryInf$ \lexists[x][\varphi(x)], \Gamma \fCenter \Delta$
\DisplayProof
\hfill
\Axiom$ \Gamma \fCenter \Delta, \varphi(t) $
\RightLabel{\RightR{\lexists}}
\UnaryInf$ \Gamma \fCenter \Delta, \lexists[x][\varphi(x)]$
\DisplayProof
\end{defish}

पुन्हा, $t$~हे बंद पद आहे आणि $a$~हा असा स्थिरांक आहे की तो
\LeftR{\lexists} नियमाच्या खालच्या क्रमवर्तीमध्ये येत नाही.
\LeftR{\lexists} अनुमानाचा $a$~हा \emph{आयगेन चर} आहे, असे आपण म्हणतो.

\RightR{\lforall} किंवा \LeftR{\lexists} अनुमानाच्या खालच्या
क्रमवर्तीमध्ये आयगेन चर येत नसावा, या अटीला \emph{आयगेन चराची अट}
म्हणतात.

\begin{explain}
ही प्रथा आठवा: $\varphi$ हे असे सूत्र असेल की त्यात
चल~$x$ मुक्त असेल, तर हे आपण~$\varphi(x)$ असे लिहून दर्शवतो.
त्याच संदर्भात, मग $\varphi(t)$ हा~$\Subst{\varphi}{t}{x}$ चा संक्षेप असतो.
त्यामुळे \RightR{\lexists} नियम पुढीलप्रमाणेही लिहिता येईल:
\begin{prooftree}
  \Axiom$\Gamma \fCenter \Delta, \Subst{\varphi}{t}{x}$
  \RightLabel{\RightR{\lexists}}
  \UnaryInf$\Gamma \fCenter \Delta, \lexists[x][\varphi]$
\end{prooftree}
लक्षात घ्या की~$\varphi$ मध्ये $t$ आधीच येऊ शकते; उदा., $\varphi$ हे
~$\Atom{\Obj P}{t,x}$ असू शकते. म्हणून,~$\Gamma \Sequent \Delta,
\Atom{\Obj P}{t,t}$ पासून $\Gamma \Sequent \Delta,
\lexists[x][\Atom{\Obj P}{t,x}]$ निष्पन्न करणे हा
\RightR{\lexists} चा योग्य उपयोग आहे---आधारविधानात~$t$ ज्या ठिकाणी
आले आहे, त्यांपैकी एक किंवा अधिक ठिकाणी (म्हणजे, त्या पदाच्या एक किंवा अधिक
आस्थितींच्या जागी) बद्ध चल~$x$ घालता येतो; सर्वच आस्थिती बदलणे आवश्यक नाही. परंतु
\RightR{\lforall} आणि~\LeftR{\lexists} मधील आयगेन चराच्या अटींनुसार
स्थिरांक~$a$ हा~$\varphi$ मध्ये येता कामा नये. म्हणून,
$\RightR{\lforall}$ वापरून $\Gamma \Sequent \Delta, \Atom{\Obj P}{a,a}$
पासून $\Gamma \Sequent \Delta, \lforall[x][\Atom{\Obj P}{a,x}]$
योग्य रीतीने निष्पन्न करता येत नाही.
\end{explain}

\begin{explain}
\RightR{\lexists} आणि \LeftR{\lforall} मध्ये पद~$t$ बंद असण्याव्यतिरिक्त
त्यावर आणखी कोणतेही निर्बंध नाहीत. याउलट, \LeftR{\lexists} आणि \RightR{\lforall}
नियमांमध्ये आयगेन चराच्या अटीनुसार वरच्या क्रमवर्तीमध्ये~$\varphi(a)$ च्या
बाहेर स्थिरांक~$a$ कोठेही येता कामा नये. ही पद्धत निर्दोष आहे,
म्हणजे ती केवळ वैध क्रमवर्तीच निष्पन्न करते, याची खात्री करण्यासाठी ही
अट आवश्यक आहे. ही अट नसती, तर पुढील निष्पत्त्या अनुज्ञात असत्या:
\begin{prooftree}
  \Axiom$\varphi(a) \fCenter \varphi(a)$
  \RightLabel{*\LeftR{\lexists}}
  \UnaryInf$\lexists[x][\varphi(x)] \fCenter \varphi(a)$
  \RightLabel{\RightR{\lforall}}
  \UnaryInf$\lexists[x][\varphi(x)] \fCenter \lforall[x][\varphi(x)]$
  \DisplayProof\bottomAlignProof
  \qquad
  \Axiom$\varphi(a) \fCenter \varphi(a)$
  \RightLabel{*\RightR{\lforall}}
  \UnaryInf$\varphi(a) \fCenter \lforall[x][\varphi(x)]$
  \RightLabel{\LeftR{\lexists}}
  \UnaryInf$\lexists[x][\varphi(x)] \fCenter \lforall[x][\varphi(x)]$
\end{prooftree}
परंतु $\lexists[x][\varphi(x)] \Sequent \lforall[x][\varphi(x)]$ ही क्रमवर्ती
वैध नाही.
\end{explain}









\section{संरचनात्मक नियम}\label{fol:seq:srl:sec}

क्रमवर्तीच्या डाव्या आणि उजव्या बाजूंवरील वाक्यांची मांडणी बदलता येईल,
असे आणखी काही नियम आपल्याला आवश्यक आहेत. तार्किक नियम ज्या वाक्ये वर
लागू होतो, ती आधारविधानात अगदी डाव्या किंवा अगदी उजव्या टोकाला असणे आवश्यक
असते. त्यामुळे वाक्ये योग्य स्थानी हलवता यावीत म्हणून आपल्याला
``अदलाबदल'' नियमाची गरज आहे. तसेच कधी कधी दोन एकसारखी वाक्ये एकात
एकत्र करता येणे आणि दोन्हींपैकी कोणत्याही बाजूला वाक्य जोडता येणे
महत्त्वाचे असते.

\subsection{दुर्बलीकरण}

\begin{defish}
\Axiom$ \Gamma \fCenter \Delta $
\RightLabel{\LeftR{\Weakening}}
\UnaryInf$ \varphi, \Gamma \fCenter \Delta$
\DisplayProof
\hfill
\Axiom$ \Gamma \fCenter \Delta$
\RightLabel{\RightR{\Weakening}}
\UnaryInf$ \Gamma \fCenter \Delta, \varphi$
\DisplayProof
\end{defish}

\subsection{आकुंचन}

\begin{defish}
\Axiom$ \varphi, \varphi, \Gamma \fCenter \Delta $
\RightLabel{\LeftR{\Contraction}}
\UnaryInf$ \varphi, \Gamma \fCenter \Delta$
\DisplayProof
\hfill
\Axiom$ \Gamma \fCenter \Delta, \varphi, \varphi$
\RightLabel{\RightR{\Contraction}}
\UnaryInf$ \Gamma \fCenter \Delta, \varphi$
\DisplayProof
\end{defish}

\subsection{अदलाबदल}

\begin{defish}
\Axiom$ \Gamma, \varphi, \psi, \Pi \fCenter \Delta $
\RightLabel{\LeftR{\Exchange}}
\UnaryInf$ \Gamma, \psi, \varphi, \Pi \fCenter \Delta$
\DisplayProof
\hfill
\Axiom$ \Gamma \fCenter \Delta, \varphi, \psi, \Lambda$
\RightLabel{\RightR{\Exchange}}
\UnaryInf$ \Gamma \fCenter \Delta, \psi, \varphi, \Lambda$
\DisplayProof
\end{defish}

दुर्बलीकरण, आकुंचन आणि अदलाबदल या अनुमानांची मालिका अनेकदा दुहेरी
अनुमानरेषांनी दर्शवली जाते.

पुढील ``कट'' नावाचा नियम काटेकोर अर्थाने आवश्यक नाही; मात्र त्यामुळे
निष्पत्त्या पुन्हा वापरणे आणि एकत्र करणे पुष्कळ सोपे होते.
\begin{defish}
\[
\Axiom$ \Gamma \fCenter \Delta, \varphi$
\Axiom$ \varphi, \Pi \fCenter \Lambda $
\RightLabel{\Cut}
\BinaryInf$ \Gamma, \Pi \fCenter \Delta, \Lambda$
\DisplayProof
\]
\end{defish}









\section{निष्पत्ती}\label{fol:seq:der:sec}

\begin{explain}
आरंभीची क्रमवर्ती कशी असते हे आपण सांगितले आहे आणि निगमन नियम दिले आहेत.
क्रमवर्ती कलनातील निष्पत्त्या यांपासून विगमनाधारित रीतीने निर्माण केल्या
जातात: प्रत्येक निष्पत्ती ही एकतर स्वतःच आरंभीची क्रमवर्ती असते, किंवा
एका अनुमानाच्या वर असलेल्या एक किंवा दोन निष्पत्त्या ने बनलेली असते.
\end{explain}

\begin{defn}[$\Log{LK}$ निष्पत्ती]
क्रमवर्ती~$S$ ची \emph{$\Log{LK}$-निष्पत्ती} म्हणजे पुढील अटी पूर्ण
करणारा क्रमवर्तींचा सांत वृक्ष होय:
\begin{enumerate}
\item वृक्षातील सर्वांत वरच्या क्रमवर्ती आरंभीच्या क्रमवर्ती असतात.
\item वृक्षातील सर्वांत खालची क्रमवर्ती~$S$ असते.
\item वृक्षातील $S$ व्यतिरिक्त प्रत्येक क्रमवर्ती ही निगमन नियमाच्या योग्य
  उपयोजनाचे आधारविधान असते आणि त्या उपयोजनाचा निष्कर्ष वृक्षात त्या
  क्रमवर्तीच्या थेट खाली असतो.
\end{enumerate}
यानंतर $S$ ला त्या निष्पत्तीची \emph{अंतिम क्रमवर्ती} म्हणतो आणि
$S$ ही \emph{$\Log{LK}$ मध्ये निष्पन्न करता येण्याजोगे} आहे (किंवा ती
$\Log{LK}$-निष्पन्न करता येण्याजोगे आहे) असे म्हणतो.
\end{defn}

\begin{ex}
प्रत्येक आरंभीची क्रमवर्ती, उदा., $\chi \Sequent \chi$, ही निष्पत्ती असते.
उदाहरणार्थ, $\LeftR{\Weakening}$ नियम लागू करून आपण तिच्यापासून नवी
निष्पत्ती मिळवू शकतो:
\begin{prooftree}
\Axiom$ \Gamma \fCenter \Delta $
\RightLabel{\LeftR{\Weakening}}
\UnaryInf$ \varphi, \Gamma \fCenter \Delta$
\end{prooftree}
मात्र हा नियम सर्वसाधारण स्वरूपाचा आहे: नियमातील $\varphi$ च्या जागी कोणतेही
वाक्य, उदा., $\theta$ सुद्धा, घेता येते. आधारविधान आपल्या आरंभीच्या
क्रमवर्ती $\chi \Sequent \chi$ शी जुळत असेल, तर $\Gamma$ आणि $\Delta$ ही दोन्ही
केवळ~$\chi$ असतात आणि निष्कर्ष $\theta, \chi \Sequent \chi$ असा असेल. म्हणून पुढील
ही निष्पत्ती आहे:
\begin{prooftree}
\Axiom$ \chi \fCenter \chi $
\RightLabel{\LeftR{\Weakening}}
\UnaryInf$ \theta, \chi \fCenter \chi$
\end{prooftree}
आता आपण दुसरा नियम, समजा $\LeftR{\Exchange}$, लागू करू शकतो; त्यामुळे
डावीकडील दोन वाक्यांची अदलाबदल करता येते. म्हणून पुढीलही योग्य
निष्पत्ती आहे:
\begin{prooftree}
\Axiom$ \chi \fCenter \chi $
\RightLabel{\LeftR{\Weakening}}
\UnaryInf$ \theta, \chi \fCenter \chi$
\RightLabel{\LeftR{\Exchange}}
\UnaryInf$ \chi, \theta \fCenter \chi$
\end{prooftree}
नियमाच्या या उपयोजनात—मूळ नियम पुढीलप्रमाणे दिला होता—
\begin{prooftree}
\Axiom$ \Gamma, \varphi, \psi, \Pi \fCenter \Delta $
\RightLabel{\LeftR{\Exchange}}
\UnaryInf$ \Gamma, \psi, \varphi, \Pi \fCenter \Delta,$
\end{prooftree}
$\Gamma$ आणि $\Pi$ दोन्ही रिक्त ठेवले होते, $\Delta$ हे $\chi$ होते, आणि
$\varphi$ व $\psi$ यांच्या भूमिका अनुक्रमे $\theta$ आणि~$\chi$ यांनी बजावल्या. अगदी
याच प्रकारे पुढीलही निष्पत्ती आहे हे आपल्याला दिसते:
\begin{prooftree}
\Axiom$ \theta \fCenter \theta $
\RightLabel{\LeftR{\Weakening}}
\UnaryInf$ \chi, \theta \fCenter \theta$
\end{prooftree}
आता या दोन निष्पत्त्या घेऊन $\RightR{\land}$ च्या साहाय्याने त्यांना
एकत्र करता येते. तो नियम असा होता:
\begin{prooftree}
\Axiom$\Gamma \fCenter \Delta, \varphi$
\Axiom$ \Gamma \fCenter \Delta, \psi$
\RightLabel{\RightR{\land}}
\BinaryInf$ \Gamma \fCenter \Delta, \varphi \land \psi $
\end{prooftree}
आपल्या उदाहरणात आधारविधाने त्यांवर संपणाऱ्या निष्पत्त्यांच्या अंतिम
क्रमवर्तींशी जुळली पाहिजेत. याचा अर्थ $\Gamma$ हे $\chi, \theta$ आहे, $\Delta$
रिक्त आहे, $\varphi$ हे $\chi$ आहे आणि $\psi$ हे $\theta$ आहे. म्हणून अनुमान योग्य
असेल, तर निष्कर्ष $\chi, \theta \Sequent \chi \land \theta$ हा आहे.
\begin{prooftree}
\Axiom$ \chi \fCenter \chi $
\RightLabel{\LeftR{\Weakening}}
\UnaryInf$ \theta, \chi \fCenter \chi$
\RightLabel{\LeftR{\Exchange}}
\UnaryInf$ \chi, \theta \fCenter \chi$
\Axiom$ \theta \fCenter \theta $
\RightLabel{\LeftR{\Weakening}}
\UnaryInf$ \chi, \theta \fCenter \theta$
\RightLabel{\RightR{\land}}
\BinaryInf$ \chi, \theta \fCenter \chi \land \theta $
\end{prooftree}
अर्थात, आधारविधानांचा क्रम उलटाही करता येतो; मग $\varphi$ हे $\theta$ आणि $\psi$
हे~$\chi$ असेल.
\begin{prooftree}
\Axiom$ \theta \fCenter \theta $
\RightLabel{\LeftR{\Weakening}}
\UnaryInf$ \chi, \theta \fCenter \theta$
\Axiom$ \chi \fCenter \chi $
\RightLabel{\LeftR{\Weakening}}
\UnaryInf$ \theta, \chi \fCenter \chi$
\RightLabel{\LeftR{\Exchange}}
\UnaryInf$ \chi, \theta \fCenter \chi$
\RightLabel{\RightR{\land}}
\BinaryInf$ \chi, \theta \fCenter \theta \land \chi $
\end{prooftree}
\end{ex}









\section{निष्पत्ती: उदाहरणे}\label{fol:seq:pro:sec}

\begin{ex}
क्रमवर्ती $\varphi \land \psi \Sequent \varphi$ ची एक $\Log{LK}$-निष्पत्ती द्या.

हवी असलेली अंतिम क्रमवर्ती निष्पत्तीच्या तळाशी लिहून आपण सुरुवात करतो.
\begin{prooftree}
\AxiomC{}
\UnaryInf$\varphi\land \psi \fCenter \varphi$
\end{prooftree}
यानंतर, या स्वरूपाची खालची क्रमवर्ती कोणत्या प्रकारच्या अनुमानामुळे मिळू शकते
हे शोधावे लागते. ते एखाद्या संरचनात्मक नियमाचे उपयोजन असू शकते; पण आधी
तार्किक नियम शोधणे योग्य ठरते. खालच्या क्रमवर्तीत येणारा एकमेव तार्किक
संयोजक $\land$ आहे; म्हणून आपण $\land$ चा नियम शोधत आहोत. $\land$ हे चिन्ह
पूर्वांगात येत असल्यामुळे आपण \LeftR{\land} नियम पाहतो.
\begin{prooftree}
\AxiomC{}
\RightLabel{\LeftR{\land}}
\UnaryInf$\varphi\land \psi \fCenter \varphi$
\end{prooftree}
\LeftR{\land} अनुमानाची वरची क्रमवर्ती काय असू शकते यासाठी दोन पर्याय आहेत:
वरची क्रमवर्ती $\varphi \Sequent \varphi$ किंवा $\psi \Sequent \varphi$ असू शकते. स्पष्टपणे,
$\varphi \Sequent \varphi$ ही आरंभीची क्रमवर्ती आहे (ही आपल्या दृष्टीने चांगली बाब
आहे), तर $\psi \Sequent \varphi$ ही सर्वसाधारणपणे निष्पन्न करता येत नाही. आपण वरची
क्रमवर्ती भरू:
\begin{prooftree}
\Axiom$\varphi \fCenter \varphi$
\RightLabel{\LeftR{\land}}
\UnaryInf$\varphi\land \psi \fCenter \varphi$
\end{prooftree}
आता आपल्याकडे क्रमवर्ती $\varphi \land \psi \Sequent \varphi$ ची योग्य
$\Log{LK}$-निष्पत्ती आहे.
\end{ex}

\begin{ex}
क्रमवर्ती $\lnot \varphi \lor \psi \Sequent \varphi \lif \psi$ ची एक
$\Log{LK}$-निष्पत्ती द्या.

हवी असलेली अंतिम क्रमवर्ती निष्पत्तीच्या तळाशी लिहून सुरुवात करा.
\begin{prooftree}
\AxiomC{}
\UnaryInf$\lnot \varphi \lor \psi \fCenter \varphi \lif \psi$
\end{prooftree}
ही अंतिम क्रमवर्ती देऊ शकेल असा तार्किक नियम शोधण्यासाठी तिच्यातील तार्किक
संयोजक पाहू: $\lnot$, $\lor$ आणि $\lif$. सध्या आपल्याला फक्त $\lor$ आणि
$\lif$ यांचा विचार करायचा आहे, कारण ती अंतिम क्रमवर्तीतील वाक्यांची
मुख्य संकारक आहेत; $\lnot$ दुसऱ्या संयोजकाच्या व्याप्तीच्या आत आहे,
म्हणून त्याचा विचार आपण नंतर करू. शेवटच्या अनुमानासाठी आपल्याकडे
\LeftR{\lor} नियम आणि \RightR{\lif} नियम हे तार्किक नियमांचे पर्याय आहेत.
खरे तर कोणताही नियम निवडता येईल; पण आपण \RightR{\lif} नियम निवडू (दुसरे
काही कारण नसले तरी त्यामुळे दोन शाखांत विभागणे काही काळ पुढे ढकलता येते).
ही खालची क्रमवर्ती देऊ शकणाऱ्या \RightR{\lif} अनुमानांच्या स्वरूपानुसार ते
पुढीलप्रमाणे असले पाहिजे:
\begin{prooftree}
\AxiomC{}
\UnaryInf$ \varphi, \lnot \varphi \lor \psi \fCenter \psi $
\RightLabel{\RightR{\lif}} \UnaryInf$ \lnot \varphi \lor \psi \fCenter \varphi \lif \psi $
\end{prooftree}

$\lnot \varphi \lor \psi$ हे पूर्वांगाच्या बाहेरील टोकाला हलवले, तर आपण
\LeftR{\lor} नियम लागू करू शकतो. रूपबंधानुसार त्याचे पुढीलप्रमाणे दोन वरच्या
क्रमवर्तींत विभाजन झाले पाहिजे:
\begin{prooftree}
\AxiomC{}
\UnaryInf$\lnot \varphi, \varphi \fCenter \psi$
\AxiomC{}
\UnaryInf$\psi, \varphi \fCenter \psi$
\RightLabel{\LeftR{\lor}}
\BinaryInf$ \lnot \varphi \lor \psi, \varphi \fCenter \psi $
\RightLabel{\LeftR{\Exchange}}
\UnaryInf$ \varphi, \lnot \varphi \lor \psi \fCenter \psi $
\RightLabel{\RightR{\lif}}
\UnaryInf$ \lnot \varphi \lor \psi \fCenter \varphi \lif \psi $
\end{prooftree}
आपण वरच्या दिशेने आरंभीच्या क्रमवर्तींपर्यंत वाट काढण्याचा प्रयत्न करीत आहोत
हे लक्षात ठेवा; आपण बरेच जवळ आलो आहोत असे दिसते!{} उजवी शाखा आरंभीच्या
क्रमवर्तीपासून केवळ एक दुर्बलीकरण आणि एक अदलाबदल दूर आहे; ते केल्यावर ती
पूर्ण होते:
\begin{prooftree}
\AxiomC{}
\UnaryInf$\lnot \varphi, \varphi \fCenter \psi$
\Axiom$\psi \fCenter \psi$
\RightLabel{\LeftR{\Weakening}}
\UnaryInf$\varphi, \psi \fCenter \psi$
\RightLabel{\LeftR{\Exchange}}
\UnaryInf$\psi, \varphi \fCenter \psi$
\RightLabel{\LeftR{\lor}}
\BinaryInf$\lnot \varphi \lor \psi, \varphi \fCenter \psi $
\RightLabel{\LeftR{\Exchange}}
\UnaryInf$ \varphi, \lnot \varphi \lor \psi \fCenter \psi $
\RightLabel{\RightR{\lif}}
\UnaryInf$ \lnot \varphi \lor \psi \fCenter \varphi \lif \psi $
\end{prooftree}

आता डावी शाखा पाहू. कोणत्याही वाक्य मधील एकमेव तार्किक संयोजक हे
पूर्वांगातील वाक्ये मधील $\lnot$ चिन्ह आहे; म्हणून आपण
\LeftR{\lnot} नियमाचे एक उपयोजन शोधत आहोत.
\begin{prooftree}
\AxiomC{}
\UnaryInf$ \varphi \fCenter \psi, \varphi$
\RightLabel{\LeftR{\lnot}}
\UnaryInf$\lnot \varphi, \varphi \fCenter \psi$
\Axiom$\psi \fCenter \psi$
\RightLabel{\LeftR{\Weakening}}
\UnaryInf$\varphi, \psi \fCenter \psi$
\RightLabel{\LeftR{\Exchange}}
\UnaryInf$\psi, \varphi \fCenter \psi$
\RightLabel{\LeftR{\lor}}
\BinaryInf$\lnot \varphi \lor \psi, \varphi \fCenter \psi $
\RightLabel{\LeftR{\Exchange}}
\UnaryInf$ \varphi, \lnot \varphi \lor \psi \fCenter \psi $
\RightLabel{\RightR{\lif}}
\UnaryInf$ \lnot \varphi \lor \psi \fCenter \varphi \lif \psi $
\end{prooftree}
उजवी शाखा जशी पूर्ण केली, त्याचप्रमाणे ही डावी शाखासुद्धा पूर्ण होण्यासाठी
आता केवळ एक दुर्बलीकरण आणि एक अदलाबदल करायची आहे.
\begin{prooftree}
\Axiom$\varphi \fCenter \varphi$
\RightLabel{\RightR{\Weakening}}
\UnaryInf$ \varphi \fCenter \varphi, \psi$
\RightLabel{\RightR{\Exchange}}
\UnaryInf$ \varphi \fCenter \psi, \varphi$
\RightLabel{\LeftR{\lnot}}
\UnaryInf$\lnot \varphi, \varphi \fCenter \psi$
\Axiom$\psi \fCenter \psi$
\RightLabel{\LeftR{\Weakening}}
\UnaryInf$\varphi, \psi \fCenter \psi$
\RightLabel{\LeftR{\Exchange}}
\UnaryInf$\psi, \varphi \fCenter \psi$
\RightLabel{\LeftR{\lor}}
\BinaryInf$\lnot \varphi \lor \psi, \varphi \fCenter \psi $
\RightLabel{\LeftR{\Exchange}}
\UnaryInf$ \varphi, \lnot \varphi \lor \psi \fCenter \psi $
\RightLabel{\RightR{\lif}}
\UnaryInf$ \lnot \varphi \lor \psi \fCenter \varphi \lif \psi $
\end{prooftree}
\end{ex}

\begin{ex}
क्रमवर्ती $\lnot \varphi \lor \lnot \psi \Sequent \lnot (\varphi \land \psi)$ ची एक
$\Log{LK}$-निष्पत्ती द्या.

वरील तंत्रे वापरून, हवी असलेली अंतिम क्रमवर्ती तळाशी लिहून आपण सुरुवात करतो.
\begin{prooftree}
\AxiomC{}
\UnaryInf$ \lnot \varphi \lor \lnot \psi \fCenter \lnot (\varphi \land \psi) $
\end{prooftree}
अंतिम क्रमवर्तीतील वाक्यांची उपलब्ध मुख्य संक्रियके $\lor$ चिन्ह आणि
$\lnot$ चिन्ह ही आहेत. येथे \LeftR{\lor} किंवा \RightR{\lnot} यांपैकी कोणताही
नियम लागू केला तरी चालेल; पण दोन शाखांत विभागणे क्षणभर टाळता यावे म्हणून आपण
\RightR{\lnot} नियमापासून सुरुवात करू:
\begin{prooftree}
\AxiomC{}
\UnaryInf$\varphi \land \psi, \lnot \varphi \lor \lnot \psi \fCenter $
\RightLabel{\RightR{\lnot}}
\UnaryInf$\lnot \varphi \lor \lnot \psi \fCenter \lnot (\varphi \land \psi)$
\end{prooftree}
आता \LeftR{\land} नियम पाहायचा की \LeftR{\lor} नियम, असा पर्याय आहे.
\LeftR{\land} नियम लागू केल्यावर काय होते ते पाहू: $\varphi, \lnot \varphi \lor \lnot \psi
\Sequent \quad$ या क्रमवर्तीपासून किंवा $\psi, \lnot \varphi \lor \lnot \psi \Sequent
\quad$ या क्रमवर्तीपासून सुरुवात करण्याचा पर्याय आपल्याकडे आहे. ही
निष्पत्ती $\varphi$ आणि $\psi$ यांच्या दृष्टीने सममित असल्यामुळे आपण पहिला पर्याय
निवडू:
\begin{prooftree}
\AxiomC{}
\UnaryInf$\varphi, \lnot \varphi \lor \lnot \psi \fCenter $
\RightLabel{\LeftR{\land}}
\UnaryInf$\varphi \land \psi, \lnot \varphi \lor \lnot \psi \fCenter $
\RightLabel{\RightR{\lnot}}
\UnaryInf$\lnot \varphi \lor \lnot \psi \fCenter \lnot (\varphi \land \psi)$
\end{prooftree}
निष्पत्ती पुढे भरत राहिल्यावर आपल्यासमोर एक अडचण येते:
\begin{prooftree}
\Axiom$\varphi \fCenter \varphi$
\RightLabel{\LeftR{\lnot}}
\UnaryInf$ \lnot \varphi, \varphi \fCenter$
\AxiomC{}
\RightLabel{?}
\UnaryInf$\varphi \fCenter \psi$
\RightLabel{\LeftR{\lnot}}
\UnaryInf$ \lnot \psi, \varphi \fCenter$
\RightLabel{\LeftR{\lor}}
\BinaryInf$\lnot \varphi \lor \lnot \psi, \varphi \fCenter $
\RightLabel{\LeftR{\Exchange}}
\UnaryInf$\varphi, \lnot \varphi \lor \lnot \psi \fCenter $
\RightLabel{\LeftR{\land}}
\UnaryInf$\varphi \land \psi, \lnot \varphi \lor \lnot \psi \fCenter $
\RightLabel{\RightR{\lnot}}
\UnaryInf$\lnot \varphi \lor \lnot \psi \fCenter \lnot (\varphi \land \psi)$
\end{prooftree}
उजव्या शाखेच्या वरच्या टोकाचे आणखी लघुकरण करता येत नाही आणि संरचनात्मक
अनुमानांच्या साहाय्याने तिला आरंभीची क्रमवर्ती बनवताही येत नाही; म्हणून हा
योग्य मार्ग नाही. त्यामुळे वर \LeftR{\land} नियम लागू करणे ही स्पष्टपणे चूक
होती. मागील अवस्थेकडे परत जाऊन त्याऐवजी \LeftR{\lor} नियम लागू केल्यावर
आपल्याला पुढील निष्पत्ती मिळते:
\begin{prooftree}
\AxiomC{}
\UnaryInf$\lnot \varphi, \varphi \land \psi \fCenter $

\AxiomC{}
\UnaryInf$\lnot \psi, \varphi \land \psi \fCenter $

\RightLabel{\LeftR{\lor}}
\BinaryInf$\lnot \varphi \lor \lnot \psi, \varphi \land \psi \fCenter $
\RightLabel{\LeftR{\Exchange}}
\UnaryInf$\varphi \land \psi, \lnot \varphi \lor \lnot \psi \fCenter $
\RightLabel{\RightR{\lnot}}
\UnaryInf$\lnot \varphi \lor \lnot \psi \fCenter \lnot (\varphi \land \psi)$
\end{prooftree}
पूर्वी केल्याप्रमाणे प्रत्येक शाखा पूर्ण केल्यावर आपल्याला पुढील निष्पत्ती
मिळते:
\begin{prooftree}
\Axiom$ \varphi \fCenter\varphi$
\RightLabel{\LeftR{\land}}
\UnaryInf$\varphi \land \psi \fCenter \varphi$
\RightLabel{\LeftR{\lnot}}
\UnaryInf$\lnot \varphi, \varphi \land \psi \fCenter $

\Axiom$ \psi \fCenter \psi$
\RightLabel{\LeftR{\land}}
\UnaryInf$\varphi \land \psi \fCenter \psi$
\RightLabel{\LeftR{\lnot}}
\UnaryInf$\lnot \psi, \varphi \land \psi \fCenter $

\RightLabel{\LeftR{\lor}}
\BinaryInf$\lnot \varphi \lor \lnot \psi, \varphi \land \psi \fCenter $
\RightLabel{\LeftR{\Exchange}}
\UnaryInf$\varphi \land \psi, \lnot \varphi \lor \lnot \psi \fCenter $
\RightLabel{\RightR{\lnot}}
\UnaryInf$\lnot \varphi \lor \lnot \psi \fCenter \lnot (\varphi \land \psi)$
\end{prooftree}
(या पायऱ्यांत $\lnot$ नियमांपेक्षा खाली $\land$ नियम लागू केले असते, तरीही
आपल्याला योग्य निष्पत्ती मिळाली असती).
\end{ex}

\begin{ex}
आतापर्यंत आपण आकुंचन नियम वापरलेला नाही; परंतु कधी कधी तो आवश्यक असतो. असे
घडणारे एक उदाहरण पाहू. आपल्याला $\quad \Sequent \varphi \lor \lnot \varphi$ सिद्ध
करायचे आहे असे समजा. $\RightR{\lor}$ उलट्या दिशेने लागू केल्यास पुढील दोन
निष्पत्त्या पैकी एक मिळेल:
\begin{prooftree}
\AxiomC{}
\UnaryInf$ \fCenter \varphi$
\RightLabel{\RightR{\lor}}
\UnaryInf$ \fCenter \varphi \lor \lnot \varphi$
\DisplayProof\qquad\bottomAlignProof
\AxiomC{}
\UnaryInf$\varphi \fCenter $
\RightLabel{\RightR{\lnot}}
\UnaryInf$ \fCenter \lnot \varphi$
\RightLabel{\RightR{\lor}}
\UnaryInf$ \fCenter \varphi \lor \lnot \varphi$
\end{prooftree}
यांपैकी कोणतीही अर्थातच आरंभीच्या क्रमवर्तीवर संपत नाही. यासाठीची युक्ती अशी
की आकुंचन नियम आपल्याला वाक्याच्या दोन प्रती एकात एकत्र करू देतो—आणि
सिद्धता शोधताना, म्हणजे तळापासून वर जाताना, आपण आधारविधानात $\varphi \lor \lnot
\varphi$ ची एक प्रत राखून ठेवू शकतो, उदा.,
\begin{prooftree}
\AxiomC{}
\UnaryInf$ \fCenter \varphi \lor \lnot \varphi, \varphi$
\RightLabel{\RightR{\lor}}
\UnaryInf$ \fCenter \varphi \lor \lnot \varphi, \varphi \lor \lnot \varphi$
\RightLabel{\RightR{\Contraction}}
\UnaryInf$ \fCenter \varphi \lor \lnot \varphi$
\end{prooftree}
आता आपण $\RightR{\lor}$ दुसऱ्यांदा लागू करून~$\lnot \varphi$ सुद्धा मिळवू शकतो;
त्यातून एक पूर्ण निष्पत्ती मिळते.
\begin{prooftree}
\Axiom$\varphi \fCenter \varphi$
\RightLabel{\RightR{\lnot}}
\UnaryInf$\fCenter \varphi, \lnot \varphi$
\RightLabel{\RightR{\lor}}
\UnaryInf$\fCenter \varphi, \varphi \lor \lnot \varphi$
\RightLabel{\RightR{\Exchange}}
\UnaryInf$ \fCenter \varphi \lor \lnot \varphi, \varphi$
\RightLabel{\RightR{\lor}}
\UnaryInf$ \fCenter \varphi \lor \lnot \varphi, \varphi \lor \lnot \varphi$
\RightLabel{\RightR{\Contraction}}
\UnaryInf$ \fCenter \varphi \lor \lnot \varphi$
\end{prooftree}
\end{ex}

\begin{prob}
पुढील क्रमवर्तींसाठी निष्पत्त्या तयार करा:
\begin{enumerate}
\item $\varphi \land (\psi \land \chi) \Sequent (\varphi \land \psi) \land \chi$.
\item $\varphi \lor (\psi \lor \chi) \Sequent (\varphi \lor \psi) \lor \chi$.
\item $\varphi \lif (\psi \lif \chi) \Sequent \psi \lif (\varphi \lif \chi)$.
\item $\varphi \Sequent \lnot\lnot \varphi$.
\end{enumerate}
\end{prob}

\begin{prob}
पुढील क्रमवर्तींसाठी निष्पत्त्या तयार करा:
\begin{enumerate}
\item $(\varphi \lor \psi) \lif \chi \Sequent \varphi \lif \chi$.
\item $(\varphi \lif \chi) \land (\psi \lif \chi) \Sequent (\varphi \lor \psi) \lif \chi$.
\item $\Sequent \lnot(\varphi \land \lnot \varphi)$.
\item $\psi \lif \varphi \Sequent \lnot \varphi \lif \lnot \psi$.
\item $\Sequent (\varphi \lif \lnot \varphi) \lif \lnot \varphi$.
\item $\Sequent \lnot(\varphi \lif \psi) \lif \lnot \psi$.
\item $\varphi \lif \chi \Sequent \lnot (\varphi \land \lnot \chi)$.
\item $\varphi \land \lnot \chi \Sequent \lnot (\varphi \lif \chi)$.
\item $\varphi \lor \psi, \lnot \psi \Sequent \varphi$.
\item $\lnot \varphi \lor \lnot \psi \Sequent \lnot(\varphi \land \psi)$.
\item $\Sequent (\lnot \varphi \land \lnot \psi) \lif\lnot(\varphi \lor \psi)$.
\item $\Sequent \lnot(\varphi \lor \psi) \lif (\lnot \varphi \land \lnot \psi)$.
\end{enumerate}
\end{prob}

\begin{prob}
पुढील क्रमवर्तींसाठी निष्पत्त्या तयार करा:
\begin{enumerate}
\item $\lnot(\varphi \lif \psi) \Sequent \varphi$.
\item $\lnot(\varphi \land \psi) \Sequent \lnot \varphi \lor \lnot \psi$.
\item $\varphi \lif \psi \Sequent \lnot \varphi \lor \psi$.
\item $\Sequent \lnot \lnot \varphi \lif \varphi$.
\item $\varphi \lif \psi, \lnot \varphi \lif \psi \Sequent \psi$.
\item $(\varphi \land \psi) \lif \chi \Sequent (\varphi \lif \chi) \lor (\psi \lif \chi)$.
\item $(\varphi \lif \psi) \lif \varphi \Sequent \varphi$.
\item $\Sequent (\varphi \lif \psi) \lor (\psi \lif \chi)$.
\end{enumerate}
(या सर्वांसाठी $\RightR{\Contraction}$~नियम आवश्यक आहे.)
\end{prob}








\section{संख्यापकांसह निष्पत्ती}\label{fol:seq:prq:sec}

\begin{ex}
क्रमवर्ती $\lexists[x][\lnot \varphi(x)]
\Sequent \lnot \lforall[x][\varphi(x)]$ ची एक $\Log{LK}$-निष्पत्ती द्या.

संख्यापक हाताळताना आयगेन चराच्या अटीचा भंग होणार नाही याची खात्री करावी
लागते; काही वेळा त्यासाठी ठरावीक अनुमाने कोणत्या क्रमाने करायची, यात फेरबदल
करावा लागतो. सामान्यतः आयगेन चराची अट लागू असलेले नियम आधी हाताळण्याचा
प्रयत्न केल्यास मदत होते (पूर्ण झालेल्या सिद्धतेत ते खालच्या बाजूस असतील).
तसेच पुढे काय येईल याचा विचार करून आरंभीची क्रमवर्ती कशी असेल याचा अंदाज
लावणे उपयुक्त ठरते. आपल्या बाबतीत ती $\varphi(a) \Sequent \varphi(a)$ यासारखी असावी
लागेल. याचा अर्थ असा की संख्यापकांचे नियम ``उलट्या दिशेने'' लागू करताना
$\lforall$ नियम आणि $\lexists$ नियम या दोन्हींसाठी एकच पद---ज्याला आपण $a$
म्हणू---निवडावे लागेल. प्रत्येक नियमासाठी वेगवेगळी पदे निवडली, तर शेवटी
$\varphi(a) \Sequent \varphi(b)$ यासारखी क्रमवर्ती मिळेल; ती अर्थातच निष्पन्न करता
येत नाही.

नेहमीप्रमाणे आपण पुढील क्रमवर्ती लिहून सुरुवात करतो:
\begin{prooftree}
\AxiomC{}
\UnaryInf$\lexists[x][\lnot \varphi(x)] \fCenter \lnot \lforall[x][\varphi(x)]$
\end{prooftree}
आपण \LeftR{\exists} नियम किंवा \RightR{\lnot} नियम यांपैकी कोणताही लागू करू
शकतो. \LeftR{\exists} नियमाला आयगेन चराची अट लागू असल्यामुळे तो उशिरा
हाताळण्यापेक्षा लवकर हाताळणे योग्य ठरेल; म्हणून तोच नियम आधी लागू करू.
\begin{prooftree}
\AxiomC{}
\UnaryInf$ \lnot \varphi(a) \fCenter \lnot \lforall[x][\varphi(x)]$
\RightLabel{\LeftR{\lexists}}
\UnaryInf$ \lexists[x][\lnot \varphi(x)] \fCenter \lnot \lforall[x][\varphi(x)]$
\end{prooftree}
\LeftR{\lnot} आणि \RightR{\lnot} नियम उलट्या दिशेने लागू केल्यावर पुढील
निष्पत्ती मिळते:
\begin{prooftree}
\AxiomC{}
\UnaryInf$\lforall[x][\varphi(x)] \fCenter \varphi(a)$
\RightLabel{\LeftR{\lnot}}
\UnaryInf$\lnot \varphi(a), \lforall[x][\varphi(x)] \fCenter $
\RightLabel{\LeftR{\Exchange}}
\UnaryInf$\lforall[x][\varphi(x)], \lnot \varphi(a) \fCenter $
\RightLabel{\RightR{\lnot}}
\UnaryInf$ \lnot \varphi(a) \fCenter \lnot \lforall[x] \varphi(x)$
\RightLabel{\LeftR{\lexists}}
\UnaryInf$ \lexists[x] \lnot \varphi(x) \fCenter \lnot \lforall[x] \varphi(x)$
\end{prooftree}
या टप्प्यावर \LeftR{\forall} नियम लागू करणे हाच आपल्यापुढील एकमेव पर्याय
आहे. या नियमाला आयगेन चराचे बंधन लागू नसल्यामुळे आपला मार्ग मोकळा आहे.
आपल्याला आरंभीची क्रमवर्ती (जिचे स्वरूप $\varphi(a) \Sequent \varphi(a)$ असे आहे)
मिळवण्याचा प्रयत्न करायचा आहे, हे आठवा; म्हणून नियम लागू करताना $\varphi$ मध्ये
घालायचे पद म्हणून $a$ निवडले पाहिजे.
\begin{prooftree}
\Axiom$\varphi(a) \fCenter \varphi(a)$
\RightLabel{\LeftR{\lforall}}
\UnaryInf$\lforall[x][\varphi(x)] \fCenter \varphi(a)$
\RightLabel{\LeftR{\lnot}}
\UnaryInf$\lnot \varphi(a), \lforall[x][\varphi(x)] \fCenter $
\RightLabel{\LeftR{\Exchange}}
\UnaryInf$\lforall[x][\varphi(x)], \lnot \varphi(a) \fCenter $
\RightLabel{\RightR{\lnot}}
\UnaryInf$ \lnot \varphi(a) \fCenter \lnot \lforall[x][\varphi(x)]$
\RightLabel{\LeftR{\lexists}}
\UnaryInf$ \lexists[x][ \lnot \varphi(x)] \fCenter \lnot \lforall[x][\varphi(x)]$
\end{prooftree}
विशेषतः संख्यापक हाताळताना, या टप्प्यावर आयगेन चराच्या अटीचा भंग झालेला नाही
ना, हे पुन्हा तपासणे महत्त्वाचे असते. आयगेन चराची अट लागू असलेला आपण वापरलेला
एकमेव नियम \LeftR{\exists} होता आणि आयगेन चर~$a$ त्याच्या खालच्या क्रमवर्तीत
(म्हणजे अंतिम क्रमवर्तीत) येत नाही; म्हणून ही योग्य निष्पत्ती आहे.
\end{ex}

\begin{prob}
पुढील क्रमवर्तींसाठी निष्पत्त्या तयार करा:
\begin{enumerate}
\item $\Sequent (\lforall[x][\varphi(x)] \land \lforall[y][\psi(y)]) \lif
\lforall[z][(\varphi(z) \land \psi(z))]$.
\item $\Sequent (\lexists[x][\varphi(x)] \lor \lexists[y][\psi(y)]) \lif
\lexists[z][(\varphi(z) \lor \psi(z))]$.
\item $\lforall[x][(\varphi(x) \lif \psi)] \Sequent \lexists[y][\varphi(y)] \lif \psi$.
\item $\lforall[x][\lnot \varphi(x)] \Sequent \lnot\lexists[x][\varphi(x)]$.
\item $\Sequent \lnot\lexists[x][\varphi(x)] \lif \lforall[x][\lnot \varphi(x)]$.
\item $\Sequent \lnot\lexists[x][\lforall[y][((\varphi(x,y) \lif \lnot
\varphi(y,y)) \land (\lnot \varphi(y,y) \lif \varphi(x,y)))]]$.
\end{enumerate}
\end{prob}

\begin{prob}
पुढील क्रमवर्तींसाठी निष्पत्त्या तयार करा:
\begin{enumerate}
\item $\Sequent \lnot\lforall[x][\varphi(x)] \lif \lexists[x][\lnot\varphi(x)]$.
\item $(\lforall[x][\varphi(x)] \lif \psi) \Sequent \lexists[y][(\varphi(y) \lif \psi)]$.
\item $\Sequent \lexists[x][(\varphi(x) \lif \lforall[y][\varphi(y)])]$.
\end{enumerate}
(या सर्वांसाठी $\RightR{\Contraction}$~नियम आवश्यक आहे.)
\end{prob}







\begin{editorial}
  या विभागात नैसर्गिक निगमनासाठी सिद्धतायोग्यता संबंध आणि सुसंगतता यांच्या
  व्याख्या एकत्रित केल्या आहेत.
\end{editorial}



\section{सिद्धता-उपपत्तीय संकल्पना}\label{fol:seq:ptn:sec}

\begin{explain}
जशा आपण अनेक महत्त्वाच्या चिन्हार्थविषयक संकल्पना (वैधता, तार्किक निष्पन्नता,
पूर्ततायोग्यता) परिभाषित केल्या आहेत, तशाच त्यांना अनुरूप
\emph{सिद्धता-उपपत्तीय संकल्पना} आता परिभाषित करू. संरचना मध्ये
वाक्यांची पूर्ति होते की नाही, याचा आधार घेऊन या संकल्पना परिभाषित
केल्या जात नाहीत; त्याऐवजी ठरावीक क्रमवर्तींची निष्पन्नता किंवा
अनिष्पन्नता यांचा आधार घेतला जातो. या दोन प्रकारच्या संकल्पना
एकमेकांशी जुळतात, हा एक महत्त्वाचा शोध होता. त्या जुळतात, हाच
\emph{निर्दोषता प्रमेय} आणि \emph{संपूर्णता प्रमेय} यांचा आशय आहे.
\end{explain}

\begin{defn}[प्रमेये]
वाक्य~$\varphi$ हे \emph{प्रमेय} असते, जर $\Log{LK}$ मध्ये
$\quad \Sequent \varphi$ या क्रमवर्तीची निष्पत्ती असेल. $\varphi$ हे प्रमेय
असेल तर आपण $\Proves \varphi$ लिहितो आणि ते प्रमेय नसेल तर $\Proves/ \varphi$ लिहितो.
\end{defn}

\begin{defn}[निष्पन्नता]
वाक्य~$\varphi$ हे वाक्यांच्या संच~$\Gamma$ \emph{पासून
निष्पन्न करता येण्याजोगे} असते, म्हणजे $\Gamma \Proves \varphi$, तेव्हा आणि केवळ तेव्हाच,
जेव्हा एखादा सांत उपसंच~$\Gamma_0 \subseteq \Gamma$ आणि $\Gamma_0$ मधील
वाक्यांची एखादी क्रमिका $\Gamma_0'$ अशी असते की $\Log{LK}$
$\Gamma_0' \Sequent \varphi$ हे निष्पन्न करते. $\varphi$ हे $\Gamma$ पासून
निष्पन्न करता येण्याजोगे नसेल, तर आपण $\Gamma \Proves/ \varphi$ लिहितो.
\end{defn}

आकुंचन, दुर्बलीकरण आणि अदलाबदल नियमांमुळे $\Gamma_0'$ मधील
वाक्यांचा क्रम आणि त्यांची पुनरावृत्ती कितीदा होते हे महत्त्वाचे नसते:
$\Gamma_0' \Sequent \varphi$ ही क्रमवर्ती निष्पन्न करता येण्याजोगे असेल, तर
$\Gamma_0'$ मध्ये असलेली तीच वाक्ये असणाऱ्या कोणत्याही
$\Gamma_0''$ साठी $\Gamma_0'' \Sequent \varphi$ ही क्रमवर्तीसुद्धा निष्पन्न करता
येते. उदाहरणार्थ, $\Gamma_0 = \{\psi, \chi\}$ असेल, तर
$\Gamma_0' = \tuple{\psi, \psi, \chi}$ आणि $\Gamma_0'' =
\tuple{\chi, \chi, \psi}$ या दोन्ही क्रमिकांमध्ये $\Gamma_0$ मधीलच
वाक्ये आहेत. यांपैकी एक क्रमिका असणारी क्रमवर्ती निष्पन्न करता येण्याजोगे असेल,
तर दुसरी असणारी क्रमवर्तीसुद्धा निष्पन्न करता येते; उदा.:
\begin{prooftree}
  \AxiomC{}
  \Deduce$\psi, \psi, \chi \fCenter \varphi$
  \RightLabel{\LeftR{\Contraction}}
  \UnaryInf$\psi, \chi \fCenter \varphi$
  \RightLabel{\LeftR{\Exchange}}
  \UnaryInf$\chi, \psi \fCenter \varphi$
  \RightLabel{\LeftR{\Weakening}}
  \UnaryInf$\chi, \chi, \psi \fCenter \varphi$
\end{prooftree}
यापुढे $\Gamma_0$ हा वाक्यांचा सांत संच असेल, तर
$\Gamma_0 \Sequent \varphi$ याचा अर्थ असा कोणताही क्रमवर्ती घेऊ की ज्याच्या
पूर्वांगात $\Gamma_0$ मधील वाक्यांची क्रमिका आहे; तसेच आवश्यक असल्यास
आकुंचन, अदलाबदल आणि दुर्बलीकरणे गृहीत धरू.

\begin{defn}[सुसंगतता]
वाक्यांचा संच~$\Gamma$ हा \emph{विसंगत} असतो तेव्हा आणि केवळ तेव्हाच, जेव्हा
एखादा सांत उपसंच~$\Gamma_0 \subseteq \Gamma$ असा असतो की $\Log{LK}$
$\Gamma_0 \Sequent \quad$ हे निष्पन्न करते. $\Gamma$ विसंगत नसेल,
म्हणजे प्रत्येक सांत $\Gamma_0 \subseteq \Gamma$ साठी $\Log{LK}$
$\Gamma_0 \Sequent \quad$ हे निष्पन्न करत नसेल, तर त्याला
\emph{सुसंगत} म्हणतो.
\end{defn}

\begin{prop}[परावर्तिता]
\label{fol:seq:ptn:prop:reflexivity}
$\varphi \in \Gamma$ असेल, तर $\Gamma \Proves \varphi$.
\end{prop}

\begin{proof}
$\varphi \Sequent \varphi$ ही आरंभीची क्रमवर्ती निष्पन्न करता येण्याजोगे आहे आणि $\{\varphi\}
\subseteq \Gamma$.
\end{proof}

\begin{prop}[एकस्वनिकता]
\label{fol:seq:ptn:prop:monotonicity}
$\Gamma \subseteq \Delta$ आणि $\Gamma \Proves \varphi$ असेल, तर $\Delta
\Proves \varphi$.
\end{prop}

\begin{proof}
$\Gamma \Proves \varphi$ असे समजा; म्हणजे, एखादा सांत $\Gamma_0
\subseteq \Gamma$ असा आहे की $\Gamma_0 \Sequent \varphi$ ही क्रमवर्ती
निष्पन्न करता येण्याजोगे आहे. $\Gamma \subseteq \Delta$ असल्यामुळे $\Gamma_0$ हा
$\Delta$ चासुद्धा सांत उपसंच आहे. म्हणून $\Gamma_0 \Sequent \varphi$ ची
निष्पत्ती ही $\Delta \Proves \varphi$ हेसुद्धा दाखवते.
\end{proof}

\begin{prop}[संक्रमकता]
\label{fol:seq:ptn:prop:transitivity}
$\Gamma \Proves \varphi$ आणि $\{\varphi\} \cup \Delta \Proves
\psi$ असेल, तर $\Gamma \cup \Delta \Proves \psi$.
\end{prop}

\begin{proof}
$\Gamma \Proves \varphi$ असेल, तर एखादा सांत $\Gamma_0 \subseteq \Gamma$
आणि $\Gamma_0 \Sequent \varphi$ ची निष्पत्ती~$\pi_0$ असते.
$\{\varphi\} \cup \Delta \Proves \psi$ असेल, तर एखाद्या सांत उपसंचासाठी
$\Delta_0 \subseteq \Delta$, $\varphi, \Delta_0 \Sequent \psi$ ची
निष्पत्ती~$\pi_1$ असते. पुढील निष्पत्ती विचारात घ्या:
\begin{prooftree}
  \AxiomC{}
  \RightLabel{$\pi_0$}
  \Deduce$\Gamma_0 \fCenter \varphi$
  \AxiomC{}
  \RightLabel{$\pi_1$}
  \Deduce$\varphi, \Delta_0 \fCenter \psi$
  \RightLabel{\Cut}
  \BinaryInf$\Gamma_0, \Delta_0 \fCenter \psi$
\end{prooftree}
$\Gamma_0 \cup \Delta_0 \subseteq \Gamma \cup \Delta$ असल्यामुळे यावरून
$\Gamma \cup \Delta \Proves \psi$ हे दिसते.
\end{proof}

विशेषतः, याचा अर्थ असा की $\Gamma \Proves \varphi$ आणि $\varphi
\Proves \psi$ असेल, तर $\Gamma \Proves \psi$. तसेच $\varphi_1,
\dots, \varphi_n \Proves \psi$ असेल आणि प्रत्येक~$i$ साठी $\Gamma \Proves \varphi_i$
असेल, तर $\Gamma \Proves \psi$, हेसुद्धा यावरून मिळते.

\begin{prop}
\label{fol:seq:ptn:prop:incons}
$\Gamma$ विसंगत असतो तेव्हा आणि केवळ तेव्हाच, जेव्हा प्रत्येक
  वाक्य~$\varphi$ साठी $\Gamma \Proves {\varphi}$.
\end{prop}

\begin{proof}
सराव.
\end{proof}

\begin{prob}
\ref{fol:seq:ptn:prop:incons} सिद्ध करा.
\end{prob}

\begin{prop}[संहतता]
  \label{fol:seq:ptn:prop:proves-compact}
  \begin{enumerate}
  \item $\Gamma \Proves \varphi$ असेल, तर एखादा सांत उपसंच $\Gamma_0
    \subseteq \Gamma$ असा आहे की $\Gamma_0 \Proves \varphi$.
  \item $\Gamma$ चा प्रत्येक सांत उपसंच सुसंगत असेल, तर $\Gamma$ सुसंगत
    असतो.
  \end{enumerate}
\end{prop}

\begin{proof}
  \begin{enumerate}
    \item $\Gamma \Proves \varphi$ असेल, तर एखादा सांत उपसंच
      $\Gamma_0 \subseteq \Gamma$ असा आहे की $\Gamma_0
      \Sequent \varphi$ या क्रमवर्तीची निष्पत्ती आहे. परिणामी,
      $\Gamma_0 \Proves \varphi$.
    \item $\Gamma$ विसंगत असेल, तर एखादा सांत
      उपसंच~$\Gamma_0 \subseteq \Gamma$ असा आहे की $\Log{LK}$
      $\Gamma_0 \Sequent \quad$ हे निष्पन्न करते. पण मग $\Gamma_0$ हा
      $\Gamma$ चा विसंगत सांत उपसंच आहे.
  \end{enumerate}
\end{proof}









\section{निष्पन्नता आणि सुसंगतता}\label{fol:seq:prv:sec}

आता आपण निष्पन्नता संबंधाचे अनेक गुणधर्म सिद्ध करू. ते स्वतंत्रपणेही
महत्त्वाचे आहेत, पण संपूर्णता प्रमेयाच्या सिद्धतेत प्रत्येकाची भूमिका असेल.

\begin{prop}\label{fol:seq:prv:prop:provability-contr}
  जर $\Gamma \Proves \varphi$ आणि $\Gamma \cup \{\varphi\}$ विसंगत असेल, तर
  $\Gamma$ विसंगत आहे.
\end{prop}

\begin{proof}
अशी सांत $\Gamma_0$ आणि $\Gamma_1 \subseteq \Gamma$ असतात की
$\Log{LK}$ निष्पन्न $\Gamma_0 \Sequent \varphi$ आणि $\varphi, \Gamma_1
\Sequent \quad$. $\Gamma_0 \Sequent \varphi$ ची $\Log{LK}$-निष्पत्ती
~$\pi_0$ आणि $\Gamma_1, \varphi \Sequent \quad$ ची $\Log{LK}$-निष्पत्ती
~$\pi_1$ अशी नावे देऊ. मग आपण निष्पन्न करू शकतो
\begin{prooftree}
\AxiomC{}
\RightLabel{$\pi_0$}
\Deduce$ \Gamma_0 \fCenter \varphi $
\AxiomC{}
\RightLabel{$\pi_1$}
\Deduce$\varphi, \Gamma_1 \fCenter $
\RightLabel{\Cut}
\BinaryInf$ \Gamma_0,\Gamma_1 \fCenter $
\end{prooftree}

कारण $\Gamma_0 \subseteq \Gamma$ आणि $\Gamma_1 \subseteq \Gamma$,
$\Gamma_0 \cup \Gamma_1 \subseteq \Gamma$; म्हणून $\Gamma$ विसंगत आहे.
\end{proof}

\begin{prop}
\label{fol:seq:prv:prop:prov-incons}
$\Gamma \Proves \varphi$ तेव्हा आणि केवळ तेव्हाच, जेव्हा $\Gamma \cup \{\lnot \varphi\}$
विसंगत असते.
\end{prop}

\begin{proof}
प्रथम $\Gamma \Proves \varphi$ असे गृहीत धरा; म्हणजे $\Gamma \Sequent \varphi$ ची
निष्पत्ती~$\pi_0$ आहे. $\LeftR{\lnot}$ नियम जोडल्यावर आपल्याला
$\lnot \varphi, \Gamma
\Sequent \quad$ ची निष्पत्ती मिळते; म्हणजे $\Gamma \cup \{\lnot \varphi\}$
विसंगत आहे.

जर $\Gamma \cup \{\lnot \varphi\}$ विसंगत असेल, तर $\lnot \varphi, \Gamma
\Sequent \quad$ ची निष्पत्ती~$\pi_1$ असते. खालील $\Gamma \Sequent \varphi$
ची निष्पत्ती आहे:
\begin{prooftree}
  \Axiom$\varphi \fCenter \varphi$
  \RightLabel{\RightR{\lnot}}
  \UnaryInf$\fCenter \varphi, \lnot \varphi$
  \AxiomC{}
  \RightLabel{$\pi_1$}
  \Deduce$\lnot \varphi, \Gamma \fCenter$
  \RightLabel{\Cut}
  \BinaryInf$\Gamma \fCenter \varphi$
\end{prooftree}
\end{proof}

\begin{prob}
$\Gamma \Proves \lnot \varphi$ तेव्हा आणि केवळ तेव्हाच, जेव्हा $\Gamma \cup \{\varphi\}$
विसंगत असते, हे सिद्ध करा.
\end{prob}

\begin{prop}\label{fol:seq:prv:prop:explicit-inc}
  जर $\Gamma \Proves \varphi$ आणि $\lnot \varphi \in \Gamma$ असेल, तर $\Gamma$
  विसंगत आहे.
\end{prop}

\begin{proof}
  समजा $\Gamma \Proves \varphi$ आणि $\lnot \varphi \in \Gamma$. मग $\Gamma_0 \Sequent
  \varphi$ ची निष्पत्ती~$\pi$ असते. $\lnot \varphi, \Gamma_0 \Sequent \quad$
  ही क्रमवर्तीही निष्पन्न करता येण्याजोगे आहे:
  \begin{prooftree}
    \AxiomC{}
    \RightLabel{$\pi$}
    \Deduce$\Gamma_0 \fCenter \varphi$
    \Axiom$\varphi \fCenter \varphi$
    \RightLabel{\LeftR{\lnot}}
    \UnaryInf$\lnot \varphi, \varphi \fCenter$
    \RightLabel{\LeftR{\Exchange}}
    \UnaryInf$\varphi, \lnot \varphi \fCenter$
    \RightLabel{\Cut}
    \BinaryInf$\Gamma_0, \lnot \varphi \fCenter$
  \end{prooftree}
  कारण $\lnot \varphi \in \Gamma$ आणि $\Gamma_0 \subseteq \Gamma$, यावरून
  $\Gamma$ विसंगत असल्याचे दिसते.
\end{proof}

\begin{prop}\label{fol:seq:prv:prop:provability-exhaustive}
  जर $\Gamma \cup \{\varphi\}$ आणि $\Gamma \cup \{\lnot \varphi\}$ दोन्ही
  विसंगत असतील, तर $\Gamma$ विसंगत आहे.
\end{prop}

\begin{proof}
अशी सांत संच $\Gamma_0 \subseteq \Gamma$ आणि $\Gamma_1
\subseteq \Gamma$, तसेच $\Log{LK}$-निष्पत्त्या $\pi_0$ आणि $\pi_1$
$\varphi, \Gamma_0 \Sequent \quad$ आणि $\lnot \varphi, \Gamma_1 \Sequent
\quad$ यांचे अनुक्रमे, असतात. मग आपण निष्पन्न करू शकतो
\begin{prooftree}
\AxiomC{}
\RightLabel{$\pi_0$}
\Deduce$ \varphi, \Gamma_0 \fCenter $
\RightLabel{\RightR{\lnot}}
\UnaryInf$ \Gamma_0 \fCenter \lnot \varphi$
\AxiomC{}
\RightLabel{$\pi_1$}
\Deduce$\lnot \varphi, \Gamma_1 \fCenter  $
\RightLabel{\Cut}
\BinaryInf$ \Gamma_0, \Gamma_1 \fCenter $
\end{prooftree}
कारण $\Gamma_0 \subseteq \Gamma$ आणि $\Gamma_1 \subseteq \Gamma$,
$\Gamma_0 \cup \Gamma_1 \subseteq \Gamma$. म्हणून $\Gamma$ विसंगत आहे.
\end{proof}










\section{निष्पन्नता आणि विधानीय संयोजक}\label{fol:seq:ppr:sec}

\begin{explain}
  क्रमवर्ती कलनातील निष्पन्नता संबंध~$\Proves$ हा विधानीय संयोजकांशी
  संबंधित काही मूलभूत तथ्ये सिद्ध करण्यासाठी पुरेसा सामर्थ्यवान आहे; उदाहरणार्थ,
  $\varphi \land \psi
  \Proves \varphi$ आणि $\varphi, \varphi \lif \psi \Proves \psi$ (विध्यात्मक अनुमान). ही
  तथ्ये संपूर्णता प्रमेयाच्या सिद्धतेसाठी आवश्यक आहेत.
\end{explain}

\begin{prop}\label{fol:seq:ppr:prop:provability-land}
  \begin{enumerate}
  \item \label{fol:seq:ppr:prop:provability-land-left} $\varphi \land \psi \Proves
    \varphi$ आणि $\varphi \land \psi \Proves \psi$ ही दोन्ही तथ्ये सत्य आहेत.
  \item \label{fol:seq:ppr:prop:provability-land-right} $\varphi, \psi \Proves \varphi \land
    \psi$.
  \end{enumerate}
\end{prop}

\begin{proof}
  \begin{enumerate}
  \item $\varphi \land \psi \Sequent \varphi$ आणि $\varphi \land \psi \Sequent
    \psi$ या दोन्ही क्रमवर्ती निष्पन्न करता येण्याजोगे आहेत:
    \begin{prooftree}
      \Axiom$\varphi \fCenter \varphi$
      \RightLabel{\LeftR{\land}}
      \UnaryInf$\varphi \land \psi \fCenter \varphi$
      \DisplayProof\qquad\bottomAlignProof
      \Axiom$\psi \fCenter \psi$
      \RightLabel{\LeftR{\land}}
      \UnaryInf$\varphi \land \psi \fCenter \psi$
    \end{prooftree}
    \item $\varphi, \psi \Sequent \varphi \land \psi$ या क्रमवर्तीची निष्पत्ती येथे आहे:
    \begin{prooftree}
      \Axiom$\varphi \fCenter \varphi$
      \Axiom$\psi \fCenter \psi$
      \RightLabel{\RightR{\land}}
      \BinaryInf$\varphi, \psi \fCenter \varphi \land \psi$
    \end{prooftree}
  \end{enumerate}
\end{proof}

\begin{prop}\label{fol:seq:ppr:prop:provability-lor}
  \begin{enumerate}
  \item $\varphi \lor \psi, \lnot \varphi, \lnot \psi$ विसंगत आहे.
  \item $\varphi \Proves \varphi \lor \psi$ आणि $\psi \Proves \varphi \lor \psi$ ही दोन्ही तथ्ये आहेत.
  \end{enumerate}
\end{prop}

\begin{proof}
  \begin{enumerate}
  \item $\varphi \lor \psi, \lnot \varphi,
    \lnot \psi \Sequent$ या क्रमवर्तीची निष्पत्ती देऊ:
    \begin{prooftree}
      \Axiom$\varphi \fCenter \varphi$
      \RightLabel{\LeftR{\lnot}}
      \UnaryInf$\lnot \varphi, \varphi \fCenter$
      \doubleLine
      \UnaryInf$\varphi, \lnot \varphi, \lnot \psi \fCenter$
      \Axiom$\psi \fCenter \psi$
      \RightLabel{\LeftR{\lnot}}
      \UnaryInf$\lnot \psi, \psi \fCenter$
      \doubleLine
      \UnaryInf$\psi, \lnot \varphi, \lnot \psi \fCenter$
      \RightLabel{\LeftR{\lor}}
      \BinaryInf$ \varphi \lor \psi, \lnot \varphi, \lnot \psi \fCenter $
    \end{prooftree}
    (दुहेरी अनुमान-रेषा दुर्बलीकरण, आकुंचन आणि अदलाबदल यांच्या अनेक अनुमानांना
    दर्शवतात, हे लक्षात ठेवा.)
  \item   $\varphi \Sequent \varphi \lor \psi$ आणि $\psi \Sequent \varphi
    \lor \psi$ या दोन्ही क्रमवर्तींना निष्पत्त्या आहेत:
    \begin{prooftree}
      \Axiom$\varphi \fCenter \varphi$
      \RightLabel{\RightR{\lor}}
      \UnaryInf$\varphi \fCenter \varphi \lor \psi$
      \DisplayProof\qquad\bottomAlignProof
      \Axiom$\psi \fCenter \psi$
      \RightLabel{\RightR{\lor}}
      \UnaryInf$\psi \fCenter \varphi \lor \psi$
    \end{prooftree}
  \end{enumerate}
\end{proof}

\begin{prop}\label{fol:seq:ppr:prop:provability-lif}
  \begin{enumerate}
  \item \label{fol:seq:ppr:prop:provability-lif-left}  $\varphi, \varphi \lif \psi \Proves \psi$.
  \item \label{fol:seq:ppr:prop:provability-lif-right}
    $\lnot \varphi \Proves \varphi \lif \psi$ आणि $\psi \Proves \varphi \lif \psi$ ही दोन्ही तथ्ये आहेत.
  \end{enumerate}
\end{prop}

\begin{proof}
  \begin{enumerate}
    \item $\varphi \lif \psi, \varphi \Sequent \psi$ ही क्रमवर्ती निष्पन्न करता येण्याजोगे आहे:
      \begin{prooftree}
        \Axiom$\varphi \fCenter \varphi$
        \Axiom$\psi \fCenter \psi$
        \RightLabel{\LeftR{\lif}}
        \BinaryInf$\varphi \lif \psi, \varphi  \fCenter \psi$
      \end{prooftree}
    \item $\lnot \varphi \Sequent \varphi \lif \psi$ आणि $\psi
      \Sequent \varphi \lif \psi$ या दोन्ही क्रमवर्ती निष्पन्न करता येण्याजोगे आहेत:
      \begin{prooftree}
        \Axiom$\varphi \fCenter \varphi$
        \RightLabel{\LeftR{\lnot}}
        \UnaryInf$\lnot \varphi, \varphi \fCenter$
        \RightLabel{\LeftR{\Exchange}}
        \UnaryInf$\varphi, \lnot \varphi \fCenter$
        \RightLabel{\RightR{\Weakening}}
        \UnaryInf$\varphi, \lnot \varphi \fCenter \psi$
        \RightLabel{\RightR{\lif}}
        \UnaryInf$\lnot \varphi \fCenter \varphi \lif \psi$
        \DisplayProof\qquad\bottomAlignProof
        \Axiom$\psi \fCenter \psi$
        \RightLabel{\LeftR{\Weakening}}
        \UnaryInf$\varphi, \psi \fCenter \psi$
        \RightLabel{\RightR{\lif}}
        \UnaryInf$\psi \fCenter \varphi \lif \psi$
      \end{prooftree}
  \end{enumerate}
\end{proof}








\section{निष्पन्नता आणि संख्यापक}\label{fol:seq:qpr:sec}

\begin{explain}
  संपूर्णता प्रमेयासाठी या विभागात स्थापित केलेली~$\Proves$ संबंधी तथ्ये
  क्रमवर्ती कलनातील नियमांनी निष्पन्न होतात, हेही आवश्यक आहे.
\end{explain}

\begin{thm}
\label{fol:seq:qpr:thm:strong-generalization} जर $c$ हे $\Gamma$ किंवा $\varphi(x)$ मध्ये
न येणारे स्थिर असेल आणि $\Gamma \Proves \varphi(c)$, तर $\Gamma
\Proves \lforall[x][\varphi(x)]$.
\end{thm}

\begin{proof}
$\pi_0$ ही $\Gamma_0 \Sequent \varphi(c)$ या क्रमवर्तीची $\Log{LK}$-निष्पत्ती
असलेली एखादी सांत $\Gamma_0 \subseteq \Gamma$ घ्या.  $\RightR{\lforall}$
अनुमान जोडल्यावर आपल्याला $\Gamma_0 \Sequent
\lforall[x][\varphi(x)]$ ची निष्पत्ती मिळते, कारण $c$ हे $\Gamma$ किंवा
$\varphi(x)$ मध्ये येत नाही आणि म्हणून आयगेन चराची अट पूर्ण होते.
\end{proof}

\begin{prop}
\label{fol:seq:qpr:prop:provability-quantifiers}
\begin{enumerate}
\item $\varphi(t) \Proves \lexists[x][\varphi(x)]$.
\item $\lforall[x][\varphi(x)] \Proves \varphi(t)$.
\end{enumerate}
\end{prop}

\begin{proof}
\begin{enumerate}
\item $\varphi(t) \Sequent \lexists[x][\varphi(x)]$ ही क्रमवर्ती
  निष्पन्न करता येण्याजोगे आहे:
  \begin{prooftree}
    \Axiom$\varphi(t) \fCenter \varphi(t)$
    \RightLabel{\RightR{\lexists}}
    \UnaryInf$\varphi(t) \fCenter \lexists[x][\varphi(x)]$
    \end{prooftree}
\item $\lforall[x][\varphi(x)] \Sequent \varphi(t)$ ही क्रमवर्ती
  निष्पन्न करता येण्याजोगे आहे:
  \begin{prooftree}
    \Axiom$\varphi(t) \fCenter \varphi(t)$
    \RightLabel{\LeftR{\lforall}}
    \UnaryInf$\lforall[x][\varphi(x)] \fCenter \varphi(t)$
    \end{prooftree}
\end{enumerate}
\end{proof}









\section{निर्दोषता}\label{fol:seq:sou:sec}

\begin{explain}
क्रमवर्ती कलनासारखी निष्पत्ती-पद्धत \emph{निर्दोष} असते, जर ती
प्रत्यक्षात धारण न करणाऱ्या गोष्टी निष्पन्न करू शकत नसेल. त्यामुळे
निर्दोषता हा निष्पत्ती-पद्धतींसाठी हमी दिलेल्या सुरक्षिततेचा एक प्रकार
आहे. कोणता सिद्धता-उपपत्तीय गुणधर्म विचारात आहे यावर अवलंबून, उदाहरणार्थ,
आपल्याला पुढील गोष्टी हव्या असतात:
\begin{enumerate}
\item प्रत्येक निष्पन्न करता येण्याजोगे~$\varphi$ हे वैध असते;
\item एखादे वाक्य इतर काही वाक्यांपासून निष्पन्न करता येण्याजोगे असेल, तर ते
  त्यांचे तार्किक परिणामही असते;
\item वाक्यांचा संच विसंगत असेल, तर तो अपूर्ततायोग्य असतो.
\end{enumerate}
हे निष्पत्ती-पद्धतीचे महत्त्वाचे गुणधर्म आहेत. यांपैकी एखादा गुणधर्म
खरा नसेल, तर निष्पत्ती-पद्धतीत दोष आहे—ती खूप जास्त निष्पत्ती
निष्पन्न करेल. म्हणून निष्पत्ती-पद्धतीची निर्दोषता सिद्ध करणे अत्यंत
महत्त्वाचे आहे.

वरील सर्व सिद्धता-उपपत्तीय गुणधर्म क्रमवर्ती कलनातील ठरावीक क्रमवर्तींच्या
निष्पन्नता द्वारे परिभाषित केलेले असल्याने, (1)--(3) सिद्ध करण्यासाठी
निष्पन्न करता येण्याजोगे क्रमवर्तींच्या चिन्हार्थविषयक गुणधर्मांबद्दल काहीतरी सिद्ध करावे
लागते. प्रथम क्रमवर्ती \emph{वैध} असणे म्हणजे काय ते परिभाषित करू आणि मग
प्रत्येक निष्पन्न करता येण्याजोगे क्रमवर्ती वैध असते हे दाखवू. त्यानंतर (1)--(3) या
परिणामाच्या उपपरिणामांप्रमाणे मिळतील.
\end{explain}

\begin{defn}
संरचना~$\Struct
  M$ क्रमवर्ती
$\Gamma \Sequent \Delta$ \emph{पूर्ण करते} तेव्हा आणि केवळ तेव्हाच, जेव्हा
$\Sat/{M}{\varphi}$ हे एखाद्या $\varphi \in \Gamma$ साठी किंवा
$\Sat{M}{\varphi}$ हे एखाद्या $\varphi \in \Delta$ साठी असते.

क्रमवर्ती \emph{वैध} असते तेव्हा आणि केवळ तेव्हाच, जेव्हा प्रत्येक
संरचना~$\Struct
  M$ ते पूर्ण करते.
\end{defn}

\begin{thm}[निर्दोषता]
\label{fol:seq:sou:thm:sequent-soundness} जर $\Log{LK}$ निष्पन्न $\Theta
\Sequent \Xi$, तर $\Theta \Sequent \Xi$ वैध आहे.
\end{thm}

\begin{proof}
  $\Theta \Sequent \Xi$ ची निष्पत्ती~$\pi$ असू द्या. $\pi$ मधील
अनुमानांची संख्या~$n$ यावर विगमनाने पुढे जाऊ.

अनुमानांची संख्या~$0$ असेल, तर $\pi$ मध्ये फक्त आरंभीची क्रमवर्ती असते.
प्रत्येक आरंभीची क्रमवर्ती $\varphi \Sequent \varphi$ स्पष्टपणे वैध असते, कारण प्रत्येक
$\Struct M$ साठी एकतर
  $\Sat/{M}{\varphi}$ किंवा
$\Sat{M}{\varphi}$.

  अनुमानांची संख्या~0 पेक्षा जास्त असेल, तर सर्वांत खालच्या अनुमानाच्या
प्रकारानुसार प्रकरणे वेगळी करू. विगमन गृहीतकानुसार त्या अनुमानाची आधारविधाने
वैध आहेत असे गृहीत धरू, कारण कोणत्याही आधारविधानाच्या निष्पत्तीमधील
अनुमानांची संख्या~$n$ पेक्षा कमी असते.

प्रथम, एकच आधारविधान असलेली संभाव्य अनुमाने पाहू.
\begin{enumerate}
\item शेवटचे अनुमान दुर्बलीकरण आहे. मग $\Theta \Sequent \Xi$ हे एकतर
  $\varphi, \Gamma \Sequent \Delta$ (शेवटचे अनुमान
  \LeftR{\Weakening} असेल) किंवा $\Gamma \Sequent \Delta, \varphi$ (ते
  \RightR{\Weakening} असेल), आणि निष्पत्ती पुढीलपैकी एका रूपात संपते:
  \begin{prooftree}
    \AxiomC{}
    \Deduce$\Gamma \fCenter \Delta$
    \RightLabel{\LeftR{\Weakening}}
    \UnaryInf$\varphi, \Gamma \fCenter \Delta$
    \DisplayProof\qquad\bottomAlignProof
    \AxiomC{}
    \Deduce$\Gamma \fCenter \Delta$
    \RightLabel{\RightR{\Weakening}}
    \UnaryInf$\Gamma \fCenter \Delta, \varphi$
  \end{prooftree}
  विगमन गृहीतकानुसार $\Gamma \Sequent \Delta$ वैध आहे; म्हणजे प्रत्येक
  संरचना~$\Struct{M}$,
  एकतर एखादे $\chi \in \Gamma$ असे आहे की
  $\Sat/{M}{\chi}$, किंवा एखादे $\chi \in
  \Delta$ असे आहे की $\Sat{M}{\chi}$.

  जर $\Sat/{M}{\chi}$ हे एखाद्या $\chi \in \Gamma$ साठी
  असेल, तर $\Theta = \varphi, \Gamma$ असल्याने $\chi \in \Theta$ देखील आहे आणि
  म्हणून $\Sat/{M}{\chi}$ हे एखाद्या $\chi \in \Theta$ साठी
  आहे. त्याचप्रमाणे, जर $\Sat{M}{\chi}$ हे एखाद्या
  $\chi \in \Delta$ साठी असेल, तर $\chi \in \Xi$ आणि
  $\Sat{M}{\chi}$ हे एखाद्या $\chi \in \Xi$ साठी आहे.
  म्हणून $\Theta \Sequent \Xi$ वैध आहे.
\item शेवटचे अनुमान \LeftR{\lnot} आहे. मग त्याचे आधारविधान
  $\Gamma \Sequent \Delta, \varphi$ आणि निष्कर्ष $\lnot \varphi, \Gamma \Sequent
  \Delta$ असतो; म्हणजे निष्पत्ती पुढील रूपात संपते:
  \begin{prooftree}
    \AxiomC{}
    \Deduce$\Gamma \fCenter \Delta, \varphi$
    \RightLabel{\LeftR{\lnot}}
    \UnaryInf$\lnot \varphi, \Gamma \fCenter \Delta$
  \end{prooftree}
  आणि $\Theta = \lnot \varphi, \Gamma$, तर $\Xi = \Delta$.

  विगमन गृहीतकानुसार $\Gamma \Sequent \Delta, \varphi$ वैध आहे; म्हणजे प्रत्येक
  $\Struct{M}$ साठी एकतर (a) एखाद्या
  $\chi \in \Gamma$ साठी
  $\Sat/{M}{\chi}$, किंवा (b) एखाद्या $\chi \in
  \Delta$ साठी, $\Sat{M}{\chi}$, किंवा (c)
  $\Sat{M}{\varphi}$. $\Theta
  \Sequent \Xi$ देखील वैध आहे हे दाखवायचे आहे. $\Struct{M}$
    ही संरचना घ्या. (a) असल्यास,
  एखादे $\chi \in \Gamma$ असे आहे की
  $\Sat/{M}{\chi}$, आणि $\chi \in \Theta$ मध्ये
  देखील आहे. (b) असल्यास, एखादे $\chi \in \Delta$ असे आहे की
  $\Sat{M}{\chi}$, आणि $\chi \in \Xi$ मध्ये
  देखील आहे. शेवटी, जर $\Sat{M}{\varphi}$ असेल, तर
  $\Sat/{M}{\lnot \varphi}$. कारण $\lnot \varphi \in
  \Theta$ मध्ये असल्याने, एखादे $\chi \in \Theta$ असे आहे की
  $\Sat/{M}{\chi}$. म्हणून $\Theta
  \Sequent \Xi$ वैध आहे.
\item शेवटचे अनुमान \RightR{\lnot} आहे: सरावासाठी.
\item शेवटचे अनुमान \LeftR{\land} आहे. याचे दोन प्रकार आहेत: आधारविधानाच्या
  डाव्या बाजूला असलेल्या $\varphi$ किंवा $\psi$ पासून $\varphi
  \land \psi$ डावीकडे निष्पन्न केले जाऊ शकते. पहिल्या प्रकारात $\pi$ पुढील रूपात
  संपते:
  \begin{prooftree}
    \AxiomC{}
    \Deduce$\varphi, \Gamma \fCenter \Delta$
    \RightLabel{\LeftR{\land}}
    \UnaryInf$\varphi \land \psi, \Gamma \fCenter \Delta$
  \end{prooftree}
  आणि $\Theta = \varphi \land \psi, \Gamma$, तर $\Xi = \Delta$. आता
  संरचना~$\Struct
  M$ घ्या. विगमन गृहीतकानुसार
  $\varphi, \Gamma \Sequent \Delta$ वैध असल्याने, (a)
  $\Sat/{M}{\varphi}$, (b)
  $\Sat/{M}{\chi}$ हे एखाद्या $\chi \in \Gamma$ साठी, किंवा
  (c)~$\Sat{M}{\chi}$ हे एखाद्या $\chi \in \Delta$ साठी.
  (a) प्रकरणात $\Sat/{M}{\varphi \land \psi}$, म्हणून
  $\chi \in \Theta$ (म्हणजे $\varphi \land \psi$) असे आहे की
  $\Sat/{M}{\chi}$. (b) प्रकरणात एखादे $\chi
  \in \Gamma$ असे आहे की $\Sat/{M}{\chi}$ आणि
  $\chi \in \Theta$ देखील आहे. (c) प्रकरणात एखादे $\chi \in \Delta$ असे आहे
  की $\Sat{M}{\chi}$ आणि $\Xi = \Delta$ असल्याने
  $\chi \in \Xi$ देखील आहे. म्हणून प्रत्येक प्रकरणात
  $\Struct
    M$ $\varphi \land \psi, \Gamma \Sequent
  \Delta$ पूर्ण करते. $\Struct{M}$ स्वेच्छ
  असल्याने $\Gamma \Sequent \Delta$ वैध आहे. $\varphi \land \psi$ हे $\psi$ पासून
  निष्पन्न होण्याचे प्रकरणही असेच हाताळता येते; फक्त $\varphi$ ऐवजी $\psi$ घ्या.
\item शेवटचे अनुमान \RightR{\lor} आहे. याचे दोन प्रकार आहेत: आधारविधानाच्या
  उजव्या बाजूला असलेल्या $\varphi$ किंवा $\psi$ पासून $\varphi
  \lor \psi$ उजवीकडे निष्पन्न केले जाऊ शकते. पहिल्या प्रकारात $\pi$ पुढील रूपात
  संपते:
  \begin{prooftree}
    \AxiomC{}
    \Deduce$\Gamma \fCenter \Delta, \varphi$
    \RightLabel{\RightR{\lor}}
    \UnaryInf$\Gamma \fCenter \Delta, \varphi \lor \psi$
  \end{prooftree}
  येथे $\Theta = \Gamma$ आणि $\Xi = \Delta, \varphi \lor \psi$. आता
  संरचना~$\Struct{M}$ घ्या.
  $\Gamma \Sequent \Delta, \varphi$ वैध असल्याने, (a)
  $\Sat{M}{\varphi}$, (b)
  $\Sat/{M}{\chi}$ हे एखाद्या $\chi \in \Gamma$ साठी, किंवा
  (c)~$\Sat{M}{\chi}$ हे एखाद्या $\chi \in \Delta$ साठी.
  (a) प्रकरणात $\Sat{M}{\varphi \lor \psi}$. (b) प्रकरणात
  एखादे $\chi \in \Gamma$ असे आहे की
  $\Sat/{M}{\chi}$. (c) प्रकरणात एखादे $\chi
  \in \Delta$ असे आहे की $\Sat{M}{\chi}$. म्हणून
  प्रत्येक प्रकरणात $\Struct{M}$
  $\Gamma \Sequent \Delta, \varphi \lor \psi$, म्हणजेच $\Theta \Sequent \Xi$ पूर्ण
  करते. $\Struct{M}$ स्वेच्छ असल्याने
  $\Theta \Sequent \Xi$ वैध आहे. $\varphi \lor \psi$ हे $\psi$ पासून निष्पन्न होण्याचे
  प्रकरणही असेच हाताळता येते; फक्त $\varphi$ ऐवजी $\psi$ घ्या.
\item शेवटचे अनुमान \RightR{\lif} आहे. मग $\pi$ पुढील रूपात संपते:
  \begin{prooftree}
    \AxiomC{}
    \Deduce$\varphi, \Gamma \fCenter \Delta, \psi$
    \RightLabel{\RightR{\lif}}
    \UnaryInf$\Gamma \fCenter \Delta, \varphi \lif \psi$
  \end{prooftree}
  पुन्हा, विगमन गृहीतकानुसार आधारविधान वैध आहे; निष्कर्षही वैध आहे हे
  दाखवायचे आहे. $\Struct{M}$ स्वेच्छ घ्या.
  $\varphi, \Gamma \Sequent \Delta, \psi$ वैध असल्याने पुढीलपैकी किमान एक प्रकरण
  लागू होते: (a) $\Sat/{M}{\varphi}$, (b)
  $\Sat{M}{\psi}$, (c)
  $\Sat/{M}{\chi}$ हे एखाद्या~$~\chi \in \Gamma$ साठी, किंवा
  (d) $\Sat{M}{\chi}$ हे एखाद्या $\chi \in \Delta$ साठी.
  (a) आणि (b) प्रकरणांत $\Sat{M}{\varphi \lif \psi}$
  आणि म्हणून $\chi \in \Delta, \varphi \lif \psi$ असे आहे की
  $\Sat{M}{\chi}$. (c) प्रकरणात एखाद्या $\chi \in
  \Gamma$ साठी $\Sat/{M}{\chi}$. (d) प्रकरणात
  एखाद्या $\chi \in \Delta$ साठी $\Sat{M}{\chi}$. प्रत्येक
  प्रकरणात $\Struct{M}$ $\Gamma
  \Sequent \Delta, \varphi \lif \psi$ पूर्ण करते. $\Struct{M}$
  स्वेच्छ असल्याने $\Gamma \Sequent \Delta, \varphi \lif \psi$ वैध आहे.
\item शेवटचे अनुमान \LeftR{\lforall} आहे. मग सूत्र~$\varphi(x)$ आणि
  बंद पद~$t$ असे आहेत की $\pi$ पुढील रूपात संपते:
  \begin{prooftree}
    \AxiomC{}
    \Deduce$\varphi(t), \Gamma \fCenter \Delta$
    \RightLabel{\LeftR{\lforall}}
    \UnaryInf$\lforall[x][\varphi(x)], \Gamma \fCenter \Delta$
  \end{prooftree}
  निष्कर्ष $\lforall[x][\varphi(x)], \Gamma
  \Sequent \Delta$ वैध आहे हे दाखवायचे आहे. संरचना~$\Struct M$ घ्या.
  आधारविधान $\varphi(t), \Gamma \Sequent \Delta$ वैध असल्याने, (a)
  $\Sat/{M}{\varphi(t)}$, (b) $\Sat/{M}{\chi}$ हे एखाद्या $\chi \in \Gamma$ साठी, किंवा
  (c)~$\Sat{M}{\chi}$ हे एखाद्या $\chi \in \Delta$ साठी असते. (a) प्रकरणात,
  \readerexternalref{fol:syn:sem:prop:quant-terms} नुसार, जर
  $\Sat{M}{\lforall[x][\varphi(x)]}$ असेल, तर $\Sat{M}{\varphi(t)}$. म्हणून
  $\Sat/{M}{\varphi(t)}$ असल्याने $\Sat/{M}{\lforall[x][\varphi(x)]}$. (b) आणि
  (c) प्रकरणांत $\Struct{M}$ देखील $\lforall[x][\varphi(x)], \Gamma
  \Sequent \Delta$ पूर्ण करते. $\Struct M$ स्वेच्छ असल्याने,
  $\lforall[x][\varphi(x)], \Gamma \Sequent \Delta$ वैध आहे.
\item शेवटचे अनुमान \RightR{\lexists} आहे: सरावासाठी.
\item शेवटचे अनुमान \RightR{\lforall} आहे. मग
  सूत्र~$\varphi(x)$ आणि स्थिरांक~$a$ असे आहेत की $\pi$ पुढील रूपात
  संपते:
  \begin{prooftree}
    \AxiomC{}
    \Deduce$\Gamma \fCenter \Delta, \varphi(a)$
    \RightLabel{\RightR{\lforall}}
    \UnaryInf$\Gamma \fCenter \Delta, \lforall[x][\varphi(x)]$
  \end{prooftree}
  येथे आयगेन चराची अट पूर्ण आहे, म्हणजे $a$ हे $\varphi(x)$, $\Gamma$ किंवा
  $\Delta$ मध्ये येत नाही. विगमन गृहीतकानुसार शेवटच्या अनुमानाचे आधारविधान
  वैध आहे. निष्कर्षही वैध आहे हे दाखवायचे आहे; म्हणजे कोणत्याही
  संरचना~$\Struct M$ साठी (a) $\Sat{M}{\lforall[x][\varphi(x)]}$, (b)
  एखाद्या $\chi \in \Gamma$ साठी $\Sat/{M}{\chi}$, किंवा (c) एखाद्या
  $\chi \in \Delta$ साठी $\Sat{M}{\chi}$.

  $\Struct{M}$ ही स्वेच्छ संरचना घ्या. (b) किंवा (c) लागू असेल तर
  काम झाले; म्हणून दोन्ही लागू नाहीत असे गृहीत धरा: प्रत्येक $\chi \in
  \Gamma$ साठी $\Sat{M}{\chi}$ आणि प्रत्येक $\chi \in \Delta$ साठी
  $\Sat/{M}{\chi}$. आता (a), म्हणजेच $\Sat{M}{\lforall[x][\varphi(x)]}$, हे दाखवायचे
  आहे. \readerexternalref{fol:syn:ass:prop:sat-quant} नुसार सर्व चल-विन्यास~$s$ साठी
  $\Sat{M}{\varphi(x)}[s]$ दाखवणे पुरेसे आहे. म्हणून $s$ हा स्वेच्छ चल-विन्यास घ्या.
  $\Struct{M'}$ ही $\Struct{M}$ सारखीच रचना घ्या, फक्त
  $\Assign{a}{M'} = s(x)$ असेल. \readerexternalref{fol:syn:ext:cor:extensionality-sent}
  नुसार प्रत्येक $\chi \in \Gamma$ साठी $a$ हे $\Gamma$ मध्ये येत नसल्यामुळे
  $\Sat{M'}{\chi}$, आणि प्रत्येक $\chi \in \Delta$ साठी $\Sat/{M'}{\chi}$.
  पण आधारविधान वैध असल्याने $\Sat{M'}{\varphi(a)}$. \readerexternalref{fol:syn:ass:prop:sentence-sat-true}
  नुसार $\varphi(a)$ हे वाक्य असल्यामुळे $\Sat{M'}{\varphi(a)}[s]$.
  आता $\varAssign{s}{s}{x}$ आणि $s(x) = \Value{a}{M'}[s]$, कारण
  $\Struct{M'}$ याच प्रकारे परिभाषित केले आहे. म्हणून
  \readerexternalref{fol:syn:ext:prop:ext-formulas} लागू होते आणि $\Sat{M'}{\varphi(x)}[s]$
  मिळते. $a$ हे $\varphi(x)$ मध्ये येत नसल्याने
  \readerexternalref{fol:syn:ext:prop:extensionality} नुसार $\Sat{M}{\varphi(x)}[s]$.
  $s$ स्वेच्छ असल्याने सर्व चल-विन्यासांसाठी $\Sat{M}{\varphi(x)}[s]$ सिद्ध झाले.
\item शेवटचे अनुमान \LeftR{\lexists} आहे: सरावासाठी.

\end{enumerate}
आता दोन आधारविधाने असलेली संभाव्य अनुमाने पाहू.
\begin{enumerate}
\item शेवटचे अनुमान कट आहे; म्हणून $\pi$ पुढील रूपात संपते:
  \begin{prooftree}
    \AxiomC{}
    \Deduce$\Gamma \fCenter \Delta, \varphi$
    \AxiomC{}
    \Deduce$\varphi, \Pi \fCenter \Lambda$
    \RightLabel{\Cut}
    \BinaryInf$\Gamma, \Pi \fCenter \Delta, \Lambda$
  \end{prooftree}
  $\Struct{M}$ ही संरचना घ्या. विगमन गृहीतकानुसार आधारविधाने वैध असल्याने
  $\Struct{M}$ दोन्ही आधारविधाने पूर्ण करते.
  दोन प्रकरणे पाहू: (a) $\Sat/{M}{\varphi}$ आणि (b)
  $\Sat{M}{\varphi}$. (a) प्रकरणात डावे आधारविधान पूर्ण
  करण्यासाठी $\Struct{M}$ ने
  $\Gamma \Sequent \Delta$ पूर्ण केले पाहिजे; मग ते निष्कर्षही पूर्ण करते.
  (b) प्रकरणात उजवे आधारविधान पूर्ण करण्यासाठी
  $\Struct{M}$ ने $\Pi \setminus \Lambda$ पूर्ण
  केले पाहिजे. पुन्हा, $\Struct{M}$ निष्कर्ष पूर्ण करते.
\item शेवटचे अनुमान \RightR{\land} आहे. मग $\pi$ पुढील रूपात संपते:
  \begin{prooftree}
    \AxiomC{}
    \Deduce$\Gamma \fCenter \Delta, \varphi$
    \AxiomC{}
    \Deduce$\Gamma \fCenter \Delta, \psi$
    \RightLabel{\RightR{\land}}
    \BinaryInf$\Gamma \fCenter \Delta, \varphi \land \psi$
  \end{prooftree}
  संरचना~$\Struct
    M$ घ्या. $\Struct{M}$
  ने $\Gamma \Sequent \Delta$ पूर्ण केले तर काम झाले. तसे नसल्याचे गृहीत धरा.
  विगमन गृहीतकानुसार $\Gamma \fCenter
  \Delta, \varphi$ वैध असल्याने,
  $\Sat{M}{\varphi}$. त्याचप्रमाणे, कारण $\Gamma
  \Sequent \Delta, \psi$ वैध असल्याने,
  $\Sat{M}{\psi}$. म्हणून
  $\Sat{M}{\varphi \land \psi}$.
\item शेवटचे अनुमान \LeftR{\lor} आहे: सरावासाठी.
\item शेवटचे अनुमान \LeftR{\lif} आहे. मग $\pi$ पुढील रूपात संपते:
  \begin{prooftree}
    \AxiomC{}
    \Deduce$\Gamma \fCenter \Delta, \varphi$
    \AxiomC{}
    \Deduce$\psi, \Pi \fCenter \Lambda$
    \RightLabel{\LeftR{\lif}}
    \BinaryInf$\varphi \lif \psi, \Gamma, \Pi \fCenter \Delta, \Lambda$
  \end{prooftree}
  पुन्हा,
  संरचना~$\Struct{M}$
  $\Struct{M}$ ने $\Gamma, \Pi \Sequent
  \Delta, \Lambda$ पूर्ण केलेले नाही असे गृहीत धरा. आपल्याला
  $\Sat/{M}{\varphi \lif \psi}$ दाखवायचे आहे.
  $\Struct{M}$ ने $\Gamma,
  \Pi \Sequent \Delta, \Lambda$ पूर्ण केलेले नसल्याने ते $\Gamma \Sequent
  \Delta$ किंवा $\Pi \Sequent \Lambda$ यांपैकी कोणतेही पूर्ण करत नाही. $\Gamma
  \Sequent \Delta, \varphi$ वैध असल्याने $\Sat{M}{\varphi}$.
  $\psi, \Pi \Sequent \Lambda$ वैध असल्याने
  $\Sat/{M}{\psi}$. म्हणून
  $\Sat/{M}{\varphi \lif \psi}$, जे दाखवायचे होते तेच.
\end{enumerate}
\end{proof}


\begin{prob}
\ref{fol:seq:sou:thm:sequent-soundness} ची सिद्धता पूर्ण करा.
\end{prob}




\begin{cor}
\label{fol:seq:sou:cor:weak-soundness}
जर $\Proves \varphi$, तर $\varphi$ हे वैध आहे.
\end{cor}

\begin{cor}
\label{fol:seq:sou:cor:entailment-soundness}
जर $\Gamma \Proves \varphi$, तर $\Gamma \Entails \varphi$.
\end{cor}

\begin{proof}
जर $\Gamma \Proves \varphi$ असेल, तर $\Gamma_0 \subseteq \Gamma$ असा काही
मर्यादित संच आणि $\Gamma_0 \Sequent \varphi$ ची निष्पत्ती असते. \ref{fol:seq:sou:thm:sequent-soundness}
नुसार, प्रत्येक
संरचना~$\Struct{M}$
किंवा $\psi \in \Gamma_0$ असत्य करते, किंवा $\varphi$ सत्य करते. म्हणून,
जर $\Sat{M}{\Gamma}$ असेल, तर
$\Sat{M}{\varphi}$ देखील असेल.
\end{proof}

\begin{cor}
\label{fol:seq:sou:cor:consistency-soundness}
जर $\Gamma$ पूर्ण करण्यायोग्य असेल, तर तो सुसंगत आहे.
\end{cor}

\begin{proof}
प्रतिलोम विधान सिद्ध करू. $\Gamma$ सुसंगत नाही असे गृहीत धरा. मग
$\Gamma_0 \subseteq \Gamma$ असा काही मर्यादित संच आणि
$\Gamma_0 \Sequent \quad$ ची निष्पत्ती असते. \ref{fol:seq:sou:thm:sequent-soundness}
नुसार $\Gamma_0 \Sequent \quad$ वैध आहे. दुसऱ्या शब्दांत, प्रत्येक
संरचना~$\Struct{M}$ साठी
$\chi \in \Gamma_0$ असे आहे की
$\Sat/{M}{\chi}$; आणि $\Gamma_0 \subseteq \Gamma$
असल्याने ते $\chi$ $\Gamma$ मध्येही आहे. त्यामुळे कोणतेही
$\Struct{M}$ $\Gamma$ पूर्ण करत नाही, आणि
$\Gamma$ पूर्ण करण्यायोग्य नाही.
\end{proof}









\section{निष्पत्ती सह एकरूपता}\label{fol:seq:ide:sec}

निष्पत्त्या आणि एकरूपता साठी अतिरिक्त आरंभीच्या क्रमवर्ती
आणि अनुमाननियमांची आवश्यकता असते.

\begin{defn}[$\eq$ साठी आरंभीच्या क्रमवर्ती]
जर $t$ हे बंद पद असेल, तर ${} \Sequent \eq[t][t]$ ही आरंभीची क्रमवर्ती असते.
\end{defn}

$\eq$ साठीचे नियम पुढीलप्रमाणे आहेत ($t_1$ आणि $t_2$ ही बंद पदे आहेत):

\begin{defish}
\Axiom$ \eq[t_1][t_2], \Gamma \fCenter \Delta, \varphi(t_1) $
\RightLabel{$\eq$}
\UnaryInf$\eq[t_1][t_2], \Gamma \fCenter \Delta, \varphi(t_2)$
\DisplayProof
\hfill
\Axiom$\eq[t_1][t_2], \Gamma \fCenter \Delta, \varphi(t_2) $
\RightLabel{$\eq$}
\UnaryInf$\eq[t_1][t_2], \Gamma  \fCenter \Delta, \varphi(t_1)$
\DisplayProof
\end{defish}

\begin{ex}
जर $s$ आणि $t$ ही बंद पदे असतील, तर $\eq[s][t], \varphi(s)
\Proves \varphi(t)$:
\begin{prooftree}
\Axiom$ \varphi(s) \fCenter \varphi(s)$
\RightLabel{\LeftR{\Weakening}}
\UnaryInf$\eq[s][t], \varphi(s)  \fCenter \varphi(s)$
\RightLabel{$\eq$}
\UnaryInf$\eq[s][t], \varphi(s)  \fCenter \varphi(t)$
\end{prooftree}
याला एकरूप वस्तूंच्या प्रतिस्थापनाचे तत्त्व किंवा लाइब्निझचा नियम
म्हणून ओळखले जाऊ शकते.

$\Log{LK}$ सिद्ध करते की $\eq$ सममित आणि संक्रमक आहे:
\begin{prooftree}
\Axiom$ \fCenter \eq[t_1][t_1] $
\RightLabel{\LeftR{\Weakening}}
\UnaryInf$ \eq[t_1][t_2] \fCenter \eq[t_1][t_1] $
\RightLabel{$\eq$}
\UnaryInf$ \eq[t_1][t_2] \fCenter \eq[t_2][t_1]$
\DisplayProof\qquad\bottomAlignProof
\Axiom$ \eq[t_1][t_2] \fCenter \eq[t_1][t_2] $
\RightLabel{\LeftR{\Weakening}}
\UnaryInf$\eq[t_2][t_3], \eq[t_1][t_2]  \fCenter \eq[t_1][t_2] $
\RightLabel{$\eq$}
\UnaryInf$\eq[t_2][t_3], \eq[t_1][t_2]  \fCenter \eq[t_1][t_3]$
\RightLabel{\LeftR{\Exchange}}
\UnaryInf$\eq[t_1][t_2], \eq[t_2][t_3]  \fCenter \eq[t_1][t_3]$
\end{prooftree}
डावीकडील निष्पत्तीमध्ये सूत्र~$\eq[x][t_1]$ हे आपले
$\varphi(x)$ आहे. उजवीकडे $\varphi(x)$ हे~$\eq[t_1][x]$ असे घेतो.
\end{ex}

\begin{prob}
पुढील क्रमवर्तींच्या निष्पत्त्या द्या:
\begin{enumerate}
\item $\Sequent \lforall[x][\lforall[y][((x = y \land \varphi(x)) \lif \varphi(y))]]$
\item $\lexists[x][\varphi(x)] \land \lforall[y][\lforall[z][((\varphi(y) \land
    \varphi(z)) \lif y = z)]] \Sequent
\lexists[x][(\varphi(x) \land \lforall[y][(\varphi(y) \lif y = x)])]$
\end{enumerate}
\end{prob}









\section{एकरूपता सह निर्दोषता}\label{fol:seq:sid:sec}

\begin{prop}
$\Log{LK}$ मधील आरंभीच्या क्रमवर्ती आणि एकरूपतेच्या नियमांसहची पद्धत निर्दोष आहे.
\end{prop}

\begin{proof}
${} \Sequent \eq[t][t]$ या रूपातील आरंभीच्या क्रमवर्ती वैध आहेत, कारण
प्रत्येक संरचना~$\Struct M$ साठी $\Sat{M}{\eq[t][t]}$ असते. (लक्षात घ्या की
$t$ हे बंद पद आहे असे आपण गृहीत धरतो, म्हणजे त्यात कोणतेही चल नाही;
म्हणून चल-विन्यास अप्रासंगिक असतात.)

निष्पत्तीमधील शेवटचे अनुमान $=$ आहे असे गृहीत धरा. मग आधारविधान
$\eq[t_1][t_2], \Gamma \Sequent \Delta, \varphi(t_1)$ आणि निष्कर्ष
$\eq[t_1][t_2], \Gamma \Sequent \Delta, \varphi(t_2)$ असतो. संरचना~$\Struct M$
घ्या. निष्कर्ष वैध आहे हे दाखवायचे आहे; म्हणजे $\Sat{M}{\eq[t_1][t_2]}$ आणि
$\Sat{M}{\Gamma}$ असल्यास, एकतर $\Sat{M}{\chi}$ हे एखाद्या $\chi \in \Delta$ साठी
असते किंवा $\Sat{M}{\varphi(t_2)}$ असते.

विगमन गृहीतकानुसार आधारविधान वैध आहे. म्हणजे $\Sat{M}{\eq[t_1][t_2]}$ आणि
$\Sat{M}{\Gamma}$ असल्यास, एकतर (a) एखाद्या $\chi \in \Delta$ साठी
$\Sat{M}{\chi}$ किंवा (b) $\Sat{M}{\varphi(t_1)}$. (a) प्रकरणात काम झाले.
(b) प्रकरणाचा विचार करू. $s$ हा चल-विन्यास घ्या, ज्यासाठी
$s(x) = \Value{t_1}{M}$. \readerexternalref{fol:syn:ass:prop:sentence-sat-true} नुसार,
$\Sat{M}{\varphi(t_1)}[s]$. $\varAssign{s}{s}{x}$ असल्याने,
\readerexternalref{fol:syn:ext:prop:ext-formulas} नुसार $\Sat{M}{\varphi(x)}[s]$. कारण
$\Sat{M}{\eq[t_1][t_2]}$, आपल्याला $\Value{t_1}{M} = \Value{t_2}{M}$,
आणि म्हणून $s(x) = \Value{t_2}{M}$ मिळते. \readerexternalref{fol:syn:ext:prop:ext-formulas}
पुन्हा लागू केल्याने $\Sat{M}{\varphi(t_2)}[s]$ देखील मिळते.
\readerexternalref{fol:syn:ass:prop:sentence-sat-true} नुसार,
$\Sat{M}{\varphi(t_2)}$.
\end{proof}



\chapter{नैसर्गिक निगमन}\label{fol:ntd::chap}
\begin{editorial}
हे प्रकरण Gentzen/Prawitz शैलीतील नैसर्गिक निगमन पद्धत मांडते.

नैसर्गिक निगमनाशी सिद्धता-पद्धती म्हणून संबंधित सामग्री समाविष्ट किंवा
वगळण्यासाठी ``prfND'' हा टॅग वापरा.
\end{editorial}






\section{नियम आणि निष्पत्ती}\label{fol:ntd:rul:sec}

\begin{explain}
नैसर्गिक निगमन पद्धतींचा उद्देश गणितीय सिद्धतेत वापरल्या जाणाऱ्या अनौपचारिक
तर्कपद्धतीशी जवळचे साम्य राखणे हा आहे (म्हणूनच ती काहीशी ``नैसर्गिक'' आहे).
नैसर्गिक निगमनातील सिद्धता गृहीतकांपासून सुरू होतात.
त्यानंतर अनुमाननियम लागू केले जातात. \Intro{\lnot}, \Intro{\lif},
\Elim{\lor} आणि \Elim{\lexists} अनुमाननियमांद्वारे
गृहीतके ``मुक्त'' केली जातात; स्पष्टतेसाठी मुक्त केलेल्या
गृहीतकाचा लेबल त्या अनुमानाजवळ ठेवला जातो.
\end{explain}

\begin{defn}[गृहीतक]
\emph{गृहीतक} म्हणजे कोणत्याही शाखेच्या सर्वांत वरच्या स्थानावर असलेले
कोणतेही वाक्य होय.
\end{defn}

नैसर्गिक निगमनातील निष्पत्त्या ही वाक्यांची विशिष्ट वृक्षरचना असते;
त्यातील सर्वांत वरच्या वाक्ये गृहीतके असतात. एखादे वाक्य एक,
दोन किंवा तीन इतर क्रमवर्तींच्या खाली असल्यास, ते अनुमाननियमाने योग्य
प्रकारे आलेले असले पाहिजे. अनुमानाच्या शीर्षस्थानी असलेल्या वाक्यांना
\emph{आधारविधाने} म्हणतात आणि खालील वाक्याला अनुमानाचा
\emph{निष्कर्ष} म्हणतात. प्रत्येक संकारक साठी नियम जोड्यांमध्ये येतात:
एक प्रवेशन नियम आणि एक विलोपन नियम. ते निष्कर्षात संकारक आणतात
किंवा नियमाच्या आधारविधानातून संकारक काढतात. काही नियम विशिष्ट
प्रकारचे गृहीतक \emph{मुक्त} करण्याची परवानगी देतात. कोणते
गृहीतक कोणत्या अनुमानाने मुक्त केले हे दर्शवण्यासाठी गृहीतक आणि
अनुमान दोन्हींना लेबल देतो. हे गृहीतक ``$\Discharge{\varphi}{n}$.` असे लिहून दर्शवले जाते.


सर्व संकारक $\land$, $\lor$, $\lif$, $\lnot$ आणि $\lfalse$ यांच्यासाठीचे नियम विचारात घेणे प्रचलित आहे, जरी त्यांपैकी काही परिभाषित केलेले असले तरी.











\section{विधानीय नियम}\label{fol:ntd:prl:sec}


\subsection{$\land$ साठीचे नियम}

\begin{defish}
\AxiomC{$\varphi$}
\AxiomC{$\psi$}
\RightLabel{\Intro{\land}}
\BinaryInfC{$\varphi \land \psi$}
\DisplayProof
\hfill
\begin{tabular}{r}
\AxiomC{$\varphi \land \psi$}
\RightLabel{\Elim{\land}}
\UnaryInfC{$\varphi$}
\DisplayProof
\\[3ex]
\AxiomC{$\varphi \land \psi$}
\RightLabel{\Elim{\land}}
\UnaryInfC{$\psi$}
\DisplayProof
\end{tabular}
\end{defish}

\subsection{$\lor$ साठीचे नियम}

\begin{defish}
\begin{tabular}{r}
\AxiomC{$\varphi$}
\RightLabel{\Intro{\lor}}
\UnaryInfC{$\varphi \lor \psi$}
\DisplayProof
\\[3ex]
\AxiomC{$\psi$}
\RightLabel{\Intro{\lor}}
\UnaryInfC{$\varphi \lor \psi$}
\DisplayProof
\end{tabular}
\hfill
\AxiomC{$\varphi \lor \psi$}
\AxiomC{$\Discharge{\varphi}{n}$}
\DeduceC{$\chi$}
\AxiomC{$\Discharge{\psi}{n}$}
\DeduceC{$\chi$}
\DischargeRule{\Elim{\lor}}{n}
\TrinaryInfC{$\chi$}
\DisplayProof
\end{defish}

\subsection{$\lif$ साठीचे नियम}

\begin{defish}
\AxiomC{$\Discharge{\varphi}{n}$}
\DeduceC{$\psi$}
\DischargeRule{\Intro{\lif}}{n}
\UnaryInfC{$\varphi \lif \psi$}
\DisplayProof
\hfill
\AxiomC{$\varphi \lif \psi$}
\AxiomC{$\varphi$}
\RightLabel{\Elim{\lif}}
\BinaryInfC{$\psi$}
\DisplayProof
\end{defish}

\subsection{$\lnot$ साठीचे नियम}

\begin{defish}
\AxiomC{$\Discharge{\varphi}{n}$}
\noLine
\DeduceC{$\lfalse$}
\DischargeRule{\Intro{\lnot}}{n}
\UnaryInfC{$\lnot \varphi$}
\DisplayProof
\hfill
\AxiomC{$\lnot \varphi$}
\AxiomC{$\varphi$}
\RightLabel{\Elim{\lnot}}
\BinaryInfC{$\lfalse$}
\DisplayProof
\end{defish}

\subsection{$\lfalse$ साठीचे नियम}

\begin{defish}
\AxiomC{$\lfalse$}
\RightLabel{\FalseInt}
\UnaryInfC{$\varphi$}
\DisplayProof
\hfill
\AxiomC{$\Discharge{\lnot \varphi}{n}$}
\DeduceC{$\lfalse$}
\DischargeRule{\FalseCl}{n}
\UnaryInfC{$\varphi$}
\DisplayProof
\end{defish}

लक्षात घ्या की $\Intro{\lnot}$ आणि $\FalseCl$ हे अगदी सारखे आहेत: फरक
असा की $\Intro{\lnot}$ हा नकारयुक्त वाक्य~$\lnot \varphi$ निष्पन्न
करतो, तर $\FalseCl$ हा सकारात्मक वाक्य~$\varphi$ निष्पन्न करतो.

जेव्हा एखादा नियम एखादे गृहीतक मुक्त करता येते असे दर्शवतो, तेव्हा आपण ती
परवानगी मानतो, आवश्यकता नाही. उदा., $\Intro{\lif}$ नियमात आधारविधान~$\psi$
च्या निष्पत्तीमधील $\varphi$ या रूपाची कितीही गृहीतके—शून्यसुद्धा—आपण
मुक्त करू शकतो.









\section{संख्यापकांचे नियम}\label{fol:ntd:qrl:sec}

\subsection{$\lforall$ साठीचे नियम}

\begin{defish}
\AxiomC{$\varphi(a)$}
\RightLabel{\Intro{\lforall}}
\UnaryInfC{$\lforall[x][\Atom{\varphi}{x}]$}
\DisplayProof
\hfill
\AxiomC{$\lforall[x][\Atom{\varphi}{x}]$}
\RightLabel{\Elim{\lforall}}
\UnaryInfC{$\varphi(t)$}
\DisplayProof
\end{defish}

$\lforall$ च्या नियमांमध्ये, $t$ हे बंद पद आहे (म्हणजे, कोणतेही चर
नसलेले पद), आणि $a$~हा असा स्थिरांक आहे की तो
निष्कर्ष~$\lforall[x][\varphi(x)]$ मध्ये किंवा आधारविधान~$\varphi(a)$ ने समाप्त
होणाऱ्या निष्पत्तीमध्ये मुक्त न केलेली असलेल्या कोणत्याही
गृहीतकात येत नाही. आपण $a$ ला \Intro{\lforall} अनुमानाचा
\emph{आयगेन चर} म्हणतो.\footnote{वरील नियमात $a$ हा अचर असला तरी आपण
``आयगेन चर'' ही संज्ञा वापरतो. यामागे ऐतिहासिक कारणे आहेत.}

\subsection{$\lexists$ साठीचे नियम}

\begin{defish}
\AxiomC{$\Atom{\varphi}{t}$}
\RightLabel{\Intro{\lexists}}
\UnaryInfC{$\lexists[x][\Atom{\varphi}{x}]$}
\DisplayProof
\hfill
\AxiomC{$\lexists[x][\Atom{\varphi}{x}]$}
\AxiomC{[$\Atom{\varphi}{a}$]$^n$}
\DeduceC{$\chi$}
\DischargeRule{\Elim{\lexists}}{n}
\BinaryInfC{$\chi$}
\DisplayProof
\end{defish}

पुन्हा, $t$ हे बंद पद आहे आणि $a$ हा असा स्थिरांक आहे की तो
आधारविधान $\lexists[x][\varphi(x)]$ मध्ये, निष्कर्ष~$\chi$ मध्ये किंवा दोन
आधारविधानांनी समाप्त होणाऱ्या निष्पत्त्या मधील कोणत्याही
मुक्त न केलेली गृहीतकात ($\varphi(a)$ ही गृहीतके वगळता) येत नाही. आपण
$a$ ला \Elim{\lexists} अनुमानाचा \emph{आयगेन चर} म्हणतो.

\Intro{\lforall} आणि \Elim{\lexists} अनुमानांसाठी आयगेन चरावर लागू होणारे
निर्बंध वर स्वतंत्रपणे सांगितले आहेत. त्यांत संबंधित आधारविधान व निष्कर्ष,
तसेच आधारविधानांकडे नेणाऱ्या निष्पत्त्या मधील मुक्त न केलेली गृहीतके,
यांबाबत वेगवेगळे निर्बंध आहेत. या निर्बंधांना \emph{आयगेन चराची अट} म्हणतात.

\begin{explain}
ही प्रथा आठवा: $\varphi$ हे असे सूत्र असेल की त्यात
चल~$x$ मुक्त असेल, तर हे आपण~$\varphi(x)$ असे लिहून दर्शवतो.
त्याच संदर्भात, मग $\varphi(t)$ हा~$\Subst{\varphi}{t}{x}$ चा संक्षेप असतो.
त्यामुळे $\Intro\lexists$ नियम पुढीलप्रमाणेही लिहिता येईल:
\begin{prooftree}
\AxiomC{$\Subst{\varphi}{t}{x}$}
\RightLabel{\Intro{\lexists}}
\UnaryInfC{$\lexists[x][\varphi]$}
\end{prooftree}
लक्षात घ्या की~$\varphi$ मध्ये $t$ आधीपासून येऊ शकते; उदा., $\varphi$ हे
$\Atom{\Obj P}{t,x}$ असू शकते. म्हणून,~$\Atom{\Obj P}{t,t}$ पासून
$\lexists[x][\Atom{\Obj P}{t,x}]$ निष्पन्न करणे हे~$\Intro\lexists$
चे योग्य उपयोजन आहे---आधारविधानात~$t$ जेथे येते त्यांपैकी एक किंवा
अधिक ठिकाणी त्याऐवजी बद्ध चल~$x$ घालता येतो; सर्वच ठिकाणी
असा ``बदल'' करणे आवश्यक नाही. परंतु $\Intro\lforall$ आणि
$\Elim\lexists$ मधील आयगेन चराच्या अटीनुसार स्थिरांक~$a$ हा~$\varphi$
मध्ये येता कामा नये. म्हणून, $\Intro\lforall$ वापरून
$\Atom{\Obj P}{a,a}$ पासून $\lforall[x][\Atom{\Obj P}{a,x}]$ योग्य
रीतीने निष्पन्न करता येत नाही.
\end{explain}

\begin{explain}
\Intro{\lexists} आणि \Elim{\lforall} मध्ये पद~$t$ बंद
असण्याव्यतिरिक्त त्यावर आणखी कोणतेही निर्बंध नाहीत; म्हणून इतर कोणत्याही
अटीची काळजी करावी लागत नाही. याउलट, \Elim{\lexists} आणि
\Intro{\lforall} नियमांमध्ये आयगेन चराच्या अटीनुसार स्थिरांक~$a$ हा
निष्कर्षात किंवा कोणत्याही मुक्त न केलेली गृहीतकात कोठेही येता कामा नये.
ही पद्धत निर्दोष आहे, म्हणजे ती केवळ अशा मुक्त न केलेली गृहीतकांपासूनच
वाक्ये निष्पन्न करते की ज्यांच्यापासून ती तार्किकरीत्या निष्पन्न
होतात, याची खात्री करण्यासाठी ही अट आवश्यक आहे. ही अट नसती, तर पुढील
निष्पत्तीला परवानगी मिळाली असती:
\begin{prooftree}
  \AxiomC{$\lexists[x][\varphi(x)]$}
  \AxiomC{$\Discharge{\varphi(a)}{1}$}
  \RightLabel{*\Intro{\lforall}}
  \UnaryInfC{$\lforall[x][\varphi(x)]$}
  \RightLabel{\Elim{\lexists}}
  \BinaryInfC{$\lforall[x][\varphi(x)]$}
\end{prooftree}
परंतु, $\lexists[x][\varphi(x)] \Entails/ \lforall[x][\varphi(x)]$.

संख्यापकांच्या विलोपन नियमांमध्ये चलांच्या जागी केवळ बंद पदे
ठेवण्याची परवानगी असल्यामुळे, वाक्यांच्या संचापासून निष्पन्न करता
येणारे कोणतेही सूत्र हे स्वतः वाक्य असते, हे निष्पन्न होते.
\end{explain}










\section{निष्पत्ती}\label{fol:ntd:der:sec}

\begin{explain}
गृहीतक म्हणजे काय हे आपण सांगितले आहे आणि निगमन नियम दिले आहेत.
नैसर्गिक निगमनातील निष्पत्त्या यांपासून आगमनाने निर्माण होतात:
प्रत्येक निष्पत्ती ही एकतर स्वतंत्र गृहीतक असते किंवा एक, दोन वा तीन
निष्पत्त्या नंतर केलेले योग्य अनुमान असते.
\end{explain}

\begin{defn}[निष्पत्ती]
\emph{निष्पत्ती} म्हणजे गृहीतके~$\Gamma$
यांपासून वाक्य~$\varphi$ ची पुढील अटी पूर्ण करणारी वाक्यांची
सांत वृक्षरचना होय:
\begin{enumerate}
\item वृक्षातील सर्वांत वरची वाक्ये ही एकतर $\Gamma$ मध्ये असतात
  किंवा वृक्षातील एखाद्या अनुमानाने मुक्त केलेली असतात.
\item वृक्षातील सर्वांत खालचे वाक्य हे~$\varphi$ असते.
\item वृक्षाच्या तळाशी असलेले वाक्य~$\varphi$ वगळता, वृक्षातील प्रत्येक
  वाक्य हे एखाद्या निगमन नियमाच्या योग्य उपयोजनाचे आधारविधान असते
  आणि त्या उपयोजनाचा निष्कर्ष वृक्षात त्या वाक्याच्या थेट खाली येतो.
\end{enumerate}
मग $\varphi$ हा त्या निष्पत्तीचा \emph{निष्कर्ष} आणि $\Gamma$ ही तिची
मुक्त न केलेली गृहीतके आहेत, असे आपण म्हणतो.

जर~$\Gamma$ या गृहीतकांपासून $\varphi$ ची निष्पत्ती अस्तित्वात असेल, तर
$\varphi$ हे~$\Gamma$ पासून \emph{निष्पन्न करता येण्याजोगे} आहे असे म्हणतो; चिन्हांत:
$\Gamma \Proves \varphi$. प्रत्येक गृहीतक मुक्त केलेले असणारी $\varphi$
ची निष्पत्ती असेल, तर आपण~$\Proves \varphi$ असे लिहितो.
\end{defn}

\begin{ex}
प्रत्येक गृहीतक स्वतःच निष्पत्ती असते. म्हणून, उदा., एकटे $\varphi$ ही
निष्पत्ती आहे आणि एकटे $\psi$ सुद्धा. यांवर, समजा,
$\Intro{\land}$ नियम लागू करून नवी निष्पत्ती मिळवता येते:
\begin{prooftree}
\AxiomC{$\varphi$}
\AxiomC{$\psi$}
\RightLabel{\Intro{\land}}
\BinaryInfC{$\varphi \land \psi$}
\end{prooftree}
हे नियम सर्वसाधारण असतात: त्यांतील $\varphi$ आणि~$\psi$ यांच्या जागी कोणतीही
वाक्ये, उदा., $\chi$ आणि~$\theta$, ठेवता येतात. मग निष्कर्ष
$\chi \land \theta$ असेल, आणि म्हणून
\begin{prooftree}
\AxiomC{$\chi$}
\AxiomC{$\theta$}
\RightLabel{\Intro{\land}}
\BinaryInfC{$\chi \land \theta$}
\end{prooftree}
ही योग्य निष्पत्ती आहे. अर्थात, गृहीतकांची अदलाबदल करून $\theta$ ला
$\varphi$ ची आणि $\chi$ ला~$\psi$ ची भूमिका देता येते. त्यामुळे,
\begin{prooftree}
\AxiomC{$\theta$}
\AxiomC{$\chi$}
\RightLabel{\Intro{\land}}
\BinaryInfC{$\theta \land \chi$}
\end{prooftree}
हीसुद्धा योग्य निष्पत्ती आहे.

आता आपण दुसरा नियम, समजा, $\Intro{\lif}$ लागू करू शकतो. त्यामुळे सोपाधिक
विधानाचा निष्कर्ष काढता येतो आणि त्या विधानाच्या पूर्वांगाशी एकरूप असलेले
कोणतेही गृहीतक मुक्त करता येते. म्हणून पुढील दोन्या निष्पत्त्या योग्य आहेत:
\begin{prooftree}
\AxiomC{$\Discharge{\chi}{1}$}
\AxiomC{$\theta$}
\RightLabel{\Intro{\land}}
\BinaryInfC{$\chi \land \theta$}
\DischargeRule{\Intro{\lif}}{1}
\UnaryInfC{$\chi \lif (\chi \land \theta)$}
\DisplayProof\bottomAlignProof
\AxiomC{$\chi$}
\AxiomC{$\Discharge{\theta}{1}$}
\RightLabel{\Intro{\land}}
\BinaryInfC{$\chi \land \theta$}
\DischargeRule{\Intro{\lif}}{1}
\UnaryInfC{$\theta \lif (\chi \land \theta)$}
\end{prooftree}
त्या अनुक्रमे $\theta \Proves \chi \lif (\chi \land \theta)$ आणि
$\chi \Proves \theta \lif (\chi \land \theta)$ हे दाखवतात.

गृहीतके मुक्त करणे ही परवानगी आहे, आवश्यकता नाही, हे आठवा: गृहीतके मुक्त
करणे बंधनकारक नसते. विशेषतः, निष्पत्तीमध्ये ती गृहीतके नसली तरीही
नियम लागू करता येतो. उदाहरणार्थ, मुक्त करण्यासाठी~$\varphi$ हे गृहीतक
नसले तरी पुढील उपयोजन अनुज्ञात आहे:
\begin{prooftree}
  \AxiomC{$\psi$}
  \DischargeRule{\Intro{\lif}}{1}
  \UnaryInfC{$\varphi \lif \psi$}
\end{prooftree}
\end{ex}









\section{निष्पत्ती: उदाहरणे}\label{fol:ntd:pro:sec}

\begin{ex}
वाक्य $(\varphi \land \psi) \lif \varphi$ ची निष्पत्ती देऊ.

हवा असलेला निष्कर्ष निष्पत्तीच्या तळाशी लिहून आपण सुरुवात करतो.
\begin{prooftree}
\AxiomC{}
\UnaryInfC{$(\varphi\land \psi) \lif \varphi$}
\end{prooftree}

पुढे, या रूपाचे वाक्य कोणत्या प्रकारच्या अनुमानामुळे मिळू शकते हे
शोधावे लागते. निष्कर्षाचा मुख्य संकारक हा $\lif$ आहे; म्हणून
\Intro{\lif} नियम वापरून निष्कर्षापर्यंत पोहोचण्याचा प्रयत्न करू. काम पुढे
जाताना संबंधित गृहीतके लिहून ठेवणे आणि निगमन नियमांना लेबले देणे उत्तम;
त्यामुळे सिद्धतेच्या शेवटी सर्व गृहीतके मुक्त झाली आहेत का हे सहज
दिसते.
\begin{prooftree}
\AxiomC{$\Discharge{\varphi \land \psi}{1}$}
\DeduceC{$\varphi$}
\DischargeRule{\Intro{\lif}}{1}
\UnaryInfC{$(\varphi\land \psi) \lif \varphi$}
\end{prooftree}

आता गृहीतक $\varphi \land \psi$ पासून $\varphi$ पर्यंतच्या पायऱ्या भराव्या लागतील.
येथे हाताळण्यासाठी फक्त एकच संयोजक, $\land$, असल्यामुळे $\land$ चा विलोपन
नियम वापरलाच पाहिजे. त्यामुळे पुढील सिद्धता मिळते:
\begin{prooftree}
\AxiomC{$\Discharge{\varphi \land \psi}{1}$}
\RightLabel{\Elim{\land}}
\UnaryInfC{$\varphi$}
\DischargeRule{\Intro{\lif}}{1}
\UnaryInfC{$(\varphi\land \psi) \lif \varphi$}
\end{prooftree}
आता आपल्याकडे $(\varphi \land \psi) \lif \varphi$ ची योग्य निष्पत्ती आहे.
\end{ex}

\begin{ex}
आता $(\lnot \varphi \lor \psi) \lif (\varphi \lif \psi)$ ची निष्पत्ती देऊ.

हवा असलेला निष्कर्ष निष्पत्तीच्या तळाशी लिहून आपण सुरुवात करतो.
\begin{prooftree}
\AxiomC{}
\UnaryInfC{$(\lnot \varphi \lor \psi) \lif (\varphi \lif \psi)$}
\end{prooftree}
हा निष्कर्ष देऊ शकणारा तार्किक नियम शोधण्यासाठी निष्कर्षातील $\lnot$,
$\lor$ आणि $\lif$ हे तार्किक संयोजक पाहतो. सध्या $\lif$ ची पहिली आस्थितीच
महत्त्वाची आहे, कारण अंतिम क्रमवर्तीतील वाक्याचा तो मुख्य संकारक
आहे; $\lnot$, $\lor$ आणि $\lif$ ची दुसरी आस्थिती दुसऱ्या संयोजकाच्या
व्याप्तीत असल्यामुळे त्यांचा विचार नंतर करू. म्हणून आपण \Intro{\lif}
नियमाने सुरुवात करतो. त्याचे योग्य उपयोजन पुढीलप्रमाणे असले पाहिजे:
\begin{prooftree}
\AxiomC{$\Discharge{\lnot \varphi \lor \psi}{1}$}
\DeduceC{$\varphi \lif \psi$}
\DischargeRule{\Intro{\lif}}{1}
\UnaryInfC{$(\lnot \varphi \lor \psi) \lif (\varphi \lif \psi)$}
\end{prooftree}
पुढे जाण्यासाठी आता दोन पर्याय आहेत. एकतर तळापासून वरच्या दिशेने काम
चालू ठेवून \Intro{\lif} नियमाचे आणखी एक उपयोजन शोधता येईल, किंवा वरून
खाली काम करून \Elim{\lor} नियम लागू करता येईल. दुसरा पर्याय घेऊ.
\Elim{\lor} चे सर्वांत डावे आधारविधान म्हणून $\lnot \varphi \lor \psi$ हे
गृहीतक वापरू. \Elim{\lor} च्या योग्य उपयोजनासाठी इतर दोन्ही आधारविधाने
$\varphi \lif \psi$ या निष्कर्षाशी एकरूप असली पाहिजेत; मात्र प्रत्येकाला
अनुक्रमे आणखी एका गृहीतकापासून, म्हणजे $\lnot \varphi \lor \psi$ च्या दोन
विकल्पघटकांपैकी एकापासून, निष्पन्न करता येते. त्यामुळे आपली
निष्पत्ती पुढीलप्रमाणे दिसेल:
\begin{prooftree}
\AxiomC{$\Discharge{\lnot \varphi \lor \psi}{1}$}
\AxiomC{$\Discharge{\lnot \varphi}{2}$}
\DeduceC{$\varphi \lif \psi$}
\AxiomC{$\Discharge{\psi}{2}$}
\DeduceC{$\varphi \lif \psi$}
\DischargeRule{\Elim{\lor}}{2}
\TrinaryInfC{$\varphi \lif \psi$}
\DischargeRule{\Intro{\lif}}{1}
\UnaryInfC{$(\lnot \varphi \lor \psi) \lif (\varphi \lif \psi)$}
\end{prooftree}

उजवीकडील दोन्ही शाखांत आपल्याला $\varphi \lif \psi$ हे निष्पन्न करायचे आहे;
त्यासाठी \Intro{\lif} वापरणे सर्वांत योग्य आहे.
\begin{prooftree}
\AxiomC{$\Discharge{\lnot \varphi \lor \psi}{1}$}
\AxiomC{$\Discharge{\lnot \varphi}{2}, \Discharge{\varphi}{3}$}
\DeduceC{$\psi$}
\DischargeRule{\Intro{\lif}}{3}
\UnaryInfC{$\varphi \lif \psi$}
\AxiomC{$\Discharge{\psi}{2}, \Discharge{\varphi}{4}$}
\DeduceC{$\psi$}
\DischargeRule{\Intro{\lif}}{4}
\UnaryInfC{$\varphi \lif \psi$}
\DischargeRule{\Elim{\lor}}{2}
\TrinaryInfC{$\varphi \lif \psi$}
\DischargeRule{\Intro{\lif}}{1}
\UnaryInfC{$(\lnot \varphi \lor \psi) \lif (\varphi \lif \psi)$}
\end{prooftree}

निष्पत्तीच्या उरलेल्या दोन भागांसाठी, मधल्या भागात $\lnot \varphi$ आणि
$\varphi$ यांपासून, तर उजवीकडील भागात $\varphi$ आणि $\psi$ यांपासून, $\psi$ मिळवणाऱ्या
निष्पत्त्या हवी आहेत. यांपैकी पहिली आधी घेऊ. $\lnot \varphi$ आणि $\varphi$ ही
\Elim{\lnot} ची दोन आधारविधाने आहेत:
\begin{prooftree}
\AxiomC{$\Discharge{\lnot \varphi}{2}$}
\AxiomC{$\Discharge{\varphi}{3}$}
\RightLabel{\Elim{\lnot}}
\BinaryInfC{$\lfalse$}
\DeduceC{$\psi$}
\end{prooftree}
\FalseInt वापरून $\psi$ हा निष्कर्ष मिळवता येतो आणि शाखा पूर्ण करता येते.
\begin{prooftree}
\AxiomC{$\Discharge{\lnot \varphi \lor \psi}{1}$}
\AxiomC{$\Discharge{\lnot \varphi}{2}$}
\AxiomC{$\Discharge{\varphi}{3}$}
\RightLabel{\Intro{\lfalse}}
\BinaryInfC{$\lfalse$}
\RightLabel{\FalseInt}
\UnaryInfC{$\psi$}
\DischargeRule{\Intro{\lif}}{3}
\UnaryInfC{$\varphi \lif \psi$}
\AxiomC{$\Discharge{\psi}{2}, \Discharge{\varphi}{4}$}
\DeduceC{$\psi$}
\DischargeRule{\Intro{\lif}}{4}
\UnaryInfC{$\varphi \lif \psi$}
\DischargeRule{\Elim{\lor}}{2}
\TrinaryInfC{$\varphi \lif \psi$}
\DischargeRule{\Intro{\lif}}{1}
\UnaryInfC{$(\lnot \varphi \lor \psi) \lif (\varphi \lif \psi)$}
\end{prooftree}

आता सर्वांत उजवी शाखा पाहू. येथे हे लक्षात घेणे महत्त्वाचे आहे की
निष्पत्तीची व्याख्या
\emph{गृहीतके मुक्त करण्याची परवानगी देते}, पण
\emph{तसे करणे आवश्यक ठरवत नाही}. दुसऱ्या शब्दांत, $\varphi$ आणि $\psi$ या
गृहीतकांपैकी एकाचा उपयोग न करता दुसऱ्यापासून $\psi$ निष्पन्न करता येत असेल,
तर ते योग्य आहे. $\psi$ पासून $\psi$ हे निष्पन्न करणे तर क्षुल्लक आहे:
एकटे $\psi$ ही अशी निष्पत्ती असून कोणत्याही अनुमानाची गरज नाही. म्हणून
गृहीतक~$\varphi$ सरळ वगळता येते.
\begin{prooftree}
\AxiomC{$\Discharge{\lnot \varphi \lor \psi}{1}$}
\AxiomC{$\Discharge{\lnot \varphi}{2}$}
\AxiomC{$\Discharge{\varphi}{3}$}
\RightLabel{\Elim{\lnot}}
\BinaryInfC{$\lfalse$}
\RightLabel{\FalseInt}
\UnaryInfC{$\psi$}
\DischargeRule{\Intro{\lif}}{3}
\UnaryInfC{$\varphi \lif \psi$}
\AxiomC{$\Discharge{\psi}{2}$}
\RightLabel{\Intro{\lif}}
\UnaryInfC{$\varphi \lif \psi$}
\DischargeRule{\Elim{\lor}}{2}
\TrinaryInfC{$\varphi \lif \psi$}
\DischargeRule{\Intro{\lif}}{1}
\UnaryInfC{$(\lnot \varphi \lor \psi) \lif (\varphi \lif \psi)$}
\end{prooftree}
लक्षात घ्या की पूर्ण झालेल्या निष्पत्तीमध्ये सर्वांत उजवीकडील
\Intro{\lif} अनुमान प्रत्यक्षात कोणतेही गृहीतक मुक्त करत नाही.
\end{ex}

\begin{ex}
आतापर्यंत आपल्याला \FalseCl{} नियमाची गरज पडलेली नाही. हा नियम विशेष आहे,
कारण नियमाच्या निष्कर्षाचे उप-सूत्र नसलेले गृहीतक तो मुक्त करू देतो.
त्याचा \FalseInt{} नियमाशी जवळचा संबंध आहे. खरे तर \FalseInt{} नियम हे
\FalseCl{} नियमाचे विशेष प्रकरण आहे---``अंतःप्रज्ञावादी तर्कशास्त्र''
नावाचे एक तर्कशास्त्र आहे आणि त्यात केवळ \FalseInt{} अनुज्ञात आहे. इतर
कोणताही मार्ग चालत नसताना \FalseCl{} नियम हा अखेरचा उपाय असतो. उदाहरणार्थ,
आपल्याला $\varphi \lor \lnot \varphi$ हे निष्पन्न करायचे आहे असे समजा. नेहमीच्या
पद्धतीनुसार $\Intro{\lor}$ वापरून $\varphi \lor \lnot \varphi$ हे निष्पन्न
करण्याचा प्रयत्न करू. पण त्यासाठी कोणत्याही गृहीतकांशिवाय $\varphi$ किंवा
$\lnot \varphi$ यांपैकी एक निष्पन्न करावे लागेल, आणि ते शक्य नाही. अशा वेळी
\FalseCl{} उपयोगी पडतो!
\begin{prooftree}
  \AxiomC{$\Discharge{\lnot(\varphi \lor \lnot \varphi)}{1}$}
  \DeduceC{$\lfalse$}
  \DischargeRule{\FalseCl}{1}
  \UnaryInfC{$\varphi \lor \lnot \varphi$}
\end{prooftree}
आता $\lnot(\varphi \lor \lnot \varphi)$ पासून $\lfalse$ ची निष्पत्ती शोधायची
आहे. $\lfalse$ हा $\Elim{\lnot}$ चा निष्कर्ष असल्यामुळे तो नियम वापरून
पाहू:
\begin{prooftree}
  \AxiomC{$\Discharge{\lnot(\varphi \lor \lnot \varphi)}{1}$}
  \DeduceC{$\lnot \varphi$}
  \AxiomC{$\Discharge{\lnot(\varphi \lor \lnot \varphi)}{1}$}
  \DeduceC{$\varphi$}
  \RightLabel{\Elim{\lnot}}
  \BinaryInfC{$\lfalse$}
  \DischargeRule{\FalseCl}{1}
  \UnaryInfC{$\varphi \lor \lnot \varphi$}
\end{prooftree}
$\lnot \varphi$ ची निष्पत्ती शोधण्याच्या आपल्या पद्धतीनुसार
$\Intro{\lnot}$ चे उपयोजन करावे लागेल:
\begin{prooftree}
  \AxiomC{$\Discharge{\lnot(\varphi \lor \lnot \varphi)}{1}, \Discharge{\varphi}{2}$}
  \DeduceC{$\lfalse$}
  \DischargeRule{\Intro{\lnot}}{2}
  \UnaryInfC{$\lnot \varphi$}
  \AxiomC{$\Discharge{\lnot(\varphi \lor \lnot \varphi)}{1}$}
  \DeduceC{$\varphi$}
  \RightLabel{\Elim{\lnot}}
  \BinaryInfC{$\lfalse$}
  \DischargeRule{\FalseCl}{1}
  \UnaryInfC{$\varphi \lor \lnot \varphi$}
\end{prooftree}
येथे $\lnot(\varphi \lor \lnot \varphi)$ हे गृहीतक आणि आपल्या नव्या $\varphi$
गृहीतकापासून~$\Intro{\lor}$ ने निष्पन्न होणारे $\varphi \lor \lnot \varphi$
यांना $\Elim{\lnot}$ लावून $\lfalse$ सहज मिळवता येते:
\begin{prooftree}
  \AxiomC{$\Discharge{\lnot(\varphi \lor \lnot \varphi)}{1}$}
  \AxiomC{$\Discharge{\varphi}{2}$}
  \RightLabel{\Intro{\lor}}
  \UnaryInfC{$\varphi \lor \lnot \varphi$}
  \RightLabel{\Elim{\lnot}}
  \BinaryInfC{$\lfalse$}
  \DischargeRule{\Intro{\lnot}}{2}
  \UnaryInfC{$\lnot \varphi$}
  \AxiomC{$\Discharge{\lnot(\varphi \lor \lnot \varphi)}{1}$}
  \DeduceC{$\varphi$}
  \RightLabel{\Elim{\lnot}}
  \BinaryInfC{$\lfalse$}
  \DischargeRule{\FalseCl}{1}
  \UnaryInfC{$\varphi \lor \lnot \varphi$}
\end{prooftree}
उजव्या बाजूला हीच पद्धत वापरतो; फक्त तेथे $\varphi$ हे~\FalseCl ने मिळवतो:
\begin{prooftree}
  \AxiomC{$\Discharge{\lnot(\varphi \lor \lnot \varphi)}{1}$}
  \AxiomC{$\Discharge{\varphi}{2}$}
  \RightLabel{\Intro{\lor}}
  \UnaryInfC{$\varphi \lor \lnot \varphi$}
  \RightLabel{\Elim{\lnot}}
  \BinaryInfC{$\lfalse$}
  \DischargeRule{\Intro{\lnot}}{2}
  \UnaryInfC{$\lnot \varphi$}
  \AxiomC{$\Discharge{\lnot(\varphi \lor \lnot \varphi)}{1}$}
  \AxiomC{$\Discharge{\lnot \varphi}{3}$}
  \RightLabel{\Intro{\lor}}
  \UnaryInfC{$\varphi \lor \lnot \varphi$}
  \RightLabel{\Elim{\lnot}}
  \BinaryInfC{$\lfalse$}
  \DischargeRule{\FalseCl}{3}
  \UnaryInfC{$\varphi$}
  \RightLabel{\Elim{\lnot}}
  \BinaryInfC{$\lfalse$}
  \DischargeRule{\FalseCl}{1}
  \UnaryInfC{$\varphi \lor \lnot \varphi$}
\end{prooftree}
\end{ex}

\begin{prob}
पुढील गोष्टी दाखवणाऱ्या निष्पत्त्या द्या:
\begin{enumerate}
\item $\varphi \land (\psi \land \chi) \Proves (\varphi \land \psi) \land \chi$.
\item $\varphi \lor (\psi \lor \chi) \Proves (\varphi \lor \psi) \lor \chi$.
\item $\varphi \lif (\psi \lif \chi) \Proves \psi \lif (\varphi \lif \chi)$.
\item $\varphi \Proves \lnot\lnot \varphi$.
\end{enumerate}
\end{prob}

\begin{prob}
पुढील गोष्टी दाखवणाऱ्या निष्पत्त्या द्या:
\begin{enumerate}
\item $(\varphi \lor \psi) \lif \chi \Proves \varphi \lif \chi$.
\item $(\varphi \lif \chi) \land (\psi \lif \chi) \Proves (\varphi \lor \psi) \lif \chi$.
\item $\Proves \lnot(\varphi \land \lnot \varphi)$.
\item $\psi \lif \varphi \Proves \lnot \varphi \lif \lnot \psi$.
\item $\Proves (\varphi \lif \lnot \varphi) \lif \lnot \varphi$.
\item $\Proves \lnot(\varphi \lif \psi) \lif \lnot \psi$.
\item $\varphi \lif \chi \Proves \lnot (\varphi \land \lnot \chi)$.
\item $\varphi \land \lnot \chi \Proves \lnot (\varphi \lif \chi)$.
\item $\varphi \lor \psi, \lnot \psi \Proves \varphi$.
\item $\lnot \varphi \lor \lnot \psi \Proves \lnot(\varphi \land \psi)$.
\item $\Proves (\lnot \varphi \land \lnot \psi) \lif\lnot(\varphi \lor \psi)$.
\item $\Proves \lnot(\varphi \lor \psi) \lif (\lnot \varphi \land \lnot \psi)$.
\end{enumerate}
\end{prob}

\begin{prob}
पुढील गोष्टी दाखवणाऱ्या निष्पत्त्या द्या:
\begin{enumerate}
\item $\lnot(\varphi \lif \psi) \Proves \varphi$.
\item $\lnot(\varphi \land \psi) \Proves \lnot \varphi \lor \lnot \psi$.
\item $\varphi \lif \psi \Proves \lnot \varphi \lor \psi$.
\item $\Proves \lnot \lnot \varphi \lif \varphi$.
\item $\varphi \lif \psi, \lnot \varphi \lif \psi \Proves \psi$.
\item $(\varphi \land \psi) \lif \chi \Proves (\varphi \lif \chi) \lor (\psi \lif \chi)$.
\item $(\varphi \lif \psi) \lif \varphi \Proves \varphi$.
\item $\Proves (\varphi \lif \psi) \lor (\psi \lif \chi)$.
\end{enumerate}
(या सर्वांसाठी $\FalseCl$~नियम आवश्यक आहे.)
\end{prob}









\section{संख्यापकांसह निष्पत्ती}\label{fol:ntd:prq:sec}

\begin{ex}
संख्यापक हाताळताना आयगेन चराच्या अटीचा भंग होणार नाही याची खात्री करावी
लागते; त्यासाठी कधीकधी काही अनुमाने करण्याचा क्रम बदलून पाहावा लागतो.
सामान्यतः आयगेन चराच्या अटीच्या अधीन असलेले नियम आधी हाताळण्याचा प्रयत्न
उपयोगी ठरतो (पूर्ण सिद्धतेत ते अधिक खाली असतील).

सूत्र $\lexists[x][\lnot \varphi(x)] \lif \lnot \lforall[x][\varphi(x)]$
ची निष्पत्ती कशी द्यायची ते पाहू. नेहमीप्रमाणे सुरुवात करून लिहू:
\begin{prooftree}
\AxiomC{}
\UnaryInfC{$\lexists[x][\lnot \varphi(x)]\lif \lnot \lforall[x][\varphi(x)]$}
\end{prooftree}
ही शेवटची पायरी \Intro{\lif} नियमाने समर्थित होण्यासाठी काय आवश्यक आहे ते
प्रथम लिहू.
\begin{prooftree}
\AxiomC{$\Discharge{\lexists[x][\lnot \varphi(x)]}{1}$}
\DeduceC{$\lnot \lforall[x][\varphi(x)]$}
\DischargeRule{\Intro{\lif}}{1}
\UnaryInfC{$\lexists[x][\lnot \varphi(x)]\lif \lnot \lforall[x][\varphi(x)]$}
\end{prooftree}
$\lnot \lforall[x][\varphi(x)]$ वर लागू करता येईल असा कोणताही स्पष्ट नियम
नसल्यामुळे, \Elim{\lexists} नियम वापरता येईल अशा रीतीने निष्पत्ती
मांडू. येथे आयगेन चराच्या अटीकडे लक्ष देऊन, प्रमुख आधारविधान
$\lexists[x][\lnot \varphi(x)]$ किंवा निष्कर्ष यांत, तसेच त्यांच्याकडे नेणाऱ्या
निष्पत्तीतील इतर कोणत्याही मुक्त न केलेल्या गृहीतकात न येणारा अचर निवडावा
लागतो. (परंतु येथे कोणतेही स्थिरांक येत नसल्यामुळे कोणताही अचर चालेल.)
\begin{prooftree}
\AxiomC{$\Discharge{\lexists[x][\lnot \varphi(x)]}{1}$}
\AxiomC{$\Discharge{\lnot \varphi(a)}{2}$}
\DeduceC{$\lnot \lforall[x][\varphi(x)]$}
\DischargeRule{\Elim{\lexists}}{2}
\BinaryInfC{$\lnot \lforall[x][\varphi(x)]$}
\DischargeRule{\Intro{\lif}}{1}
\UnaryInfC{$\lexists[x][\lnot \varphi(x)] \lif \lnot \lforall[x][\varphi(x)]$}
\end{prooftree}
$\lnot \lforall[x][\varphi(x)]$ हे निष्पन्न करण्यासाठी \Intro{\lnot} नियम
वापरून पाहू: यासाठी, आवश्यक असल्यास $\lforall[x][\varphi(x)]$ हे अतिरिक्त गृहीतक वापरून,
व्याघात निष्पन्न करावा लागेल. अर्थात या व्याघातात $\lnot \varphi(a)$ हे गृहीतक
वापरले जाऊ शकते; ते \Elim{\lexists} अनुमानाने मुक्त केले जाईल. ही रचना
पुढीलप्रमाणे करता येते:
\begin{prooftree}
\AxiomC{$\Discharge{\lexists[x][\lnot \varphi(x)]}{1}$}
\AxiomC{$\Discharge{\lnot \varphi(a)}{2}, \Discharge{\lforall[x][\varphi(x)]}{3}$}
\DeduceC{$\lfalse$}
\DischargeRule{\Intro{\lnot}}{3}
\UnaryInfC{$\lnot \lforall[x][\varphi(x)]$}
\DischargeRule{\Elim{\lexists}}{2}
\BinaryInfC{$\lnot \lforall[x][\varphi(x)]$}
\DischargeRule{\Intro{\lif}}{1}
\UnaryInfC{$\lexists[x][\lnot \varphi(x)]\lif \lnot \lforall[x][\varphi(x)]$}
\end{prooftree}
व्याघात मिळण्याच्या आपण अगदी जवळ आलो आहोत. आयगेन चराची अट नसलेला
\Elim{\lforall} नियम येथे लागू करणे सर्वांत सोपे आहे. सार्वत्रिक
संख्यापकाने बद्ध असलेल्या~$x$ च्या जागी हवे ते कोणतेही बंद पद ठेवता येते;
व्याघातापर्यंत पोहोचण्यासाठी~$a$ चाच पुढे वापर करणे सर्वांत सयुक्तिक आहे.
\begin{prooftree}
\AxiomC{$\Discharge{\lexists[x][\lnot \varphi(x)]}{1}$}
\AxiomC{$\Discharge{\lnot \varphi(a)}{2}$}
\AxiomC{$\Discharge{\lforall[x][\varphi(x)]}{3}$}
\RightLabel{\Elim{\lforall}}
\UnaryInfC{$\varphi(a)$}
\RightLabel{\Elim{\lnot}}
\BinaryInfC{$\lfalse$}
\DischargeRule{\Intro{\lnot}}{3}
\UnaryInfC{$\lnot \lforall[x][\varphi(x)]$}
\DischargeRule{\Elim{\lexists}}{2}
\BinaryInfC{$\lnot \lforall[x][\varphi(x)]$}
\DischargeRule{\Intro{\lif}}{1}
\UnaryInfC{$\lexists[x][\lnot \varphi(x)]\lif \lnot \lforall[x][\varphi(x)]$}
\end{prooftree}

विशेषतः संख्यापक हाताळताना, आयगेन चराच्या अटीचा भंग झालेला नाही याची या
टप्प्यावर पुन्हा तपासणी करणे महत्त्वाचे आहे. आपण लावलेल्या नियमांपैकी फक्त
\Elim{\exists} हा नियम त्या अटीच्या अधीन होता. त्याचा आयगेन चर~$a$ प्रमुख
आधारविधानात किंवा निष्कर्षात येत नाही आणि नियमाने मुक्त केलेले तात्पुरते
गृहीतक वगळता संबंधित कोणत्याही मुक्त न केलेल्या गृहीतकातही येत नाही. म्हणून
ही योग्य निष्पत्ती आहे.
\end{ex}

\begin{ex}
कधीकधी इतर सूत्रांपासून सूत्र निष्पन्न करायचे असते. अशा वेळी
काही गृहीतके मुक्त न केलेली राहू शकतात. अंतिम ध्येयाबरोबरच गृहीतकांचीही
नोंद ठेवणे महत्त्वाचे आहे.

गृहीतके $\lexists[x][(\varphi(x) \land \psi(x))]$ आणि
$\lforall[x][(\psi(x) \lif \chi(x,b))]$ यांपासून सूत्र
$\lexists[x][\chi(x,b)]$ ची निष्पत्ती कशी द्यायची ते पाहू.
नेहमीप्रमाणे सुरुवात करून निष्कर्ष तळाशी लिहू.
\begin{prooftree}
\AxiomC{}
\UnaryInfC{$\lexists[x][\chi(x,b)]$}
\end{prooftree}

आपल्याकडे वापरण्यासाठी दोन आधारविधाने आहेत. पहिले वापरण्यासाठी, म्हणजे
$\lexists[x][(\varphi(x) \land \psi(x))]$ पासून
$\lexists[x][\chi(x, b)]$ ची निष्पत्ती शोधण्याचा प्रयत्न करण्यासाठी,
\Elim{\lexists} नियम वापरू. त्याला आयगेन चराची अट असल्यामुळे तो नियम आधी
लावू. मग पुढील मिळते:
\begin{prooftree}
\AxiomC{$\lexists[x][(\varphi(x) \land \psi(x))]$}
\AxiomC{$\Discharge{\varphi(a) \land \psi(a)}{1}$}
\DeduceC{$\lexists[x][\chi(x,b)]$}
\DischargeRule{\Elim{\lexists}}{1}
\BinaryInfC{$\lexists[x][\chi(x,b)]$}
\end{prooftree}
आपण वापरत असलेल्या दोन्ही गृहीतकांत~$\psi$ समान आहे. या टप्प्यावर
\Elim{\land} लावून~$\psi(a)$ वेगळे काढणे उपयोगी ठरेल.
\begin{prooftree}
\AxiomC{$\lexists[x][(\varphi(x) \land \psi(x)])$}
\AxiomC{$\Discharge{\varphi(a)
\land \psi(a)}{1}$}
\RightLabel{\Elim{\land}}
\UnaryInfC{$\psi(a)$}
\DeduceC{$\lexists[x][\chi(x,b)]$}
\DischargeRule{\Elim{\lexists}}{1}
\BinaryInfC{$\lexists[x][\chi(x,b)]$}
\end{prooftree}

वापरण्यासाठी असलेले दुसरे गृहीतक म्हणजे~$\lforall[x][(\psi(x) \lif
  \chi(x,b))]$. येथे आयगेन चराची अट नसल्यामुळे \Elim{\lforall} वापरून
स्थिरांक~$a$ हे $x$ च्या जागी ठेवता येते आणि
$\psi(a) \lif \chi(a, b)$ मिळते. आता आपल्याकडे $\psi(a) \lif \chi(a,b)$ आणि
$\psi(a)$ दोन्ही आहेत. यानंतर \Elim{\lif} नियमाचे सरळ उपयोजन करावे.
\begin{prooftree}
\AxiomC{$\lexists[x][(\varphi(x) \land \psi(x))]$}
\AxiomC{$\lforall[x][(\psi(x) \lif \chi(x,b))]$}
\RightLabel{\Elim{\lforall}}
\UnaryInfC{$\psi(a) \lif \chi(a,b)$}
\AxiomC{$\Discharge{\varphi(a)
\land \psi(a)}{1}$}
\RightLabel{\Elim{\land}}
\UnaryInfC{$\psi(a)$}
\RightLabel{\Elim{\lif}}
\BinaryInfC{$\chi(a,b)$}
\DeduceC{$\lexists[x][\chi(x,b)]$}
\DischargeRule{\Elim{\lexists}}{1}
\insertBetweenHyps{\hspace{-5em}}
\BinaryInfC{$\lexists[x][\chi(x,b)]$}
\end{prooftree}
आपण ध्येयाच्या अगदी जवळ आलो आहोत!{} \Intro{\lexists} चे एक उपयोजन केले
की ध्येय गाठले.
\begin{prooftree}
\AxiomC{$\lexists[x][(\varphi(x) \land \psi(x))]$}
\AxiomC{$\lforall[x][(\psi(x) \lif \chi(x,b))]$}
\RightLabel{\Elim{\lforall}}
\UnaryInfC{$\psi(a) \lif \chi(a,b)$}
\AxiomC{$\Discharge{\varphi(a)
\land \psi(a)}{1}$}
\RightLabel{\Elim{\land}}
\UnaryInfC{$\psi(a)$}
\RightLabel{\Elim{\lif}}
\BinaryInfC{$\chi(a,b)$}
\RightLabel{\Intro{\lexists}}
\UnaryInfC{$\lexists[x][\chi(x,b)]$}
\DischargeRule{\Elim{\lexists}}{1}
\insertBetweenHyps{\hspace{-5em}}
\BinaryInfC{$\lexists[x][\chi(x,b)]$}
\end{prooftree}
प्रत्येक पायरीवर आयगेन चराच्या अटींचा भंग झालेला नाही याची खात्री केल्यामुळे
ही योग्य निष्पत्ती आहे असा विश्वास ठेवता येतो.
\end{ex}

\begin{ex}
गृहीतके $\lforall[x][\varphi(x)] \lif \lexists[y][\psi(y)]$ आणि
$\lnot\lexists[y][\psi(y)]$ यांपासून सूत्र
$\lnot\lforall[x][\varphi(x)]$ ची निष्पत्ती द्या.
नेहमीप्रमाणे सुरुवात करून अपेक्षित सूत्र तळाशी लिहू.
\begin{prooftree}
\AxiomC{}
\UnaryInfC{$\lnot\lforall[x][\varphi(x)]$}
\end{prooftree}
निष्पत्तीची शेवटची ओळ नकरण आहे; म्हणून \Intro{\lnot} वापरून पाहू.
त्यासाठी व्याघात कसा निष्पन्न करायचा ते शोधावे लागेल.
\begin{prooftree}
\AxiomC{$\Discharge{\lforall[x][\varphi(x)]}{1}$}
\DeduceC{$\lfalse$}
\DischargeRule{\Intro{\lnot}}{1}
\UnaryInfC{$\lnot\lforall[x][\varphi(x)]$}
\end{prooftree}
येथपर्यंत सर्व ठीक आहे. \Elim{\lforall} वापरता येईल, पण त्याने ध्येयापर्यंत
पोहोचण्यास मदत होईल का हे स्पष्ट नाही. त्याऐवजी आपले एक गृहीतक वापरू.
$\lforall[x][\varphi(x)] \lif \lexists[y][\psi(y)]$ आणि
$\lforall[x][\varphi(x)]$ यांच्या आधारे \Elim{\lif} नियम वापरता येईल.
\begin{prooftree}
\AxiomC{$\lforall[x][\varphi(x)] \lif \lexists[y][\psi(y)]$}
\AxiomC{$\Discharge{\lforall[x][\varphi(x)]}{1}$}
\RightLabel{\Elim{\lif}}
\BinaryInfC{$\lexists[y][\psi(y)]$}
\DeduceC{$\lfalse$}
\DischargeRule{\Intro{\lnot}}{1}
\UnaryInfC{$\lnot\lforall[x][\varphi(x)]$}
\end{prooftree}
आता वापरण्यासाठी एक अखेरचे गृहीतक उरले आहे. \Elim{\lnot} वापरून
व्याघातापर्यंत पोहोचण्यास ते मदत करेल असे दिसते.
\begin{prooftree}
\AxiomC{$\lnot\lexists[y][\psi(y)]$}
\AxiomC{$\lforall[x][\varphi(x)] \lif \lexists[y][\psi(y)]$}
\AxiomC{$\Discharge{\lforall[x][\varphi(x)]}{1}$}
\RightLabel{\Elim{\lif}}
\BinaryInfC{$\lexists[y][\psi(y)]$}
\RightLabel{\Elim{\lnot}}
\BinaryInfC{$\lfalse$}
\DischargeRule{\Intro{\lnot}}{1}
\UnaryInfC{$\lnot\lforall[x][\varphi(x)]$}
\end{prooftree}
\end{ex}

\begin{prob}
पुढील गोष्टी दाखवणाऱ्या निष्पत्त्या द्या:
\begin{enumerate}
\item $\Proves (\lforall[x][\varphi(x)] \land \lforall[y][\psi(y)]) \lif
\lforall[z][(\varphi(z) \land \psi(z))]$.
\item $\Proves (\lexists[x][\varphi(x)] \lor \lexists[y][\psi(y)]) \lif
\lexists[z][(\varphi(z) \lor \psi(z))]$.
\item $\lforall[x][(\varphi(x) \lif \psi)] \Proves \lexists[y][\varphi(y)] \lif \psi$.
\item $\lforall[x][\lnot \varphi(x)] \Proves \lnot\lexists[x][\varphi(x)]$.
\item $\Proves \lnot\lexists[x][\varphi(x)] \lif \lforall[x][\lnot \varphi(x)]$.
\item $\Proves \lnot\lexists[x][\lforall[y][((\varphi(x,y) \lif \lnot
\varphi(y,y)) \land (\lnot \varphi(y,y) \lif \varphi(x,y)))]]$.
\end{enumerate}
\end{prob}

\begin{prob}
पुढील गोष्टी दाखवणाऱ्या निष्पत्त्या द्या:
\begin{enumerate}
\item $\Proves \lnot\lforall[x][\varphi(x)] \lif \lexists[x][\lnot\varphi(x)]$.
\item $(\lforall[x][\varphi(x)] \lif \psi) \Proves \lexists[y][(\varphi(y) \lif \psi)]$.
\item $\Proves \lexists[x][(\varphi(x) \lif \lforall[y][\varphi(y)])]$.
\end{enumerate}
(या सर्वांसाठी $\FalseCl$~नियम आवश्यक आहे.)
\end{prob}









\section{सिद्धता-उपपत्तीय संकल्पना}\label{fol:ntd:ptn:sec}

\begin{editorial}
या विभागात नैसर्गिक निगमनासाठी सिद्धतायोग्यता संबंध आणि सुसंगतता यांच्या
व्याख्या एकत्रित केल्या आहेत.
\end{editorial}

\begin{explain}
जशा आपण अनेक महत्त्वाच्या चिन्हार्थविषयक संकल्पना (वैधता, तार्किक निष्पन्नता,
पूर्ततायोग्यता) परिभाषित केल्या आहेत, तशाच आता त्यांना अनुरूप
\emph{सिद्धता-उपपत्तीय संकल्पना} परिभाषित करू. या संकल्पना संरचना मध्ये
वाक्यांची पूर्ति होते की नाही याच्या आधारे परिभाषित केलेल्या नसून,
काही वाक्यांची इतरांपासूनची निष्पन्नता किंवा
अनिष्पन्नता यांचा आधार घेऊन परिभाषित केल्या आहेत. या दोन प्रकारच्या
संकल्पना एकमेकांशी जुळतात, हा एक महत्त्वाचा शोध होता. त्या जुळतात, हाच
\emph{निर्दोषता प्रमेय} आणि \emph{संपूर्णता प्रमेय} यांचा आशय आहे.
\end{explain}


\begin{defn}[प्रमेये]
एखादे वाक्य~$\varphi$ हे \emph{प्रमेय} असते, जर नैसर्गिक निगमनात $\varphi$ ची
अशी निष्पत्ती असेल की तिच्यातील सर्व गृहीतके मुक्त आहेत.
$\varphi$ हे प्रमेय असेल तर आपण $\Proves \varphi$ लिहितो आणि ते प्रमेय नसेल तर
$\Proves/ \varphi$ लिहितो.
\end{defn}

\begin{defn}[निष्पन्नता]
वाक्य $\varphi$ हे वाक्यांच्या संच~$\Gamma$ \emph{पासून
निष्पन्न करता येण्याजोगे} असते, म्हणजे $\Gamma \Proves \varphi$, तेव्हा आणि केवळ तेव्हाच,
जेव्हा निष्कर्ष~$\varphi$ असणारी अशी निष्पत्ती असेल की तिच्यातील प्रत्येक
गृहीतक एकतर मुक्त आहे किंवा $\Gamma$ मध्ये आहे. $\varphi$ हे $\Gamma$
पासून निष्पन्न करता येण्याजोगे नसेल, तर आपण $\Gamma \Proves/ \varphi$ लिहितो.
\end{defn}

\begin{defn}[सुसंगतता]
वाक्यांचा संच~$\Gamma$ हा \emph{विसंगत} असतो तेव्हा आणि केवळ तेव्हाच,
जेव्हा $\Gamma \Proves \lfalse$. $\Gamma$ विसंगत नसेल, म्हणजे
$\Gamma \Proves/ \lfalse$ असेल, तर त्याला \emph{सुसंगत} म्हणतो.
\end{defn}

\begin{prop}[परावर्तिता]
\label{fol:ntd:ptn:prop:reflexivity}
$\varphi \in \Gamma$ असेल, तर $\Gamma \Proves \varphi$.
\end{prop}

\begin{proof}
केवळ $\varphi$ हे गृहीतक स्वतःच $\varphi$ ची अशी निष्पत्ती आहे की तिच्यातील
प्रत्येक मुक्त न केलेली गृहीतक (म्हणजे $\varphi$) $\Gamma$ मध्ये आहे.
\end{proof}


\begin{prop}[एकस्वनिकता]
\label{fol:ntd:ptn:prop:monotonicity}
$\Gamma \subseteq \Delta$ आणि $\Gamma \Proves \varphi$ असेल, तर $\Delta
\Proves \varphi$.
\end{prop}

\begin{proof}
$\Gamma$ पासून $\varphi$ ची कोणतीही निष्पत्ती ही $\Delta$ पासून $\varphi$ चीही
निष्पत्ती असते.
\end{proof}

\begin{prop}[संक्रमकता]
\label{fol:ntd:ptn:prop:transitivity}
$\Gamma \Proves \varphi$ आणि $\{\varphi\} \cup \Delta \Proves
\psi$ असेल, तर $\Gamma \cup \Delta \Proves \psi$.
\end{prop}

\begin{proof}
$\Gamma \Proves \varphi$ असेल, तर $\varphi$ ची अशी निष्पत्ती~$\delta_0$ असते
की तिच्यातील सर्व मुक्त न केलेली गृहीतके $\Gamma$ मध्ये असतात. $\{\varphi\}
\cup \Delta \Proves \psi$ असेल, तर $\psi$ ची अशी निष्पत्ती~$\delta_1$
असते की तिच्यातील सर्व मुक्त न केलेली गृहीतके $\{\varphi\} \cup \Delta$ मध्ये
असतात. आता पुढील रचना विचारात घ्या:
\begin{prooftree}
  \AxiomC{$\Delta, \Discharge{\varphi}{1}$}
  \RightLabel{$\delta_1$}
  \DeduceC{$\psi$}
  \DischargeRule{\Intro{\lif}}{1}
  \UnaryInfC{$\varphi \lif \psi$}
  \AxiomC{$\Gamma$}
  \RightLabel{$\delta_0$}
  \DeduceC{$\varphi$}
  \RightLabel{\Elim{\lif}}
  \BinaryInfC{$\psi$}
\end{prooftree}
आता सर्व मुक्त न केलेली गृहीतके $\Gamma \cup \Delta$ मध्ये आहेत; म्हणून
यावरून $\Gamma \cup \Delta \Proves \psi$ हे दिसते.
\end{proof}

$\Gamma = \{\varphi_1, \varphi_2, \ldots, \varphi_k\}$ हा सांत संच असेल, तेव्हा
$\Gamma \Proves \psi$ ऐवजी आपण $\varphi_1,\varphi_2,\ldots,\varphi_k \Proves \psi$ हे
सरलीकृत संकेतलेखन वापरू शकतो. विशेषतः, $\varphi \Proves \psi$ याचा अर्थ
$\{\varphi\} \Proves \psi$ असा होतो.

लक्षात घ्या की $\Gamma \Proves \varphi$ आणि $\varphi
\Proves \psi$ असेल, तर $\Gamma \Proves \psi$. तसेच $\varphi_1,
\dots, \varphi_n \Proves \psi$ असेल आणि प्रत्येक~$i$ साठी $\Gamma \Proves \varphi_i$
असेल, तर $\Gamma \Proves \psi$, हेसुद्धा यावरून मिळते.

\begin{prop}
\label{fol:ntd:ptn:prop:incons}
पुढील विधाने सममूल्य आहेत.
    \begin{enumerate}
      \item \( \Gamma \) विसंगत आहे.
      \item प्रत्येक वाक्य~\( {\varphi} \) साठी \( \Gamma \Proves {\varphi} \).
      \item एखाद्या वाक्य~\( {\varphi} \) साठी \( \Gamma \Proves {\varphi} \) आणि \( \Gamma \Proves \lnot {\varphi} \).
    \end{enumerate}
\end{prop}

\begin{proof}
सराव.
\end{proof}


\begin{prob}
\ref{fol:ntd:ptn:prop:incons} सिद्ध करा.
\end{prob}




\begin{prop}[संहतता]
  \label{fol:ntd:ptn:prop:proves-compact}
  \begin{enumerate}
  \item $\Gamma \Proves \varphi$ असेल, तर एखादा सांत उपसंच $\Gamma_0
    \subseteq \Gamma$ असा आहे की $\Gamma_0 \Proves \varphi$.
  \item $\Gamma$ चा प्रत्येक सांत उपसंच सुसंगत असेल, तर $\Gamma$ सुसंगत
    असतो.
  \end{enumerate}
\end{prop}

\begin{proof}
  \begin{enumerate}
    \item $\Gamma \Proves \varphi$ असेल, तर $\Gamma$ पासून $\varphi$ ची
      निष्पत्ती~$\delta$ असते. $\Gamma_0$ हा $\delta$ मधील
      मुक्त न केलेली गृहीतकांचा संच असू द्या. प्रत्येक निष्पत्ती सांत
      असल्यामुळे $\Gamma_0$ मधील वाक्यांची संख्या सांतच असते.
      म्हणून $\delta$ ही सांत~$\Gamma_0 \subseteq \Gamma$ पासून $\varphi$ ची
      निष्पत्ती आहे.
    \item (1) मध्ये $\varphi
      \ident \lfalse$ हे विशेष प्रकरण घेतल्यावर मिळणारे हे निषेधात्मक उलट
      विधान आहे.
  \end{enumerate}
\end{proof}









\section{निष्पन्नता आणि सुसंगतता}\label{fol:ntd:prv:sec}

आता आपण निष्पन्नता संबंधाचे अनेक गुणधर्म सिद्ध करू. ते स्वतंत्रपणेही
महत्त्वाचे आहेत, पण संपूर्णता प्रमेयाच्या सिद्धतेत प्रत्येकाची भूमिका असेल.

\begin{prop}\label{fol:ntd:prv:prop:provability-contr}
  जर $\Gamma \Proves \varphi$ आणि $\Gamma \cup \{\varphi\}$ विसंगत असेल, तर
  $\Gamma$ विसंगत आहे.
\end{prop}

\begin{proof}
$\Gamma$ पासून $\varphi$ ची निष्पत्ती~$\delta_1$ आणि $\Gamma \cup \{\varphi\}$
पासून $\lfalse$ ची निष्पत्ती~$\delta_2$ असू द्या. मग आपण पुढीलप्रमाणे
निष्पन्न करू शकतो:
\begin{prooftree}
\AxiomC{$\Gamma, \Discharge{\varphi}{1}$}
\RightLabel{$\delta_2$}
\DeduceC{$\lfalse$}
\DischargeRule{\Intro{\lnot}}{1}
\UnaryInfC{$\lnot \varphi$}
\AxiomC{$\Gamma$}
\RightLabel{$\delta_1$}
\DeduceC{$\varphi$}
\RightLabel{\Elim{\lnot}}
\BinaryInfC{$\lfalse$}
\end{prooftree}
नव्या निष्पत्तीमध्ये गृहीतक~$\varphi$ हे मुक्त आहे; म्हणून ही
$\Gamma$ पासूनची निष्पत्ती आहे.
\end{proof}

\begin{prop}
\label{fol:ntd:prv:prop:prov-incons}
$\Gamma \Proves \varphi$ तेव्हा आणि केवळ तेव्हाच, जेव्हा $\Gamma \cup \{\lnot \varphi\}$
विसंगत असते.
\end{prop}

\begin{proof}
प्रथम $\Gamma \Proves \varphi$ असे समजा; म्हणजे $\Gamma$ मधील
मुक्त न केलेली गृहीतकांपासून $\varphi$ ची निष्पत्ती~$\delta_0$ असते.
पुढीलप्रमाणे $\Gamma \cup \{\lnot \varphi\}$ पासून $\lfalse$ ची
निष्पत्ती मिळते:
\begin{prooftree}
  \AxiomC{$\lnot \varphi$}
  \AxiomC{$\Gamma$}
  \RightLabel{$\delta_0$}
  \DeduceC{$\varphi$}
  \RightLabel{\Elim{\lnot}}
  \BinaryInfC{$\lfalse$}
\end{prooftree}

आता $\Gamma \cup \{\lnot \varphi\}$ विसंगत आहे असे समजा आणि $\Gamma \cup
\{\lnot \varphi\}$ मधील मुक्त न केलेली गृहीतकांपासून $\lfalse$ ची संबंधित
निष्पत्ती~$\delta_1$ असू द्या. $\FalseCl$ वापरून पुढीलप्रमाणे केवळ
$\Gamma$ पासून $\varphi$ ची निष्पत्ती मिळते:
\begin{prooftree}
  \AxiomC{$\Gamma, \Discharge{\lnot \varphi}{1}$}
  \RightLabel{$\delta_1$}
  \DeduceC{$\lfalse$}
  \RightLabel{\FalseCl}
  \DischargeRule{\FalseCl}{1}
  \UnaryInfC{$\varphi$}
\end{prooftree}
\end{proof}

\begin{prob}
$\Gamma \Proves \lnot \varphi$ तेव्हा आणि केवळ तेव्हाच, जेव्हा $\Gamma \cup \{\varphi\}$
विसंगत असते, हे सिद्ध करा.
\end{prob}

\begin{prop}\label{fol:ntd:prv:prop:explicit-inc}
  जर $\Gamma \Proves \varphi$ आणि $\lnot \varphi \in \Gamma$ असेल, तर $\Gamma$
  विसंगत आहे.
\end{prop}

\begin{proof}
  समजा $\Gamma \Proves \varphi$ आणि $\lnot \varphi \in \Gamma$. मग $\Gamma$ पासून
  $\varphi$ ची निष्पत्ती~$\delta$ असते. $\Elim{\lnot}$ नियमाचा पुढील सोपा
  उपयोग विचारात घ्या:
  \begin{prooftree}
    \AxiomC{$\lnot \varphi$}
    \AxiomC{$\Gamma$}
    \RightLabel{$\delta$}
    \DeduceC{$\varphi$}
    \RightLabel{\Elim{\lnot}}
    \BinaryInfC{$\lfalse$}
  \end{prooftree}
  $\lnot \varphi \in \Gamma$ असल्यामुळे सर्व मुक्त न केलेली गृहीतके
  $\Gamma$ मध्ये आहेत. म्हणून यावरून $\Gamma \Proves \lfalse$ हे दिसते.
\end{proof}

\begin{prop}\label{fol:ntd:prv:prop:provability-exhaustive}
  $\Gamma \cup \{\varphi\}$ आणि $\Gamma \cup \{\lnot \varphi\}$ हे दोन्ही संच
  विसंगत असतील, तर $\Gamma$ विसंगत आहे.
\end{prop}

\begin{proof}
$\delta_1$ आणि $\delta_2$ या अनुक्रमे $\Gamma \cup \{ \varphi \}$ पासून
$\lfalse$ ची आणि $\Gamma \cup \{ \lnot \varphi \}$ पासून $\lfalse$ ची
निष्पत्त्या असू द्या. मग आपण पुढीलप्रमाणे निष्पन्न करू शकतो:
\begin{prooftree}
\AxiomC{$\Gamma, \Discharge{\lnot \varphi}{2}$}
\RightLabel{$\delta_2$}
\DeduceC{$\lfalse$}
\DischargeRule{\Intro{\lnot}}{2}
\UnaryInfC{$\lnot \lnot \varphi$}
\AxiomC{$\Gamma, \Discharge{\varphi}{1}$}
\RightLabel{$\delta_1$}
\DeduceC{$\lfalse$}
\DischargeRule{\Intro{\lnot}}{1}
\UnaryInfC{$\lnot \varphi$}
\RightLabel{\Elim{\lnot}}
\BinaryInfC{$\lfalse$}
\end{prooftree}
$\varphi$ आणि $\lnot \varphi$ ही गृहीतके मुक्त असल्यामुळे ही केवळ
$\Gamma$ पासून $\lfalse$ ची निष्पत्ती आहे. म्हणून $\Gamma$ विसंगत
आहे.
\end{proof}










\section{निष्पन्नता आणि विधानीय संयोजक}\label{fol:ntd:ppr:sec}

\begin{explain}
  नैसर्गिक निगमनाचा निष्पन्नता संबंध~$\Proves$ हा विधानीय संयोजकांशी
  संबंधित काही मूलभूत तथ्ये सिद्ध करण्यासाठी पुरेसा सामर्थ्यवान आहे; उदाहरणार्थ,
  $\varphi \land \psi
  \Proves \varphi$ आणि $\varphi, \varphi \lif \psi \Proves \psi$ (विध्यात्मक अनुमान). ही
  तथ्ये संपूर्णता प्रमेयाच्या सिद्धतेसाठी आवश्यक आहेत.
\end{explain}

\begin{prop}\label{fol:ntd:ppr:prop:provability-land}
  \begin{enumerate}
  \item \label{fol:ntd:ppr:prop:provability-land-left} $\varphi \land \psi \Proves
    \varphi$ आणि $\varphi \land \psi \Proves \psi$ ही दोन्ही तथ्ये सत्य आहेत.
  \item \label{fol:ntd:ppr:prop:provability-land-right} $\varphi, \psi \Proves \varphi \land \psi$.
  \end{enumerate}
\end{prop}

\begin{proof}
  \begin{enumerate}
  \item आपण पुढील दोन्ही निष्पन्न करू शकतो:
    \begin{prooftree}
      \AxiomC{$\varphi \land \psi$}
      \RightLabel{\Elim{\land}}
      \UnaryInfC{$\varphi$}
      \DisplayProof\qquad\bottomAlignProof
      \AxiomC{$\varphi \land \psi$}
      \RightLabel{\Elim{\land}}
      \UnaryInfC{$\psi$}
    \end{prooftree}
  \item आपण पुढील निष्पन्न करू शकतो:
    \begin{prooftree}
      \AxiomC{$\varphi$}
      \AxiomC{$\psi$}
      \RightLabel{\Intro{\land}}
      \BinaryInfC{$\varphi \land \psi$}
    \end{prooftree}
  \end{enumerate}
\end{proof}


\begin{prop}\label{fol:ntd:ppr:prop:provability-lor}
  \begin{enumerate}
  \item $\varphi \lor \psi, \lnot \varphi, \lnot \psi$ विसंगत आहे.
  \item $\varphi \Proves \varphi \lor \psi$ आणि $\psi \Proves \varphi \lor \psi$ ही दोन्ही
    तथ्ये सत्य आहेत.
  \end{enumerate}
\end{prop}

\begin{proof}
  \begin{enumerate}
  \item पुढील निष्पत्ती विचारात घ्या:
    \begin{prooftree}
      \AxiomC{$\varphi \lor \psi$}
      \AxiomC{$\lnot \varphi$}
      \AxiomC{$\Discharge{\varphi}{1}$}
      \RightLabel{\Elim{\lnot}}
      \BinaryInfC{$\lfalse$}
      \AxiomC{$\lnot \psi$}
      \AxiomC{$\Discharge{\psi}{1}$}
      \RightLabel{\Elim{\lnot}}
      \BinaryInfC{$\lfalse$}
      \DischargeRule{\Elim{\lor}}{1}
      \TrinaryInfC{$\lfalse$}
    \end{prooftree}
    ही $\varphi \lor \psi$, $\lnot \varphi$ आणि $\lnot \psi$ या मुक्त न केलेली
    गृहीतकांपासून $\lfalse$ ची निष्पत्ती आहे.
  \item आपण पुढील दोन्ही निष्पन्न करू शकतो:
    \begin{prooftree}
      \AxiomC{$\varphi$}
      \RightLabel{\Intro{\lor}}
      \UnaryInfC{$\varphi \lor \psi$}
      \DisplayProof\qquad\bottomAlignProof
      \AxiomC{$\psi$}
      \RightLabel{\Intro{\lor}}
      \UnaryInfC{$\varphi \lor \psi$}
    \end{prooftree}
  \end{enumerate}
\end{proof}

\begin{prop}\label{fol:ntd:ppr:prop:provability-lif}
  \begin{enumerate}
  \item \label{fol:ntd:ppr:prop:provability-lif-left} $\varphi, \varphi \lif \psi \Proves \psi$.
  \item \label{fol:ntd:ppr:prop:provability-lif-right}
    $\lnot \varphi \Proves \varphi \lif \psi$ आणि $\psi \Proves \varphi \lif \psi$ ही दोन्ही
    तथ्ये सत्य आहेत.
  \end{enumerate}
\end{prop}

\begin{proof}
  \begin{enumerate}
  \item आपण पुढील निष्पन्न करू शकतो:
    \begin{prooftree}
      \AxiomC{$\varphi \lif \psi$}
      \AxiomC{$\varphi$}
      \RightLabel{\Elim{\lif}}
      \BinaryInfC{$\psi$}
    \end{prooftree}

  \item हे पुढील दोन निष्पत्त्या दाखवतात:
    \begin{prooftree}
      \AxiomC{$\lnot \varphi$}
      \AxiomC{$\Discharge{\varphi}{1}$}
      \RightLabel{\Elim{\lnot}}
      \BinaryInfC{$\lfalse$}
      \RightLabel{\FalseInt}
      \UnaryInfC{$\psi$}
      \DischargeRule{\Intro{\lif}}{1}
      \UnaryInfC{$\varphi \lif \psi$}
      \DisplayProof\qquad\bottomAlignProof
      \AxiomC{$\psi$}
      \RightLabel{\Intro{\lif}}
      \UnaryInfC{$\varphi \lif \psi$}
    \end{prooftree}
    लक्षात घ्या की $\Intro{\lif}$ गृहीतक~$\varphi$ हे मुक्त करू शकतो;
    पण ते करणे आवश्यक नाही.
  \end{enumerate}
\end{proof}










\section{निष्पन्नता आणि संख्यापक}\label{fol:ntd:qpr:sec}

\begin{explain}
  संपूर्णता प्रमेयासाठी या विभागात स्थापित केलेली~$\Proves$ संबंधी तथ्ये
  नैसर्गिक निगमनाच्या नियमांनी निष्पन्न होतात, हेही आवश्यक आहे.
\end{explain}

\begin{thm}
\label{fol:ntd:qpr:thm:strong-generalization} जर $c$ हे $\Gamma$ किंवा $\varphi(x)$ मध्ये
न येणारे स्थिर असेल आणि $\Gamma \Proves \varphi(c)$, तर $\Gamma
\Proves \lforall[x][\varphi(x)]$.
\end{thm}

\begin{proof}
$\delta$ ही $\Gamma$ पासून $\varphi(c)$ ची निष्पत्ती असू द्या.
\Intro{\lforall} अनुमान जोडल्यावर आपल्याला $\lforall[x][\varphi(x)]$ ची
निष्पत्ती मिळते. $c$ हे $\Gamma$ किंवा $\varphi(x)$ मध्ये येत नसल्यामुळे
आयगेन चराची अट पूर्ण होते.
\end{proof}

\begin{prop}
\label{fol:ntd:qpr:prop:provability-quantifiers}
\begin{enumerate}
\item $\varphi(t) \Proves \lexists[x][\varphi(x)]$.

\item $\lforall[x][\varphi(x)] \Proves \varphi(t)$.
\end{enumerate}
\end{prop}

\begin{proof}
\begin{enumerate}
\item पुढील ही~$\lexists[x][\varphi(x)]$ ची $\varphi(t)$ पासूनची
  निष्पत्ती आहे:
\begin{prooftree}
\AxiomC{$\varphi(t)$}
\RightLabel{\Intro{\lexists}}
\UnaryInfC{$\lexists[x][\varphi(x)]$}
\end{prooftree}

\item पुढील ही~$\varphi(t)$ ची $\lforall[x][\varphi(x)]$ पासूनची
  निष्पत्ती आहे:
\begin{prooftree}
\AxiomC{$\lforall[x][\varphi(x)]$}
\RightLabel{\Elim{\lforall}}
\UnaryInfC{$\varphi(t)$}
\end{prooftree}
\end{enumerate}
\end{proof}









\section{निर्दोषता}\label{fol:ntd:sou:sec}

\begin{explain}
नैसर्गिक निगमनासारखी निष्पत्ती-पद्धत \emph{निर्दोष} असते, जर ती
प्रत्यक्षात निष्पन्न न होणाऱ्या गोष्टी निष्पन्न करू शकत नसेल. त्यामुळे
निर्दोषता हा निष्पत्ती-पद्धतींसाठी हमी दिलेल्या सुरक्षिततेचा एक प्रकार
आहे. कोणता सिद्धता-उपपत्तीय गुणधर्म विचारात आहे यावर अवलंबून, उदाहरणार्थ,
आपल्याला पुढील गोष्टी हव्या असतात:
\begin{enumerate}
\item प्रत्येक निष्पन्न करता येण्याजोगे वाक्य हे वैध असते;
\item एखादे वाक्य इतर काही वाक्यांपासून निष्पन्न करता येण्याजोगे असेल, तर ते
  त्यांचा तार्किक परिणामही असते;
\item वाक्यांचा संच विसंगत असेल, तर तो अपूर्ततायोग्य असतो.
\end{enumerate}
हे निष्पत्ती-पद्धतीचे महत्त्वाचे गुणधर्म आहेत. यांपैकी एखादा गुणधर्म
खरा नसेल, तर निष्पत्ती-पद्धतीत दोष आहे---ती खूप जास्त गोष्टी
निष्पन्न करेल. म्हणून निष्पत्ती-पद्धतीची निर्दोषता सिद्ध करणे अत्यंत
महत्त्वाचे आहे.
\end{explain}

\begin{thm}[निर्दोषता]
\label{fol:ntd:sou:thm:soundness}
$\varphi$ हे मुक्त न केलेली गृहीतके~$\Gamma$ यांच्यापासून निष्पन्न करता येण्याजोगे असेल, तर
$\Gamma \Entails \varphi$.
\end{thm}

\begin{proof}
$\delta$ ही $\varphi$ ची निष्पत्ती असू द्या. $\delta$ मधील अनुमानांच्या
संख्येवर विगमनाने पुढे जाऊ.

विगमनाच्या आधारपायरीसाठी अनुमानांची संख्या~$0$ असेल तेव्हा दावा सिद्ध करू.
या प्रकरणात $\delta$ मध्ये फक्त एकच वाक्य~$\varphi$, म्हणजे एक गृहीतक,
असते. ते गृहीतक मुक्त न केलेली असते, कारण गृहीतके फक्त अनुमानांद्वारे
मुक्त होऊ शकतात आणि येथे कोणतेही अनुमान नाही. म्हणून सिद्धतेतील
सर्व मुक्त न केलेली गृहीतके पूर्ण करणारी कोणतीही
संरचना~$\Struct{M}$
$\varphi$ हेदेखील पूर्ण करते.

आता विगमन पायरी पाहू. $\delta$ मध्ये $n$ अनुमाने आहेत असे समजा. सर्वांत
खालच्या अनुमानाची आधारविधाने अशा उप-निष्पत्त्या वापरून निष्पन्न केली
आहेत की प्रत्येकामध्ये $n$ पेक्षा कमी अनुमाने आहेत. विगमन गृहीतक असे धरू:
सर्वांत खालच्या अनुमानाची आधारविधाने, त्या आधारविधानांवर संपणाऱ्या
उप-निष्पत्त्यांच्या मुक्त न केलेली गृहीतकांपासून निष्पन्न होतात. संपूर्ण
सिद्धतेच्या मुक्त न केलेली गृहीतकांपासून निष्कर्ष~$\varphi$ निष्पन्न होतो, हे
आपल्याला दाखवायचे आहे.

सर्वांत खालच्या अनुमानाच्या प्रकारानुसार प्रकरणे वेगळी करू. प्रथम एकच
आधारविधान असलेली संभाव्य अनुमाने पाहू.

\begin{enumerate}
\item शेवटचे अनुमान \Intro{\lnot} आहे असे समजा: निष्पत्तीचे रूप
  पुढीलप्रमाणे आहे:
  \begin{prooftree}
    \AxiomC{$\Gamma, \Discharge{\varphi}{n}$}
    \RightLabel{$\delta_1$}
    \DeduceC{$\lfalse$}
    \DischargeRule{\Intro{\lnot}}{n}
    \UnaryInfC{$\lnot \varphi$}
  \end{prooftree}
  विगमन गृहीतकानुसार $\lfalse$ हे $\delta_1$ च्या मुक्त न केलेली
  गृहीतकांपासून~$\Gamma \cup \{\varphi\}$ निष्पन्न होते. एखादी
  संरचना~$\Struct{M}$
  विचारात घ्या. जर $\Sat{M}{\Gamma}$, तर
  $\Sat{M}{\lnot \varphi}$, हे दाखवायचे आहे.
  विसंगतीद्वारे सिद्धतेसाठी असे समजा की
  $\Sat{M}{\Gamma}$, पण
  $\Sat/{M}{\lnot \varphi}$, म्हणजे
  $\Sat{M}{\varphi}$. याचा अर्थ
  $\Sat{M}{\Gamma \cup \{\varphi\}}$ असा होईल. हे आपल्या
  विगमन गृहीतकाच्या विरुद्ध आहे. म्हणून
  $\Sat{M}{\lnot \varphi}$.

\item शेवटचे अनुमान \Elim{\land} आहे. याचे दोन प्रकार आहेत: $\varphi \land
  \psi$ या आधारविधानापासून $\varphi$ किंवा $\psi$ निष्पन्न केले जाऊ शकते. पहिले
  प्रकरण पाहू. निष्पत्ती~$\delta$ पुढीलप्रमाणे दिसते:
  \begin{prooftree}
    \AxiomC{$\Gamma$}
    \RightLabel{$\delta_1$}
    \DeduceC{$\varphi \land \psi$}
    \RightLabel{\Elim{\land}}
    \UnaryInfC{$\varphi$}
  \end{prooftree}
  विगमन गृहीतकानुसार $\varphi \land \psi$ हे $\delta_1$ च्या
  मुक्त न केलेली गृहीतकांपासून~$\Gamma$ निष्पन्न होते. एखादी
  संरचना~$\Struct{M}$ विचारात घ्या.
  जर $\Sat{M}{\Gamma}$, तर
  $\Sat{M}{\varphi}$, हे दाखवायचे आहे.
  $\Sat{M}{\Gamma}$ असे समजा. आपल्या विगमन
  गृहीतकानुसार ($\Gamma \Entails \varphi \land \psi$),
  $\Sat{M}{\varphi \land \psi}$ हे आपल्याला माहीत आहे.
  व्याख्येनुसार, $\Sat{M}{\varphi \land \psi}$ तेव्हा आणि
  केवळ तेव्हाच, जेव्हा $\Sat{M}{\varphi}$ आणि
  $\Sat{M}{\psi}$. ($\varphi \land \psi$ पासून $\psi$
  निष्पन्न करण्याचे प्रकरणही असेच हाताळता येते.)

\item शेवटचे अनुमान \Intro{\lor} आहे. याचे दोन प्रकार आहेत: $\varphi$ या
  आधारविधानापासून किंवा $\psi$ या आधारविधानापासून $\varphi \lor \psi$ निष्पन्न
  केले जाऊ शकते. पहिले प्रकरण पाहू. निष्पत्तीचे रूप पुढीलप्रमाणे आहे:
  \begin{prooftree}
    \AxiomC{$\Gamma$}
    \RightLabel{$\delta_1$}
    \DeduceC{$\varphi$}
    \RightLabel{\Intro{\lor}}
    \UnaryInfC{$\varphi \lor \psi$}
  \end{prooftree}
  विगमन गृहीतकानुसार $\varphi$ हे $\delta_1$ च्या मुक्त न केलेली
  गृहीतकांपासून~$\Gamma$ निष्पन्न होते. एखादी
  संरचना~$\Struct{M}$
  विचारात घ्या. जर $\Sat{M}{\Gamma}$, तर
  $\Sat{M}{\varphi \lor \psi}$, हे दाखवायचे आहे.
  $\Sat{M}{\Gamma}$ असे समजा; मग
  $\Gamma \Entails \varphi$ (विगमन गृहीतक) असल्यामुळे
  $\Sat{M}{\varphi}$. म्हणून
  $\Sat{M}{\varphi \lor \psi}$ हेदेखील असलेच पाहिजे.
  ($\psi$ पासून $\varphi \lor \psi$ निष्पन्न करण्याचे प्रकरणही असेच हाताळता येते.)

\item शेवटचे अनुमान \Intro{\lif} आहे. $\varphi$ हे गृहीतक आणि $\psi$ हा निष्कर्ष
  असलेल्या उपसिद्धतेपासून $\varphi \lif \psi$ निष्पन्न केले जाते, म्हणजे:
  \begin{prooftree}
    \AxiomC{$\Gamma, \Discharge{\varphi}{n}$}
    \RightLabel{$\delta_1$}
    \DeduceC{$\psi$}
    \DischargeRule{\Intro{\lif}}{n}
    \UnaryInfC{$\varphi \lif \psi$}
  \end{prooftree}
  विगमन गृहीतकानुसार $\psi$ हे $\delta_1$ च्या मुक्त न केलेली
  गृहीतकांपासून निष्पन्न होते, म्हणजे $\Gamma \cup \{\varphi\} \Entails \psi$.
  एखादी
  संरचना~$\Struct{M}$
  विचारात घ्या. शेवटच्या अनुमानात $\varphi$ चे विसर्जन झाल्यामुळे $\delta$
  ची मुक्त न केलेली गृहीतके फक्त $\Gamma$ आहेत. म्हणून
  $\Gamma \Entails \varphi \lif \psi$ हे दाखवायचे आहे. विसंगतीद्वारे सिद्धतेसाठी
  असे समजा की एखाद्या
  संरचना~$\Struct{M}$
  साठी $\Sat{M}{\Gamma}$, पण
  $\Sat/{M}{\varphi \lif \psi}$. मग
  $\Sat{M}{\varphi}$ आणि
  $\Sat/{M}{\psi}$. पण गृहीतकानुसार $\psi$ हा
  $\Gamma \cup \{\varphi\}$ चा तार्किक परिणाम आहे, म्हणजे
  $\Sat{M}{\psi}$; ही विसंगती आहे. म्हणून
  $\Gamma \Entails \varphi \lif \psi$.

\item शेवटचे अनुमान \FalseInt आहे. येथे $\delta$ पुढीलप्रमाणे संपते:
  \begin{prooftree}
    \AxiomC{$\Gamma$}
    \RightLabel{$\delta_1$}
    \DeduceC{$\lfalse$}
    \RightLabel{\FalseInt}
    \UnaryInfC{$\varphi$}
  \end{prooftree}
  विगमन गृहीतकानुसार $\Gamma \Entails \lfalse$. आपल्याला
  $\Gamma \Entails \varphi$ हे दाखवायचे आहे. तसे नसल्याचे समजा; मग एखाद्या
  $\Struct{M}$ साठी
    $\Sat{M}{\Gamma}$ आणि
    $\Sat/{M}{\varphi}$. पण नेहमीच
    $\Sat/{M}{\lfalse}$ असल्यामुळे याचा अर्थ
    $\Gamma \Entails/ \lfalse$ असा होईल, जो विगमन गृहीतकाच्या विरुद्ध आहे.

\item शेवटचे अनुमान \FalseCl आहे: सराव.


\item शेवटचे अनुमान \Intro{\lforall} आहे. मग $\delta$ चे रूप पुढीलप्रमाणे आहे:
  \begin{prooftree}
    \AxiomC{$\Gamma$}
    \RightLabel{$\delta_1$}
    \DeduceC{$\varphi(a)$}
    \RightLabel{\Intro{\lforall}}
    \UnaryInfC{$\lforall[x][\varphi(x)]$}
  \end{prooftree}
  विगमन गृहीतकानुसार आधारविधान~$\varphi(a)$ हे मुक्त न केलेली
  गृहीतके~$\Gamma$ यांचा तार्किक परिणाम आहे.
  $\Sat{M}{\Gamma}$ अशी एखादी रचना~$\Struct{M}$ विचारात घ्या. आपल्याला
  $\Sat{M}{\lforall[x][\varphi(x)]}$ हे दाखवायचे आहे.
  $\lforall[x][\varphi(x)]$ हे वाक्य असल्यामुळे याचा अर्थ प्रत्येक
  चल-विन्यास~$s$ साठी $\Sat{M}{\varphi(x)}[s]$ दाखवायचे आहे
  (\readerexternalref{fol:syn:ass:prop:sat-quant}). $\Gamma$ मध्ये केवळ वाक्ये असल्यामुळे
  \readerexternalref{fol:syn:sat:defn:satisfaction} नुसार प्रत्येक $\psi \in \Gamma$ साठी
  $\Sat{M}{\psi}[s]$. $\Struct{M'}$ ही $\Struct{M}$ सारखीच रचना असू द्या,
  फक्त $\Assign{a}{M'} = s(x)$. $a$ हे $\Gamma$ मध्ये येत नसल्यामुळे
  \readerexternalref{fol:syn:ext:cor:extensionality-sent} नुसार
  $\Sat{M'}{\Gamma}$. $\Gamma \Entails \varphi(a)$ असल्यामुळे
  $\Sat{M'}{\varphi(a)}$. $\varphi(a)$ हे वाक्य असल्यामुळे
  \readerexternalref{fol:syn:ass:prop:sentence-sat-true} नुसार
  $\Sat{M'}{\varphi(a)}[s]$. \readerexternalref{fol:syn:ext:prop:ext-formulas} नुसार
  $\Sat{M'}{\varphi(x)}[s]$ तेव्हा आणि केवळ तेव्हाच, जेव्हा
  $\Sat{M'}{\varphi(a)}$ (लक्षात घ्या की $\varphi(a)$ म्हणजे
  $\Subst{\varphi(x)}{a}{x}$). म्हणून $\Sat{M'}{\varphi(x)}[s]$. $a$ हे $\varphi(x)$
  मध्ये येत नसल्यामुळे \readerexternalref{fol:syn:ext:prop:extensionality} नुसार
  $\Sat{M}{\varphi(x)}[s]$. पण $s$ हा स्वेच्छ चल-विन्यास होता, म्हणून
  $\Sat{M}{\lforall[x][\varphi(x)]}$.

\item शेवटचे अनुमान \Intro{\lexists} आहे: सराव.

\item शेवटचे अनुमान \Elim{\forall} आहे: सराव.

\end{enumerate}

आता अनेक आधारविधाने असलेली संभाव्य अनुमाने पाहू:
\Elim{\lor}, \Intro{\land}, \Elim{\lif} आणि
  \Elim{\lexists}.
\begin{enumerate}

\item शेवटचे अनुमान \Intro{\land} आहे. $\varphi$ आणि $\psi$ या आधारविधानांपासून
  $\varphi \land \psi$ निष्पन्न केले जाते आणि $\delta$ चे रूप पुढीलप्रमाणे आहे:
  \begin{prooftree}
    \AxiomC{$\Gamma_1$}
    \RightLabel{$\delta_1$}
    \DeduceC{$\varphi$}
    \AxiomC{$\Gamma_2$}
    \RightLabel{$\delta_2$}
    \DeduceC{$\psi$}
    \RightLabel{\Intro{\land}}
    \BinaryInfC{$\varphi \land \psi$}
  \end{prooftree}
  विगमन गृहीतकानुसार $\varphi$ हे $\delta_1$ च्या मुक्त न केलेली
  गृहीतकांपासून~$\Gamma_1$ निष्पन्न होते आणि $\psi$ हे $\delta_2$ च्या
  मुक्त न केलेली गृहीतकांपासून~$\Gamma_2$ निष्पन्न होते. $\delta$ ची
  मुक्त न केलेली गृहीतके $\Gamma_1 \cup \Gamma_2$ आहेत; म्हणून
  $\Gamma_1 \cup \Gamma_2 \Entails \varphi \land \psi$ हे दाखवायचे आहे.
  $\Sat{M}{\Gamma_1 \cup \Gamma_2}$ असलेली एखादी
  संरचना~$\Struct{M}$
  विचारात घ्या. $\Sat{M}{\Gamma_1}$ आणि
  $\Gamma_1 \Entails \varphi$ असल्यामुळे
  $\Sat{M}{\varphi}$ असलेच पाहिजे; तसेच
  $\Sat{M}{\Gamma_2}$ आणि
  $\Gamma_2 \Entails \psi$ असल्यामुळे
  $\Sat{M}{\psi}$. या दोन्हींवरून
  $\Sat{M}{\varphi \land \psi}$.

\item शेवटचे अनुमान \Elim{\lor} आहे: सराव.

\item शेवटचे अनुमान \Elim{\lif} आहे. $\varphi \lif \psi$ आणि $\varphi$ या
  आधारविधानांपासून $\psi$ निष्पन्न केले जाते. निष्पत्ती~$\delta$
  पुढीलप्रमाणे दिसते:
  \begin{prooftree}
    \AxiomC{$\Gamma_1$}
    \RightLabel{$\delta_1$}
    \DeduceC{$\varphi \lif \psi$}
    \AxiomC{$\Gamma_2$}
    \RightLabel{$\delta_2$}
    \DeduceC{$\varphi$}
    \RightLabel{\Elim{\lif}}
    \BinaryInfC{$\psi$}
  \end{prooftree}
  विगमन गृहीतकानुसार $\varphi \lif \psi$ हे $\delta_1$ च्या
  मुक्त न केलेली गृहीतकांपासून~$\Gamma_1$ निष्पन्न होते आणि $\varphi$ हे
  $\delta_2$ च्या मुक्त न केलेली गृहीतकांपासून~$\Gamma_2$ निष्पन्न होते.
  एखादी
  संरचना~$\Struct{M}$
  विचारात घ्या. जर $\Sat{M}{\Gamma_1 \cup
    \Gamma_2}$, तर $\Sat{M}{\psi}$, हे दाखवायचे आहे.
  $\Sat{M}{\Gamma_1 \cup \Gamma_2}$ असे समजा.
  $\Gamma_1 \Entails \varphi \lif \psi$ असल्यामुळे
  $\Sat{M}{\varphi \lif \psi}$. तसेच
  $\Gamma_2 \Entails \varphi$ असल्यामुळे
  $\Sat{M}{\varphi}$. याचा अर्थ
  $\Sat{M}{\psi}$ असा होतो (कारण
  $\Sat/{M}{\psi}$ असेल, तर
  $\Sat{M}{\varphi}$ असल्यामुळे
  $\Sat/{M}{\varphi \lif \psi}$ असे होईल, जे
  $\Sat{M}{\varphi \lif \psi}$ च्या विरुद्ध आहे).

\item शेवटचे अनुमान \Elim{\lnot} आहे: सराव.

\item शेवटचे अनुमान \Elim{\lexists} आहे: सराव.
\end{enumerate}
\end{proof}


\begin{prob}
\ref{fol:ntd:sou:thm:soundness} ची सिद्धता पूर्ण करा.
\end{prob}




\begin{cor}
\label{fol:ntd:sou:cor:weak-soundness}
जर $\Proves \varphi$, तर $\varphi$ हे वैध आहे.
\end{cor}

\begin{cor}
\label{fol:ntd:sou:cor:consistency-soundness}
जर $\Gamma$ पूर्ततायोग्य असेल, तर तो सुसंगत आहे.
\end{cor}

\begin{proof}
प्रतिलोम विधान सिद्ध करू. $\Gamma$ सुसंगत नाही असे समजा. मग
$\Gamma \Proves \lfalse$, म्हणजे $\Gamma$ मधील मुक्त न केलेली
गृहीतकांपासून $\lfalse$ ची निष्पत्ती आहे. \ref{fol:ntd:sou:thm:soundness}
नुसार $\Gamma$ पूर्ण करणारी कोणतीही
संरचना~$\Struct{M}$
$\lfalse$ देखील पूर्ण करते. प्रत्येक
संरचना~$\Struct{M}$
साठी $\Sat/{M}{\lfalse}$ असल्यामुळे कोणतीही
$\Struct{M}$ $\Gamma$ पूर्ण करू शकत नाही,
म्हणजे $\Gamma$ पूर्ततायोग्य नाही.
\end{proof}









\section{निष्पत्ती सह एकरूपता}\label{fol:ntd:ide:sec}

एकरूपता असलेल्या निष्पत्त्या साठी अतिरिक्त अनुमाननियम आवश्यक असतात.

\begin{defish}
\AxiomC{}
\RightLabel{\Intro{\eq}}
\UnaryInfC{$\eq[t][t]$}
\DisplayProof
\hfill
\begin{tabular}{r}
\AxiomC{$\eq[t_1][t_2]$}
\AxiomC{$\varphi(t_1)$}
\RightLabel{\Elim{\eq}}
\BinaryInfC{$\varphi(t_2)$}
\DisplayProof
\\[3ex]
\AxiomC{$\eq[t_1][t_2]$}
\AxiomC{$\varphi(t_2)$}
\RightLabel{\Elim{\eq}}
\BinaryInfC{$\varphi(t_1)$}
\DisplayProof
\end{tabular}
\end{defish}

वरील नियमांमध्ये $t$, $t_1$ आणि $t_2$ ही बंद पदे आहेत. \Intro{\eq} नियम
आपल्याला कोणत्याही गृहीतकांशिवाय, थेट $\eq[t][t]$ या रूपाचे कोणतेही
एकरूपता-विधान निष्पन्न करू देतो.

\begin{ex}
$s$ आणि $t$ ही बंद पदे असतील, तर $\varphi(s), \eq[s][t] \Proves \varphi(t)$:
\begin{prooftree}
\AxiomC{$\eq[s][t]$}
\AxiomC{$\varphi(s)$}
\RightLabel{$\Elim{\eq}$}
\BinaryInfC{$\varphi(t)$}
\end{prooftree}
याला “एकरूप वस्तूंच्या प्रतिस्थापनाचे तत्त्व” किंवा लाइब्निझचा नियम
म्हणून ओळखले जाऊ शकते.
\end{ex}

\begin{prob}
$=$ हे सममित आणि संक्रमक दोन्ही आहे हे सिद्ध करा; म्हणजे
$\lforall[x][\lforall[y][(\eq[x][y] \lif
    \eq[y][x])]]$ आणि $\lforall[x][\lforall[y][\lforall[z]((\eq[x][y]
    \land \eq[y][z]) \lif \eq[x][z])]]$ यांच्या
निष्पत्त्या द्या.
\end{prob}

\begin{ex}
आपण पुढील वाक्य निष्पन्न करतो:
\begin{align*}
& \lforall[x][\lforall[y][((\varphi(x) \land \varphi(y)) \lif \eq[x][y])]]
\intertext{या वाक्य पासून:}
& \lexists[x][\lforall[y][(\varphi(y) \lif \eq[y][x])]]
\end{align*}
निष्पत्ती उलट्या दिशेने विकसित करू:
\begin{prooftree}
\AxiomC{$\lexists[x][\lforall[y][(\varphi(y) \lif \eq[y][x])]]
\quad \Discharge{\varphi(a) \land \varphi(b)}{1}$}
\DeduceC{$\eq[a][b]$}
\DischargeRule{\Intro{\lif}}{1}
\UnaryInfC{$((\varphi(a) \land \varphi(b)) \lif \eq[a][b])$}
\RightLabel{\Intro{\lforall}}
\UnaryInfC{$\lforall[y][((\varphi(a) \land \varphi(y)) \lif \eq[a][y])]$}
\RightLabel{\Intro{\lforall}}
\UnaryInfC{$\lforall[x][\lforall[y][((\varphi(x) \land \varphi(y)) \lif \eq[x][y])]]$}
\end{prooftree}
आता मुख्य गृहीतक वापरावे लागेल: ते अस्तित्ववाची सूत्र असल्यामुळे
\Elim{\lexists} वापरून मध्यवर्ती निष्कर्ष~$\eq[a][b]$ निष्पन्न करू.
\begin{prooftree}
\AxiomC{$\lexists[x][\lforall[y][(\varphi(y) \lif \eq[y][x])]]$}
\AxiomC{$\Discharge{\lforall[y][(\varphi(y) \lif \eq[y][c]])}{2}$}
\noLine
\UnaryInfC{$\Discharge{\varphi(a) \land \varphi(b)}{1}$}
\DeduceC{$\eq[a][b]$}
\DischargeRule{\Elim{\lexists}}{2}
\BinaryInfC{$\eq[a][b]$}
\DischargeRule{\Intro{\lif}}{1}
\UnaryInfC{$((\varphi(a) \land \varphi(b)) \lif \eq[a][b])$}
\RightLabel{\Intro{\lforall}}
\UnaryInfC{$\lforall[y][((\varphi(a) \land \varphi(y)) \lif \eq[a][y])]$}
\RightLabel{\Intro{\lforall}}
\UnaryInfC{$\lforall[x][\lforall[y][((\varphi(x) \land \varphi(y)) \lif \eq[x][y])]]$}
\end{prooftree}
वरच्या उजव्या बाजूची उप-निष्पत्ती तिच्या गृहीतकांचा वापर करून
$\eq[a][c]$ आणि $\eq[b][c]$ दाखवल्यावर पूर्ण होते. यासाठी दोन स्वतंत्र
निष्पत्त्या आवश्यक आहेत. $\eq[a][c]$ साठीची निष्पत्ती
पुढीलप्रमाणे आहे:
\begin{prooftree}
\AxiomC{$\Discharge{\lforall[y][(\varphi(y) \lif \eq[y][c]])}{2}$}
\RightLabel{\Elim{\lforall}}
\UnaryInfC{$\varphi(a) \lif \eq[a][c]$}
\AxiomC{$\Discharge{\varphi(a) \land \varphi(b)}{1}$}
\RightLabel{\Elim{\land}}
\UnaryInfC{$\varphi(a)$}
\RightLabel{\Elim{\lif}}
\BinaryInfC{$\eq[a][c]$}
\end{prooftree}
$\eq[a][c]$ आणि $\eq[b][c]$ यांच्यापासून \Elim{\eq} वापरून
$\eq[a][b]$ निष्पन्न करतो.
\end{ex}

\begin{prob}
पुढील सूत्रांच्या निष्पत्त्या द्या:
\begin{enumerate}
\item $\lforall[x][\lforall[y][((\eq[x][y] \land \varphi(x)) \lif \varphi(y))]]$
\item $\lexists[x][\varphi(x)] \land \lforall[y][\lforall[z][((\varphi(y) \land
    \varphi(z)) \lif \eq[y][z])]] \lif \lexists[x][(\varphi(x) \land
  \lforall[y][(\varphi(y) \lif \eq[y][x])])]$
\end{enumerate}
\end{prob}









\section{एकरूपता सह निर्दोषता}\label{fol:ntd:sid:sec}

\begin{prop}
$\eq$ साठीच्या नियमांसह नैसर्गिक निगमन निर्दोष आहे.
\end{prop}

\begin{proof}
$\eq[t][t]$ या रूपातील कोणतेही सूत्र वैध आहे, कारण प्रत्येक
संरचना~$\Struct M$ साठी $\Sat{M}{\eq[t][t]}$. (लक्षात घ्या की
$t$ हे बंद पद आहे असे आपण गृहीत धरतो, म्हणजे त्यात कोणतेही चल नाही;
म्हणून चल-विन्यास अप्रासंगिक असतात.)

निष्पत्तीमधील शेवटचे अनुमान \Elim{\eq} आहे असे समजा, म्हणजे
निष्पत्तीचे रूप पुढीलप्रमाणे आहे:
\begin{prooftree}
  \AxiomC{$\Gamma_1$}
  \RightLabel{$\delta_1$}
  \DeduceC{$\eq[t_1][t_2]$}
  \AxiomC{$\Gamma_2$}
  \RightLabel{$\delta_2$}
  \DeduceC{$\varphi(t_1)$}
  \RightLabel{\Elim{\eq}}
  \BinaryInfC{$\varphi(t_2)$}
\end{prooftree}
आधारविधाने~$\eq[t_1][t_2]$ आणि $\varphi(t_1)$ ही अनुक्रमे
मुक्त न केलेली गृहीतके~$\Gamma_1$ आणि $\Gamma_2$ यांच्यापासून
निष्पन्न केली आहेत. $\varphi(t_2)$ हे $\Gamma_1 \cup \Gamma_2$ पासून
निष्पन्न होते, हे दाखवायचे आहे. $\Sat{M}{\Gamma_1 \cup
  \Gamma_2}$ असलेली संरचना~$\Struct{M}$ विचारात घ्या. विगमन
गृहीतकानुसार $\Sat{M}{\varphi(t_1)}$ आणि $\Sat{M}{\eq[t_1][t_2]}$.
त्यामुळे $\Value{t_1}{M} = \Value{t_2}{M}$. $s$ हा कोणताही चल-विन्यास
आणि $m = \Value{t_1}{M} = \Value{t_2}{M}$ असू द्या.
\readerexternalref{fol:syn:ext:prop:ext-formulas} नुसार
$\Sat{M}{\varphi(t_1)}[s]$ तेव्हा आणि केवळ तेव्हाच, जेव्हा
$\Sat{M}{\varphi(x)}[\Subst{s}{m}{x}]$, आणि हे तेव्हा आणि केवळ तेव्हाच, जेव्हा
$\Sat{M}{\varphi(t_2)}[s]$. $\Sat{M}{\varphi(t_1)}$ असल्यामुळे
$\Sat{M}{\varphi(t_2)}$.
\end{proof}



\chapter{टॅब्लो}\label{fol:tab::chap}\begingroup\setlength{\emergencystretch}{3em}\NewDocumentCommand{\sFmlaWide}{m m o}{\ensuremath{\IfNoValueTF{#3}{}{#3\,}\hbox to 1.2em{\ensuremath{#1}\hfil} #2}}
\begin{editorial}
हे प्रकरण चिन्हांकित सूत्रांवर आधारलेली विश्लेषणात्मक टॅब्लो पद्धत मांडते.

टॅब्लोशी सिद्धता-पद्धती म्हणून संबंधित सामग्री समाविष्ट किंवा
वगळण्यासाठी ``prfTab'' टॅग वापरा.
\end{editorial}









\section{नियम आणि टॅब्लो}\label{fol:tab:rul:sec}

टॅब्लो म्हणजे वाक्य संरचना मध्ये ज्या संभाव्य प्रकारे
सत्य अथवा असत्य असू शकते त्यांचा पद्धतशीर आढावा होय. टॅब्लोचे मूलभूत
घटक चिन्हांकित सूत्रे असतात: वाक्ये सोबत सत्यतामूल्याचे
``चिन्ह,'' म्हणजे $\True$ किंवा~$\False$. ही चिन्हांकित सूत्रे
(खाली वाढणाऱ्या) वृक्षात मांडली जातात.

\begin{defn}
  एक \emph{चिन्हांकित सूत्र} म्हणजे सत्यतामूल्य आणि वाक्य यांची जोडी,
  म्हणजे पुढीलपैकी एक:
  \[
  \sFmla{\True}{\varphi} \text{ किंवा } \sFmla{\False}{\varphi}.
  \]
\end{defn}

सहज अर्थानुसार, $\sFmla{\True}{\varphi}$ चे वाचन आपण ``$\varphi$ सत्य असू
शकते'' आणि $\sFmla{\False}{\varphi}$ चे वाचन ``$\varphi$ असत्य असू शकते''
(एखाद्या संरचना मध्ये) असे करू शकतो.

वृक्षातील प्रत्येक चिन्हांकित सूत्र एकतर \emph{गृहीतक} असते (ती
वृक्षाच्या अगदी वरच्या टोकाला नोंदवलेली असतात), किंवा तिच्या वर असलेल्या
चिन्हांकित सूत्राला अनेक निगमन नियमांपैकी एखादा नियम लावून ती मिळते.
पूर्ववर्ती सूत्राच्या प्रत्येक संभाव्य मुख्य संकारक साठी दोन नियम
असतात: चिन्ह~$\True$ असेल त्या स्थितीसाठी एक आणि चिन्ह~$\False$ असेल त्या
स्थितीसाठी एक. काही नियम वृक्षाची शाखा फोडू देतात, तर काही नियम शाखेत फक्त
चिन्हांकित सूत्रे जोडतात. एखादा नियम फक्त अगदी आधी येणाऱ्या
चिन्हांकित सूत्र वरच लागू करावा असे नाही; मुळापासून नियम लागू करण्याच्या
स्थानापर्यंतच्या शाखेतील कोणत्याही चिन्हांकित सूत्र वर तो लागू करता येतो
(आणि अनेकदा आधीच्या एखाद्यावर लावावाच लागतो).

एखाद्या शाखेत $\sFmla{\True}{\varphi}$ आणि $\sFmla{\False}{\varphi}$ ही दोन्ही
असतील तेव्हा ती शाखा \emph{बंद} असते. ज्याची प्रत्येक शाखा बंद असते तो
टॅब्लो बंद असतो. सहज अर्थानुसार, कोणतीही शाखा एक संयुक्त शक्यता
वर्णन करते; परंतु $\sFmla{\True}{\varphi}$ आणि $\sFmla{\False}{\varphi}$ एकाच वेळी
शक्य नाहीत. दुसऱ्या शब्दांत, एखादी शाखा बंद असेल तर तिने वर्णन केलेली शक्यता
बाद होते. विशेषतः याचा अर्थ असा की, बंद टॅब्लो हा
$\sFmla{\True}{\varphi}$ स्वरूपाचे प्रत्येक गृहीतक सत्य आणि
$\sFmla{\False}{\varphi}$ स्वरूपाचे प्रत्येक गृहीतक असत्य करणाऱ्या सर्व शक्यता
बाद करतो.

बंद टॅब्लो \emph{$\varphi$ साठीचा} म्हणजे ज्याचे
मूळ~$\sFmla{\False}{\varphi}$ आहे असा बंद टॅब्लो. असा बंद टॅब्लो
अस्तित्वात असल्यास, $\varphi$ असत्य असण्याच्या सर्व शक्यता बाद होतात; म्हणजेच
$\varphi$ प्रत्येक संरचना मध्ये सत्य असले पाहिजे.












\section{विधानीय नियम}\label{fol:tab:prl:sec}

\subsection{$\lnot$ साठीचे नियम}

\begin{defish}
\AxiomC{\sFmla{\True}{\lnot \varphi}}
\RightLabel{\TRule{\True}{\lnot}}
\UnaryInfC{\sFmla{\False}{\varphi}}
\DisplayProof
\hfill
\AxiomC{\sFmla{\False}{\lnot \varphi}}
\RightLabel{\TRule{\False}{\lnot}}
\UnaryInfC{\sFmla{\True}{\varphi}}
\DisplayProof
\end{defish}

\subsection{$\land$ साठीचे नियम}

\begin{defish}\noindent
\AxiomC{\sFmla{\True}{\varphi \land \psi}}
\RightLabel{\TRule{\True}{\land}}
\UnaryInfC{\sFmla{\True}{\varphi}}
\noLine
\UnaryInfC{\sFmla{\True}{\psi}}
\DisplayProof
\hfill
\AxiomC{\sFmla{\False}{\varphi \land \psi}}
\RightLabel{\TRule{\False}{\land}}
\UnaryInfC{$\sFmla{\False}{\varphi} \quad \mid \quad \sFmla{\False}{\psi}$}
\DisplayProof
\end{defish}

\subsection{$\lor$ साठीचे नियम}

\begin{defish}
\AxiomC{\sFmla{\True}{\varphi \lor \psi}}
\RightLabel{\TRule{\True}{\lor}}
\UnaryInfC{$\sFmla{\True}{\varphi} \quad \mid \quad \sFmla{\True}{\psi}$}
\DisplayProof
\hfill
\AxiomC{\sFmla{\False}{\varphi \lor \psi}}
\RightLabel{\TRule{\False}{\lor}}
\UnaryInfC{\sFmla{\False}{\varphi}}
\noLine
\UnaryInfC{\sFmla{\False}{\psi}}
\DisplayProof
\end{defish}

\subsection{$\lif$ साठीचे नियम}

\begin{defish}
\AxiomC{\sFmla{\True}{\varphi \lif \psi}}
\RightLabel{\TRule{\True}{\lif}}
\UnaryInfC{$\sFmla{\False}{\varphi} \quad \mid \quad \sFmla{\True}{\psi}$}
\DisplayProof
\hfill
\AxiomC{\sFmla{\False}{\varphi \lif \psi}}
\RightLabel{\TRule{\False}{\lif}}
\UnaryInfC{\sFmla{\True}{\varphi}}
\noLine
\UnaryInfC{\sFmla{\False}{\psi}}
\DisplayProof
\end{defish}

\subsection{कट नियम}

\begin{defish}
\AxiomC{}
\RightLabel{\Cut}
\UnaryInfC{$\sFmla{\True}{\varphi} \quad \mid \quad \sFmla{\False}{\varphi}$}
\DisplayProof
\end{defish}

\Cut{} नियम पूर्वीच्या चिन्हांकित सूत्र वर ``लावला'' जात नाही; त्याऐवजी
तो टॅब्लो मधील प्रत्येक शाखा दोन शाखांत फोडू देतो: एका शाखेत
$\sFmla{\True}{\varphi}$ आणि दुसरीत~$\sFmla{\False}{\varphi}$ असते. हा नियम आवश्यक
नाही---ज्या चिन्हांकित सूत्रांच्या संचासाठी बंद टॅब्लो आहे, त्या
संचासाठी \Cut{} न वापरणाराही एक बंद टॅब्लो असतो---परंतु त्यामुळे
टॅब्लो सोयीस्कर रीतीने एकत्र जोडता येतात.












\section{संख्यापकांचे नियम}\label{fol:tab:qrl:sec}

\subsection{$\lforall$ साठीचे नियम}

\begin{defish}
\AxiomC{\sFmla{\True}{\lforall[x][\varphi(x)]}}
\RightLabel{\TRule{\True}{\forall}}
\UnaryInfC{\sFmla{\True}{\varphi(t)}}
\DisplayProof
\hfill
\AxiomC{\sFmla{\False}{\lforall[x][\varphi(x)]}}
\RightLabel{\TRule{\False}{\lforall}}
\UnaryInfC{\sFmla{\False}{\varphi(a)}}
\DisplayProof
\end{defish}

\TRule{\True}{\lforall} मध्ये, $t$ हे बंद पद आहे (म्हणजे, चरे नसलेले
पद). \TRule{\False}{\lforall} मध्ये, $a$~हा असा स्थिरांक आहे की तो
\TRule{\False}{\lforall} नियमाच्या वरच्या शाखेत कोठेही येता कामा नये.
\TRule{\False}{\forall} अनुमानाचा $a$ हा \emph{आयगेन चर} आहे, असे आपण
म्हणतो.\footnote{वरील नियमात $a$ हा स्थिरांक असला तरी आपण
``आयगेन चर'' ही संज्ञा वापरतो. यामागे ऐतिहासिक कारणे आहेत.}

\subsection{$\lexists$ साठीचे नियम}

\begin{defish}
\AxiomC{\sFmla{\True}{\lexists[x][\varphi(x)]}}
\RightLabel{\TRule{\True}{\lexists}}
\UnaryInfC{\sFmla{\True}{\varphi(a)}}
\DisplayProof
\hfill
\AxiomC{\sFmla{\False}{\lexists[x][\varphi(x)]}}
\RightLabel{\TRule{\False}{\lexists}}
\UnaryInfC{\sFmla{\False}{\varphi(t)}}
\DisplayProof
\end{defish}

पुन्हा, $t$~हे बंद पद आहे आणि $a$~हा असा स्थिरांक आहे की तो
\TRule{\True}{\lexists} नियमाच्या वरच्या शाखेत येत नाही. आपण $a$ ला
\TRule{\True}{\lexists} अनुमानाचा \emph{आयगेन चर} म्हणतो.

\TRule{\False}{\lforall} किंवा \TRule{\True}{\lexists} अनुमानाच्या वरच्या
शाखेत आयगेन चर येत नसावा, या अटीला \emph{आयगेन चराची अट} म्हणतात.

\begin{explain}
ही प्रथा आठवा: $\varphi$ हे असे सूत्र असेल की त्यात
चल~$x$ मुक्त असेल, तर हे आपण~$\varphi(x)$ असे लिहून दर्शवतो.
त्याच संदर्भात, मग $\varphi(t)$ हा~$\Subst{\varphi}{t}{x}$ चा संक्षेप असतो.
त्यामुळे \TRule{\False}{\lexists} नियम पुढीलप्रमाणेही लिहिता येईल:
\begin{prooftree}
  \AxiomC{\sFmla{\False}{\lexists[x][\varphi]}}
  \RightLabel{\TRule{\False}{\lexists}}
  \UnaryInfC{\sFmla{\False}{\Subst{\varphi}{t}{x}}}
\end{prooftree}
लक्षात घ्या की~$\varphi$ मध्ये $t$ आधीच येऊ शकते; उदा., $\varphi$ हे
~$\Atom{\Obj P}{t,x}$ असू शकते. म्हणून,
$\sFmla{\False}{\lexists[x][\Atom{\Obj P}{t,x}]}$ पासून
$\sFmla{\False}{\Atom{\Obj P}{t,t}}$ निष्पन्न करणे हा
\TRule{\False}{\lexists} चा योग्य उपयोग आहे. परंतु
\TRule{\False}{\lforall} आणि~\TRule{\True}{\lexists} मधील आयगेन चराच्या
अटींनुसार स्थिरांक~$a$ हा~$\varphi$ मध्ये येता कामा नये. म्हणून,
$\sFmla{\False}{\lforall[x][\Atom{\Obj P}{a,x}]}$ पासून
$\sFmla{\False}{\Atom{\Obj P}{a,a}}$ हे
~$\TRule{\False}{\lforall}$ वापरून योग्य रीतीने निष्पन्न करता येत नाही.
\end{explain}

\begin{explain}
\TRule{\True}{\lforall} आणि \TRule{\False}{\lexists} मध्ये पद~$t$ बंद
असण्याव्यतिरिक्त त्यावर आणखी कोणतेही निर्बंध नाहीत. याउलट,
\TRule{\True}{\lexists} आणि \TRule{\False}{\lforall} नियमांमध्ये आयगेन
चराच्या अटीनुसार संबंधित अनुमानांच्या वरच्या शाखांमध्ये स्थिरांक~$a$
कोठेही येता कामा नये. ही पद्धत निर्दोष राहील याची खात्री करण्यासाठी ही अट
आवश्यक आहे. ही अट नसती, तर पुढील आकृती
$\lexists[x][\varphi(x)] \lif \lforall[x][\varphi(x)]$ साठीचा
बंद टॅब्लो ठरली असती:
\begin{center}
\begin{tableau}{}
  [\sFmla{\False}{\lexists[x][\varphi(x)] \lif \lforall[x][\varphi(x)]}, just=\TAss
    [\sFmla{\True}{\lexists[x][\varphi(x)]},
      just={\TRule{\False}{\lif}[1]}
      [\sFmla{\False}{\lforall[x][\varphi(x)]},
        just={\TRule{\False}{\lif}[1]}
        [\sFmla{\True}{\varphi(a)},
          just={\TRule{\True}{\lexists}[2]}
          [\sFmla{\False}{\varphi(a)},
            just={\TRule{\False}{\lforall}[3]}, close]
        ]
      ]
    ]
  ]
\end{tableau}
\end{center}
परंतु $\lexists[x][\varphi(x)] \lif
\lforall[x][\varphi(x)]$ हे वैध नाही.
\end{explain}












\section{टॅब्लो}\label{fol:tab:der:sec}

\begin{explain}
गृहीतक म्हणजे काय हे आपण सांगितले आहे आणि निगमन नियम दिले आहेत.
टॅब्लो यांपासून विगमनाधारित रीतीने निर्माण होतात: प्रत्येक
टॅब्लो ही एकतर एक किंवा अधिक गृहीतकांनी बनलेली एकच शाखा असते, किंवा
टॅब्लोच्या एखाद्या शाखेवर निगमन नियमांपैकी एक नियम लावल्याने ती मिळते.
\end{explain}

\begin{defn}[टॅब्लो]
$\sFmlaWide{S_1}{\varphi_1}$, \dots,\linebreak $\sFmlaWide{S_n}{\varphi_n}$ या गृहीतकांसाठीचा
टॅब्लो (येथे प्रत्येक $S_i$ हे $\True$ किंवा~$\False$ यांपैकी एक
आहे) म्हणजे पुढील अटी पूर्ण करणारा चिन्हांकित सूत्रांचा सांत वृक्ष होय:
\begin{enumerate}
\item वृक्षातील सर्वांत वरची $n$ चिन्हांकित सूत्रे म्हणजे एकाखाली एक
  असलेली\linebreak $\sFmlaWide{S_i}{\varphi_i}$ ही सूत्रे होत.
\item वृक्षातील गृहीतक नसलेली प्रत्येक चिन्हांकित सूत्र तिच्या वरच्या
  शाखेतील चिन्हांकित सूत्राला निगमन नियम योग्य रीतीने लावल्याने मिळते.
\end{enumerate}
टॅब्लोची एखादी शाखा $\sFmla{\True}{\varphi}$ आणि
~$\sFmla{\False}{\varphi}$ ही दोन्ही सूत्रे धारण करत असेल तर आणि तरच ती
\emph{बंद} असते; अन्यथा ती \emph{खुली} असते. ज्याची प्रत्येक शाखा बंद
असते तो टॅब्लो (त्याच्या गृहीतकांच्या संचासाठीचा)
\emph{बंद टॅब्लो} असतो. टॅब्लो बंद नसेल, म्हणजे त्यात किमान
एक खुली शाखा असेल, तर तो \emph{खुला} असतो.
\end{defn}

\begin{ex}
  गृहीतकांचा प्रत्येक संच स्वतःच टॅब्लो असतो, परंतु साधारणतः तो बंद
  नसतो. (स्पष्टपणे, तो बंद असेल तर गृहीतकांत $\sFmla{\True}{\varphi}$ आणि
  ~$\sFmla{\False}{\varphi}$ ही चिन्हांकित सूत्रांची जोडी आधीच असली पाहिजे.)

  एखाद्या टॅब्लो मधील चिन्हांकित सूत्र~$\varphi$ ला निगमन नियमांपैकी
  एक नियम लावून त्या टॅब्लोपासून (तो खुला असो वा बंद) नवा आणि अधिक मोठा
  टॅब्लो मिळवता येतो. नियम लागू केलेली~$\varphi$ ची आस्थिती ज्या ज्या शाखेत
  आहे, त्या प्रत्येक शाखेच्या शेवटी तो नियम एक किंवा अधिक
  चिन्हांकित सूत्रे जोडतो.

  उदाहरणार्थ, $\sFmla{\True}{\varphi \land \lnot \varphi}$ हे गृहीतक विचारात घ्या.
  केवळ त्या गृहीतकाने बनलेला (खुला) टॅब्लो पुढीलप्रमाणे आहे:
  \begin{center}
    \begin{tableau}{}
      [\sFmla{\True}{\varphi \land \lnot \varphi}, just=\TAss]
    \end{tableau}{}
  \end{center}
  त्या गृहीतकाला $\TRule{\True}{\land}$ नियम लावून त्यापासून नवा
  टॅब्लो मिळतो. हा नियम टॅब्लोमध्ये $\sFmla{\True}{\varphi}$ आणि
  $\sFmla{\True}{\lnot \varphi}$ या दोन नव्या ओळी जोडू देतो:
  \begin{center}
    \begin{tableau}{}
      [\sFmla{\True}{\varphi \land \lnot \varphi}, just=\TAss
        [\sFmla{\True}{\varphi}, just={\TRule{\True}{\land}[1]},
          [\sFmla{\True}{\lnot \varphi}, just={\TRule{\True}{\land}[1]}
          ]
        ]
      ]
    \end{tableau}{}
  \end{center}
  टॅब्लो लिहिताना आपण लावलेले नियम उजवीकडे नोंदवतो (उदा.,
  $\TRule{\True}{\land} 1$ याचा अर्थ त्या ओळीवरील चिन्हांकित सूत्र ही
  ओळ~$1$ वरील चिन्हांकित सूत्राला $\TRule{\True}{\land}$ नियम लावल्याने
  मिळाली आहे). या नव्या टॅब्लोमध्ये आता अधिक चिन्हांकित सूत्रे
  आहेत, परंतु त्यांपैकी फक्त एकीला ($\sFmla{\True}{\lnot \varphi}$) नियम लावता
  येतो (या स्थितीत $\TRule{\True}{\lnot}$ नियम). त्यामुळे पुढील बंद
  टॅब्लो मिळतो:
  \begin{center}
    \begin{tableau}{}
      [\sFmla{\True}{\varphi \land \lnot \varphi}, just=\TAss
        [\sFmla{\True}{\varphi}, just={\TRule{\True}{\land}[1]}
          [\sFmla{\True}{\lnot \varphi}, just={\TRule{\True}{\land}[1]}
            [\sFmla{\False}{\varphi}, just={\TRule{\True}{\lnot}[3]}, close]
          ]
        ]
      ]
    \end{tableau}{}
  \end{center}
\end{ex}












\section{टॅब्लोची उदाहरणे}\label{fol:tab:pro:sec}

\begin{ex}
$(\varphi \land \psi) \lif \varphi$ या वाक्य साठी बंद टॅब्लो शोधू या.

संबंधित गृहीतक टॅब्लोच्या सर्वांत वर लिहून आपण सुरुवात करतो.
\begin{oltableau}
  [\sFmla{\False}{(\varphi \land \psi) \lif \varphi},
    just = \TAss]
\end{oltableau}

एकच गृहीतक असल्यामुळे ज्याला नियम लावता येईल अशी एकच चिन्हांकित सूत्र
आहे. (प्रत्येक चिन्हांकित सूत्राला नेहमी जास्तीत जास्त एकच नियम लागू
करता येतो: तो म्हणजे वाक्याच्या संबंधित चिन्हाचा आणि
मुख्य संकारक चा नियम.) या स्थितीत याचा अर्थ असा की आपण
$\TRule{\False}{\lif}$ लावला पाहिजे.
\begin{oltableau}
  [\sFmla{\False}{(\varphi \land \psi) \lif \varphi},
    checked, just = \TAss
    [\sFmla{\True}{\varphi \land \psi},
      just={\TRule{\False}{\lif}[1]}
      [\sFmla{\False}{\varphi}, just={\TRule{\False}{\lif}[1]}]
    ]
  ]
\end{oltableau}
कोणत्या चिन्हांकित सूत्रे ना त्यांचे संबंधित नियम लावले आहेत याची नोंद
ठेवण्यासाठी आपण त्या वाक्याशेजारी खूण करतो. मात्र नियम सर्व खुल्या शाखांवर
लावला असेल \emph{तरच} खूण करा. एखाद्या चिन्हांकित सूत्राला प्रत्येक
खुल्या शाखेत संबंधित नियम लावल्यानंतर तिच्याकडे परत जाऊन तो नियम पुन्हा
लावावा लागणार नाही. येथे एकच शाखा असल्यामुळे नियम एकदाच लावावा लागतो.
(या खुणा केवळ टॅब्लो तयार करण्याची सोय आहेत आणि टॅब्लोच्या विन्यासाचा
अधिकृत भाग नाहीत, हे लक्षात घ्या.)

नियम लावता येईल अशी एक नवी चिन्हांकित सूत्र आहे: ओळ~$2$ वरील
$\sFmla{\True}{\varphi \land \psi}$. $\TRule{\True}{\land}$ नियम लावल्यावर
पुढील परिणाम मिळतो:
\begin{oltableau}
  [\sFmla{\False}{(\varphi \land \psi) \lif \varphi},
    checked, just = \TAss
    [\sFmla{\True}{\varphi \land \psi},
      just={\TRule{\False}{\lif}[1]}, checked
      [\sFmla{\False}{\varphi}, just={\TRule{\False}{\lif}[1]}
        [\sFmla{\True}{\varphi}, just={\TRule{\True}{\land}[2]}
          [\sFmla{\True}{\psi}, just={\TRule{\True}{\land}[2]}, close
          ]
        ]
      ]
    ]
  ]
\end{oltableau}
शाखेत आता $\sFmla{\True}{\varphi}$ (ओळ~$4$ वर) आणि
$\sFmla{\False}{\varphi}$ (ओळ~$3$ वर) ही दोन्ही असल्यामुळे शाखा बंद आहे.
ती एकमेव शाखा असल्यामुळे टॅब्लो बंद आहे. आपण
~$(\varphi \land \psi) \lif \varphi$ साठी बंद टॅब्लो शोधला आहे.
\end{ex}

\begin{ex}
आता $(\lnot \varphi \lor \psi) \lif (\varphi \lif \psi)$ साठी बंद टॅब्लो शोधू या.

संबंधित गृहीतकाने आपण सुरुवात करतो:
\begin{oltableau}
  [\sFmla{\False}{(\lnot \varphi \lor \psi) \lif
      (\varphi \lif \psi)}, just=\TAss]
\end{oltableau}
या टॅब्लो मधील एकमेव चिन्हांकित सूत्राचा मुख्य संकारक~$\lif$
आणि चिन्ह~$\False$ आहे; म्हणून तिला $\TRule{\False}{\lif}$ नियम लावल्यावर
पुढील परिणाम मिळतो:
\begin{oltableau}
  [\sFmla{\False}{(\lnot \varphi \lor \psi)
      \lif (\varphi \lif \psi)}, just=\TAss, checked
    [\sFmla{\True}{\lnot \varphi \lor \psi},
      just={\TRule{\False}{\lif}[1]}
      [\sFmla{\False}{(\varphi \lif \psi)},
        just={\TRule{\False}{\lif}[1]}
      ]
    ]
  ]
\end{oltableau}
आता ओळ~$2$ ला~$\TRule{\True}{\lor}$ लावायचा की ओळ~$3$ ला
$\TRule{\False}{\lif}$ लावायचा, असा पर्याय आपल्यापुढे आहे. प्रत्येक
चिन्हांकित सूत्राला तिचा संबंधित नियम प्रत्येक शाखेत लावला, तर आपण कोणता
क्रम निवडतो याने प्रत्यक्षात फरक पडत नाही. म्हणून पहिला पर्याय निवडू या.
$\TRule{\True}{\lor}$ नियमामुळे टॅब्लोची शाखा फुटू शकते आणि नियमाचे
दोन निष्कर्ष त्या दोन नव्या शाखांना जोडलेल्या नव्या चिन्हांकित सूत्रे
असतील. त्यामुळे पुढील परिणाम मिळतो:
\begin{oltableau}
  [\sFmla{\False}{(\lnot \varphi \lor \psi) \lif
      (\varphi \lif \psi)}, just=\TAss, checked
    [\sFmla{\True}{\lnot \varphi \lor \psi},
      just={\TRule{\False}{\lif}[1]}, checked
      [\sFmla{\False}{(\varphi \lif \psi)},
        just={\TRule{\False}{\lif}[1]}
        [\sFmla{\True}{\lnot \varphi}, just={\TRule{\True}{\lor}[2]}]
        [\sFmla{\True}{\psi}, just={\TRule{\True}{\lor}[2]}]
      ]
    ]
  ]
\end{oltableau}
ओळ~$3$ ला आपण अद्याप $\TRule{\False}{\lif}$ नियम लावलेला नाही; तो आता
लावू या. वेळ वाचवण्यासाठी तो दोन्ही शाखांवर लावू. प्रत्येक खुल्या शाखेत
संबंधित नियम लावला असेल तरच चिन्हांकित सूत्र शेजारी खूण करतो, हे आठवा.
म्हणून नियम ज्याला लागू होतो ती चिन्हांकित सूत्र असलेल्या प्रत्येक
शाखेच्या शेवटी तो नियम लावणे उचित ठरते. त्यामुळे विविध शाखांत खाली जाऊन त्या
चिन्हांकित सूत्र कडे पुन्हा परतावे लागत नाही.
\begin{oltableau}
  [\sFmla{\False}{(\lnot \varphi \lor \psi) \lif
      (\varphi \lif \psi)}, just=\TAss, checked
    [\sFmla{\True}{\lnot \varphi \lor \psi},
      just={\TRule{\False}{\lif}[1]}, checked
      [\sFmla{\False}{(\varphi \lif \psi)},
        just={\TRule{\False}{\lif}[1]}, checked
        [\sFmla{\True}{\lnot \varphi}, just={\TRule{\True}{\lor}[2]}
          [\sFmla{\True}{\varphi}, just={\TRule{\False}{\lif}[3]}
            [\sFmla{\False}{\psi}, just={\TRule{\False}{\lif}[3]}]
          ]
        ]
        [\sFmla{\True}{\psi}, just={\TRule{\True}{\lor}[2]}
          [\sFmla{\True}{\varphi}, just={\TRule{\False}{\lif}[3]}
            [\sFmla{\False}{\psi}, just={\TRule{\False}{\lif}[3]}, close]
          ]
        ]
      ]
    ]
  ]
\end{oltableau}
उजवी शाखा आता बंद आहे. डाव्या शाखेत आपण अजूनही ओळ~$4$ ला
$\TRule{\True}{\lnot}$ नियम लावू शकतो. त्यामुळे
~$\sFmla{\False}{\varphi}$ मिळते आणि डावी शाखा बंद होते:
\begin{oltableau}
  [\sFmla{\False}{(\lnot \varphi \lor \psi) \lif
      (\varphi \lif \psi)}, just=\TAss, checked
    [\sFmla{\True}{\lnot \varphi \lor \psi},
      just={\TRule{\False}{\lif}[1]}, checked
      [\sFmla{\False}{(\varphi \lif \psi)},
        just={\TRule{\False}{\lif}[1]}, checked
        [\sFmla{\True}{\lnot \varphi}, just={\TRule{\True}{\lor}[2]}
          [\sFmla{\True}{\varphi}, just={\TRule{\False}{\lif}[3]}
            [\sFmla{\False}{\psi}, just={\TRule{\False}{\lif}[3]}
              [\sFmla{\False}{\varphi},
                just={\TRule{\True}{\lnot}[4]}, close]
            ]
          ]
        ]
        [\sFmla{\True}{\psi}, just={\TRule{\True}{\lor}[2]}
          [\sFmla{\True}{\varphi}, just={\TRule{\False}{\lif}[3]}
            [\sFmla{\False}{\psi},
              just={\TRule{\False}{\lif}[3]}, close]
          ]
        ]
      ]
    ]
  ]
\end{oltableau}
\end{ex}

\begin{ex}
कितीही चिन्हांकित सूत्रे गृहीतके म्हणून घेऊन त्यांच्यासाठी टॅब्लो
देता येतात. शाखा फोडू देणारे एकाहून अधिक नियम लावणेही अनेकदा आवश्यक असते;
आणि सर्वसाधारणपणे टॅब्लोमध्ये कितीही शाखा असू शकतात. उदाहरणार्थ,
पुढील सूत्रांसाठीचा टॅब्लो विचारात घ्या:
$\{\sFmla{\True}{\varphi \lor (\psi \land \chi)}, \sFmla{\False}{(\varphi \lor \psi) \land (\varphi
  \lor \chi)}\}$. पहिल्या गृहीतकाला $\TRule{\True}{\lor}$ लावून आपण
सुरुवात करतो:
\begin{oltableau}
  [\sFmla{\True}{\varphi \lor (\psi \land \formula{C})},
    just=\TAss, checked
    [\sFmla{\False}{(\varphi \lor \psi) \land
        (\varphi \lor \formula{C})}, just=\TAss
      [\sFmla{\True}{\varphi}, just={\TRule{\True}{\lor}[1]}]
      [\sFmla{\True}{\psi \land \formula{C}},
        just={\TRule{\True}{\lor}[1]}]
    ]
  ]
\end{oltableau}
आता ओळ~$2$ ला $\TRule{\False}{\land}$ नियम लावता येतो. दोन्ही शाखांवर
तो एकाच वेळी लावल्यामुळे आपण ओळ~$2$ वर खूण करू शकतो:
\begin{oltableau}
  [\sFmla{\True}{\varphi \lor (\psi \land \formula{C})},
    just=\TAss, checked
    [\sFmla{\False}{(\varphi \lor \psi) \land
        (\varphi \lor \formula{C})}, just=\TAss, checked
      [\sFmla{\True}{\varphi}, just={\TRule{\True}{\lor}[1]}
        [\sFmla{\False}{\varphi \lor \psi},
          just={\TRule{\False}{\land}[2]}]
        [\sFmla{\False}{\varphi \lor \formula{C}},
          just={\TRule{\False}{\land}[2]}]
      ]
      [\sFmla{\True}{\psi \land \formula{C}},
        just={\TRule{\True}{\lor}[1]}
        [\sFmla{\False}{\varphi \lor \psi},
          just={\TRule{\False}{\land}[2]}]
        [\sFmla{\False}{\varphi \lor \formula{C}},
          just={\TRule{\False}{\land}[2]}]
      ]
    ]
  ]
\end{oltableau}
आता $\varphi \lor \psi$ असलेल्या सर्व शाखांवर $\TRule{\False}{\lor}$ लावता
येतो:
\begin{oltableau}
  [\sFmla{\True}{\varphi \lor (\psi \land \formula{C})},
    just=\TAss, checked
    [\sFmla{\False}{(\varphi \lor \psi) \land
        (\varphi \lor \formula{C})}, just=\TAss, checked
      [\sFmla{\True}{\varphi}, just={\TRule{\True}{\lor}[1]}
        [\sFmla{\False}{\varphi \lor \psi},
          just={\TRule{\False}{\land}[2]}, checked
          [\sFmla{\False}{\varphi}, just={\TRule{\False}{\lor}[4]}
            [\sFmla{\False}{\psi},
              just={\TRule{\False}{\lor}[4]}, close]
          ]
        ]
        [\sFmla{\False}{\varphi \lor \formula{C}},
          just={\TRule{\False}{\land}[2]}]
      ]
      [\sFmla{\True}{\psi \land \formula{C}},
        just={\TRule{\True}{\lor}[1]}
        [\sFmla{\False}{\varphi \lor \psi},
          just={\TRule{\False}{\land}[2]}, checked
          [\sFmla{\False}{\varphi}, just={\TRule{\False}{\lor}[4]}
            [\sFmla{\False}{\psi},
              just={\TRule{\False}{\lor}[4]}]
          ]
        ]
        [\sFmla{\False}{\varphi \lor \formula{C}},
          just={\TRule{\False}{\land}[2]}]
      ]
    ]
  ]
\end{oltableau}
सर्वांत डावी शाखा आता बंद आहे. आता $\varphi \lor \chi$ ला
$\TRule{\False}{\lor}$ लावू या:
\begin{oltableau}
  [\sFmla{\True}{\varphi \lor (\psi \land \formula{C})},
    just=\TAss, checked
    [\sFmla{\False}{(\varphi \lor \psi) \land
        (\varphi \lor \formula{C})}, just=\TAss, checked
      [\sFmla{\True}{\varphi}, just={\TRule{\True}{\lor}[1]}
        [\sFmla{\False}{\varphi \lor \psi},
          just={\TRule{\False}{\land}[2]}, checked
          [\sFmla{\False}{\varphi},
            just={\TRule{\False}{\lor}[4]}
            [\sFmla{\False}{\psi},
              just={\TRule{\False}{\lor}[4]}, close]
          ]
        ]
        [\sFmla{\False}{\varphi \lor \formula{C}},
          just={\TRule{\False}{\land}[2]}, checked
          [\sFmla{\False}{\varphi},
            just={\TRule{\False}{\lor}[4]}, move by=2
            [\sFmla{\False}{\formula{C}},
              just={\TRule{\False}{\lor}[4]}, close]
          ]
        ]
      ]
      [\sFmla{\True}{\psi \land \formula{C}},
        just={\TRule{\True}{\lor}[1]}
        [\sFmla{\False}{\varphi \lor \psi},
          just={\TRule{\False}{\land}[2]}, checked
          [\sFmla{\False}{\varphi},
            just={\TRule{\False}{\lor}[4]}
            [\sFmla{\False}{\psi},
              just={\TRule{\False}{\lor}[4]}
            ]
          ]
        ]
        [\sFmla{\False}{\varphi \lor \formula{C}},
          just={\TRule{\False}{\land}[2]}, checked
          [\sFmla{\False}{\varphi},
            just={\TRule{\False}{\lor}[4]}, move by=2
            [\sFmla{\False}{\formula{C}},
              just={\TRule{\False}{\lor}[4]}]
          ]
        ]
      ]
    ]
  ]
\end{oltableau}
स्पष्टतेसाठी $\TRule{\False}{\lor}$ दुसऱ्यांदा लावल्याचा परिणाम आपण खाली
हलवला आहे, हे लक्षात घ्या. येथे समर्थन समान असते, त्यामुळे तसे करण्याची
गरज नव्हती.

दोन शाखा खुल्या राहतात आणि ओळ~$3$ वरील $\sFmla{\True}{\psi \land \chi}$
वर खूण झालेली नाही. तिला $\TRule{\True}{\land}$ लावून पुढील बंद
टॅब्लो मिळतो:
\begin{oltableau}
  [\sFmla{\True}{\varphi \lor (\psi \land \formula{C})},
    just=\TAss, checked
    [\sFmla{\False}{(\varphi \lor \psi) \land
        (\varphi \lor \formula{C})}, just=\TAss, checked
      [\sFmla{\True}{\varphi}, just={\TRule{\True}{\lor}[1]}
        [\sFmla{\False}{\varphi \lor \psi},
          just={\TRule{\False}{\land}[2]}, checked
          [\sFmla{\False}{\varphi},
            just={\TRule{\False}{\lor}[4]}
            [\sFmla{\False}{\psi},
              just={\TRule{\False}{\lor}[4]}, close]
          ]
        ]
        [\sFmla{\False}{\varphi \lor \formula{C}},
          just={\TRule{\False}{\land}[2]}, checked
          [\sFmla{\False}{\varphi},
            just={\TRule{\False}{\lor}[4]}
            [\sFmla{\False}{\formula{C}},
              just={\TRule{\False}{\lor}[4]}, close]
          ]
        ]
      ]
      [\sFmla{\True}{\psi \land \formula{C}},
        just={\TRule{\True}{\lor}[1]}, checked
        [\sFmla{\False}{\varphi \lor \psi},
          just={\TRule{\False}{\land}[2]}, checked
          [\sFmla{\False}{\varphi},
            just={\TRule{\False}{\lor}[4]}
            [\sFmla{\False}{\psi},
              just={\TRule{\False}{\lor}[4]}
              [\sFmla{\True}{\psi},
                just={\TRule{\True}{\land}[3]}
                [\sFmla{\True}{\formula{C}},
                  just={\TRule{\True}{\land}[3]},close
                ]
              ]
            ]
          ]
        ]
        [\sFmla{\False}{\varphi \lor \formula{C}},
          just={\TRule{\False}{\land}[2]}, checked
          [\sFmla{\False}{\varphi},
            just={\TRule{\False}{\lor}[4]}
            [\sFmla{\False}{\formula{C}},
              just={\TRule{\False}{\lor}[4]}
              [\sFmla{\True}{\psi},
                just={\TRule{\True}{\land}[3]}
                [\sFmla{\True}{\formula{C}},
                  just={\TRule{\True}{\land}[3]},close
                ]
              ]
            ]
          ]
        ]
      ]
    ]
  ]
\end{oltableau}

तुलनेसाठी, त्याच गृहीतकांच्या संचासाठी नियम वेगळ्या क्रमाने लावलेला बंद
टॅब्लो पुढे दिला आहे:
\begin{oltableau}
  [\sFmla{\True}{\varphi \lor (\psi \land \formula{C})},
    just=\TAss, checked
    [\sFmla{\False}{(\varphi \lor \psi) \land
        (\varphi \lor \formula{C})}, just=\TAss, checked
      [\sFmla{\False}{\varphi \lor \psi},
        just={\TRule{\False}{\land}[2]}, checked
        [\sFmla{\False}{\varphi},
          just={\TRule{\False}{\lor}[3]}
          [\sFmla{\False}{\psi},
            just={\TRule{\False}{\lor}[3]}
            [\sFmla{\True}{\varphi},
              just={\TRule{\True}{\lor}[1]},close
            ]
            [\sFmla{\True}{\psi \land \formula{C}},
              just={\TRule{\True}{\lor}[1]}, checked
              [\sFmla{\True}{\psi},
                just={\TRule{\True}{\land}[6]}
                [\sFmla{\True}{\formula{C}},
                  just={\TRule{\True}{\land}[6]},close
                ]
              ]
            ]
          ]
        ]
      ]
      [\sFmla{\False}{\varphi \lor \formula{C}},
        just={\TRule{\False}{\land}[2]}, checked
        [\sFmla{\False}{\varphi},
          just={\TRule{\False}{\lor}[3]}
          [\sFmla{\False}{\formula{C}},
            just={\TRule{\False}{\lor}[3]}
            [\sFmla{\True}{\varphi},
              just={\TRule{\True}{\lor}[1]},close
            ]
            [\sFmla{\True}{\psi \land \formula{C}},
              just={\TRule{\True}{\lor}[1]}, checked
              [\sFmla{\True}{\psi},
                just={\TRule{\True}{\land}[6]}
                [\sFmla{\True}{\formula{C}},
                  just={\TRule{\True}{\land}[6]},close
                ]
              ]
            ]
          ]
        ]
      ]
    ]
  ]
\end{oltableau}
\end{ex}

\begin{prob}
पुढील सूत्रांचे बंद टॅब्लो द्या:
\begin{enumerate}
\item $\sFmla{\True}{\varphi \land (\psi \land \chi)}, \sFmla{\False}{(\varphi \land \psi) \land \chi}$.
\item $\sFmla{\True}{\varphi \lor (\psi \lor \chi)}, \sFmla{\False}{(\varphi \lor \psi) \lor \chi}$.
\item $\sFmla{\True}{\varphi \lif (\psi \lif \chi)}, \sFmla{\False}{\psi \lif (\varphi \lif \chi)}$.
\item $\sFmla{\True}{\varphi}, \sFmla{\False}{\lnot\lnot \varphi}$.
\end{enumerate}
\end{prob}

\begin{prob}
पुढील सूत्रांचे बंद टॅब्लो द्या:
\begin{enumerate}
\item $\sFmla{\True}{(\varphi \lor \psi) \lif \chi}, \sFmla{\False}{\varphi \lif \chi}$.
\item $\sFmla{\True}{(\varphi \lif \chi) \land (\psi \lif \chi)}, \sFmla{\False}{(\varphi \lor \psi) \lif \chi}$.
\item $\sFmla{\False}{\lnot(\varphi \land \lnot \varphi)}$.
\item $\sFmla{\True}{\psi \lif \varphi}, \sFmla{\False}{\lnot \varphi \lif \lnot \psi}$.
\item $\sFmla{\False}{(\varphi \lif \lnot \varphi) \lif \lnot \varphi}$.
\item $\sFmla{\False}{\lnot(\varphi \lif \psi) \lif \lnot \psi}$.
\item $\sFmla{\True}{\varphi \lif \chi}, \sFmla{\False}{\lnot (\varphi \land \lnot \chi)}$.
\item $\sFmla{\True}{\varphi \land \lnot \chi}, \sFmla{\False}{\lnot (\varphi \lif \chi)}$.
\item $\sFmla{\True}{\varphi \lor \psi, \lnot \psi}, \sFmla{\False}{\varphi}$.
\item $\sFmla{\True}{\lnot \varphi \lor \lnot \psi}, \sFmla{\False}{\lnot(\varphi \land \psi)}$.
\item $\sFmla{\False}{(\lnot \varphi \land \lnot \psi) \lif\lnot(\varphi \lor \psi)}$.
\item $\sFmla{\False}{\lnot(\varphi \lor \psi) \lif (\lnot \varphi \land \lnot \psi)}$.
\end{enumerate}
\end{prob}

\begin{prob}
पुढील सूत्रांचे बंद टॅब्लो द्या:
\begin{enumerate}
\item $\sFmla{\True}{\lnot(\varphi \lif \psi)}, \sFmla{\False}{\varphi}$.
\item $\sFmla{\True}{\lnot(\varphi \land \psi)}, \sFmla{\False}{\lnot \varphi \lor \lnot \psi}$.
\item $\sFmla{\True}{\varphi \lif \psi}, \sFmla{\False}{\lnot \varphi \lor \psi}$.
\item $\sFmla{\False}{\lnot \lnot \varphi \lif \varphi}$.
\item $\sFmla{\True}{\varphi \lif \psi}, \sFmla{\True}{\lnot \varphi \lif \psi}, \sFmla{\False}{\psi}$.
\item $\sFmla{\True}{(\varphi \land \psi) \lif \chi}, \sFmla{\False}{(\varphi \lif \chi) \lor (\psi \lif \chi)}$.
\item $\sFmla{\True}{(\varphi \lif \psi) \lif \varphi}, \sFmla{\False}{\varphi}$.
\item $\sFmla{\False}{(\varphi \lif \psi) \lor (\psi \lif \chi)}$.
\end{enumerate}
\end{prob}












\section{संख्यापकांसह टॅब्लो}\label{fol:tab:prq:sec}

\begin{ex}
संख्यापक हाताळताना आयगेन चराच्या अटीचा भंग होणार नाही याची खात्री करावी
लागते; त्यासाठी कधीकधी काही अनुमाने करण्याचा क्रम बदलून पाहावा लागतो.
सामान्यतः आयगेन चराच्या अटीच्या अधीन असलेले नियम आधी हाताळण्याचा प्रयत्न
उपयोगी ठरतो (पूर्ण टॅब्लोमध्ये ते अधिक वर असतील).

$\lexists[x][\lnot \varphi(x)] \lif \lnot \lforall[x][\varphi(x)]$ या
वाक्य साठी टॅब्लो कसा द्यायचा ते पाहू. नेहमीप्रमाणे संबंधित
गृहीतक नोंदवून आपण सुरुवात करतो:
\begin{oltableau}
[\sFmla{\False}{\lexists[x][\lnot \varphi(x)]\lif \lnot
    \lforall[x][\varphi(x)]}, just=\TAss]
\end{oltableau}
मुख्य संकारक $\lif$ असल्यामुळे आपण $\TRule{\False}{\lif}$ लावतो:
\begin{oltableau}
[\sFmla{\False}{\lexists[x][\lnot \varphi(x)]\lif \lnot
    \lforall[x][\varphi(x)]}, just=\TAss, checked
  [\sFmla{\True}{\lexists[x][\lnot \varphi(x)]},
    just={\TRule{\False}{\lif}[1]}
    [\sFmla{\False}{\lnot\lforall[x][\varphi(x)]},
      just={\TRule{\False}{\lif}[1]}]
  ]
]
\end{oltableau}
यानंतर ओळ~$2$ हाताळायची आहे. आपण $\TRule{\True}{\lexists}$ वापरतो.
त्यासाठी नवा स्थिरांक आवश्यक आहे; अद्याप कोणतेही स्थिरांक आलेले
नसल्यामुळे कोणताही एक, समजा~$a$, निवडता येतो.
\begin{oltableau}
[\sFmla{\False}{\lexists[x][\lnot \varphi(x)]\lif \lnot
    \lforall[x][\varphi(x)]}, just=\TAss, checked
  [\sFmla{\True}{\lexists[x][\lnot \varphi(x)]},
    just={\TRule{\False}{\lif}[1]}, checked
    [\sFmla{\False}{\lnot\lforall[x][\varphi(x)]},
      just={\TRule{\False}{\lif}[1]}
      [\sFmla{\True}{\lnot\varphi(a)}, just={\TRule{\True}{\lexists}[2]}]
    ]
  ]
]
\end{oltableau}
आता ओळ~$3$ ला $\TRule{\False}{\lnot}$ लावतो:
\begin{oltableau}
[\sFmla{\False}{\lexists[x][\lnot \varphi(x)]\lif \lnot
    \lforall[x][\varphi(x)]}, just=\TAss, checked
  [\sFmla{\True}{\lexists[x][\lnot \varphi(x)]},
    just={\TRule{\False}{\lif}[1]}, checked
    [\sFmla{\False}{\lnot\lforall[x][\varphi(x)]},
      just={\TRule{\False}{\lif}[1]}, checked
      [\sFmla{\True}{\lnot\varphi(a)}, just={\TRule{\True}{\lexists}[2]}
        [\sFmla{\True}{\lforall[x][\varphi(x)]},
          just = {\TRule{\False}{\lnot}[3]}
        ]
      ]
    ]
  ]
]
\end{oltableau}
ओळ~$4$ ला $\TRule{\True}{\lnot}$ आणि त्यानंतर ओळ~$5$ ला
$\TRule{\True}{\lforall}$ लावल्यावर बंद टॅब्लो मिळतो.
\begin{oltableau}
[\sFmla{\False}{\lexists[x][\lnot \varphi(x)]\lif \lnot
    \lforall[x][\varphi(x)]}, just=\TAss, checked
  [\sFmla{\True}{\lexists[x][\lnot \varphi(x)]},
    just={\TRule{\False}{\lif}[1]}, checked
    [\sFmla{\False}{\lnot\lforall[x][\varphi(x)]},
      just={\TRule{\False}{\lif}[1]}, checked
      [\sFmla{\True}{\lnot\varphi(a)}, just={\TRule{\True}{\lexists}[2]}
        [\sFmla{\True}{\lforall[x][\varphi(x)]},
          just = {\TRule{\False}{\lnot}[3]}
          [\sFmla{\False}{\varphi(a)}, just={\TRule{\True}{\lnot}[4]}
            [\sFmla{\True}{\varphi(a)},
              just={\TRule{\True}{\lforall}[5]}, close
            ]
          ]
        ]
      ]
    ]
  ]
]
\end{oltableau}
\end{ex}

\begin{ex}
पुढील संचासाठी टॅब्लो कसा द्यायचा ते पाहू:
\[
\sFmla{\False}{\lexists[x][\chi(x,b)]},
\sFmla{\True}{\lexists[x][(\varphi(x) \land \psi(x))]},
\sFmla{\True}{\lforall[x][(\psi(x) \lif \chi(x,b))]}.
\]
नेहमीप्रमाणे गृहीतकांनी सुरुवात करतो:
\begin{oltableau}
  [\sFmla{\False}{\lexists[x][\formula{C}(x,b)]}, just=\TAss
    [\sFmla{\True}{\lexists[x][(\varphi(x) \land \psi(x))]},
      just=\TAss
      [\sFmla{\True}{\lforall[x][(\psi(x) \lif \formula{C}(x,b))]},
        just=\TAss]
    ]
  ]
\end{oltableau}
आयगेन चराची अट असलेला नियम आपण नेहमी आधी लावला पाहिजे; येथे तो म्हणजे
ओळ~$2$ ला $\TRule{\True}{\lexists}$. गृहीतकांत स्थिरांक~$b$ असल्यामुळे
वेगळा अचर वापरावा लागतो; पुन्हा~$a$ निवडू या.
\begin{oltableau}
  [\sFmla{\False}{\lexists[x][\formula{C}(x,b)]}, just=\TAss
    [\sFmla{\True}{\lexists[x][(\varphi(x) \land \psi(x))]},
      just=\TAss, checked
      [\sFmla{\True}{\lforall[x][(\psi(x) \lif \formula{C}(x,b))]},
        just=\TAss
        [\sFmla{\True}{\varphi(a) \land \psi(a)},
          just={\TRule{\True}{\lexists}[2]}]
      ]
    ]
  ]
\end{oltableau}
आता ओळ~$1$ ला $\TRule{\False}{\lexists}$ किंवा ओळ~$3$ ला
$\TRule{\True}{\lforall}$ लावल्यास, $x$ च्या जागी कोणते बंद पद~$t$ ठेवायचे
ते ठरवावे लागते. या नियमांना आयगेन चराची अट नसल्यामुळे हवे ते कोणतेही बंद
पद निवडता येते. काही स्थितींमध्ये वेगवेगळी~$t$ वापरून नियम अनेकदा लावावा
लागू शकतो. मात्र सर्वसाधारणपणे टॅब्लोमध्ये आधीच आलेल्या पदांपैकी
एक---येथे $a$ किंवा~$b$---निवडणे फायदेशीर ठरते; आणि येथे $a$ मुळे बंद
शाखा मिळण्याची शक्यता अधिक आहे असा अंदाज करता येतो.
\begin{oltableau}
  [\sFmla{\False}{\lexists[x][\formula{C}(x,b)]}, just=\TAss
    [\sFmla{\True}{\lexists[x][(\varphi(x) \land \psi(x))]},
      just=\TAss, checked
      [\sFmla{\True}{\lforall[x][(\psi(x) \lif \formula{C}(x,b))]}, just=\TAss
        [\sFmla{\True}{\varphi(a) \land \psi(a)}, just={\TRule{\True}{\lexists}[2]}
          [\sFmla{\False}{\formula{C}(a,b)}, just={\TRule{\False}{\lexists}[1]}
            [\sFmla{\True}{\psi(a) \lif \formula{C}(a,b)},
              just={\TRule{\True}{\lforall}[3]}
            ]
          ]
        ]
      ]
    ]
  ]
\end{oltableau}
ओळी $1$ आणि~$3$ वरील चिन्हांकित सूत्रे पुन्हा वापराव्या लागू शकतात,
म्हणून त्यांवर खूण करत नाही. आता ओळ~$4$ ला $\TRule{\True}{\land}$ लावा:
\begin{oltableau}
  [\sFmla{\False}{\lexists[x][\formula{C}(x,b)]}, just=\TAss
    [\sFmla{\True}{\lexists[x][(\varphi(x) \land \psi(x))]},
      just=\TAss, checked
      [\sFmla{\True}{\lforall[x][(\psi(x) \lif \formula{C}(x,b))]}, just=\TAss
        [\sFmla{\True}{\varphi(a) \land \psi(a)},
          just={\TRule{\True}{\lexists}[2]}, checked
          [\sFmla{\False}{\formula{C}(a,b)}, just={\TRule{\False}{\lexists}[1]}
            [\sFmla{\True}{\psi(a) \lif \formula{C}(a,b)},
              just={\TRule{\True}{\lforall}[3]}
              [\sFmla{\True}{\varphi(a)}, just={\TRule{\True}{\land}[4]}
                [\sFmla{\True}{\psi(a)}, just={\TRule{\True}{\land}[4]}
                ]
              ]
            ]
          ]
        ]
      ]
    ]
  ]
\end{oltableau}
आता ओळ~$6$ ला $\TRule{\True}{\lif}$ लावल्यास टॅब्लो बंद होतो:
\begin{oltableau}
  [\sFmla{\False}{\lexists[x][\formula{C}(x,b)]}, just=\TAss
    [\sFmla{\True}{\lexists[x][(\varphi(x) \land \psi(x))]},
      just=\TAss, checked
      [\sFmla{\True}{\lforall[x][(\psi(x) \lif \formula{C}(x,b))]},
        just=\TAss
        [\sFmla{\True}{\varphi(a) \land \psi(a)},
          just={\TRule{\True}{\lexists}[2]}, checked
          [\sFmla{\False}{\formula{C}(a,b)},
            just={\TRule{\False}{\lexists}[1]}
            [\sFmla{\True}{\psi(a) \lif \formula{C}(a,b)},
              just={\TRule{\True}{\lforall}[3]}, checked
              [\sFmla{\True}{\varphi(a)},
                just={\TRule{\True}{\land}[4]}
                [\sFmla{\True}{\psi(a)},
                  just={\TRule{\True}{\land}[4]}
                  [\sFmla{\False}{\psi(a)},
                    just={\TRule{\True}{\lif}[6]},close
                  ]
                  [\sFmla{\True}{\formula{C}(a,b)},
                    just={\TRule{\True}{\lif}[6]},close
                  ]
                ]
              ]
            ]
          ]
        ]
      ]
    ]
  ]
\end{oltableau}


\end{ex}

\begin{ex}
पुढील संचासाठी आपण टॅब्लो तयार करतो:
\[
\sFmla{\True}{\lforall[x][\varphi(x)]}, \sFmla{\True}{\lforall[x][\varphi(x)]
  \lif \lexists[y][\psi(y)]}, \sFmla{\True}{\lnot\lexists[y][\psi(y)]}.
\]
नेहमीप्रमाणे गृहीतके लिहून सुरुवात करतो:
\begin{oltableau}
  [\sFmla{\True}{\lforall[x][\varphi(x)]}, just=\TAss
    [\sFmla{\True}{\lforall[x][\varphi(x)] \lif
        \lexists[y][\psi(y)]}, just=\TAss
      [\sFmla{\True}{\lnot\lexists[y][\psi(y)]}, just=\TAss
      ]
    ]
  ]
\end{oltableau}
ओळ~$3$ ला $\TRule{\True}{\lnot}$ नियम लावून सुरुवात करतो. ``आयगेन
चराची अट असलेले नियम नेहमी आधी लावा'' या नियमाचा एक परिणाम असा की
``आयगेन चराची अट नसलेले संख्यापक-नियम आवश्यक होईपर्यंत पुढे ढकला.''
शाखा फोडणारे नियमही पुढे ढकला.
\begin{oltableau}
  [\sFmla{\True}{\lforall[x][\varphi(x)]}, just=\TAss
    [\sFmla{\True}{\lforall[x][\varphi(x)] \lif
        \lexists[y][\psi(y)]}, just=\TAss
      [\sFmla{\True}{\lnot\lexists[y][\psi(y)]}, just=\TAss, checked
        [\sFmla{\False}{\lexists[y][\psi(y)]}, just={\TRule{\True}{\lnot}[3]}]
      ]
    ]
  ]
\end{oltableau}
नव्या ओळ~$4$ साठी $\TRule{\False}{\lexists}$ आवश्यक आहे आणि हा आयगेन
चराची अट नसलेला संख्यापक-नियम आहे. म्हणून तो पुढे ढकलून त्याऐवजी ओळ~$2$
वर $\TRule{\True}{\lif}$ वापरतो.
\begin{oltableau}
  [\sFmla{\True}{\lforall[x][\varphi(x)]}, just=\TAss
    [\sFmla{\True}{\lforall[x][\varphi(x)] \lif
        \lexists[y][\psi(y)]}, just=\TAss, checked
      [\sFmla{\True}{\lnot\lexists[y][\psi(y)]}, just=\TAss, checked
        [\sFmla{\False}{\lexists[y][\psi(y)]},
          just={\TRule{\True}{\lnot}[3]},
          [\sFmla{\False}{\lforall[x][\varphi(x)]}, just={\TRule{\True}{\lif}[2]}]
          [\sFmla{\True}{\lexists[y][\psi(y)]}, just={\TRule{\True}{\lif}[2]}]
        ]
      ]
    ]
  ]
\end{oltableau}
दोन्ही नव्या चिन्हांकित सूत्रे ना आयगेन चराच्या अटी असलेले नियम आवश्यक
आहेत; म्हणून ते नियम यानंतर लावावेत:
\begin{oltableau}
  [\sFmla{\True}{\lforall[x][\varphi(x)]}, just=\TAss
    [\sFmla{\True}{\lforall[x][\varphi(x)] \lif
        \lexists[y][\psi(y)]}, just=\TAss, checked
      [\sFmla{\True}{\lnot\lexists[y][\psi(y)]}, just=\TAss, checked
        [\sFmla{\False}{\lexists[y][\psi(y)]},
          just={\TRule{\True}{\lnot}[3]}
          [\sFmla{\False}{\lforall[x][\varphi(x)]}, just={\TRule{\True}{\lif}[2]},checked
            [\sFmla{\False}{\varphi(b)}, just={\TRule{\False}{\lforall}[5]}]
          ]
          [\sFmla{\True}{\lexists[y][\psi(y)]}, just={\TRule{\True}{\lif}[2]},checked
            [\sFmla{\True}{\psi(c)}, just={\TRule{\True}{\lexists}[5]}]
          ]
        ]
      ]
    ]
  ]
\end{oltableau}
शाखा बंद करण्यासाठी ओळी $1$ आणि~$4$ वरील चिन्हांकित सूत्रे वापराव्या
लागतात. त्यांच्या संबंधित नियमांना (\TRule{\True}{\lforall} आणि
\TRule{\False}{\lexists}) आयगेन चराची अट नाही; म्हणून योग्य असलेली कोणतीही
बंद पदे निवडण्याचे स्वातंत्र्य आहे. येथे ती अनुक्रमे $b$ आणि~$c$ आहेत.
\begin{oltableau}
  [\sFmla{\True}{\lforall[x][\varphi(x)]}, just=\TAss
    [\sFmla{\True}{\lforall[x][\varphi(x)] \lif
        \lexists[y][\psi(y)]}, just=\TAss, checked
      [\sFmla{\True}{\lnot\lexists[y][\psi(y)]}, just=\TAss, checked
        [\sFmla{\False}{\lexists[y][\psi(y)]},
          just={\TRule{\True}{\lnot}[3]}
          [\sFmla{\False}{\lforall[x][\varphi(x)]},
            just={\TRule{\True}{\lif}[2]},checked
            [\sFmla{\False}{\varphi(b)},
              just={\TRule{\False}{\lforall}[5]}
              [\sFmla{\True}{\varphi(b)},
                just={\TRule{\True}{\lforall}[1]},close
              ]
            ]
          ]
          [\sFmla{\True}{\lexists[y][\psi(y)]},
            just={\TRule{\True}{\lif}[2]},checked
            [\sFmla{\True}{\psi(c)},
              just={\TRule{\True}{\lexists}[5]}
              [\sFmla{\False}{\psi(c)},
                just={\TRule{\False}{\lexists}[4]},close
              ]
            ]
          ]
        ]
      ]
    ]
  ]
\end{oltableau}
\end{ex}

\begin{prob}
पुढील सूत्रांचे बंद टॅब्लो द्या:
\begin{enumerate}
\item $\sFmla{\False}{(\lforall[x][\varphi(x)] \land \lforall[y][\psi(y)])
\lif \lforall[z][(\varphi(z) \land \psi(z))]}$.
\item $\sFmla{\False}{(\lexists[x][\varphi(x)] \lor \lexists[y][\psi(y)])
\lif \lexists[z][(\varphi(z) \lor \psi(z))]}$.
\item $\sFmla{\True}{\lforall[x][(\varphi(x) \lif \psi)]},
\sFmla{\False}{\lexists[y][\varphi(y)] \lif \psi}$.
\item $\sFmla{\True}{\lforall[x][\lnot \varphi(x)]},
\sFmla{\False}{\lnot\lexists[x][\varphi(x)]}$.
\item $\sFmla{\False}{\lnot\lexists[x][\varphi(x)] \lif \lforall[x][\lnot
\varphi(x)]}$.
\item $\sFmla{\False}{\lnot\lexists[x][\lforall[y][((\varphi(x,y) \lif
\lnot \varphi(y,y)) \land (\lnot \varphi(y,y) \lif \varphi(x,y)))]]}$.
\end{enumerate}
\end{prob}

\begin{prob}
पुढील सूत्रांचे बंद टॅब्लो द्या:
\begin{enumerate}
\item $\sFmla{\False}{\lnot\lforall[x][\varphi(x)] \lif \lexists[x][\lnot\varphi(x)]}$.
\item $\sFmla{\True}{(\lforall[x][\varphi(x)] \lif \psi)}, \sFmla{\False}{\lexists[y][(\varphi(y) \lif \psi)]}$.
\item $\sFmla{\False}{\lexists[x][(\varphi(x) \lif \lforall[y][\varphi(y)])]}$.
\end{enumerate}
\end{prob}












\section{सिद्धता-उपपत्तीय संकल्पना}\label{fol:tab:ptn:sec}

\begin{editorial}
या विभागात टॅब्लोसाठी सिद्धतायोग्यता संबंध आणि सुसंगतता यांच्या व्याख्या
एकत्रित केल्या आहेत.
\end{editorial}

\begin{explain}
जशा आपण अनेक महत्त्वाच्या चिन्हार्थविषयक संकल्पना (वैधता, तार्किक निष्पन्नता,
पूर्ततायोग्यता) परिभाषित केल्या आहेत, तशाच त्यांना अनुरूप
\emph{सिद्धता-उपपत्तीय संकल्पना} आता परिभाषित करू. संरचना मध्ये
वाक्यांची पूर्ति होते की नाही याचा आधार घेऊन या संकल्पना परिभाषित
केल्या जात नाहीत; त्याऐवजी ठरावीक बंद टॅब्लोx च्या अस्तित्वाचा आधार
घेतला जातो. या दोन प्रकारच्या संकल्पना एकमेकांशी जुळतात, हा एक महत्त्वाचा
शोध होता. त्या जुळतात, हाच \emph{निर्दोषता प्रमेय} आणि
\emph{संपूर्णता प्रमेय} यांचा आशय आहे.
\end{explain}


\begin{defn}[प्रमेये]
एखादे वाक्य~$\varphi$ हे \emph{प्रमेय} असते, जर
$\sFmla{\False}{\varphi}$ साठी बंद टॅब्लो असेल. $\varphi$ हे प्रमेय असेल
तर आपण $\Proves \varphi$ लिहितो आणि ते प्रमेय नसेल तर $\Proves/ \varphi$ लिहितो.
\end{defn}

\begin{defn}[निष्पन्नता]
वाक्य $\varphi$ हे वाक्यांच्या संच~$\Gamma$ \emph{पासून
निष्पन्न करता येण्याजोगे} असते, म्हणजे $\Gamma \Proves \varphi$, तेव्हा आणि केवळ तेव्हाच,
जेव्हा एखादा सांत संच $\{\psi_1, \dots, \psi_n\} \subseteq \Gamma$ असा असतो
आणि पुढील संचासाठी बंद टॅब्लो असतो:
\[
\{
  \sFmla{\False}{\varphi},
  \sFmla{\True}{\psi_1}, \dots,
  \sFmla{\True}{\psi_n}
  \}.
\]
$\varphi$ हे $\Gamma$ पासून निष्पन्न करता येण्याजोगे नसेल, तर आपण $\Gamma \Proves/
\varphi$ लिहितो.
\end{defn}

\begin{defn}[सुसंगतता]
वाक्यांचा संच~$\Gamma$ हा \emph{विसंगत} असतो तेव्हा आणि केवळ तेव्हाच,
जेव्हा एखादा सांत संच $\{\psi_1, \dots, \psi_n\} \subseteq \Gamma$ असा असतो
आणि पुढील संचासाठी बंद टॅब्लो असतो:
\[
\{
  \sFmla{\True}{\psi_1}, \dots,
  \sFmla{\True}{\psi_n}
  \}.
\]
$\Gamma$ विसंगत नसेल, तर त्याला \emph{सुसंगत} म्हणतो.
\end{defn}

\begin{prop}[परावर्तिता]
\label{fol:tab:ptn:prop:reflexivity}
$\varphi \in \Gamma$ असेल, तर $\Gamma \Proves \varphi$.
\end{prop}

\begin{proof}
$\varphi \in \Gamma$ असेल, तर $\{\varphi\}$ हा $\Gamma$ चा सांत उपसंच आहे आणि पुढील टॅब्लो
\begin{oltableau}
  [\sFmla{\False}{\varphi}, just = \TAss
    [\sFmla{\True}{\varphi}, just = \TAss,close]
  ]
\end{oltableau}
बंद आहे.
\end{proof}


\begin{prop}[एकस्वनिकता]
\label{fol:tab:ptn:prop:monotonicity}
$\Gamma \subseteq \Delta$ आणि $\Gamma \Proves \varphi$ असेल, तर $\Delta
\Proves \varphi$.
\end{prop}

\begin{proof}
$\Gamma$ चा कोणताही सांत उपसंच हा $\Delta$ चासुद्धा सांत उपसंच असतो.
\end{proof}

\begin{prop}[संक्रमकता]
\label{fol:tab:ptn:prop:transitivity}
$\Gamma \Proves \varphi$ आणि $\{\varphi\} \cup
\Delta \Proves \psi$ असेल, तर $\Gamma \cup \Delta \Proves \psi$.
\end{prop}

\begin{proof}
$\{\varphi\} \cup \Delta \Proves \psi$ असेल, तर $\Delta_0 =
\{\chi_1, \dots, \chi_n\} \subseteq \Delta$ असा सांत उपसंच असतो आणि पुढील संचासाठी
\begin{align*}
\{\sFmla{\False}{\psi}, & \sFmla{\True}{\varphi}, \sFmla{\True}{\chi_1},
\dots, \sFmla{\True}{\chi_n}\}
\intertext{बंद टॅब्लो असतो. जर $\Gamma \Proves \varphi$ असेल, तर
  $\{\theta_1, \dots, \theta_m\} \subseteq \Gamma$ असा संच असतो आणि पुढील संचासाठीही}
\{\sFmla{\False}{\varphi}, & \sFmla{\True}{\theta_1},
\dots, \sFmla{\True}{\theta_m}\}
\end{align*}
बंद टॅब्लो असतो.

आता पुढील गृहीतके असलेला टॅब्लो विचारात घ्या:
\[
\sFmla{\False}{\psi},
\sFmla{\True}{\chi_1}, \dots, \sFmla{\True}{\chi_n},
\sFmla{\True}{\theta_1}, \dots, \sFmla{\True}{\theta_m}.
\]
त्यात~$\varphi$ वर \Cut{} नियम लावा. यामुळे दोन शाखा तयार होतात: एका शाखेत
$\sFmla{\True}{\varphi}$, तर दुसऱ्या शाखेत $\sFmla{\False}{\varphi}$ असते. त्यामुळे
पहिल्या शाखेत पुढील सर्व गृहीतके उपलब्ध असतात:
\[
\{\sFmla{\False}{\psi}, \sFmla{\True}{\varphi},
\sFmla{\True}{\chi_1}, \dots, \sFmla{\True}{\chi_n}\}
\]
या गृहीतकांसाठी बंद टॅब्लो असल्यामुळे तो त्या शाखेला जोडता येतो;
$\sFmla{\True}{\varphi}$ मधून जाणारी प्रत्येक शाखा बंद होते. दुसऱ्या शाखेत
पुढील सर्व गृहीतके उपलब्ध असतात:
\[
\{\sFmla{\False}{\varphi}, \sFmla{\True}{\theta_1}, \dots,
\sFmla{\True}{\theta_m}\}
\]
म्हणून दुसरी शाखासुद्धा पूर्ण करून बंद टॅब्लो मिळवता येतो. यावरून
$\Gamma \cup \Delta \Proves \psi$ हे दिसते.
\end{proof}

विशेषतः, याचा अर्थ असा की $\Gamma \Proves \varphi$ आणि $\varphi
\Proves \psi$ असेल, तर $\Gamma \Proves \psi$. तसेच $\varphi_1,
\dots, \varphi_n \Proves \psi$ असेल आणि प्रत्येक~$i$ साठी $\Gamma \Proves \varphi_i$
असेल, तर $\Gamma \Proves \psi$, हेसुद्धा यावरून मिळते.

\begin{prop}
\label{fol:tab:ptn:prop:incons}
$\Gamma$ विसंगत असतो तेव्हा आणि केवळ तेव्हाच, जेव्हा प्रत्येक
  वाक्य~$\varphi$ साठी $\Gamma \Proves \varphi$.
\end{prop}

\begin{proof}
सराव.
\end{proof}


\begin{prob}
\ref{fol:tab:ptn:prop:incons} सिद्ध करा.
\end{prob}




\begin{prop}[संहतता]
  \label{fol:tab:ptn:prop:proves-compact}
  \begin{enumerate}
  \item $\Gamma \Proves \varphi$ असेल, तर एखादा सांत उपसंच $\Gamma_0
    \subseteq \Gamma$ असा आहे की $\Gamma_0 \Proves \varphi$.
  \item $\Gamma$ चा प्रत्येक सांत उपसंच सुसंगत असेल, तर $\Gamma$ सुसंगत
    असतो.
  \end{enumerate}
\end{prop}

\begin{proof}
  \begin{enumerate}
    \item $\Gamma \Proves \varphi$ असेल, तर $\Gamma_0 =
      \{\psi_1, \dots, \psi_n\}$ असा त्याचा सांत उपसंच आणि पुढील संचासाठी बंद टॅब्लो असतो:
      \[
      \{\sFmla{\False}{\varphi}, \sFmla{\True}{\psi_1}, \dots, \sFmla{\True}{\psi_n}\}
      \]
      हा टॅब्लो $\Gamma_0 \Proves \varphi$ हेसुद्धा दाखवतो.
    \item $\Gamma$ विसंगत असेल, तर
      $\Gamma_0 = \{\psi_1, \dots, \psi_n\}$ असा त्याचा एखादा सांत उपसंच असून पुढील
      संचाचा बंद टॅब्लो असतो:
      \[
      \{\sFmla{\True}{\psi_1}, \dots, \sFmla{\True}{\psi_n}\}
      \]
      हा बंद टॅब्लो $\Gamma_0$ विसंगत असल्याचे दाखवतो.
  \end{enumerate}
\end{proof}












\section{निष्पन्नता आणि सुसंगतता}\label{fol:tab:prv:sec}

आता आपण निष्पन्नता संबंधाचे अनेक गुणधर्म सिद्ध करू. ते स्वतंत्रपणेही
महत्त्वाचे आहेत, पण संपूर्णता प्रमेयाच्या सिद्धतेत प्रत्येकाची भूमिका असेल.

\begin{prop}\label{fol:tab:prv:prop:provability-contr}
  जर $\Gamma \Proves \varphi$ आणि $\Gamma \cup \{\varphi\}$ विसंगत असेल, तर
  $\Gamma$ विसंगत आहे.
\end{prop}

\begin{proof}
$\Gamma_0 = \{\psi_1, \dots, \psi_n\}$ आणि $\Gamma_1
  =\{\chi_1, \dots, \chi_m\} \subseteq \Gamma$ असे सांत उपसंच असतात की
  \begin{align*}
    \{\sFmla{\False}{\varphi}, &
    \sFmla{\True}{\psi_1}, \dots, \sFmla{\True}{\psi_n}\} \\
    \{\sFmla{\True}{\varphi}, &
    \sFmla{\True}{\chi_1}, \dots, \sFmla{\True}{\chi_m}\}
  \end{align*}
  यांसाठी बंद टॅब्लो असतात. $\varphi$ वर \Cut{} नियम वापरून हे दोन्ही
  एकाच बंद टॅब्लोमध्ये एकत्र करता येतात; त्यावरून $\Gamma_0
  \cup \Gamma_1$ विसंगत असल्याचे दिसते. $\Gamma_0
  \subseteq \Gamma$ आणि $\Gamma_1 \subseteq \Gamma$ असल्यामुळे $\Gamma_0 \cup
  \Gamma_1 \subseteq \Gamma$; म्हणून $\Gamma$ विसंगत आहे.
\end{proof}

\begin{prop}
\label{fol:tab:prv:prop:prov-incons}
$\Gamma \Proves \varphi$ तेव्हा आणि केवळ तेव्हाच, जेव्हा $\Gamma \cup \{\lnot \varphi\}$
विसंगत असते.
\end{prop}

\begin{proof}
प्रथम $\Gamma \Proves \varphi$ असे समजा; म्हणजे पुढील संचासाठी बंद
टॅब्लो असतो:
\[
\{\sFmla{\False}{\varphi},
\sFmla{\True}{\psi_1}, \dots, \sFmla{\True}{\psi_n}\}
\]
$\TRule{\True}{\lnot}$ नियम वापरून त्याचे पुढील संचासाठीच्या बंद
टॅब्लोमध्ये रूपांतर करता येते:
\[
\{\sFmla{\True}{\lnot \varphi},
\sFmla{\True}{\psi_1}, \dots, \sFmla{\True}{\psi_n}\}.
\]

दुसरीकडे, दुसऱ्या संचासाठी बंद टॅब्लो असेल, तर त्यातील पहिले
गृहीतक~$\sFmla{\True}{\lnot \varphi}$ आणि त्या गृहीतकाला
\TRule{\True}{\lnot} लावून मिळणारी सर्व सूत्रे काढून, तसेच
$\sFmla{\False}{\varphi}$ हे गृहीतक जोडून, त्याचे पहिल्या संचासाठीच्या बंद
टॅब्लोमध्ये रूपांतर करता येते. कारण एखादी शाखा पूर्वी
$\sFmla{\True}{\lnot \varphi}$ ला \TRule{\True}{\lnot} लावून मिळणारा
निष्कर्ष, म्हणजे $\sFmla{\False}{\varphi}$, असल्यामुळे बंद झाली असेल, तर नव्या
टॅब्लो मधील संबंधित शाखाही बंद असते. जुन्या टॅब्लोतील एखादी शाखा
$\sFmla{\True}{\lnot \varphi}$ हे गृहीतक आणि $\sFmla{\False}{\lnot \varphi}$
ही दोन्ही सूत्रे असल्यामुळे बंद झाली असेल, तर $\sFmla{\False}{\lnot \varphi}$
ला $\TRule{\False}{\lnot}$ लावून $\sFmla{\True}{\varphi}$ मिळवता येते आणि
ती शाखा बंद करता येते. आपण $\sFmla{\False}{\varphi}$ हे गृहीतक जोडलेले
असल्यामुळे ही शाखा बंद होते.
\end{proof}

\begin{prob}
$\Gamma \Proves \lnot \varphi$ तेव्हा आणि केवळ तेव्हाच, जेव्हा $\Gamma \cup \{\varphi\}$
विसंगत असते, हे सिद्ध करा.
\end{prob}

\begin{prop}\label{fol:tab:prv:prop:explicit-inc}
  जर $\Gamma \Proves \varphi$ आणि $\lnot \varphi \in \Gamma$ असेल, तर $\Gamma$
  विसंगत आहे.
\end{prop}

\begin{proof}
  समजा $\Gamma \Proves \varphi$ आणि $\lnot \varphi \in \Gamma$. मग
  $\psi_1$, \dots, $\psi_n \in \Gamma$ अशी वाक्ये असतात की
  \[
  \{\sFmla{\False}{\varphi}, \sFmla{\True}{\psi_1}, \dots, \sFmla{\True}{\psi_n}\}
  \]
  यासाठी बंद टॅब्लो असतो. \sFmla{\False}{\varphi} हे गृहीतक
  \sFmla{\True}{\lnot \varphi} ने बदला आणि \sFmla{\True}{\lnot \varphi} ला
  \TRule{\True}{\lnot} लावून मिळणारा निष्कर्ष गृहीतकांनंतर घाला. जुन्या
  टॅब्लोमध्ये ओळ~$1$ च्या आधारे समर्थित कोणतेही वाक्य आता
  ओळ~$n+2$ च्या आधारे समर्थित होते. म्हणून जुना टॅब्लो बंद असेल, तर
  नवा टॅब्लोसुद्धा बंद असतो. सर्व गृहीतके~$\Gamma$ मध्ये असल्यामुळे हा टॅब्लो
  $\Gamma$ विसंगत असल्याचे दाखवतो.
\end{proof}

\begin{prop}\label{fol:tab:prv:prop:provability-exhaustive}
  $\Gamma \cup \{\varphi\}$ आणि $\Gamma \cup \{\lnot \varphi\}$ हे दोन्ही संच
  विसंगत असतील, तर $\Gamma$ विसंगत आहे.
\end{prop}

\begin{proof}
  $\psi_1$, \dots, $\psi_n \in \Gamma$ आणि $\chi_1$, \dots,
  $\chi_m \in \Gamma$ अशी वाक्ये असतील आणि
  \begin{align*}
    \{\sFmla{\True}{\varphi}, &
    \sFmla{\True}{\psi_1}, \dots, \sFmla{\True}{\psi_n}\} \text{ आणि}\\
    \{\sFmla{\True}{\lnot \varphi}, &
    \sFmla{\True}{\chi_1}, \dots, \sFmla{\True}{\chi_m}\}
  \end{align*}
  या दोन्हींसाठी बंद टॅब्लो असतील, तर
  $\sFmla{\True}{\psi_1}$, \dots, $\sFmla{\True}{\psi_n}$ आणि
  $\sFmla{\True}{\chi_1}$, \dots, $\sFmla{\True}{\chi_m}$ ही गृहीतके घेऊन,
  त्यानंतर \Cut{} नियम लावून, $\Gamma$ विसंगत असल्याचे दाखवणारा एकच
  संयुक्त टॅब्लो तयार करता येतो. त्यामुळे दोन शाखा मिळतात: एक
  $\sFmla{\True}{\varphi}$ ने, तर दुसरी $\sFmla{\False}{\varphi}$ ने सुरू होते.

  डाव्या बाजूस पहिल्या टॅब्लो मधील त्याच्या गृहीतकांखालचा भाग जोडा.
  येथे प्रत्येक नियमाचा उपयोग अजूनही योग्य आहे, कारण
  $\sFmla{\True}{\varphi}$ सहित पहिल्या टॅब्लो मधील प्रत्येक गृहीतक
  उपलब्ध आहे. त्यामुळे $\sFmla{\True}{\varphi}$ च्या खालील प्रत्येक शाखा बंद होते.

  उजव्या बाजूस दुसऱ्या टॅब्लो मधील त्याच्या गृहीतकांखालचा भाग जोडा;
  मात्र $\sFmla{\True}{\lnot \varphi}$ ला~$\TRule{\True}{\lnot}$ लावून मिळणारे
  सर्व परिणाम काढून टाका. $\sFmla{\True}{\lnot \varphi}$ ला
  $\TRule{\True}{\lnot}$ लावल्याचा निष्कर्ष $\sFmla{\False}{\varphi}$ हा आहे;
  तरीही तो उपलब्ध असतो, कारण तो संयुक्त टॅब्लोच्या उजव्या बाजूवरील
  \Cut{} नियमाचा निष्कर्ष आहे.

  दुसऱ्या टॅब्लोतील एखादी शाखा $\sFmla{\True}{\lnot \varphi}$ हे गृहीतक
  (जे आता संयुक्त टॅब्लोमध्ये गृहीतक म्हणून येत नाही) आणि
  $\sFmla{\False}{\lnot \varphi}$ ही दोन्ही सूत्रे असल्यामुळे बंद झाली असेल,
  तर $\sFmla{\False}{\lnot \varphi}$ ला $\TRule{\False}{\lnot}$ लावून
  $\sFmla{\True}{\varphi}$ मिळवता येते. मग संयुक्त टॅब्लो मधील संबंधित
  शाखासुद्धा बंद होते, कारण तिच्यात \Cut{} नियमाचा उजवीकडील निष्कर्ष
  $\sFmla{\False}{\varphi}$ असतो. दुसऱ्या टॅब्लो मधील एखादी शाखा अन्य
  कोणत्याही कारणाने बंद झाली असेल, तर संयुक्त टॅब्लो मधील संबंधित
  शाखासुद्धा बंद होते; कारण जुन्या दुसऱ्या टॅब्लो मधील त्या शाखेवर
  येणारी $\sFmla{\True}{\lnot \varphi}$ व्यतिरिक्त इतर सर्व चिन्हांकित सूत्रे
  संयुक्त टॅब्लो मधील संबंधित शाखेवरही येतात.
  \end{proof}















\section{निष्पन्नता आणि विधानीय संयोजक}\label{fol:tab:ppr:sec}

\begin{explain}
  टॅब्लोचा निष्पन्नता संबंध~$\Proves$ हा विधानीय संयोजकांशी संबंधित
  काही मूलभूत तथ्ये सिद्ध करण्यासाठी पुरेसा सामर्थ्यवान आहे; उदाहरणार्थ,
  $\varphi \land \psi \Proves \varphi$ आणि $\varphi, \varphi \lif \psi \Proves \psi$
  (विध्यात्मक अनुमान). ही तथ्ये संपूर्णता प्रमेयाच्या सिद्धतेसाठी आवश्यक आहेत.
\end{explain}

\begin{prop}\label{fol:tab:ppr:prop:provability-land}
  \begin{enumerate}
  \item \label{fol:tab:ppr:prop:provability-land-left} $\varphi \land \psi \Proves
    \varphi$ आणि $\varphi \land \psi \Proves \psi$ ही दोन्ही तथ्ये सत्य आहेत.
  \item \label{fol:tab:ppr:prop:provability-land-right} $\varphi, \psi \Proves \varphi \land
    \psi$.
  \end{enumerate}
\end{prop}

\begin{proof}
  \begin{enumerate}
  \item $\{\sFmla{\False}{\varphi}, \sFmla{\True}{\varphi \land \psi}\}$ आणि
    $\{\sFmla{\False}{\psi}, \sFmla{\True}{\varphi \land \psi}\}$ या दोन्हींसाठी
    बंद टॅब्लो आहेत:
    \begin{oltableau}{}
      [\sFmla{\False}{\varphi}, just=\TAss
        [\sFmla{\True}{\varphi \land \psi}, just=\TAss
          [\sFmla{\True}{\varphi},just={\TRule{\True}{\land}[2]}
            [\sFmla{\True}{\psi},just={\TRule{\True}{\land}[2]}, close
            ]
          ]
        ]
      ]
    \end{oltableau}
    \begin{oltableau}{}
      [\sFmla{\False}{\psi}, just=\TAss
        [\sFmla{\True}{\varphi \land \psi}, just=\TAss
          [\sFmla{\True}{\varphi},just={\TRule{\True}{\land}[2]}
            [\sFmla{\True}{\psi},just={\TRule{\True}{\land}[2]}, close
            ]
          ]
        ]
      ]
    \end{oltableau}
    \item $\{\sFmla{\True}{\varphi}, \sFmla{\True}{\psi},
      \sFmla{\False}{\varphi \land \psi}\}$ साठीचा बंद टॅब्लो पुढीलप्रमाणे आहे:
      \begin{oltableau}
        [\sFmla{\False}{\varphi \land \psi}, just = \TAss
          [\sFmla{\True}{\varphi}, just = \TAss
            [\sFmla{\True}{\psi}, just=\TAss
              [\sFmla{\False}{\varphi}, just = {\TRule{\False}{\land}[1]}, close]
              [\sFmla{\False}{\psi}, just = {\TRule{\False}{\land}[1]}, close]
            ]
          ]
        ]
      \end{oltableau}
  \end{enumerate}
\end{proof}

\begin{prop}\label{fol:tab:ppr:prop:provability-lor}
  \begin{enumerate}
  \item $\{\varphi \lor \psi, \lnot \varphi, \lnot \psi\}$ विसंगत आहे.
  \item $\varphi \Proves \varphi \lor \psi$ आणि $\psi \Proves \varphi \lor \psi$ ही दोन्ही
    तथ्ये सत्य आहेत.
  \end{enumerate}
\end{prop}

\begin{proof}
  \begin{enumerate}
  \item $\{\sFmla{\True}{\varphi \lor \psi}, \sFmla{\True}{\lnot \varphi},
    \sFmla{\True}{\lnot \psi}\}$ साठी बंद टॅब्लो देऊ:
    \begin{oltableau}
      [\sFmla{\True}{\varphi \lor \psi}, just = \TAss
        [\sFmla{\True}{\lnot \varphi}, just = \TAss
          [\sFmla{\True}{\lnot \psi}, just = \TAss
            [\sFmla{\False}{\varphi}, just = {\TRule{\True}{\lnot}[2]}
              [\sFmla{\False}{\psi}, just = {\TRule{\True}{\lnot}[3]}
                [\sFmla{\True}{\varphi}, just = {\TRule{\True}{\lor}[1]}, close]
                [\sFmla{\True}{\psi}, just = {\TRule{\True}{\lor}[1]}, close]
              ]
            ]
          ]
        ]
      ]
    \end{oltableau}
  \item $\{\sFmla{\False}{\varphi \lor \psi}, \sFmla{\True}{\varphi}\}$ आणि
    $\{\sFmla{\False}{\varphi \lor \psi}, \sFmla{\True}{\psi}\}$ या दोन्हींसाठी
    बंद टॅब्लो आहेत:
    \begin{oltableau}{}
      [\sFmla{\False}{\varphi \lor \psi}, just=\TAss
        [\sFmla{\True}{\varphi}, just=\TAss
          [\sFmla{\False}{\varphi},just={\TRule{\False}{\lor}[1]}
            [\sFmla{\False}{\psi},just={\TRule{\False}{\lor}[1]}, close
            ]
          ]
        ]
      ]
    \end{oltableau}
    \begin{oltableau}{}
      [\sFmla{\False}{\varphi \lor \psi}, just=\TAss
        [\sFmla{\True}{\psi}, just=\TAss
          [\sFmla{\False}{\varphi},just={\TRule{\False}{\lor}[1]}
            [\sFmla{\False}{\psi},just={\TRule{\False}{\lor}[1]}, close
            ]
          ]
        ]
      ]
    \end{oltableau}
  \end{enumerate}
\end{proof}

\begin{prop}\label{fol:tab:ppr:prop:provability-lif}
  \begin{enumerate}
  \item \label{fol:tab:ppr:prop:provability-lif-left}  $\varphi, \varphi \lif \psi \Proves \psi$.
  \item \label{fol:tab:ppr:prop:provability-lif-right}
    $\lnot \varphi \Proves \varphi \lif \psi$ आणि $\psi \Proves \varphi \lif \psi$ ही दोन्ही
    तथ्ये सत्य आहेत.
  \end{enumerate}
\end{prop}

\begin{proof}
  \begin{enumerate}
    \item $\{\sFmla{\False}{\psi}, \sFmla{\True}{\varphi \lif \psi},
      \sFmla{\True}{\varphi}\}$ साठी बंद टॅब्लो आहे:
      \begin{oltableau}
        [\sFmla{\False}{\psi}, just=\TAss
          [\sFmla{\True}{\varphi \lif \psi}, just=\TAss
            [\sFmla{\True}{\varphi}, just=\TAss
              [\sFmla{\False}{\varphi}, just = {\TRule{\True}{\lif}[2]}, close]
              [\sFmla{\True}{\psi}, just = {\TRule{\True}{\lif}[2]}, close]
            ]
          ]
        ]
      \end{oltableau}
    \item $\{\sFmla{\False}{\varphi \lif \psi}, \sFmla{\True}{\lnot \varphi}\}$ आणि
      $\{\sFmla{\False}{\varphi \lif \psi}, \sFmla{\True}{\psi}\}$ या दोन्हींसाठी
      बंद टॅब्लो आहेत:
      \begin{oltableau}
        [\sFmla{\False}{\varphi \lif \psi}, just = \TAss
          [\sFmla{\True}{\lnot \varphi}, just = \TAss
            [\sFmla{\True}{\varphi}, just = {\TRule{\False}{\lif}[1]}
              [\sFmla{\False}{\psi}, just = {\TRule{\False}{\lif}[1]}
                [\sFmla{\False}{\varphi}, just = {\TRule{\True}{\lnot}[2]}, close]
              ]
            ]
          ]
        ]
      \end{oltableau}
      \begin{oltableau}
        [\sFmla{\False}{\varphi \lif \psi}, just = \TAss
          [\sFmla{\True}{\psi}, just = \TAss
            [\sFmla{\True}{\varphi}, just = {\TRule{\False}{\lif}[1]}
              [\sFmla{\False}{\psi}, just = {\TRule{\False}{\lif}[1]},close
              ]
            ]
          ]
        ]
      \end{oltableau}
  \end{enumerate}
\end{proof}











\section{निष्पन्नता आणि संख्यापक}\label{fol:tab:qpr:sec}

\begin{explain}
  संपूर्णता प्रमेयासाठी या विभागात स्थापित केलेली~$\Proves$ संबंधी तथ्ये
  टॅब्लोच्या नियमांनी निष्पन्न होतात, हेही आवश्यक आहे.
\end{explain}

\begin{thm}
\label{fol:tab:qpr:thm:strong-generalization} जर $c$ हे $\Gamma$ किंवा~$\varphi(x)$ मध्ये
न येणारे अचर असेल आणि $\Gamma \Proves \varphi(c)$, तर $\Gamma
\Proves \lforall[x][\varphi(x)]$.
\end{thm}

\begin{proof}
$\Gamma \Proves \varphi(c)$ असे समजा; म्हणजे $\psi_1$, \dots, $\psi_n
\in \Gamma$ अशी वाक्ये आणि पुढील संचासाठी बंद टॅब्लो असतो:
\begin{align*}
\{ \sFmla{\False}{\varphi(c)}, &
\sFmla{\True}{\psi_1}, \dots, \sFmla{\True}{\psi_n} \}.
\intertext{पुढील संचासाठीही बंद टॅब्लो असतो, हे दाखवायचे आहे:}
\{ \sFmla{\False}{\lforall[x][\varphi(x)]}, &
\sFmla{\True}{\psi_1}, \dots, \sFmla{\True}{\psi_n} \}.
\end{align*}
बंद टॅब्लो घेऊन त्यातील पहिले गृहीतक
$\sFmla{\False}{\lforall[x][\varphi(x)]}$ ने बदला आणि गृहीतकांनंतर
$\sFmla{\False}{\varphi(c)}$ घाला.
\begin{center}
\begin{tableau}{not line numbering}
  [\sFmla{\False}{\varphi(c)}
    [\sFmla{\True}{\psi_1}
      [\vdots
      [\sFmla{\True}{\psi_n},
        [][][]]]]]
\end{tableau}\qquad
\begin{tableau}{not line numbering}
  [\sFmla{\False}{\lforall[x][\varphi(x)]}
    [\sFmla{\True}{\psi_1}
      [\vdots
        [\sFmla{\True}{\psi_n},
          [\sFmla{\False}{\varphi(c)}
        [][][]]]]]]
\end{tableau}
\end{center}
हा टॅब्लो अजूनही बंद आहे, कारण पूर्वी गृहीतके म्हणून उपलब्ध असलेली सर्व
वाक्ये नव्या टॅब्लोच्या वरच्या भागात अजूनही उपलब्ध आहेत. घातलेली
ओळ ही $\TRule{\False}{\lforall}$ च्या योग्य उपयोगाचा परिणाम आहे, कारण
स्थिरांक~$c$ हे $\psi_1$, \dots, $\psi_n$ किंवा
$\lforall[x][\varphi(x)]$ मध्ये येत नाही; म्हणजे ते नव्या टॅब्लोमध्ये घातलेल्या
ओळीच्या वर येत नाही.
\end{proof}

\begin{prop}
\label{fol:tab:qpr:prop:provability-quantifiers}
\begin{enumerate}
\item $\varphi(t) \Proves \lexists[x][\varphi(x)]$.
\item $\lforall[x][\varphi(x)] \Proves \varphi(t)$.
\end{enumerate}
\end{prop}

\begin{proof}
\begin{enumerate}
\item
  $\sFmla{\False}{\lexists[x][\varphi(x)]}, \sFmla{\True}{\varphi(t)}$ साठीचा बंद
  टॅब्लो पुढीलप्रमाणे आहे:
  \begin{oltableau}
    [\sFmla{\False}{\lexists[x][\varphi(x)]}, just = \TAss
      [\sFmla{\True}{\varphi(t)}, just = \TAss
        [\sFmla{\False}{\varphi(t)}, just = {\TRule{\False}{\lexists}[1]}, close]]]
    \end{oltableau}
\item
  $\sFmla{\False}{\varphi(t)}, \sFmla{\True}{\lforall[x][\varphi(x)]}, $
  साठीचा बंद टॅब्लो पुढीलप्रमाणे आहे:
  \begin{oltableau}
    [\sFmla{\False}{\varphi(t)}, just = \TAss
      [\sFmla{\True}{\lforall[x][\varphi(x)]}, just = \TAss
        [\sFmla{\True}{\varphi(t)}, just = {\TRule{\True}{\lforall}[2]}, close]]]
    \end{oltableau}
\end{enumerate}
\end{proof}












\section{निर्दोषता}\label{fol:tab:sou:sec}

\begin{explain}
टॅब्लोसारखी निष्पत्ती-पद्धत \emph{निर्दोष} असते, जर ती प्रत्यक्षात
धारण न करणाऱ्या गोष्टी निष्पन्न करू शकत नसेल. त्यामुळे निर्दोषता हा
निष्पत्ती-पद्धतींसाठी हमी दिलेल्या सुरक्षिततेचा एक प्रकार आहे. कोणता
सिद्धता-उपपत्तीय गुणधर्म विचारात आहे यावर अवलंबून, उदाहरणार्थ, आपल्याला
पुढील गोष्टी हव्या असतात:
\begin{enumerate}
\item प्रत्येक निष्पन्न करता येण्याजोगे~$\varphi$ हे वैध असते;
\item एखादे वाक्य इतर काही वाक्यांपासून निष्पन्न करता येण्याजोगे असेल, तर ते
  त्यांचा तार्किक परिणामही असते;
\item वाक्यांचा संच विसंगत असेल, तर तो अपूर्ततायोग्य असतो.
\end{enumerate}
हे निष्पत्ती-पद्धतीचे महत्त्वाचे गुणधर्म आहेत. यांपैकी एखादा गुणधर्म
खरा नसेल, तर निष्पत्ती-पद्धतीत दोष आहे---ती खूप जास्त गोष्टी
निष्पन्न करेल. म्हणून निष्पत्ती-पद्धतीची निर्दोषता सिद्ध करणे अत्यंत
महत्त्वाचे आहे.

हे सर्व सिद्धता-उपपत्तीय गुणधर्म कोणत्या ना कोणत्या प्रकारच्या बंद
टॅब्लो द्वारे परिभाषित केलेले असल्यामुळे, वरील (1)--(3) सिद्ध करण्यासाठी
बंद टॅब्लोच्या चिन्हार्थविषयक गुणधर्मांबद्दल काहीतरी सिद्ध करावे लागते.
प्रथम एखाद्या चिन्हांकित सूत्राची संरचनेत पूर्ती होणे म्हणजे काय ते
परिभाषित करू; मग टॅब्लो बंद असेल, तर कोणतीही संरचना त्याची सर्व गृहीतके
पूर्ण करत नाही, हे दाखवू. त्यानंतर (1)--(3) या परिणामाचे उपपरिणाम म्हणून
मिळतील.
\end{explain}

\begin{defn}
संरचना~$\Struct{M}$
चिन्हांकित सूत्र $\sFmla{\True}{\varphi}$ \emph{पूर्ण करते} तेव्हा आणि केवळ तेव्हाच, जेव्हा
$\Sat{M}{\varphi}$; आणि ते
$\sFmla{\False}{\varphi}$ पूर्ण करते तेव्हा आणि केवळ तेव्हाच, जेव्हा
$\Sat/{M}{\varphi}$. $\Struct{M}$
हा चिन्हांकित सूत्रांचा संच~$\Gamma$ पूर्ण करतो तेव्हा आणि केवळ तेव्हाच,
जेव्हा तो प्रत्येक $\sFmla{S}{\varphi} \in \Gamma$ पूर्ण करतो. $\Gamma$ पूर्ण
करणारी एखादी संरचना असेल, तर तो
\emph{पूर्ततायोग्य} आणि अन्यथा \emph{अपूर्ततायोग्य} असतो.
\end{defn}

\begin{thm}[निर्दोषता]
  \label{fol:tab:sou:thm:tableau-soundness}
  $\Gamma$ साठी बंद टॅब्लो असेल, तर $\Gamma$ अपूर्ततायोग्य असतो.
\end{thm}

\begin{proof}
टॅब्लोच्या शाखेवरील चिन्हांकित सूत्रांचा संच पूर्ततायोग्य असेल, तर
आणि केवळ तेव्हाच त्या शाखेला पूर्ततायोग्य म्हणू; तसेच टॅब्लोमध्ये
किमान एक पूर्ततायोग्य शाखा असेल, तर त्याला पूर्ततायोग्य म्हणू.

पुढील गोष्ट दाखवू: एखाद्या पूर्ततायोग्य टॅब्लोचा अनुमानाच्या कोणत्याही
एका नियमाने विस्तार केल्यास मिळणारा टॅब्लो नेहमी पूर्ततायोग्य असतो.
यावरून प्रमेय सिद्ध होईल: $\Gamma$ मधील केवळ गृहीतके असलेल्या टॅब्लो ला
अनुमानाचे नियम लावून कोणताही बंद टॅब्लो मिळतो. म्हणून $\Gamma$
पूर्ततायोग्य असेल, तर त्याच्यासाठीचा कोणताही टॅब्लो पूर्ततायोग्य असेल.
पण बंद टॅब्लो स्पष्टपणे पूर्ततायोग्य नसतो: प्रत्येक शाखेत
$\sFmla{\True}{\varphi}$ आणि $\sFmla{\False}{\varphi}$ ही दोन्ही सूत्रे असतात, आणि
कोणतीही संरचना $\varphi$ ची पूर्ती आणि अपूर्ती या दोन्ही करू शकत नाही.

आपल्याकडे पूर्ततायोग्य टॅब्लो, म्हणजे किमान एक पूर्ततायोग्य शाखा असलेला
टॅब्लो, आहे असे समजा. अनुमानाचा नियम लावल्यावर एकतर शाखेत
चिन्हांकित सूत्रे जोडली जातात, किंवा शाखा दोन भागांत विभागली जाते. त्या
नियमाच्या उपयोगाने विस्तारली न गेलेली एखादी पूर्ततायोग्य शाखा टॅब्लोमध्ये
असेल, तर विस्तारित टॅब्लोमध्येही ती पूर्ततायोग्य शाखा राहते; म्हणून
विस्तारित टॅब्लोसुद्धा पूर्ततायोग्य असतो. त्यामुळे पूर्ततायोग्य शाखेलाच नियम
लावला जातो ते प्रकरण पाहणे पुरेसे आहे.

त्या शाखेवरील चिन्हांकित सूत्रांचा संच~$\Gamma$ असू द्या, आणि ज्या
चिन्हांकित सूत्राला नियम लावला जातो ते $\sFmla{S}{\varphi} \in \Gamma$ असू
द्या. नियम लावल्यावर शाखा विभागली जात नसेल, तर विस्तारित शाखा, म्हणजे
$\Gamma$ आणि नियमाचे निष्कर्ष मिळून बनलेला संच, अजूनही पूर्ततायोग्य आहे हे
दाखवावे लागेल. नियम लावल्यावर शाखा विभागली जात असेल, तर मिळणाऱ्या दोन
शाखांपैकी किमान एक पूर्ततायोग्य आहे हे दाखवावे लागेल.

प्रथम, शाखा न विभागणारी संभाव्य अनुमाने पाहू.
\begin{enumerate}
\item $\sFmla{\True}{\lnot \psi} \in \Gamma$ ला
  $\TRule{\True}{\lnot}$ लावून शाखेचा विस्तार केला आहे असे समजा. मग
  विस्तारित शाखेवरील चिन्हांकित सूत्रांचा संच $\Gamma \cup
  \{\sFmla{\False}{\psi}\}$ असतो. समजा
  $\Sat{M}{\Gamma}$. विशेषतः,
  $\Sat{M}{\lnot \psi}$. त्यामुळे
  $\Sat/{M}{\psi}$, म्हणजे
  $\Struct{M}$ हे
  $\sFmla{\False}{\psi}$ पूर्ण करते.
\item $\sFmla{\False}{\lnot \psi} \in \Gamma$ ला
  $\TRule{\False}{\lnot}$ लावून शाखेचा विस्तार केला आहे: सराव.
\item $\sFmla{\True}{\psi \land \chi} \in \Gamma$ ला
  $\TRule{\True}{\land}$ लावून शाखेचा विस्तार केला आहे. यामुळे शाखेवर
  $\sFmla{\True}{\psi}$ आणि $\sFmla{\True}{\chi}$ ही दोन नवी
  चिन्हांकित सूत्रे मिळतात. समजा
  $\Sat{M}{\Gamma}$; विशेषतः,
  $\Sat{M}{\psi \land \chi}$. मग
  $\Sat{M}{\psi}$ आणि
  $\Sat{M}{\chi}$. म्हणजे
  $\Struct{M}$ हे
  $\sFmla{\True}{\psi}$ आणि $\sFmla{\True}{\chi}$ ही दोन्ही सूत्रे पूर्ण करते.
\item $\sFmla{\False}{\psi \lor \chi} \in \Gamma$ ला
  $\TRule{\False}{\lor}$ लावून शाखेचा विस्तार केला आहे: सराव.
\item $\sFmla{\False}{\psi \lif \chi} \in \Gamma$ ला
  $\TRule{\False}{\lif}$ लावून शाखेचा विस्तार केला आहे. यामुळे शाखेवर
  $\sFmla{\True}{\psi}$ आणि $\sFmla{\False}{\chi}$ ही दोन नवी
  चिन्हांकित सूत्रे मिळतात. समजा
  $\Sat{M}{\Gamma}$; विशेषतः,
  $\Sat/{M}{\psi \lif \chi}$. मग
  $\Sat{M}{\psi}$ आणि
  $\Sat/{M}{\chi}$. म्हणजे
  $\Struct{M}$ हे
  $\sFmla{\True}{\psi}$ आणि $\sFmla{\False}{\chi}$ ही दोन्ही सूत्रे पूर्ण करते.


\item $\sFmla{\True}{\lforall[x][\varphi(x)]} \in \Gamma$ ला
  $\TRule{\True}{\lforall}$ लावून शाखेचा विस्तार केला आहे. यामुळे शाखेवर
  नवीन चिन्हांकित सूत्र~$\sFmla{\True}{\varphi(t)}$ मिळते.
  समजा $\Sat{M}{\Gamma}$; विशेषतः,
  $\Sat{M}{\lforall[x][\varphi(x)]}$.
  \readerexternalref{fol:syn:sem:prop:quant-terms} नुसार $\Sat{M}{\varphi(t)}$. परिणामी,
  $\Struct{M}$ हे $\sFmla{\True}{\varphi(t)}$ पूर्ण करते.

\item $\sFmla{\False}{\lforall[x][\varphi(x)]} \in \Gamma$ ला
  $\TRule{\False}{\lforall}$ लावून शाखेचा विस्तार केला आहे. यामुळे शाखेवर
  नवीन चिन्हांकित सूत्र~$\sFmla{\False}{\varphi(a)}$ मिळते, जेथे $a$ हे
  $\Gamma$ मध्ये न येणारे स्थिरांक आहे. $\Gamma$ पूर्ततायोग्य असल्यामुळे
  $\Sat{M}{\Gamma}$ अशी एखादी $\Struct{M}$ असते; विशेषतः,
  $\Sat/{M}{\lforall[x][\varphi(x)]}$. आता
  $\Gamma \cup \{\sFmla{\False}{\varphi(a)}\}$ पूर्ततायोग्य आहे हे दाखवायचे आहे.
  त्यासाठी पुढीलप्रमाणे योग्य~$\Struct{M'}$ परिभाषित करू.

  \readerexternalref{fol:syn:ass:prop:sat-quant} नुसार,
  $\Sat/{M}{\lforall[x][\varphi(x)]}$ तेव्हा आणि केवळ तेव्हाच, जेव्हा काही $s$
  साठी $\Sat/{M}{\varphi(x)}[s]$. आता $\Struct{M'}$ ही $\Struct{M}$ सारखीच
  असू द्या, फक्त $\Assign{a}{M'} = s(x)$. $a$ हे $\Gamma$ मध्ये येत
  नसल्यामुळे, \readerexternalref{fol:syn:ext:cor:extensionality-sent} नुसार कोणत्याही
  $\sFmla{\True}{\chi} \in \Gamma$ साठी $\Sat{M'}{\chi}$, आणि कोणत्याही
  $\sFmla{\False}{\chi} \in \Gamma$ साठी $\Sat/{M'}{\chi}$.

  \readerexternalref{fol:syn:ext:prop:extensionality} नुसार $\Sat/{M'}{\varphi(x)}[s]$.
  \readerexternalref{fol:syn:ext:prop:ext-formulas} नुसार $\Sat/{M'}{\varphi(a)}[s]$.
  $\varphi(a)$ हे वाक्य असल्यामुळे,
  \readerexternalref{fol:syn:ass:prop:sentence-sat-true} नुसार $\Sat/{M'}{\varphi(a)}$;
  म्हणजे $\Struct{M'}$ हे $\sFmla{\False}{\varphi(a)}$ पूर्ण करते.
\item $\sFmla{\True}{\lexists[x][\psi(x)]} \in \Gamma$ ला
  $\TRule{\True}{\lexists}$ लावून शाखेचा विस्तार केला आहे: सराव.
\item $\sFmla{\False}{\lexists[x][\psi(x)]} \in \Gamma$ ला
  $\TRule{\False}{\lexists}$ लावून शाखेचा विस्तार केला आहे: सराव.
\end{enumerate}
आता शाखा विभागणारी संभाव्य अनुमाने पाहू.
\begin{enumerate}
\item $\sFmla{\False}{\psi \land \chi} \in \Gamma$ ला
  $\TRule{\False}{\land}$ लावून शाखेचा विस्तार केला आहे. यामुळे दोन शाखा
  मिळतात: डावीकडील शाखा $\sFmla{\False}{\psi}$ मधून आणि उजवीकडील शाखा
  $\sFmla{\False}{\chi}$ मधून पुढे जाते. समजा
  $\Sat{M}{\Gamma}$; विशेषतः,
  $\Sat/{M}{\psi \land \chi}$. मग एकतर
  $\Sat/{M}{\psi}$ किंवा
  $\Sat/{M}{\chi}$. पहिल्या प्रकरणात
  $\Struct{M}$ हे
  $\sFmla{\False}{\psi}$ पूर्ण करते; म्हणजे
  $\Struct{M}$ डावीकडील शाखेवरील सूत्रे पूर्ण करते.
  दुसऱ्या प्रकरणात $\Struct{M}$ हे
  $\sFmla{\False}{\chi}$ पूर्ण करते; म्हणजे
  $\Struct{M}$ उजवीकडील शाखेवरील सूत्रे पूर्ण करते.
\item $\sFmla{\True}{\psi \lor \chi} \in \Gamma$ ला
  $\TRule{\True}{\lor}$ लावून शाखेचा विस्तार केला आहे: सराव.
\item $\sFmla{\True}{\psi \lif \chi} \in \Gamma$ ला
  $\TRule{\True}{\lif}$ लावून शाखेचा विस्तार केला आहे: सराव.
\item \Cut{} लावून शाखेचा विस्तार केला आहे. यामुळे दोन शाखा मिळतात: एका
  शाखेत $\sFmla{\True}{\psi}$ आणि दुसऱ्या शाखेत $\sFmla{\False}{\psi}$ असते.
  $\Sat{M}{\Gamma}$ आणि एकतर
  $\Sat{M}{\psi}$ किंवा
  $\Sat/{M}{\psi}$ असल्यामुळे,
  $\Struct{M}$ हे डावीकडील किंवा उजवीकडील शाखा
  पूर्ण करते.
\end{enumerate}
\end{proof}


\begin{prob}
\ref{fol:tab:sou:thm:tableau-soundness} ची सिद्धता पूर्ण करा.
\end{prob}




\begin{cor}
\label{fol:tab:sou:cor:weak-soundness}
जर $\Proves \varphi$, तर $\varphi$ हे वैध आहे.
\end{cor}

\begin{cor}
\label{fol:tab:sou:cor:entailment-soundness}
जर $\Gamma \Proves \varphi$, तर $\Gamma \Entails \varphi$.
\end{cor}

\begin{proof}
  जर $\Gamma \Proves \varphi$, तर काही $\psi_1$, \dots, $\psi_n \in
  \Gamma$ साठी $\{\sFmla{\False}{\varphi}, \sFmla{\True}{\psi_1}, \dots,
  \sFmla{\True}{\psi_n}\}$ चा बंद टॅब्लो असतो.
  \ref{fol:tab:sou:thm:tableau-soundness} नुसार, प्रत्येक
  संरचना~$\Struct{M}$
  काही $\psi_i$ असत्य करते किंवा $\varphi$ सत्य करते. म्हणून
  $\Sat{M}{\Gamma}$ असेल, तर
  $\Sat{M}{\varphi}$ सुद्धा असेल.
\end{proof}

\begin{cor}
\label{fol:tab:sou:cor:consistency-soundness}
जर $\Gamma$ पूर्ततायोग्य असेल, तर तो सुसंगत आहे.
\end{cor}

\begin{proof}
प्रतिलोम विधान सिद्ध करू. $\Gamma$ सुसंगत नाही असे समजा. मग
$\psi_1$, \dots, $\psi_n \in \Gamma$ अशी वाक्ये आणि
$\{\sFmla{\True}{\psi_1}, \dots, \sFmla{\True}{\psi_n}\}$ चा बंद
टॅब्लो असतो. \ref{fol:tab:sou:thm:tableau-soundness} नुसार, प्रत्येक
$i=1$, \dots,~$n$ साठी $\Sat{M}{\psi_i}$ अशी
कोणतीही $\Struct{M}$ नसते. पण मग $\Gamma$
पूर्ततायोग्य नसतो.
\end{proof}












\section{टॅब्लो सह एकरूपता}\label{fol:tab:ide:sec}

टॅब्लोमध्ये एकरूपता असल्यास अतिरिक्त निगमन नियम आवश्यक असतात.
$\eq$ चे नियम पुढीलप्रमाणे आहेत ($t$, $t_1$ आणि $t_2$ ही बंद पदे आहेत):

\begin{defish}
\AxiomC{}
\RightLabel{$\eq$}
\UnaryInfC{\sFmla{\True}{\eq[t][t]}}
\DisplayProof
\hfill
\AxiomC{\sFmla{\True}{\eq[t_1][t_2]}}
\noLine
\UnaryInfC{\sFmla{\True}{\varphi(t_1)}}
\RightLabel{$\TRule{\True}{\eq}$}
\UnaryInfC{\sFmla{\True}{\varphi(t_2)}}
\DisplayProof
\hfill
\AxiomC{\sFmla{\True}{\eq[t_1][t_2]}}
\noLine
\UnaryInfC{\sFmla{\False}{\varphi(t_1)}}
\RightLabel{$\TRule{\False}{\eq}$}
\UnaryInfC{\sFmla{\False}{\varphi(t_2)}}
\DisplayProof
\end{defish}
लक्षात घ्या की इतर सर्व नियमांपेक्षा वेगळे, $\TRule{\True}{\eq}$ आणि
$\TRule{\False}{\eq}$ यांना शाखेवर आधीपासून \emph{दोन} चिन्हांकित
सूत्रे असणे आवश्यक आहे, म्हणजे $\sFmla{\True}{\eq[t_1][t_2]}$
आणि $\sFmla{S}{\varphi(t_1)}$ ही दोन्ही.

\begin{ex}
जर $s$ आणि $t$ ही बंद पदे असतील, तर $\eq[s][t], \varphi(s)
\Proves \varphi(t)$:
\begin{oltableau}
  [\sFmla{\False}{\varphi(t)}, just = \TAss
    [\sFmla{\True}{\eq[s][t]}, just = \TAss
      [\sFmla{\True}{\varphi(s)}, just = \TAss
        [\sFmla{\True}{\varphi(t)}, just={\TRule{\True}{\eq}[2, 3]}, close]
      ]
    ]
  ]
\end{oltableau}
याला एकरूप वस्तूंच्या प्रतिस्थापनाचे तत्त्व किंवा लाइब्निझचा नियम म्हणून
ओळखले जाऊ शकते.

टॅब्लो द्वारे $\eq$ सममित आहे, म्हणजे $\eq[s_1][s_2]
\Proves \eq[s_2][s_1]$, हे सिद्ध होते:
\begin{oltableau}
  [\sFmla{\False}{\eq[s_2][s_1]}, just = \TAss
    [\sFmla{\True}{\eq[s_1][s_2]}, just = \TAss
      [\sFmla{\True}{\eq[s_1][s_1]}, just = {$\eq$}
        [\sFmla{\True}{\eq[s_2][s_1]}, just = {\TRule{\True}{\eq}[2, 3]}, close]
      ]
    ]
  ]
\end{oltableau}
येथे ओळ~$2$ ही $\TRule{\True}{\eq}$ साठीचे पहिले आधारभूत
सूत्र, $\sFmla{\True}{\eq[s_1][s_2]}$, आहे. ओळ~$3$ हे
$\sFmla{\True}{\varphi(s_1)}$ या रूपातील दुसरे आधारविधान आहे---$\varphi(x)$ हे
$\eq[x][s_1]$ आहे असे विचार करा; मग $\varphi(s_1)$ म्हणजे
$\eq[s_1][s_1]$ आणि $\varphi(s_2)$ म्हणजे $\eq[s_2][s_1]$.

तसेच $\eq$ संक्रमक आहे, म्हणजे $\eq[s_1][s_2],
\eq[s_2][s_3] \Proves \eq[s_1][s_3]$, हेही ते सिद्ध करतात:
\begin{oltableau}
  [\sFmla{\False}{\eq[s_1][s_3]}, just = \TAss
    [\sFmla{\True}{\eq[s_1][s_2]}, just = \TAss
      [\sFmla{\True}{\eq[s_2][s_3]}, just = \TAss
        [\sFmla{\True}{\eq[s_1][s_3]}, just = {\TRule{\True}{\eq}[3, 2]}, close]
      ]
    ]
  ]
\end{oltableau}
या टॅब्लोमध्ये $\TRule{\True}{\eq}$ चे पहिले आधारभूत सूत्र
ओळ~$3$ वरील $\sFmla{\True}{\eq[s_2][s_3]}$ आहे ($s_2$ ही~$t_1$ ची,
आणि $s_3$ ही~$t_2$ ची भूमिका बजावतात). $\sFmla{\True}{\varphi(s_2)}$
या रूपातील दुसरे आधारविधान ओळ~$2$ आहे. येथे $\varphi(x)$ हे
$\eq[s_1][x]$ आहे असे विचार करा; त्यामुळे $\varphi(s_2)$ हे
$\eq[s_1][s_2]$ (म्हणजे ओळ~$2$) आणि निष्कर्षातील
सूत्र~$\eq[s_1][s_3]$ हे $\varphi(s_3)$ होते.
\end{ex}

\begin{prob}
पुढील गोष्टींसाठी बंद टॅब्लो द्या:
\begin{enumerate}
\item $\sFmla{\False}{\lforall[x][\lforall[y][((x = y \land \varphi(x))
      \lif \varphi(y))]]}$
\item $\sFmla{\False}{\lexists[x][(\varphi(x) \land
    \lforall[y][(\varphi(y) \lif y = x)])]}$,\\
  $\sFmla{\True}{\lexists[x][\varphi(x)] \land
    \lforall[y][\lforall[z][((\varphi(y) \land \varphi(z)) \lif y = z)]]}$
\end{enumerate}
\end{prob}












\section{एकरूपता सह निर्दोषता}\label{fol:tab:sid:sec}

\begin{prop}
एकरूपतेचे नियम असलेले टॅब्लो निर्दोष आहेत: कोणताही बंद टॅब्लो
  पूर्ततायोग्य नसतो.
\end{prop}

\begin{proof}
मागीलप्रमाणे एवढेच दाखवायचे आहे की टॅब्लोची एखादी शाखा पूर्ततायोग्य
असेल, तर तिला~$\eq$ चा कोणताही नियम लावून मिळणारी शाखासुद्धा पूर्ततायोग्य
असते. शाखेवरील चिन्हांकित सूत्रांचा संच~$\Gamma$ असू द्या आणि
$\Struct{M}$ ही $\Gamma$ पूर्ण करणारी संरचना असू द्या.

शाखेचा~$\eq$ वापरून विस्तार केला आहे, म्हणजे चिन्हांकित
सूत्र~$\sFmla{\True}{\eq[t][t]}$ जोडले आहे असे समजा. स्पष्टपणे
$\Sat{M}{\eq[t][t]}$, म्हणून $\Struct{M}$ हे
$\Gamma \cup \{\sFmla{\True}{\eq[t][t]}\}$ सुद्धा पूर्ण करते.

$\TRule{\True}{\eq}$ वापरून शाखेचा विस्तार केला असेल, तर
सूत्र~$\sFmla{S}{\varphi(t_2)}$ जोडतो; परंतु $\Gamma$ मध्ये
$\sFmla{\True}{\eq[t_1][t_2]}$ आणि $\sFmla{\True}{\varphi(t_1)}$ ही दोन्ही
असतात. म्हणून $\Sat{M}{\eq[t_1][t_2]}$ आणि $\Sat{M}{\varphi(t_1)}$.
$s$ हा $s(x) = \Value{t_1}{M}$ असलेला चल-विन्यास असू द्या.
\readerexternalref{fol:syn:ass:prop:sentence-sat-true} नुसार
$\Sat{M}{\varphi(t_1)}[s]$. $\varAssign{s}{s}{x}$ असल्यामुळे,
\readerexternalref{fol:syn:ext:prop:ext-formulas} नुसार $\Sat{M}{\varphi(x)}[s]$.
$\Sat{M}{\eq[t_1][t_2]}$ असल्यामुळे
$\Value{t_1}{M} = \Value{t_2}{M}$, आणि म्हणून
$s(x) = \Value{t_2}{M}$. \readerexternalref{fol:syn:ext:prop:ext-formulas} पुन्हा
लावल्याने $\Sat{M}{\varphi(t_2)}[s]$ सुद्धा मिळते.
\readerexternalref{fol:syn:ass:prop:sentence-sat-true} नुसार
$\Sat{M}{\varphi(t_2)}$. $\TRule{\False}{\eq}$ चे प्रकरण याचप्रमाणे आहे.
\end{proof}



\endgroup
\chapter{स्वयंसिद्धकीय निष्पत्ती}\label{fol:axd::chap}
\begin{editorial}
  या प्रकरणातील मजकूर कोणते संयोजक आणि संख्यापक आदिम किंवा परिभाषित
  आहेत हे दर्शवणाऱ्या विविध टॅग्जचे पालन करतो याची खात्री करण्याचा अद्याप
  कोणताही प्रयत्न केलेला नाही: यांपैकी $\liff$ परिभाषित आहे, तर बाकी
  सर्व आदिम आहेत असे गृहीत धरले आहे. FOL टॅग सत्य
  असल्यास आपण संख्यापकांसह आवृत्ती तयार करतो; अन्यथा संख्यापकांविना
  आवृत्ती तयार करतो.
\end{editorial}









\section{नियम आणि निष्पत्ती}\label{fol:axd:rul:sec}

\begin{explain}
  स्वयंसिद्धकीय निष्पत्त्या ही तर्कशास्त्रासाठी कदाचित सर्वांत सोपी
  निष्पत्ती-पद्धत आहे. निष्पत्ती म्हणजे केवळ सूत्रांचा
  अनुक्रम. निष्पत्ती म्हणून गणले जाण्यासाठी अनुक्रमातील प्रत्येक
  सूत्र एकतर स्वयंसिद्धकाचे उदाहरण असले पाहिजे, किंवा अनुक्रमात
  त्याच्या आधी येणाऱ्या एक वा अधिक सूत्रांपासून निगमन नियमाने
  निष्पन्न झाले पाहिजे. निष्पत्ती तिचे शेवटचे सूत्र
  निष्पन्न करते.
\end{explain}

\begin{defn}[निष्पन्नता]
जर $\Gamma$ हा $\Lang L$ मधील सूत्रांचा संच असेल, तर
$\Gamma$ पासूनची \emph{निष्पत्ती} म्हणजे सूत्रांचा सांत अनुक्रम
$\varphi_1$, \dots,~$\varphi_n$ असा की प्रत्येक $i \le n$ साठी पुढीलपैकी एक अट
पूर्ण होते:
\begin{enumerate}
\item $\varphi_i \in \Gamma$; किंवा
\item $\varphi_i$ हे स्वयंसिद्धक आहे; किंवा
\item $j < i$ (आणि $k < i$) असलेल्या एखाद्या $\varphi_j$ (आणि $\varphi_k$)
  पासून निगमन नियमाने $\varphi_i$ निष्पन्न होते.
\end{enumerate}
\end{defn}

योग्य निष्पत्ती म्हणून कशाची गणना होते हे आपण कोणते निगमन नियम
मान्य करतो (आणि अर्थातच कोणती सूत्रे स्वयंसिद्धके मानतो) यावर अवलंबून
असते. निगमन नियम म्हणजे असे जर-तर विधान जे विशिष्ट अटींमध्ये
निष्पत्तीमधील पायरी~$A_i$ ही निगमनाची योग्य पायरी आहे असे सांगते.

\begin{defn}[निगमन नियम]
\emph{निगमन नियम} हा $\Gamma$ पासूनच्या निष्पत्तीमध्ये निगमनाची
योग्य पायरी म्हणून कशाची गणना होते यासाठी पुरेशी अट देतो.
\end{defn}

उदाहरणार्थ, $\varphi \in \Gamma$ असलेला कोणताही एक-घटकी अनुक्रम $\varphi$ हा
आपोआपच निष्पत्ती म्हणून गणला जात असल्यामुळे, पुढील नियम हा
अतिशय सोपा निगमन नियम असू शकतो:
\begin{quote}
  जर $\varphi \in \Gamma$, तर $\Gamma$ पासूनच्या कोणत्याही निष्पत्तीमध्ये $\varphi$ ही नेहमीच निगमनाची योग्य पायरी असते.
\end{quote}
त्याचप्रमाणे, $\varphi$ हे स्वयंसिद्धकांपैकी एक असेल, तर एकटे $\varphi$ ही
निष्पत्ती असते; म्हणून पुढील विधानही निगमन नियम आहे:
\begin{quote}
  जर $\varphi$ हे स्वयंसिद्धक असेल, तर $\varphi$ ही निगमनाची योग्य पायरी असते.
\end{quote}
विचाराधीन पायरीच्या आधी दिसणाऱ्या सूत्रांचा निगमन नियम आधार घेतो
तेव्हा गोष्ट अधिक रोचक होते. पुढील नियमाला \emph{विध्यात्मक अनुमान}
म्हणतात:
\begin{quote}
  जर $\psi \lif \varphi$ आणि $\psi$ हे निष्पत्तीमध्ये यापूर्वी आढळत असतील,
  तर~$\varphi$ ही निगमनाची योग्य पायरी असते.
\end{quote}
हा एकमेव निगमन नियम असेल, तर वर दिलेल्या निष्पत्तीच्या व्याख्येचा
अर्थ असा होतो: प्रत्येक $i \le n$ साठी पुढीलपैकी एक अट पूर्ण होत असेल
तेव्हा आणि केवळ तेव्हाच $\varphi_1$, \dots,~$\varphi_n$ ही निष्पत्ती असते:
\begin{enumerate}
\item $\varphi_i \in \Gamma$; किंवा
\item $\varphi_i$ हे स्वयंसिद्धक आहे; किंवा
\item एखाद्या $j < i$ साठी $\varphi_j$ हे $\psi \lif \varphi_i$ आहे, आणि एखाद्या
  $k < i$ साठी $\varphi_k$ हे~$\psi$ आहे.
\end{enumerate}
शेवटचे कलम असे सांगते की~$\varphi_j$ ($\psi \lif \varphi_i$) आणि $\varphi_k$ ($\psi$)
यांपासून विध्यात्मक अनुमानाने $\varphi_i$ निष्पन्न होते. आपण $1$ पासून~$n$
पर्यंत जाऊ शकलो आणि प्रत्येक वेळी~$\Gamma$ मध्ये असलेले, स्वयंसिद्धक
असलेले किंवा निगमन नियमाने योग्य निगमन-पायरी ठरणारे सूत्र~$\varphi_i$
मिळाले, तर संपूर्ण अनुक्रम योग्य निष्पत्ती म्हणून गणला जातो.

\begin{defn}[निष्पन्नता]
जर $\Gamma$ पासून सुरू होऊन $\varphi$ वर संपणारी निष्पत्ती असेल, तर
सूत्र~$\varphi$ हे $\Gamma$ पासून \emph{निष्पन्न करता येण्याजोगे} असते; हे आपण
$\Gamma \Proves \varphi$ असे लिहितो.
\end{defn}

\begin{defn}[प्रमेये]
रिक्त संचापासून $\varphi$ ची निष्पत्ती असेल, तर सूत्र~$\varphi$ हे
\emph{प्रमेय} असते. $\varphi$ हे प्रमेय असेल तर आपण $\Proves \varphi$ लिहितो आणि
ते प्रमेय नसेल तर $\Proves/ \varphi$ लिहितो.
\end{defn}













\section{विधानीय संयोजकांसाठी स्वयंसिद्धके आणि नियम}\label{fol:axd:prp:sec}

\begin{defn}[स्वयंसिद्धके]
विधानीय संयोजकांसाठीच्या \emph{स्वयंसिद्धकांचा} संच $\PAx$ हा पुढील रूपांच्या
सर्व सूत्रांचा बनलेला आहे:
\begin{align}
  & (\varphi \land \psi) \lif \varphi \label{fol:axd:prp:ax:land1}\\
  & (\varphi \land \psi) \lif \psi \label{fol:axd:prp:ax:land2}\\
  & \varphi \lif (\psi \lif (\varphi \land \psi)) \label{fol:axd:prp:ax:land3}\\
  & \varphi \lif (\varphi \lor \psi) \label{fol:axd:prp:ax:lor1}\\
  & \varphi \lif (\psi \lor \varphi) \label{fol:axd:prp:ax:lor2}\\
  & (\varphi \lif \chi) \lif ((\psi \lif \chi) \lif ((\varphi \lor \psi) \lif \chi)) \label{fol:axd:prp:ax:lor3}\\
  & \varphi \lif (\psi \lif \varphi) \label{fol:axd:prp:ax:lif1}\\
  & (\varphi \lif (\psi \lif \chi)) \lif ((\varphi \lif \psi) \lif (\varphi \lif \chi)) \label{fol:axd:prp:ax:lif2}\\
  & (\varphi \lif \psi) \lif ((\varphi \lif \lnot \psi) \lif \lnot \varphi) \label{fol:axd:prp:ax:lnot1}\\
  & \lnot \varphi \lif (\varphi \lif \psi) \label{fol:axd:prp:ax:lnot2}\\
  & \ltrue \label{fol:axd:prp:ax:ltrue}\\
  & \lfalse \lif \varphi \label{fol:axd:prp:ax:lfalse1}\\
  & (\varphi \lif \lfalse) \lif \lnot \varphi \label{fol:axd:prp:ax:lfalse2}\\
  & \lnot\lnot \varphi \lif \varphi \label{fol:axd:prp:ax:dne}
\end{align}
\end{defn}

\begin{defn}[विध्यात्मक अनुमान]
  जर $\psi$ आणि $\psi \lif \varphi$ हे निष्पत्तीमध्ये आधीच आलेले असतील,
  तर $\varphi$ ही निगमनाची योग्य पायरी असते.
\end{defn}

विध्यात्मक अनुमान या नियमाचा संक्षेप आपण ``\MP'' असा करू.












\section{संख्यापकांसाठी स्वयंसिद्धके आणि नियम}\label{fol:axd:qua:sec}

\begin{defn}[संख्यापकांसाठी स्वयंसिद्धके]
संख्यापकांचे नियमन करणारी \emph{स्वयंसिद्धके} म्हणजे पुढील रूपांचे सर्व
दृष्टांत होत:
\begin{align}
\label{fol:axd:qua:ax:q1} & \lforall[x][\psi] \lif \psi(t), \\
\label{fol:axd:qua:ax:q2} & \psi(t) \lif \lexists[x][\psi].
\end{align}
कोणत्याही बंद पद~$t$ साठी.
\end{defn}

\begin{defn}[संख्यापकांसाठी नियम]
\item जर $\psi \lif \varphi(a)$ हे निष्पत्तीमध्ये आधीच आलेले असेल आणि
  $a$ ची~$\Gamma$ किंवा~$\psi$ मध्ये आस्थिति नसेल, तर $\psi \lif
  \lforall[x][\varphi(x)]$ ही निगमनाची योग्य पायरी असते.
\item जर $\varphi(a) \lif \psi$ हे निष्पत्तीमध्ये आधीच आलेले असेल आणि
  $a$ ची~$\Gamma$ किंवा~$\psi$ मध्ये आस्थिति नसेल, तर
  $\lexists[x][\varphi(x)] \lif \psi$ ही निगमनाची योग्य पायरी असते.
\end{defn}

यांपैकी कोणत्याही नियमाचा संक्षेप आपण ``\QR'' असा करू.












\section{निष्पत्ती: उदाहरणे}\label{fol:axd:pro:sec}


\begin{ex}
आपल्याला $(\lnot \theta \lor !E) \lif (\theta \lif
!E)$ हे सिद्ध करायचे आहे असे समजा. स्पष्टपणे, हे आपल्या कोणत्याही
स्वयंसिद्धकाचे उदाहरण नाही; म्हणून ते निष्पन्न करण्यासाठी \MP{} नियम
वापरावा लागेल. आपला एकमेव नियम~MP आहे; $\varphi$ आणि $\varphi \lif \psi$ दिले असता
या नियमाने~$\psi$ ला समर्थित करता येते. एक युक्ती अशी असेल: $\varphi$ म्हणून
$\lnot \theta$, $\psi$ म्हणून $!E$ आणि $\chi$ म्हणून $\theta \lif !E$ घेऊन
\ref{fol:axd:prp:ax:lor3} वापरणे, म्हणजे पुढील उदाहरण वापरणे:
\[
(\lnot \theta \lif (\theta \lif !E)) \lif ((!E \lif (\theta \lif !E)) \lif ((\lnot
\theta \lor !E) \lif (\theta \lif !E))).
\]
का? MP चे दोन उपयोग केल्यावर शेवटचा भाग मिळतो, आणि आपल्याला नेमके तेच
हवे आहे. तसेच $\lnot \theta \lif (\theta \lif !E)$ हे
\ref{fol:axd:prp:ax:lnot2} चे उदाहरण आणि $!E \lif (\theta \lif !E)$ हे
\ref{fol:axd:prp:ax:lif1} चे उदाहरण आहे, हे सहज दिसते. म्हणून आपली
निष्पत्ती अशी आहे:
\begin{derivation}
  1. &  $\lnot \theta \lif (\theta \lif !E)$ & \ref{fol:axd:prp:ax:lnot2} \\
  2. & $(\lnot \theta \lif (\theta \lif !E)) \lif {}$\\
  &\qquad $((!E \lif (\theta \lif !E)) \lif ((\lnot \theta \lor !E) \lif (\theta \lif !E)))$ & \ref{fol:axd:prp:ax:lor3}\\
  3. & $(!E \lif (\theta \lif !E)) \lif ((\lnot \theta \lor !E) \lif (\theta \lif !E))$ & 1, 2, \MP\\
  4. & $!E \lif (\theta \lif !E)$ & \ref{fol:axd:prp:ax:lif1}\\
  5. & $(\lnot \theta \lor !E) \lif (\theta \lif !E)$ & 3, 4, \MP
\end{derivation}
\end{ex}

\begin{ex}\label{fol:axd:pro:ex:identity}
$\theta \lif \theta$ ची निष्पत्ती शोधून पाहू. हे स्वयंसिद्धकाचे उदाहरण
नाही; म्हणून ते निष्पन्न करण्यासाठी आपल्याला \MP{} वापरावा लागेल.
\ref{fol:axd:prp:ax:lif1} हे~$\varphi \lif \psi$ या रूपाचे स्वयंसिद्धक आहे आणि त्याला
आपण~\MP लावू शकतो. त्याचा उपयोग व्हायचा असेल, तर अर्थातच या वेळी
\MP{} ने योग्य पायरी म्हणून समर्थित होणारे $\psi$ हे~$\theta \lif \theta$ असले
पाहिजे, कारण आपल्याला तेच निष्पन्न करायचे आहे. त्यामुळे $\varphi$ सुद्धा
$\theta$ असले पाहिजे; म्हणजे आपण \ref{fol:axd:prp:ax:lif1} चे पुढील उदाहरण पाहू शकतो:
\[
\theta \lif (\theta \lif \theta)
\]
\MP लावण्यासाठी आपल्याला संबंधित दुसरे आधारविधान, म्हणजे~$\varphi$, सुद्धा
समर्थित करावे लागेल. पण आपल्या प्रकरणात ते~$\theta$ असेल, आणि एकटे~$\theta$
निष्पन्न करता येणार नाही. म्हणून आपल्याला वेगळी युक्ती हवी.

केवळ~$\lif$ असलेले दुसरे स्वयंसिद्धक म्हणजे \ref{fol:axd:prp:ax:lif2}, म्हणजे
\[
(\varphi \lif (\psi \lif \chi)) \lif ((\varphi \lif \psi) \lif (\varphi \lif \chi))
\]
\MP{} दोनदा लावून आपण शेवटच्या अंतःस्थ सशर्त विधानापर्यंत पोहोचू शकतो.
पुन्हा, याचा अर्थ असा की \ref{fol:axd:prp:ax:lif2} चे असे उदाहरण आपल्याला हवे
ज्यात $\varphi \lif \chi$ हे $\theta \lif \theta$, म्हणजे आपण साध्य करू पाहत असलेले
सूत्र, असेल. मग अर्थातच $\varphi$ आणि $\chi$ ही दोन्ही~$\theta$ असतात.
$\varphi \lif (\psi \lif \chi)$ आणि $\varphi \lif \psi$, म्हणजे आपल्या प्रकरणात
$\theta \lif (\psi \lif \theta)$ आणि $\theta \lif \psi$, ही दोन्हीही निष्पन्न करता येण्याजोगे
होतील अशा रीतीने~$\psi$ कसे निवडावे? आपण $\psi$ काहीही निवडले, तरी यांपैकी
पहिले हे आधीच \ref{fol:axd:prp:ax:lif1} चे उदाहरण आहे. आणि $\psi$ हे
$(\theta \lif \theta)$ असेल, तर $\theta \lif \psi$ हे \ref{fol:axd:prp:ax:lif1} चे आणखी
एक उदाहरण असेल. म्हणून आपली निष्पत्ती अशी आहे:
\begin{derivation}
  1. & $\theta \lif ((\theta \lif \theta) \lif \theta)$ & \ref{fol:axd:prp:ax:lif1}\\
  2. & $(\theta \lif ((\theta \lif \theta) \lif \theta)) \lif {}$\\
  & \qquad $((\theta \lif (\theta \lif \theta)) \lif (\theta \lif \theta))$ & \ref{fol:axd:prp:ax:lif2}\\
  3. & $(\theta \lif (\theta \lif \theta)) \lif (\theta \lif \theta)$ & 1, 2, \MP\\
  4. & $\theta \lif (\theta \lif \theta)$ & \ref{fol:axd:prp:ax:lif1}\\
  5. & $\theta \lif \theta$ & 3, 4, \MP
\end{derivation}
\end{ex}

\begin{ex}\label{fol:axd:pro:ex:chain}
कधी कधी काही इतर सूत्रे~$\Gamma$ पासून एखाद्या सूत्राची
निष्पत्ती आहे हे आपल्याला दाखवायचे असते. उदाहरणार्थ,
$\Gamma = \{\varphi \lif \psi, \psi \lif \chi\}$ पासून $\varphi \lif \chi$ हे
निष्पन्न करता येते, हे दाखवू.
\begin{derivation}
  1. & $\varphi \lif \psi$ & \Hyp\\
  2. & $\psi \lif \chi$ & \Hyp\\
  3. & $(\psi \lif \chi) \lif (\varphi \lif (\psi \lif \chi))$ & \ref{fol:axd:prp:ax:lif1} \\
  4. & $\varphi \lif (\psi \lif \chi)$ & 2, 3, \MP\\
  5. & $(\varphi \lif (\psi \lif \chi)) \lif {}$\\
  & \qquad $((\varphi \lif \psi) \lif (\varphi \lif \chi))$ & \ref{fol:axd:prp:ax:lif2}\\
  6. & $((\varphi \lif \psi) \lif (\varphi \lif \chi))$ & 4, 5, \MP\\
  7. & $\varphi \lif \chi$ & 1, 6, \MP
\end{derivation}
``\Hyp'' (``hypothesis'', म्हणजे ``गृहीतक'') असे लेबल असलेल्या ओळी
त्या ओळीवरील सूत्र हा~$\Gamma$ चा घटक असल्याचे दर्शवतात.
\end{ex}

\begin{prop}
  \label{fol:axd:pro:prop:chain} जर $\Gamma \Proves \varphi \lif \psi$ आणि $\Gamma
  \Proves \psi \lif \chi$, तर $\Gamma \Proves \varphi \lif \chi$
\end{prop}

\begin{proof}
  $\Gamma \Proves \varphi \lif \psi$ आणि $\Gamma \Proves \psi \lif
  \chi$ असे समजा. मग~$\Gamma$ पासून $\varphi \lif \psi$ ची निष्पत्ती
  असते; तसेच~$\Gamma$ पासून $\psi \lif \chi$ चीही निष्पत्ती असते.
  त्यांची जोडणी करून त्या एकाच निष्पत्तीमध्ये एकत्र करा. आता मागील
  उदाहरणातील निष्पत्तीच्या ओळी 3--7 जोडा. ही~$\Gamma$ पासूनची
  निष्पत्ती आहे, आणि नव्या निष्पत्तीची शेवटची ओळ
  $\varphi \lif \chi$ आहे. लक्षात घ्या की ओळी 4 आणि~7 साठी दिलेली
  समर्थने पुढील बदलांनंतरही वैध राहतात: ओळ क्रमांक~1 चा निर्देश $\varphi
  \lif \psi$ च्या निष्पत्तीच्या शेवटच्या ओळीच्या निर्देशाने आणि ओळ
  क्रमांक~2 चा निर्देश $\psi \lif \chi$ च्या निष्पत्तीच्या शेवटच्या
  ओळीच्या निर्देशाने बदला.
\end{proof}

\begin{prob}
स्वयंसिद्धकांपासूनच्या निष्पत्त्या मांडून पुढील विधाने सिद्ध करा:
\begin{enumerate}
\item $(\varphi \land \psi) \lif (\psi \land \varphi)$
\item $((\varphi \land \psi) \lif \chi) \lif (\varphi \lif (\psi \lif \chi))$
\item $\lnot(\varphi \lor \psi) \lif \lnot \varphi$
\end{enumerate}
\end{prob}














\section{संख्यापकांसह निष्पत्ती}\label{fol:axd:prq:sec}

\begin{ex}
$(\lforall[x][\varphi(x)] \land
\lforall[y][\psi(y)]) \lif \lforall[x][(\varphi(x) \land \psi(x))]$ ची
निष्पत्ती देऊ.

प्रथम, लक्षात घ्या की
\begin{align*}
  (\lforall[x][\varphi(x)] \land \lforall[y][\psi(y)]) & \lif \lforall[x][\varphi(x)]\\
  \intertext{हे \ref{fol:axd:prp:ax:land1} चे उदाहरण आहे, आणि}
  \lforall[x][\varphi(x)] & \lif \varphi(a) \\
  \intertext{हे \ref{fol:axd:qua:ax:q1} चे उदाहरण आहे. म्हणून \ref{fol:axd:pro:prop:chain} वरून आपल्याला माहीत आहे की}
  (\lforall[x][\varphi(x)] \land \lforall[y][\psi(y)]) & \lif \varphi(a) \\
  \intertext{हे निष्पन्न करता येण्याजोगे आहे. त्याचप्रमाणे,}
  (\lforall[x][\varphi(x)] \land \lforall[y][\psi(y)]) & \lif \lforall[y][\psi(y)] \qquad\text{आणि}\\
    \lforall[y][\psi(y)] & \lif \psi(a)\\
    \intertext{ही अनुक्रमे \ref{fol:axd:prp:ax:land2} आणि \ref{fol:axd:qua:ax:q1} यांची उदाहरणे असल्यामुळे,}
    (\lforall[x][\varphi(x)] \land \lforall[y][\psi(y)]) & \lif \psi(a)\\
    \intertext{हे \ref{fol:axd:pro:prop:chain} मुळे सिद्धतायोग्य आहे. \ref{fol:axd:prp:ax:land3} चे योग्य उदाहरण आणि~\MP चे दोन उपयोग करून आपल्याला दिसते की}
    (\lforall[x][\varphi(x)] \land \lforall[y][\psi(y)]) & \lif (\varphi(a) \land \psi(a))\\
    \intertext{हे सिद्धतायोग्य आहे. आता \QR{} लावून पुढील सूत्र मिळते:}
(\lforall[x][\varphi(x)] \land \lforall[y][\psi(y)]) & \lif \lforall[x][(\varphi(x) \land \psi(x))].
\end{align*}
\end{ex}












\section{सिद्धता-उपपत्तीय संकल्पना}\label{fol:axd:ptn:sec}

\begin{explain}
जशा आपण अनेक महत्त्वाच्या चिन्हार्थविषयक संकल्पना
(वैधता, तार्किक निष्पन्नता, पूर्ततायोग्यता)
परिभाषित केल्या आहेत, तशाच आता त्यांना अनुरूप
\emph{सिद्धता-उपपत्तीय संकल्पना} परिभाषित करू. या संकल्पना संरचना
मध्ये वाक्यांची पूर्ति होते की नाही याचा आधार घेऊन परिभाषित केल्या
जात नाहीत; त्याऐवजी विशिष्ट सूत्रांची निष्पन्नता किंवा
अनिष्पन्नता यांचा आधार घेतला जातो. या दोन प्रकारच्या संकल्पना
एकमेकांशी जुळतात, हा एक महत्त्वाचा शोध होता. त्या जुळतात, हाच
\emph{निर्दोषता प्रमेय} आणि \emph{संपूर्णता प्रमेय} यांचा आशय आहे.
\end{explain}

\begin{defn}[निष्पन्नता]
सूत्र~$\varphi$ हे $\Gamma$ \emph{पासून निष्पन्न करता येण्याजोगे} असते, म्हणजे
$\Gamma \Proves \varphi$, जर~$\Gamma$ पासून सुरू होऊन~$\varphi$ वर संपणारी
निष्पत्ती असेल.
\end{defn}

\begin{defn}[प्रमेये]
रिक्त संचापासून $\varphi$ ची निष्पत्ती असेल, तर सूत्र~$\varphi$ हे
\emph{प्रमेय} असते. $\varphi$ हे प्रमेय असेल तर आपण $\Proves \varphi$ लिहितो आणि
ते प्रमेय नसेल तर $\Proves/ \varphi$ लिहितो.
\end{defn}

\begin{defn}[सुसंगतता]
सूत्रांचा संच~$\Gamma$ हा $\Gamma\Proves/ \lfalse$ असेल तेव्हा आणि
केवळ तेव्हाच \emph{सुसंगत} असतो; अन्यथा तो \emph{विसंगत} असतो.
\end{defn}

\begin{prop}[परावर्तिता]
\label{fol:axd:ptn:prop:reflexivity}
$\varphi \in \Gamma$ असेल, तर $\Gamma \Proves \varphi$.
\end{prop}

\begin{proof}
  एकटे सूत्र~$\varphi$ ही~$\Gamma$ पासून $\varphi$ ची निष्पत्ती आहे.
\end{proof}

\begin{prop}[एकस्वनिकता]
\label{fol:axd:ptn:prop:monotonicity}
$\Gamma \subseteq \Delta$ आणि $\Gamma \Proves \varphi$ असेल, तर $\Delta
\Proves \varphi$.
\end{prop}

\begin{proof}
$\Gamma$ पासून $\varphi$ ची कोणतीही निष्पत्ती ही~$\Delta$ पासून $\varphi$ चीही
निष्पत्ती असते.
\end{proof}

\begin{prop}[संक्रमकता]
\label{fol:axd:ptn:prop:transitivity}
$\Gamma \Proves \varphi$ आणि $\{\varphi\} \cup \Delta \Proves
\psi$ असेल, तर $\Gamma \cup \Delta \Proves \psi$.
\end{prop}

\begin{proof}
  $\{\varphi\} \cup \Delta \Proves \psi$ असे समजा. मग~$\{\varphi\} \cup
  \Delta$ पासून $\psi_1$, \dots, $\psi_l = \psi$ ही निष्पत्ती असते.
  त्या निष्पत्तीमधील काही पायऱ्या योग्य असतात, कारण एखादा नियम
  त्यांपूर्वीच्या~$\psi_i = \varphi$ या ओळीचा निर्देश करतो. गृहीतकानुसार,
  $\Gamma$ पासून~$\varphi$ ची निष्पत्ती असते, म्हणजे निष्पत्ती
  $\varphi_1$, \dots, $\varphi_k = \varphi$ अशी असते की प्रत्येक $\varphi_i$ हे स्वयंसिद्धक,
  $\Gamma$ चा घटक, किंवा निगमन नियमाने योग्य ठरलेले असते. आता
  पुढील अनुक्रम विचारात घ्या:
  \[
  \varphi_1, \dots, \varphi_k = \varphi, \psi_1, \dots, \psi_l = \psi.
  \]
  ही $\Gamma \cup \Delta$ पासून~$\psi$ ची योग्य निष्पत्ती आहे, कारण
  प्रत्येक $B_i = \varphi$ ला आता~$\varphi_k = \varphi$ ला समर्थित करणाऱ्या त्याच नियमाने
  समर्थन मिळते.
\end{proof}

लक्षात घ्या की विशेषतः याचा अर्थ असा होतो: $\Gamma \Proves \varphi$ आणि $\varphi
\Proves \psi$ असेल, तर $\Gamma \Proves \psi$. तसेच $\varphi_1,
\dots, \varphi_n \Proves \psi$ असेल आणि प्रत्येक~$i$ साठी $\Gamma \Proves \varphi_i$
असेल, तर $\Gamma \Proves \psi$, हेसुद्धा यावरून मिळते.

\begin{prop}
\label{fol:axd:ptn:prop:incons}
$\Gamma$ विसंगत असतो तेव्हा आणि केवळ तेव्हाच, जेव्हा प्रत्येक~$\varphi$ साठी
$\Gamma \Proves \varphi$.
\end{prop}

\begin{proof}
सराव.
\end{proof}


\begin{prob}
\ref{fol:axd:ptn:prop:incons} सिद्ध करा.
\end{prob}




\begin{prop}[संहतता]
\label{fol:axd:ptn:prop:proves-compact}
  \begin{enumerate}
  \item $\Gamma \Proves \varphi$ असेल, तर एखादा सांत उपसंच $\Gamma_0
    \subseteq \Gamma$ असा असतो की $\Gamma_0 \Proves \varphi$.
  \item $\Gamma$ चा प्रत्येक सांत उपसंच सुसंगत असेल, तर $\Gamma$ सुसंगत
    असतो.
  \end{enumerate}
\end{prop}

\begin{proof}
  \begin{enumerate}
    \item $\Gamma \Proves \varphi$ असेल, तर सूत्रांचा सांत अनुक्रम
      $\varphi_1$, \dots,~$\varphi_n$ असा असतो की $\varphi \ident \varphi_n$ आणि प्रत्येक
      $\varphi_i$ हे एकतर तर्कशास्त्रीय स्वयंसिद्धक, $\Gamma$ चा घटक,
      किंवा आधीच्या सूत्रांपासून विध्यात्मक अनुमानाने निष्पन्न झालेले
      असते. $\Gamma$ मध्ये असलेल्या त्या $\varphi_i$ चा संच~$\Gamma_0$ म्हणून
      घ्या. मग ही निष्पत्ती त्याचप्रमाणे~$\Gamma_0$ पासूनची
      निष्पत्ती असते, आणि म्हणून $\Gamma_0 \Proves \varphi$.
    \item (1) मध्ये $\varphi
      \ident \lfalse$ हे विशेष प्रकरण घेतल्यावर मिळणारे हे निषेधात्मक उलट
      विधान आहे.
\end{enumerate}
\end{proof}















\section{निगमन प्रमेय}\label{fol:axd:ded:sec}

आपण पाहिल्याप्रमाणे, स्वयंसिद्धकीय पद्धतीत निष्पत्त्या देणे किचकट
असते आणि निष्पत्त्या शोधणे कठीण असू शकते. सूत्रांच्या लांब
याद्या प्रत्यक्ष लिहिण्याऐवजी काही सोपे निष्कर्ष वापरून अशा
निष्पत्त्या अस्तित्वात आहेत असा युक्तिवाद करणे सहसा सोपे असते. असे तीन
निष्कर्ष आपण आधीच प्रस्थापित केले आहेत: $\varphi \in \Gamma$ हे माहीत असेल,
तर $\Gamma \Proves \varphi$ असे नेहमी म्हणता येते, हे
\ref{fol:axd:ptn:prop:reflexivity} सांगतो. $\Gamma \Proves \varphi$ असेल, तर
$\Gamma \cup \{\psi\} \Proves \varphi$ सुद्धा असते, हे
\ref{fol:axd:ptn:prop:monotonicity} सांगतो. आणि $\Gamma \Proves \varphi$ व
$\varphi \Proves \psi$ असेल, तर $\Gamma \Proves \psi$, असे
\ref{fol:axd:ptn:prop:transitivity} वरून मिळते. पुढे आणखी एक सोपा निष्कर्ष,
विध्यात्मक अनुमानाची ``अधि''-आवृत्ती, दिली आहे:

\begin{prop}
  \label{fol:axd:ded:prop:mp} जर $\Gamma \Proves \varphi$ आणि $\Gamma \Proves \varphi \lif
  \psi$, तर $\Gamma \Proves \psi$.
\end{prop}

\begin{proof}
  $\{\varphi, \varphi \lif \psi\} \Proves \psi$ हे आपल्याकडे आहे:
  \begin{derivation}
    1. & $\varphi$ & Hyp.\\
    2. & $\varphi \lif \psi$ & Hyp.\\
    3. & $\psi$ & 1, 2, MP
  \end{derivation}
  \ref{fol:axd:ptn:prop:transitivity} वरून $\Gamma \Proves \psi$.
\end{proof}

या संदर्भात आपण वापरणार असलेला सर्वांत महत्त्वाचा निष्कर्ष म्हणजे निगमन
प्रमेय:

\begin{thm}[निगमन प्रमेय]
\label{fol:axd:ded:thm:deduction-thm} $\Gamma \cup \{\varphi\} \Proves \psi$ तेव्हा आणि
केवळ तेव्हाच, जेव्हा $\Gamma \Proves \varphi \lif \psi$.
\end{thm}

\begin{proof}
``जर'' ही दिशा तात्काळ मिळते. $\Gamma \Proves \varphi \lif \psi$ असेल, तर
\ref{fol:axd:ptn:prop:monotonicity} वरून $\Gamma \cup \{\varphi\} \Proves \varphi \lif \psi$
सुद्धा असते. तसेच \ref{fol:axd:ptn:prop:reflexivity} वरून
$\Gamma \cup \{\varphi\} \Proves \varphi$. म्हणून \ref{fol:axd:ded:prop:mp} वरून
$\Gamma \cup \{\varphi\} \Proves \psi$.

``केवळ जर'' या दिशेसाठी, $\Gamma \cup \{\varphi\}$ पासून $\psi$ च्या
निष्पत्तीच्या लांबीवर विगमन करू.

विगमनाच्या आधारपायरीसाठी, लांबी~$1$ असलेल्या प्रत्येक निष्पत्ती साठी
दावा सिद्ध करू. $\Gamma \cup \{\varphi\}$ पासून~$\psi$ ची लांबी~$1$ असलेली
निष्पत्ती म्हणजे एकटे $\psi$; ती योग्य असेल, तर $\psi$ हे एकतर
$\in \Gamma \cup \{\varphi\}$ असते किंवा स्वयंसिद्धक असते. $\psi \in \Gamma$
असेल किंवा ते स्वयंसिद्धक असेल, तर $\Gamma \Proves \psi$. तसेच
\ref{fol:axd:prp:ax:lif1} वरून $\Gamma \Proves \psi \lif (\varphi \lif \psi)$, आणि
\ref{fol:axd:ded:prop:mp} वरून $\Gamma \Proves \varphi \lif \psi$. $\psi \in \{ \varphi\}$
असेल, तर $\Gamma \Proves \varphi \lif \psi$, कारण शेवटचे वाक्य~$\varphi
\lif \psi$ हे $\varphi \lif \varphi$ सारखेच आहे आणि ते आपण
\ref{fol:axd:pro:ex:identity} मध्ये निष्पन्न केले आहे.

विगमन पायरीसाठी, $\Gamma \cup \{\varphi\}$ पासून~$\psi$ ची एखादी
निष्पत्ती विध्यात्मक अनुमानाने समर्थित असलेल्या~$\psi$ या पायरीने
संपते असे समजा. (तिला विध्यात्मक अनुमानाने समर्थन मिळत नसेल, तर
$\psi \in \Gamma$, $\psi \ident \varphi$, किंवा $\psi$ हे स्वयंसिद्धक असते, आणि
विगमनाच्या आधारपायरीतील युक्तिवादच लागू होतो.) मग एखाद्या
सूत्र~$\chi$ साठी त्या निष्पत्तीमधील आधीच्या काही पायऱ्या
$\chi \lif \psi$ आणि $\chi$ असतात; म्हणजे $\Gamma \cup \{\varphi\} \Proves \chi \lif \psi$
आणि $\Gamma \cup \{\varphi\} \Proves \chi$. त्यांच्याशी संबंधित
निष्पत्त्या लहान असल्यामुळे त्यांना विगमन गृहीतक लागू होते. म्हणून
आपल्याकडे पुढील दोन्ही आहेत:
\begin{align*}
  & \Gamma \Proves \varphi \lif (\chi \lif \psi); \\
  & \Gamma \Proves \varphi \lif \chi.
\end{align*}
पण \ref{fol:axd:prp:ax:lif2} वरून पुढीलही मिळते:
\[
\Gamma \Proves (\varphi \lif (\chi \lif \psi)) \lif
((\varphi\lif \chi)  \lif (\varphi \lif \psi)),
\]
आणि \ref{fol:axd:ded:prop:mp} चे दोन उपयोग केल्यावर आवश्यकतेनुसार
$\Gamma \Proves \varphi \lif \psi$ मिळते.
\end{proof}

निगमन प्रमेय खरे ठरावे यासाठीच \ref{fol:axd:prp:ax:lif1} आणि
\ref{fol:axd:prp:ax:lif2} यांची नेमकी अशी निवड केली होती, हे लक्षात घ्या.

निष्पन्नता विषयीची पुढील काही उपयुक्त तथ्ये सरावासाठी सोडली आहेत.

\begin{prop}
  \label{fol:axd:ded:prop:derivfacts}
\begin{enumerate}
\item $\Proves (\varphi \lif \psi) \lif ((\psi \lif \chi)
  \lif (\varphi \lif \chi))$; \label{fol:axd:ded:derivfacts:a}

\item $\Gamma \cup \{ \lnot \varphi\}
  \Proves \lnot \psi$ असेल, तर $\Gamma \cup \{ \psi\} \Proves
  \varphi$ (प्रतिपरिवर्तन); \label{fol:axd:ded:derivfacts:b}
\item  $\{ \varphi, \lnot\varphi\} \Proves
    \psi$ (Ex Falso Quodlibet, म्हणजे असत्यापासून काहीही; स्फोट); \label{fol:axd:ded:derivfacts:c}
\item  $\{ \lnot\lnot\varphi\} \Proves
  \varphi$ (द्विनकरण विलोपन);\label{fol:axd:ded:derivfacts:d}
\item $\Gamma \Proves \lnot\lnot\varphi$ असेल, तर $\Gamma \Proves
  \varphi$;\label{fol:axd:ded:derivfacts:e}
\end{enumerate}
\end{prop}


\begin{prob}
  \ref{fol:axd:ded:prop:derivfacts} सिद्ध करा.
\end{prob}















\section{संख्यापकांसह निगमन प्रमेय}\label{fol:axd:ddq:sec}

\begin{thm}[निगमन प्रमेय]
\label{fol:axd:ddq:thm:deduction-thm-q} जर $\Gamma \cup \{\varphi\} \Proves \psi$ असेल,
तर $\Gamma \Proves \varphi \lif \psi$.
\end{thm}

\begin{proof}
$\Gamma \cup \{\varphi\}$ पासून $\psi$ च्या निष्पत्तीच्या लांबीवर
आपण पुन्हा विगमन करू.

विगमनाच्या आधारपायरीचा पुरावा~\ref{fol:axd:ded:thm:deduction-thm} च्या
पुराव्यातील आधारपायरीच्या पुराव्यासारखाच आहे.

विगमन पायरीसाठी, $\Gamma \cup \{\varphi\}$ पासून $\psi$ ची निष्पत्ती
एखाद्या निगमन नियमाने समर्थित असलेल्या~$\psi$ या पायरीने संपते असे पुन्हा
समजा. निगमन नियम विध्यात्मक अनुमान असेल, तर आपण
\ref{fol:axd:ded:thm:deduction-thm} च्या पुराव्याप्रमाणेच पुढे जाऊ. निगमन
नियम \QR असेल, तर $\psi \ident \chi \lif \lforall[x][\theta(x)]$ आहे आणि
$\chi \lif \theta(a)$ या रूपाचे सूत्र त्या निष्पत्तीमध्ये आधी
येते, जेथे $a$ हे~$\chi$, $\varphi$, किंवा $\Gamma$ मध्ये येत नाही. म्हणून
आपल्याकडे पुढील आहे:
\begin{align*}
  \Gamma \cup \{\varphi\} & \Proves \chi \lif \theta(a),\\
  \intertext{आणि विगमन गृहीतक लागू होते; म्हणजे पुढील मिळते:}
    \Gamma & \Proves \varphi \lif (\chi \lif \theta(a)).\\
  \intertext{पुढील सूत्रावरून}
  & \Proves (\varphi \lif (\chi \lif \theta(a))) \lif ((\varphi \land \chi) \lif \theta(a))\\
  \intertext{आणि विध्यात्मक अनुमानामुळे पुढील मिळते:}
  \Gamma & \Proves (\varphi \land \chi) \lif \theta(a).\\
  \intertext{आयगेन चराची अट अजूनही लागू असल्यामुळे, या निष्पत्तीमध्ये \QR ने समर्थित एक पायरी जोडून पुढील मिळवता येते:}
    \Gamma & \Proves (\varphi \land \chi) \lif \lforall[x][\theta(x)].\\
    \intertext{आपल्याकडे पुढीलही आहे:}
    & \Proves ((\varphi \land \chi) \lif \lforall[x][\theta(x)]) \lif (\varphi \lif (\chi \lif \lforall[x][\theta(x)])),\\
    \intertext{म्हणून विध्यात्मक अनुमानाने,}
    \Gamma & \Proves \varphi \lif (\chi \lif \lforall[x][\theta(x)]),
\end{align*}

म्हणजे $\Gamma \Proves \varphi \lif \psi$.


$\psi$ ला \QR या नियमाने समर्थन मिळते, पण त्याचे रूप
$\lexists[x][\theta(x)] \lif \chi$ आहे, हे प्रकरण आपण सरावासाठी सोडतो.
\end{proof}

\begin{prob}
  \ref{fol:axd:ddq:thm:deduction-thm-q} चा पुरावा पूर्ण करा.
\end{prob}












\section{निष्पन्नता आणि सुसंगतता}\label{fol:axd:prv:sec}

आता आपण निष्पन्नता संबंधाचे अनेक गुणधर्म सिद्ध करू. ते स्वतंत्रपणेही
महत्त्वाचे आहेत, पण संपूर्णता प्रमेयाच्या सिद्धतेत प्रत्येकाची भूमिका असेल.

\begin{prop}\label{fol:axd:prv:prop:provability-contr}
  जर $\Gamma \Proves \varphi$ आणि $\Gamma \cup \{\varphi\}$ विसंगत असेल, तर
  $\Gamma$ विसंगत आहे.
\end{prop}

\begin{proof}
  जर $\Gamma \cup \{\varphi\}$ विसंगत असेल, तर $\Gamma \cup \{\varphi\}
  \Proves \lfalse$. \ref{fol:axd:ptn:prop:reflexivity} वरून प्रत्येक
  $\psi \in \Gamma$ साठी $\Gamma \Proves \psi$. गृहीतकानुसार
  $\Gamma \Proves \varphi$ सुद्धा असल्यामुळे प्रत्येक $\psi \in \Gamma \cup
  \{\varphi\}$ साठी $\Gamma \Proves \psi$. \ref{fol:axd:ptn:prop:transitivity}
  वरून $\Gamma \Proves \lfalse$, म्हणजे $\Gamma$ विसंगत आहे.
\end{proof}

\begin{prop}
\label{fol:axd:prv:prop:prov-incons}
$\Gamma \Proves \varphi$ तेव्हा आणि केवळ तेव्हाच, जेव्हा $\Gamma \cup \{\lnot \varphi\}$
विसंगत असते.
\end{prop}

\begin{proof}
प्रथम $\Gamma \Proves \varphi$ असे समजा. मग \ref{fol:axd:ptn:prop:monotonicity}
वरून $\Gamma \cup \{\lnot \varphi\} \Proves \varphi$. \ref{fol:axd:ptn:prop:reflexivity}
वरून $\Gamma \cup \{\lnot \varphi\} \Proves \lnot \varphi$. तसेच
\ref{fol:axd:prp:ax:lnot2} वरून $\Proves \lnot \varphi \lif (\varphi \lif \lfalse)$.
म्हणून \ref{fol:axd:ded:prop:mp} चे दोन उपयोग करून $\Gamma \cup \{\lnot
\varphi\} \Proves \lfalse$ मिळते.

आता $\Gamma \cup \{\lnot \varphi\}$ विसंगत आहे असे गृहीत धरा, म्हणजे
$\Gamma \cup \{\lnot \varphi\} \Proves \lfalse$. निगमन प्रमेयावरून $\Gamma
\Proves \lnot \varphi \lif \lfalse$. \ref{fol:axd:prp:ax:lfalse2} वरून $\Gamma
\Proves (\lnot \varphi \lif \lfalse) \lif \lnot\lnot \varphi$, म्हणून
\ref{fol:axd:ded:prop:mp} वरून $\Gamma \Proves \lnot\lnot \varphi$. आपल्याकडे
$\Gamma \Proves \lnot\lnot \varphi \lif \varphi$ हेही असल्यामुळे
(\ref{fol:axd:prp:ax:dne}), \ref{fol:axd:ded:prop:mp} च्या आणखी एका उपयोगाने
$\Gamma \Proves \varphi$ मिळते.
\end{proof}

\begin{prob}
$\Gamma \Proves \lnot \varphi$ तेव्हा आणि केवळ तेव्हाच, जेव्हा $\Gamma \cup \{\varphi\}$
विसंगत असते, हे सिद्ध करा.
\end{prob}

\begin{prop}\label{fol:axd:prv:prop:explicit-inc}
  जर $\Gamma \Proves \varphi$ आणि $\lnot \varphi \in \Gamma$ असेल, तर $\Gamma$
  विसंगत आहे.
\end{prop}

\begin{proof}
  \ref{fol:axd:prp:ax:lnot2} वरून $\Gamma \Proves \lnot \varphi \lif (\varphi \lif \lfalse)$.
  \ref{fol:axd:ded:prop:mp} चे दोन उपयोग करून $\Gamma \Proves \lfalse$ मिळते.
\end{proof}


\begin{prop}\label{fol:axd:prv:prop:provability-exhaustive}
  $\Gamma \cup \{\varphi\}$ आणि $\Gamma \cup \{\lnot \varphi\}$ हे दोन्ही संच
  विसंगत असतील, तर $\Gamma$ विसंगत आहे.
\end{prop}

\begin{proof}
सराव.
\end{proof}


\begin{prob}
  \ref{fol:axd:prv:prop:provability-exhaustive} सिद्ध करा
\end{prob}
















\section{निष्पन्नता आणि विधानीय संयोजक}\label{fol:axd:ppr:sec}

\begin{explain}
  स्वयंसिद्धकीय निगमनाचा निष्पन्नता संबंध~$\Proves$ हा विधानीय
  संयोजकांशी संबंधित काही मूलभूत तथ्ये सिद्ध करण्यासाठी पुरेसा सामर्थ्यवान
  आहे; उदाहरणार्थ, $\varphi \land \psi \Proves \varphi$ आणि $\varphi, \varphi \lif \psi
  \Proves \psi$ (विध्यात्मक अनुमान). ही तथ्ये संपूर्णता प्रमेयाच्या
  सिद्धतेसाठी आवश्यक आहेत.
\end{explain}

\begin{prop}\label{fol:axd:ppr:prop:provability-land}
  \begin{enumerate}
  \item \label{fol:axd:ppr:prop:provability-land-left} $\varphi \land \psi \Proves
    \varphi$ आणि $\varphi \land \psi \Proves \psi$ ही दोन्ही तथ्ये सत्य आहेत.
  \item \label{fol:axd:ppr:prop:provability-land-right} $\varphi, \psi \Proves \varphi \land \psi$.
  \end{enumerate}
\end{prop}

\begin{proof}
  \begin{enumerate}
    \item \ref{fol:axd:prp:ax:land1} आणि \ref{fol:axd:prp:ax:land2} यांपासून,
      अनुक्रमे विध्यात्मक अनुमान लावून.

  \item \ref{fol:axd:prp:ax:land3} पासून विध्यात्मक अनुमान दोनदा लावून.
  \end{enumerate}
\end{proof}


\begin{prop}\label{fol:axd:ppr:prop:provability-lor}
  \begin{enumerate}
  \item $\varphi \lor \psi, \lnot \varphi, \lnot \psi$ विसंगत आहे.
  \item $\varphi \Proves \varphi \lor \psi$ आणि $\psi \Proves \varphi \lor \psi$ ही दोन्ही
    तथ्ये सत्य आहेत.
  \end{enumerate}
\end{prop}

\begin{proof}
  \begin{enumerate}
  \item \ref{fol:axd:prp:ax:lnot2} पासून $\Proves \lnot \varphi \lif (\varphi
    \lif \lfalse)$ आणि $\Proves \lnot \psi \lif (\psi \lif \lfalse)$
    मिळतात. म्हणून निगमन प्रमेयाने $\{\lnot \varphi\} \Proves \varphi \lif
    \lfalse$ आणि $\{\lnot \psi\} \Proves \psi \lif \lfalse$ मिळतात.
    \ref{fol:axd:prp:ax:lor3} पासून $\{\lnot \varphi, \lnot \psi\} \Proves (\varphi
    \lor \psi) \lif \lfalse$ मिळते. निगमन प्रमेयाने $\{\varphi \lor \psi,
    \lnot \varphi, \lnot \psi\} \Proves \lfalse$.

  \item \ref{fol:axd:prp:ax:lor1} आणि \ref{fol:axd:prp:ax:lor2} यांपासून विध्यात्मक
    अनुमानाने.
  \end{enumerate}
\end{proof}

\begin{prop}\label{fol:axd:ppr:prop:provability-lif}
  \begin{enumerate}
  \item \label{fol:axd:ppr:prop:provability-lif-left} $\varphi, \varphi \lif \psi \Proves \psi$.
  \item \label{fol:axd:ppr:prop:provability-lif-right}
    $\lnot \varphi \Proves \varphi \lif \psi$ आणि $\psi \Proves \varphi \lif \psi$ ही दोन्ही
    तथ्ये सत्य आहेत.
  \end{enumerate}
\end{prop}

\begin{proof}
  \begin{enumerate}
  \item आपण पुढील निष्पन्न करू शकतो:
    \begin{derivation}
      1. & $\varphi$ & \Hyp\\
      2. & $\varphi \lif \psi$ & \Hyp\\
      3. & $\psi$ & 1, 2, \MP
    \end{derivation}

  \item अनुक्रमे \ref{fol:axd:prp:ax:lnot2}, \ref{fol:axd:prp:ax:lif1} आणि निगमन
    प्रमेय वापरून.
  \end{enumerate}
\end{proof}















\section{निष्पन्नता आणि संख्यापक}\label{fol:axd:qpr:sec}

\begin{explain}
  संपूर्णता प्रमेयासाठी या विभागात स्थापित केलेली~$\Proves$ संबंधी तथ्ये
  स्वयंसिद्धकीय निगमनातून निष्पन्न होतात, हेही आवश्यक आहे.
\end{explain}

\begin{thm}
\label{fol:axd:qpr:thm:strong-generalization} जर $c$ हे $\Gamma$ किंवा $\varphi(x)$ मध्ये
न येणारे स्थिरांक असेल आणि $\Gamma \Proves \varphi(c)$, तर $\Gamma
\Proves \lforall[x][\varphi(x)]$.
\end{thm}

\begin{proof}
निगमन प्रमेयाने $\Gamma \Proves \ltrue \lif \varphi(c)$. $c$ हे
$\Gamma$ किंवा~$\top$ मध्ये येत नसल्यामुळे $\Gamma \Proves
\ltrue \lif \lforall[x][\varphi(x)]$ मिळते. त्यानंतर सत्य स्थिरांकाचे
स्वयंसिद्धक आणि विध्यात्मक अनुमान वापरून $\Gamma \Proves
\lforall[x][\varphi(x)]$ मिळते.

\end{proof}

\begin{prop}
\label{fol:axd:qpr:prop:provability-quantifiers}
\begin{enumerate}
\item $\varphi(t) \Proves \lexists[x][\varphi(x)]$.

\item $\lforall[x][\varphi(x)] \Proves \varphi(t)$.
\end{enumerate}
\end{prop}

\begin{proof}
\begin{enumerate}
\item \ref{fol:axd:qua:ax:q2} आणि निगमन प्रमेय वापरून.
\item \ref{fol:axd:qua:ax:q1} आणि निगमन प्रमेय वापरून.
\end{enumerate}
\end{proof}












\section{निर्दोषता}\label{fol:axd:sou:sec}

\begin{explain}
स्वयंसिद्धकीय निगमनासारखी निष्पत्ती-पद्धत
\emph{निर्दोष} असते, जर ती प्रत्यक्षात निष्पन्न न होणाऱ्या गोष्टी
निष्पन्न करू शकत नसेल. त्यामुळे निर्दोषता हा निष्पत्ती-पद्धतींसाठी
हमी दिलेल्या सुरक्षिततेचा एक प्रकार आहे. कोणता सिद्धता-उपपत्तीय गुणधर्म
विचारात आहे यावर अवलंबून, उदाहरणार्थ, आपल्याला पुढील गोष्टी हव्या असतात:
\begin{enumerate}
\item प्रत्येक निष्पन्न करता येण्याजोगे~$\varphi$ वैध असते;
\item $\varphi$ हे इतर काही सूत्रे~$\Gamma$ यांच्यापासून निष्पन्न करता येण्याजोगे असेल,
  तर ते त्यांचा तार्किक परिणामही असते;
\item सूत्रांचा संच~$\Gamma$ विसंगत असेल, तर तो अपूर्ततायोग्य असतो.
\end{enumerate}
हे निष्पत्ती-पद्धतीचे महत्त्वाचे गुणधर्म आहेत. यांपैकी एखादा गुणधर्म
खरा नसेल, तर निष्पत्ती-पद्धतीत दोष आहे---ती खूप जास्त गोष्टी
निष्पन्न करेल. म्हणून निष्पत्ती-पद्धतीची निर्दोषता सिद्ध करणे अत्यंत
महत्त्वाचे आहे.
\end{explain}

\begin{prop}
  जर $\varphi$ हे स्वयंसिद्धक असेल, तर
  प्रत्येक संरचना~$\Struct{M}$ आणि चल-विन्यास~$s$ साठी $\Sat{M}{\varphi}[s]$.
\end{prop}

\begin{proof}
सर्व स्वयंसिद्धके वैध आहेत हे तपासावे लागते. उदाहरणार्थ,
  \ref{fol:axd:qua:ax:q1} चे प्रकरण पाहू: $t$ हे $\varphi$ मध्ये $x$ साठी
  आदेशनासाठी मुक्त आहे असे समजा आणि
  $\Sat{M}{\lforall[x][\varphi]}[s]$ असे गृहीत धरा. मग पूर्ततेच्या व्याख्येनुसार
  प्रत्येक $\varAssign{s'}{s}{x}$ साठी $\Sat{M}{\varphi}[s']$, आणि विशेषतः
  $s'(x) = \Value{t}{M}[s]$ असताना हे खरे आहे.
  \ref{fol:syn:ext:prop:ext-formulas} नुसार
  $\Sat{M}{\Subst{\varphi}{t}{x}}[s]$. यावरून
  $\Sat{M}{(\lforall[x][\varphi] \lif \Subst{\varphi}{t}{x})}[s]$ हे सिद्ध होते.
\end{proof}

\begin{thm}[निर्दोषता]
\label{fol:axd:sou:thm:soundness}
जर $\Gamma \Proves \varphi$, तर $\Gamma \Entails \varphi$.
\end{thm}

\begin{proof}
$\Gamma$ पासून $\varphi$ च्या निष्पत्तीच्या लांबीवर विगमन करू. अनुमानांनी
समर्थित कोणत्याही पायऱ्या नसतील, तर निष्पत्तीमधील सर्व सूत्रे
ही स्वयंसिद्धकांची दृष्टांते असतात किंवा~$\Gamma$ मध्ये असतात. आधीच्या
विधानानुसार सर्व स्वयंसिद्धके वैध आहेत;
म्हणून $\varphi$ हे स्वयंसिद्धक असेल, तर $\Gamma \Entails \varphi$. जर
$\varphi \in \Gamma$, तर स्पष्टपणे $\Gamma \Entails \varphi$.

जर~$\varphi$ च्या निष्पत्तीमधील शेवटची पायरी विध्यात्मक अनुमानाने समर्थित
असेल, तर त्या निष्पत्तीमध्ये $\psi$ आणि $\psi \lif \varphi$ ही सूत्रे
असतात; आणि त्या सूत्रे वर संपणाऱ्या निष्पत्तीच्या भागांना विगमन
गृहीतक लागू होते (कारण त्यांच्यात अनुमानाने समर्थित किमान एक पायरी कमी
असते). म्हणून विगमन गृहीतकानुसार $\Gamma \Entails \psi$ आणि
$\Gamma \Entails \psi \lif \varphi$. मग
\ref{fol:syn:sem:thm:sem-deduction}
नुसार $\Gamma \Entails \varphi$.

आता शेवटची पायरी \QR ने समर्थित आहे असे समजा. मग त्या पायरीचे
रूप $\chi \lif \lforall[x][\psi(x)]$ असते आणि तिच्याआधीची पायरी
$\chi \lif \psi(c)$ असते; येथे $c$ हे $\Gamma$, $\chi$ किंवा
$\lforall[x][B(x)]$ यांमध्ये येत नाही. विगमन गृहीतकानुसार
$\Gamma \Entails \chi \lif \psi(c)$.
\ref{fol:syn:sem:thm:sem-deduction} नुसार $\Gamma \cup \{\chi\} \Entails
\psi(c)$.


$\Sat{M}{\Gamma \cup
  \{\chi\}}$ अशी एखादी रचना~$\Struct{M}$ विचारात घ्या. आपल्याला
$\Sat{M}{\lforall[x][\psi(x)]}$ हे दाखवायचे आहे.
$\lforall[x][\psi(x)]$ हे वाक्य असल्यामुळे याचा अर्थ प्रत्येक
चल-विन्यास~$s$ साठी $\Sat{M}{\psi(x)}[s]$ दाखवायचे आहे
(\ref{fol:syn:ass:prop:sat-quant}). $\Gamma \cup \{\chi\}$ मध्ये केवळ
वाक्ये असल्यामुळे \ref{fol:syn:sat:defn:satisfaction} नुसार प्रत्येक
$\theta \in \Gamma$ साठी $\Sat{M}{\theta}[s]$. $\Struct{M'}$ ही
$\Struct{M}$ सारखीच रचना असू द्या, फक्त $\Assign{c}{M'} = s(x)$.
$c$ हे~$\Gamma$ किंवा~$\chi$ मध्ये येत नसल्यामुळे
\ref{fol:syn:ext:cor:extensionality-sent} नुसार
$\Sat{M'}{\Gamma \cup \{\chi\}}$. $\Gamma \cup \{\chi\}
\Entails \psi(c)$ असल्यामुळे $\Sat{M'}{\psi(c)}$. $\psi(c)$ हे वाक्य
असल्यामुळे \ref{fol:syn:ass:prop:sentence-sat-true} नुसार
$\Sat{M'}{\psi(c)}[s]$. \ref{fol:syn:ext:prop:ext-formulas} नुसार
$\Sat{M'}{\psi(x)}[s]$ तेव्हा आणि केवळ तेव्हाच, जेव्हा
$\Sat{M'}{\psi(c)}[s]$ (लक्षात घ्या की $\psi(c)$ म्हणजे
$\Subst{\psi(x)}{c}{x}$). म्हणून $\Sat{M'}{\psi(x)}[s]$. $c$ हे~$\psi(x)$
मध्ये येत नसल्यामुळे \ref{fol:syn:ext:prop:extensionality} नुसार
$\Sat{M}{\psi(x)}[s]$. पण $s$ हा स्वेच्छ चल-विन्यास होता, म्हणून
$\Sat{M}{\lforall[x][\psi(x)]}$. त्यामुळे $\Gamma \cup \{\chi\} \Entails
\lforall[x][\psi(x)]$. \ref{fol:syn:sem:thm:sem-deduction} नुसार
$\Gamma \Entails \chi \lif \lforall[x][\psi(x)]$.


$\varphi$ हे \QR ने समर्थित पण $\lexists[x][\psi(x)] \lif \chi$ या रूपाचे असलेले
प्रकरण सरावासाठी सोडले आहे.
\end{proof}


\begin{prob}
  \ref{fol:axd:sou:thm:soundness} ची सिद्धता पूर्ण करा.
\end{prob}


\begin{cor}
\label{fol:axd:sou:cor:weak-soundness}
जर $\Proves \varphi$, तर $\varphi$ हे वैध आहे.
\end{cor}

\begin{cor}
\label{fol:axd:sou:cor:consistency-soundness}
जर $\Gamma$ पूर्ततायोग्य असेल, तर तो सुसंगत आहे.
\end{cor}

\begin{proof}
प्रतिलोम विधान सिद्ध करू. $\Gamma$ सुसंगत नाही असे समजा. मग
$\Gamma \Proves \lfalse$, म्हणजे~$\Gamma$ पासून $\lfalse$ ची
निष्पत्ती असते. \ref{fol:axd:sou:thm:soundness} नुसार $\Gamma$ पूर्ण करणाऱ्या
प्रत्येक
संरचना~$\Struct{M}$
ने~$\lfalse$ पूर्ण केले पाहिजे. प्रत्येक
संरचना~$\Struct{M}$ साठी
  $\Sat/{M}{\lfalse}$ असल्यामुळे कोणतीही
$\Struct{M}$ $\Gamma$ पूर्ण करू शकत नाही;
म्हणजे $\Gamma$ अपूर्ततायोग्य आहे.
\end{proof}














\section{निष्पत्ती सह एकरूपता}\label{fol:axd:ide:sec}

निष्पत्त्यांमध्ये $\eq$ समाविष्ट करण्यासाठी आपण फक्त नवीन
स्वयंसिद्धक-आकार जोडतो. निष्पत्ती आणि $\Proves$ यांच्या व्याख्या
तशाच राहतात; आता आपण नवीन स्वयंसिद्धकेही मान्य करतो.

\begin{defn}[एकरूपता साठीची स्वयंसिद्धके]
\begin{align}
\label{fol:axd:ide:ax:id1} & \eq[t][t], \\
\label{fol:axd:ide:ax:id2} & \eq[t_1][t_2] \lif (\psi(t_1) \lif
  \psi(t_2)),
\end{align}
येथे $t$, $t_1$, $t_2$ ही कोणतीही बंद पदे आहेत.
\end{defn}

\begin{prop}\label{fol:axd:ide:prop:sound}
  \ref{fol:axd:ide:ax:id1} आणि \ref{fol:axd:ide:ax:id2} ही स्वयंसिद्धके वैध आहेत.
\end{prop}

\begin{proof}
  सराव.
\end{proof}

\begin{prob}
\ref{fol:axd:ide:prop:sound} सिद्ध करा.
\end{prob}

\begin{prop}\label{fol:axd:ide:prop:iden1}
 कोणतेही पद $t$ आणि कोणताही संच~$\Gamma$ यांसाठी
 $\Gamma \Proves \eq[t][t]$.
\end{prop}

\begin{prop}\label{fol:axd:ide:prop:iden2}
  जर $\Gamma \Proves \varphi(t_1)$ आणि $\Gamma \Proves
  \eq[t_1][t_2]$, तर $\Gamma \Proves \varphi(t_2)$.
\end{prop}

\begin{proof}
पुढील सूत्र
\[
(\eq[t_1][t_2] \lif (\varphi(t_1) \lif \varphi(t_2)))
\]
हा~\ref{fol:axd:ide:ax:id2} चा दृष्टांत आहे. निष्कर्ष~\MP चे दोन उपयोग केल्यावर
मिळतो.
\end{proof}



\chapter{संपूर्णता प्रमेय}\label{fol:com::chap}









\section{परिचय}\label{fol:com:int:sec}

संपूर्णता प्रमेय हा तर्कशास्त्राविषयीच्या सर्वांत मूलभूत निष्कर्षांपैकी एक
आहे. त्याच्या दोन मांडण्या आहेत; त्यांची सममूल्यता आपण सिद्ध करू. त्याची
पहिली मांडणी चिन्हार्थी परिणाम आणि आपल्या निष्पत्ती-पद्धतीतील संबंधाबद्दल
एक मूलभूत विधान करते: जर वाक्य~$\varphi$ हे काही वाक्ये
$\Gamma$ यांच्यापासून निष्पन्न होत असेल, तर $\Gamma \Proves \varphi$ हे
प्रस्थापित करणारी निष्पत्ती सुद्धा असते. त्यामुळे प्रत्यक्षात निष्पन्न
न होणाऱ्या गोष्टी सिद्ध न करता निष्पत्ती-पद्धत जितकी सामर्थ्यवान असू
शकते तितकी ती असते.

दुसऱ्या मांडणीत ते प्रतिमानाच्या अस्तित्वाविषयीचा निष्कर्ष म्हणून सांगता येते:
वाक्यांचा प्रत्येक सुसंगत संच पूर्ततायोग्य असतो. सुसंगतता ही
सिद्धता-उपपत्तीय संकल्पना आहे: आपली निष्पत्ती-पद्धत विशिष्ट
निष्पत्त्या निर्माण करू शकत नाही, असे ती सांगते. पण~$\Gamma$ पासून
विशिष्ट प्रकारच्या निष्पत्त्यांनाहीत एवढ्यावरूनच
संरचना~$\Struct{M}$
अस्तित्वात असून $\Sat{M}{\Gamma}$ आहे, याची हमी
आहे, असे कोण म्हणेल? संपूर्णता प्रमेय प्रथम सिद्ध होण्यापूर्वी---
खरे तर आज आपल्याकडे असलेल्या निष्पत्ती-पद्धती निर्माण होण्यापूर्वी---
महान जर्मन गणितज्ञ डाव्हीट हिल्बर्ट यांचे मत होते की गणिती सिद्धांतांची
सुसंगतता त्यांच्या विषयातील वस्तूंच्या अस्तित्वाची हमी देते. गोटलोप फ्रेग
यांना लिहिलेल्या पत्रात त्यांनी ते पुढीलप्रमाणे मांडले:
\begin{quote}
  स्वेच्छपणे दिलेली स्वयंसिद्धके आणि त्यांचे सर्व परिणाम परस्परविरोधी नसतील,
  तर ती सत्य असतात आणि स्वयंसिद्धकांनी परिभाषित केलेल्या वस्तू अस्तित्वात
  असतात. माझ्यासाठी सत्य आणि अस्तित्वाचा निकष हाच आहे.
\end{quote}
फ्रेग यांनी याला तीव्र विरोध केला. संपूर्णता प्रमेयाची दुसरी मांडणी दाखवते
की निदान पुढील अर्थाने हिल्बर्ट बरोबर होते: स्वयंसिद्धके सुसंगत असतील, तर
त्या सर्वांना सत्य करणारी \emph{एखादी}
संरचना अस्तित्वात असते.

संपूर्णता प्रमेय---किंवा अधिक नेमके सांगायचे तर त्याची सिद्धता---महत्त्वाची
असण्याची हीच काही एकमेव कारणे नाहीत. त्याचे अनेक महत्त्वाचे परिणाम आहेत;
त्यांपैकी काहींची स्वतंत्र चर्चा करू. उदाहरणार्थ, $\Gamma \Proves \varphi$ हे
दाखवणारी कोणतीही निष्पत्ती सांत असल्यामुळे ती~$\Gamma$ मधील सांतच
वाक्ये वापरू शकते. म्हणून संपूर्णता प्रमेयावरून असे निष्पन्न होते की
$\varphi$ हा~$\Gamma$ चा परिणाम असेल, तर तो~$\Gamma$ च्या एखाद्या सांत उपसंचाचाही
परिणाम असतो. याला \emph{संहतता} म्हणतात. सममूल्य रीतीने, $\Gamma$ चा प्रत्येक
सांत उपसंच सुसंगत असेल, तर $\Gamma$ स्वतः सुसंगत असला पाहिजे.

संहतता प्रमेय हे संपूर्णता प्रमेयापासून निष्पत्त्या मधून जाणाऱ्या वळणाने
मिळत असले, तरी संपूर्णता प्रमेयाची \emph{सिद्धता} वापरून ते थेट प्रस्थापित
करणेही शक्य आहे. कारण त्या सिद्धतेत विशिष्ट गुणधर्म---सुसंगतता---असलेला
वाक्यांचा संच घेऊन त्यातून विशिष्ट गुणधर्म असलेली संरचना
रचली जाते (या प्रकरणात ती त्या संचाची पूर्ति करते). सुसंगत संचांऐवजी
``सांततः पूर्ततायोग्य'' वाक्यांच्या संचांपासून सुरुवात केल्यास जवळजवळ
हीच रचना वापरून संहतता थेट प्रस्थापित करता येते.
या रचनेतून इतर परिणामही मिळतात; उदाहरणार्थ, वाक्यांच्या
  कोणत्याही पूर्ततायोग्य संचाला सांत किंवा गणनीय अनंत असे प्रतिमान असते.
  (या निष्कर्षाला ल्योव्हेनहाइम--स्कोलेम प्रमेय म्हणतात.) सर्वसाधारणपणे,
  वाक्यांच्या संचांपासून संरचना रचण्याची पद्धत तर्कशास्त्रात
  वारंवार आणि कधी कधी तत्त्वज्ञानातही वापरली जाते.












\section{सिद्धतेचा आराखडा}\label{fol:com:out:sec}

संपूर्णता प्रमेयाची सिद्धता काहीशी गुंतागुंतीची आहे आणि ती प्रथम वाचताना
वाट चुकणे सोपे आहे. म्हणून सिद्धतेचा आराखडा पाहू. पहिली पायरी म्हणजे
दृष्टिकोनातील बदल; त्यामुळे सिद्धतेकडे जाणारा मार्ग दिसू लागतो. संपूर्णतेचा
अर्थ “जेव्हा जेव्हा $\Gamma \Entails \varphi$, तेव्हा $\Gamma \Proves \varphi$”
असा घेतल्यास कल्पनाही सुचणे कठीण जाऊ शकते: कारण $\Gamma \Proves \varphi$ हे
दाखवण्यासाठी आपल्याला निष्पत्ती शोधावी लागते, आणि $\Gamma \Entails
\varphi$ हे गृहीतक त्यासाठी कोणत्याही प्रकारे मदत करते असे दिसत नाही. काही
सिद्धता-पद्धतींसाठी निष्पत्ती थेट रचता येते; पण आपण थोडा वेगळा मार्ग
घेऊ. आवश्यक दृष्टिकोन-बदल असा आहे: संपूर्णता पुढीलप्रमाणेही मांडता येते:
“$\Gamma$ सुसंगत असेल, तर तो पूर्ततायोग्य असतो.” कदाचित~$\Gamma$ मधील
माहिती आणि तो सुसंगत असल्याचे गृहीतक वापरून आपण~$\Gamma$ मधील प्रत्येक
वाक्य पूर्ण करणारी
संरचना रचू शकू. शेवटी, आपल्याला
कोणत्या प्रकारची संरचना हवी आहे हे
माहीत आहे:~$\Gamma$ जिचे वर्णन करतो अशीच!{}

$\Gamma$ मध्ये फक्त आण्विक
  वाक्ये असतील, तर त्याच्यासाठी
प्रतिमान रचणे सोपे आहे. सर्व आण्विक वाक्ये ही
  $\Atom{P}{a_1,\dots,a_n}$ या रूपाची आहेत असे समजा; येथे $a_i$ ही
  स्थिरांक आहेत. आपल्याला फक्त

  प्रांत~$\Domain{M}$ आणि~$P$ साठी असे निर्देशन ठरवायचे आहे की
  $\Sat{M}{\Atom{P}{a_1,\ldots,a_n}}$. पण ते फार कठीण नाही:
  $\Domain{M} = \Nat$, $\Assign{\Obj c_i}{M} = i$ असे ठेवा आणि प्रत्येक
  $\Atom{P}{a_1,\ldots,a_n} \in \Gamma$ साठी $\tuple{k_1,
    \dots, k_n}$ हे क्रमित बहुपद $\Assign{P}{M}$ मध्ये घाला; येथे $k_i$
  हा स्थिर चिन्ह~$a_i$ चा निर्देशांक आहे (म्हणजे $a_i \ident \Obj
  c_{k_i}$).

आता $\Gamma$ मध्ये~$\lnot \psi$ हे एखादे सूत्र असून $\psi$ आण्विक आहे
असे समजा. $\Struct{M}$ ची रचना $\lnot \psi$ सत्य
करण्याच्या शक्यतेत अडथळा आणेल, अशी काळजी वाटू शकते. पण येथेच~$\Gamma$ ची
सुसंगतता उपयोगी पडते: $\lnot \psi \in \Gamma$ असेल, तर $\psi \notin \Gamma$;
अन्यथा $\Gamma$ विसंगत असेल. आणि $\psi \notin \Gamma$ असेल, तर आपल्या
$\Struct{M}$ च्या रचनेनुसार
$\Sat/{M}{\psi}$; म्हणून
$\Sat{M}{\lnot \psi}$. आतापर्यंत सर्व ठीक आहे.

$\Gamma$ मध्ये गुंतागुंतीची, अनाण्विक सूत्रे असतील तर काय? उदाहरणार्थ,
त्यात $\varphi \land \psi$ आहे असे समजा. ते सत्य करण्यासाठी $\varphi$ आणि $\psi$ ही
दोन्ही~$\Gamma$ मध्ये असल्याप्रमाणे आपण पुढे गेले पाहिजे. आणि $\varphi \lor
\psi \in \Gamma$ असेल, तर त्यांपैकी किमान एक सत्य करावे लागेल; म्हणजे
त्यांपैकी एक~$\Gamma$ मध्ये असल्याप्रमाणे पुढे जावे लागेल.

यावरून पुढील कल्पना सुचते: आपण~$\Gamma$ मध्ये अतिरिक्त सूत्रे
अशा रीतीने भर घालतो की (a)~मिळणारा संच सुसंगत राहील, (b)~प्रत्येक
संभाव्य आण्विक वाक्य~$\varphi$ साठी मिळणाऱ्या संचात $\varphi$ किंवा
$\lnot \varphi$ यांपैकी एक असेल, आणि (c)~संचात $\varphi \land \psi$ असेल तेव्हा
$\varphi$ व $\psi$ ही दोन्ही असतील, तर $\varphi \lor \psi$ असेल तेव्हा $\varphi$ किंवा
$\psi$ यांपैकी किमान एकही असेल, इत्यादी. हे आपण करत राहतो (कदाचित
अनंतकाळपर्यंत). अशा रीतीने भर घातलेल्या सर्व सूत्रांच्या संचाला
$\Gamma^*$ म्हणू. मग वरील रचनेतून आपल्याला अशी
संरचना~$\Struct{M}$
मिळेल की ती~$\Gamma^*$ मधील सर्व वाक्ये पूर्ण करते, हे विगमनाने सिद्ध
करता येईल. $\Gamma \subseteq \Gamma^*$ असल्यामुळे ती~$\Gamma$ मधीलही
सर्व वाक्ये पूर्ण करेल. प्रत्यक्षात (a) आणि~(b) यांची हमी देणे पुरेसे
ठरते. ज्या वाक्यसंचासाठी (b) खरे असते त्याला \emph{संपूर्ण} म्हणतात.
त्यामुळे सुसंगत संच~$\Gamma$ चा विस्तार करून सुसंगत आणि संपूर्ण संच
$\Gamma^*$ मिळवणे हे आपले उद्दिष्ट असेल.


या योजनेत एक अडचण आहे: $\lexists[x][\varphi(x)] \in \Gamma$ असेल, तर या
प्रक्रियेत एखादे स्थिरांक~$c$ निवडून $\varphi(c)$ ची भर घालता यावी अशी
आपली अपेक्षा आहे. पण ते नेहमी करता येईल हे कसे कळणार? कदाचित आपल्या
भाषेत मोजकीच स्थिरांक असतील आणि त्यांपैकी प्रत्येकासाठी
$\lnot \varphi(c) \in \Gamma$ असेल. आपण $\varphi(c)$ चीही भर घालू शकत नाही,
कारण त्यामुळे संच विसंगत होईल; आणि $\Struct{M}$ ने $\varphi(c)$ सत्य करावे की
$\lnot \varphi(c)$ हे आपल्याला कळणार नाही. शिवाय, $\Gamma$ मध्ये अशा भाषेतील
वाक्येच असू शकतात जिच्यात एकही स्थिर चिन्ह नाही (उदा., संचसिद्धान्ताची
भाषा).

या समस्येचा उपाय म्हणजे सुरुवातीलाच अनंत स्थिरे आणि त्यांना संख्यापकांशी
योग्य प्रकारे जोडणारी वाक्ये यांची भर घालणे. (अर्थात, त्यामुळे विसंगती
निर्माण होऊ शकत नाही, हे आपल्याला तपासावे लागेल.)

आण्विक वाक्यांमध्ये फक्त स्थिरांक असतील, तर आपली मूळ रचना चांगली
चालते. पण भाषेत फलने सुद्धा असू शकतात. अशा वेळी सर्व काही योग्य
रीतीने चालण्यासाठी या फलने ना नेमून द्यायची~$\Nat$ वरील योग्य
फलने शोधणे अवघड ठरू शकते. म्हणून आणखी एक युक्ती वापरू: $\Obj c_i$ चा
अर्थ लावण्यासाठी $i$ वापरण्याऐवजी स्थिरांकाचा संचच प्रांत म्हणून
घेऊ. मग $\Struct M$ प्रत्येक स्थिरांक ला तेच स्थिर नेमून देऊ शकते:
$\Assign{\Obj c_i}{M} = \Obj c_i$. पण पूर्ण मार्गच का घेऊ नये:
$\Domain{M}$ हा भाषेतील सर्व \emph{पदांचा} संच असू द्या!{} असे केल्यास
फलने ना फलनांचे स्पष्ट निर्देशन मिळते (या फलनांची प्रचले पदे असतात
आणि त्यांची मूल्येही पदे असतात): फलन~$\Obj f^n_i$ ला आपण असे
फलन नेमून देतो की $n$ पदे $t_1$, \dots,~$t_n$ आदान म्हणून दिल्यावर ते
$\Obj f^n_i(t_1, \dots, t_n)$ हे पद मूल्य म्हणून देते.

कोड्याचा शेवटचा भाग म्हणजे~$\eq$ चे काय करायचे. विधेय~$\eq$ चा
अर्थ स्थिर आहे: $\Sat{M}{\eq[t][t']}$ तेव्हा आणि केवळ तेव्हाच, जेव्हा
$\Value{t}{M} = \Value{t'}{M}$. आता पद~$t$ चे मूल्य हे $t$ स्वतःच
असेल अशी मांडणी केली, तर $t$ आणि $t'$ हे अगदी एकच पद असल्याखेरीज ही
संरचना $\eq[t][t']$ या रूपातील \emph{एकही} वाक्य सत्य करणार नाही.
आणि ही अर्थातच समस्या आहे; कारण फलने असलेल्या भाषेतील जवळजवळ
प्रत्येक रोचक उपपत्तीमध्ये $t$ व $t'$ ही वेगळी पदे असूनही $\eq[t][t']$
या रूपाची वाक्ये प्रमेये असतील (उदा., अंकगणिताच्या उपपत्तींमध्ये
$\eq[(\Obj 0+ \Obj 0)][\Obj 0]$). ही समस्या सोडवण्यासाठी आपण
$\Struct M$ चा प्रांत बदलतो: $\Domain{M}$ मधील वस्तू म्हणून पदे वापरण्याऐवजी
पदांचे संच वापरतो; आणि प्रत्येक संचात~$\Gamma$ मधील वाक्यांनी समान असणे
आवश्यक ठरवलेली सर्व पदे असतात. उदाहरणार्थ, $\Gamma$ ही अंकगणिताची उपपत्ती
असेल, तर अशा संचांपैकी एकात $\Obj 0$, $(\Obj 0 + \Obj 0)$, $(\Obj
0 \times \Obj 0)$, इत्यादी असतील. हाच संच आपण $\Obj 0$ ला नेमून देऊ;
आणि तो त्यातील सर्व पदांचे मूल्यही आहे, उदा., $(\Obj 0 + \Obj 0)$ चेही,
हे निष्पन्न होईल. म्हणून सुधारित संरचनांमध्ये
$\eq[(\Obj 0+ \Obj 0)][\Obj 0]$ हे वाक्य सत्य असेल.

मग आपण पुढीलप्रमाणे काम करू. प्रथम संपूर्ण सुसंगत संचांच्या गुणधर्मांचा
अभ्यास करू. विशेषतः, संपूर्ण सुसंगत संचात $\varphi \land \psi$ तेव्हा आणि
केवळ तेव्हाच असते, जेव्हा त्यात $\varphi$ व~$\psi$ ही दोन्ही असतात; $\varphi \lor \psi$
तेव्हा आणि केवळ तेव्हाच असते, जेव्हा त्यांपैकी किमान एक असते; इत्यादी,
हे सिद्ध करू (\ref{fol:com:ccs:prop:ccs}). त्यानंतर वाक्यांच्या
  “संतृप्त” संचांची व्याख्या करून त्यांचा अभ्यास करू. प्रत्येक संख्यकीकृत
  वाक्याला तिच्या दृष्टांतांशी जोडणाऱ्या सशर्तांचा समावेश असलेला
  संच म्हणजे संतृप्त संच (\ref{fol:com:hen:defn:henkin-exp}). कोणत्याही सुसंगत
  संच~$\Gamma$ चा विस्तार करून संतृप्त संच~$\Gamma'$ नेहमी मिळवता येतो,
  हे आपण दाखवू (\ref{fol:com:hen:lem:henkin}). एखादा संच सुसंगत, संतृप्त आणि
  संपूर्ण असेल, तर त्यात $\lexists[x][\varphi(x)]$ तेव्हा
    आणि केवळ तेव्हाच असते, जेव्हा एखाद्या बंद पद~$t$ साठी त्यात $\varphi(t)$
    असते आणि
  $\lforall[x][\varphi(x)]$ तेव्हा आणि केवळ तेव्हाच असते,
    जेव्हा प्रत्येक बंद पद~$t$ साठी त्यात $\varphi(t)$ असते
  (\ref{fol:com:hen:prop:saturated-instances}), हाही गुणधर्म त्याला असतो.
त्यानंतर संतृप्त सुसंगत संच
$\Gamma'$ घेऊन त्याचा विस्तार करून
संतृप्त, सुसंगत आणि संपूर्ण संच~$\Gamma^*$
मिळवता येतो, हे दाखवू (\ref{fol:com:lin:lem:lindenbaum}). याच~$\Gamma^*$ चा
वापर आपण आपले पद-प्रतिमान~$\Struct M(\Gamma^*)$ परिभाषित करण्यासाठी करू.
पद-प्रतिमानाचा प्रांत हा बंद पदांचा संच असतो आणि त्यातील
  विधेयांचा अर्थ आण्विक वाक्ये यांच्या~$\Gamma^*$ मधील सदस्यत्वाने ठरते
(\ref{fol:com:mod:defn:termmodel}). संतृप्त, संपूर्ण सुसंगत
संचांचे गुणधर्म वापरून
$\Sat{M(\Gamma^*)}{\varphi}$ तेव्हा आणि
केवळ तेव्हाच, जेव्हा $\varphi \in \Gamma^*$, हे दाखवू
(\ref{fol:com:mod:lem:truth}); म्हणून विशेषतः
$\Sat{M(\Gamma^*)}{\Gamma}$.
शेवटी, $\Gamma$ मध्ये~$\eq$ सुद्धा असेल तर पद-प्रतिमानाची
  व्याख्या कशी करायची याचा विचार करू
  (\ref{fol:com:ide:defn:term-model-factor}) आणि ते~$\Gamma^*$ ची पूर्ति करते,
  हे दाखवू (\ref{fol:com:ide:lem:truth}).
















\section{वाक्याचे संपूर्ण सुसंगत संच}\label{fol:com:ccs:sec}

\begin{defn}[संपूर्ण संच]
\label{fol:com:ccs:def:complete-set} वाक्यांचा संच~$\Gamma$
\emph{संपूर्ण} असतो तेव्हा आणि केवळ तेव्हाच, जेव्हा कोणत्याही
वाक्य~$\varphi$ साठी $\varphi \in \Gamma$ किंवा $\lnot \varphi \in \Gamma$.
\end{defn}

\begin{explain}
संपूर्ण वाक्यसंच कोणताही प्रश्न अनुत्तरित ठेवत नाहीत. कोणत्याही
वाक्य~$\varphi$ साठी $\varphi$ सत्य आहे की असत्य हे~$\Gamma$ “सांगतो.”
संपूर्ण संचांचे महत्त्व संपूर्णता प्रमेयाच्या सिद्धतेपलीकडेही आहे.
उदाहरणार्थ, जी उपपत्ती संपूर्ण असून स्वयंसिद्धकीय रीतीने मांडता
येते, तिच्याविषयी नेहमी निर्णय घेता येतो.
\end{explain}

\begin{explain}
संपूर्णतेच्या सिद्धतेत संपूर्ण सुसंगत संच महत्त्वाचे असतात, कारण
वाक्यांचा प्रत्येक सुसंगत संच~$\Gamma$ हा एखाद्या संपूर्ण
सुसंगत संच~$\Gamma^*$ मध्ये समाविष्ट असतो, याची हमी आपण देऊ शकतो.
संपूर्ण सुसंगत संचात प्रत्येक वाक्य~$\varphi$ साठी $\varphi$ किंवा
त्याचे नकरण $\lnot \varphi$ यांपैकी एक असते, पण दोन्ही नसतात. हे विशेषतः
आण्विक वाक्ये यांच्यासाठी खरे असते.
म्हणून स्थिरांकाची योग्य भर घालून विस्तारित केलेल्या
भाषेतील संपूर्ण सुसंगत संचापासून आपण अशी
संरचना रचू शकतो की तिच्यातील
विधेयांचा अर्थ
हा~$\Gamma^*$ मध्ये कोणती आण्विक
वाक्ये आहेत त्यानुसार परिभाषित
केला जातो. त्यानंतर ही संरचना
$\Gamma^*$ मधील सर्व वाक्ये (आणि म्हणून~$\Gamma$ मधीलही सर्व)
सत्य करते, हे दाखवता येते. या शेवटच्या तथ्याच्या सिद्धतेसाठी
$\lnot \varphi \in \Gamma^*$ तेव्हा आणि केवळ तेव्हाच, जेव्हा $\varphi \notin
\Gamma^*$; $(\varphi \lor \psi) \in \Gamma^*$ तेव्हा आणि केवळ तेव्हाच, जेव्हा
$\varphi \in \Gamma^*$ किंवा $\psi \in \Gamma^*$; इत्यादी तथ्ये आवश्यक असतात.
\end{explain}

पुढे आपण $\Proves$ ची परावर्तिता, एकस्वनिकता आणि संक्रमकता हे गुणधर्म
अनेकदा अध्याहृतपणे वापरू (पाहा

  \ref{fol:seq:ptn:sec}, \ref{fol:ntd:ptn:sec}, \ref{fol:axd:ptn:sec}, \ref{fol:tab:ptn:sec}).

\begin{prop}
\label{fol:com:ccs:prop:ccs}
$\Gamma$ हा संपूर्ण आणि सुसंगत आहे असे समजा. मग:
\begin{enumerate}
\item \label{fol:com:ccs:prop:ccs-prov-in} जर $\Gamma \Proves \varphi$, तर $\varphi \in
  \Gamma$.

\item \label{fol:com:ccs:prop:ccs-and} $\varphi \land \psi \in \Gamma$
  तेव्हा आणि केवळ तेव्हाच, जेव्हा $\varphi \in \Gamma$ आणि $\psi \in \Gamma$
  ही दोन्ही खरे आहेत.

\item \label{fol:com:ccs:prop:ccs-or} $\varphi \lor \psi \in \Gamma$ तेव्हा
  आणि केवळ तेव्हाच, जेव्हा $\varphi \in \Gamma$ किंवा $\psi \in \Gamma$.

\item \label{fol:com:ccs:prop:ccs-if} $\varphi \lif \psi \in \Gamma$ तेव्हा
  आणि केवळ तेव्हाच, जेव्हा $\varphi \notin \Gamma$ किंवा $\psi \in \Gamma$.
\end{enumerate}
\end{prop}

\begin{proof}
पुढील सर्व प्रकरणांत $\Gamma$ हा संपूर्ण आणि सुसंगत आहे असे समजू.
\begin{enumerate}
\item जर $\Gamma \Proves \varphi$, तर $\varphi \in \Gamma$.

$\Gamma \Proves \varphi$ असे समजा. उलटपक्षी $\varphi \notin \Gamma$ असे गृहीत
धरा. $\Gamma$ हा संपूर्ण असल्यामुळे $\lnot \varphi \in \Gamma$. पुढील
संदर्भांनुसार

  \ref{fol:seq:prv:prop:explicit-inc}, \ref{fol:ntd:prv:prop:explicit-inc}, \ref{fol:axd:prv:prop:explicit-inc}, \ref{fol:tab:prv:prop:explicit-inc},
$\Gamma$ विसंगत आहे. हे $\Gamma$ सुसंगत असल्याच्या गृहीतकाशी विसंगत आहे.
म्हणून $\varphi \notin \Gamma$ असणे शक्य नाही; त्यामुळे $\varphi \in \Gamma$.

\item

$\varphi \land \psi \in \Gamma$ तेव्हा आणि केवळ तेव्हाच, जेव्हा $\varphi \in
\Gamma$ आणि $\psi \in \Gamma$ ही दोन्ही खरे आहेत:

पुढे जाणाऱ्या दिशेसाठी $\varphi \land \psi \in \Gamma$ असे समजा. मग
पुढील संदर्भांतील

  \ref{fol:seq:ppr:prop:provability-land}, \ref{fol:ntd:ppr:prop:provability-land}, \ref{fol:axd:ppr:prop:provability-land}, \ref{fol:tab:ppr:prop:provability-land}, बाब~(1) नुसार
$\Gamma \Proves \varphi$ आणि $\Gamma \Proves \psi$. \ref{fol:com:ccs:prop:ccs-prov-in}
नुसार $\varphi \in \Gamma$ आणि $\psi \in \Gamma$, अपेक्षेप्रमाणे.

उलट दिशेसाठी $\varphi \in \Gamma$ आणि $\psi \in \Gamma$ असे समजा. पुढील
संदर्भांतील

  \ref{fol:seq:ppr:prop:provability-land}, \ref{fol:ntd:ppr:prop:provability-land}, \ref{fol:axd:ppr:prop:provability-land}, \ref{fol:tab:ppr:prop:provability-land}, बाब~(2) नुसार
$\Gamma \Proves \varphi \land \psi$. \ref{fol:com:ccs:prop:ccs-prov-in} नुसार $\varphi \land
\psi \in \Gamma$.

\item

प्रथम $\varphi \lor \psi \in \Gamma$ असेल, तर $\varphi \in \Gamma$ किंवा $\psi \in
\Gamma$, हे दाखवू. $\varphi \lor \psi \in \Gamma$, पण $\varphi \notin \Gamma$ आणि
$\psi \notin \Gamma$ असे समजा. $\Gamma$ हा संपूर्ण असल्यामुळे
$\lnot \varphi \in \Gamma$ आणि $\lnot \psi \in \Gamma$. पुढील संदर्भांतील

  \ref{fol:seq:ppr:prop:provability-lor}, \ref{fol:ntd:ppr:prop:provability-lor}, \ref{fol:axd:ppr:prop:provability-lor}, \ref{fol:tab:ppr:prop:provability-lor},
बाब~(1) नुसार $\Gamma$ विसंगत आहे; हा विरोधाभास आहे. म्हणून $\varphi \in
\Gamma$ किंवा $\psi \in \Gamma$.

उलट दिशेसाठी $\varphi \in \Gamma$ किंवा $\psi \in \Gamma$ असे समजा. पुढील
संदर्भांतील

  \ref{fol:seq:ppr:prop:provability-lor}, \ref{fol:ntd:ppr:prop:provability-lor}, \ref{fol:axd:ppr:prop:provability-lor}, \ref{fol:tab:ppr:prop:provability-lor}, बाब~(2) नुसार
$\Gamma \Proves \varphi \lor \psi$. \ref{fol:com:ccs:prop:ccs-prov-in} नुसार $\varphi \lor
\psi \in \Gamma$, अपेक्षेप्रमाणे.

\item

पुढे जाणाऱ्या दिशेसाठी $\varphi \lif \psi \in \Gamma$ असे समजा आणि उलटपक्षी
$\varphi \in \Gamma$ व $\psi \notin \Gamma$ असे गृहीत धरा. या गृहीतकांनुसार
$\varphi \lif \psi \in \Gamma$ आणि $\varphi \in \Gamma$. पुढील संदर्भांतील

  \ref{fol:seq:ppr:prop:provability-lif}, \ref{fol:ntd:ppr:prop:provability-lif}, \ref{fol:axd:ppr:prop:provability-lif}, \ref{fol:tab:ppr:prop:provability-lif}, बाब~(1) नुसार
$\Gamma \Proves \psi$. पण मग \ref{fol:com:ccs:prop:ccs-prov-in} नुसार $\psi \in
\Gamma$; हे $\psi \notin \Gamma$ या गृहीतकाशी विसंगत आहे.

उलट दिशेसाठी प्रथम $\varphi \notin \Gamma$ हे प्रकरण पाहू. $\Gamma$ हा
संपूर्ण असल्यामुळे $\lnot \varphi \in \Gamma$. पुढील संदर्भांतील

  \ref{fol:seq:ppr:prop:provability-lif}, \ref{fol:ntd:ppr:prop:provability-lif}, \ref{fol:axd:ppr:prop:provability-lif}, \ref{fol:tab:ppr:prop:provability-lif}, बाब~(2) नुसार
$\Gamma \Proves \varphi \lif \psi$. पुन्हा \ref{fol:com:ccs:prop:ccs-prov-in} नुसार
$\varphi \lif \psi \in \Gamma$, अपेक्षेप्रमाणे.

आता $\psi \in \Gamma$ हे प्रकरण पाहू. पुढील संदर्भांतील

  \ref{fol:seq:ppr:prop:provability-lif}, \ref{fol:ntd:ppr:prop:provability-lif}, \ref{fol:tab:ppr:prop:provability-lif}, \ref{fol:axd:ppr:prop:provability-lif},
पुन्हा बाब~(2) नुसार $\Gamma \Proves \varphi \lif \psi$.
\ref{fol:com:ccs:prop:ccs-prov-in} नुसार $\varphi \lif \psi \in \Gamma$.
\end{enumerate}
\end{proof}


\begin{prob}
\ref{fol:com:ccs:prop:ccs} ची सिद्धता पूर्ण करा.
\end{prob}















\section{हेंकिन विस्तार}\label{fol:com:hen:sec}

\begin{explain}
संपूर्णता प्रमेय सिद्ध करण्यातील एक आव्हान असे आहे की संपूर्ण सुसंगत
संच~$\Gamma$ पासून आपण रचलेल्या प्रतिमानाने~$\Gamma$ मधील सर्व संख्यकीकृत
सूत्रे सत्य केली पाहिजेत. याची हमी देण्यासाठी लिऑन हेंकिन यांची एक
युक्ती वापरतो. थोडक्यात, भाषेत अनंत स्थिरांकाची भर घालून तिचा विस्तार
करणे आणि एक मुक्त चल असलेल्या प्रत्येक सूत्र~$\varphi(x)$ साठी
पुढील रूपाचे सूत्र जोडणे, ही ती युक्ती:
$\lexists[x][\varphi(x)] \lif \varphi(c)$,
येथे $c$ हे नवीन स्थिरांक पैकी एक आहे. $\Gamma$ पूर्ण करणारी
संरचना आपण रचू तेव्हा, यामुळे प्रत्येक
सत्य अस्तित्ववाची वाक्याचा साक्षी
नवीन स्थिरांमध्ये असेल, याची हमी मिळेल.
\end{explain}

\begin{prop}
\label{fol:com:hen:prop:lang-exp}
$\Gamma$ हा $\Lang L$ मध्ये सुसंगत असेल आणि $\Lang L$ मध्ये नवीन
स्थिरांकाचा गणनीय अनंत संच $\Obj d_0$, $\Obj d_1$, \dots
मिळवून $\Lang L'$ तयार केली असेल, तर $\Gamma$ हा~$\Lang L'$ मध्ये
सुसंगत असतो.
\end{prop}

\begin{defn}[संतृप्त संच]
भाषा $\Lang {L}$ मधील सूत्रांचा संच $\Gamma$
\emph{संतृप्त} असतो तेव्हा आणि केवळ तेव्हाच, जेव्हा एक मुक्त
चल~$x$ असलेल्या प्रत्येक सूत्र~$\varphi(x) \in \Frm[L]$
साठी असे स्थिरांक~$c \in \Lang{L}$ असते की
$\lexists[x][\varphi(x)] \lif \varphi(c) \in \Gamma$.
\end{defn}

पुढील प्रमेयाच्या सिद्धतेत खालील व्याख्या वापरली जाईल.

\begin{defn}
\label{fol:com:hen:defn:henkin-exp}
$\Lang L'$ ही \ref{fol:com:hen:prop:lang-exp} मधीलप्रमाणे असू द्या. ज्या
$\Lang L'$ मधील सूत्रे~$\varphi_i(x_i)$ मध्ये एक चर ($x_i$) मुक्तपणे
येतो, त्या सर्वांचे $\varphi_0(x_0)$, $\varphi_1(x_1)$, \dots असे प्रगणन निश्चित
करा. आपण~$n$ वर विगमनाने वाक्ये~$\theta_n$ परिभाषित करतो.

$\Lang{L}$ मध्ये आपण भर घातलेल्या $\Obj d_i$ पैकी~$\varphi_0(x_0)$ मध्ये
न येणारे पहिले स्थिरांक म्हणजे $c_0$ असू द्या. $\theta_0$,
\dots,~$\theta_{n-1}$ आधीच परिभाषित केली आहेत असे गृहीत धरून, नवीन
स्थिरांक~$\Obj d_i$ पैकी~$\theta_0$, \dots,~$\theta_{n-1}$ यांमध्ये किंवा
$\varphi_n(x_n)$ मध्ये न येणारे पहिले स्थिर म्हणजे $c_n$ असू द्या.

आता $\theta_{n}$ हे पुढील सूत्र असू द्या:
$\lexists[x_{n}][\varphi_{n}(x_{n})] \lif
  \varphi_{n}(c_{n})$.
\end{defn}

\begin{lem}
\label{fol:com:hen:lem:henkin}
प्रत्येक सुसंगत संच~$\Gamma$ चा विस्तार करून संतृप्त सुसंगत
संच~$\Gamma'$ मिळवता येतो.
\end{lem}

\begin{proof}
भाषा~$\Lang{L}$ मधील वाक्यांचा सुसंगत संच~$\Gamma$ दिला असता,
नवीन स्थिरांकाचा गणनीय अनंत संच जोडून भाषेचा विस्तार
$\Lang{L'}$ पर्यंत करा. \ref{fol:com:hen:prop:lang-exp} नुसार $\Gamma$ या अधिक
समृद्ध भाषेतही सुसंगत असतो. तसेच $\theta_i$ ही
\ref{fol:com:hen:defn:henkin-exp} मधीलप्रमाणे असू द्या. पुढील व्याख्या करा:
\begin{align*}
\Gamma_0 & = \Gamma \\
\Gamma_{n+1} & = \Gamma_n \cup \{\theta_n \}
\end{align*}
म्हणजे $\Gamma_{n+1} = \Gamma \cup \{ \theta_0, \dots, \theta_n \}$; आणि
$\Gamma' = \bigcup_{n} \Gamma_n$ असू द्या. $\Gamma'$ स्पष्टपणे संतृप्त आहे.

$\Gamma'$ विसंगत असेल, तर कोणत्यातरी $n$ साठी $\Gamma_n$ विसंगत असेल
(सराव: का ते स्पष्ट करा). म्हणून $\Gamma'$ सुसंगत आहे हे दाखवण्यासाठी
प्रत्येक संच~$\Gamma_n$ सुसंगत आहे हे~$n$ वर विगमनाने दाखवणे पुरेसे आहे.

विगमनाचा आधार म्हणजे $\Gamma_0 = \Gamma$ सुसंगत आहे हे विधानच; ते
प्रमेयाचे गृहीतक आहे. विगमन-पायरीसाठी $\Gamma_{n}$ सुसंगत आहे, पण
$\Gamma_{n+1} = \Gamma_n \cup \{\theta_n\}$ विसंगत आहे असे समजा. लक्षात
घ्या की $\theta_n$ हे
$\lexists[x_{n}][\varphi_{n}(x_n)] \lif \varphi_{n}(c_{n})$
आहे; येथे $\varphi_n(x_n)$ हे फक्त~$x_n$ चर मुक्त असलेले $\Lang{L'}$ चे
सूत्र आहे. $c_n$ निवडण्याच्या पद्धतीनुसार
(\ref{fol:com:hen:defn:henkin-exp} पाहा), $c_n$ हे~$\varphi_n(x_n)$ किंवा~$\Gamma_n$
मध्ये येत नाही.

$\Gamma_n \cup \{\theta_n\}$ विसंगत असेल, तर $\Gamma_n
\Proves \lnot \theta_n$; म्हणून पुढील दोन्ही विधाने खरी आहेत:

\[
\Gamma_n \Proves \lexists[x_n][\varphi_n(x_n)]
\qquad
\Gamma_n \Proves \lnot \varphi_n(c_n)
\]
$c_n$ हे $\Gamma_n$ किंवा~$\varphi_n(x_n)$ मध्ये येत नसल्यामुळे
\ref{fol:axd:qpr:thm:strong-generalization}, \ref{fol:seq:qpr:thm:strong-generalization}, \ref{fol:ntd:qpr:thm:strong-generalization}, \ref{fol:tab:qpr:thm:strong-generalization}
लागू होते.
$\Gamma_n \Proves \lnot \varphi_n(c_n)$
यापासून अनुक्रमे
$\Gamma_n \Proves \lforall[x_n][\lnot \varphi_n(x_n)]$
मिळते. अशा रीतीने आपल्याकडे पुढील दोन्ही विधाने आहेत:
$\Gamma_n \Proves \lexists[x_n][\varphi_n(x_n)]$ आणि
$\Gamma_n \Proves \lforall[x_n][\lnot \varphi_n(x_n)]$;
म्हणून $\Gamma_n$ स्वतः विसंगत आहे.
(लक्षात घ्या की
$\lforall[x_n][\lnot \varphi_n(x_n)] \Proves
  \lnot\lexists[x_n][\varphi_n(x_n)]$.)
हा विरोधाभास आहे: $\Gamma_n$ सुसंगत असल्याचे गृहीत होते. म्हणून
$\Gamma_n \cup \{ \theta_n\}$ सुसंगत आहे.

\end{proof}

\begin{explain}
आता आपण दाखवू की संतृप्त असलेल्या \emph{संपूर्ण}, सुसंगत संचांना पुढील
गुणधर्म असतो: त्यात सर्ववाची संख्यकीकृत वाक्य
  तेव्हा आणि केवळ तेव्हाच असते, जेव्हा त्याचे सर्व दृष्टांत त्यात असतात
   आणि त्यात अस्तित्ववाची
  संख्यकीकृत वाक्य तेव्हा आणि केवळ तेव्हाच असते, जेव्हा त्याचा
  किमान एक दृष्टांत त्यात असतो. संपूर्ण, सुसंगत, संतृप्त संचापासून
आपण निर्माण करू ती संरचना त्यातील सर्व संख्यकीकृत वाक्ये सत्य करते,
हे दाखवण्यासाठी हा गुणधर्म वापरू.
\end{explain}

\begin{prop}\label{fol:com:hen:prop:saturated-instances}
$\Gamma$ हा संपूर्ण, सुसंगत आणि संतृप्त आहे असे समजा.
\begin{enumerate}
\item $\lexists[x][\varphi(x)] \in \Gamma$ तेव्हा आणि केवळ तेव्हाच,
  जेव्हा किमान एका बंद पद~$t$ साठी $\varphi(t) \in \Gamma$.
\item $\lforall[x][\varphi(x)] \in \Gamma$ तेव्हा आणि केवळ तेव्हाच,
  जेव्हा सर्व बंद पदे~$t$ यांच्यासाठी $\varphi(t) \in \Gamma$.
\end{enumerate}
\end{prop}

\begin{proof}
  \begin{enumerate}
  \item
    प्रथम $\lexists[x][\varphi(x)]
      \in \Gamma$ असे समजा. $\Gamma$ संतृप्त असल्यामुळे कोणत्यातरी
      स्थिरांक~$c$ साठी
      $(\lexists[x][\varphi(x)] \lif \varphi(c)) \in \Gamma$. पुढील संदर्भांतील
      \ref{fol:axd:ppr:prop:provability-lif}, \ref{fol:seq:ppr:prop:provability-lif}, \ref{fol:ntd:ppr:prop:provability-lif}, \ref{fol:tab:ppr:prop:provability-lif},
      बाब~(1), आणि
      \ref{fol:com:ccs:prop:ccs}\ref{fol:com:ccs:prop:ccs-prov-in} नुसार $\varphi(c)
      \in \Gamma$.

    उलट दिशेसाठी संतृप्तता आवश्यक नाही: $\varphi(t) \in \Gamma$ असे समजा.
    मग पुढील संदर्भांतील
      \ref{fol:axd:qpr:prop:provability-quantifiers}, \ref{fol:seq:qpr:prop:provability-quantifiers}, \ref{fol:ntd:qpr:prop:provability-quantifiers}, \ref{fol:tab:qpr:prop:provability-quantifiers},
      बाब~(1) नुसार $\Gamma \Proves \lexists[x][\varphi(x)]$.
    \ref{fol:com:ccs:prop:ccs}\ref{fol:com:ccs:prop:ccs-prov-in} नुसार
    $\lexists[x][\varphi(x)] \in \Gamma$.
  \item
    सर्व बंद पदे~$t$ यांच्यासाठी
      $\varphi(t) \in \Gamma$ असे समजा. विरोधासाठी
      $\lforall[x][\varphi(x)] \notin \Gamma$ असे गृहीत धरा. $\Gamma$ संपूर्ण
      असल्यामुळे $\lnot\lforall[x][\varphi(x)] \in \Gamma$. संतृप्ततेनुसार
      कोणत्यातरी स्थिरांक~$c$ साठी
      $(\lexists[x][\lnot \varphi(x)] \lif \lnot \varphi(c)) \in
        \Gamma$. $c$ हे बंद पद असल्यामुळे गृहीतकानुसार $\varphi(c) \in
      \Gamma$. पण यामुळे $\Gamma$ विसंगत होईल. (सराव: पुढील संच विसंगत
      असल्याचे दाखवणारी निष्पत्ती द्या:
      \[
      \lnot \lforall[x][\varphi(x)],
      \lexists[x][\lnot \varphi(x)] \lif \lnot
        \varphi(c), \varphi(c)
      \]
      .)

      उलट दिशेसाठी संतृप्तता आवश्यक नाही:
      $\lforall[x][\varphi(x)] \in \Gamma$ असे समजा. मग पुढील संदर्भांतील
      \ref{fol:axd:qpr:prop:provability-quantifiers}, \ref{fol:seq:qpr:prop:provability-quantifiers}, \ref{fol:ntd:qpr:prop:provability-quantifiers}, \ref{fol:tab:qpr:prop:provability-quantifiers},
      बाब~(2) नुसार $\Gamma \Proves \varphi(t)$. \ref{fol:com:ccs:prop:ccs}
      नुसार $\varphi(t) \in \Gamma$.
  \end{enumerate}
\end{proof}












\section{लिंडेनबाउमचे पूर्वप्रमेय}\label{fol:com:lin:sec}

\begin{explain}
आता आपण असे पूर्वप्रमेय सिद्ध करू की वाक्यांचा कोणताही सुसंगत
संच केवळ सुसंगतच नव्हे, तर संपूर्ण असलेल्या एखाद्या वाक्यसंचात
समाविष्ट असतो. या सिद्धतेत प्रत्येक वेळी एक वाक्य जोडले जाते
आणि प्रत्येक पायरीवर संच सुसंगत राहील याची हमी दिली जाते. हे अशा रीतीने
करतो की प्रत्येक $\varphi$ साठी कोणत्यातरी टप्प्यावर $\varphi$ किंवा $\lnot
\varphi$ यांपैकी एक जोडले जाईल. मग त्या रचनेतील सर्व टप्प्यांचा संघ $\varphi$
किंवा त्याचे नकरण~$\lnot \varphi$ यांपैकी एक समाविष्ट करतो आणि म्हणून तो
संपूर्ण असतो. प्रत्येक टप्प्यावर विसंगती येणार नाही याची खात्री केल्यामुळे
तो सुसंगतही असतो.
\end{explain}

\begin{lem}[लिंडेनबाउमचे पूर्वप्रमेय]
\label{fol:com:lin:lem:lindenbaum} भाषा~$\Lang{L}$ मधील प्रत्येक सुसंगत
संच~$\Gamma$ चा विस्तार करून संपूर्ण आणि सुसंगत असा
संच~$\Gamma^*$ मिळवता येतो.
\end{lem}

\begin{proof}
$\Gamma$ सुसंगत असू द्या. $\Lang L$ मधील सर्व वाक्यांचे
$\varphi_0$, $\varphi_1$, \dots{} असे प्रगणन असू द्या. $\Gamma_0 = \Gamma$
अशी व्याख्या करा, आणि
\[
\Gamma_{n+1} =
\begin{cases}
\Gamma_n \cup \{ \varphi_n \} & \textrm{जर $\Gamma_n \cup \{\varphi_n\}$
  सुसंगत असेल;} \\
\Gamma_n \cup \{ \lnot \varphi_n \} & \textrm{अन्यथा.}
\end{cases}
\]
$\Gamma^* = \bigcup_{n \geq 0} \Gamma_n$ असू द्या.

प्रत्येक $\Gamma_n$ सुसंगत आहे: व्याख्येनुसार $\Gamma_0$ सुसंगत आहे.
$\Gamma_{n+1} = \Gamma_n \cup \{\varphi_n\}$ असेल, तर याचे कारण हा
दुसरा संच सुसंगत असणे हेच आहे. तो सुसंगत नसेल, तर $\Gamma_{n+1} =
\Gamma_n \cup \{\lnot \varphi_n\}$. $\Gamma_n \cup \{\lnot \varphi_n\}$
सुसंगत आहे हे आपल्याला पडताळावे लागेल. तो सुसंगत नाही असे समजा. मग
$\Gamma_n \cup \{\varphi_n\}$ आणि $\Gamma_n \cup \{\lnot \varphi_n\}$ हे
\emph{दोन्ही} विसंगत आहेत. याचा अर्थ

  \ref{fol:axd:prv:prop:provability-exhaustive}, \ref{fol:seq:prv:prop:provability-exhaustive}, \ref{fol:ntd:prv:prop:provability-exhaustive}, \ref{fol:tab:prv:prop:provability-exhaustive}
नुसार $\Gamma_n$ विसंगत असेल; पण हे विगमन-गृहीतकाच्या विरुद्ध आहे.

प्रत्येक~$n$ आणि प्रत्येक $i < n$ साठी $\Gamma_i \subseteq \Gamma_n$.
हे~$n$ वरील सोप्या विगमनाने मिळते. $n=0$ असेल, तर $i < 0$ असा
कोणताही निर्देशांक नसतो; म्हणून विधान आपोआप लागू होते. विगमन-पायरीसाठी ते~$n$
साठी खरे आहे असे समजा. $i < n+1$ असेल, तर $\Gamma_i \subseteq
\Gamma_{n+1}$ हे दाखवू. रचनेनुसार $\Gamma_{n+1} = \Gamma_n \cup
\{\varphi_n\}$ किंवा $= \Gamma_n \cup \{\lnot \varphi_n\}$. म्हणून $\Gamma_n
\subseteq \Gamma_{n+1}$. $i < n+1$ असेल, तर विगमन-गृहीतकानुसार
$\Gamma_i \subseteq \Gamma_n$ ($i < n$ असल्यास), किंवा $\Gamma_n
\subseteq \Gamma_n$ हे क्षुल्लक तथ्य लागू होते ($i = n$ असल्यास).
म्हणून~$\subseteq$ च्या संक्रमणशीलतेने $\Gamma_i \subseteq
\Gamma_{n+1}$ मिळते.

यावरून $\Gamma^*$ सुसंगत आहे हे मिळते. ते कसे ते पाहू: $\Gamma'
\subseteq \Gamma^*$ हा सांत संच असू द्या. प्रत्येक $\psi \in \Gamma'$
हे कोणत्यातरी~$i$ साठी~$\Gamma_i$ मध्येही असते. या निर्देशांकांतील
सर्वांत मोठा निर्देशांक $n$ असू द्या. $i \le n$ असेल, तर $\Gamma_i
\subseteq \Gamma_n$ असल्यामुळे प्रत्येक $\psi \in \Gamma'$ हे
$\in \Gamma_n$ देखील आहे, म्हणजे $\Gamma' \subseteq \Gamma_n$, आणि
$\Gamma_n$ सुसंगत आहे. म्हणून प्रत्येक सांत उपसंच $\Gamma' \subseteq
\Gamma^*$ सुसंगत आहे.

  \ref{fol:axd:ptn:prop:proves-compact}, \ref{fol:seq:ptn:prop:proves-compact}, \ref{fol:ntd:ptn:prop:proves-compact}, \ref{fol:tab:ptn:prop:proves-compact} नुसार $\Gamma^*$
सुसंगत आहे.

$\Frm[L]$ मधील प्रत्येक वाक्य हा~$\Gamma^*$ ची व्याख्या
करण्यासाठी वापरलेल्या यादीत येतो. $\varphi_n \notin \Gamma^*$ असेल, तर
$\Gamma_n \cup \{\varphi_n\}$ विसंगत होते हेच त्याचे कारण आहे. पण मग
$\lnot \varphi_n \in \Gamma^*$; म्हणून $\Gamma^*$ हा संपूर्ण आहे.
\end{proof}












\section{प्रतिमानाची रचना}\label{fol:com:mod:sec}

\begin{explain}
सध्या आपण~$\eq$ चा विचार करत नाही; म्हणजे~$\eq$ नसलेला
  वाक्यांचा सुसंगत संच~$\Gamma$ पूर्ततायोग्य आहे एवढेच आपल्याला
  दाखवायचे आहे. प्रथम~$\Gamma$ चा विस्तार करून सुसंगत, संपूर्ण
  आणि संतृप्त असा संच~$\Gamma^*$ मिळवतो. या बाबतीत
  प्रतिमान~$\Struct{M(\Gamma^*)}$ ची व्याख्या सोपी आहे: आपण~$\Lang{L'}$
  मधील बंद पदांचा संच प्रांत म्हणून घेतो. प्रत्येक स्थिरांक ला तेच
  स्थिर नेमतो आणि, अधिक व्यापकपणे, प्रत्येक बंद पद~$t$ साठी
  $\Value{t}{M(\Gamma^*)} = t$ याची खात्री करतो. विधेये ना
  असे विस्तार नेमतो की एखादे आण्विक वाक्य हे
  $\Struct{M(\Gamma^*)}$ मध्ये तेव्हा आणि केवळ तेव्हाच सत्य असते, जेव्हा
  ते~$\Gamma^*$ मध्ये असते. त्यामुळे~$\Gamma^*$ मधील सर्व आण्विक
  वाक्ये हे~$\Struct{M(\Gamma^*)}$ मध्ये सत्य होतील, हे उघड आहे.
  आपण सुरुवात केलेला~$\Gamma^*$ सुसंगत, संपूर्ण आणि संतृप्त असेल, तर
  उरलेली वाक्येही सत्य असतात.
\end{explain}


\begin{defn}[पद-प्रतिमान]
\label{fol:com:mod:defn:termmodel}
$\Gamma^*$ हा भाषा~$\Lang L$ मधील वाक्यांचा संपूर्ण,
सुसंगत आणि संतृप्त संच असू द्या. $\Gamma^*$ चे
\emph{पद-प्रतिमान}~$\Struct M(\Gamma^*)$ ही पुढीलप्रमाणे परिभाषित केलेली
संरचना आहे:
\begin{enumerate}
\item प्रांत~$\Domain{M(\Gamma^*)}$ हा~$\Lang L$ मधील सर्व बंद
  पदांचा संच आहे.
\item स्थिरांक~$c$ चा निर्वचनार्थ स्वतः~$c$ आहे:
  $\Assign{c}{M(\Gamma^*)} = c$.
\item फलन~$f$ ला असे फलन नेमले जाते की बंद पदे $t_1$, \dots,
  $t_n$ ही वितर्के दिली असता त्याचे मूल्य बंद पद $f(t_1, \dots, t_n)$
  असते:
\[
\Assign{f}{M(\Gamma^*)}(t_1, \dots, t_n) = f(t_1,\dots, t_n)
\]
\item $R$ हे $n$-स्थानी विधेय असेल, तर
  \[
  \tuple{t_1, \dots,
  t_n} \in \Assign{R}{M(\Gamma^*)} \text{ तेव्हा आणि केवळ तेव्हाच } \Atom{R}{t_1, \dots,
    t_n} \in \Gamma^*.
  \]
\end{enumerate}
\end{defn}



आपल्याकडे खरोखरच $\Value{t}{M(\Gamma^*)} = t$ आहे, हे आता पडताळू.

\begin{lem}
\label{fol:com:mod:lem:val-in-termmodel} \ref{fol:com:mod:defn:termmodel} मधील
$\Struct M(\Gamma^*)$ हे पद-प्रतिमान असू द्या; मग
$\Value{t}{M(\Gamma^*)} = t$.
\end{lem}
\begin{proof}
 ही सिद्धता~$t$ वरील विगमनाने मिळते. $t$ हे स्थिरांक असण्याचा
 आधार-प्रसंग पद-प्रतिमानाच्या व्याख्येवरून थेट मिळतो. विगमन-पायरीसाठी
 $t_1, \ldots, t_n$ ही अशी बंद पदे आहेत की
 $\Value{t_i}{M(\Gamma^*)} = t_i$, आणि $f$ हे $n$-स्थानी
 फलन आहे, असे समजा. मग
\begin{align*}
\Value{f(t_1,\ldots,t_n)}{M(\Gamma^*)} &= \Assign{f}{M(\Gamma^*)}(\Value{t_1}
 {M(\Gamma^*)},
\ldots, \Value{t_n}{M(\Gamma^*)}) \\
&= \Assign{f}{M(\Gamma^*)}(t_1, \dots, t_n) \\
&= f(t_1,\dots, t_n),
\end{align*}
आणि म्हणून विगमनाने हे प्रत्येक बंद पद~$t$ साठी लागू होते.
\end{proof}



\begin{explain}
  एखादी संरचना~$\Struct{M}$ अस्तित्ववाची संख्यकीकृत
    वाक्य~$\lexists[x][\varphi(x)]$ सत्य करू शकते, पण ती सत्य करते
    असा कोणताही दृष्टांत~$\varphi(t)$ नसेल.
  एखादी संरचना~$\Struct{M}$ सर्ववाची संख्यकीकृत
    वाक्य~$\lforall[x][\varphi(x)]$ चे सर्व दृष्टांत~$\varphi(t)$ सत्य करू
    शकते, पण~$\lforall[x][\varphi(x)]$ सत्य करणार नाही. याचे कारण सामान्यतः
  प्रांत~$\Domain{M}$ मधील प्रत्येक घटक हे एखाद्या बंद पदाचे
  मूल्य नसते ($\Struct{M}$ बंद पदांनी आच्छादित नसेल). म्हणूनच
  पूर्ति-संबंधाची व्याख्या चर-मूल्यांकनांच्या साहाय्याने केली जाते. पण आपल्या
  पद-प्रतिमान~$\Struct{M(\Gamma^*)}$ साठी त्याची आवश्यकता नसती---कारण
  ते बंद पदांनी आच्छादित आहे. पुढील निष्पत्तीचा आशय हाच आहे.
\end{explain}

\begin{prop}
\label{fol:com:mod:prop:quant-termmodel}
\ref{fol:com:mod:defn:termmodel} मधील $\Struct M(\Gamma^*)$ हे पद-प्रतिमान असू द्या.
\begin{enumerate}
\item $\Sat{M(\Gamma^*)}{\lexists[x][\varphi(x)]}$ तेव्हा आणि केवळ तेव्हाच,
  जेव्हा किमान एका बंद पद~$t$ साठी $\Sat{M(\Gamma^*)}{\varphi(t)}$.
\item $\Sat{M(\Gamma^*)}{\lforall[x][\varphi(x)]}$ तेव्हा आणि केवळ तेव्हाच,
  जेव्हा सर्व बंद पदे~$t$ यांच्यासाठी $\Sat{M(\Gamma^*)}{\varphi(t)}$.
\end{enumerate}
\end{prop}

\begin{proof}
\begin{enumerate}
\item
  \ref{fol:syn:ass:prop:sat-quant} नुसार
    $\Sat{M(\Gamma^*)}{\lexists[x][\varphi(x)]}$ तेव्हा आणि केवळ तेव्हाच,
    जेव्हा किमान एका चर-मूल्यांकन~$s$ साठी
    $\Sat{M(\Gamma^*)}{\varphi(x)}[s]$. $\Domain{M(\Gamma^*)}$ मध्ये
    $\Lang{L}$ मधील बंद पदे असल्यामुळे, हे तेव्हा आणि केवळ तेव्हाच लागू
    होते, जेव्हा असे किमान एक बंद पद~$t$ असते की $s(x) = t$ आणि
    $\Sat{M(\Gamma^*)}{\varphi(x)}[s]$. \ref{fol:syn:ext:prop:ext-formulas}
    नुसार, $s(x) = t$ असेल तेथे $\Sat{M(\Gamma^*)}{\varphi(x)}[s]$ तेव्हा आणि
    केवळ तेव्हाच, जेव्हा $\Sat{M(\Gamma^*)}{\varphi(t)}[s]$. तसेच
    \ref{fol:syn:ass:prop:sentence-sat-true} नुसार,
    $\varphi(t)$ हे वाक्य असल्यामुळे $\Sat{M(\Gamma^*)}{\varphi(t)}[s]$ तेव्हा आणि
    केवळ तेव्हाच, जेव्हा $\Sat{M(\Gamma^*)}{\varphi(t)}$.
\item
  \ref{fol:syn:ass:prop:sat-quant} नुसार
    $\Sat{M(\Gamma^*)}{\lforall[x][\varphi(x)]}$ तेव्हा आणि केवळ तेव्हाच,
    जेव्हा प्रत्येक चर-मूल्यांकन $s$ साठी
    $\Sat{M(\Gamma^*)}{\varphi(x)}[s]$. लक्षात घ्या की
    $\Domain{M(\Gamma^*)}$ मध्ये~$\Lang{L}$ मधील बंद पदे असतात. त्यामुळे
    प्रत्येक बंद पद~$t$ साठी $s(x) = t$ असे चर-मूल्यांकन मिळते, आणि
    कोणत्याही चर-मूल्यांकनासाठी $s(x)$ हे एखादे बंद पद~$t$ असते.
    \ref{fol:syn:ext:prop:ext-formulas} नुसार, $s(x) = t$ असेल तेथे
    $\Sat{M(\Gamma^*)}{\varphi(x)}[s]$ तेव्हा आणि केवळ तेव्हाच, जेव्हा
    $\Sat{M(\Gamma^*)}{\varphi(t)}[s]$. तसेच
    \ref{fol:syn:ass:prop:sentence-sat-true} नुसार,
    $\varphi(t)$ हे वाक्य असल्यामुळे $\Sat{M(\Gamma^*)}{\varphi(t)}[s]$ तेव्हा आणि
    केवळ तेव्हाच, जेव्हा $\Sat{M(\Gamma^*)}{\varphi(t)}$.
\end{enumerate}
\end{proof}




\begin{lem}[सत्यता पूर्वप्रमेय]
\label{fol:com:mod:lem:truth} $\varphi$ मध्ये~$\eq$ येत नाही असे समजा. मग
$\Sat{M(\Gamma^*)}{\varphi}$ तेव्हा आणि केवळ तेव्हाच, जेव्हा $\varphi \in \Gamma^*$.
\end{lem}

\begin{proof}
दोन्ही दिशा एकाच वेळी, $\varphi$ वरील विगमनाने सिद्ध करू.
\begin{enumerate}
\item \indcase{\varphi}{\lfalse}
  {$\Sat/{M(\Gamma^*)}{\lfalse}$
    हे पूर्तिच्या व्याख्येनुसार मिळते. दुसरीकडे, $\Gamma^*$ सुसंगत
    असल्यामुळे $\lfalse \notin \Gamma^*$.}

\item \indcase{\varphi}{\ltrue}
  {$\Sat{M(\Gamma^*)}{\ltrue}$
    हे पूर्तिच्या व्याख्येनुसार मिळते. दुसरीकडे, $\Gamma^*$ सुसंगत व
    संपूर्ण असल्यामुळे आणि $\Gamma^* \Proves \ltrue$ असल्यामुळे
    $\ltrue \in \Gamma^*$.}

\item \indcase{\varphi}{R(t_1, \dots, t_n)}
  {$\Sat{M(\Gamma^*)}{\Atom{R}{t_1, \dots, t_n}}$ तेव्हा आणि केवळ तेव्हाच,
    जेव्हा $\tuple{t_1, \dots, t_n} \in \Assign{R}{M(\Gamma^*)}$
    (पूर्तिच्या व्याख्येनुसार); आणि हे तेव्हा आणि केवळ तेव्हाच, जेव्हा
    $R(t_1, \dots, t_n) \in \Gamma^*$ ($\Struct M(\Gamma^*)$ च्या
    रचनेनुसार).}

\item
  \indcase{\varphi}{\lnot \psi}
    {$\Sat{M(\Gamma^*)}{\indfrm}$ तेव्हा
    आणि केवळ तेव्हाच, जेव्हा
    $\Sat/{M(\Gamma^*)}{\psi}$
    (पूर्तिच्या व्याख्येनुसार). विगमन-गृहीतकानुसार,
    $\Sat/{M(\Gamma^*)}{\psi}$ तेव्हा
    आणि केवळ तेव्हाच, जेव्हा $\psi \notin \Gamma^*$. $\Gamma^*$ सुसंगत
    आणि संपूर्ण असल्यामुळे $\psi \notin \Gamma^*$ तेव्हा आणि केवळ
    तेव्हाच, जेव्हा $\lnot \psi \in \Gamma^*$.}

\item

  \indcase{\varphi}{\psi \land
    \chi}{$\Sat{M(\Gamma^*)}{\indfrm}$
    तेव्हा आणि केवळ तेव्हाच, जेव्हा
    $\Sat{M(\Gamma^*)}{\psi}$ आणि
    $\Sat{M(\Gamma^*)}{\chi}$ हे दोन्ही
    लागू होतात (पूर्तिच्या व्याख्येनुसार); आणि हे तेव्हा आणि केवळ तेव्हाच,
    जेव्हा $\psi \in \Gamma^*$ आणि $\chi \in \Gamma^*$ हे दोन्ही लागू
    होतात (विगमन-गृहीतकानुसार). तसेच
    \ref{fol:com:ccs:prop:ccs}\ref{fol:com:ccs:prop:ccs-and} नुसार, हे तेव्हा आणि
    केवळ तेव्हाच लागू होते, जेव्हा $(\psi \land \chi) \in \Gamma^*$.}

\item

  \indcase{\varphi}{\psi \lor \chi}
    {$\Sat{M(\Gamma^*)}{\indfrm}$
    तेव्हा आणि केवळ तेव्हाच, जेव्हा
    $\Sat{M(\Gamma^*)}{\psi}$ किंवा
    $\Sat{M(\Gamma^*)}{\chi}$
    (पूर्तिच्या व्याख्येनुसार); आणि हे तेव्हा आणि केवळ तेव्हाच, जेव्हा
    $\psi \in \Gamma^*$ किंवा $\chi \in \Gamma^*$ (विगमन-गृहीतकानुसार).
    \ref{fol:com:ccs:prop:ccs}\ref{fol:com:ccs:prop:ccs-or} नुसार, हे तेव्हा आणि
    केवळ तेव्हाच लागू होते, जेव्हा $(\psi \lor \chi) \in \Gamma^*$.}

\item

  \indcase{\varphi}{\psi \lif \chi}
    {$\Sat{M(\Gamma^*)}{\indfrm}$
    तेव्हा आणि केवळ तेव्हाच, जेव्हा
    $\Sat/{M(\Gamma^*)}{\psi}$ किंवा $\Sat{M (\Gamma^*)}{\chi}$
    (पूर्तिच्या व्याख्येनुसार); आणि हे तेव्हा आणि केवळ तेव्हाच, जेव्हा
    $\psi \notin \Gamma^*$ किंवा $\chi \in \Gamma^*$ (विगमन-गृहीतकानुसार).
    \ref{fol:com:ccs:prop:ccs}\ref{fol:com:ccs:prop:ccs-if} नुसार, हे तेव्हा आणि
    केवळ तेव्हाच लागू होते, जेव्हा $(\psi \lif \chi) \in \Gamma^*$.}


\item

    \indcase{\varphi}{\lforall[x][\psi(x)]}{$\Sat{M(\Gamma^*)}{\indfrm}$ तेव्हा
      आणि केवळ तेव्हाच, जेव्हा सर्व बंद पदे~$t$ यांच्यासाठी
      $\Sat{M(\Gamma^*)}{\psi(t)}$ (\ref{fol:com:mod:prop:quant-termmodel}).
      विगमन-गृहीतकानुसार हे तेव्हा आणि केवळ तेव्हाच लागू होते, जेव्हा सर्व
      बंद पदे~$t$ यांच्यासाठी $\psi(t) \in \Gamma^*$; आणि
      \ref{fol:com:hen:prop:saturated-instances} नुसार हे पुन्हा तेव्हा आणि
      केवळ तेव्हाच लागू होते, जेव्हा $\lforall[x][\psi(x)] \in \Gamma^*$.}


\item

    \indcase{\varphi}{\lexists[x][\psi(x)]}{$\Sat{M(\Gamma^*)}{\indfrm}$ तेव्हा
      आणि केवळ तेव्हाच, जेव्हा किमान एका बंद पद~$t$ साठी
      $\Sat{M(\Gamma^*)}{\psi(t)}$ (\ref{fol:com:mod:prop:quant-termmodel}).
      विगमन-गृहीतकानुसार हे तेव्हा आणि केवळ तेव्हाच लागू होते, जेव्हा किमान
      एका बंद पद~$t$ साठी $\psi(t) \in \Gamma^*$. तसेच
      \ref{fol:com:hen:prop:saturated-instances} नुसार हे पुन्हा तेव्हा आणि
      केवळ तेव्हाच लागू होते, जेव्हा $\lexists[x][\psi(x)] \in \Gamma^*$.}


\end{enumerate}
\end{proof}
















\section{एकरूपता}\label{fol:com:ide:sec}

\begin{explain}
मागील विभागात दिलेली पद-प्रतिमानाची रचना ही~$\eq$ नसलेल्या
संच~$\Gamma$ साठी प्रथम-कोटी तर्कशास्त्राची संपूर्णता प्रस्थापित करण्यास
पुरेशी आहे. पद-प्रतिमान हे~$\eq$ नसलेले प्रत्येक $\varphi \in \Gamma^*$
पूर्त करते (आणि म्हणून सर्व $\varphi \in \Gamma$ पूर्त करते). पण~$\eq$
उपस्थित असेल, तर ही रचना चालत नाही. कारण मग~$\Gamma^*$ मध्ये
वाक्य~$\eq[t][t']$ असू शकते, पण पद-प्रतिमानात कोणत्याही पदाचे
मूल्य ते पद स्वतःच असते. म्हणून $t$ आणि $t'$ ही भिन्न पदे असतील, तर
पद-प्रतिमानातील त्यांची मूल्ये---म्हणजे अनुक्रमे $t$ आणि $t'$---भिन्न
असतात; त्यामुळे $\eq[t][t']$ असत्य असते. मात्र ``भागाकार घेणे'' म्हणून
ओळखली जाणारी रचना वापरून हे दुरुस्त करता येते.
\end{explain}

\begin{defn}
  $\Gamma^*$ हा~$\Lang L$ मधील वाक्यांचा सुसंगत आणि संपूर्ण संच
  असू द्या. $\Lang L$ मधील बंद पदांच्या संचावर आपण~$\approx$ हा संबंध
  पुढीलप्रमाणे परिभाषित करतो:
  \[
  t \approx t' \text{\quad तेव्हा आणि केवळ तेव्हाच \quad} \eq[t][t'] \in \Gamma^*
  \]
\end{defn}

\begin{prop}
\label{fol:com:ide:prop:approx-equiv}
$\approx$ या संबंधाला पुढील गुणधर्म आहेत:
\begin{enumerate}
\item $\approx$ परावर्ती आहे.
\item $\approx$ सममित आहे.
\item $\approx$ संक्रमक आहे.
\item $t \approx t'$, $f$ हे फलन, आणि $t_1$, \dots,
  $t_{i-1}$, $t_{i+1}$, \dots, $t_n$ ही बंद पदे असतील, तर
\[
\Atom{f}{t_1,\dots, t_{i-1}, t, t_{i+1}, \dots, t_n} \approx
\Atom{f}{t_1,\dots, t_{i-1}, t', t_{i+1}, \dots, t_n}.
\]
\item $t \approx t'$, $R$ हे विधेय, आणि $t_1$, \dots,
  $t_{i-1}$, $t_{i+1}$, \dots, $t_n$ ही बंद पदे असतील, तर
\begin{multline*}
\Atom{R}{t_1,\dots, t_{i-1}, t, t_{i+1}, \dots, t_n} \in \Gamma^* \text{ तेव्हा आणि केवळ तेव्हाच } \\
\Atom{R}{t_1,\dots, t_{i-1}, t', t_{i+1}, \dots, t_n} \in \Gamma^*.
\end{multline*}
\end{enumerate}
\end{prop}

\begin{proof}
$\Gamma^*$ सुसंगत आणि संपूर्ण असल्यामुळे $\eq[t][t'] \in
\Gamma^*$ तेव्हा आणि केवळ तेव्हाच, जेव्हा $\Gamma^* \Proves
\eq[t][t']$. म्हणून पुढील विधाने दाखवणे पुरेसे आहे:
\begin{enumerate}
\item प्रत्येक बंद पद~$t$ साठी $\Gamma^* \Proves \eq[t][t]$.
\item $\Gamma^* \Proves \eq[t][t']$ असेल, तर $\Gamma^* \Proves \eq[t'][t]$.
\item $\Gamma^* \Proves \eq[t][t']$ आणि $\Gamma^* \Proves
  \eq[t'][t'']$ असेल, तर $\Gamma^* \Proves \eq[t][t'']$.
\item $\Gamma^* \Proves \eq[t][t']$ असेल, तर
\[
\Gamma^* \Proves
\eq[\Atom{f}{t_1,\dots,t_{i-1},t,t_{i+1},\dots,t_n}][\Atom{f}{t_1,\dots,t_{i-1},t',t_{i+1},\dots,t_n}]
\]
प्रत्येक $n$-स्थानी फलन~$f$ आणि बंद पदे $t_1$, \dots,
$t_{i-1}$, $t_{i+1}$, \dots,~$t_n$ यांच्यासाठी.
\item $\Gamma^* \Proves \eq[t][t']$ आणि
$\Gamma^* \Proves
\Atom{R}{t_1,\dots,t_{i-1},t,t_{i+1},\dots,t_n}$ असेल, तर
\[
\Gamma^* \Proves
\Atom{R}{t_1,\dots,t_{i-1},t',t_{i+1},\dots,t_n}
\]
प्रत्येक $n$-स्थानी विधेय~$R$ आणि बंद पदे $t_1$, \dots,
$t_{i-1}$, $t_{i+1}$, \dots,~$t_n$ यांच्यासाठी.
\end{enumerate}

\end{proof}

\begin{prob}
\ref{fol:com:ide:prop:approx-equiv} ची सिद्धता पूर्ण करा.
\end{prob}

\begin{defn}
$\Gamma^*$ हा भाषा~$\Lang L$ मधील सुसंगत आणि संपूर्ण संच आहे,
$t$ हे बंद पद आहे, आणि~$\approx$ हा मागील व्याख्येतील संबंध आहे असे
समजा. मग:
\[
\equivrep{t}{\approx} = \Setabs{t'}{t'\in \Trm[L], t \approx t'}
\]
आणि $\equivclass{\Trm[L]}{\approx} = \Setabs{\equivrep{t}{\approx}}{t \in \Trm[L]}$.
\end{defn}

\begin{defn}
\label{fol:com:ide:defn:term-model-factor}
$\Struct M = \Struct M(\Gamma^*)$ हे \ref{fol:com:mod:defn:termmodel} मधील
$\Gamma^*$ चे पद-प्रतिमान असू द्या. मग
$\Struct{\equivclass{M}{\approx}}$ ही पुढील संरचना आहे:
\begin{enumerate}
\item $\Domain{\equivclass{M}{\approx}} = \equivclass{\Trm[L]}{\approx}$.
\item $\Assign{c}{\equivclass{M}{\approx}} = \equivrep{c}{\approx}$
\item $\Assign{f}{\equivclass{M}{\approx}}(\equivrep{t_1}{\approx}, \dots,
  \equivrep{t_n}{\approx}) = \equivrep{\Atom{f}{t_1,\dots, t_n}}{\approx}$
\item $\tuple{\equivrep{t_1}{\approx}, \dots, \equivrep{t_n}{\approx}} \in
  \Assign{R}{\equivclass{M}{\approx}}$ तेव्हा आणि केवळ तेव्हाच, जेव्हा
  $\Sat{M}{\Atom{R}{t_1,\dots, t_n}}$; म्हणजे तेव्हा आणि केवळ तेव्हाच,
  जेव्हा $\Atom{R}{t_1,\dots, t_n} \in \Gamma^*$.
\end{enumerate}
\end{defn}

\begin{explain}
लक्षात घ्या की~$\equivclass{\Trm[L]}{\approx}$ च्या घटकांना
$\equivrep{t}{\approx}$ असे संबोधून, म्हणजेच
\emph{प्रतिनिधी}~$t \in \equivrep{t}{\approx}$ यांच्या साहाय्याने,
आपण त्यांच्यासाठी $\Assign{f}{\equivclass{M}{\approx}}$ आणि
$\Assign{R}{\equivclass{M}{\approx}}$ परिभाषित केली आहेत. या व्याख्या
प्रतिनिधींच्या निवडीवर अवलंबून नाहीत याची आपल्याला खात्री केली पाहिजे.
म्हणजे, तेच तुल्यतावर्ग निश्चित करणारी दुसरी~$t'$ निवडली
($\equivrep{t}{\approx} = \equivrep{t'}{\approx}$), तरी व्याख्यांतून तेच
फल मिळाले पाहिजे. उदाहरणार्थ, $R$ हे एकस्थानी विधेय असेल, तर
व्याख्येच्या शेवटच्या बाबीनुसार $\equivrep{t}{\approx} \in
\Assign{R}{\equivclass{M}{\approx}}$ तेव्हा आणि केवळ तेव्हाच, जेव्हा
$\Sat{M}{\Atom{R}{t}}$. आता $t \approx t'$ असलेल्या दुसऱ्या
पद~$t'$ साठी $\Sat/{M}{\Atom{R}{t'}}$ असेल, तर व्याख्येनुसार
$\equivrep{t'}{\approx} \notin \Assign{R}{\equivclass{M}{\approx}}$ असावे
लागेल. $t \approx t'$ असेल, तर $\equivrep{t}{\approx} =
\equivrep{t'}{\approx}$; पण $\equivrep{t}{\approx} \in
\Assign{R}{\equivclass{M}{\approx}}$ आणि $\equivrep{t}{\approx}
\notin \Assign{R}{\equivclass{M}{\approx}}$ ही दोन्ही विधाने लागू होऊ
शकत नाहीत. मात्र \ref{fol:com:ide:prop:approx-equiv} मुळे असे होऊ शकत नाही, याची
हमी मिळते.

\end{explain}

\begin{prop}
$\Struct{\equivclass{M}{\approx}}$ सुव्याख्यायित आहे; म्हणजे $t_1$,
  \dots, $t_n$, $t_1'$, \dots, $t_n'$ ही बंद पदे असतील आणि
  $t_i \approx t_i'$ असेल, तर
\begin{enumerate}
\item $\equivrep{\Atom{f}{t_1,\dots, t_n}}{\approx} =
    \equivrep{\Atom{f}{t_1',\dots, t_n'}}{\approx}$, म्हणजे
  \[
  \Atom{f}{t_1,\dots, t_n} \approx \Atom{f}{t_1',\dots, t_n'}
  \]
  आणि
\item $\Sat{M}{\Atom{R}{t_1,\dots, t_n}}$ तेव्हा आणि केवळ तेव्हाच,
  जेव्हा $\Sat{M}{\Atom{R}{t_1',\dots, t_n'}}$; म्हणजे
  \[
    \Atom{R}{t_1,\dots, t_n} \in \Gamma^* \text{ तेव्हा आणि केवळ तेव्हाच }
    \Atom{R}{t_1',\dots, t_n'} \in \Gamma^*.
  \]
\end{enumerate}
\end{prop}

\begin{proof}
\ref{fol:com:ide:prop:approx-equiv} वरून~$n$ वरील विगमनाने मिळते.
\end{proof}

पद-प्रतिमानाच्या बाबतीत होते त्याप्रमाणेच, सत्यता पूर्वप्रमेय सिद्ध
करण्यापूर्वी आपल्याला पुढील पूर्वप्रमेयाची गरज आहे.

\begin{lem}
\label{fol:com:ide:lem:val-in-termmodel-factored} $\Struct M = \Struct M
 (\Gamma^*)$ असू द्या; मग $\Value{t}{\equivclass{M}{\approx}} = \equivrep{t}
 {\approx}$.
\end{lem}
\begin{proof}
ही सिद्धता \ref{fol:com:mod:lem:val-in-termmodel} च्या सिद्धतेसारखीच आहे.
\end{proof}

\begin{prob}
\ref{fol:com:ide:lem:val-in-termmodel-factored} ची सिद्धता पूर्ण करा.
\end{prob}

\begin{lem}
\label{fol:com:ide:lem:truth}
सर्व वाक्ये~$\varphi$ यांच्यासाठी $\Sat{\equivclass{M}{\approx}}{\varphi}$ तेव्हा आणि
केवळ तेव्हाच, जेव्हा $\varphi \in \Gamma^*$.
\end{lem}

\begin{proof}
\ref{fol:com:mod:lem:truth} च्या सिद्धतेप्रमाणेच~$\varphi$ वरील विगमनाने सिद्धता
मिळते. $\varphi \ident \eq[t][t']$ असण्याच्या प्रसंगाकडेच फक्त अधिक लक्ष
देण्याची गरज आहे.
\begin{align*}
\Sat{\equivclass{M}{\approx}}{\eq[t][t']} & \text{ तेव्हा आणि केवळ तेव्हाच } \equivrep{t}{\approx} = \equivrep{t'}{\approx}
\text{ ($\Struct{\equivclass{M}{\approx}}$ च्या व्याख्येनुसार)}\\
& \text{ तेव्हा आणि केवळ तेव्हाच } t \approx t' \text{ ($\equivrep{t}{\approx}$ च्या व्याख्येनुसार)}\\
& \text{ तेव्हा आणि केवळ तेव्हाच } \eq[t][t'] \in \Gamma^* \text{ ($\approx$ च्या व्याख्येनुसार).}
\end{align*}
\end{proof}

\begin{digress}
लक्षात घ्या की $\Struct{M(\Gamma^*)}$ हे नेहमी गणनीय आणि अनंत
असले, तरी $\Struct{\equivclass{M}{\approx}}$ सांत असू शकते; कारण
$\equivrep{t}{\approx}$ असे फक्त सांत अनेक वर्ग असू शकतात. हे अपेक्षितच
आहे, कारण~$\Gamma$ मध्ये अशी वाक्ये असू शकतात की ती सत्य असणारी
कोणतीही संरचना सांत असणे आवश्यक असते. उदाहरणार्थ,
$\lforall[x][\lforall[y][x = y]]$ हे सुसंगत वाक्य आहे, पण ज्याच्या
प्रांत मध्ये नेमका एकच घटक आहे अशाच संरचनांमध्ये
त्याची पूर्ति होते.
\end{digress}












\section{संपूर्णता प्रमेय}\label{fol:com:cth:sec}

\begin{explain}
आपल्या निष्पत्ती एकत्र करू: आपल्याला संपूर्णता प्रमेय मिळते.
\end{explain}

\begin{thm}[संपूर्णता प्रमेय]
\label{fol:com:cth:thm:completeness}
$\Gamma$ हा वाक्यांचा संच असू द्या. $\Gamma$ सुसंगत असेल, तर
तो पूर्ततायोग्य आहे.
\end{thm}

\begin{proof}
$\Gamma$ सुसंगत आहे असे समजा.
  \ref{fol:com:hen:lem:henkin} नुसार असा संतृप्त सुसंगत संच
  $\Gamma' \supseteq \Gamma$ असतो. \ref{fol:com:lin:lem:lindenbaum} नुसार
असा $\Gamma^* \supseteq \Gamma'$ असतो की तो
सुसंगत आणि संपूर्ण असतो.

  $\Gamma' \subseteq \Gamma^*$ असल्यामुळे प्रत्येक
  सूत्र~$\varphi(x)$ साठी~$\Gamma^*$ मध्ये पुढील रूपाचे वाक्य
  असते:
  $\lexists[x][\varphi(x)] \lif \varphi(c)$.
  म्हणून~$\Gamma^*$ संतृप्त आहे. $\Gamma$ मध्ये~$\eq$ येत नसेल, तर
\ref{fol:com:mod:lem:truth} नुसार
$\Sat{M(\Gamma^*)}{\varphi}$ तेव्हा आणि केवळ तेव्हाच, जेव्हा $\varphi \in
\Gamma^*$. यावरून विशेषतः प्रत्येक $\varphi \in \Gamma$ साठी
$\Sat{M(\Gamma^*)}{\varphi}$; म्हणून
$\Gamma$ पूर्ततायोग्य आहे.
  $\Gamma$ मध्ये~$\eq$ येत असेल, तर \ref{fol:com:ide:lem:truth} नुसार
  सर्व वाक्ये~$\varphi$ यांच्यासाठी
  $\Sat{\equivclass{M}{\approx}}{\varphi}$ तेव्हा आणि केवळ तेव्हाच, जेव्हा
  $\varphi \in \Gamma^*$. विशेषतः, प्रत्येक $\varphi \in \Gamma$ साठी
  $\Sat{\equivclass{M}{\approx}}{\varphi}$; म्हणून $\Gamma$ पूर्ततायोग्य आहे.
\end{proof}

\begin{cor}[संपूर्णता प्रमेय, दुसरी आवृत्ती]
\label{fol:com:cth:cor:completeness}
प्रत्येक $\Gamma$ आणि प्रत्येक वाक्ये~$\varphi$ साठी: जर
$\Gamma \Entails \varphi$, तर $\Gamma \Proves \varphi$.
\end{cor}

\begin{proof}
लक्षात घ्या की \ref{fol:com:cth:cor:completeness} आणि
\ref{fol:com:cth:thm:completeness} यांमधील $\Gamma$ वर सर्ववाची संख्यकीकरण केलेले
आहे. आपला गोंधळ होऊ नये म्हणून \ref{fol:com:cth:thm:completeness} हे वेगळा चर
वापरून पुन्हा मांडू: वाक्यांचा कोणताही संच~$\Delta$ दिला असता,
जर $\Delta$ सुसंगत असेल, तर तो पूर्ततायोग्य आहे. प्रतिपरिवर्तनाने,
$\Delta$ पूर्ततायोग्य नसेल, तर $\Delta$ विसंगत आहे. अनुप्रमेय सिद्ध
करण्यासाठी याचा वापर करू.

$\Gamma \Entails \varphi$ असे समजा. मग \ref{fol:syn:sem:prop:entails-unsat}
नुसार $\Gamma \cup \{\lnot \varphi\}$ अपूर्ततायोग्य आहे. $\Gamma \cup
\{\lnot \varphi\}$ हाच आपला $\Delta$ घेतल्यास,
\ref{fol:com:cth:thm:completeness} च्या मागील आवृत्तीनुसार $\Gamma \cup
\{\lnot \varphi\}$ विसंगत आहे. पुढील संदर्भांनुसार,

  \ref{fol:axd:prv:prop:prov-incons}, \ref{fol:seq:prv:prop:prov-incons}, \ref{fol:ntd:prv:prop:prov-incons}, \ref{fol:tab:prv:prop:prov-incons},
$\Gamma \Proves \varphi$.
\end{proof}


\begin{prob}
\ref{fol:com:cth:cor:completeness} वापरून
\ref{fol:com:cth:thm:completeness} सिद्ध करा आणि अशा रीतीने
संपूर्णता प्रमेयाच्या दोन्ही मांडण्या समतुल्य आहेत हे दाखवा.
\end{prob}





\begin{prob}
एखादी निष्पत्ती प्रणाली संपूर्ण असण्यासाठी तिचे नियम इतके समर्थ
असले पाहिजेत की प्रत्येक अपूर्ततायोग्य संच विसंगत असल्याचे ते सिद्ध करू
शकतील. संपूर्णता सिद्ध करण्यासाठी निष्पत्तीचे कोणते नियम आवश्यक
होते? यांपैकी कोणताही नियम सिद्धतेत कोठेच वापरलेला नाही का? या प्रश्नांची
उत्तरे देण्यासाठी, \ref{fol:com:cth:thm:completeness} च्या सिद्धतेपर्यंत
नेणाऱ्या कोणत्या निष्पत्तींमध्ये निष्पत्तीचे कोणते नियम वापरले, हे
दाखवणारी यादी किंवा आकृती तयार करा. या सिद्धतांमधील नियमांच्या कोणत्याही
अध्याहृत वापरांची नोंद अवश्य करा.
\end{prob}















\section{संहतता प्रमेय}\label{fol:com:com:sec}

संपूर्णता प्रमेयाची एक महत्त्वाची निष्पत्ती म्हणजे संहतता प्रमेय.
वाक्यांच्या एखाद्या संचाचा प्रत्येक \emph{सांत} उपसंच पूर्ततायोग्य असेल,
तर संपूर्ण संच पूर्ततायोग्य असतो---तो संच स्वतः अनंत असला तरीही, असे
संहतता प्रमेय सांगते. हे मुळीच उघड नाही. प्रथमदर्शनी तरी अशी शक्यता
नाकारली जाईल असे काही दिसत नाही की वाक्यांचे अनंत संच व्याघाती
असतील, पण व्याघात, जणू काही, अनंत संख्येमुळेच उद्भवत असेल. अशी परिस्थिती
नाकारता येते, असे संहतता प्रमेय सांगते: ज्यांचा प्रत्येक सांत उपसंच
पूर्ततायोग्य आहे असे वाक्यांचे अपूर्ततायोग्य अनंत संच नसतात.
संपूर्णता प्रमेयाप्रमाणे याची तार्किक निष्पन्नतेशी संबंधित आवृत्तीही आहे:
वाक्यांच्या अनंत संचापासून काही निष्पन्न होत असेल, तर ते आधीच
एखाद्या सांत उपसंचापासून निष्पन्न होते.

\begin{defn}
  सूत्रांचा संच~$\Gamma$ \emph{सांततः पूर्ततायोग्य} असतो तेव्हा आणि
  केवळ तेव्हाच, जेव्हा प्रत्येक सांत $\Gamma_0 \subseteq \Gamma$
  पूर्ततायोग्य असतो.
\end{defn}

\begin{thm}[संहतता प्रमेय]
\label{fol:com:com:thm:compactness}
वाक्यांचा कोणताही संच~$\Gamma$ आणि कोणतेही वाक्य~$\varphi$ यांच्यासाठी पुढील
विधाने लागू होतात:
\begin{enumerate}
  \item $\Gamma \Entails \varphi$ तेव्हा आणि केवळ तेव्हाच, जेव्हा असा सांत
    $\Gamma_0 \subseteq \Gamma$ असतो की $\Gamma_0 \Entails \varphi$.
  \item $\Gamma$ पूर्ततायोग्य असतो तेव्हा आणि केवळ तेव्हाच, जेव्हा तो
    सांततः पूर्ततायोग्य असतो.
\end{enumerate}
\end{thm}

\begin{proof}
आपण (2) सिद्ध करू. $\Gamma$ पूर्ततायोग्य असेल, तर अशी
संरचना~$\Struct{M}$
असते की प्रत्येक $\varphi \in \Gamma$ साठी
$\Sat{M}{\varphi}$. अर्थात, ही
$\Struct{M}$ देखील~$\Gamma$ चा प्रत्येक सांत
उपसंच पूर्त करते; म्हणून~$\Gamma$ सांततः पूर्ततायोग्य आहे.

आता $\Gamma$ सांततः पूर्ततायोग्य आहे असे समजा. मग प्रत्येक सांत
उपसंच~$\Gamma_0 \subseteq \Gamma$ पूर्ततायोग्य आहे. निर्दोषतेनुसार
(
  \ref{fol:axd:sou:cor:consistency-soundness}, \ref{fol:seq:sou:cor:consistency-soundness}, \ref{fol:ntd:sou:cor:consistency-soundness}, \ref{fol:tab:sou:cor:consistency-soundness}),
प्रत्येक सांत उपसंच सुसंगत आहे. मग पुढील संदर्भांनुसार $\Gamma$ स्वतः
सुसंगत असला पाहिजे:

  \ref{fol:axd:ptn:prop:proves-compact}, \ref{fol:seq:ptn:prop:proves-compact}, \ref{fol:ntd:ptn:prop:proves-compact}, \ref{fol:tab:ptn:prop:proves-compact}.
संपूर्णतेनुसार (\ref{fol:com:cth:thm:completeness}), $\Gamma$ सुसंगत
असल्यामुळे तो पूर्ततायोग्य आहे.
\end{proof}


\begin{prob}
\ref{fol:com:com:thm:compactness} मधील (1) सिद्ध करा.
\end{prob}





\begin{ex}
उपपत्ती~$\Gamma$ च्या प्रत्येक प्रतिमान~$\Struct{M}$ मध्ये प्रत्येक
पद~$t$ हे अर्थातच~$\Domain{M}$ मधील घटक निर्देशित करते.
$\Domain{M}$ मधील प्रत्येक घटक देखील कोणत्यातरी पदाने निर्देशित
केला जातो, याची हमी देता येते का? दुसऱ्या शब्दांत, ज्यांची सर्व प्रतिमाने
बंद पदांनी आच्छादित आहेत अशा उपपत्ती~$\Gamma$ असतात का? $\Gamma$ ला
अनंत प्रतिमाने असतील, तर असे नसते, हे संहतता प्रमेय दाखवते. ते पुढीलप्रमाणे
दिसते: $\Struct{M}$ हे~$\Gamma$ चे अनंत प्रतिमान असू द्या आणि $c$ हे
$\Gamma$ च्या भाषेत नसलेले स्थिरांक असू द्या. $\Gamma$ ची
भाषा~$\Lang{L}$ मधील प्रत्येक पद~$t$ साठी असलेल्या सर्व
$\eq/[c][t]$ या वाक्यांचा~$\Delta$ हा संच असू द्या, म्हणजे
  \[
  \Delta = \Setabs{\eq/[c][t]}{t \in \Trm[L]}.
  \]
$\Gamma \cup \Delta$ चा सांत उपसंच $\Gamma' \cup \Delta'$ असा लिहिता
येतो; येथे $\Gamma' \subseteq \Gamma$ आणि $\Delta' \subseteq \Delta$.
$\Delta'$ सांत असल्यामुळे त्यात फक्त सांत अनेक पदे असू शकतात.
त्यांपैकी कोणत्याही पदाने निर्देशित न केलेला $\Domain{M}$ मधील
घटक $a \in \Domain{M}$ असू द्या आणि $\Struct{M'}$ ही
$\Struct{M}$ सारखीच, पण $\Assign{c}{M'} = a$ देखील असलेली संरचना
असू द्या. $\Delta'$ मध्ये येणाऱ्या प्रत्येक~$t$ साठी
$a \neq \Value{t}{M}$ असल्यामुळे $\Sat{M'}{\Delta'}$. तसेच
$\Sat{M}{\Gamma}$, $\Gamma' \subseteq \Gamma$, आणि~$\Gamma$ मध्ये $c$
येत नसल्यामुळे $\Sat{M'}{\Gamma'}$ देखील. म्हणून $\Gamma \cup \Delta$
च्या प्रत्येक सांत उपसंच $\Gamma' \cup \Delta'$ साठी
$\Sat{M'}{\Gamma' \cup \Delta'}$. त्यामुळे $\Gamma \cup \Delta$ चा
प्रत्येक सांत उपसंच पूर्ततायोग्य आहे. संहततेनुसार $\Gamma \cup \Delta$
स्वतः पूर्ततायोग्य आहे. म्हणून $\Sat{M}{\Gamma \cup \Delta}$ अशी
प्रतिमाने असतात. असे प्रत्येक $\Struct{M}$ हे~$\Gamma$ चे प्रतिमान आहे,
पण ते बंद पदांनी आच्छादित नाही; कारण~$\Lang{L}$ मधील सर्व पदे~$t$
यांच्यासाठी $\Value{c}{M} \neq \Value{t}{M}$.
\end{ex}

\begin{ex}
विधेय~$<$, स्थिरांक~$\Obj{0}$, $\Obj{1}$, आणि
फलने~$+$, $\times$, व~$-$ असलेल्या भाषा~$\Lang{L}$ चा
विचार करा. प्रांत~$\Rat$ आणि उघड निर्वचनार्थ असलेल्या
संरचना~$\Struct{Q}$ मध्ये सत्य असलेल्या या भाषा मधील सर्व
वाक्यांचा संच~$\Gamma$ असू द्या. $\Gamma$ म्हणजे परिमेय
संख्यांविषयी सत्य असलेल्या~$\Lang{L}$ मधील सर्व वाक्यांचा संच.
अर्थात~$\Rat$ मध्ये (आणि~$\Real$ मध्येही) अशा संख्या~$r$ नसतात की त्या
$0$ पेक्षा मोठ्या, पण प्रत्येक $k \in \PosInt$ साठी $1/k$ पेक्षा लहान
असतात. अशी संख्या अस्तित्वात असती, तर ती \emph{अतिसूक्ष्म संख्या} असती:
शून्येतर, पण अनंतपटीने लहान. $\Gamma$ ची अतिसूक्ष्म संख्या असलेली प्रतिमाने
अस्तित्वात आहेत, हे दाखवण्यासाठी संहतता प्रमेय वापरता येते. आपल्या भाषेत
भागाकाराचे फलन नाही (शून्याने भागाकार अपरिभाषित असतो आणि
फलनांचे निर्वचन सर्वत्र परिभाषित फलनांनी केले पाहिजे). तरीही
$r < 1/k$ हे व्यक्त करता येते; कारण हे तेव्हा आणि केवळ तेव्हाच लागू होते,
जेव्हा $r \cdot k < 1$. आता $c$ हे नवीन स्थिरांक असू द्या आणि
$\Delta$ पुढीलप्रमाणे असू द्या:
\[
\{\Obj{0} < c\}
\cup \Setabs{ c \times \num k < \Obj{1} }{k \in \PosInt}
\]
(येथे $\num{k} = (\Obj{1} + (\Obj{1} + \dots + (\Obj{1} +
\Obj{1})\dots))$ मध्ये $k$ इतके~$\Obj{1}$ आहेत). $\Delta$ च्या कोणत्याही
सांत उपसंच~$\Delta_0$ साठी असा~$K$ असतो की~$\Delta_0$ मधील सर्व
वाक्ये $c \times \num{k} < \Obj{1}$ यांमध्ये $k < K$ असते.
$\Assign{c}{Q'} = 1/K$ असे ठेवून $\Struct{Q}$ चा विस्तार
$\Struct{Q'}$ पर्यंत केला, तर कोणत्याही सांत $\Gamma_0 \subseteq \Gamma$
साठी $\Sat{Q'}{\Gamma_0 \cup \Delta_0}$, आणि म्हणून $\Gamma \cup \Delta$
सांततः पूर्ततायोग्य आहे (सराव: हे सविस्तर सिद्ध करा). संहततेनुसार
$\Gamma \cup \Delta$ पूर्ततायोग्य आहे. $\Gamma \cup \Delta$ च्या
कोणत्याही प्रतिमान~$\Struct{S}$ मध्ये एक अतिसूक्ष्म संख्या असते, ती म्हणजे
$\Assign{c}{S}$.
\end{ex}



\begin{prob}
अंकगणिताच्या प्रमाण प्रतिमान~$\Struct{N}$ मध्ये प्रत्येक सूत्र~$\num{n}
< x$ ची पूर्ति करणारा कोणताही घटक~$k \in \Domain{N}$ नाही
(येथे $\num{n}$ हे $n$ इतकी~$\prime$ चिन्हे असलेले
$\Obj{0}^{\prime\dots\prime}$ आहे). अंकगणिताच्या प्रमाण
प्रतिमान~$\Struct{N}$ मध्ये सत्य असलेली अंकगणिताच्या भाषेतील
वाक्ये ही प्रत्येक सूत्र~$\num{n} < x$ ची पूर्ति \emph{करणारा}
घटक असलेल्या संरचना~$\Struct{N'}$ मध्येही सत्य आहेत,
हे दाखवण्यासाठी संहतता प्रमेय वापरा.
\end{prob}



\begin{ex}
एकरूपता असलेले प्रथम-कोटी तर्कशास्त्र हे प्रांताच्या आकाराला
एखादी किमान मर्यादा असली पाहिजे, असे व्यक्त करू शकते हे आपल्याला माहीत
आहे: वाक्य~$\varphi_{\ge n}$ (जे ``किमान $n$ भिन्न वस्तू आहेत'' असे सांगते)
फक्त अशा रचनांमध्ये सत्य असते ज्यांत $\Domain{M}$ मध्ये किमान~$n$ वस्तू
असतात. म्हणून
  \[
  \Delta = \Setabs{\varphi_{\ge n}}{n \ge 1}
  \]
असे घेतले, तर~$\Delta$ चे कोणतेही प्रतिमान अनंत असले पाहिजे. म्हणून
एखाद्या उपपत्तीला फक्त अनंत प्रतिमानेच आहेत याची हमी तिच्यात~$\Delta$
मिळवून देता येते: $\Gamma \cup \Delta$ ची प्रतिमाने म्हणजे~$\Gamma$ ची
सर्व आणि फक्त अनंत प्रतिमाने.

अशा रीतीने प्रथम-कोटी तर्कशास्त्र अनंतता व्यक्त करू शकते. पण ते सांतता
व्यक्त करू शकत नाही, हे संहतता प्रमेय दाखवते. कारण वाक्यांचा
एखादा संच~$\Lambda$ सर्व आणि फक्त सांत संरचनांमध्ये पूर्त होतो
असे समजा. मग $\Delta \cup \Lambda$ सांततः पूर्ततायोग्य आहे. का? सांत
$\Delta' \cup \Lambda' \subseteq \Delta \cup \Lambda$ असा संच असू
द्या; येथे $\Delta' \subseteq \Delta$ आणि $\Lambda' \subseteq \Lambda$.
$\varphi_{\ge n} \in \Delta'$ असेल अशा संख्यांपैकी सर्वांत मोठी संख्या~$n$
असू द्या. $\Lambda$ ची सर्व सांत रचनांमध्ये पूर्ति होत असल्यामुळे,
त्याला सांत अनेक, पण $\ge n$ इतके घटक असलेले
प्रतिमान~$\Struct{M}$ असते. पण मग $\Sat{M}{\Delta' \cup \Lambda'}$.
संहततेनुसार $\Delta \cup \Lambda$ ला अनंत प्रतिमान असते; पण हे
$\Lambda$ ची पूर्ति फक्त सांत संरचनांमध्ये होते या गृहीतकाच्या
विरुद्ध आहे.
\end{ex}













\section{संहतता प्रमेयाची प्रत्यक्ष सिद्धता}\label{fol:com:cpd:sec}

संपूर्णता प्रमेयाच्या सिद्धतेतील कल्पनाच वापरून, संपूर्णता प्रमेयाचा आधार
न घेता संहतता प्रमेय प्रत्यक्ष सिद्ध करता येते. संपूर्णता प्रमेयाच्या
सिद्धतेत आपण वाक्यांच्या सुसंगत संच~$\Gamma$ पासून सुरुवात केली,
त्याचा विस्तार करून वाक्यांचा सुसंगत, संतृप्त आणि
संपूर्ण संच~$\Gamma^*$ मिळवला; आणि मग~$\Gamma^*$ पासून रचलेल्या
पद-प्रतिमान~$\Struct{M(\Gamma^*)}$
मध्ये~$\Gamma$ मधील सर्व वाक्ये सत्य आहेत, म्हणजे~$\Gamma$
पूर्ततायोग्य आहे, हे दाखवले.

सांततः पूर्ततायोग्य वाक्यसंच पूर्ततायोग्य असतो हे दाखवण्यासाठी हीच पद्धत
वापरता येते. फक्त सत्यता पूर्वप्रमेयापर्यंत नेणाऱ्या निष्पत्तींच्या अशा
समांतर आवृत्त्या सिद्ध कराव्या लागतील की त्यांत ``सुसंगत'' ऐवजी ``सांततः
पूर्ततायोग्य'' वापरलेले असेल.

\begin{prop}
\label{fol:com:cpd:prop:fsat-ccs}
$\Gamma$ हा संपूर्ण आणि सांततः पूर्ततायोग्य आहे असे समजा. मग:
\begin{enumerate}
\item $(\varphi \land \psi) \in \Gamma$
  तेव्हा आणि केवळ तेव्हाच, जेव्हा $\varphi \in \Gamma$ आणि $\psi \in \Gamma$
  ही दोन्ही विधाने लागू होतात.

\item $(\varphi \lor \psi) \in \Gamma$ तेव्हा आणि केवळ तेव्हाच,
  जेव्हा $\varphi \in \Gamma$ किंवा $\psi \in \Gamma$ यांपैकी एक लागू होते.

\item $(\varphi \lif \psi) \in \Gamma$ तेव्हा आणि केवळ तेव्हाच,
  जेव्हा $\varphi \notin \Gamma$ किंवा $\psi \in \Gamma$ यांपैकी एक लागू होते.
\end{enumerate}
\end{prop}


\begin{prob}
\ref{fol:com:cpd:prop:fsat-ccs} सिद्ध करा. $\Proves$ चा वापर टाळा.
\end{prob}





\begin{lem}
\label{fol:com:cpd:lem:fsat-henkin} प्रत्येक सांततः पूर्ततायोग्य संच~$\Gamma$ चा
विस्तार करून संतृप्त आणि सांततः पूर्ततायोग्य संच~$\Gamma'$ मिळवता येतो.
\end{lem}



\begin{prob}
\ref{fol:com:cpd:lem:fsat-henkin} सिद्ध करा. (सूचना: महत्त्वाची
पायरी म्हणजे निष्पत्त्या किंवा सुसंगततेचा अजिबात आधार न घेता,
$\Gamma_n$ सांततः पूर्ततायोग्य असेल, तर $\Gamma_n \cup \{\theta_n\}$
देखील सांततः पूर्ततायोग्य आहे हे दाखवणे.)
\end{prob}



\begin{prop}\label{fol:com:cpd:prop:fsat-instances}
$\Gamma$ हा संपूर्ण, सांततः पूर्ततायोग्य आणि संतृप्त आहे असे समजा.
\begin{enumerate}
\item $\lexists[x][\varphi(x)] \in \Gamma$ तेव्हा आणि केवळ तेव्हाच,
  जेव्हा किमान एका बंद पद~$t$ साठी $\varphi(t) \in \Gamma$.
\item $\lforall[x][\varphi(x)] \in \Gamma$ तेव्हा आणि केवळ तेव्हाच,
  जेव्हा सर्व बंद पदे~$t$ यांच्यासाठी $\varphi(t) \in \Gamma$.
\end{enumerate}
\end{prop}



\begin{prob}
\ref{fol:com:cpd:prop:fsat-instances} सिद्ध करा.
\end{prob}


\begin{lem}
\label{fol:com:cpd:lem:fsat-lindenbaum} प्रत्येक सांततः पूर्ततायोग्य संच~$\Gamma$
चा विस्तार करून संपूर्ण आणि सांततः पूर्ततायोग्य संच~$\Gamma^*$
मिळवता येतो.
\end{lem}


\begin{prob}
\ref{fol:com:cpd:lem:fsat-lindenbaum} सिद्ध करा. (सूचना: महत्त्वाची
पायरी म्हणजे $\Gamma_n$ सांततः पूर्ततायोग्य असेल, तर $\Gamma_n \cup
\{\varphi_n\}$ किंवा $\Gamma_n \cup \{\lnot \varphi_n\}$ यांपैकी एक संच सांततः
पूर्ततायोग्य असतो हे दाखवणे.)
\end{prob}




\begin{thm}[संहतता]
\label{fol:com:cpd:thm:compactness-direct} $\Gamma$ पूर्ततायोग्य असतो तेव्हा आणि
केवळ तेव्हाच, जेव्हा तो सांततः पूर्ततायोग्य असतो.
\end{thm}

\begin{proof}
$\Gamma$ पूर्ततायोग्य असेल, तर अशी
संरचना~$\Struct{M}$
असते की प्रत्येक $\varphi \in \Gamma$ साठी
$\Sat{M}{\varphi}$. अर्थात, ही
$\Struct{M}$ देखील~$\Gamma$ चा प्रत्येक
सांत उपसंच पूर्त करते; म्हणून~$\Gamma$ सांततः पूर्ततायोग्य आहे.

आता $\Gamma$ सांततः पूर्ततायोग्य आहे असे समजा.
  \ref{fol:com:cpd:lem:fsat-henkin} नुसार असा सांततः पूर्ततायोग्य, संतृप्त संच
  $\Gamma' \supseteq \Gamma$ असतो. \ref{fol:com:cpd:lem:fsat-lindenbaum} नुसार
$\Gamma'$ चा विस्तार करून संपूर्ण आणि सांततः
पूर्ततायोग्य संच~$\Gamma^*$ मिळवता येतो, आणि $\Gamma^*$
  अजूनही संतृप्त असतो. \ref{fol:com:mod:defn:termmodel} मधीलप्रमाणे
पद-प्रतिमान~$\Struct{M(\Gamma^*)}$
रचा. लक्षात घ्या की \ref{fol:com:mod:prop:quant-termmodel} हे
$\Gamma^*$ सुसंगत आहे (किंवा त्या दृष्टीने संपूर्ण वा संतृप्त आहे)
या तथ्यावर अवलंबून नव्हते; ते फक्त $\Struct{M(\Gamma^*)}$ हे बंद पदांनी
आच्छादित आहे या तथ्यावर अवलंबून होते.
सत्यता पूर्वप्रमेयाची सिद्धता (\ref{fol:com:mod:lem:truth}) तशीच लागू होते,
जर \ref{fol:com:ccs:prop:ccs} च्या संदर्भाऐवजी
\ref{fol:com:cpd:prop:fsat-ccs} चा संदर्भ वापरला आणि
\ref{fol:com:hen:prop:saturated-instances} च्या संदर्भाऐवजी
\ref{fol:com:cpd:prop:fsat-instances} चा संदर्भ वापरला.

\end{proof}


\begin{prob}
\ref{fol:com:cpd:thm:compactness-direct} च्या सिद्धतेला आवश्यक असलेली
सत्यता पूर्वप्रमेयाची (\ref{fol:com:mod:lem:truth}) संपूर्ण सिद्धता
लिहा.
\end{prob}














\section{लोव्हेनहाइम--स्कोलेम प्रमेय}\label{fol:com:dls:sec}

एखाद्या उपपत्तीला अनंत प्रतिमान असेल, तर तिला फार तर गणनीय अनंत
असलेले प्रतिमानही असते, असे लोव्हेनहाइम--स्कोलेम प्रमेय सांगते. या
तथ्याची तात्काळ निष्पत्ती अशी की एखाद्या संरचनेचा आकार
अगणनीय आहे हे प्रथम-कोटी तर्कशास्त्र व्यक्त करू शकत नाही: सर्व
अगणनीय संरचनांमध्ये पूर्त होणारे कोणतेही वाक्य
किंवा वाक्यांचा कोणताही संच हा एखाद्या गणनीय रचनेतही
पूर्त होतो.

\begin{thm}
\label{fol:com:dls:thm:downward-ls} $\Gamma$ सुसंगत असेल, तर त्याला
गणनीय प्रतिमान असते; म्हणजे ज्याचा प्रांत सांत किंवा
गणनीय अनंत आहे अशा संरचनांमध्ये तो पूर्ततायोग्य असतो.
\end{thm}

\begin{proof}
$\Gamma$ सुसंगत असेल, तर संपूर्णता प्रमेयाच्या सिद्धतेतून मिळणाऱ्या
संरचना~$\Struct M$ चा प्रांत~$\Domain{M}$ हा भाषा~$\Lang L$ मधील
पदांच्या संचापेक्षा मोठा नसतो. म्हणून $\Struct M$ फार तर
गणनीय अनंत आहे.
\end{proof}

\begin{thm}
\label{fol:com:dls:noidentity-ls} $\Gamma$ हा एकरूपतेविना प्रथम-कोटी तर्कशास्त्राच्या
भाषेतील वाक्यांचा सुसंगत संच असेल, तर त्याला गणनीय अनंत
प्रतिमान असते; म्हणजे ज्याचा प्रांत अनंत आणि गणनीय आहे अशा
संरचनांमध्ये तो पूर्ततायोग्य असतो.
\end{thm}

\begin{proof}
$\Gamma$ सुसंगत असेल आणि त्यात एकरूपता येणारी वाक्ये नसतील, तर संपूर्णता
प्रमेयाच्या सिद्धतेतून मिळणाऱ्या संरचना~$\Struct M$ चा
प्रांत~$\Domain{M}$ हा भाषा~$\Lang L'$ मधील पदांच्या संचाशी एकरूप असतो.
म्हणून $\Trm[L']$ हे गणनीय अनंत असल्यामुळे $\Struct{M}$ देखील
गणनीय अनंत आहे.
\end{proof}

\begin{ex}[स्कोलेमची विरोधापत्ती]
त्सेर्मेलो--फ्रेंकेल संचसिद्धान्त~$\Log{ZFC}$ ही एक अतिशय समर्थ चौकट
आहे; संचांच्या आकारांविषयीच्या तथ्यांसह जवळजवळ सर्व गणिती विधाने तिच्यात
व्यक्त करता येतात. म्हणून, उदाहरणार्थ, वास्तव संख्यांचा संच~$\Real$ हा
अगणनीय आहे हे~$\Log{ZFC}$ सिद्ध करू शकते; कोणत्याही संचाचा
घातसंच त्या संचापेक्षा मोठा असतो हे कँटर यांचे प्रमेयही ती सिद्ध करू शकते;
आणि असेच पुढे. $\Log{ZFC}$ सुसंगत असेल, तर तिची सर्व प्रतिमाने अनंत
असतात. शिवाय, त्या सर्व प्रतिमानांत असे घटक असतात की उपपत्तीच्या
मते ते अगणनीय असतात; उदाहरणार्थ, नैसर्गिक संख्यांचा घातसंच
अस्तित्वात आहे हे~$\Log{ZFC}$ चे प्रमेय सत्य करणारा घटक. लोव्हेनहाइम--स्कोलेम
प्रमेयानुसार $\Log{ZFC}$ ला गणनीय प्रतिमानेही असतात---ज्यांत
``अगणनीय'' संच असतात, पण जी स्वतः गणनीय असतात.
\end{ex}



\chapter{प्रथम-क्रम तर्कशास्त्राचा परिचय}\label{fol:int::chap}









\section{प्रथम-क्रम तर्कशास्त्र}\label{fol:int:fol:sec}

आकारिक तर्कशास्त्राशी तुमची पहिली ओळख झाली तेव्हापासून प्रथम-क्रम
तर्कशास्त्र बहुधा तुम्हाला परिचित असेल.\footnote{खरे तर, तुम्ही त्याच्याशी
परिचित आहात असे आम्ही जवळजवळ गृहीतच धरतो!{} तसे नसल्यास,
\emph{forall x} (Magnus आणि सहकारी, 2021) यासारख्या अधिक प्राथमिक
पाठ्यपुस्तकाची उजळणी करू शकता.} ते तुम्हाला ``संख्यापक तर्कशास्त्र''
किंवा ``विधेय तर्कशास्त्र'' या नावांनी माहीत असू शकते. सर्वप्रथम,
प्रथम-क्रम तर्कशास्त्र ही एक आकारिक भाषा आहे. म्हणजे तिचा एक विशिष्ट
शब्दसंग्रह असतो आणि तिच्या अभिव्यक्ती या त्या शब्दसंग्रहातील चिन्हमाला
असतात. पण प्रत्येक चिन्हमाला अनुमत नसते. अनुमत अभिव्यक्तींचे विविध
प्रकार आहेत: पदे, सूत्रे आणि वाक्ये. आपल्याला मुख्यतः
प्रथम-क्रम तर्कशास्त्रातील वाक्ये मध्ये रस आहे: ती आपल्याला
इंग्रजीतील वाक्यांचे आकारिक प्रतिरूप देतात आणि त्यांच्याविषयी
तर्कशास्त्रज्ञाला नेहमी रस असणारे प्रश्न आपण विचारू शकतो. उदाहरणार्थ:
\begin{itemize}
    \item $\psi$ हे~$\varphi$ पासून तार्किकरीत्या निष्पन्न होते का?
    \item $\varphi$ तार्किकरीत्या सत्य, तार्किकरीत्या असत्य, की
परिस्थितीवश आहे?
    \item $\varphi$ आणि $\psi$ सममूल्य आहेत का?
\end{itemize}

हे प्रश्न मुख्यतः प्रथम-क्रम तर्कशास्त्रातील वाक्यांच्या
``अर्था''विषयीचे आहेत. उदाहरणार्थ, $\psi$ हे~$\varphi$ पासून तार्किकरीत्या
निष्पन्न होते का या प्रश्नाचे तत्त्वज्ञ पुढीलप्रमाणे विश्लेषण करेल:
$\varphi$ सत्य पण~$\psi$ असत्य असणारा काही प्रसंग आहे का ($\psi$ हे~$\varphi$
पासून निष्पन्न होत नाही), की $\varphi$ सत्य करणारा प्रत्येक प्रसंग~$\psi$ ही
सत्य करतो ($\psi$ हे~$\varphi$ पासून निष्पन्न होते)? पण ``प्रसंग'' म्हणजे काय
हे अजून सांगितलेले नाही---ते \emph{चिन्हार्थमीमांसेचे} काम आहे.
प्रथम-क्रम तर्कशास्त्राची चिन्हार्थमीमांसा तत्त्वज्ञाच्या ``प्रसंग'' या
अंतःप्रज्ञ कल्पनेचे गणितीयदृष्ट्या काटेकोर प्रतिमान देते आणि---हे
महत्त्वाचे आहे---एखाद्या प्रसंगात वाक्य~$\varphi$ \emph{सत्य असणे}
म्हणजे काय याचेही प्रतिमान देते. आपण विकसित करणार असलेल्या या
गणितीयदृष्ट्या काटेकोर प्रतिमानाला संरचना म्हणतो. ``यात सत्य''
याला काटेकोर अर्थ देणाऱ्या संबंधाला \emph{पूर्ति}-संबंध म्हणतात. म्हणून
वाक्ये~$\varphi$ आणि संरचना~$\Struct{M}$ यांच्यासाठी आपण
``$\varphi$ ची~$\Struct{M}$ मध्ये पूर्ति होते'' (चिन्हांत:
$\Sat{M}{\varphi}$) हे परिभाषित करणार आहोत. हे झाल्यावर ``पासून निष्पन्न
होते'' किंवा ``तार्किकरीत्या सत्य आहे'' यांसारख्या इतर चिन्हार्थविषयक
संज्ञांच्याही काटेकोर व्याख्या देता येतील. उदाहरणार्थ,
$\lforall[x][(\varphi(x) \lif \psi(x))],
\lexists[x][\varphi(x)] \Entails \lexists[x][\psi(x)]$ आहे की नाही हे पुन्हा
गणितीय काटेकोरपणे ठरवणे या व्याख्यांमुळे शक्य होईल. अर्थातच उत्तर
``होय'' असेल. इंग्रजी वाक्यांचे प्रथम-क्रम तर्कशास्त्रात चिन्हांकन
करण्याचे प्रशिक्षण तुम्ही आधीच घेतले असेल, तर ही, उदाहरणार्थ, ``सर्व
मुंग्या कीटक आहेत, मुंग्या अस्तित्वात आहेत, म्हणून कीटक अस्तित्वात
आहेत'' यांची चिन्हांकने आहेत हे तुम्ही ओळखाल. हा उघडपणे वैध युक्तिवाद
आहे आणि म्हणून आपल्या आकारिक भाषेतील ``पासून निष्पन्न होते'' याच्या
गणितीय प्रतिमानानेही तेच उत्तर दिले पाहिजे.


आकारिक तर्कशास्त्राशी तुमची पहिली ओळख झाली तेव्हापासून आठवत असणारा
आणखी एक विषय म्हणजे \emph{निष्पत्त्या}. तुम्ही आकारिक
तर्कशास्त्राचा पहिला अभ्यासक्रम घेतला असेल, तर तुमच्या शिक्षकांनी अशा
निष्पत्त्या शोधण्याचा सराव तुमच्याकडून करून घेतला असेल; कदाचित
वरील तार्किक निष्पन्नता लागू असल्याचे दाखवणारी निष्पत्ती देखील.
निष्पत्त्या देण्याच्या अनेक पद्धती आहेत: तुम्ही ``नैसर्गिक निगमन''
किंवा ``सत्यता-वृक्ष'' नावाची पद्धत केली असेल; पण इतर अनेक पद्धतीही
आहेत. वरील तर्कशास्त्रीय प्रश्नांची उत्तरे देता येतील अशी साधने पुरवणे
हे निष्पत्ती-पद्धतींचे उद्दिष्ट आहे. उदाहरणार्थ, ज्या नैसर्गिक
निगमनातील निष्पत्तीमध्ये $\lforall[x][(\varphi(x) \lif \psi(x))]$ आणि
$\lexists[x][\varphi(x)]$ ही आधारविधाने आणि $\lexists[x][\psi(x)]$ हा
निष्कर्ष (शेवटची ओळ) आहे, ती $\lexists[x][\psi(x)]$ हे
$\lforall[x][(\varphi(x) \lif \psi(x))]$ आणि $\lexists[x][\varphi(x)]$ यांपासून
तार्किकरीत्या निष्पन्न होते याची \emph{पुष्टी करते}.

पण असे का? वरवर पाहता, निष्पत्ती-पद्धतींचा चिन्हार्थमीमांसेशी
काहीच संबंध नाही: आकारिक निष्पत्ती देणे म्हणजे केवळ काही
नियमांनी नियंत्रित रीतीने चिन्हांची मांडणी करणे; त्यात ``प्रसंगां''चा
किंवा ``यात सत्य''चा मुळीच उल्लेख नसतो. निष्पत्ती-पद्धती आणि
चिन्हार्थमीमांसा यांच्यातील संबंध अधितार्किक अन्वेषणाने प्रस्थापित
करावा लागतो. उदाहरणार्थ, आधारविधाने
$\lforall[x][(\varphi(x) \lif \psi(x))]$ आणि $\lexists[x][\varphi(x)]$ यांपासून
$\lexists[x][\psi(x)]$ ची आकारिक निष्पत्ती शक्य असते तेव्हा आणि
केवळ तेव्हाच $\lforall[x][(\varphi(x) \lif \psi(x))]$ आणि
$\lexists[x][\varphi(x)]$ ही दोन्ही मिळून
$\lexists[x][\psi(x)]$ तार्किकरीत्या निष्पन्न करतात, याची गणितीय
सिद्धता आवश्यक आहे. मात्र, हे करता येण्यापूर्वी विन्यासमीमांसा आणि
चिन्हार्थमीमांसा यांच्या व्याख्या अचूक करण्याचे बरेच कष्टप्रद काम
करावे लागते.












\section{विन्यासमीमांसा}\label{fol:int:syn:sec}

प्रथम-क्रम तर्कशास्त्रातील कोणत्या चिन्हमाला वाक्ये म्हणून गणल्या
जातात हे प्रथम काटेकोरपणे सांगावे लागेल. हे आपण नंतर करू; सध्या
उदाहरणांच्या साहाय्याने पुढे जाऊ. प्रथम-क्रम तर्कशास्त्राचे मूलभूत
रचनाघटक---म्हणजे त्याचा शब्दसंग्रह---दोन भागांत विभागलेले आहेत. पहिल्या
भागात विशिष्ट गोष्टींबद्दल काही सांगण्यासाठी किंवा विशिष्ट गोष्टी
निर्देशित करण्यासाठी आपण वापरतो ती चिन्हे येतात. गोष्टी निर्देशित
करण्यासाठी आपण स्थिरांक वापरतो आणि निर्देशित केलेल्या गोष्टींबद्दल
काही सांगण्यासाठी विधेये वापरतो. उदाहरणार्थ, एखादी एक गोष्ट
निर्देशित करण्यासाठी $\Obj a$ हे स्थिरांक म्हणून वापरू शकतो आणि मग
वाक्य~$\Atom{\Obj P}{\Obj a}$ वापरून तिच्याबद्दल काही सांगू शकतो.
``$\Obj a$'' आणि ``$\Obj P$'' यांचे काही अर्थ तुम्ही मनात धरले असतील,
तर $\Atom{\Obj P}{\Obj a}$ हे इंग्रजीतील वाक्य म्हणून वाचू शकता (आणि
आकारिक तर्कशास्त्र प्रथम शिकताना तुम्ही बहुधा तसे केलेही असेल).
प्रथम-क्रम तर्कशास्त्रातील अशी साधी वाक्ये मिळाल्यानंतर,
शब्दसंग्रहातील दुसरा भाग---तार्किक चिन्हे (संयोजक आणि संख्यापक)---वापरून
त्यांपासून अधिक जटिल वाक्ये रचता येतात. म्हणून, उदाहरणार्थ,
$(\Atom{\Obj P}{\Obj a} \land \Atom{\Obj Q}{\Obj b})$ किंवा
~$\lexists[\Obj x][\Atom{\Obj P}{\Obj x}]$ यांसारख्या अभिव्यक्ती
रचता येतात.

काटेकोर अधितार्किक अभ्यासासाठी आवश्यक असलेल्या चिन्हार्थमीमांसेच्या
नेमक्या व्याख्या आणि आपल्या निष्पत्ती-पद्धतींचे नियम देता यावेत,
यासाठी प्रथम-क्रम तर्कशास्त्रातील वाक्य म्हणून नेमके काय गणले
जाते याची काटेकोर व्याख्या सर्वप्रथम द्यावी लागेल. मूलभूत कल्पना समजणे
पुरेसे सोपे आहे: फक्त विधेये आणि स्थिरांक यांपासून आपण
$\Atom{\Obj P}{\Obj a}$ यासारखी काही साधी वाक्ये रचू शकतो. मग
संयोजक आणि संख्यापक वापरून त्यांपासून अधिक जटिल वाक्ये रचतो. पण अधिक
जटिल वाक्ये रचण्याची अनुमती देणारे नियम नेमके कोणते? ते स्पष्ट
करावेच लागतील; अन्यथा ``प्रथम-क्रम तर्कशास्त्रातील वाक्य'' याची
आपण पुरेशी काटेकोर व्याख्या केलेली नाही. येथे काही अडचणी आहेत. पहिली
अडचण म्हणजे योग्य चिन्हमालाच वाक्ये म्हणून गणल्या जातील याची
खात्री करणे. दुसरी म्हणजे हे अशा रीतीने करणे की
\emph{सर्व}~वाक्ये विषयी गणितीय सिद्धता देता येतील. शेवटी,
``$\varphi$ मधील प्रत्येक $\Obj x$ च्या जागी~$\Obj b$ ठेवा'' यासारख्या
वाक्ये वरील काही प्राथमिक क्रियांच्याही काटेकोर व्याख्या द्याव्या
लागतील. प्रथम-क्रम तर्कशास्त्रातील संख्यापक आणि चले यांमुळे हे
कसे करावे हे पूर्णपणे स्पष्ट नसते, ही अडचण आहे. उदाहरणार्थ,
$\lexists[\Obj x][\Atom{\Obj P}{\Obj a}]$ हे वाक्य म्हणून
गणले जावे का? $\lexists[\Obj x][\lexists[\Obj x][\Atom{\Obj P}{\Obj
x}]]$ चे काय? आणि ``$(\Atom{\Obj P}{\Obj x} \land \lexists[\Obj
x][\Atom{\Obj P}{\Obj x}])$ मध्ये $\Obj x$ च्या जागी~$\Obj b$ ठेवा''
याचा परिणाम काय असावा?












\section{सूत्र}\label{fol:int:fml:sec}

प्रथम-क्रम तर्कशास्त्रातील वाक्ये काटेकोरपणे निर्दिष्ट करण्यासाठी
आणि चलांच्या वापरामुळे निर्माण होणाऱ्या अडचणी हाताळण्यासाठी
आपण पुढील पद्धत वापरणार आहोत. प्रथम आपण अभिव्यक्तींचा एक
\emph{वेगळा} संच परिभाषित करतो: सूत्रे. ते केल्यानंतर त्यांत
चले कोणती भूमिका बजावतात याचा विचार करता येतो---आणि
``मुक्त'' व ``बद्ध'' चले या काही इतर कल्पनांच्या आधारे
वाक्य म्हणजे काय हे परिभाषित करता येते (म्हणजे, मुक्त
चले नसलेले सूत्र). यामुळे ``वाक्य'' ची व्याख्या
अधिक हाताळण्याजोगी होते एवढ्याचसाठी आपण हे करत नाही; पूर्ति ही
चिन्हार्थविषयक संकल्पना ज्या रीतीने परिभाषित करू त्यासाठीही ते
निर्णायक ठरेल.

फक्त एक विधेय~$\Obj P$, एक स्थिरांक~$\Obj a$ आणि
$\lnot$, $\land$ व~$\lexists$ ही तार्किक चिन्हे असलेल्या एका साध्या
प्रथम-क्रम भाषेसाठी ``सूत्र'' परिभाषित करू या. आपल्या संपूर्ण
व्याख्या यापेक्षा खूप व्यापक असतील: आपण अनंतपणे अनेक विधेये
आणि स्थिरांक स्वीकारू. किंबहुना, स्थिरांक आणि चले
यांच्याशी जोडून ``पदे'' रचता येतील अशा फलनांचाही विचार करू.
सध्या $\Obj a$ आणि चरे एवढीच आपली पदे असतील. आपल्याला अनंतपणे अनेक
चले मात्र आवश्यक आहेत. अधिकृतपणे आपण $\Obj v_0$, $\Obj
v_1$, \dots, ही चिन्हे चरे म्हणून वापरू.

\begin{defn}
\emph{सूत्रे} चा संच~$\Frm$ पुढीलप्रमाणे परिभाषित केला आहे:
\begin{enumerate}
\item\label{fol:int:fml:fmls-atom} $\Atom{\Obj P}{\Obj a}$ आणि $\Atom{\Obj
  P}{\Obj v_i}$ ही सूत्रे आहेत ($i \in \Nat$).

\item \label{fol:int:fml:fmls-not}$\varphi$ हे सूत्र असेल, तर $\lnot \varphi$ हे
  सूत्र आहे.

\item $\varphi$ आणि $\psi$ ही सूत्रे असतील, तर $(\varphi \land
  \psi)$ हे सूत्र आहे.

\item \label{fol:int:fml:fmls-ex}$\varphi$ हे सूत्र आणि $x$ हे चल
  असेल, तर $\lexists[x][\varphi]$ हे सूत्र आहे.

\item \label{fol:int:fml:fmls-limit}दुसरे काहीही सूत्र नाही.
\end{enumerate}
\end{defn}

कोणत्याही $i \in \Nat$ साठी $\Atom{\Obj P}{\Obj a}$ आणि
$\Atom{\Obj P}{\Obj v_i}$ ही सूत्रे आहेत, असे
\ref{fol:int:fml:fmls-atom} सांगतो. त्यांना तथाकथित \emph{आण्विक} सूत्रे
म्हणतात. त्यांपासून सुरुवात करता येते. इतर उपवाक्ये आपण आधीच
रचलेल्यांपासून नवी सूत्रे रचण्याचे मार्ग देतात. म्हणून,
उदाहरणार्थ, $\Atom{\Obj P}{\Obj v_2}$ हे \ref{fol:int:fml:fmls-atom} नुसार आधीच
सूत्र असल्यामुळे, $\lnot \Atom{\Obj P}{\Obj v_2}$ हे
\ref{fol:int:fml:fmls-not} नुसार सूत्र आहे. मग
\ref{fol:int:fml:fmls-ex} नुसार $\lexists[\Obj v_2][\lnot \Atom{\Obj P}{\Obj
v_2}]$ हे आणखी एक सूत्र आहे; आणि असेच पुढे जाता येते. \emph{फक्त} याच
रीतीने रचता येणाऱ्या चिन्हमाला सूत्रे म्हणून गणल्या जातात, असे
\ref{fol:int:fml:fmls-limit} सांगतो. विशेषतः,
$\lexists[\Obj v_0][\Atom{\Obj P}{\Obj a}]$ आणि $\lexists[\Obj
  v_0][\lexists[\Obj v_0][\Atom{\Obj P}{\Obj a}]]$ ही
\emph{खरोखर} सूत्रे म्हणून गणली जातात; परंतु अनावश्यक बाहेरील
कंसांमुळे $(\lnot \Atom{\Obj P}{\Obj a})$ हे गणले जात नाही.

सूत्रे परिभाषित करण्याच्या या पद्धतीला \emph{विगमनाधारित
व्याख्या} म्हणतात आणि त्यामुळे \emph{संरचनात्मक विगमन} म्हणतात त्या
विगमनाने सिद्धता करण्याच्या प्रकाराचा वापर करून सूत्रे विषयी
गोष्टी सिद्ध करता येतात. यांची सर्वसाधारण चर्चा
\readerexternalref{mth:ind:idf:sec} आणि \readerexternalref{mth:ind:sti:sec} येथे केली
आहे; पुढील सिद्धतांचा सखोल अभ्यास करण्यापूर्वी तिची उजळणी करावी. मूलतः,
सर्व सूत्रे साठी एखादी गोष्ट सत्य आहे याची सिद्धता द्यायची असेल,
तर ती आण्विक सूत्रे साठी सत्य आहे हे प्रथम दाखवता आणि नंतर ती
कोणत्याही सूत्र~$\varphi$ साठी (आणि~$\psi$ साठी) सत्य \emph{असेल}, तर
ती $\lnot \varphi$, $(\varphi \land \psi)$ आणि $\lexists[x][\varphi]$ यांच्यासाठीही
\emph{सत्य असते}, हे दाखवता. उदाहरणार्थ, यावरून ती
$\lexists[\Obj v_2][\lnot \Atom{\Obj P}{\Obj v_2}]$ साठी सत्य आहे हे
सिद्ध होते: पहिल्या भागावरून ती आण्विक सूत्र~$\Atom{\Obj
P}{\Obj v_2}$ साठी सत्य आहे हे माहीत असते. मग दुसऱ्या भागावरून ती
$\lnot \Atom{\Obj P}{\Obj v_2}$ साठी सत्य आहे आणि त्याच भागावरून पुन्हा
$\lexists[\Obj v_2][\lnot \Atom{\Obj P}{\Obj v_2}]$ याच्यासाठी सत्य
आहे, हे मिळते. सर्व सूत्रे आण्विक सूत्रांपासून
विगमनाधारित रीतीने निर्माण केली जात असल्यामुळे, ही पद्धत त्यांपैकी
कोणत्याही सूत्रासाठी चालते.












\section{पूर्ति}\label{fol:int:sat:sec}

सूत्रे म्हणजे काय हे माहीत झाल्यावर आपण लगेच प्रथम-क्रम
तर्कशास्त्राच्या चिन्हार्थमीमांसेकडे वळू शकतो: येथे मूलभूत व्याख्या
संरचनेची आहे. आपल्या साध्या भाषेसाठी संरचना~$\Struct M$
चे फक्त तीन घटक असतात: \emph{प्रांत} म्हणतात तो अरिक्त संच
$\Domain{M}$, $\Struct M$ मध्ये $\Obj a$ ने निर्देशित केलेली गोष्ट, आणि
$\Struct M$ मध्ये $\Obj P$ ज्या गोष्टींबाबत सत्य आहे त्या गोष्टी.
$\Obj a$ ने निर्देशित केलेली वस्तू~$\Assign{\Obj a}{M}$ ने, तर
$\Obj P$ ज्या गोष्टींबाबत सत्य आहे त्यांचा संच~$\Assign{\Obj P}{M}$ ने
दर्शवतात. संरचना~$\Struct{M}$ मध्ये याच तीन गोष्टी असतात:
$\Domain{M}$, $\Assign{\Obj a}{M} \in \Domain{M}$ आणि
$\Assign{\Obj P}{M} \subseteq \Domain{M}$. सर्वसाधारण स्थिती अधिक
गुंतागुंतीची असेल, कारण त्यात अनेक विधेये आणि स्थिरांक
असतील, विधेये ना एकाहून अधिक स्थाने असू शकतील, आणि
फलने ही असतील.


ज्या सूत्रे मध्ये चले नाहीत त्यांच्यासाठी पूर्तिची
व्याख्या देण्यास हे पुरेसे आहे. आपण सूत्रे ज्या रीतीने परिभाषित
केली तिचे प्रतिबिंब दाखवणारी विगमनाधारित व्याख्या देण्याची कल्पना आहे.
एखाद्या आण्विक सूत्राची~$\Struct{M}$ मध्ये पूर्ति केव्हा होते हे आपण
सांगतो आणि मग, उदाहरणार्थ, $\varphi$ ची~$\Struct{M}$ मध्ये पूर्ति होते की
नाही याच्या आधारे $\lnot \varphi$ ची~$\Struct{M}$ मध्ये पूर्ति केव्हा होते
हे सांगतो. उदाहरणार्थ, पुढीलप्रमाणे व्याख्या करता येईल:
\begin{enumerate}
  \item $\Atom{\Obj P}{\Obj a}$ ची~$\Struct{M}$ मध्ये पूर्ति तेव्हा आणि
  केवळ तेव्हाच होते, जेव्हा $\Assign{\Obj a}{M} \in \Assign{\Obj P}{M}$.
  \item $\lnot \varphi$ ची~$\Struct{M}$ मध्ये पूर्ति तेव्हा आणि केवळ तेव्हाच
  होते, जेव्हा $\varphi$ ची~$\Struct{M}$ मध्ये पूर्ति होत नाही.
  \item $(\varphi \land \psi)$ ची~$\Struct{M}$ मध्ये पूर्ति तेव्हा आणि केवळ
  तेव्हाच होते, जेव्हा $\varphi$ ची~$\Struct{M}$ मध्ये पूर्ति होते आणि
  $\psi$ ची~$\Struct{M}$ मध्येही पूर्ति होते.
\end{enumerate}
$\Domain{M} = \{0, 1, 2\}$, $\Assign{\Obj a}{M} = 1$ आणि
$\Assign{\Obj P}{M} = \{1, 2\}$ असे समजू. ही व्याख्या आपल्याला
$\Atom{\Obj P}{\Obj a}$ ची~$\Struct{M}$ मध्ये पूर्ति होते असे सांगेल
(कारण $\Assign{\Obj a}{M} =  1 \in \{1,2\} = \Assign{\Obj P}{M}$).
ती पुढे असे सांगते की $\lnot \Atom{\Obj P}{\Obj a}$ ची~$\Struct{M}$
मध्ये पूर्ति होत नाही; म्हणून $\lnot\lnot \Atom{\Obj P}{\Obj a}$ ची
पूर्ति होते आणि $(\lnot \Atom{\Obj P}{\Obj a} \land \Atom{\Obj P}{\Obj a})$
ची पूर्ति होत नाही; आणि असेच पुढे जाता येते.

संख्यापकांसाठी व्याख्या द्यायची झाल्यावर अडचण निर्माण होते: आपल्याला
``$\lexists[\Obj v_0][\Atom{\Obj P}{\Obj v_0}]$ ची पूर्ति तेव्हा आणि
केवळ तेव्हाच होते, जेव्हा $\Atom{\Obj P}{\Obj v_0}$ ची पूर्ति होते''
असे काहीतरी म्हणायचे आहे. पण संरचना~$\Struct{M}$ हे
चले विषयी काय करायचे ते सांगत नाही. आपल्याला प्रत्यक्षात
$\Atom{\Obj P}{\Obj v_0}$ ची \emph{$\Obj v_0$ च्या एखाद्या मूल्यासाठी}
पूर्ति होते असे म्हणायचे आहे. हे काटेकोर करण्यासाठी फक्त $\Obj a$ लाच
नव्हे, तर~$\Obj v_0$ लाही~$\Domain{M}$ मधील घटक नेमून देण्याचा
मार्ग हवा. यासाठी आपण चल ची \emph{मूल्यांकने} आणतो.
चल चे मूल्यांकन म्हणजे चले ना~$\Domain{M}$ मधील
घटक कडे पाठवणारे फलन~$s$ (आपल्या उदाहरणात, $0$, $1$ किंवा~$2$
यांपैकी एका घटकाकडे). एखाद्या सूत्रामध्ये कोणती चले
येतील हे आपल्याला आधी माहीत नसल्यामुळे, $s$ कोणत्या चले ना
मूल्ये नेमून देते यावर मर्यादा घालता येत नाही. सोपा उपाय म्हणजे $s$ ने
\emph{सर्व} चले~$\Obj v_0$, $\Obj v_1$, \dots\@ यांना मूल्ये
नेमून द्यावीत अशी अट ठेवणे. आपल्याला लागतील तीच मूल्ये आपण वापरू.


सूत्रांची पूर्ति फक्त संरचनेच्या सापेक्ष परिभाषित
करण्याऐवजी, आपण ती संरचना~$\Struct{M}$ \emph{आणि}
चल चे मूल्यांकन~$s$ या दोन्हींच्या सापेक्ष परिभाषित करू आणि
संक्षेपाने $\Sat{M}{\varphi}[s]$ असे लिहू. चले असलेली आण्विक
सूत्रे हाताळण्यासाठी आपल्या व्याख्येत आता एक अतिरिक्त उपवाक्य
असेल:
\begin{enumerate}
  \item $\Sat{M}{\Atom{\Obj P}{\Obj a}}[s]$ तेव्हा आणि केवळ तेव्हाच,
  जेव्हा $\Assign{\Obj a}{M} \in \Assign{\Obj P}{M}$.
  \item $\Sat{M}{\Atom{\Obj P}{\Obj v_i}}[s]$ तेव्हा आणि केवळ तेव्हाच,
  जेव्हा $s(\Obj v_i) \in \Assign{\Obj P}{M}$.
  \item $\Sat{M}{\lnot \varphi}[s]$ तेव्हा आणि केवळ तेव्हाच, जेव्हा
  $\Sat{M}{\varphi}[s]$ नाही.
  \item $\Sat{M}{(\varphi \land \psi)}[s]$ तेव्हा आणि केवळ तेव्हाच, जेव्हा
  $\Sat{M}{\varphi}[s]$ आणि $\Sat{M}{\psi}[s]$ दोन्ही आहेत.
\end{enumerate}
ठीक आहे, यामुळे एक अडचण सुटते: $s(\Obj v_0)$ या मूल्यासाठी
$\Struct{M}$ हे $\Atom{\Obj P}{\Obj v_0}$ ची पूर्ति केव्हा करते हे आता
सांगता येते. $\lexists[\Obj v_0][\Atom{\Obj P}{\Obj v_0}]$ साठी योग्य
व्याख्या मिळवण्यासाठी आणखी एक गोष्ट करावी लागेल:
$\Sat{M}{\lexists[\Obj v_0][\Atom{\Obj P}{\Obj v_0}]}[s]$ तेव्हा आणि
केवळ तेव्हाच असावे, जेव्हा~$\Obj v_0$ ला मूल्य नेमून देणाऱ्या
\emph{एखाद्या} पद्धती~$s'$ साठी
$\Sat{M}{\Atom{\Obj P}{\Obj v_0}}[s']$ असते. पण~$\Obj v_0$ ला नेमलेले
मूल्य $s(\Obj v_0)$ ने निर्देशित केलेले मूल्यच असणे आवश्यक नाही.
त्यासाठी आपण एक संकेतपद्धती आणू: $m \in \Domain{M}$ असेल, तर
$\Subst{s}{m}{\Obj v_0}$ हे~$s$ सारखेच असणारे (फक्त~$\Obj v_0$ सोडून
इतर सर्व चले साठी) पण~$\Obj v_0$ ला~$m$ नेमून देणारे
मूल्यांकन मानू. आता आपली व्याख्या अशी देता येईल:
\begin{enumerate}\setcounter{enumi}{4}
  \item $\Sat{M}{\lexists[\Obj v_i][\varphi]}[s]$ तेव्हा आणि केवळ तेव्हाच,
  जेव्हा एखाद्या~$m \in \Domain{M}$ साठी
  $\Sat{M}{\varphi}[\Subst{s}{m}{\Obj v_i}]$.
\end{enumerate}
हे योग्य रीतीने चालते का? प्रत्येक~$i \in \Nat$ साठी
$s(\Obj v_i) = 0$ ठेवू या. एखादा $m \in \Domain{M}$ असा असेल की
$\Sat{M}{\Atom{\Obj P}{\Obj v_0}}[\Subst{s}{m}{\Obj v_0}]$, तेव्हा आणि
केवळ तेव्हाच $\Sat{M}{\lexists[\Obj v_0][\Atom{\Obj P}{\Obj v_0}]}[s]$.
आणि असा घटक आहे: आपण $m = 1$ किंवा $m = 2$ निवडू शकतो. $s$ नेच
~$\Obj v_0$ ला नेमून दिलेले $s(\Obj v_0)$ हे मूल्य---या उदाहरणात,
$0$---हे काम करत नसले तरी हे सत्य आहे, हे लक्षात घ्या. आपल्याकडे
$\Sat{M}{\Atom{\Obj P}{\Obj v_0}}[\Subst{s}{1}{\Obj v_0}]$ आहे, पण
$\Sat{M}{\Atom{\Obj P}{\Obj v_0}}[s]$ नाही.

हे गोंधळात टाकणारे आणि अवजड वाटत असेल, तर ते तसे आहे. पण सर्व
सूत्रे साठी पूर्तिची काटेकोर, विगमनाधारित व्याख्या देण्यासाठी ही
अतिरिक्त गुंतागुंत आवश्यक आहे आणि चिन्हार्थविषयक संकल्पना काटेकोरपणे
परिभाषित करण्यासाठी आपल्याला अशाच काहीतरी पद्धतीची गरज आहे. हे
करण्याचे इतरही मार्ग आहेत, पण ते सर्व तितकेच (अन)सुबक आहेत.












\section{वाक्य}\label{fol:int:snt:sec}

ठीक आहे, आता आपल्याकडे संरचना आणि सूत्रे यांच्यासाठी
पूर्तिच्या (``यात सत्य'') व्याख्येचा आराखडा आहे. पण तिला हा अतिरिक्त
भाग---चल चे मूल्यांकन---आवश्यक आहे आणि आपल्याला हवी होती ती
वाक्यांची व्याख्या. मूल्यांकने कशी काढून टाकायची आणि
वाक्ये म्हणजे काय?

आकारिक तर्कशास्त्राशी तुमची पहिली ओळख झाली तेव्हा चले आणि
संख्यापक यांच्यातील संबंधाची केलेली चर्चा तुम्हाला बहुधा आठवत असेल.
संख्यापकानंतर नेहमी चल येते आणि मग त्या संख्यापकाचा
वाक्य मधील ज्या भागाला उपयोग होतो (त्याची ``व्याप्ती''), त्या
भागात ते चल त्या संख्यापकाने ``बद्ध'' केले आहे असे आपण
समजतो. सूत्रे मध्ये प्रत्येक चल ला जुळणारा संख्यापक
असणे आवश्यक नव्हते आणि जुळणारे संख्यापक नसलेली चले ही
``मुक्त'' किंवा ``अबद्ध'' असतात. मुक्त चले नसलेली सर्व
सूत्रे म्हणजे वाक्ये असे आपण मानू.

सूत्र~$\varphi$ मधील चल ची एखादी आस्थिती केव्हा बद्ध असते,
तिला कोणता संख्यापक बद्ध करतो आणि ती केव्हा मुक्त असते, याची अंतःप्रज्ञ
कल्पना समजणे पुन्हा अवघड नाही. हे तपासण्याची, कदाचित कंस मोजण्याची,
एखादी पद्धत तुम्ही शिकला असाल. पण आपण काटेकोर व्याख्येचा आग्रह धरला
पाहिजे---आणि सूत्रांची व्याख्या आपण विगमनाने केली असल्यामुळे,
सूत्र~$\varphi$ मधील चल~$x$ च्या मुक्त आणि बद्ध आस्थितींची
व्याख्याही विगमनाने देता येते. उदाहरणार्थ, आपल्या सोप्या भाषेसाठी ती
अशी दिसू शकेल:
\begin{enumerate}
  \item $\varphi$ हे आण्विक असेल, तर त्यातील $x$ च्या सर्व आस्थिती मुक्त
  असतात (म्हणजे, $\Atom{\Obj P}{x}$ मधील $x$ ची आस्थिती मुक्त असते).
  \item $\varphi$ हे $\lnot \psi$ या रूपाचे असेल, तर $\lnot \psi$ मधील~$x$ ची
  आस्थिती तेव्हा आणि केवळ तेव्हाच मुक्त असते, जेव्हा~$\psi$ मधील~$x$ ची
  संबंधित आस्थिती मुक्त असते (म्हणजे,~$\psi$ मधील चरांच्या मुक्त
  आस्थिती म्हणजे नेमक्या $\lnot \psi$ मधील संबंधित आस्थिती होत).
  \item $\varphi$ हे $(\psi \land \chi)$ या रूपाचे असेल, तर $(\psi \land \chi)$
  मधील~$x$ ची आस्थिती तेव्हा आणि केवळ तेव्हाच मुक्त असते, जेव्हा~$\psi$
  किंवा~$\chi$ मधील~$x$ ची संबंधित आस्थिती मुक्त असते.
  \item $\varphi$ हे $\lexists[x][\psi]$ या रूपाचे असेल, तर $\varphi$ मधील $x$ ची
  कोणतीही आस्थिती मुक्त नसते; ते $\lexists[y][\psi]$ या रूपाचे असेल आणि
  $y$ हे~$x$ पेक्षा वेगळे चल असेल, तर
  $\lexists[y][\psi]$ मधील~$x$ ची आस्थिती तेव्हा आणि केवळ तेव्हाच मुक्त
  असते, जेव्हा~$\psi$ मधील~$x$ ची संबंधित आस्थिती मुक्त असते.
\end{enumerate}

चरांच्या मुक्त आणि बद्ध आस्थितींची काटेकोर व्याख्या मिळाल्यानंतर आपण
सरळ असे म्हणू शकतो: मुक्त चलांची आस्थिती नसलेले कोणतेही
सूत्र म्हणजे वाक्य होय.












\section{चिन्हार्थविषयक संकल्पना}\label{fol:int:sem:sec}

$\Sat{M}{\varphi}[s]$ लागू आहे का याचा विचार करताना, सोयीसाठी आपण $s$ कडून
सर्व चले ना मूल्ये नेमून घेतो, पण~$\varphi$ मधील चले ना
त्याने नेमून दिलेली मूल्येच फक्त वापरली जातात, असे आपण वर म्हटले.
खरे तर, $\varphi$ मधील \emph{मुक्त} चरांची मूल्येच महत्त्वाची असतात. आपण
काटेकोर असल्यामुळे अर्थातच हे तथ्य सिद्ध करणार आहोत. वाक्ये मध्ये
मुक्त चरे नसल्यामुळे, संरचनांमध्ये त्यांची पूर्ति होते की नाही
यासाठी $s$ चा काहीही फरक पडत नाही. म्हणून, $\varphi$ हे वाक्य असेल,
तेव्हा ``प्रत्येक~$s$ साठी $\Sat{M}{\varphi}[s]$'' असा $\Sat{M}{\varphi}$ चा अर्थ
परिभाषित करता येतो; आणि तो योगायोगाने किमान एका~$s$ साठी
$\Sat{M}{\varphi}[s]$ असते तेव्हा आणि केवळ तेव्हाच सत्य असतो. सूत्रे
साठी पूर्तिची कार्यक्षम व्याख्या मिळवण्यासाठी आपल्याला चल ची
मूल्यांकने आणावी लागतात, पण वाक्ये साठी पूर्ति ही चल
च्या मूल्यांकनांपासून स्वतंत्र असते.

``$\Sat{M}{\varphi}$'' ची व्याख्या मिळाल्यानंतर प्रथम-क्रम तर्कशास्त्रातील
वाक्यांच्या संदर्भात ``प्रसंग'' आणि ``यात सत्य'' यांचा अर्थ काय
हे आपल्याला माहीत होते. मग वाक्ये साठीच्या $\Sat{M}{\varphi}$ च्या
व्याख्येच्या आधारे वैधता, तार्किक निष्पन्नता आणि पूर्ततायोग्यता या मूलभूत
चिन्हार्थविषयक संकल्पना परिभाषित करता येतात. प्रत्येक संरचना
एखाद्या वाक्याची पूर्ति करत असेल, तर ते वाक्य वैध, $\Entails \varphi$, असते.
~$\Gamma$ मधील सर्व वाक्यांची पूर्ति करणारे प्रत्येक
संरचना हे~$\varphi$ चीही पूर्ति करत असेल, तर वाक्यांच्या संचातून
ते तार्किकरीत्या निष्पन्न होते, $\Gamma \Entails \varphi$. आणि एखादे
संरचना एखाद्या वाक्यांच्या संचातील सर्व वाक्यांची एकाच
वेळी पूर्ति करत असेल, तर तो संच पूर्ततायोग्य असतो.

सूत्रे विगमनाधारित रीतीने परिभाषित केली आहेत आणि पूर्ति हीही
सूत्रांच्या संरचनेवरील विगमनाने परिभाषित केली आहे; म्हणून आपल्या
चिन्हार्थमीमांसेचे गुणधर्म सिद्ध करण्यासाठी आणि परिभाषित
चिन्हार्थविषयक संकल्पनांचा परस्परसंबंध लावण्यासाठी विगमन वापरता येते.
यांपैकी काही गुणधर्म आपण एकत्र करून सिद्ध करू; काही अंशी प्रत्येक
गुणधर्म स्वतंत्रपणे रोचक असल्यामुळे, पण मुख्यतः चिन्हार्थमीमांसा आणि
निष्पत्ती-पद्धती यांच्यातील संबंधाचा अभ्यास पुढे करताना त्यांपैकी
अनेक उपयोगी पडणार असल्यामुळे. हे करण्यासाठी आपण अजून उल्लेख न केलेल्या
काही विन्यासविषयक संकल्पना आणि क्रियाही (काटेकोरपणे, म्हणजे विगमनाने)
परिभाषित कराव्या लागतील.












\section{आदेशन}\label{fol:int:sub:sec}

सर्व गोष्टींचा परस्परसंबंध कसा आहे आणि विन्यासमीमांसा व
चिन्हार्थमीमांसा यांचा विकास पुढील अधिक प्रगत अभ्यासाचा पाया कसा घालतो,
हे दाखवण्यासाठी एका उदाहरणाची चर्चा करू. आपल्या निष्पत्ती-पद्धतींनी
आपल्याला $\lforall[\Obj v_0][\Atom{\Obj P}{\Obj v_0}]$ पासून
$\Atom{\Obj P}{\Obj a}$ निष्पन्न करू दिले पाहिजे. कदाचित हे निगमन नियम
म्हणूनही मांडायचे असेल. पण त्यासाठी ते शक्य तितक्या सर्वसाधारण शब्दांत
मांडता आले पाहिजे: फक्त $\Obj P$, $\Obj a$ आणि $\Obj v_0$ यांच्यासाठी
नव्हे, तर कोणत्याही सूत्र~$\varphi$, पद~$t$ आणि चल~$x$
यांच्यासाठी. (स्थिरांक ही पदे असतात हे आठवा; पण स्थिरांक आणि
फलने यांच्यापासून रचलेल्या अधिक जटिल पदांचाही आपण विचार करू.)
म्हणून आपल्याला असे काहीतरी म्हणता आले पाहिजे: ``जेव्हा तुम्ही
$\lforall[x][\varphi(x)]$ निष्पन्न केले असेल, तेव्हा~$\varphi(t)$ चे अनुमान
करणे समर्थित असते---$\lforall[x]$ काढून आणि~$x$ च्या जागी~$t$ ठेवून
मिळणारे हे सूत्र आहे.'' पण ``$x$ च्या जागी~$t$ ठेवणे'' म्हणजे नेमके
काय? $\varphi(x)$ आणि~$\varphi(t)$ यांच्यातील संबंध कोणता? हे नेहमी चालते का?


हे काटेकोर करण्यासाठी आपण \emph{आदेशन} ही क्रिया परिभाषित करतो.
आदेशन प्रत्यक्षात गुंतागुंतीचे आहे, कारण~$\varphi$ मधील सर्व~$x$ फक्त~$t$ ने
बदलता येत नाहीत आणि प्रत्येक~$t$ चे कोणत्याही~$x$ साठी आदेशन करता येत
नाही. आपण हे पुन्हा विगमनाधारित व्याख्या वापरून हाताळू. पण हे
केल्यानंतर ``$\lforall[x][\varphi(x)]$ पासून $\varphi(t)$ चे अनुमान करा'' असा
निगमन नियम निर्दिष्ट करणे ही काटेकोर व्याख्या बनते. शिवाय,
$\lforall[x][\varphi(x)]$ पासून~$\varphi(t)$ तार्किकरीत्या निष्पन्न होते या अर्थाने
हा चांगला निगमन नियम आहे, हे दाखवता येईल. पण हे सिद्ध करण्यासाठी प्रथम
दर्शनी ``आपण हे का करत आहोत?'' असा प्रश्न पडेल अशी आणखी एक गोष्ट पुन्हा
सिद्ध करावी लागेल. $\lforall[x][\varphi(x)]$ पासून~$\varphi(t)$ तार्किकरीत्या
निष्पन्न होते हे पुढील तथ्यावर अवलंबून आहे: $\Sat{M}{\varphi(t)}$ लागू आहे
की नाही हे फक्त पद~$t$ च्या मूल्यावर अवलंबून असते. म्हणजे, $t$ ने
निर्देशित होणारा~$\Domain{M}$ मधील कोणताही घटक आपण $m$ मानला,
तर $\Sat{M}{\varphi(t)}[s]$ तेव्हा आणि केवळ तेव्हाच, जेव्हा
$\Sat{M}{\varphi(x)}[\Subst{s}{m}{x}]$. $t$ मध्ये चले असली तरीही
हे लागू असते; पण हा परिणाम नेमका कसा मांडायचा याबाबत काळजी घ्यावी
लागेल.












\section{प्रतिमाने आणि उपपत्ती}\label{fol:int:mod:sec}

प्रथम-क्रम तर्कशास्त्राची विन्यासमीमांसा आणि चिन्हार्थमीमांसा परिभाषित
केल्यानंतर संरचनेचे गुणधर्म आणि चिन्हार्थविषयक संकल्पना यांचा
अभ्यास सुरू करता येतो. निष्पत्ती-पद्धतीही परिभाषित करून त्यांचा
अभ्यास करता येतो. वाक्यांचा एखादा संच घेतल्यावर आपण विचारू
शकतो: त्या संचातील सर्व वाक्ये कोणती संरचना सत्य करतात?
वाक्यांचा संच~$\Gamma$ दिला असता, त्यांची पूर्ति करणाऱ्या
संरचना~$\Struct{M}$ ला \emph{$\Gamma$ चे प्रतिमान} म्हणतात.
~$\Gamma$ पासून सुरुवात करून आपण त्याची प्रतिमाने शोधू शकतो---ती कशी
दिसतात? ती किती मोठी किंवा लहान असावी लागतात? पण एखादे एक
संरचना किंवा संरचनेचा संग्रह घेऊनही विचारता येते: त्यांत
कोणती वाक्ये सत्य आहेत? या संरचनांमध्ये ती आणि फक्त तीच
सत्य असतात या अर्थाने त्यांचे \emph{अभिलक्षण} करणारी वाक्ये आहेत
का? अशा प्रकारचे प्रश्न \emph{प्रतिमान उपपत्ती} या क्षेत्रातील आहेत.
त्यांवरच \emph{स्वयंसिद्धकीय पद्धत} ही आधारित असते: एखाद्या
संरचनेच्या संग्रहाचे उपपत्तीच्या स्वयंसिद्धकांचा, म्हणजे
वाक्यांच्या संचाचा, वापर करून वर्णन करणे. स्वयंसिद्धकांपासून
प्रथम-क्रम तर्कशास्त्रात तार्किकरीत्या निष्पन्न होणारी नेमकी तीच
वाक्ये स्वयंसिद्धकांच्या सर्व प्रतिमानांत सत्य असतात, या
निरीक्षणामुळे हे शक्य होते.

अतिशय सोपे उदाहरण म्हणून पूर्वक्रमांचा विचार करा. एखाद्या संच~$A$ वरील
परावर्ती आणि संक्रमक असा संबंध~$R$ म्हणजे पूर्वक्रम होय. द्विस्थानी
संबंध $R \subseteq A \times A$ असलेला संच~$A$ म्हणजे एकच द्विस्थानी
संबंधचिन्ह~$\Obj P$ असलेल्या प्रथम-क्रम भाषेला संरचना देण्यासाठी
नेमके जे आवश्यक आहे ते: आपण $\Domain{M} = A$ आणि
$\Assign{\Obj P}{M} = R$ ठेवू. $R$ हा पूर्वक्रम असल्यामुळे तो परावर्ती
आणि संक्रमक आहे आणि हे सांगणारा प्रथम-क्रम तर्कशास्त्रातील
वाक्यांचा संच~$\Gamma$ आपल्याला मिळू शकतो:
\begin{align*}
  & \lforall[\Obj v_0][\Atom{\Obj P}{\Obj v_0,\Obj v_0}]\\
  & \lforall[\Obj v_0][\lforall[\Obj v_1][\lforall[\Obj v_2][((\Atom{\Obj P}{\Obj v_0,\Obj v_1} \land \Atom{\Obj P}{\Obj v_1,\Obj v_2}) \lif \Atom{\Obj P}{\Obj v_0,\Obj v_2})]]]
\end{align*}
ही वाक्ये म्हणजे फक्त ``कोणत्याही~$x$ साठी $Rxx$'' ($R$ हा
परावर्ती आहे) आणि ``जेव्हा $Rxy$ आणि $Ryz$, तेव्हा $Rxz$ ही'' ($R$ हा
संक्रमक आहे) यांची चिन्हांकने आहेत. संरचना~$\Struct{M}$ हे या
दोन वाक्ये~$\Gamma$ चे प्रतिमान तेव्हा आणि केवळ तेव्हाच असते,
जेव्हा $R$ (म्हणजे $\Assign{\Obj P}{M}$) हा~$A$ (म्हणजे
$\Domain{M}$) वरील पूर्वक्रम असतो. दुसऱ्या शब्दांत, $\Gamma$ ची प्रतिमाने
म्हणजे नेमके पूर्वक्रम होत. फक्त~$\Obj P$ हे विधेय असलेल्या
प्रथम-क्रम भाषेत व्यक्त करता येणारा सर्व पूर्वक्रमांचा कोणताही गुणधर्म
(वरील परावर्तिता आणि संक्रमणता यांसारखा) हा~$\Gamma$ मधील दोन
वाक्यांपासून तार्किकरीत्या निष्पन्न होतो आणि उलटही तसेच होते.
म्हणून~$\Gamma$ च्या प्रतिमानांविषयी आपण जे काही सिद्ध करू शकतो, ते सर्व
पूर्वक्रमांविषयी सिद्ध केलेले असते.

कोणत्याही विशिष्ट उपपत्तीसाठी आणि प्रतिमानांच्या वर्गासाठी (जसे
$\Gamma$ आणि सर्व पूर्वक्रम) संबंधित प्रथम-क्रम भाषेत काय व्यक्त करता
येते आणि काय व्यक्त करता येत नाही याविषयी रोचक प्रश्न असतील. सर्व भाषा
आणि प्रतिमानवर्गांसाठी रोचक असणारे संरचनेचे काही गुणधर्म आहेत;
विशेषतः प्रांताच्या आकारमानाशी संबंधित गुणधर्म. उदाहरणार्थ,
कोणत्याही~$n \in \PosInt$ साठी प्रांत मध्ये नेमके
$n$~घटक आहेत असे नेहमी व्यक्त करता येते. अनंतपणे अनेक
वाक्यांचा संच वापरून प्रांत अनंत आहे असेही व्यक्त करता येते.
पण प्रांत सांत आहे किंवा प्रांत अगणनीय आहे असे व्यक्त करता
येत नाही. प्रथम-क्रम भाषांच्या मर्यादांविषयीचे हे परिणाम संहतता आणि
लोव्हेनहाइम--स्कोलेम प्रमेयांचे परिणाम आहेत.












\section{निर्दोषता आणि संपूर्णता}\label{fol:int:scp:sec}

प्रथम-क्रम तर्कशास्त्रासाठीच्या निष्पत्ती-पद्धतींचाही आपण परिचय करून
घेऊ. तर्कशास्त्रज्ञांनी अनेक निष्पत्ती-पद्धती विकसित केल्या आहेत,
पण त्या सर्व वाक्ये दरम्यान तोच निष्पन्नता संबंध परिभाषित
करतात. विशिष्ट आणि काटेकोरपणे परिभाषित प्रकारची निष्पत्ती असेल,
तर~$\Gamma$ हे~$\varphi$ \emph{निष्पन्न} करते, $\Gamma \Proves \varphi$, असे
आपण म्हणतो. निष्पत्त्या म्हणजे नेहमी चिन्हांच्या सांत मांडण्या
असतात---कदाचित वाक्यांची यादी किंवा त्याहून गुंतागुंतीची काही
संरचना. एखाद्या संच~$\Gamma$ पासून वाक्य तार्किकरीत्या निष्पन्न
होते का हे ठरवण्याचे साधन पुरवणे हे निष्पत्ती-पद्धतींचे उद्दिष्ट
आहे. ते उद्दिष्ट साध्य करण्यासाठी $\Gamma \Entails \varphi$ तेव्हा आणि केवळ
तेव्हाच $\Gamma \Proves \varphi$ असले पाहिजे.

$\Gamma \Proves \varphi$ पण $\Gamma \Entails \varphi$ नसेल, तर आपली
निष्पत्ती-पद्धत अतिप्रबल असेल, म्हणजे ती गरजेपेक्षा जास्त सिद्ध
करेल. $\Gamma \Proves \varphi$ असेल तर $\Gamma \Entails \varphi$ असते या
गुणधर्माला \emph{निर्दोषता} म्हणतात आणि कोणत्याही चांगल्या
निष्पत्ती-पद्धतीसाठी ही किमान आवश्यकता आहे. दुसरीकडे,
$\Gamma \Entails \varphi$ पण $\Gamma \Proves \varphi$ नसेल, तर आपली
निष्पत्ती-पद्धत अतिदुर्बल असेल, म्हणजे ती पुरेसे सिद्ध करणार नाही.
$\Gamma \Entails \varphi$ असेल तर $\Gamma \Proves \varphi$ असते या गुणधर्माला
\emph{संपूर्णता} म्हणतात. निर्दोषता सिद्ध करणे सहसा तुलनेने सोपे असते
(विगमनाधारित रीतीने परिभाषित केलेल्या निष्पत्त्यांच्या संरचनेवर
विगमन करून). संपूर्णता सिद्ध करणे अधिक कठीण असते.

निर्दोषता आणि संपूर्णता यांचे अनेक महत्त्वाचे परिणाम आहेत.
वाक्यांचा एखादा संच~$\Gamma$ हा व्याघात (जसे
$\varphi \land \lnot \varphi$) निष्पन्न करत असेल, तर त्याला \emph{विसंगत}
म्हणतात. विसंगत~$\Gamma$ ला कोणतेही प्रतिमान असू शकत नाही; तो
अपूर्ततायोग्य असतो. संपूर्णतेवरून याचा उलट निष्कर्ष मिळतो: विसंगत
नसलेल्या---किंवा, आपण म्हणू त्याप्रमाणे, \emph{सुसंगत}---प्रत्येक
$\Gamma$ ला प्रतिमान असते. खरे तर, हे संपूर्णतेशी सममूल्य आहे आणि
संपूर्णतेचे हेच रूप आपण प्रत्यक्षात सिद्ध करू. हा सखोल आणि कदाचित
आश्चर्यकारक परिणाम आहे: $\Gamma$ पासून $\varphi \land \lnot \varphi$ सिद्ध करता
येत नाही एवढ्यामुळेच~$\Gamma$ वर्णन करते तसे संरचना अस्तित्वात
असण्याची हमी मिळते. म्हणून संपूर्णता पुढील प्रश्नाचे उत्तर देते:
वाक्यांच्या कोणत्या संचांना प्रतिमाने असतात? उत्तर: सर्व आणि फक्त
सुसंगत संचांना.

निर्दोषता आणि संपूर्णता प्रमेयांचे दोन महत्त्वाचे परिणाम आहेत: संहतता
प्रमेय आणि लोव्हेनहाइम--स्कोलेम प्रमेय. प्रतिमानांच्या उपपत्तीतील हे
महत्त्वाचे परिणाम आहेत आणि अनेक रोचक निष्कर्ष प्रस्थापित करण्यासाठी
त्यांचा उपयोग करता येतो. आपण आधीच दोन निष्कर्षांचा उल्लेख केला आहे:
प्रथम-क्रम तर्कशास्त्रात संरचनेचा प्रांत सांत आहे किंवा तो
अगणनीय आहे असे व्यक्त करता येत नाही.

ऐतिहासिकदृष्ट्या हे सर्व---प्रथम-क्रम तर्कशास्त्राची विन्यासमीमांसा आणि
चिन्हार्थमीमांसा कशी परिभाषित करायची, चांगल्या निष्पत्ती-पद्धती
कशा परिभाषित करायच्या, त्या निर्दोष आणि संपूर्ण आहेत हे कसे सिद्ध
करायचे, प्रथम-क्रम भाषांत काय व्यक्त करता येते आणि काय व्यक्त करता येत
नाही हे स्पष्ट करणे---शोधून योग्य रीतीने ठरवण्यास बराच काळ लागला. हे
कसे करायचे हे आता माहीत असले, तरी सर्व तपशील समजून घेणे अजूनही
गोंधळात टाकणारे आणि कंटाळवाणे ठरू शकते. पण ते महत्त्वाचेही आहे, कारण
प्रथम-क्रम तर्कशास्त्राच्या आकारिक भाषेसाठी येथे विकसित केलेल्या पद्धती
तर्कशास्त्र, संगणकशास्त्र आणि भाषाविज्ञान यांत सर्वत्र वापरल्या जातात.
म्हणून सर्व तपशीलांचा अभ्यास करण्याचे दीर्घकाळात फळ मिळते.



\chapter{प्रथम-क्रम तर्कशास्त्राची विन्यासमीमांसा}\label{fol:syn::chap}









\section{परिचय}\label{fol:syn:itx:sec}

प्रथम-क्रम तर्कशास्त्राची उपपत्ती आणि अधिउपपत्ती विकसित करण्यासाठी,
आपण प्रथम त्याच्या व्यंजकांची विन्यासमीमांसा आणि चिन्हार्थमीमांसा
परिभाषित केली पाहिजे. प्रथम-क्रम तर्कशास्त्राची व्यंजके म्हणजे पदे आणि
सूत्रे. पदे चले, स्थिरांक आणि फलने
यांपासून तयार होतात. सूत्रे मात्र विधेये आणि पदे
यांपासून तयार होतात (यांतून सर्वांत लहान, ``आण्विक'' सूत्रे
तयार होतात); मग तार्किक संयोजके आणि संख्यापक वापरून आण्विक
सूत्रांपासून अधिक जटिल सूत्रे तयार करता येतात. रचनानियम
मांडण्याचे अनेक निरनिराळे मार्ग आहेत; आपण त्यांपैकी केवळ एक शक्य
मार्ग देतो. इतर पद्धती वेगळी चिन्हे निवडतील, संयोजकांचे वेगळे संच
आदिम म्हणून निवडतील आणि कंसांचा वेगळ्या रीतीने वापर करतील (किंवा
तथाकथित पोलिश संकेतपद्धतीप्रमाणे कंस अजिबात वापरणार नाहीत). तरी सर्व
दृष्टिकोनांत सामाईक बाब ही की रचनानियम पदांचा आणि सूत्रांचा
संच \emph{विगमनाने} परिभाषित करतात. हे योग्य रीतीने केल्यास,
रचनानियमांनुसार प्रत्येक व्यंजक मूलतः एकाच रीतीने तयार होऊ शकते.
त्यामुळे तयार होणारी व्यंजके \emph{एकमेव रीतीने वाचनीय} असतात. परिणामी,
त्यांची ही विगमनाधारित व्याख्या आपल्याला त्याच पद्धतीने---विगमनाधारित
व्याख्येने---त्या व्यंजकांना अर्थ देता येतो, याची खात्री देते.












\section{प्रथम-क्रम भाषा}\label{fol:syn:fol:sec}


प्रथम-क्रम तर्कशास्त्राची व्यंजके अशा मूलभूत शब्दसंग्रहापासून तयार
केली जातात ज्यात \emph{चले}, \emph{स्थिरांक},
\emph{विधेये} आणि कधीकधी \emph{फलने} असतात. त्यांपासून,
तार्किक संयोजके, संख्यापक आणि कंस व स्वल्पविराम यांसारखी विरामचिन्हे
यांच्या साहाय्याने \emph{पदे} आणि \emph{सूत्रे} तयार केली जातात.

\begin{explain}
अनौपचारिकपणे, विधेये ही गुणधर्मांची आणि संबंधांची नावे,
स्थिरांक ही स्वतंत्र वस्तूंची नावे आणि फलने ही
संगतींची नावे आहेत. एकरूपता~$\eq$ सोडल्यास हीच \emph{अतार्किक
चिन्हे} होत आणि ती मिळून एक भाषा बनते. कोणतीही प्रथम-क्रम
भाषा~$\Lang L$ तिच्या अतार्किक चिन्हांनी निर्धारित होते. सर्वांत
सामान्य बाबतीत $\Lang L$ मध्ये प्रत्येक प्रकारची अनंत चिन्हे असतात.
\end{explain}

सामान्य बाबतीत प्रथम-क्रम तर्कशास्त्रात आपण पुढील चिन्हे वापरतो:

\begin{enumerate}
\item तार्किक चिन्हे
\begin{enumerate}
\item तार्किक संयोजके:
  \startycommalist
  \ycomma $\lnot$ (नकरण)
  \ycomma $\land$ (संधियोग)
  \ycomma $\lor$ (विकल्पयोग)
  \ycomma $\lif$ (शर्त)
  \ycomma $\liff$ (द्विशर्त)
  \ycomma $\lforall$ (वैश्विक संख्यापक)
  \ycomma $\lexists$ (अस्तित्ववाची संख्यापक).
\item असत्यता~$\lfalse$ साठीचे विधानीय स्थिर.
\item सत्यता~$\ltrue$ साठीचे विधानीय स्थिर.
\item द्विस्थानी एकरूपता~$\eq$.
\item चलांचा एक गणनीय अनंत संच: $\Obj v_0$, $\Obj v_1$, $\Obj
  v_2$, \dots
\end{enumerate}
\item प्रथम-क्रम तर्कशास्त्राची \emph{मानक भाषा} बनवणारी अतार्किक चिन्हे
\begin{enumerate}
\item प्रत्येक $n>0$ साठी $n$-स्थानी विधेयांचा एक गणनीय अनंत संच: $\Obj
  A^n_0$, $\Obj A^n_1$, $\Obj A^n_2$, \dots
\item स्थिरांकाचा एक गणनीय अनंत संच: $\Obj c_0$, $\Obj c_1$, $\Obj
  c_2$, \dots.
\item प्रत्येक $n>0$ साठी $n$-स्थानी फलनांचा एक गणनीय अनंत संच:
  $\Obj f^n_0$, $\Obj f^n_1$, $\Obj f^n_2$, \dots
\end{enumerate}
\item विरामचिन्हे: (, ) आणि स्वल्पविराम.
\end{enumerate}

आपल्या बहुतेक व्याख्या आणि निष्कर्ष प्रथम-क्रम तर्कशास्त्राच्या संपूर्ण
मानक भाषेसाठी मांडले जातील. मात्र, उपयोगानुसार आपण भाषा केवळ काही
विधेये, स्थिरांक आणि फलने पुरती मर्यादितही करू शकतो.

\begin{ex}
अंकगणिताची भाषा~$\Lang L_A$ हिच्यात एक द्विस्थानी विधेय~$<$,
एक स्थिरांक~$\Obj 0$, एक एकस्थानी फलन~$\prime$ आणि दोन
द्विस्थानी फलने~$+$ व~$\times$ असतात.
\end{ex}

\begin{ex}
संचसिद्धान्ताची भाषा~$\Lang L_Z$ हिच्यात केवळ एक द्विस्थानी
विधेय~$\in$ असते.
\end{ex}

\begin{ex}
क्रमांची भाषा~$\Lang L_\le$ हिच्यात केवळ द्विस्थानी
विधेय~$\le$ असते.
\end{ex}

पुन्हा, हे संकेत आहेत: औपचारिकदृष्ट्या ही केवळ पर्यायी नावे आहेत;
उदा., $<$, $\in$ आणि $\le$ ही $\Obj A^2_0$ ची, $\Obj 0$ हे
$\Obj c_0$ चे, $\prime$ हे $\Obj f^1_0$ चे, $+$ हे $\Obj f^2_0$ चे आणि $\times$ हे
$\Obj f^2_1$ चे पर्यायी नाव आहे.

.





\begin{intro}
वर आपण वापरलेल्या संज्ञांपेक्षा आणि चिन्हांपेक्षा वेगळ्या संज्ञा व
चिन्हे तुम्हाला परिचित असू शकतात. तर्कशास्त्राच्या ग्रंथांत (आणि
अध्यापनात) ``नकरण'' यासाठी सामान्यतः $\sim$, $\neg$ किंवा~!\ यांपैकी
एखादे, तर ``संधियोग'' यासाठी $\wedge$, $\cdot$ किंवा $\&$ यांपैकी
एखादे चिन्ह वापरतात. ``सशर्त (संकेतार्थक)'' किंवा ``अभिव्यंजन'' यांसाठी
$\rightarrow$, $\Rightarrow$ आणि $\supset$ ही चिन्हे सामान्यतः वापरली जातात.
``द्विसशर्त,'' ``द्वि-अभिव्यंजन'' किंवा
  ``(वास्तविक) सममूल्यता'' यांसाठीची चिन्हे $\leftrightarrow$,
  $\Leftrightarrow$ आणि $\equiv$ ही आहेत.
$\lfalse$ या चिन्हाला निरनिराळ्या ठिकाणी
  ``असत्यता,'' ``फाल्सम,'' ``विसंगती'' किंवा ``तळ'' म्हणतात.
$\ltrue$ या चिन्हाला निरनिराळ्या ठिकाणी
  ``सत्यता,'' ``वेरम'' किंवा ``शिखर'' म्हणतात.

स्थिरांक (ज्यांना कधीकधी नावे म्हणतात) यांच्यासाठी लॅटिन
वर्णमालेच्या सुरुवातीची लघु अक्षरे (उदा., $a$, $b$, $c$) आणि
चले साठी तिच्या शेवटची लघु अक्षरे (उदा., $x$, $y$, $z$)
वापरण्याचा संकेत आहे. संख्यापक चले सोबत जोडले जातात, उदा.,
$x$; वैश्विक संख्यापकासाठी $\forall x$, $(\forall x)$, $(x)$, $\Pi x$,
$\bigwedge_x$ आणि अस्तित्ववाची संख्यापकासाठी $\exists x$, $(\exists
x)$, $(Ex)$, $\Sigma x$, $\bigvee_x$ ही संकेतलेखनातील रूपांतरे आहेत.
\end{intro}

\begin{explain}
आपण सर्व विधानीय कारक आणि दोन्ही संख्यापक यांना भाषेची आदिम चिन्हे
मानू शकतो. त्याऐवजी आदिम चिन्हांचा लहान संच निवडून उरलेले
संकारक परिभाषित मानू शकतो. बूलीय कारकांच्या
``सत्यता-फलनात्मकदृष्ट्या संपूर्ण'' संचांत $\{ \lnot, \lor \}$,
$\{ \lnot, \land \}$ आणि $\{ \lnot,
\lif\}$ यांचा समावेश होतो---यांपैकी प्रत्येक संच कोणत्याही एका
संख्यापकासोबत जोडल्यास अभिव्यक्तीदृष्ट्या संपूर्ण प्रथम-क्रम भाषा मिळते.

इतर दोन संकारक तुम्हाला परिचित असू शकतात: शेफर स्ट्रोक~$|$
(हे नाव हेन्री शेफर यांच्यावरून दिले आहे) आणि पिअर्स बाण~$\downarrow$,
ज्याला क्वाइनचा डॅगर असेही म्हणतात. अनुक्रमे ``नँड'' आणि ``नॉर'' हे
त्यांचे नेहमीचे वाचन दिल्यास, हे कारक स्वतःच सत्यता-फलनात्मकदृष्ट्या
संपूर्ण असतात.
\end{explain}












\section{पदे आणि सूत्र}\label{fol:syn:frm:sec}

प्रथम-क्रम भाषा~$\Lang L$ दिली की, $\Lang L$ च्या मूलभूत
शब्दसंग्रहापासून तयार होणारी व्यंजके आपण परिभाषित करू शकतो. त्यांत
विशेषतः \emph{पदे} आणि \emph{सूत्रे} यांचा समावेश होतो.

\begin{defn}[पदे]
\label{fol:syn:frm:defn:terms}
$\Lang L$ मधील \emph{पदांचा} संच~$\Trm[L]$ पुढीलप्रमाणे विगमनाने
परिभाषित केला जातो:
\begin{enumerate}
\item प्रत्येक चल हे पद आहे.
\item $\Lang L$ मधील प्रत्येक स्थिरांक हे पद आहे.
\item $f$ हे $n$-स्थानी फलन आणि $t_1$, \dots, $t_n$ ही
  पदे असतील, तर $\Atom{f}{t_1, \ldots, t_n}$ हे पद आहे.
\item
  इतर काहीही पद नाही.
\end{enumerate}
ज्या पदात एकही चले नाही ते \emph{बंद पद} होय.
\end{defn}

\begin{explain}
भाषेचे विनिर्देशन करताना स्थिरांक ही पदांपासून वेगळी चिन्हांची
श्रेणी म्हणून दिसतात; पण त्याऐवजी ती शून्यस्थानी फलने मध्ये
समाविष्ट करता आली असती. मग पदांच्या व्याख्येतील दुसरी बाब आवश्यक
राहिली नसती. $n = 0$ असेल तेव्हा $\Atom{f}{t_1, \ldots, t_n}$ याचा
अर्थ केवळ स्वतंत्र $f$ असा घ्यावा लागेल.
\end{explain}

\begin{defn}[सूत्रे]
\label{fol:syn:frm:defn:formulas}
$\Lang L$ या भाषेतील \emph{सूत्रे} चा संच~$\Frm[L]$
पुढीलप्रमाणे विगमनाने परिभाषित केला जातो:
\begin{enumerate}
\item $\lfalse$ हे आण्विक सूत्र आहे.

\item $\ltrue$ हे आण्विक सूत्र आहे.

\item $R$ हे $\Lang L$ मधील $n$-स्थानी विधेय आणि $t_1$, \dots,
  $t_n$ ही $\Lang L$ मधील पदे असतील, तर $\Atom{R}{t_1,\ldots, t_n}$ हे
  आण्विक सूत्र आहे.

\item $t_1$ आणि $t_2$ ही $\Lang L$ मधील पदे असतील, तर $\Atom{\eq}{t_1, t_2}$
  हे आण्विक सूत्र आहे.

\item $\varphi$ हे सूत्र असेल, तर $\lnot \varphi$ हे
  सूत्र आहे.

\item $\varphi$ आणि $\psi$ ही सूत्रे असतील, तर $(\varphi \land
  \psi)$ हे सूत्र आहे.

\item $\varphi$ आणि $\psi$ ही सूत्रे असतील, तर $(\varphi \lor \psi)$
  हे सूत्र आहे.

\item $\varphi$ आणि $\psi$ ही सूत्रे असतील, तर $(\varphi \lif \psi)$
  हे सूत्र आहे.

\item $\varphi$ आणि $\psi$ ही सूत्रे असतील, तर $(\varphi \liff \psi)$
  हे सूत्र आहे.

\item $\varphi$ हे सूत्र आणि $x$ हे चल असेल,
  तर $\lforall[x][\varphi]$ हे सूत्र आहे.

\item $\varphi$ हे सूत्र आणि $x$ हे चल असेल,
  तर $\lexists[x][\varphi]$ हे सूत्र आहे.

\item इतर काहीही सूत्र नाही.
\end{enumerate}
\end{defn}

\begin{explain}
पदांच्या संचाची आणि सूत्रांच्या संचाची व्याख्या या
\emph{विगमनाधारित व्याख्या} आहेत. मूलतः, आपण सूत्रांचा संच
अनंतपणे अनेक टप्प्यांत तयार करतो. आरंभीच्या टप्प्यात सर्व आण्विक सूत्रे
ही सूत्रे आहेत असे घोषित करतो; हे व्याख्येतील पहिल्या काही बाबींशी,
म्हणजे
$\ltrue$,
$\lfalse$,
$\Atom{R}{t_1,\dots,t_n}$ आणि $\Atom{\eq}{t_1,t_2}$ यांच्या बाबींशी जुळते.
म्हणून ``आण्विक सूत्र'' म्हणजे या रूपांपैकी कोणतेही सूत्र.

व्याख्येतील इतर बाबी आधीच तयार केलेल्या सूत्रांपासून नवी
सूत्रे तयार करण्याचे नियम देतात. दुसऱ्या टप्प्यात या नियमांनी
आण्विक सूत्रांपासून सूत्रे तयार करता येतात. तिसऱ्या
टप्प्यात आण्विक सूत्रे आणि दुसऱ्या टप्प्यात मिळालेली सूत्रे यांपासून
नवी सूत्रे तयार करतो, आणि ही प्रक्रिया पुढे चालू राहते. अशा एखाद्या
टप्प्यात शेवटी तयार होणारी प्रत्येक गोष्ट सूत्र असते; इतर
काहीही सूत्र नसते.
\end{explain}

संकेताप्रमाणे, आपण $\eq$ हे चिन्ह त्याच्या युक्तिवादांच्या मध्ये लिहून
कंस वगळतो: $\eq[t_1][t_2]$ हा $\Atom{\eq}{t_1,t_2}$ चा संक्षेप आहे.
तसेच, $\lnot \Atom{\eq}{t_1,t_2}$ चा संक्षेप $\eq/[t_1][t_2]$ असा
लिहितो. द्विस्थानी संयोजक~$\ast$ वापरून $\psi$, $\chi$ पासून तयार केलेले
सूत्र $(\psi \ast \chi)$ लिहिताना आपण अनेकदा सर्वांत बाहेरची कंसांची
जोडी वगळून केवळ~$\psi \ast \chi$ असे लिहू.

\begin{intro}
काही तर्कशास्त्राच्या ग्रंथांत
$\lexists[x][\varphi]$
आणि
$\lforall[x][\varphi]$
यांची गणना
सूत्रे म्हणून करण्यासाठी
चल~$x$ हे~$\varphi$ मध्ये आलेले असणे आवश्यक मानतात. ही अट न
घातल्याने कोणतीही अडचण येत नाही आणि काम अधिक सोपे होते.
\end{intro}



विशिष्ट उपयोगासाठीच्या भाषेत काम करताना आपण अनेकदा द्विस्थानी
विधेये आणि फलने त्यांच्या संबंधित पदांच्या मध्ये
लिहितो; उदा., अंकगणिताच्या भाषेत $t_1 < t_2$ आणि $(t_1 + t_2)$, तर
संचसिद्धान्ताच्या भाषेत $t_1 \in t_2$. अंकगणिताच्या भाषेतील उत्तराधिकारी
फलन तर संकेताप्रमाणे त्याच्या युक्तिवादाच्या \emph{नंतर} लिहिले
जाते:~$t'$. मात्र, औपचारिकदृष्ट्या हे अनुक्रमे
$\Atom{\Obj A^2_0}{t_1, t_2}$, $\Obj f^2_0(t_1, t_2)$,
$\Atom{\Obj A^2_0}{t_1, t_2}$ आणि $f^1_0(t)$ यांचे केवळ रूढ संक्षेप आहेत.

\begin{defn}[विन्यासात्मक अनन्यता]
$\ident$ हे चिन्ह चिन्हमालांमधील विन्यासात्मक अनन्यता व्यक्त करते;
म्हणजे $\varphi \ident \psi$ तेव्हा आणि केवळ तेव्हाच, जेव्हा $\varphi$ आणि $\psi$
या समान लांबीच्या चिन्हमाला असतात आणि प्रत्येक स्थानी त्यांच्यात तेच
चिन्ह असते.
\end{defn}

$\ident$ या चिन्हाच्या दोन्ही बाजूंना जोडणीने मिळालेल्या चिन्हमाला
येऊ शकतात. उदाहरणार्थ, $\varphi \ident (\psi \lor \chi)$ याचा अर्थ असा:
आरंभीचा कंस, चिन्हमाला~$\psi$, $\lor$ हे चिन्ह, चिन्हमाला~$\chi$ आणि
शेवटचा कंस या क्रमाने जोडून मिळणारी चिन्हमाला आणि चिन्हमाला~$\varphi$
या एकच आहेत. असे असेल, तर $\varphi$ चे पहिले चिन्ह आरंभीचा कंस आहे;
$\varphi$ मध्ये दुसऱ्या चिन्हापासून सुरू होणारी $\psi$ ही उपचिन्हमाला आहे;
तिच्यानंतर $\lor$ येते; इत्यादी बाबी आपल्याला माहीत होतात.

पदे आणि सूत्रे ही मूलभूत घटकांपासून विगमनाधारित व्याख्यांनी तयार
होत असल्यामुळे, त्यांच्याविषयीच्या गोष्टी सिद्ध करण्यासाठी पुढील विगमन
तत्त्वे वापरता येतात.

\begin{lem}[\emph{पदांवरील विगमनाचे तत्त्व}]
    \label{fol:syn:frm:lem:trmind}
    $\Lang L$ ही प्रथम-क्रम भाषा असू द्या.
    एखादा गुणधर्म~$P$ असा असेल की

    \begin{enumerate}
        \item तो प्रत्येक चल~$v$ ला लागू होतो,

        \item तो $\Lang L$ मधील प्रत्येक स्थिरांक~$a$ ला लागू होतो, आणि

        \item तो $t_1$,~\dots, $t_n$ यांना लागू होत असेल आणि $f$ हे
        $\Lang L$ मधील $n$-स्थानी फलन असेल, तर तो
        $f(t_1,\dotsc,t_n)$ लाही लागू होतो
    \end{enumerate}
    (येथे $t_1$,~\dots, $t_n$ ही $\Lang{L}$ मधील पदे आहेत असे गृहीत धरले आहे),
    तर $P$ हा $\Trm[L]$ मधील प्रत्येक पदाला लागू होतो.
\end{lem}

\begin{prob}
    \ref{fol:syn:frm:lem:trmind} सिद्ध करा.
\end{prob}

\begin{lem}[\emph{सूत्रे वरील विगमनाचे तत्त्व}]
    \label{fol:syn:frm:thm:frmind}
    $\Lang L$ ही प्रथम-क्रम भाषा असू द्या.
    एखादा गुणधर्म~$P$ सर्व आण्विक सूत्रे ला लागू होत असेल आणि तो असा असेल की

    \begin{enumerate}
        \item तो $\varphi$ ला लागू होत असेल, तर तो $\lnot \varphi$ ला
            लागू होतो;
        \item तो $\varphi$ आणि~$\psi$ ला लागू होत असेल, तर तो
            $(\varphi \land \psi)$ ला लागू होतो;
        \item तो $\varphi$ आणि~$\psi$ ला लागू होत असेल, तर तो
            $(\varphi \lor \psi)$ ला लागू होतो;
        \item तो $\varphi$ आणि~$\psi$ ला लागू होत असेल, तर तो
            $(\varphi \lif \psi)$ ला लागू होतो;
        \item तो $\varphi$ आणि~$\psi$ ला लागू होत असेल, तर तो
            $(\varphi \liff \psi)$ ला लागू होतो;
        \item तो $\varphi$ ला लागू होत असेल, तर तो
            $\lexists[x][\varphi]$ ला लागू होतो;
        \item तो $\varphi$ ला लागू होत असेल, तर तो
            $\lforall[x][\varphi]$ ला लागू होतो;
  \end{enumerate}
  (येथे $\varphi$ आणि $\psi$ ही $\Lang{L}$ मधील सूत्रे आहेत असे गृहीत धरले आहे),
  तर $P$ हा $\Frm[L]$ मधील सर्व सूत्रांना लागू होतो.
\end{lem}












\section{एकमेव वाचनीयता}\label{fol:syn:unq:sec}

\begin{explain}
आपण सूत्रांची ज्या रीतीने व्याख्या केली आहे त्यामुळे प्रत्येक
सूत्राचे \emph{एकमेव वाचन} असते; म्हणजे सूत्रे साठीच्या
आपल्या रचनानियमांनुसार ते तयार करण्याचा मूलतः एकच मार्ग आणि त्याची
``अर्थलावणी'' करण्याचा एकच मार्ग असतो. तसे नसते, तर आपल्याला संदिग्ध
सूत्रे मिळाली असती, म्हणजे एकाहून अधिक वाचने किंवा अर्थलावण्या
असलेली सूत्रे---आणि हे आपल्याला टाळायचे आहे हे स्पष्ट आहे. पण
त्याहून महत्त्वाचे म्हणजे, या गुणधर्माविना आपण पुढे देणार असलेल्या
बहुतेक व्याख्या आणि सिद्धता चालणार नाहीत.

एकमेव वाचनीयतेची खात्री न देणारे सूत्रे तयार करण्याचे सदोष नियम
दिले असते तर काय झाले असते हे पाहणे, हा मुद्दा स्पष्ट करण्याचा बहुधा
सर्वोत्तम मार्ग आहे. उदाहरणार्थ, संयोजकांच्या रचनानियमांत आपण कंस
विसरलो असतो; उदा., पुढील गोष्टीला परवानगी दिली असती:
\begin{quote}
$\varphi$ आणि $\psi$ ही सूत्रे असतील, तर $\varphi \lif \psi$ हेही तसेच आहे.
\end{quote}
$\theta$ या आण्विक सूत्रापासून सुरुवात करून यामुळे $\theta \lif \theta$ तयार करता
आले असते. यापासून, $\theta$ सोबत, $\theta \lif \theta \lif \theta$ मिळाले असते. पण
हे करण्याचे दोन मार्ग आहेत:
\begin{enumerate}
\item $\theta$ हे $\varphi$ आणि $\theta \lif \theta$ हे $\psi$ मानणे.
\item $\varphi$ हे $\theta \lif \theta$ आणि $\psi$ हे~$\theta$ मानणे.
\end{enumerate}
त्याप्रमाणे, सूत्र~$\theta \lif \theta \lif \theta$ ``वाचण्याचे'' दोन मार्ग
आहेत. ते $\psi \lif \chi$ या रूपात असून $\psi$ हे $\theta$ आणि $\chi$ हे $\theta
\lif \theta$ आहे; पण ते $\psi \lif \chi$ या रूपात \emph{देखील आहे}, जेथे $\psi$
हे $\theta \lif \theta$ आणि $\chi$ हे~$\theta$ आहे.

असे झाले, तर आपल्या व्याख्या नेहमी चालणार नाहीत. उदाहरणार्थ, एखाद्या
सूत्राचा मुख्य संकारक परिभाषित करताना आपण म्हणतो: $\psi \lif \chi$
या रूपातील सूत्रात~$\lif$ ची दर्शवलेली आस्थिति हाच मुख्य संकारक
आहे. पण $\theta \lif \theta \lif \theta$ या सूत्राची वर सांगितलेल्या दोन वेगळ्या
मार्गांनी $\psi \lif \chi$ शी जुळणी करता आली, तर एका बाबतीत $\lif$ ची
पहिली आस्थिति मुख्य संकारक मिळते आणि दुसऱ्या बाबतीत दुसरी
आस्थिति मिळते. पण मुख्य संकारक हे सूत्राचे \emph{फलन} असावे
असा आपला उद्देश आहे; म्हणजे प्रत्येक सूत्रामध्ये मुख्य संकारक ची नेमकी एक आस्थिति असली पाहिजे.
\end{explain}

\begin{lem}
सूत्र~$\varphi$ मधील डाव्या आणि उजव्या कंसांची संख्या समान असते.
\end{lem}

\begin{proof}
$\varphi$ ज्या रीतीने तयार होते तिच्यावर विगमन करून हे सिद्ध करू. यासाठी
दोन गोष्टी आवश्यक आहेत: (a) प्रथम, सर्व आण्विक सूत्रांना संबंधित
गुणधर्म लागू होतो हे सिद्ध करावे लागते (विगमनाचा आधार). (b) नंतर,
दिलेल्या सूत्रांपासून नवी सूत्रे तयार केली आणि जुन्या सूत्रांना तो
गुणधर्म लागू होत असेल, तर नव्या सूत्रांनाही तो लागू होतो हे सिद्ध
करावे लागते.

$l(\varphi)$ ही $\varphi$ मधील डाव्या कंसांची संख्या आणि $r(\varphi)$ ही उजव्या
कंसांची संख्या असू द्या; त्याचप्रमाणे $l(t)$ आणि $r(t)$ या पद~$t$
मधील अनुक्रमे डाव्या आणि उजव्या कंसांच्या संख्या असू द्या.

\begin{prob}
    कोणत्याही पद~$t$ साठी $l(t) = r(t)$ हे सिद्ध करा.
\end{prob}

\begin{enumerate}
\item \indcase{\varphi}{\lfalse}{$\indfrm$ मध्ये $0$ डावे आणि $0$
    उजवे कंस आहेत.}

\item \indcase{\varphi}{\ltrue}{$\indfrm$ मध्ये $0$ डावे आणि $0$
    उजवे कंस आहेत.}

\item \indcase{\varphi}{\Atom{R}{t_1,\dots,t_n}}{$l(\indfrm) = 1 + l(t_1) +
  \dots + l(t_n) = 1 + r(t_1) + \dots + r(t_n) = r(\indfrm)$. येथे
  कोणत्याही पद~$t$ साठी $l(t) = r(t)$ ही, सरावासाठी सोडलेली, बाब
  वापरली आहे.}

\item \indcase{\varphi}{\eq[t_1][t_2]}{$l(\indfrm) = l(t_1) + l(t_2) =
  r(t_1) + r(t_2) = r(\indfrm)$.}

\item \indcase{\varphi}{\lnot \psi}{विगमन गृहीतकानुसार,
    $l(\psi) = r(\psi)$. म्हणून $l(\indfrm) = l(\psi) = r(\psi) =
    r(\indfrm)$.}

\item \indcase{\varphi}{(\psi \ast \chi)}{विगमन गृहीतकानुसार, $l(\psi) =
  r(\psi)$ आणि $l(\chi) = r(\chi)$. म्हणून $l(\indfrm) = 1 + l(\psi) + l(\chi) =
  1 + r(\psi) + r(\chi) = r(\indfrm)$.}

\item \indcase{\varphi}{\lforall[x][\psi]}{विगमन
    गृहीतकानुसार, $l(\psi) = r(\psi)$. म्हणून $l(\indfrm) = l(\psi) = r(\psi) =
    r(\indfrm)$.}





\item \indcase{\varphi}{\lexists[x][\psi]}{त्याचप्रमाणे.}
\end{enumerate}
\end{proof}

\begin{defn}[उचित पूर्वखंड]
$\psi$ ही चिन्हमाला आणि एक रिकामी नसलेली चिन्हमाला यांची जोडणी केल्यावर
चिन्हमाला~$\varphi$ मिळत असेल, तर चिन्हमाला $\psi$ ही चिन्हमाला~$\varphi$ चा
\emph{उचित पूर्वखंड} आहे.
\end{defn}

\begin{lem}\label{fol:syn:unq:lem:no-prefix}
$\varphi$ हे सूत्र आणि $\psi$ हा $\varphi$ चा उचित पूर्वखंड असेल, तर $\psi$
हे सूत्र नाही.
\end{lem}

\begin{proof}
सराव.
\end{proof}

\begin{prob}
\ref{fol:syn:unq:lem:no-prefix} सिद्ध करा.
\end{prob}

\begin{prop}
\label{fol:syn:unq:prop:unique-atomic}
$\varphi$ हे आण्विक सूत्र असेल, तर पुढीलपैकी एक आणि फक्त एकच अट त्याला
लागू होते.
\begin{enumerate}
\item $\varphi \ident \lfalse$.
\item $\varphi \ident \ltrue$.
\item $\varphi \ident \Atom{R}{t_1,\dots,t_n}$, जेथे $R$ हे $n$-स्थानी
  विधेय, $t_1$, \dots, $t_n$ ही पदे आणि $R$, $t_1$, \dots,
  $t_n$ यांपैकी प्रत्येक एकमेव रीतीने निर्धारित आहे.
\item $\varphi \ident \eq[t_1][t_2]$, जेथे $t_1$ आणि $t_2$ ही एकमेव रीतीने
  निर्धारित पदे आहेत.
\end{enumerate}
\end{prop}

\begin{proof}
सराव.
\end{proof}

\begin{prob}
\ref{fol:syn:unq:prop:unique-atomic} सिद्ध करा (सूचना: पदांसाठी
\ref{fol:syn:unq:lem:no-prefix} ची एक आवृत्ती मांडा आणि सिद्ध करा.)
\end{prob}

\begin{prop}[एकमेव वाचनीयता]
प्रत्येक सूत्राला पुढीलपैकी एक आणि फक्त एकच अट लागू होते.
\begin{enumerate}
\item $\varphi$ हे आण्विक आहे.

\item $\varphi$ हे $\lnot \psi$ या रूपात आहे.

\item $\varphi$ हे $(\psi \land \chi)$ या रूपात आहे.

\item $\varphi$ हे $(\psi \lor \chi)$ या रूपात आहे.

\item $\varphi$ हे $(\psi \lif \chi)$ या रूपात आहे.

\item $\varphi$ हे $(\psi \liff \chi)$ या रूपात आहे.

\item $\varphi$ हे $\lforall[x][\psi]$ या रूपात आहे.

\item $\varphi$ हे $\lexists[x][\psi]$ या रूपात आहे.
\end{enumerate}
शिवाय, प्रत्येक बाबतीत $\psi$, किंवा $\psi$ आणि $\chi$, एकमेव रीतीने
निर्धारित असतात. म्हणजे, उदाहरणार्थ, $\psi$, $\chi$ आणि $\psi'$, $\chi'$ अशा
दोन वेगळ्या जोड्या अस्तित्वात नाहीत की $\varphi$ हे एकाच वेळी
$(\psi \lif \chi)$ आणि $(\psi' \lif \chi')$ या दोन्ही रूपांत असेल.
\end{prop}

\begin{proof}
रचनानियमांनुसार, सूत्र आण्विक नसेल, तर त्याची सुरुवात आरंभीच्या
कंसाने~(, $\lnot$ ने, किंवा एखाद्या संख्यापकाने झाली
पाहिजे. दुसरीकडे, पुढीलपैकी एखाद्या चिन्हाने सुरू होणारे प्रत्येक
सूत्र आण्विक असले पाहिजे: विधेय, फलन,
स्थिरांक, $\lfalse$, $\ltrue$.

म्हणून आपल्याला प्रत्यक्षात एवढेच दाखवायचे आहे: $\varphi$ हे $(\psi \ast
\chi)$ या रूपात आणि $(\psi' \mathbin{\ast'} \chi')$ या रूपातही असेल, तर
$\psi \ident \psi'$, $\chi \ident \chi'$ आणि $\ast = {\ast'}$.

म्हणून $\varphi \ident (\psi \ast \chi)$ आणि $\varphi \ident (\psi'
\mathbin{\ast'} \chi')$ हे दोन्ही गृहीत धरा. मग एकतर $\psi \ident \psi'$
किंवा तसे नाही. तसे असेल, तर स्पष्टपणे $\ast = {\ast'}$ आणि $\chi
\ident \chi'$, कारण त्या स्थितीत त्या $\varphi$ च्या एकाच स्थानी सुरू होणाऱ्या
आणि समान लांबीच्या उपचिन्हमाला आहेत. उरलेली बाब म्हणजे $\psi
\not\ident \psi'$. $\psi$ आणि $\psi'$ या दोन्ही $\varphi$ च्या एकाच स्थानी सुरू
होणाऱ्या उपचिन्हमाला असल्यामुळे, त्यांपैकी एक दुसरीचा उचित पूर्वखंड
असला पाहिजे. पण \ref{fol:syn:unq:lem:no-prefix} नुसार हे अशक्य आहे.
\end{proof}













\section{सूत्राचा मुख्य संकारक}\label{fol:syn:mai:sec}

\begin{explain}
सूत्र~$\varphi$ तयार करताना सर्वांत शेवटी वापरलेल्या कारकाविषयी
बोलणे अनेकदा उपयुक्त असते. या कारकाला $\varphi$ चा \emph{मुख्य कारक}
म्हणतात. अंतःप्रेरणेने तो $\varphi$ चा ``सर्वांत बाहेरचा'' कारक होय.
उदाहरणार्थ, $\lnot \varphi$ चा मुख्य कारक $\lnot$, $(\varphi \lor \psi)$ चा
मुख्य कारक $\lor$, इत्यादी.
\end{explain}


\begin{defn}[मुख्य संकारक]
\label{fol:syn:mai:def:main-op}
सूत्र~$\varphi$ चा \emph{मुख्य संकारक} पुढीलप्रमाणे परिभाषित केला जातो:
\begin{enumerate}
\item \indcase*{\varphi}{\varphi}{$\indfrm$ ला मुख्य संकारक नाही.}

\item \indcase{\varphi}{\lnot \psi}{$\indfrm$ चा मुख्य संकारक
  हा~$\lnot$ आहे.}

\item \indcase{\varphi}{(\psi \land \chi)}{$\indfrm$ चा
  मुख्य संकारक हा~$\land$ आहे.}

\item \indcase{\varphi}{(\psi \lor \chi)}{$\indfrm$ चा
  मुख्य संकारक हा~$\lor$ आहे.}

\item \indcase{\varphi}{(\psi \lif \chi)}{$\indfrm$ चा
  मुख्य संकारक हा~$\lif$ आहे.}

\item \indcase{\varphi}{(\psi \liff \chi)}{$\indfrm$ चा
  मुख्य संकारक हा~$\liff$ आहे.}

\item \indcase{\varphi}{\lforall[x][\psi]}{$\indfrm$ चा
  मुख्य संकारक हा~$\lforall$ आहे.}

\item \indcase{\varphi}{\lexists[x][\psi]}{$\indfrm$ चा
  मुख्य संकारक हा~$\lexists$ आहे.}
\end{enumerate}
\end{defn}

प्रत्येक बाबतीत सूत्रातील मुख्य संकारक ची नेमकी दर्शवलेली
\emph{आस्थिति} आपल्याला अभिप्रेत आहे. उदाहरणार्थ,
$((\theta \lif !E) \lif (!E \lif \theta))$ हे सूत्र $(\psi \lif \chi)$ या रूपात
असून $\psi$ हे $(\theta \lif !E)$ आणि $\chi$ हे $(!E \lif \theta)$ असल्यामुळे,
$\lif$ ची दुसरी आस्थिति मुख्य संकारक आहे.

\begin{explain}
ही अशा फलनाची \emph{पुनरावर्ती} व्याख्या आहे जे सर्व अनाण्विक
सूत्रे ना त्यांच्या मुख्य संकारक च्या आस्थितीकडे पाठवते.
सूत्रांची व्याख्या विगमनाने केलेली असल्यामुळे प्रत्येक
सूत्र~$\varphi$ ला \ref{fol:syn:mai:def:main-op} मधील एखादी बाब लागू होते.
त्यामुळे प्रत्येक अनाण्विक सूत्र~$\varphi$ साठी मुख्य संकारक
अस्तित्वात आहे याची खात्री मिळते. प्रत्येक सूत्राला यांपैकी
फक्त एकच अट लागू होते आणि प्रत्येक बाबतीत $\varphi$ ज्यांपासून तयार केले
आहे ती लहान सूत्रे एकमेव रीतीने निर्धारित असतात. त्यामुळे $\varphi$
मधील मुख्य संकारक ची आस्थिति एकमेव असते आणि म्हणून आपण फलन
परिभाषित केले आहे.
\end{explain}

सूत्रांचा मुख्य संकारक कोणते चिन्ह आहे यानुसार आपण त्यांना
\ref{fol:syn:mai:tab:main-op} मधील नावे देतो.

\begin{table}[!h]
\centering
\begin{tabular}{c | c | c}
मुख्य संकारक & सूत्राचा प्रकार & उदाहरण\\
\hline
काही नाही & आण्विक (सूत्र) &
$\lfalse$,
$\ltrue$,
$\Atom{R}{t_1, \dots, t_n}$,
$\eq[t_1][t_2]$\\
$\lnot$ & नकरण & $\lnot \varphi$ \\
$\land$ & संधियोग & $(\varphi \land \psi$) \\
$\lor$ & विकल्पयोग & $(\varphi \lor \psi$) \\
$\lif$ & शर्त & $(\varphi \lif \psi$) \\
$\liff$ & द्विशर्त & $(\varphi \liff \psi)$ \\
$\lforall[][]$ & वैश्विक (सूत्र)& $\lforall[x][\varphi]$ \\
$\lexists[][]$ & अस्तित्ववाची (सूत्र)& $\lexists[x][\varphi]$
\end{tabular}
\caption{सूत्रांचे मुख्य कारक आणि नावे}
\label{fol:syn:mai:tab:main-op}
\end{table}












\section{उपसूत्र}\label{fol:syn:sbf:sec}

\begin{explain}
दिलेल्या सूत्राचे ``घटक'' असलेल्या सूत्रे विषयी बोलणे अनेकदा
उपयुक्त असते. त्यांना आपण त्याची \emph{उपसूत्रे} म्हणतो. कोणतेही
सूत्र हे स्वतःचे उपसूत्र मानले जाते; $\varphi$ स्वतः सोडून
$\varphi$ चे इतर उपसूत्र म्हणजे \emph{उचित उपसूत्र} होय.
\end{explain}

\begin{defn}[तात्काळ उपसूत्र]
$\varphi$ हे सूत्र असेल, तर $\varphi$ ची \emph{तात्काळ उपसूत्रे}
पुढीलप्रमाणे विगमनाने परिभाषित केली जातात:
\begin{enumerate}
\item आण्विक सूत्रे ना तात्काळ उपसूत्रे नसतात.

\item \indcase{\varphi}{\lnot \psi}{$\indfrm$ चे एकमेव तात्काळ
    उपसूत्र हे~$\psi$ आहे.}

\item \indcase{\varphi}{(\psi \ast \chi)}{$\indfrm$ ची तात्काळ उपसूत्रे
  $\psi$ आणि $\chi$ आहेत ($\ast$ हे कोणतेही द्विस्थानी संयोजक आहे).}

\item \indcase{\varphi}{\lforall[x][\psi]}{$\indfrm$ चे एकमेव तात्काळ
    उपसूत्र हे~$\psi$ आहे.}

\item \indcase{\varphi}{\lexists[x][\psi]}{$\indfrm$ चे एकमेव तात्काळ
    उपसूत्र हे~$\psi$ आहे.}
\end{enumerate}
\end{defn}

\begin{defn}[उचित उपसूत्र]
$\varphi$ हे सूत्र असेल, तर $\varphi$ ची \emph{उचित उपसूत्रे}
पुढीलप्रमाणे पुनरावर्ती रीतीने परिभाषित केली जातात:
\begin{enumerate}
\item आण्विक सूत्रे ना उचित उपसूत्रे नसतात.

\item \indcase{\varphi}{\lnot \psi}{$\indfrm$ ची उचित
    उपसूत्रे म्हणजे~$\psi$ आणि $\psi$ ची सर्व उचित उपसूत्रे.}

\item \indcase{\varphi}{(\psi \ast \chi)}{$\indfrm$ ची उचित उपसूत्रे
  म्हणजे $\psi$, $\chi$, $\psi$ ची सर्व उचित उपसूत्रे आणि~$\chi$ ची
  सर्व उचित उपसूत्रे.}

\item \indcase{\varphi}{\lforall[x][\psi]}{$\indfrm$ ची उचित
    उपसूत्रे म्हणजे~$\psi$ आणि $\psi$ ची सर्व उचित
    उपसूत्रे.}

\item \indcase{\varphi}{\lexists[x][\psi]}{$\indfrm$ ची उचित
    उपसूत्रे म्हणजे~$\psi$ आणि $\psi$ ची सर्व उचित
    उपसूत्रे.}
\end{enumerate}
\end{defn}

\begin{defn}[उपसूत्र]
$\varphi$ ची उपसूत्रे म्हणजे स्वतः $\varphi$ आणि त्याची सर्व उचित
उपसूत्रे.
\end{defn}

\begin{explain}
तात्काळ उपसूत्रे आणि उचित उपसूत्रे यांच्या व्याख्यांतील
सूक्ष्म फरक लक्षात घ्या. पहिल्या बाबतीत, $\varphi$ च्या प्रत्येक शक्य
रूपासाठी आपण सूत्र~$\varphi$ ची तात्काळ उपसूत्रे थेट परिभाषित केली
आहेत. ही बाबींनुसार केलेली स्पष्ट व्याख्या आहे आणि त्या बाबी
सूत्रांच्या संचाच्या विगमनाधारित व्याख्येला समांतर आहेत. दुसऱ्या
बाबतीतही सर्व सूत्रांचा संच ज्या रीतीने परिभाषित केला आहे तिला
आपण समांतर गेलो आहोत; पण प्रत्येक बाबतीत लहान सूत्रे $\psi$, $\chi$
यांची उचित उपसूत्रे या सूत्रे ना स्वतःलाही समाविष्ट करून
घेतली आहेत.
त्यामुळे व्याख्या \emph{पुनरावर्ती} होते. सर्वसाधारणपणे, विगमनाने
परिभाषित केलेल्या संचावरील फलनाची व्याख्या (आपल्या बाबतीत हा संच म्हणजे
सूत्रे) पुनरावर्ती असते, जर फलनाच्या व्याख्येतील बाबी स्वतः त्या
फलनाचा वापर करत असतील. मात्र ती सुनिर्धारित असण्यासाठी, आपण परिभाषित
करत असलेल्या युक्तिवादाच्या ``आधी'' येणाऱ्या युक्तिवादांसाठीचीच फलनाची
मूल्ये वापरतो याची खात्री केली पाहिजे---आपल्या बाबतीत $(\psi \ast \chi)$
ची ``उचित उपसूत्र'' परिभाषित करताना आपण केवळ ``आधीच्या''
सूत्रे $\psi$ आणि $\chi$ यांची उचित उपसूत्रे वापरतो.
\end{explain}

\begin{prop}
\label{fol:syn:sbf:prop:subfrm-trans}
$\psi$ हे $\varphi$ चे उपसूत्र आणि $\chi$ हे $\psi$ चे उपसूत्र आहे असे गृहीत धरा.
मग $\chi$ हे $\varphi$ चे उपसूत्र आहे. दुसऱ्या शब्दांत, उपसूत्र-संबंध
संक्रमक आहे.
\end{prop}

\begin{prob}
\ref{fol:syn:sbf:prop:subfrm-trans} सिद्ध करा.
\end{prob}

\begin{prop}
\label{fol:syn:sbf:prop:count-subfrms}
$\varphi$ हे $n$ संयोजके आणि संख्यापक असलेले सूत्र आहे असे गृहीत धरा.
मग $\varphi$ ची जास्तीत जास्त $2n+1$ उपसूत्रे आहेत.
\end{prop}

\begin{prob}
\ref{fol:syn:sbf:prop:count-subfrms} सिद्ध करा.
\end{prob}












\section{रचनाक्रमिका}\label{fol:syn:fseq:sec}

सूत्रांची विगमनाधारित व्याख्या देणे आणि त्याला पूरक तंत्र म्हणून
सूत्रांचे गुणधर्म विगमनाने सिद्ध करणे हा सुबक व कार्यक्षम
दृष्टिकोन आहे. तथापि, सूत्रांच्या रचनेचा तळापासून वर जाणारा,
टप्प्याटप्प्याचा दृष्टिकोनही उपयुक्त ठरू शकतो; येथे आपण तो
\emph{रचनाक्रमिका} या संकल्पनेच्या साहाय्याने मांडतो.

पदे आणि सूत्रे त्यांच्या विगमनाधारित व्याख्यांचा निर्देश न करता
या रीतीने कशी मांडता येतात हे दाखवण्यासाठी, आपण प्रथम एखाद्या
भाषा~$\Lang L$ मधील चिन्हांपासून बनलेल्या स्वेच्छ चिन्हमालेची संकल्पना
मांडतो.

\begin{defn}[चिन्हमाला]
\label{fol:syn:fseq:defn:string}
$\Lang L$ ही प्रथम-क्रम भाषा आहे असे समजा. \emph{$\Lang L$-चिन्हमाला}
म्हणजे $\Lang L$ च्या चिन्हांची सांत क्रमिका. संदर्भावरून भाषा~$\Lang L$
स्पष्टपणे ठरलेली असेल, तर $\Lang L$-चिन्हमालेला आपण अनेकदा फक्त
\emph{चिन्हमाला} म्हणू.
\end{defn}

\begin{ex}
कोणत्याही प्रथम-क्रम भाषा $\Lang L$ साठी सर्व $\Lang L$-सूत्रे
या $\Lang L$-चिन्हमाला असतात; पण याचे उलट खरे नसते. उदाहरणार्थ,
\[)(\Obj v_0\lif\lexists\] ही $\Lang L$-चिन्हमाला आहे, पण
$\Lang L$-सूत्र नाही.
\end{ex}

\begin{defn}[पदांसाठी रचनाक्रमिका]
\label{fol:syn:fseq:defn:fseq-trm}
$\Lang L$-चिन्हमालांची सांत क्रमिका $\tuple{t_0,\dotsc,t_n}$ ही पद~$t$
साठी \emph{रचनाक्रमिका} असते, जर $t \ident t_n$ आणि प्रत्येक $i \leq n$
साठी एकतर $t_i$ हे चल किंवा स्थिरांक असेल, अथवा $\Lang
L$ मध्ये एखादे $k$-स्थानी फलन~$f$ असेल आणि असे
$m_1,\dotsc,m_k < i$ अस्तित्वात असतील की
$t_i \ident f(t_{m_1},\dotsc,t_{m_k})$. फरक स्पष्ट करणे आवश्यक असेल,
तर पदांसाठीच्या रचनाक्रमिकांना आपण \emph{पद-रचनाक्रमिका} म्हणू.

\end{defn}

\begin{ex}
क्रमिका
\[
    \tuple{\Obj c_0, \Obj v_0, \Atom{\Obj f^2_0}{\Obj c_0, \Obj v_0}, \Atom{\Obj f^1_0}{\Atom{\Obj f^2_0}{\Obj c_0, \Obj v_0}}}
\]
ही पद $\Atom{\Obj f^1_0}{\Atom{\Obj
f^2_0}{\Obj c_0, \Obj v_0}}$ साठी रचनाक्रमिका आहे; त्याचप्रमाणे
\[
    \tuple{\Obj v_0, \Obj c_0, \Atom{\Obj f^2_0}{\Obj c_0, \Obj v_0}, \Atom{\Obj f^1_0}{\Atom{\Obj f^2_0}{\Obj c_0, \Obj v_0}}}.
\]
हीदेखील त्याची रचनाक्रमिका आहे.
\end{ex}

\begin{defn}[सूत्रांसाठी रचनाक्रमिका]
\label{fol:syn:fseq:defn:fseq-frm}
$\Lang L$-चिन्हमालांची सांत क्रमिका $\tuple{\varphi_0,\dotsc,\varphi_n}$ ही
$\varphi$ साठी \emph{रचनाक्रमिका} असते, जर $\varphi \ident \varphi_n$ आणि प्रत्येक
$i \leq n$ साठी एकतर $\varphi_i$ हे आण्विक सूत्र असेल, अथवा पुढीलपैकी
एखादी अट पूर्ण करणारे $j,k < i$ आणि एखादे चल~$x$ अस्तित्वात
असतील:
\begin{enumerate}
    \item $\varphi_i \ident \lnot \varphi_j$.
    \item $\varphi_i \ident (\varphi_j \land \varphi_k)$.
    \item $\varphi_i \ident (\varphi_j \lor \varphi_k)$.
    \item $\varphi_i \ident (\varphi_j \lif \varphi_k)$.
    \item $\varphi_i \ident (\varphi_j \liff \varphi_k)$.
    \item $\varphi_i \ident \lforall[x][\varphi_j]$.
    \item $\varphi_i \ident \lexists[x][\varphi_j]$.
\end{enumerate}
फरक स्पष्ट करणे आवश्यक असेल, तर सूत्रांसाठीच्या रचनाक्रमिकांना आपण
\emph{सूत्र-रचनाक्रमिका} म्हणू.
\end{defn}

\begin{ex}
\[
\tuple{
    \Atom{\Obj A^1_0}{\Obj v_0},
    \Atom{\Obj A^1_1}{\Obj c_1},
    (\Atom{\Obj A^1_1}{\Obj c_1} \land \Atom{\Obj A^1_0}{\Obj v_0}),
    \lexists[\Obj v_0][(\Atom{\Obj A^1_1}{\Obj c_1} \land \Atom{\Obj A^1_0}{\Obj v_0})]
}
\]
ही $\lexists[\Obj v_0][(\Atom{\Obj A^1_1}{\Obj c_1} \land \Atom{\Obj
A^1_0}{\Obj v_0})]$ ची एक रचनाक्रमिका आहे; त्याचप्रमाणे
\begin{multline*}
\tuple{
    \Atom{\Obj A^1_0}{\Obj v_0},
    \Atom{\Obj A^1_1}{\Obj c_1},
    (\Atom{\Obj A^1_1}{\Obj c_1} \land \Atom{\Obj A^1_0}{\Obj v_0}),
    \Atom{\Obj A^1_1}{\Obj c_1},\\
    \lforall[\Obj v_1][\Atom{\Obj A^1_0}{\Obj v_0}],
    \lexists[\Obj v_0][(\Atom{\Obj A^1_1}{\Obj c_1} \land \Atom{\Obj A^1_0}{\Obj v_0})]
}.
\end{multline*}

दुसऱ्या उदाहरणावरून दिसते की रचनाक्रमिकेत ``अडगळ'' असू शकते:
अनावश्यक असलेली किंवा रचनेत काहीही हातभार न लावणारी सूत्रे.
\end{ex}

\begin{prop}\label{fol:syn:fseq:prop:formed}
$\Frm[L]$ मधील प्रत्येक सूत्र~$\varphi$ ला एक रचनाक्रमिका असते.
\end{prop}

\begin{proof}
$\varphi$ आण्विक आहे असे समजा. मग क्रमिका~$\tuple{\varphi}$ ही $\varphi$ साठी
रचनाक्रमिका आहे.

आता $\psi$ आणि~$\chi$ यांना अनुक्रमे $\tuple{\psi_0,\dotsc,\psi_n}$ आणि
$\tuple{\chi_0,\dotsc,\chi_m}$ या रचनाक्रमिका आहेत असे समजा.

\begin{enumerate}
  \item जर $\varphi \ident \lnot \psi$, तर
      $\tuple{\psi_0,\dotsc,\psi_n,\lnot \psi_n}$ ही $\varphi$ साठी
      रचनाक्रमिका आहे.
  \item जर $\varphi \ident (\psi \land \chi)$, तर
      $\tuple{\psi_0,\dotsc,\psi_n,\chi_0,\dotsc,\chi_m,(\psi_n \land \chi_m)}$
      ही $\varphi$ साठी रचनाक्रमिका आहे.
  \item जर $\varphi \ident (\psi \lor \chi)$, तर
      $\tuple{\psi_0,\dotsc,\psi_n,\chi_0,\dotsc,\chi_m,(\psi_n \lor \chi_m)}$
      ही $\varphi$ साठी रचनाक्रमिका आहे.
  \item जर $\varphi \ident (\psi \lif \chi)$, तर
      $\tuple{\psi_0,\dotsc,\psi_n,\chi_0,\dotsc,\chi_m,(\psi_n \lif \chi_m)}$
      ही $\varphi$ साठी रचनाक्रमिका आहे.
  \item जर $\varphi \ident (\psi \liff \chi)$, तर
      $\tuple{\psi_0,\dotsc,\psi_n,\chi_0,\dotsc,\chi_m,(\psi_n \liff \chi_m)}$
      ही $\varphi$ साठी रचनाक्रमिका आहे.
  \item जर $\varphi \ident \lforall[x][\psi]$, तर
      $\tuple{\psi_0,\dotsc,\psi_n,\lforall[x][\psi_n]}$
      ही $\varphi$ साठी रचनाक्रमिका आहे.
  \item जर $\varphi \ident \lexists[x][\psi]$, तर
      $\tuple{\psi_0,\dotsc,\psi_n,\lexists[x][\psi_n]}$
      ही $\varphi$ साठी रचनाक्रमिका आहे.
\end{enumerate}
सूत्रे वरील विगमनाच्या तत्त्वानुसार प्रत्येक सूत्राला एक
रचनाक्रमिका असते.
\end{proof}

याचे उलट विधानही आपण सिद्ध करू शकतो. ते महत्त्वाचे आहे, कारण त्यामुळे
सूत्रे परिभाषित करण्याच्या आपल्या दोन्ही पद्धती सममूल्य आहेत, म्हणजे
त्या समान परिणाम देतात, हे दिसते. तसेच रचनाक्रमिकांच्या लांबीवर साधे
विगमन वापरून सूत्रांविषयीची प्रमेये सिद्ध करता येतात.

\begin{lem}
\label{fol:syn:fseq:lem:fseq-init}
$\tuple{\varphi_0,\dotsc,\varphi_n}$ ही $\varphi_n$ साठी रचनाक्रमिका आहे आणि
$k \leq n$ असे समजा. मग $\tuple{\varphi_0,\dotsc,\varphi_k}$ ही $\varphi_k$ साठी
रचनाक्रमिका आहे.
\end{lem}

\begin{proof}
अभ्यासप्रश्न.
\end{proof}

\begin{prob}
\ref{fol:syn:fseq:lem:fseq-init} सिद्ध करा.
\end{prob}

\begin{thm}
\label{fol:syn:fseq:thm:fseq-frm-equiv}
$\Frm[L]$ हा अशा सर्व $\Lang L$-चिन्हमाला~$\varphi$ चा संच आहे की $\varphi$
साठी एखादी सूत्र-रचनाक्रमिका अस्तित्वात आहे.
\end{thm}

\begin{proof}
भाषा~$\Lang L$ मधील ज्या सर्व चिन्हमालांना रचनाक्रमिका आहे, त्यांचा
संच~$F$ असू द्या. $\Frm[L] \subseteq F$ हे
\ref{fol:syn:fseq:prop:formed} मध्ये आपण पाहिले आहे; म्हणून आता उलट
समावेश सिद्ध करू.

$\varphi$ ला $\tuple{\varphi_0,\dotsc,\varphi_n}$ ही रचनाक्रमिका आहे असे समजा.
$n$ वरील प्रबल विगमनाने $\varphi \in \Frm[L]$ हे सिद्ध करू. लांबी
$m < n$ असलेली रचनाक्रमिका ज्या प्रत्येक चिन्हमालेला आहे, ती
$\Frm[L]$ मध्ये आहे, हे आपले विगमन गृहीतक आहे. रचनाक्रमिकेच्या
व्याख्येनुसार एकतर $\varphi \ident \varphi_n$ आण्विक आहे किंवा पुढीलपैकी एखादी
अट पूर्ण करणारे $j,k < n$ अस्तित्वात असले पाहिजेत:
\begin{enumerate}
    \item $\varphi \ident \lnot \varphi_j$.
    \item $\varphi \ident (\varphi_j \land \varphi_k)$.
    \item $\varphi \ident (\varphi_j \lor \varphi_k)$.
    \item $\varphi \ident (\varphi_j \lif \varphi_k)$.
    \item $\varphi \ident (\varphi_j \liff \varphi_k)$.
    \item $\varphi \ident \lforall[x][\varphi_j]$.
    \item $\varphi \ident \lexists[x][\varphi_j]$.
\end{enumerate}
आता प्रत्येक शक्यतेचा स्वतंत्र विचार करू. $\varphi$ आण्विक असेल, तर
$\varphi_n \in \Frm[L]$. त्याऐवजी $\varphi \ident (\varphi_j \land \varphi_k)$ असे समजा.


\ref{fol:syn:fseq:lem:fseq-init} नुसार
$\tuple{\varphi_0,\dotsc,\varphi_j}$ आणि $\tuple{\varphi_0,\dotsc,\varphi_k}$ या अनुक्रमे
$\varphi_j$ आणि~$\varphi_k$ साठी रचनाक्रमिका आहेत. या $\varphi$ च्या
रचनाक्रमिकेच्या उचित आरंभीच्या उपक्रमिका असल्यामुळे दोन्हींची लांबी
$n$ पेक्षा कमी आहे. म्हणून विगमन गृहीतकानुसार $\varphi_j$ आणि~$\varphi_k$ ही
$\Frm[L]$ मध्ये आहेत; त्यामुळे सूत्राच्या व्याख्येनुसार
$(\varphi_j \land \varphi_k)$ हेही त्यात आहे. इतर शक्यता समांतर युक्तिवादाने
सिद्ध होतात.
\end{proof}

पदांसाठीच्या रचनाक्रमिकांना सूत्रांच्या रचनाक्रमिकांसारखेच गुणधर्म
असतात.

\begin{prop}
\label{fol:syn:fseq:prop:fseq-trm-equiv}
$\Trm[L]$ हा अशा सर्व $\Lang L$-चिन्हमाला $t$ चा संच आहे की $t$
साठी एखादी पद-रचनाक्रमिका अस्तित्वात आहे.
\end{prop}

\begin{proof}
अभ्यासप्रश्न.
\end{proof}

\begin{prob}
\ref{fol:syn:fseq:prop:fseq-trm-equiv} सिद्ध करा.
सूचना: \ref{fol:syn:fseq:thm:fseq-frm-equiv} च्या सिद्धतेत वापरलेली
रणनीती वापरा.
\end{prob}

रचनाक्रमिकेत दोन प्रकारची ``अडगळ'' येऊ शकते: पुनरावृत्त घटक आणि सूत्र
किंवा पद यांच्या रचनेशी असंबद्ध घटक. न्यूनतम रचनाक्रमिका विचारात घेऊन
आपण दोन्ही प्रकार दूर करू शकतो.

\begin{defn}[न्यूनतम रचनाक्रमिका]
\label{fol:syn:fseq:defn:minimal-fseq}
सूत्र~$\varphi$ साठीची रचनाक्रमिका $\tuple{\varphi_0, \ldots, \varphi_n}$ ही
$\varphi$ साठी \emph{न्यूनतम रचनाक्रमिका} असते, जर $\varphi$ साठीची प्रत्येक
दुसरी रचनाक्रमिका~$s$ घेतली असता~$s$ ची लांबी $n+1$ पेक्षा अधिक किंवा
तिच्याइतकी असेल.

त्याचप्रमाणे, पद~$t$ साठीची रचनाक्रमिका $\tuple{t_0, \ldots, t_n}$ ही
$t$ साठी \emph{न्यूनतम रचनाक्रमिका} असते, जर $t$ साठीची प्रत्येक
दुसरी रचनाक्रमिका~$s$ घेतली असता~$s$ ची लांबी $n+1$ पेक्षा अधिक किंवा
तिच्याइतकी असेल.
\end{defn}

एका सूत्राला किंवा पदाला एकापेक्षा अधिक न्यूनतम रचनाक्रमिका असू शकतात;
पण त्यांत नेमक्या त्याच चिन्हमाला असतील, हे लक्षात घ्या.

\begin{prop}
\label{fol:syn:fseq:prop:subformula-equivs}
पुढील विधाने सममूल्य आहेत:
\begin{enumerate}
    \item $\psi$ हे $\varphi$ चे उप-सूत्र आहे.
    \item $\psi$ हे $\varphi$ च्या प्रत्येक रचनाक्रमिकेत येते.
    \item $\psi$ हे $\varphi$ च्या एखाद्या न्यूनतम रचनाक्रमिकेत येते.
\end{enumerate}
\end{prop}

\begin{proof}
अभ्यासप्रश्न.
\end{proof}

\begin{prob}
\ref{fol:syn:fseq:prop:subformula-equivs} सिद्ध करा.
\end{prob}

\begin{history}
रेमंड स्मलियन यांनी त्यांच्या \emph{First-Order Logic}
(Smullyan, 1968) या पाठ्यपुस्तकात रचनाक्रमिका प्रथम मांडल्या.
रचनाक्रमिकांचे आणखी गुणधर्म Zuckerman (1973) यांनी प्रस्थापित केले.
\end{history}












\section{मुक्त चल आणि वाक्य}\label{fol:syn:fvs:sec}

\begin{defn}[चल च्या मुक्त आस्थिती]
\label{fol:syn:fvs:defn:free-occ}
सूत्र मधील चल च्या \emph{मुक्त} आस्थिती पुढीलप्रमाणे
विगमनाने परिभाषित केल्या जातात:
\begin{enumerate}
\item \indcase*{\varphi}{$\varphi$ आण्विक आहे}{$\indfrm$ मधील सर्व
  चल-आस्थिती मुक्त असतात.}

\item \indcase{\varphi}{\lnot \psi}{$\indfrm$ मधील मुक्त
  चल-आस्थिती म्हणजे नेमक्या $\psi$ मधील मुक्त आस्थिती.}

\item \indcase{\varphi}{(\psi \ast \chi)}{$\indfrm$ मधील मुक्त
  चल-आस्थिती म्हणजे $\psi$ मधील आस्थिती आणि त्यांच्यासोबत
  $\chi$ मधील आस्थिती.}

\item \indcase{\varphi}{\lforall[x][\psi]}{$\indfrm$ मधील मुक्त
  चल-आस्थिती म्हणजे $\psi$ मधील $x$ च्या आस्थिती सोडून इतर
  सर्व मुक्त आस्थिती.}

\item \indcase{\varphi}{\lexists[x][\psi]}{$\indfrm$ मधील मुक्त
  चल-आस्थिती म्हणजे $\psi$ मधील $x$ च्या आस्थिती सोडून इतर
  सर्व मुक्त आस्थिती.}
\end{enumerate}
\end{defn}

\begin{defn}[बद्ध चरे]
सूत्र~$\varphi$ मधील चल ची आस्थिती मुक्त नसेल, तर ती
\emph{बद्ध} असते.
\end{defn}

\begin{prob}
\ref{fol:syn:fvs:defn:free-occ} च्या धर्तीवर बद्ध चर-आस्थितींची
विगमनाधारित व्याख्या द्या.
\end{prob}

\begin{defn}[व्याप्ती]
सूत्र~$\varphi$ मध्ये $\lforall[x][\psi]$ ची एखाद्या उपसूत्राची
  आस्थिती असेल, तर $\varphi$ मधील $\psi$ च्या संबंधित आस्थितीला
  $\lforall[x]$ च्या संबंधित आस्थितीची \emph{व्याप्ती} म्हणतात.
  $\lexists[x]$ साठीही त्याचप्रमाणे.

$\varphi$ मधील संख्यापकाच्या
$\lforall[x]$ किंवा
    $\lexists[x]$ या आस्थितीची व्याप्ती $\psi$ असेल,
तर $\psi$ मधील $x$ च्या मुक्त आस्थिती
$\lforall[x][\psi]$ आणि
$\lexists[x][\psi]$ मध्ये बद्ध असतात. या आस्थिती
त्या उल्लेख केलेल्या संख्यापक-आस्थितीने \emph{बद्ध केल्या आहेत} असे
आपण म्हणतो.
\end{defn}

\begin{ex}
पुढील सूत्राचा विचार करा:
\[
\lexists[\Obj v_0][\underbrace{\Atom{\Obj A^2_0}{\Obj v_0,\Obj v_1}}_{\psi}]
\]
$\psi$ हे $\lexists[\Obj v_0]$ ची व्याप्ती दर्शवते. हा संख्यापक $\psi$
मधील $\Obj v_0$ ची आस्थिती बद्ध करतो; पण $\Obj v_1$ ची आस्थिती बद्ध
करत नाही. म्हणून या बाबतीत $\Obj v_1$ हे मुक्त चर आहे.


अधिक गुंतागुंतीच्या सूत्र~$\varphi$ मध्ये हे कसे कार्य करते ते आता
पाहू शकतो:
\[
\lforall[\Obj v_0][\underbrace{(\Atom{\Obj A^1_0}{\Obj v_0} \lif
    \Atom{\Obj A^2_0}{\Obj v_0, \Obj v_1})}_{\psi}] \lif \lexists[\Obj
  v_1][\underbrace{(\Atom{\Obj A^2_1}{\Obj v_0, \Obj v_1} \lor \lforall[\Obj v_0][\overbrace{\lnot \Atom{\Obj A^1_1}{\Obj v_0}}^{\theta}])}_{\chi}]
\]
$\psi$ ही पहिल्या $\lforall[\Obj v_0]$ ची व्याप्ती, $\chi$ ही
$\lexists[\Obj v_1]$ ची व्याप्ती आणि $\theta$ ही दुसऱ्या
$\lforall[\Obj v_0]$ ची व्याप्ती आहे. पहिले $\lforall[\Obj v_0]$ हे
$\psi$ मधील $\Obj v_0$ च्या आस्थिती बद्ध करते, $\lexists[\Obj v_1]$ हे
$\chi$ मधील $\Obj v_1$ ची आस्थिती बद्ध करते आणि दुसरे
$\lforall[\Obj v_0]$ हे $\theta$ मधील $\Obj v_0$ ची आस्थिती बद्ध करते.
$\Obj v_1$ ची पहिली आस्थिती आणि $\Obj v_0$ ची चौथी आस्थिती $\varphi$ मध्ये
मुक्त आहेत. $\Obj v_0$ ची शेवटची आस्थिती $\theta$ मध्ये मुक्त, पण $\chi$
आणि~$\varphi$ मध्ये बद्ध आहे.
\end{ex}

\begin{defn}[वाक्य]
सूत्र~$\varphi$ मध्ये चलांच्या मुक्त आस्थिती नसतील तेव्हा
आणि केवळ तेव्हाच ते \emph{वाक्य} असते.
\end{defn}














\section{आदेशन}\label{fol:syn:sub:sec}

\begin{defn}[पदातील आदेशन]
$s$ मधील~$x$ च्या प्रत्येक आस्थितीच्या जागी $t$ \emph{आदेशित करून}
मिळणारा परिणाम $\Subst{s}{t}{x}$ आपण पुढीलप्रमाणे पुनरावर्ती रीतीने
परिभाषित करतो:
\begin{enumerate}
\item \indcase{s}{c}{$\Subst{\indfrm}{t}{x}$ म्हणजे फक्त $s$.}

\item \indcase{s}{y}{$y$ हे चर असून $y \not\ident x$ असेल, तर
  $\Subst{\indfrm}{t}{x}$ म्हणजेही फक्त~$s$.}

\item \indcase{s}{x}{$\Subst{\indfrm}{t}{x}$ म्हणजे~$t$.}

\item\indcase{s}{\Atom{f}{t_1, \dots, t_n}}{$\Subst{\indfrmp}{t}{x}$
  म्हणजे $\Atom{f}{\Subst{t_1}{t}{x}, \dots,
  \Subst{t_n}{t}{x}}$.}
\end{enumerate}
\end{defn}

\begin{defn}
$\varphi$ मधील~$x$ ची कोणतीही मुक्त आस्थिती $t$ मधील एखादे चर बद्ध करणाऱ्या
संख्यापकाच्या व्याप्तीत येत नसेल, तर पद~$t$ हे $\varphi$ मध्ये $x$ साठी
\emph{आदेशनासाठी मुक्त} असते.
\end{defn}

\begin{ex} ~
\begin{enumerate}
\item $\Obj v_8$ हे $\lexists[\Obj
  v_3]\Atom{\Obj A^2_4}{\Obj v_3,\Obj v_1}$ मध्ये $\Obj v_1$ साठी
  मुक्त आहे.

\item $\Obj f^2_1(\Obj v_1, \Obj v_2)$ हे
  $\lforall[\Obj v_2]\Atom{\Obj A^2_4}{\Obj v_0,\Obj v_2}$ मध्ये
  $\Obj v_0$ साठी \emph{मुक्त नाही}.
\end{enumerate}
\end{ex}

\begin{defn}[सूत्र मधील आदेशन]
$\varphi$ हे सूत्र, $x$ हे चल आणि $t$ हे $\varphi$ मध्ये~$x$
साठी आदेशनासाठी मुक्त असलेले पद असेल, तर $\varphi$ मधील~$x$ च्या सर्व मुक्त
आस्थितींच्या जागी $t$ आदेशित करून मिळणारा परिणाम म्हणजे
$\Subst{\varphi}{t}{x}$.
\begin{enumerate}
\item \indcase{\varphi}{\lfalse}{$\Subst{\indfrm}{t}{x}$
    म्हणजे $\lfalse$.}

\item \indcase{\varphi}{\ltrue}{$\Subst{\indfrm}{t}{x}$
    म्हणजे $\ltrue$.}

\item \indcase{\varphi}{\Atom{P}{t_1,\dots,
    t_n}}{$\Subst{\indfrm}{t}{x}$ म्हणजे
    $\Atom{P}{\Subst{t_1}{t}{x}, \dots, \Subst{t_n}{t}{x}}$.}

\item \indcase{\varphi}{\eq[t_1][t_2]}{$\Subst{\indfrmp}{t}{x}$ म्हणजे
  $\Subst{t_1}{t}{x} = \Subst{t_2}{t}{x}$.}

\item \indcase{\varphi}{\lnot \psi}{$\Subst{\indfrmp}{t}{x}$
    म्हणजे $\lnot \Subst{\psi}{t}{x}$.}

\item \indcase{\varphi}{(\psi \land
    \chi)}{$\Subst{\indfrmp}{t}{x}$ म्हणजे $(\Subst{\psi}{t}{x} \land
    \Subst{\chi}{t}{x})$.}

\item \indcase{\varphi}{(\psi \lor
    \chi)}{$\Subst{\indfrmp}{t}{x}$ म्हणजे $(\Subst{\psi}{t}{x} \lor
    \Subst{\chi}{t}{x})$.}

\item \indcase{\varphi}{(\psi \lif
    \chi)}{$\Subst{\indfrmp}{t}{x}$ म्हणजे $(\Subst{\psi}{t}{x} \lif
    \Subst{\chi}{t}{x})$.}

\item \indcase{\varphi}{(\psi \liff
    \chi)}{$\Subst{\indfrmp}{t}{x}$ म्हणजे $(\Subst{\psi}{t}{x} \liff
    \Subst{\chi}{t}{x})$.}

\item
\indcase{\varphi}{\lforall[y][\psi]}{$y$ हे $x$ पेक्षा वेगळे चर असेल, तर
    $\Subst{\indfrmp}{t}{x}$ म्हणजे
    $\lforall[y][\Subst{\psi}{t}{x}]$; अन्यथा
    $\Subst{\indfrmp}{t}{x}$ म्हणजे फक्त $\indfrm$.}

\item
\indcase{\varphi}{\lexists[y][\psi]}{$y$ हे $x$ पेक्षा वेगळे चर असेल, तर
    $\Subst{\indfrmp}{t}{x}$ म्हणजे
    $\lexists[y][\Subst{\psi}{t}{x}]$; अन्यथा
    $\Subst{\indfrmp}{t}{x}$ म्हणजे फक्त $\indfrm$.}
\end{enumerate}
\end{defn}

\begin{explain}
आदेशन निष्फळ असू शकते, हे लक्षात घ्या: $\varphi$ मध्ये $x$ मुळीच येत नसेल,
तर $\Subst{\varphi}{t}{x}$ म्हणजे फक्त~$\varphi$.

$\varphi$ मध्ये~$x$ साठी $t$ हे आदेशनासाठी मुक्त असले पाहिजे हा निर्बंध पुढील
प्रकारची प्रकरणे वगळण्यासाठी आवश्यक आहे. $\varphi \ident
\lexists[y][x < y]$ आणि $t \ident y$ असेल, तर
$\Subst{\varphi}{t}{x}$ हे $\lexists[y][y < y]$ होईल. या प्रकरणात आदेशन
करताना मुक्त चर~$y$ हे संख्यापक $\lexists[y]$ ने ``ग्रहित'' केले जाते;
हा परिणाम अनिष्ट आहे. उदाहरणार्थ, $\lforall[x][\psi]$ लागू असेल तेव्हा
$\Subst{\psi}{t}{x}$ हेही लागू असावे, अशी आपली अपेक्षा असते. पण
$\lforall[x][\lexists[y][x < y]]$ चा विचार करा (येथे $\psi$ म्हणजे
$\lexists[y][x < y]$). हे, उदाहरणार्थ, नैसर्गिक संख्यांविषयी सत्य
असलेले वाक्य आहे: प्रत्येक संख्या~$x$ पेक्षा मोठी एखादी संख्या~$y$
असते. $x$ च्या जागी $y$ चे आदेशन करण्याची परवानगी दिली, तर आपल्याला
$\Subst{\psi}{y}{x} \ident \lexists[y][y < y]$ मिळेल; हे असत्य आहे.
$t$ मधील कोणतेही मुक्त चर $\varphi$ मधील संख्यापकाने अखेरीस बद्ध होणार नाही,
अशी अट घालून आपण हे टाळतो.
\end{explain}

अवजड लेखन टाळण्यासाठी आपण पुढील संकेत अनेकदा वापरतो: सूत्र~$\varphi$
मध्ये चल~$x$ मुक्तपणे येऊ शकत असेल, तर हे दाखवण्यासाठी आपण
$\varphi(x)$ असेही लिहितो. कोणते $\varphi$ आणि~$x$ अभिप्रेत आहेत हे स्पष्ट असेल
आणि $t$ हे पद असेल ($\varphi(x)$ मध्ये $x$ साठी ते मुक्त आहे असे गृहीत
धरलेले असेल), तर $\Subst{\varphi}{t}{x}$ याचे संक्षिप्त रूप म्हणून आपण
$\varphi(t)$ लिहितो. म्हणून, उदाहरणार्थ, ``$\varphi(t)$ याला
$\lforall[x][\varphi(x)]$ चा दृष्टांत म्हणतो,'' असे आपण म्हणू शकतो. याचा
अर्थ असा: $\varphi$ हे कोणतेही सूत्र, $x$ हे चल आणि $t$ हे
$\varphi$ मध्ये~$x$ साठी मुक्त असलेले पद असेल, तर
$\Subst{\varphi}{t}{x}$ हा $\lforall[x][\varphi]$ चा दृष्टांत असतो.



\chapter{प्रथम-क्रम तर्कशास्त्राची चिन्हार्थमीमांसा}\label{fol:sem::chap}









\section{परिचय}\label{fol:syn:its:sec}

व्यंजकांना अर्थ देणे हे चिन्हार्थमीमांसेचे क्षेत्र आहे. चिन्हार्थमीमांसेतील
मध्यवर्ती संकल्पना म्हणजे संरचना मधील पूर्ति. संरचना
भाषेच्या मूलभूत घटकांना अर्थ देते: प्रांत म्हणजे वस्तूंचा अरिक्त
संच. संख्यापक या प्रांतातील घटकांवर फिरतात असा त्यांचा अर्थ लावला जातो,
स्थिरांक ना प्रांतातील घटक नेमून दिले जातात, फलने ना
प्रांत पासून त्याच प्रांताकडे जाणारी फलने नेमून दिली जातात आणि
विधेये ना प्रांत वरील संबंध नेमून दिले जातात. प्रांत
आणि मूलभूत शब्दसंग्रहातील चिन्हांना केलेल्या नेमणुका मिळून
संरचना बनते.
चले ही सूत्रे मध्ये येऊ शकतात; त्यांना चिन्हार्थ देण्यासाठी
आपल्याला प्रांत मधील घटक ही नेमून द्यावी लागतात---हेच
चर-मूल्यांकन होय. शेवटी, पूर्ति-संबंध या सर्व गोष्टी एकत्र आणतो.
सूत्राची संरचना~$\Struct{M}$ मध्ये चल च्या
मूल्यांकन~$s$ च्या सापेक्ष पूर्ति होऊ शकते; हे $\Sat{M}{\varphi}[s]$ असे
लिहितात. हा संबंधदेखील~$\varphi$ च्या संरचनेवरील विगमनाने परिभाषित केला जातो.
उदाहरणार्थ, $(\varphi \land \psi)$ ची पूर्ति परिभाषित करताना तार्किक संयोजकांचे
सत्यता-कोष्टक वापरून ती $\varphi$ आणि~$\psi$ यांची पूर्ति होते (किंवा होत नाही)
यांच्या आधारे सांगितली जाते. त्यानंतर असे निष्पन्न होते की सूत्र~$\varphi$
हे वाक्य, म्हणजे त्यात मुक्त चरे नसतील, तर चल चे
मूल्यांकन असंबद्ध ठरते; त्यामुळे संरचनांमध्ये वाक्यांची
सरळ पूर्ति होते (किंवा होत नाही) असे आपण म्हणू शकतो.

वाक्ये साठीच्या पूर्ति-संबंध~$\Sat{M}{\varphi}$ च्या आधारे आपण वैधता,
तार्किक निष्पन्नता आणि पूर्ततायोग्यता या मूलभूत चिन्हार्थविषयक संकल्पना
परिभाषित करू शकतो. वाक्याची प्रत्येक संरचनांमध्ये पूर्ति
होत असेल, तर ते वैध, $\Entails \varphi$, असते. वाक्यांच्या एखाद्या
संचातून ते तार्किकरीत्या निष्पन्न होते, $\Gamma \Entails \varphi$, जर~$\Gamma$ मधील
सर्व वाक्यांची पूर्ति करणारे प्रत्येक संरचना हे~$\varphi$ चीही
पूर्ति करत असेल. आणि एखाद्या वाक्यांच्या संचातील सर्व
वाक्यांची एकाच वेळी पूर्ति करणारे एखादे संरचना असेल, तर तो
संच पूर्ततायोग्य असतो. सूत्रे विगमनाधारित रीतीने परिभाषित केली आहेत
आणि पूर्ति हीही सूत्रांच्या संरचनेवरील विगमनाने परिभाषित केली जाते;
म्हणून आपल्या चिन्हार्थमीमांसेचे गुणधर्म सिद्ध करण्यासाठी आणि परिभाषित
चिन्हार्थविषयक संकल्पनांचा परस्परसंबंध लावण्यासाठी विगमन वापरता येते.












\section{प्रथम-क्रम भाषांसाठीच्या संरचना}\label{fol:syn:str:sec}

\begin{explain}
प्रथम-क्रम भाषांना स्वतःचा कोणताही \emph{अर्थ लावलेला नसतो:}
स्थिरांक, फलने आणि विधेये यांना कोणताही विशिष्ट
अर्थ जोडलेला नसतो. एखादी \emph{संरचना}
विनिर्दिष्ट करून अर्थ दिले जातात. ती \emph{प्रांत}, म्हणजे
स्थिरांक ज्या वस्तू निर्देशित करतात, फलने ज्यांवर क्रिया
करतात आणि संख्यापक ज्या वस्तूंवर फिरतात त्या वस्तू विनिर्दिष्ट करते.
याशिवाय, कोणते स्थिरांक कोणत्या वस्तू निर्देशित करतात,
फलन वस्तूंचे वस्तूंमध्ये प्रतिचित्रण कसे करते आणि विधेये
कोणत्या वस्तूंना लागू होतात हे ती विनिर्दिष्ट करते. संरचना या
तर्कशास्त्रातील \emph{चिन्हार्थविषयक} संकल्पनांचा---उदाहरणार्थ, तार्किक
निष्पन्नता, वैधता आणि पूर्ततायोग्यता या संकल्पनांचा---आधार आहेत. साहित्यात त्यांना
निरनिराळ्या ठिकाणी ``रचना,'' ``अर्थनिर्धारण'' किंवा ``प्रतिमान'' असे
म्हटले जाते.
\end{explain}

\begin{defn}[संरचना]
प्रथम-क्रम तर्कशास्त्राची भाषा~$\Lang{L}$ हिच्यासाठीची
\emph{संरचना}~$\Struct M$ पुढील घटकांनी बनते:
\begin{enumerate}
\item \emph{प्रांत:} अरिक्त संच~$\Domain M$.
\item \emph{स्थिरांकाचे अर्थनिर्धारण:} $\Lang{L}$ मधील प्रत्येक
  स्थिरांक~$c$ साठी घटक~$\Assign{c}{M} \in \Domain M$.
\item \emph{विधेयांचे अर्थनिर्धारण:} $\Lang{L}$ मधील प्रत्येक
  $n$-स्थानी विधेय~$R$ साठी ($\eq$ सोडून) $n$-स्थानी संबंध
  $\Assign{R}{M} \subseteq \Domain{M}^n$.
\item \emph{फलनांचे अर्थनिर्धारण:} $\Lang{L}$ मधील प्रत्येक
  $n$-स्थानी फलन~$f$ साठी $n$-स्थानी फलन
  $\Assign{f}{M} \colon \Domain{M}^n \to \Domain{M}$.
\end{enumerate}
\end{defn}

\begin{ex}
अंकगणिताच्या भाषेसाठीच्या संरचना~$\Struct M$ मध्ये एक संच,
स्थिरांक~$\Obj 0$ चे अर्थनिर्धारण म्हणून $\Domain M$ मधील
घटक~$\Assign{\Obj 0}{M}$, एकस्थानी फलन
$\Assign{\Obj \prime}{M} \colon \Domain{M} \to \Domain M$,
$\Domain M^2 \to \Domain M$ अशी दोन द्विस्थानी फलने
$\Assign{\Obj +}{M}$ आणि $\Assign{\Obj \times}{M}$, तसेच द्विस्थानी
संबंध~$\Assign{\Obj <}{M} \subseteq \Domain{M}^2$ असतात.

अशा रचनेचे एक स्वाभाविक उदाहरण पुढीलप्रमाणे आहे:
\begin{enumerate}
\item $\Domain N = \Nat$
\item $\Assign{\Obj 0}{N} = 0$
\item प्रत्येक $n \in \Nat$ साठी $\Assign{\Obj \prime}{N}(n) = n + 1$
\item प्रत्येक $n, m \in \Nat$ साठी $\Assign{\Obj +}{N}(n, m) = n + m$
\item प्रत्येक $n, m \in \Nat$ साठी $\Assign{\Obj \times}{N}(n, m) = n\cdot m$
\item $\Assign{\Obj <}{N} = \Setabs{\tuple{n, m}}{n \in \Nat, m \in
  \Nat, n < m}$
\end{enumerate}
अशा रीतीने परिभाषित केलेल्या $\Lang L_A$ साठीच्या रचनेला~$\Struct N$
\emph{अंकगणिताचे मानक प्रतिमान} म्हणतात, कारण ती~$\Lang L_A$ मधील
अतार्किक चिन्हांचा अर्थ आपण अपेक्षा करू त्याच रीतीने लावते.

तरीही, $\Lang L_A$ साठी इतर अनेक संरचना शक्य आहेत. उदाहरणार्थ,
प्रांत म्हणून~$\Nat$ ऐवजी पूर्णांकांचा संच~$\Int$ घेऊन $\Obj 0$,
$\Obj \prime$, $\Obj +$, $\Obj \times$, $\Obj <$ यांची अर्थनिर्धारणे
त्यानुसार परिभाषित करता येतील. पण संख्यांशी दुरान्वयानेही संबंध नसलेल्या
$\Lang L_A$ साठीच्या रचनाही आपण परिभाषित करू शकतो.
\end{ex}

\begin{ex}
संचसिद्धान्ताची भाषा~$\Lang L_Z$ हिच्यासाठीची रचना~$\Struct M$ हिला
फक्त एक संच आणि एकच द्विस्थानी संबंध आवश्यक असतात. म्हणून तांत्रिकदृष्ट्या,
उदाहरणार्थ, माणसांचा संच आणि ``$x$ हा~$y$ पेक्षा वयाने मोठा आहे'' हा संबंध
$\Lang L_Z$ साठी संरचना म्हणून वापरता येईल; त्याचप्रमाणे~$\Nat$
आणि $n, m \in \Nat$ साठीचा $n \ge m$ हाही वापरता येईल.


$\Lang L_Z$ साठीची एक विशेष रोचक संरचना, जिच्या प्रांतातील
घटक प्रत्यक्षात संच आहेत आणि जिच्यात $\Obj \in$ चे अर्थनिर्धारण
प्रत्यक्षात ``$x$ हा~$y$ चा घटक आहे'' हा संबंध आहे, ती म्हणजे
\emph{आनुवंशिकदृष्ट्या सांत संचांची} संरचना~$\Struct{{HF}}$:
\begin{enumerate}
\item $\Domain{{{HF}}} = \emptyset \cup \Pow{\emptyset} \cup
  \Pow{\Pow{\emptyset}} \cup \Pow{\Pow{\Pow{\emptyset}}} \cup \dots$;
\item $\Assign{\Obj \in}{{{HF}}} = \Setabs{\tuple{x, y}}{x, y \in
  \Domain{{{HF}}}, x \in y}$.
\end{enumerate}
\end{ex}

\begin{digress}
संरचना म्हणून काय गणले जाईल याविषयी आपण केलेल्या अटींचा आपल्या
तर्कशास्त्रावर परिणाम होतो. उदाहरणार्थ, प्रांत अरिक्त ठेवण्याची निवड
संख्यकीकृत वाक्यांच्या पूर्तिची (किंवा सत्यतेची) नेहमीची मांडणी गृहीत
धरल्यास $\lexists[x][(\varphi(x) \lor \lnot \varphi(x))]$ वैध---म्हणजे तार्किक
सत्य---असल्याची खात्री देते. तसेच, सर्व स्थिरांक नी प्रांतातील
वस्तूच निर्देशित केली पाहिजे ही अट अस्तित्ववाची सामान्यीकरण हा
निर्दोष अनुमान-रूपबंध असल्याची खात्री देते: $\varphi(a)$, म्हणून
$\lexists[x][\varphi(x)]$. नावांना प्रांताबाहेरील वस्तू निर्देशित करण्याची
किंवा काहीच निर्देशित न करण्याची मुभा दिली, तर आपण \emph{मुक्त
तर्कशास्त्राच्या} दिशेने जाऊ; त्यात अस्तित्ववाची सामान्यीकरणासाठी एक
अतिरिक्त आधारवाक्यालागते: $\varphi(a)$ आणि $\lexists[x][\eq[x][a]]$, म्हणून
$\lexists[x][\varphi(x)]$.
\end{digress}












\section{प्रथम-क्रम भाषांसाठी बंद पदांनी आच्छादित संरचना}\label{fol:syn:cov:sec}

\begin{explain}
एखाद्या पदात चले नसतील, तर ते \emph{बंद} असते, हे आठवा.
\end{explain}

\begin{defn}[बंद पदांचे मूल्य]
$t$ हे भाषा~$\Lang L$ हिचे बंद पद आणि $\Struct M$ ही~$\Lang L$ साठीची
संरचना असेल, तर \emph{मूल्य}~$\Value{t}{M}$ पुढीलप्रमाणे
परिभाषित केले जाते:
\begin{enumerate}
\item $t$ हे फक्त स्थिरांक~$c$ असेल, तर
  $\Value{c}{M} = \Assign{c}{M}$.
\item $t$ हे $\Atom{f}{t_1, \ldots, t_n}$ या रूपात असेल, तर
  \[
  \Value{t}{M} = \Assign{f}{M}(\Value{t_1}{M}, \ldots,
  \Value{t_n}{M}).
  \]
\end{enumerate}
\end{defn}

\begin{defn}[आच्छादित संरचना]
प्रांतातील प्रत्येक घटक हा एखाद्या बंद पदाचे मूल्य असेल, तर ती
संरचना \emph{आच्छादित} असते.
\end{defn}

\begin{ex}
स्थिरांक $\Obj{zero}$, $\Obj{one}$, $\Obj{two}$, \dots, द्विस्थानी
विधेय~$<$ आणि द्विस्थानी फलने $+$ व~$\times$ असलेली
भाषा~$\Lang L$ घेऊ. मग~$\Lang L$ साठीची संरचना~$\Struct M$ अशी
आहे: तिचा प्रांत $\Domain M = \{0, 1, 2, \ldots \}$ असून
$\Assign{\Obj{zero}}{M} = 0$, $\Assign{\Obj{one}}{M} = 1$,
$\Assign{\Obj{two}}{M} = 2$ आणि असेच पुढे ही तिची अर्थनिर्धारणे आहेत.
द्विस्थानी संबंधचिन्ह~$<$ साठी $\Assign{<}{M}$ हा अशा सर्व जोड्यांचा
संच आहे की $\tuple{c_1, c_2} \in \Domain{M}^2$ आणि $c_1$ हा~$c_2$ पेक्षा
लहान आहे: उदाहरणार्थ, $\tuple{1, 3} \in \Assign{<}{M}$, पण
$\tuple{2, 2} \notin \Assign{<}{M}$. द्विस्थानी फलन~$+$ साठी
$\Assign{+}{M}$ नेहमीच्या रीतीने परिभाषित करू---उदाहरणार्थ,
$\Assign{+}{M}(2,3)$ हे~$5$ कडे पाठवते; आणि द्विस्थानी
फलन~$\times$ साठीही त्याचप्रमाणे. म्हणून $\Obj{four}$ चे
मूल्य फक्त~$4$ आहे आणि $\times(\Obj{two},
+(\Obj{three},\Obj{zero}))$ चे मूल्य (किंवा मध्यस्थ संकेतपद्धतीत,
$\Obj{two} \times (\Obj{three} + \Obj{zero})$) पुढीलप्रमाणे आहे:

\begin{multline*}
\Value{\times(\Obj{two}, +(\Obj{three},\Obj{zero}))}{M} \\
\begin{aligned}
& =\Assign{\times}{M}(\Value{\Obj{two}}{M}, \Value{+(\Obj{three}, \Obj{zero})}{M})\\
& = \Assign{\times}{M}(\Value{\Obj{two}}{M}, \Assign{+}{M}(\Value{\Obj{three}}{M},
\Value{\Obj{zero}}{M})) \\
& = \Assign{\times}{M}(\Assign{\Obj{two}}{M}, \Assign{+}{M}(\Assign{\Obj{three}}{M},
\Assign{\Obj{zero}}{M})) \\
& = \Assign{\times}{M}(2, \Assign{+}{M}(3, 0)) \\
& = \Assign{\times}{M}(2, 3) \\
& = 6
\end{aligned}
\end{multline*}
\end{ex}

\begin{prob}
अंकगणिताचे मानक प्रतिमान~$\Struct N$ आच्छादित आहे का? स्पष्टीकरण द्या.
\end{prob}












\section{संरचना मधील
  सूत्राची पूर्ति}\label{fol:syn:sat:sec}

\begin{explain}
एकीकडे पदे आणि सूत्रे यांसारख्या अभिव्यक्ती आणि दुसरीकडे
संरचना यांना जोडणाऱ्या मूलभूत संकल्पना म्हणजे पदाचे
\emph{मूल्य} आणि सूत्राची \emph{पूर्ति}. अनौपचारिकपणे,
पदाचे मूल्य म्हणजे संरचना मधील घटक असतो---पद
नुसते स्थिरचिन्ह असेल, तर त्या स्थिरचिन्हाला संरचना ने नेमून
दिलेली वस्तू हे त्याचे मूल्य असते; आणि ते फलने वापरून
बनवलेले असेल, तर पदातील स्थिरचिन्हांची मूल्ये आणि फलनचिन्हांना
नेमून दिलेली फलने यांपासून त्याचे मूल्य काढले जाते. विधेयांना
दिलेल्या अर्थनिर्धारणामुळे संरचनेच्या प्रांतात सूत्र सत्य
होत असेल, तर सूत्राची त्या संरचनांमध्ये \emph{पूर्ति
होते}. पूर्तिची ही संकल्पना विगमनाधारित रीतीने विनिर्दिष्ट केली जाते:
आण्विक सूत्रांची पूर्ति केव्हा होते हे संरचनेचे
विनिर्देशन थेट सांगते; आणि एखाद्या जटिल सूत्राची पूर्ति केव्हा
होते हे आपण तिचा मुख्य संयोजक किंवा संख्यापक आणि तिच्या तात्काळ
उपसूत्रांची पूर्ति होते की नाही यांवरून परिभाषित करतो.

इथे संख्यापकांची स्थिती थोडी अवघड आहे, कारण संख्यकीकृत सूत्र
च्या तात्काळ उपसूत्रामध्ये एक मुक्त चल असते आणि
संरचना ही चलांची मूल्ये विनिर्दिष्ट करत नाहीत.
ही अडचण हाताळण्यासाठी आपण \emph{चर-मूल्यांकने} ही आणतो आणि पूर्तिची
व्याख्या केवळ संरचनेच्या सापेक्ष न देता, संरचना आणि
चल चे मूल्यांकन या दोहोंच्या सापेक्ष देतो.
\end{explain}

\begin{defn}[चर-मूल्यांकन]
संरचना~$\Struct{M}$ हिच्यासाठीचे \emph{चर-मूल्यांकन}~$s$
म्हणजे प्रत्येक चल चे~$\Domain M$ मधील घटक कडे
प्रतिचित्रण करणारे फलन; म्हणजेच, $s\colon \Var \to \Domain M$.
\end{defn}

\begin{explain}
संरचना प्रत्येक स्थिरांक ला मूल्य नेमून देते, तर
चर-मूल्यांकन प्रत्येक चराला मूल्य नेमून देते. पण त्यांच्यापासून बनलेली
पदेही प्रांत मधील घटक ना नावे देण्यासाठी वापरायची आहेत.
त्यासाठी आपण पदांचे मूल्य विगमनाधारित रीतीने परिभाषित करतो.
स्थिरांक आणि चरांसाठी मूल्य अनुक्रमे संरचना किंवा
चर-मूल्यांकन विनिर्दिष्ट करते तेच असते; अधिक जटिल पदांसाठी ते
संरचना ने फलने ना नेमून दिलेली फलने वापरून पुनरावर्ती
रीतीने काढले जाते.
\end{explain}

\begin{defn}[पदांचे मूल्य]
$t$ हे भाषा~$\Lang L$ हिचे पद, $\Struct M$ ही~$\Lang L$ साठीची
संरचना आणि $s$ हे~$\Struct M$ साठीचे चल चे
मूल्यांकन असेल, तर \emph{मूल्य}~$\Value{t}{M}[s]$ पुढीलप्रमाणे
परिभाषित केले जाते:
\begin{enumerate}
\item \indcase{t}{c}{$\Value{\indfrm}{M}[s] = \Assign{\indcomplex}{M}$.}
\item \indcase{t}{x}{$\Value{\indfrm}{M}[s] = s(\indcomplex)$.}
\item \indcase{t}{\Atom{f}{t_1, \ldots, t_n}}{
\[
\Value{\indfrm}{M}[s] = \Assign{f}{M}(\Value{t_1}{M}[s], \ldots,
\Value{t_n}{M}[s]).
\]}
\end{enumerate}
\end{defn}

\begin{defn}[$x$-पर्याय]
$s$ हे संरचना~$\Struct M$ साठीचे चल चे मूल्यांकन
असेल, तर~$s$ पासून फार तर~$x$ ला नेमून दिलेल्या मूल्यातच भिन्न असणारे
$\Struct M$ साठीचे कोणतेही चल चे मूल्यांकन~$s'$ याला~$s$ चा
\emph{$x$-पर्याय} म्हणतात. $s'$ हा~$s$ चा $x$-पर्याय असेल, तर आपण
$\varAssign{s'}{s}{x}$ असे लिहितो.
\end{defn}

\begin{explain}
मूल्यांकन~$s$ च्या $x$-पर्यायाने~$x$ ला काही वेगळेच नेमून दिले
\emph{पाहिजे} असे नाही, हे लक्षात घ्या. वस्तुतः, प्रत्येक मूल्यांकन
स्वतःचाच $x$-पर्याय मानला जातो.
\end{explain}

\begin{defn}
  $s$ हे संरचना~$\Struct M$ साठीचे चल चे मूल्यांकन
  आणि $m \in \Domain{M}$ असेल, तर मूल्यांकन~$\Subst{s}{m}{x}$ हे पुढील
  रीतीने परिभाषित केलेले चर-मूल्यांकन आहे:
  \[\Subst{s}{m}{x}(y) = \begin{cases}
    m & \text{जर } y \ident x\\
    s(y) & \text{अन्यथा}.
  \end{cases}\]
\end{defn}

दुसऱ्या शब्दांत, $\Subst{s}{m}{x}$ हा~$s$ चा असा विशिष्ट $x$-पर्याय
आहे, जो प्रांतातील घटक~$m$ हे~$x$ ला नेमून देतो आणि~$x$ शिवाय
इतर चले ना~$s$ ने नेमून दिलेल्याच गोष्टी नेमून देतो.

\begin{defn}[पूर्ति]
\label{fol:syn:sat:defn:satisfaction}
संरचना~$\Struct M$ मध्ये चल चे मूल्यांकन~$s$ याच्या
सापेक्ष सूत्र~$\varphi$ ची पूर्ति, संकेतांत $\Sat{M}{\varphi}[s]$, पुढीलप्रमाणे
पुनरावर्ती रीतीने परिभाषित केली जाते. (``$\Sat{M}{\varphi}[s]$ नाही'' या
अर्थाने आपण $\Sat/{M}{\varphi}[s]$ असे लिहितो.)
\begin{enumerate}
\item
  \indcase{\varphi}{\lfalse}{$\Sat/{M}{\indfrm}[s]$.}

\item
  \indcase{\varphi}{\ltrue}{$\Sat{M}{\indfrm}[s]$.}

\item \indcase{\varphi}{\Atom{R}{t_1, \dots, t_n}}{$\Sat{M}{\indfrm}[s]$
  तेव्हा आणि केवळ तेव्हाच, जेव्हा $\langle \Value{t_1}{M}[s], \dots,
  \Value{t_n}{M}[s] \rangle \in \Assign{R}{M}$.}

\item \indcase{\varphi}{\eq[t_1][t_2]}{$\Sat{M}{\indfrm}[s]$ तेव्हा आणि
  केवळ तेव्हाच, जेव्हा $\Value{t_1}{M}[s] = \Value{t_2}{M}[s]$.}

\item
  \indcase{\varphi}{\lnot \psi}{$\Sat{M}{\indfrm}[s]$ तेव्हा आणि केवळ तेव्हाच,
    जेव्हा $\Sat/{M}{\psi}[s]$.}

\item
  \indcase{\varphi}{(\psi \land \chi)}{$\Sat{M}{\indfrm}[s]$ तेव्हा आणि केवळ
    तेव्हाच, जेव्हा $\Sat{M}{\psi}[s]$ आणि $\Sat{M}{\chi}[s]$.}

\item
  \indcase{\varphi}{(\psi \lor \chi)}{$\Sat{M}{\indfrm}[s]$ तेव्हा आणि केवळ
    तेव्हाच, जेव्हा $\Sat{M}{\psi}[s]$ किंवा $\Sat{M}{\chi}[s]$ (किंवा
    दोन्ही).}

\item
  \indcase{\varphi}{(\psi \lif \chi)}{$\Sat{M}{\indfrm}[s]$ तेव्हा आणि केवळ
    तेव्हाच, जेव्हा $\Sat/{M}{\psi}[s]$ किंवा $\Sat{M}{\chi}[s]$ (किंवा
    दोन्ही).}

\item
  \indcase{\varphi}{(\psi \liff \chi)}{$\Sat{M}{\indfrm}[s]$ तेव्हा आणि केवळ
    तेव्हाच, जेव्हा एकतर $\Sat{M}{\psi}[s]$ आणि $\Sat{M}{\chi}[s]$ दोन्ही,
    किंवा $\Sat{M}{\psi}[s]$ व $\Sat{M}{\chi}[s]$ यांपैकी एकही नाही.}

\item
  \indcase{\varphi}{\lforall[x][\psi]}{$\Sat{M}{\indfrm}[s]$ तेव्हा आणि केवळ
    तेव्हाच, जेव्हा प्रत्येक घटक~$m \in \Domain M$ साठी
    $\Sat{M}{\psi}[\Subst{s}{m}{x}]$.}

\item
  \indcase{\varphi}{\lexists[x][\psi]}{$\Sat{M}{\indfrm}[s]$ तेव्हा आणि केवळ
  तेव्हाच, जेव्हा किमान एका घटक~$m \in \Domain M$ साठी
  $\Sat{M}{\psi}[\Subst{s}{m}{x}]$.}
\end{enumerate}
\end{defn}

\begin{explain}
शेवटच्या दोन उपवाक्यांत
चर-मूल्यांकने महत्त्वाची आहेत. ``प्रत्येक $m \in
\Domain{M}$ साठी $\Sat{M}{\psi(m)}$'' असे म्हणून आपण
$\lforall[x][\psi(x)]$ ची पूर्ति परिभाषित करू शकत नाही.
 ``किमान एका $m \in \Domain{M}$ साठी $\Sat{M}{\psi(m)}$''
असे म्हणून आपण $\lexists[x][\psi(x)]$ ची पूर्ति परिभाषित करू शकत नाही.
कारण $m \in \Domain M$ असेल, तर ते भाषेचे चिन्ह नसते आणि म्हणून
$\psi(m)$ हे सूत्र नसते (म्हणजे $\Subst{\psi}{m}{x}$ अपरिभाषित
असते).~$\Struct{M}$ मधील प्रत्येक घटक चे नाव देणारी
स्थिरांक किंवा पदे
आपल्याजवळ उपलब्ध आहेत असेही गृहीत धरता येत नाही, कारण संरचना
च्या व्याख्येत तशी कोणतीही अट नाही. प्रमाण भाषेतील स्थिरांकाचा
संच गणनीय अनंत आहे; म्हणून $\Domain{M}$ हा गणनीय नसेल,
तर प्रत्येक वस्तूला नाव देण्यासाठी पुरेशी स्थिरांक सुद्धा नसतात.

ही अडचण सोडवण्यासाठी आपण चल ची मूल्यांकने आणतो; त्यांमुळे
चरे प्रांतातील घटक शी थेट जोडता येतात. मग, उदाहरणार्थ,
$\lexists[x][\psi(x)]$ ची~$\Struct M$
मध्ये पूर्ति तेव्हा आणि केवळ तेव्हाच होते, जेव्हा
किमान एका $m \in \Domain{M}$ साठी $\Sat{M}{\psi(m)}$, असे न म्हणता आपण म्हणतो:
तिची~$\Struct M$ मध्ये~$s$ च्या \emph{सापेक्ष} पूर्ति तेव्हा आणि केवळ
तेव्हाच होते, जेव्हा $\psi(x)$ ची~$\Subst{s}{m}{x}$ च्या सापेक्ष
किमान एका $m \in \Domain M$ साठी पूर्ति होते.

\end{explain}

\begin{ex}
$\Lang{L} = \{a, b, f, R\}$ असू द्या, जिथे $a$ आणि $b$ ही
स्थिरांक, $f$ हे दोन-स्थानी फलन आणि $R$ हे दोन-स्थानी
विधेय आहे. पुढीलप्रमाणे परिभाषित केलेली संरचना~$\Struct{M}$
विचारात घ्या:
\begin{enumerate}
\item $\Domain M = \{1, 2, 3, 4\}$
\item $\Assign{a}{M} = 1$
\item $\Assign{b}{M} = 2$
\item $x+y \le 3$ असेल तर $\Assign{f}{M}(x, y) = x+y$, अन्यथा $= 3$.
\item $\Assign{R}{M} = \{\tuple{1, 1}, \tuple{1, 2}, \tuple{2, 3}, \tuple{2, 4}\}$
\end{enumerate}
प्रत्येक चल ला $1 \in \Domain{M}$ नेमून देणारे फलन
$s(x) = 1$ हे~$\Struct{M}$ साठीचे चर-मूल्यांकन आहे.

मग
\begin{align*}
\Value{f(a,b)}{M}[s] & = \Assign{f}{M}(\Value{a}{M}[s], \Value{b}{M}[s]).
\intertext{$a$ आणि $b$ ही स्थिरांक असल्यामुळे $\Value{a}{M}[s]
  = \Assign{a}{M} = 1$ आणि $\Value{b}{M}[s] = \Assign{b}{M} = 2$. म्हणून}
\Value{f(a,b)}{M}[s] & = \Assign{f}{M}(1, 2) = 1+2 = 3.
\intertext{$f(f(a,b),a)$ चे मूल्य काढण्यासाठी पुढील अभिव्यक्ती विचारात घ्यावी लागते:}
\Value{f(f(a,b),a)}{M}[s] & = \Assign{f}{M}(\Value{f(a, b)}{M}[s],
\Value{a}{M}[s]) = \Assign{f}{M}(3, 1) = 3,
\intertext{कारण $3+1 > 3$. $s(x) = 1$ आणि $\Value{x}{M}[s] =
  s(x)$ असल्यामुळे आपल्याकडे पुढील समीकरणही आहे:}
\Value{f(f(a,b),x)}{M}[s] & = \Assign{f}{M}(\Value{f(a, b)}{M}[s],
\Value{x}{M}[s]) = \Assign{f}{M}(3, 1) = 3,
\end{align*}

आण्विक सूत्र~$R(t_1, t_2)$ च्या फलसाधकांच्या मूल्यांची क्रमित
जोडी, म्हणजे $\tuple{\Value{t_1}{M}[s], \Value{t_2}{M}[s]}$, ही
$\Assign{R}{M}$ ची घटक असेल, तर त्या सूत्राची पूर्ति होते.
म्हणून, उदाहरणार्थ, $\tuple{\Value{b}{M}, \Value{f(a,b)}{M}} =
\tuple{2, 3} \in \Assign{R}{M}$ असल्यामुळे
$\Sat{M}{R(b,f(a,b))}[s]$; पण $\tuple{1, 3} \notin \Assign{R}{M}$
असल्यामुळे $\Sat/{M}{R(x, f(a,b))}[s]$.


अनाण्विक सूत्र~$\varphi$ ची पूर्ति होते का हे ठरवण्यासाठी त्याच्या मुख्य
संयोजकाला लागू होणारे विगमनाधारित व्याख्येतील उपवाक्य वापरावे. उदाहरणार्थ,
$R(a, a) \lif (R(b, x) \lor R(x, b))$ मधील मुख्य संयोजक~$\lif$ आहे आणि
\begin{align*}
  & \Sat{M}{R(a, a) \lif (R(b, x) \lor R(x, b))}[s] \text{ तेव्हा आणि केवळ तेव्हाच, जेव्हा }\\
  & \qquad
  \Sat/{M}{R(a,a)}[s] \text{ किंवा } \Sat{M}{R(b, x) \lor R(x, b)}[s]
  \intertext{$\Sat{M}{R(a,a)}[s]$ असल्यामुळे (कारण $\tuple{1,1} \in
    \Assign{R}{M}$), उत्तर अजून ठरवता येत नाही; प्रथम
    $\Sat{M}{R(b, x) \lor R(x, b)}[s]$ आहे का ते शोधावे लागते:}
  & \Sat{M}{R(b, x) \lor R(x, b)}[s] \text{ तेव्हा आणि केवळ तेव्हाच, जेव्हा }\\
  & \qquad \Sat{M}{R(b,x)}[s] \text{ किंवा } \Sat{M}{R(x, b)}[s]
  \intertext{आणि हे खरे आहे, कारण $\Sat{M}{R(x, b)}[s]$
    (कारण $\tuple{1,2} \in \Assign{R}{M}$).}
\end{align*}

आठवा, $s$ चा $x$-पर्याय म्हणजे~$s$ पासून फार तर~$x$ ला नेमून
दिलेल्या मूल्यातच भिन्न असणारे चर-मूल्यांकन. $\Domain{M}$ च्या प्रत्येक
घटक साठी~$s$ चा एक $x$-पर्याय आहे:
\begin{align*}
  s_1 & = \Subst{s}{1}{x}, &
  s_2 & = \Subst{s}{2}{x},\\
  s_3 & = \Subst{s}{3}{x}, &
  s_4 & = \Subst{s}{4}{x}.
\end{align*}
म्हणून, उदाहरणार्थ, $s_2(x) = 2$ आणि~$x$ शिवाय इतर सर्व चर~$y$ साठी
$s_2(y) = s(y) = 1$. $\Domain{M} = \{1, 2, 3, 4\}$ असल्यामुळे हेच
$\Struct{M}$ या रचनेसाठीचे~$s$ चे सर्व $x$-पर्याय आहेत. विशेषतः,
$s_1 = s$ हे लक्षात घ्या ($s$ हा नेहमीच स्वतःचा $x$-पर्याय असतो).

अस्तित्ववाची संख्यापक असलेल्या
  सूत्र~$\lexists[x][\varphi(x)]$ ची पूर्ति होते का हे ठरवण्यासाठी,
  किमान एका $m \in \Domain M$ साठी
  $\Sat{M}{\varphi(x)}[\Subst{s}{m}{x}]$ आहे का हे ठरवावे लागते. म्हणून,
  \[
  \Sat{M}{\lexists[x][(R(b,x) \lor R(x,b))]}[s],
  \]
  कारण $\Sat{M}{R(b,x) \lor R(x, b)}[\Subst{s}{1}{x}]$
  ($\Subst{s}{3}{x}$ हेसुद्धा चालले असते). पण,
  \[
  \Sat/{M}{\lexists[x][(R(b,x) \land R(x,b))]}[s]
  \]
  कारण कोणताही $m \in \Domain{M}$ निवडला, तरी
  $\Sat/{M}{R(b,x) \land R(x,b)}[\Subst{s}{m}{x}]$.

वैश्विक संख्यापक असलेल्या
  सूत्र~$\lforall[x][\varphi(x)]$ ची पूर्ति होते का हे ठरवण्यासाठी,
  प्रत्येक $m \in \Domain M$ साठी
  $\Sat{M}{\varphi(x)}[\Subst{s}{m}{x}]$ आहे का हे ठरवावे लागते. म्हणून,
  \[
  \Sat{M}{\lforall[x][(R(x,a) \lif R(a,x))]}[s],
  \]
  कारण प्रत्येक $m \in \Domain M$ साठी
  $\Sat{M}{R(x,a) \lif R(a,x)}[\Subst{s}{m}{x}]$. $m = 1$ साठी
  $\Sat{M}{R(a,x)}[\Subst{s}{1}{x}]$ असल्यामुळे उत्तरवर्ती सत्य आहे;
  $m = 2$, $3$ आणि~$4$ साठी $\Sat/{M}{R(x,a)}[\Subst{s}{m}{x}]$
  असल्यामुळे पूर्ववर्ती असत्य आहे. पण,
  \[
  \Sat/{M}{\lforall[x][(R(a,x) \lif R(x,a))]}[s]
  \]
  कारण $\Sat/{M}{R(a,x) \lif R(x,a)}[\Subst{s}{2}{x}]$ (कारण
  $\Sat{M}{R(a, x)}[\Subst{s}{2}{x}]$ आणि $\Sat/{M}{R(x,
  a)}[\Subst{s}{2}{x}]$).





अधिक जटिल उदाहरणासाठी पुढील सूत्र विचारात घ्या:
\[
\lforall[x][(R(a,x) \lif \lexists[y][R(x,y)])].
\]
$\Sat/{M}{R(a,x)}[\Subst{s}{3}{x}]$ आणि
$\Sat/{M}{R(a,x)}[\Subst{s}{4}{x}]$ असल्यामुळे, सशर्त सूत्राच्या
उत्तरवर्तीची काळजी घ्यावी लागणारी रंजक प्रकरणे फक्त $m = 1$ आणि
$m = 2$ ही आहेत. $\Sat{M}{\lexists[y][R(x,y)]}[\Subst{s}{1}{x}]$ आहे
का? किमान एक $n \in \Domain M$ असा असेल की
$\Sat{M}{R(x,y)}[\Subst{\Subst{s}{1}{x}}{n}{y}]$, तर ते खरे आहे.
वस्तुतः, $n = 1$ घेतल्यास $\Subst{\Subst{s}{1}{x}}{n}{y} =
\Subst{s}{1}{y} = s$. $s(x) = 1$, $s(y) = 1$ आणि
$\tuple{1,1} \in \Assign{R}{M}$ असल्यामुळे उत्तर होकारार्थी आहे.


$\Sat{M}{\lexists[y][R(x,y)]}[\Subst{s}{2}{x}]$ आहे का हे ठरवण्यासाठी
आपल्याला चल ची मूल्यांकने
$\Subst{\Subst{s}{2}{x}}{n}{y}$ तपासावी लागतात. इथे $n = 1$ साठी हे
मूल्यांकन $s_2 = \Subst{s}{2}{x}$ आहे आणि ते $R(x,y)$ ची पूर्ति करत
नाही ($s_2(x) = 2$, $s_2(y) = 1$ आणि
$\tuple{2,1}\notin \Assign{R}{M}$). पण
$\Subst{\Subst{s}{2}{x}}{3}{y} = \Subst{s_2}{3}{y}$ विचारात घ्या.
$\tuple{2,3} \in \Assign{R}{M}$ असल्यामुळे
$\Sat{M}{R(x,y)}[\Subst{s_2}{3}{y}]$ आणि म्हणून
$\Sat{M}{\lexists[y][R(x,y)]}[s_2]$.

म्हणून, प्रत्येक $m \in \Domain M$ साठी एकतर
$\Sat/{M}{R(a,x)}[\Subst{s}{m}{x}]$ ($m = 3$, $4$ असेल तर) किंवा
$\Sat{M}{\lexists[y][R(x,y)]}[\Subst{s}{m}{x}]$ ($m = 1$, $2$ असेल तर),
आणि म्हणून
\[
\Sat{M}{\lforall[x][(R(a,x) \lif \lexists[y][R(x,y)])]}[s].
\]
दुसरीकडे,
\[
\Sat/{M}{\lexists[x][(R(a,x) \land \lforall[y][R(x,y)])]}[s].
\]
फक्त $m = 1$ आणि $m = 2$ यांच्यासाठी
$\Sat{M}{R(a,x)}[\Subst{s}{m}{x}]$. पण~$m$ च्या या दोन्ही मूल्यांपैकी
प्रत्येकासाठी एक $n \in \Domain M$ आहे---पहिल्या मूल्यासाठी $n = 4$
आणि दुसऱ्यासाठी $n = 1$---ज्यामुळे
$\Sat/{M}{R(x,y)}[\Subst{\Subst{s}{m}{x}}{n}{y}]$ आणि म्हणून $m = 1$
व $m = 2$ साठी
$\Sat/{M}{\lforall[y][R(x,y)]}[\Subst{s}{m}{x}]$. एकंदरीत, असा कोणताही
$m \in \Domain M$ नाही की
$\Sat{M}{R(a,x) \land \lforall[y][R(x,y)]}[\Subst{s}{m}{x}]$.


\end{ex}









\begin{prob}
एक स्थिरांक, एक एक-स्थानी फलन आणि एक दोन-स्थानी
विधेय असलेली $\Lang L = \{c, f, A\}$ असू द्या; आणि
संरचना~$\Struct{M}$ पुढीलप्रमाणे दिलेली असू द्या:
\begin{enumerate}
\item $\Domain M = \{1, 2, 3\}$
\item $\Assign{c}{M} = 3$
\item $\Assign{f}{M}(1) = 2, \Assign{f}{M}(2) = 3, \Assign{f}{M}(3) = 2$
\item $\Assign{A}{M} = \{\tuple{1, 2}, \tuple{2, 3}, \tuple{3, 3}\}$
\end{enumerate}
(a) सर्व चले~$v$ साठी $s(v) = 1$ असू द्या. पुढील विधान
खरे आहे की नाही हे शोधा:
\[
\Sat{M}{\lexists[x][(A(f(z), c) \lif \lforall[y][(A(y, x) \lor A(f(y),
      x))])]}[s]
\]
तुमचे उत्तर का योग्य आहे ते स्पष्ट करा.

(b) ज्यात या सूत्राची पूर्ति होत नाही अशी वेगळी रचना आणि
चल चे मूल्यांकन द्या.
\end{prob}












\section{चर-मूल्यांकने}\label{fol:syn:ass:sec}

\begin{explain}
चल चे मूल्यांकन~$s$ हे \emph{प्रत्येक} चराला मूल्य देते---आणि
चरे अनंत असतात. अर्थात, हे आवश्यक नाही. आपण चल च्या
मूल्यांकनांनी सर्व चले ना मूल्ये नेमावीत अशी अट घालतो, कारण
त्यामुळे गोष्टी खूप सोप्या होतात. पद~$t$ चे मूल्य आणि सूत्र~$\varphi$
ची संरचनांमध्ये~$s$ च्या सापेक्ष पूर्ति होते की नाही, हे फक्त
$t$ मध्ये येणाऱ्या चले ना आणि~$\varphi$ च्या मुक्त चले ना
$s$ ने केलेल्या नेमणुकांवर अवलंबून असते. पुढील दोन विधानांचा आशय हाच
आहे. “अवलंबून असणे” ही कल्पना नेमकी करण्यासाठी आपण दोन गोष्टी दाखवतो:
$t$ मधील सर्व चरांवर एकमत असणारी कोणतीही दोन चर-मूल्यांकने तेच मूल्य
देतात; आणि दोन चर-मूल्यांकने~$\varphi$ च्या सर्व मुक्त चरांवर एकमत असतील, तर
$\varphi$ ची एकाच्या सापेक्ष पूर्ति होते तेव्हा आणि केवळ तेव्हाच तिची
दुसऱ्याच्या सापेक्ष पूर्ति होते.
\end{explain}

\begin{prop}\label{fol:syn:ass:prop:valindep}
पद~$t$ मधील चले ही $x_1$, \dots,~$x_n$ यांपैकी असतील आणि
$i = 1$, \dots,~$n$ साठी $s_1(x_i) = s_2(x_i)$ असेल, तर
$\Value{t}{M}[s_1] = \Value{t}{M}[s_2]$.
\end{prop}

\begin{proof}
$t$ च्या जटिलतेवरील विगमनाने सिद्ध करू. आधार बाबतीत~$t$ हे
स्थिरांक किंवा $x_1$, \dots,~$x_n$ यांपैकी एखादे चर असू शकते.
$t = c$ असेल, तर
$\Value{t}{M}[s_1] = \Assign{c}{M} = \Value{t}{M}[s_2]$. $t = x_i$
असेल, तर विधानाच्या गृहीतकानुसार $s_1(x_i) = s_2(x_i)$; म्हणून
$\Value{t}{M}[s_1] = s_1(x_i) = s_2(x_i) = \Value{t}{M}[s_2]$.

विगमन पायरीसाठी $t = \Atom{f}{t_1, \dots, t_k}$ असून दावा
$t_1$, \dots, $t_k$ साठी खरा आहे असे समजू. मग
\begin{align*}
  \Value{t}{M}[s_1] & = \Value{\Atom{f}{t_1, \dots, t_k}}{M}[s_1] \\
  & = \Assign{f}{M}(\Value{t_1}{M}[s_1], \dots, \Value{t_k}{M}[s_1]).
\intertext{$j = 1$, \dots,~$k$ साठी~$t_j$ मधील चले ही
    $x_1$, \dots,~$x_n$ यांपैकी आहेत. विगमन गृहीतकानुसार
    $\Value{t_j}{M}[s_1] = \Value{t_j}{M}[s_2]$. म्हणून,}
\Value{t}{M}[s_1] & = \Value{\Atom{f}{t_1, \dots, t_k}}{M}[s_1] \\
  & = \Assign{f}{M}(\Value{t_1}{M}[s_1], \dots, \Value{t_k}{M}[s_1]) \\
  & = \Assign{f}{M}(\Value{t_1}{M}[s_2], \dots, \Value{t_k}{M}[s_2]) \\
  & = \Value{\Atom{f}{t_1, \dots, t_k}}{M}[s_2] = \Value{t}{M}[s_2].
\end{align*}
\end{proof}

\begin{prop}\label{fol:syn:ass:prop:satindep}
$\varphi$ मधील मुक्त चले ही $x_1$, \dots,~$x_n$ यांपैकी असतील आणि
$i = 1$, \dots,~$n$ साठी $s_1(x_i) = s_2(x_i)$ असेल, तर
$\Sat{M}{\varphi}[s_1]$ तेव्हा आणि केवळ तेव्हाच, जेव्हा
$\Sat{M}{\varphi}[s_2]$.
\end{prop}

\begin{proof}
$\varphi$ च्या जटिलतेवर विगमन वापरू. आधार बाबतीत~$\varphi$ हे आण्विक असते आणि
$\varphi$ पुढीलपैकी असू शकते:
$\ltrue$,
$\lfalse$,
$k$-स्थानी विधेय~$R$ आणि पदे $t_1$, \dots,~$t_k$ यांसाठी
$\Atom{R}{t_1, \dots, t_k}$; किंवा पदे~$t_1$ व~$t_2$ यांसाठी
$\eq[t_1][t_2]$. शेवटच्या दोन बाबतींत आपण द्विशर्त ची फक्त
पुढची दिशा दाखवतो, कारण उलट दिशेची सिद्धता सममित आहे.

\begin{enumerate}
\item
  \indcase{\varphi}{\ltrue}{$\Sat{M}{\indfrm}[s_1]$ आणि
    $\Sat{M}{\indfrm}[s_2]$ दोन्ही.}

\item
  \indcase{\varphi}{\lfalse}{$\Sat/{M}{\varphi}[s_1]$ आणि
    $\Sat/{M}{\varphi}[s_2]$ दोन्ही.}

\item
  \indcase{\varphi}{\Atom{R}{t_1, \ldots, t_k}}{
    $\Sat{M}{\indfrm}[s_1]$ असे समजू. मग
    \[
    \langle \Value{t_1}{M}[s_1], \ldots, \Value{t_k}{M}[s_1] \rangle
    \in \Assign{R}{M}.
    \]
    $i = 1$, \dots,~$k$ साठी \ref{fol:syn:ass:prop:valindep} नुसार
    $\Value{t_i}{M}[s_1] = \Value{t_i}{M}[s_2]$. म्हणून आपल्याकडे
    $\langle \Value{t_1}{M}[s_2], \ldots, \Value{t_k}{M}[s_2] \rangle
    \in \Assign{R}{M}$ हेसुद्धा आहे; आणि म्हणून
    $\Sat{M}{\indfrm}[s_2]$.
    }
\item
  \indcase{\varphi}{\eq[t_1][t_2]}{$\Sat{M}{\indfrm}[s_1]$ असे समजू.
    मग $\Value{t_1}{M}[s_1] = \Value{t_2}{M}[s_1]$. म्हणून,
    \begin{align*}
      \Value{t_1}{M}[s_2] & = \Value{t_1}{M}[s_1]
      & \text{(\ref{fol:syn:ass:prop:valindep} नुसार)} \\
      & = \Value{t_2}{M}[s_1]
      & \text{(कारण $\Sat{M}{\eq[t_1][t_2]}[s_1]$)}\\
      &= \Value{t_2}{M}[s_2]
      & \text{(\ref{fol:syn:ass:prop:valindep} नुसार),}
    \end{align*}
    म्हणून $\Sat{M}{\eq[t_1][t_2]}[s_2]$.}
\end{enumerate}

आता~$\varphi$ पेक्षा कमी जटिल असलेल्या सर्व सूत्रे~$\psi$ साठी
$\Sat{M}{\psi}[s_1]$ तेव्हा आणि केवळ तेव्हाच, जेव्हा
$\Sat{M}{\psi}[s_2]$, असे समजू. विगमन पायरी~$\varphi$ च्या मुख्य कारकाने
ठरणाऱ्या बाबींनुसार पुढे जाते. प्रत्येक बाबतीत आपण द्विशर्त ची
फक्त पुढची दिशा दाखवतो; उलट दिशेची सिद्धता सममित आहे. संख्यापकांच्या
बाबी सोडून इतर सर्व बाबतींत आपण~$\varphi$ च्या उप-सूत्रे~$\psi$ ना
विगमन गृहीतक लावतो. $\psi$ ची मुक्त चरे ही~$\varphi$ च्या मुक्त चरांपैकी
असतात. त्यामुळे $s_1$ व~$s_2$ ही~$\varphi$ च्या मुक्त चरांवर एकमत असतील,
तर ती~$\psi$ च्या मुक्त चरांवरही एकमत असतात आणि विगमन गृहीतक~$\psi$ ला
लागू होते.

\begin{enumerate}
\item

    \indcase{\varphi}{\lnot \psi}{$\Sat{M}{\indfrm}[s_1]$ असेल, तर
      $\Sat/{M}{\psi}[s_1]$; म्हणून विगमन गृहीतकानुसार
      $\Sat/{M}{\psi}[s_2]$ आणि त्यामुळे $\Sat{M}{\indfrm}[s_2]$.}

\item

    \indcase{\varphi}{\psi \land \chi}{$\Sat{M}{\indfrm}[s_1]$ असेल, तर
      $\Sat{M}{\psi}[s_1]$ आणि $\Sat{M}{\chi}[s_1]$; म्हणून विगमन
      गृहीतकानुसार $\Sat{M}{\psi}[s_2]$ आणि $\Sat{M}{\chi}[s_2]$.
      त्यामुळे $\Sat{M}{\indfrm}[s_2]$.}

\item

    \indcase{\varphi}{\psi \lor \chi}{$\Sat{M}{\indfrm}[s_1]$ असेल, तर
      $\Sat{M}{\psi}[s_1]$ किंवा $\Sat{M}{\chi}[s_1]$. विगमन गृहीतकानुसार
      $\Sat{M}{\psi}[s_2]$ किंवा $\Sat{M}{\chi}[s_2]$; म्हणून
      $\Sat{M}{\indfrm}[s_2]$.}

\item

    \indcase{\varphi}{\psi \lif \chi}{$\Sat{M}{\indfrm}[s_1]$ असेल, तर
    $\Sat/{M}{\psi}[s_1]$ किंवा $\Sat{M}{\chi}[s_1]$. विगमन गृहीतकानुसार
    $\Sat/{M}{\psi}[s_2]$ किंवा $\Sat{M}{\chi}[s_2]$; म्हणून
    $\Sat{M}{\indfrm}[s_2]$.}

\item

    \indcase{\varphi}{\psi \liff \chi}{$\Sat{M}{\indfrm}[s_1]$ असेल, तर एकतर
      $\Sat{M}{\psi}[s_1]$ आणि $\Sat{M}{\chi}[s_1]$, किंवा
      $\Sat/{M}{\psi}[s_1]$ आणि $\Sat/{M}{\chi}[s_1]$. विगमन गृहीतकानुसार
      एकतर $\Sat{M}{\psi}[s_2]$ आणि $\Sat{M}{\chi}[s_2]$, किंवा
      $\Sat/{M}{\psi}[s_2]$ आणि $\Sat/{M}{\chi}[s_2]$. दोन्हींपैकी
      कोणतीही बाब असली, तरी $\Sat{M}{\indfrm}[s_2]$.}

\item

    \indcase{\varphi}{\lexists[x][\psi]}{$\Sat{M}{\indfrm}[s_1]$ असेल, तर असा
      एक $m \in \Domain{M}$ असतो की $\Sat{M}{\psi}[\Subst{s_1}{m}{x}]$.
      $s_1' = \Subst{s_1}{m}{x}$ आणि $s_2' =
      \Subst{s_2}{m}{x}$ असू द्या. $\psi$ ची मुक्त चरे
      $x_1$, \dots, $x_n$ आणि~$x$ यांपैकी आहेत. $s_1'$ आणि~$s_2'$ हे
      अनुक्रमे $s_1$ व~$s_2$ यांचे $x$-पर्याय असल्यामुळे आणि गृहीतकानुसार
      $s_1(x_i) = s_2(x_i)$ असल्यामुळे $s_1'(x_i) = s_2'(x_i)$.
      $s_1'$ आणि~$s_2'$ आपण ज्या रीतीने परिभाषित केली आहे, तिच्यामुळे
      $s_1'(x) = s_2'(x) = m$. मग~$\psi$ आणि $s_1'$, $s_2'$ यांना विगमन
      गृहीतक लागू होते; म्हणून $\Sat{M}{\psi}[s_2']$. $s_2' =
      \Subst{s_2}{m}{x}$ असल्यामुळे असा एक $m \in \Domain{M}$ आहे की
      $\Sat{M}{\psi}[\Subst{s_2}{m}{x}]$; म्हणून
      $\Sat{M}{\indfrm}[s_2]$.}

\item

    \indcase{\varphi}{\lforall[x][\psi]}{$\Sat{M}{\indfrm}[s_1]$ असेल, तर
      प्रत्येक $m \in \Domain{M}$ साठी
      $\Sat{M}{\psi}[\Subst{s_1}{m}{x}]$. प्रत्येक
      $m \in \Domain{M}$ साठी $\Sat{M}{\psi}[\Subst{s_2}{m}{x}]$
      हेसुद्धा दाखवायचे आहे. म्हणून मनमाना $m \in \Domain{M}$ घेऊन
      $s_1' = \Subst{s_1}{m}{x}$ आणि $s_2' =
      \Subst{s_2}{m}{x}$ विचारात घ्या. आपल्याकडे
      $\Sat{M}{\psi}[s_1']$. $\psi$ ची मुक्त चरे $x_1$, \dots, $x_n$
      आणि~$x$ यांपैकी आहेत. $s_1'$ आणि~$s_2'$ हे अनुक्रमे $s_1$
      व~$s_2$ यांचे $x$-पर्याय असल्यामुळे आणि गृहीतकानुसार
      $s_1(x_i) = s_2(x_i)$ असल्यामुळे $s_1'(x_i) = s_2'(x_i)$.
      $s_1'$ आणि~$s_2'$ आपण ज्या रीतीने परिभाषित केली आहे, तिच्यामुळे
      $s_1'(x) = s_2'(x) = m$. मग~$\psi$ आणि~$s_1'$, $s_2'$ यांना विगमन
      गृहीतक लागू होते आणि आपल्याकडे $\Sat{M}{\psi}[s_2']$. हे प्रत्येक
      $m \in \Domain{M}$ ला लागू होते, म्हणजे सर्व~$m \in \Domain{M}$
      साठी $\Sat{M}{\psi}[\Subst{s_2}{m}{x}]$; म्हणून
      $\Sat{M}{\indfrm}[s_2]$.
      }
\end{enumerate}
विगमनाने आपल्याला पुढील निष्कर्ष मिळतो: $\varphi$ मधील मुक्त चले
ही $x_1$, \dots, $x_n$ यांपैकी असतील आणि $i = 1$, \dots,~$n$ साठी
$s_1(x_i)=s_2(x_i)$ असेल, तेव्हा $\Sat{M}{\varphi}[s_1]$ तेव्हा आणि केवळ
तेव्हाच, जेव्हा $\Sat{M}{\varphi}[s_2]$.
\end{proof}



\begin{explain}
वाक्ये मध्ये मुक्त चरे नसतात; म्हणून कोणतीही दोन चर-मूल्यांकने
कोणत्याही वाक्याच्या सर्व (शून्य) मुक्त चरांना त्याच गोष्टी नेमून देतात.
आत्ताच सिद्ध केलेल्या विधानाचा अर्थ मग असा होतो की वाक्याची
एखाद्या रचनेत चर-मूल्यांकनाच्या सापेक्ष पूर्ति होते की नाही, हे पूर्णपणे
त्या मूल्यांकनापासून स्वतंत्र असते. हे तथ्य आपण नोंदवू. पुढे येणारी
वाक्याची संरचना मधील पूर्तिची व्याख्या (चर-मूल्यांकनाचा
उल्लेख न करता) त्यामुळे समर्थित होते.
\end{explain}

\begin{cor}
\label{fol:syn:ass:cor:sat-sentence}
$\varphi$ हे वाक्य आणि $s$ हे चर-मूल्यांकन असेल, तर प्रत्येक
चर-मूल्यांकन~$s'$ साठी $\Sat{M}{\varphi}[s]$ तेव्हा आणि केवळ तेव्हाच,
जेव्हा $\Sat{M}{\varphi}[s']$.
\end{cor}

\begin{proof}
  $s'$ हे कोणतेही चर-मूल्यांकन असू द्या. $\varphi$ हे वाक्य
  असल्यामुळे त्यात मुक्त चरे नाहीत; म्हणून प्रत्येक चर-मूल्यांकन~$s'$
  हे~$\varphi$ च्या सर्व मुक्त चरांना~$s$ सारख्याच गोष्टी आपोआप नेमून देते.
  त्यामुळे \ref{fol:syn:ass:prop:satindep} ची अट पूर्ण होते आणि
  $\Sat{M}{\varphi}[s]$ तेव्हा आणि केवळ तेव्हाच, जेव्हा
  $\Sat{M}{\varphi}[s']$.
\end{proof}

\begin{defn}
\label{fol:syn:ass:defn:satisfaction}
$\varphi$ हे वाक्य असेल, तर संरचना~$\Struct M$ ही~$\varphi$ ची
\emph{पूर्ति करते}, $\Sat{M}{\varphi}$, असे आपण तेव्हा आणि केवळ तेव्हाच
म्हणतो, जेव्हा सर्व चर-मूल्यांकने~$s$ साठी $\Sat{M}{\varphi}[s]$.
\end{defn}

$\Sat{M}{\varphi}$ असेल, तर आपण सरळ \emph{$\varphi$ हे~$\Struct{M}$ मध्ये सत्य
आहे} असेही म्हणतो. पूर्तिची संकल्पना स्वतंत्र वाक्यांपासून
वाक्यांच्या संचांपर्यंत स्वाभाविकपणे विस्तारते.

\begin{defn}
\label{fol:syn:ass:defn:sat}
$\Gamma$ हा वाक्यांचा संच असेल, तर संरचना~$\Struct M$
ही~$\Gamma$ ची \emph{पूर्ति करते}, $\Sat{M}{\Gamma}$, असे आपण तेव्हा
आणि केवळ तेव्हाच म्हणतो, जेव्हा प्रत्येक $\varphi \in \Gamma$ साठी
$\Sat{M}{\varphi}$.

\end{defn}

\begin{prop}\label{fol:syn:ass:prop:sentence-sat-true}
  $\Struct{M}$ ही संरचना, $\varphi$ हे वाक्य आणि $s$ हे
  चर-मूल्यांकन असू द्या. $\Sat{M}{\varphi}$ तेव्हा आणि केवळ तेव्हाच,
  जेव्हा $\Sat{M}{\varphi}[s]$.
\end{prop}

\begin{proof}
सराव.
\end{proof}

\begin{prob}
\ref{fol:syn:ass:prop:sentence-sat-true} सिद्ध करा.
\end{prob}

\begin{prop}\label{fol:syn:ass:prop:sat-quant}
  $\varphi(x)$ मध्ये फक्त~$x$ मुक्त येतो आणि $\Struct M$ ही
  संरचना आहे असे समजा. मग:
  \begin{enumerate}
    \item $\Sat{M}{\lexists[x][\varphi(x)]}$ तेव्हा आणि केवळ
      तेव्हाच, जेव्हा किमान एका चर-मूल्यांकन~$s$ साठी
      $\Sat{M}{\varphi(x)}[s]$.
    \item $\Sat{M}{\lforall[x][\varphi(x)]}$ तेव्हा आणि केवळ
      तेव्हाच, जेव्हा सर्व चर-मूल्यांकने~$s$ साठी
      $\Sat{M}{\varphi(x)}[s]$.
\end{enumerate}
\end{prop}

\begin{proof}
सराव.
\end{proof}

\begin{prob}
\ref{fol:syn:ass:prop:sat-quant} सिद्ध करा.
\end{prob}

\begin{prob}
\DeclareRobustCommand{\VDash}{\mathrel{||}\joinrel\Relbar}
$\Lang L$ ही फलने नसलेली भाषा आहे असे समजा. संरचना~$\Struct M$,
स्थिरांक~$c$ आणि $a \in \Domain M$ दिले असता, $\Struct M[a/c]$
ही~$\Struct M$ सारखीच संरचना आहे, फक्त
$\Assign{c}{M[a/c]} = a$, अशी व्याख्या करा. वाक्ये~$\varphi$ साठी
$\Struct M \VDash \varphi$ ची पुढीलप्रमाणे व्याख्या करा:
\begin{enumerate}
\item
  \indcase{\varphi}{\lfalse}{$\Struct{M} \VDash \indfrm$ नाही.}

\item
  \indcase{\varphi}{\ltrue}{$\Struct{M} \VDash \indfrm$.}

\item \indcase{\varphi}{\Atom{R}{d_1, \dots, d_n}}{$\Struct M \VDash \indfrm$
  तेव्हा आणि केवळ तेव्हाच, जेव्हा
  $\langle \Assign{d_1}{M}, \dots, \Assign{d_n}{M} \rangle \in
  \Assign{R}{M}$.}

\item \indcase{\varphi}{\eq[d_1][d_2]}{$\Struct M \VDash \indfrm$ तेव्हा
  आणि केवळ तेव्हाच, जेव्हा
  $\Assign{d_1}{M} = \Assign{d_2}{M}$.}

\item
  \indcase{\varphi}{\lnot \psi}{$\Struct{M} \VDash \indfrm$ तेव्हा आणि केवळ
    तेव्हाच, जेव्हा $\Struct{M} \VDash {\psi}$ नाही.}

\item
  \indcase{\varphi}{(\psi \land \chi)}{$\Struct{M} \VDash
    \indfrm$ तेव्हा आणि केवळ तेव्हाच, जेव्हा
    $\Struct{M} \VDash\psi$ आणि $\Struct{M} \VDash \chi$.}

\item
  \indcase{\varphi}{(\psi \lor \chi)}{$\Struct{M} \VDash \indfrm$ तेव्हा आणि
    केवळ तेव्हाच, जेव्हा $\Struct{M} \VDash \psi$ किंवा
    $\Struct{M} \VDash \chi$ (किंवा दोन्ही).}

\item
  \indcase{\varphi}{(\psi \lif \chi)}{$\Struct{M} \VDash \indfrm$ तेव्हा आणि
    केवळ तेव्हाच, जेव्हा $\Struct {M} \VDash \psi$ नाही किंवा
    $\Struct M \VDash \chi$ (किंवा दोन्ही).}

\item
  \indcase{\varphi}{(\psi \liff \chi)}{$\Struct{M} \VDash {\indfrm}$ तेव्हा आणि
    केवळ तेव्हाच, जेव्हा एकतर $\Struct{M} \VDash {\psi}$ आणि
    $\Struct{M} \VDash {\chi}$ दोन्ही, किंवा $\Struct{M} \VDash {\psi}$ व
    $\Struct{M} \VDash {\chi}$ यांपैकी एकही नाही.}

\item
  \indcase{\varphi}{\lforall[x][\psi]}{$\Struct{M} \VDash {\indfrm}$ तेव्हा आणि
    केवळ तेव्हाच, जेव्हा प्रत्येक $a \in \Domain{M}$ साठी
    $\Struct{M[a/c]} \VDash \Subst{\psi}{c}{x}$, जर~$\psi$ मध्ये~$c$ येत
    नसेल.}

\item
  \indcase{\varphi}{\lexists[x][\psi]}{$\Struct{M} \VDash
  {\indfrm}$ तेव्हा आणि केवळ तेव्हाच, जेव्हा असा एक $a \in \Domain M$
  आहे की $\Struct{M[a/c]} \VDash \Subst{\psi}{c}{x}$, जर~$\psi$ मध्ये~$c$
  येत नसेल.}
\end{enumerate}
$x_1$, \dots, $x_n$ ही~$\varphi$ मधील सर्व मुक्त चले,
$c_1$, \dots, $c_n$ ही~$\varphi$ मध्ये नसलेली स्थिरचिन्हे,
$a_1$, \dots, $a_n \in \Domain M$ आणि $s(x_i) = a_i$ असू द्या.

$\Sat{M}{\varphi}[s]$ तेव्हा आणि केवळ तेव्हाच, जेव्हा
$\Struct M[a_1/c_1,\dots,a_n/c_n]
\VDash \Subst{\Subst{\varphi}{c_1}{x_1}\dots}{c_n}{x_n}$, हे दाखवा.

(या प्रश्नावरून असे दिसते की प्रथम-क्रम तर्कशास्त्राची चिन्हार्थमीमांसा
चर-मूल्यांकनांविना देणे शक्य आहे.)
\end{prob}

\begin{prob}
$f$ हे~$\varphi(x,y)$ मध्ये नसलेले फलनचिन्ह आहे असे समजा. अशी
संरचना~$\Struct{M}$ आहे की
$\Sat{M}{\lforall[x][\lexists[y][\varphi(x,y)]]}$, तेव्हा आणि केवळ तेव्हाच
अशी~$\Struct M'$ आहे की
$\Sat{M'}{\lforall[x][\varphi(x,f(x))]}$, हे दाखवा.

(हा प्रश्न स्कोलेमचे प्रमेय म्हणून परिचित असलेल्या प्रमेयाची एक विशेष
बाब आहे; $\lforall[x][\varphi(x,f(x))]$ याला
$\lforall[x][\lexists[y][\varphi(x,y)]]$ चे \emph{स्कोलेम सामान्य रूप}
म्हणतात.)
\end{prob}












\section{विस्तारात्मकता}\label{fol:syn:ext:sec}

\begin{explain}
विस्तारात्मकता, जिला कधीकधी संबद्धता म्हणतात, अनौपचारिकपणे पुढीलप्रमाणे
व्यक्त करता येते: संरचना~$\Struct M$ मध्ये चल च्या
मूल्यांकन~$s$ च्या सापेक्ष सूत्र~$\varphi$ ची पूर्ति होण्यावर परिणाम
करणारे घटक म्हणजे फक्त प्रांताचा आकार आणि~$\varphi$ मध्ये प्रत्यक्ष
येणाऱ्या भाषेच्या घटकांना~$\Struct M$ व~$s$ यांनी केलेल्या नेमणुका.

विस्तारात्मकतेचा एक तात्काळ परिणाम असा आहे: दोन संरचना
$\Struct M$ आणि~$\Struct M'$ यांचा प्रांत एकच असेल आणि वाक्य~$\varphi$ मध्ये
येणाऱ्या भाषेच्या सर्व घटकांबाबत त्या एकमत असतील, तर~$\varphi$ स्वतः सत्य
आहे की नाही याबाबतही $\Struct M$ आणि~$\Struct M'$ एकमत असले पाहिजेत.
\end{explain}

\begin{prop}[विस्तारात्मकता]
  \label{fol:syn:ext:prop:extensionality}
  $\varphi$ हे सूत्र, $\Struct M_1$ व~$\Struct M_2$ या
  $\Domain{M_1} = \Domain{M_2}$ असलेल्या संरचना, आणि~$s$ हे
  $\Domain{M_1} = \Domain{M_2}$ वरील चर-मूल्यांकन असू द्या. $\varphi$ मध्ये
  येणारे प्रत्येक स्थिरांक~$c$, संबंधचिन्ह~$R$ आणि फलन~$f$
  यांसाठी $\Assign{c}{M_1} = \Assign{c}{M_2}$,
  $\Assign{R}{M_1}=\Assign{R}{M_2}$ आणि
  $\Assign{f}{M_1} = \Assign{f}{M_2}$ असेल, तर
  $\Sat{M_1}{\varphi}[s]$ तेव्हा आणि केवळ तेव्हाच, जेव्हा
  $\Sat{M_2}{\varphi}[s]$.
\end{prop}

\begin{proof}
  प्रथम (पद~$t$ वरील विगमनाने) प्रत्येक पदासाठी
  $\Value{t}{M_1}[s] = \Value{t}{M_2}[s]$ हे सिद्ध करा. मग नुकताच
  सिद्ध केलेला दावा विगमनाच्या आधार बाबतीत ($\varphi$ आण्विक असताना) वापरून,
  $\varphi$ वरील विगमनाने विधान सिद्ध करा.
\end{proof}

\begin{prob}
\ref{fol:syn:ext:prop:extensionality} ची सिद्धता सविस्तर पूर्ण करा.
\end{prob}

\begin{cor}[वाक्ये साठी विस्तारात्मकता]
  \label{fol:syn:ext:cor:extensionality-sent}
  $\varphi$ हे वाक्य आणि $\Struct{M_1}$, $\Struct{M_2}$ या
  \ref{fol:syn:ext:prop:extensionality} मधीलप्रमाणे असू द्या. मग
  $\Sat{M_1}{\varphi}$ तेव्हा आणि केवळ तेव्हाच, जेव्हा $\Sat{M_2}{\varphi}$.
\end{cor}

\begin{proof}
\ref{fol:syn:ext:prop:extensionality} आणि \ref{fol:syn:ass:cor:sat-sentence} यांवरून
हे निष्पन्न होते.
\end{proof}

शिवाय, पदाचे मूल्य आणि संरचना ही सूत्राची पूर्ति करते की
नाही, हे फक्त तिच्या उपपदांच्या मूल्यांवर अवलंबून असते.

\begin{prop}\label{fol:syn:ext:prop:ext-terms}
$\Struct M$ ही संरचना, $t$ व~$t'$ ही पदे आणि $s$ हे
चर-मूल्यांकन असू द्या. मग $\Value{\Subst{t}{t'}{x}}{M}[s] =
\Value{t}{M}[\Subst{s}{\Value{t'}{M}[s]}{x}]$.
\end{prop}

\begin{proof}
$t$ वरील विगमनाने सिद्ध करू.
\begin{enumerate}
\item $t$ हे स्थिरचिन्ह, म्हणजे $t\ident c$, असेल, तर
  $\Subst{t}{t'}{x} = c$ आणि $\Value{c}{M}[s] = \Assign{c}{M} =
  \Value{c}{M}[\Subst{s}{\Value{t'}{M}[s]}{x}]$.

\item $t$ हे~$x$ शिवाय दुसरे चर, म्हणजे $t \ident y$, असेल, तर
  $\Subst{t}{t'}{x} = y$ आणि
  $\varAssign{s}{\Subst{s}{\Value{t'}{M}[s]}{x}}{x}$ असल्यामुळे
  $\Value{y}{M}[s] =
  \Value{y}{M}[\Subst{s}{\Value{t'}{M}[s]}{x}]$.

\item $t \ident x$ असेल, तर $\Subst{t}{t'}{x} = t'$. पण
  $\Subst{s}{\Value{t'}{M}[s]}{x}$ च्या व्याख्येनुसार
  $\Value{x}{M}[\Subst{s}{\Value{t'}{M}[s]}{x}] = \Value{t'}{M}[s]$.

\item $t \ident \Atom{f}{t_1,\dots,t_n}$ असेल, तर:
\begin{multline*}
  \Value{\Subst{t}{t'}{x}}{M}[s] \\
  \begin{aligned}[b]
& = \Value{\Atom{f}{\Subst{t_1}{t'}{x}, \dots, \Subst{t_n}{t'}{x}}}{M}[s]\\
& \qquad    \text{$\Subst{t}{t'}{x}$ च्या व्याख्येनुसार}\\
& = \Assign{f}{M}(\Value{\Subst{t_1}{t'}{x}}{M}[s], \dots,
    \Value{\Subst{t_n}{t'}{x}}{M}[s])\\
    & \qquad  \text{$\Value{\Atom{f}{\dots}}{M}[s]$ च्या व्याख्येनुसार}\\
& = \Assign{f}{M}(\Value{t_1}{M}[\Subst{s}{\Value{t'}{M}[s]}{x}], \dots,
   \Value{t_n}{M}[\Subst{s}{\Value{t'}{M}[s]}{x}])\\
& \qquad    \text{विगमन गृहीतकानुसार}\\
& = \Value{t}{M}[\Subst{s}{\Value{t'}{M}[s]}{x}]
    \text{$\Value{\Atom{f}{\dots}}{M}[\Subst{s}{\Value{t'}{M}[s]}{x}]$ च्या व्याख्येनुसार}
  \end{aligned}
\end{multline*}

\end{enumerate}
\end{proof}

\begin{prop}\label{fol:syn:ext:prop:ext-formulas} $\Struct M$ ही
संरचना, $\varphi$ हे सूत्र, $t'$ हे~A मध्ये~x साठी
आदेशनासाठी मुक्त असलेले पद आणि $s$ हे चर-मूल्यांकन असू द्या. मग
$\Sat{M}{\Subst{\varphi}{t'}{x}}[s]$ तेव्हा आणि केवळ तेव्हाच, जेव्हा
$\Sat{M}{\varphi}[\Subst{s}{\Value{t'}{M}[s]}{x}]$.

\end{prop}

\begin{proof}
सराव.
\end{proof}

\begin{prob}
\ref{fol:syn:ext:prop:ext-formulas} सिद्ध करा.
\end{prob}

\begin{explain}
  \ref{fol:syn:ext:prop:ext-terms} आणि \ref{fol:syn:ext:prop:ext-formulas} यांचा
  आशय पुढीलप्रमाणे आहे. आपल्याजवळ पद~$t$ किंवा सूत्र~$\varphi$ आणि
  दुसरे पद~$t'$ असून, $\Subst{t}{t'}{x}$ चे मूल्य किंवा
  $\Subst{\varphi}{t'}{x}$ ची संरचना~$\Struct M$ मध्ये
  चल च्या मूल्यांकन~$s$ च्या सापेक्ष पूर्ति होते की नाही हे
  जाणून घ्यायचे आहे असे समजा. मग आपण प्रथम आदेशन करून नंतर
  $\Struct{M}$ आणि~$s$ यांच्या सापेक्ष मूल्य किंवा पूर्ति विचारात घेऊ
  शकतो; किंवा प्रथम~$\Struct{M}$ मध्ये~$s$ च्या सापेक्ष~$t'$ चे
  मूल्य~$m = \Value{t'}{M}[s]$ ठरवून, चल चे मूल्यांकन
  बदलून~$\Subst{s}{m}{x}$ करू शकतो, आणि नंतर~$\Struct{M}$ व
  $\Subst{s}{m}{x}$ यांच्या सापेक्ष~$t$ चे मूल्य किंवा
  $\Sat{M}{\varphi}[\Subst{s}{m}{x}]$ आहे की नाही हे विचारात घेऊ शकतो.
  आपण यांपैकी कोणतीही पद्धत वापरली, तरी उत्तर एकच असेल, याची हमी
  \ref{fol:syn:ext:prop:ext-terms} आणि \ref{fol:syn:ext:prop:ext-formulas}
  देतात.
\end{explain}











\section{चिन्हार्थविषयक संकल्पना}\label{fol:syn:sem:sec}

\begin{explain}
प्रथम-क्रम भाषांसाठीच्या संरचनेची व्याख्या मिळाल्यावर वाक्यांचे
काही मूलभूत चिन्हार्थविषयक गुणधर्म आणि वाक्यांमधील संबंध परिभाषित करता
येतात. त्यांतील सर्वांत सोपी संकल्पना म्हणजे वाक्याची \emph{वैधता}.
प्रत्येक संरचनांमध्ये एखाद्या वाक्याची पूर्ति होत असेल, तर ते
वाक्य वैध असते. त्यातील अतार्किक चिन्हांचा अर्थ कसा लावला आहे याचा
विचार न करता ज्या वाक्यांची पूर्ति होते, तीच वैध वाक्ये होत. म्हणून वैध
वाक्यांना \emph{तार्किक सत्ये} असेही म्हणतात---ती कोणत्याही
संरचनांमध्ये सत्य, म्हणजे पूर्त, असतात. त्यामुळे त्यांची सत्यता
फक्त त्यांत येणारी तार्किक चिन्हे आणि त्यांची विन्यासात्मक संरचना
यांवर अवलंबून असते; अतार्किक चिन्हे किंवा त्यांचे अर्थनिर्धारण यांवर नाही.
\end{explain}

\begin{defn}[वैधता]
वाक्य~$\varphi$ हे \emph{वैध}, $\Entails \varphi$, तेव्हा आणि केवळ तेव्हाच असते,
जेव्हा प्रत्येक संरचना~$\Struct M$ साठी $\Sat{M}{\varphi}$.
\end{defn}

\begin{defn}[तार्किक निष्पन्नता]
वाक्यांच्या संच~$\Gamma$ मधून वाक्य~$\varphi$
\emph{तार्किकरीत्या निष्पन्न होते}, $\Gamma \Entails \varphi$, तेव्हा आणि
केवळ तेव्हाच, जेव्हा $\Sat{M}{\Gamma}$ असलेल्या प्रत्येक
संरचना~$\Struct M$ साठी $\Sat{M}{\varphi}$.
\end{defn}


\begin{defn}[पूर्ततायोग्यता]
एखाद्या संरचना~$\Struct M$ साठी $\Sat{M}{\Gamma}$ असेल, तर
वाक्यांचा संच~$\Gamma$ \emph{पूर्ततायोग्य} असतो. $\Gamma$
पूर्ततायोग्य नसेल, तर त्याला \emph{अपूर्ततायोग्य} म्हणतात.
\end{defn}

\begin{prop}
वाक्य~$\varphi$ वैध असते तेव्हा आणि केवळ तेव्हाच, जेव्हा
वाक्यांच्या प्रत्येक संच~$\Gamma$ साठी $\Gamma \Entails \varphi$.
\end{prop}

\begin{proof}
अग्र दिशेसाठी, $\varphi$ वैध आणि $\Gamma$ हा वाक्यांचा संच असू द्या.
$\Sat{M}{\Gamma}$ अशी संरचना~$\Struct M$ असू द्या. $\varphi$ वैध
असल्यामुळे $\Sat{M}{\varphi}$; म्हणून $\Gamma \Entails \varphi$.

उलट दिशेच्या व्यत्यासासाठी, $\varphi$ अवैध असू द्या; म्हणजे
$\Sat/{M}{\varphi}$ अशी संरचना~$\Struct M$ आहे. $\Gamma =
\{ \ltrue \}$ असताना $\ltrue$ वैध असल्यामुळे $\Sat{M}{\Gamma}$.
त्यामुळे अशी संरचना~$\Struct M$ आहे की $\Sat{M}{\Gamma}$ पण
$\Sat/{M}{\varphi}$; म्हणून $\Gamma$ मधून~$\varphi$ तार्किकरीत्या निष्पन्न होत
नाही.
\end{proof}

\begin{prop}
\label{fol:syn:sem:prop:entails-unsat}
$\Gamma \Entails \varphi$ तेव्हा आणि केवळ तेव्हाच, जेव्हा
$\Gamma \cup \{\lnot \varphi\}$ अपूर्ततायोग्य असतो.
\end{prop}

\begin{proof}
अग्र दिशेसाठी, $\Gamma \Entails \varphi$ असे समजू आणि, याउलट, अशी
संरचना~$\Struct M$ आहे की $\Sat{M}{\Gamma
  \cup \{ \lnot \varphi \}}$, असे समजू. $\Sat{M}{\Gamma}$ आणि
$\Gamma \Entails \varphi$ असल्यामुळे $\Sat{M}{\varphi}$. तसेच
$\Sat{M}{\Gamma\cup \{ \lnot \varphi \}}$ असल्यामुळे
$\Sat{M}{\lnot \varphi}$; म्हणून $\Sat{M}{\varphi}$ आणि $\Sat/{M}{\varphi}$ दोन्ही
मिळतात, हा विरोधाभास आहे. त्यामुळे अशी कोणतीही
संरचना~$\Struct M$ असू शकत नाही; म्हणून
$\Gamma \cup \{ \lnot \varphi \}$ अपूर्ततायोग्य आहे.

उलट दिशेसाठी, $\Gamma \cup \{ \lnot \varphi \}$ अपूर्ततायोग्य आहे असे
समजू. म्हणून प्रत्येक संरचना~$\Struct M$ साठी एकतर
$\Sat/{M}{\Gamma}$ किंवा $\Sat{M}{\varphi}$. त्यामुळे
$\Sat{M}{\Gamma}$ असलेल्या प्रत्येक संरचना~$\Struct M$ साठी
$\Sat{M}{\varphi}$; म्हणून $\Gamma \Entails \varphi$.
\end{proof}

\begin{prob}
\begin{enumerate}
\item $\Gamma \Entails \bot$ तेव्हा आणि केवळ तेव्हाच, जेव्हा $\Gamma$
  अपूर्ततायोग्य आहे, हे दाखवा.
\item $\Gamma \cup \{\varphi\} \Entails \bot$ तेव्हा आणि केवळ तेव्हाच,
  जेव्हा $\Gamma \Entails \lnot \varphi$, हे दाखवा.
\item $\varphi$ किंवा $\Gamma$ मध्ये~$c$ येत नाही असे समजा.
  $\Gamma \Entails \lforall[x][\varphi]$ तेव्हा आणि केवळ तेव्हाच, जेव्हा
  $\Gamma \Entails \Subst{\varphi}{c}{x}$, हे दाखवा.
\end{enumerate}
\end{prob}

\begin{prop}
$\Gamma \subseteq \Gamma'$ आणि $\Gamma \Entails \varphi$ असेल, तर
$\Gamma' \Entails \varphi$.
\end{prop}

\begin{proof}
$\Gamma \subseteq \Gamma'$ आणि $\Gamma \Entails \varphi$ असे समजा.
$\Sat{M}{\Gamma'}$ अशी रचना~$\Struct M$ असू द्या; मग
$\Sat{M}{\Gamma}$ आणि $\Gamma \Entails \varphi$ असल्यामुळे
$\Sat{M}{\varphi}$. म्हणून जेव्हा जेव्हा $\Sat{M}{\Gamma'}$, तेव्हा
$\Sat{M}{\varphi}$; त्यामुळे $\Gamma' \Entails \varphi$.
\end{proof}


\begin{thm}[चिन्हार्थविषयक निगमन प्रमेय]
\label{fol:syn:sem:thm:sem-deduction}
$\Gamma \cup \{\varphi\} \Entails \psi$ तेव्हा आणि केवळ तेव्हाच, जेव्हा
$\Gamma \Entails \varphi \lif \psi$.
\end{thm}

\begin{proof}
अग्र दिशेसाठी, $\Gamma \cup \{ \varphi \} \Entails \psi$ असू द्या आणि
$\Sat{M}{\Gamma}$ अशी संरचना~$\Struct M$ असू द्या.
$\Sat{M}{\varphi}$ असेल, तर $\Sat{M}{\Gamma \cup \{ \varphi \} }$; आणि
$\Gamma \cup \{ \varphi \}$ मधून~$\psi$ तार्किकरीत्या निष्पन्न होत असल्यामुळे
$\Sat{M}{\psi}$. म्हणून $\Sat{M}{\varphi \lif \psi}$; त्यामुळे
$\Gamma \Entails \varphi \lif \psi$.

उलट दिशेसाठी, $\Gamma \Entails \varphi \lif \psi$ असू द्या आणि
$\Sat{M}{\Gamma \cup \{ \varphi \}}$ अशी संरचना~$\Struct M$ असू
द्या. मग $\Sat{M}{\Gamma}$; म्हणून $\Sat{M}{\varphi \lif \psi}$ आणि
$\Sat{M}{\varphi}$ असल्यामुळे $\Sat{M}{\psi}$. त्यामुळे जेव्हा जेव्हा
$\Sat{M}{\Gamma \cup \{ \varphi \} }$, तेव्हा $\Sat{M}{\psi}$; म्हणून
$\Gamma \cup \{ \varphi \} \Entails \psi$.
\end{proof}

\begin{prop}\label{fol:syn:sem:prop:quant-terms}
  $\Struct{M}$ ही संरचना, $\varphi(x)$ हे एक मुक्त चर~$x$ असलेले
  सूत्र आणि $t$ हे बंद पद असू द्या. मग:
  \begin{enumerate}
    \item $\varphi(t) \Entails \lexists[x][\varphi(x)]$
    \item $\lforall[x][\varphi(x)] \Entails \varphi(t)$
  \end{enumerate}
\end{prop}

\begin{proof}
  \begin{enumerate}
    \item
      $\Sat{M}{\varphi(t)}$ असे समजा. $s(x) =
        \Value{t}{M}$ असे चर-मूल्यांकन~$s$ असू द्या. $\varphi(t)$ हे
        वाक्य असल्यामुळे $\Sat{M}{\varphi(t)}[s]$.
        \ref{fol:syn:ext:prop:ext-formulas} नुसार
        $\Sat{M}{\varphi(x)}[s]$. \ref{fol:syn:ass:prop:sat-quant} नुसार
        $\Sat{M}{\lexists[x][\varphi(x)]}$.
    \item
      $\Sat{M}{\lforall[x][\varphi(x)]}$ असे समजा.
        $s(x) = \Value{t}{M}$ असे चर-मूल्यांकन~$s$ असू द्या.
        \ref{fol:syn:ass:prop:sat-quant} नुसार $\Sat{M}{\varphi(x)}[s]$.
        \ref{fol:syn:ext:prop:ext-formulas} नुसार
        $\Sat{M}{\varphi(t)}[s]$. $\varphi(t)$ हे वाक्य असल्यामुळे
        \ref{fol:syn:ass:prop:sentence-sat-true} नुसार
        $\Sat{M}{\varphi(t)}$.
  \end{enumerate}
\end{proof}





\chapter{उपपत्ती आणि त्यांची प्रतिमाने}\label{fol:mat::chap}








\section{परिचय}\label{fol:mat:int:sec}

\begin{explain}
स्वयंसिद्धकीय पद्धतीचा विकास हा विज्ञानाच्या इतिहासातील एक महत्त्वपूर्ण
पराक्रम आहे आणि गणिताच्या इतिहासात त्याला विशेष महत्त्व आहे. एखाद्या
क्षेत्राचा स्वयंसिद्धकीय विकास करताना अनेक प्रश्न स्पष्ट करावे लागतात: हे
क्षेत्र कशाविषयी आहे? सर्वांत मूलभूत संकल्पना कोणत्या? त्यांचे परस्पर संबंध
कसे आहेत? क्षेत्रातील सर्व संकल्पना या मूलभूत संकल्पनांच्या साहाय्याने
परिभाषित करता येतात का? या संकल्पना कोणते नियम पाळतात आणि त्यांनी कोणते
नियम पाळलेच पाहिजेत?

स्वयंसिद्धकीय पद्धत आणि तर्कशास्त्र जणू एकमेकांसाठीच निर्माण झाले आहेत.
आकारिक तर्कशास्त्र स्वयंसिद्धकीय उपपत्ती मांडण्यासाठी, उपपत्तीच्या
स्वयंसिद्धकांपासून प्रमेये नेमक्या निर्दिष्ट रीतीने सिद्ध करण्यासाठी आणि
स्वयंसिद्धकांची पूर्ति करणाऱ्या सर्व प्रणालींच्या गुणधर्मांचा सुव्यवस्थित
अभ्यास करण्यासाठी साधने पुरवते.
\end{explain}

\begin{defn}
वाक्यांचा संच~$\Gamma$ \emph{बंद} असतो तेव्हा आणि केवळ तेव्हाच,
जेव्हा $\Gamma \Entails \varphi$ असल्यास $\varphi \in \Gamma$ असते.
वाक्यांच्या संच~$\Gamma$ चे \emph{संवरण} म्हणजे
$\Setabs{\varphi}{\Gamma \Entails \varphi}$ होय.

जर~$\Gamma$ हे~$\Delta$ चे संवरण असेल, तर~$\Gamma$ हा वाक्यांच्या
संच~$\Delta$ ने \emph{स्वयंसिद्धकांद्वारे निरूपित} केला आहे असे म्हणतो.
\end{defn}

\begin{explain}
एखादी स्वयंसिद्धकीय उपपत्ती म्हणजे तिच्या स्वयंसिद्धकांच्या संच~$\Delta$
ने निरूपित केलेल्या वाक्यांचा संच, असे आपण मानू शकतो. दुसऱ्या शब्दांत,
आपल्याला अभ्यासायच्या स्वयंसिद्धकीय रीतीने विकसित केलेल्या शास्त्रातील
मूलभूत संकल्पना दर्शवणारी अतार्किक चिन्हे असलेली प्रथम-क्रम भाषा आणि त्या
शास्त्राचे मूलभूत नियम व्यक्त करणारा वाक्यांचा संच दिलेला असेल,
तर स्वयंसिद्धकांपासून तार्किकरीत्या निष्पन्न होणारी या भाषेतील सर्व
वाक्ये उपपत्तीचे प्रतिनिधित्व करतात असे आपण मानू शकतो. याची व्याप्ती,
एकच आदिम संकल्पना आणि सोपी स्वयंसिद्धके असलेल्या अंशतः क्रमांच्या
उपपत्तीसारख्या साध्या उदाहरणांपासून न्यूटनीय यामिकीसारख्या गुंतागुंतीच्या
उपपत्तींपर्यंत आहे.

स्वयंसिद्धकीय पद्धतीकडे पाहण्याचा हा आकारिक दृष्टिकोन इतका महत्त्वाचा
ठरण्यामागील मुख्य तार्किक तथ्ये पुढीलप्रमाणे आहेत. $\Gamma$ ही एखाद्या
उपपत्तीची स्वयंसिद्धक-प्रणाली, म्हणजे वाक्यांचा संच, आहे असे समजा.
\begin{enumerate}
\item एखादी स्वयंसिद्धक-प्रणाली संरचनेचा अभिप्रेत वर्ग नेमका कधी
  पकडते, हे आपण काटेकोरपणे सांगू शकतो. म्हणजे, आपल्याला संरचनेच्या
  एखाद्या विशिष्ट वर्गात रस असेल, तर स्वयंसिद्धक-प्रणाली~$\Gamma$ तो वर्ग
  तेव्हा आणि केवळ तेव्हाच यशस्वीपणे पकडते, जेव्हा त्या संरचना म्हणजे
  नेमकी ती~$\Struct M$ असतात ज्यांच्यासाठी $\Sat{M}{\Gamma}$.
\item या बाबतीत आपल्याला अपयश येऊ शकते, कारण काही~$\Struct M$ साठी
  $\Sat{M}{\Gamma}$ असते, पण~$\Struct M$ ही आपण अभिप्रेत धरलेल्या
  संरचना पैकी नसते. त्यामुळे~$\Struct M$ मध्ये सत्य नसलेली
  स्वयंसिद्धके आपल्याला जोडावी लागू शकतात.
\item निदान~$\Gamma$ ही सर्व अभिप्रेत संरचनांमध्ये सत्य आहे एवढ्या
  बाबतीत आपण यशस्वी झालो, तर $\Gamma \Entails \varphi$ असताना वाक्य~$\varphi$ हे
  सर्व अभिप्रेत संरचनांमध्ये सत्य असते. म्हणून, वाक्ये सर्व
  अभिप्रेत संरचनांमध्ये सत्य आहेत हे दाखवण्यासाठी, ती
  स्वयंसिद्धकांपासून तार्किकरीत्या निष्पन्न होतात एवढे दाखवून आपण तार्किक
  साधने (उदाहरणार्थ, निष्पत्ती-पद्धती) वापरू शकतो.
\item कधीकधी आपल्या मनात अभिप्रेत संरचना नसतात; त्याऐवजी आपण
  स्वयंसिद्धकांपासूनच सुरुवात करतो: काही विशिष्ट नियमांची पूर्ति करावी अशा
  काही आदिम संकल्पना आपण घेतो आणि ते नियम स्वयंसिद्धक-प्रणालीत संकेतबद्ध
  करतो. स्वयंसिद्धके परस्परविरोधी नाहीत, ही गोष्ट आपल्याला लगेच पडताळून
  पाहायची असते. ती परस्परविरोधी असतील, तर त्या नियमांचे पालन करणाऱ्या
  संकल्पना असूच शकत नाहीत आणि आपण विसंगत उपपत्ती मांडण्याचा प्रयत्न केलेला
  असतो. $\Gamma$ चे प्रतिमान शोधून असे घडत नाही हे आपण पडताळू शकतो. आणि
  आपल्या उपपत्तीची प्रतिमाने असतील, तर त्यांचा तपास करण्यासाठी तसेच प्रतिमाने
  रचण्यासाठी आपण तार्किक पद्धती वापरू शकतो.
\item स्वयंसिद्धकांचे स्वातंत्र्य हाही एक महत्त्वाचा प्रश्न आहे. एखादे
  स्वयंसिद्धक प्रत्यक्षात इतरांचा परिणाम असल्यामुळे अनावश्यक असू शकते.
  $\Gamma$ मधील स्वयंसिद्धक~$\varphi$ अनावश्यक आहे हे आपण
  $\Gamma \setminus \{\varphi\} \Entails \varphi$ सिद्ध करून दाखवू शकतो. तसेच,
  $(\Gamma \setminus \{\varphi\}) \cup \{\lnot \varphi\}$ पूर्ततायोग्य आहे हे
  दाखवून ते स्वयंसिद्धक अनावश्यक नाही हेही सिद्ध करू शकतो. उदाहरणार्थ,
  भूमितीचे समांतर गृहीतक इतर स्वयंसिद्धकांपासून स्वतंत्र आहे हे अशाच प्रकारे
  दाखवले गेले.
\item उपपत्तीतील संकल्पनांची परिभाष्यता हाही आणखी एक महत्त्वाचा प्रश्न आहे:
  भाषेच्या निवडीवर उपपत्तीची प्रतिमाने कशाची बनलेली असतात हे ठरते. पण
  उपपत्तीचा प्रत्येक पैलू तिच्या प्रतिमानांत स्वतंत्रपणे दर्शवावाच लागतो असे
  नाही. उदाहरणार्थ, प्रत्येक क्रमसंबंध~$\le$ त्याच्याशी संबंधित काटेकोर
  क्रमसंबंध~$<$ ठरवतो---एक दिला असेल, तर दुसरा परिभाषित करता येतो. म्हणून
  असा क्रम असलेल्या उपपत्तीच्या प्रतिमानात त्याच्याशी संबंधित काटेकोर
  क्रमसंबंध \emph{वेगळ्याने} असणे आवश्यक नाही. सर्वसाधारणपणे, एक संबंध
  इतरांच्या आधारे नेमका कधी परिभाषित करता येतो? एखादा संबंध इतरांच्या
  आधारे परिभाषित करणे कधी अशक्य असते (आणि म्हणून तो भाषेच्या आदिम
  संकल्पनांत जोडावा लागतो)?
\end{enumerate}
\end{explain}












\section{संरचनेचे गुणधर्म व्यक्त करणे}\label{fol:mat:exs:sec}

\begin{explain}
फलने आणि संबंध यांच्यावरील अटी व्यक्त करणे, किंवा अधिक सर्वसाधारणपणे,
रचनेतील फलने आणि संबंध या अटींची पूर्ति करतात हे व्यक्त करणे अनेकदा
उपयुक्त आणि महत्त्वाचे असते. उदाहरणार्थ, एखाद्या भाषेसाठीच्या ज्या
संरचना आपल्याला विधेयांचा जो अर्थ ``पकडायचा'' आहे तो
पकडतात, त्या तसे न करणाऱ्या रचनांपासून वेगळ्या ओळखण्याचे मार्ग आपल्याला
हवे असतात. अर्थात, कोणती संरचना आपल्याला ``अभिप्रेत'' आहेत हे
ठरवण्याचे पूर्ण स्वातंत्र्य आपल्याला आहे. उदाहरणार्थ, आपण असे ठरवू शकतो
की विधेय~$\le$ चे अर्थनिर्धारण हा क्रमसंबंध असलाच पाहिजे; किंवा
आपल्याला फक्त~$\Lang L$ च्या अशा अर्थनिर्धारणांत रस आहे की ज्यांचा प्रांत
संचांचा बनलेला आहे आणि ज्यांत~$\Obj \in$ चे अर्थनिर्धारण ``चा
घटक आहे'' या संबंधाने केलेले आहे. पण भाषेतील वाक्ये वापरून
आपण हे करू शकतो का? दुसऱ्या शब्दांत, संरचना~$\Struct M$ वरील
कोणत्या अटी आपण~$\Struct M$ च्या भाषेतील वाक्य (किंवा कदाचित
वाक्यांचा संच) वापरून व्यक्त करू शकतो? काही अटी आपल्याला व्यक्त करता
येणार नाहीत. उदाहरणार्थ, फक्त $\Domain M = \Nat$ असलेल्या
संरचना~$\Struct M$ मध्येच सत्य असे~$\Lang L_A$ चे कोणतेही वाक्य
नाही. ``प्रांतात फक्त नैसर्गिक संख्या आहेत'' हे आपण व्यक्त करू शकत नाही.
पण संरचनेचे काही ``रचनात्मक गुणधर्म'' आपण कदाचित व्यक्त करू शकतो.
संरचनेचे कोणते गुणधर्म आपण वाक्ये वापरून व्यक्त करू शकतो?
किंवा, दुसऱ्या प्रकारे विचारल्यास, एखादे वाक्य (किंवा
वाक्यांचा संच) सत्य करणाऱ्या रचना म्हणून आपण
संरचनेच्या कोणत्या संग्रहांचे वर्णन करू शकतो?
\end{explain}

\begin{defn}[संचाचे प्रतिमान]
भाषा~$\Lang L$ मधील वाक्यांचा संच~$\Gamma$ असू द्या. प्रत्येक
$\varphi \in \Gamma$ साठी $\Sat{M}{\varphi}$ असेल, तर संरचना~$\Struct M$
\emph{हे~$\Gamma$ चे प्रतिमान आहे} असे म्हणतो.
\end{defn}

\begin{ex}
वाक्य $\lforall[x][x \le x]$ हे~$\Struct M$ मध्ये तेव्हा आणि केवळ तेव्हाच
सत्य असते, जेव्हा $\Assign{\le}{M}$ हा परावर्ती संबंध असतो. वाक्य
$\lforall[x][\lforall[y][((x \le y \land y \le x) \lif x = y)]]$ हे
~$\Struct M$ मध्ये तेव्हा आणि केवळ तेव्हाच सत्य असते, जेव्हा
$\Assign{\le}{M}$ प्रतिसममित असतो. वाक्य
$\lforall[x][\lforall[y][\lforall[z][((x \le y \land y \le z)
      \lif x \le z)]]]$ हे~$\Struct M$ मध्ये तेव्हा आणि केवळ तेव्हाच सत्य
असते, जेव्हा $\Assign{\le}{M}$ संक्रमक असतो. म्हणून पुढील संचाची प्रतिमाने
\begin{align*}
\{\quad &\lforall[x][x \le x], \\
   & \lforall[x][\lforall[y][((x \le y \land y \le
    x) \lif x = y)]], \\
   &\lforall[x][\lforall[y][\lforall[z][((x \le y
      \land y \le z) \lif x \le z)]]] \quad \}
\end{align*}
म्हणजे नेमकी ती रचना होत ज्यांत~$\Assign{\le}{M}$ हा परावर्ती,
प्रतिसममित आणि संक्रमक, म्हणजेच अंशतः क्रम, असतो. त्यामुळे ही वाक्ये
आपण \emph{अंशतः क्रमांच्या प्रथम-क्रम उपपत्तीची} स्वयंसिद्धके म्हणून घेऊ
शकतो.
\end{ex}












\section{प्रथम-क्रम उपपत्तींची उदाहरणे}\label{fol:mat:the:sec}

\begin{ex}
भाषा~$\Lang L_<$ मधील काटेकोर रेषीय क्रमांची उपपत्ती पुढील संचाने
स्वयंसिद्धकांद्वारे निरूपित होते:
\begin{align*}
\{\quad & \lforall[x][\lnot x < x], \\
& \lforall[x][\lforall[y][((x < y \lor y <
    x) \lor x = y)]], \\
& \lforall[x][\lforall[y][\lforall[z][((x < y
      \land y < z) \lif x < z)]]] \quad \}
\end{align*}
ही उपपत्ती अभिप्रेत संरचना पूर्णपणे पकडते: प्रत्येक काटेकोर
रेषीय क्रम या स्वयंसिद्धक-प्रणालीचे प्रतिमान असतो; आणि उलट, $R$ हा
संच~$X$ वरील रेषीय क्रम असेल, तर $\Domain M = X$ आणि
$\Assign{<}{M} = R$ असलेली रचना~$\Struct M$ या उपपत्तीचे प्रतिमान असते.
\end{ex}

\begin{ex}
$\Obj 1$ (स्थिरांक) आणि $\cdot$ (द्विस्थानी फलन) असलेल्या
भाषेतील गटांची उपपत्ती पुढील वाक्यांनी स्वयंसिद्धकांद्वारे निरूपित होते:
\begin{align*}
& \lforall[x][\eq[(x \cdot \Obj 1)][x]]\\
& \lforall[x][\lforall[y][\lforall[z][\eq[(x \cdot (y \cdot z))][((x
          \cdot y) \cdot z)]]]]\\
& \lforall[x][\lexists[y][\eq[(x \cdot y)][\Obj 1]]]
\end{align*}
\end{ex}

\begin{ex}
पेआनो अंकगणिताची उपपत्ती अंकगणिताची भाषा~$\Lang L_A$ मधील पुढील
वाक्यांनी स्वयंसिद्धकांद्वारे निरूपित होते.
\begin{align*}
& \lforall[x][\lforall[y][(\eq[x'][y'] \lif \eq[x][y])]]\\
& \lforall[x][\eq/[\Obj 0][x']]\\
& \lforall[x][\eq[(x + \Obj 0)][x]]\\
& \lforall[x][\lforall[y][\eq[(x + y')][(x + y)']]]\\
& \lforall[x][\eq[(x \times \Obj 0)][\Obj 0]]\\
& \lforall[x][\lforall[y][\eq[(x \times y')][((x \times y) + x)]]]\\
& \lforall[x][\lforall[y][(x < y \liff \lexists[z][\eq[(z' + x)][y])]]]\\
\intertext{आणि पुढील रूपाची सर्व वाक्ये}
& (\varphi(\Obj 0) \land \lforall[x][(\varphi(x) \lif \varphi(x'))]) \lif \lforall[x][\varphi(x)]
\end{align*}
शेवटच्या रूपाची अनंतपणे अनेक वाक्ये असल्यामुळे ही स्वयंसिद्धक-प्रणाली
अनंत आहे. शेवटच्या रूपाला \emph{विगमन-रूपबंध} म्हणतात. (प्रत्यक्षात,
विगमन-रूपबंधाचे स्वरूप आपण येथे दाखवले आहे त्यापेक्षा थोडे अधिक गुंतागुंतीचे
आहे.)

शेवटचे स्वयंसिद्धक हे~$<$ ची \emph{स्पष्ट व्याख्या} आहे.
\end{ex}

\begin{ex}
शुद्ध संचांची उपपत्ती गणिताच्या पायामध्ये (आणि गणिताच्या तत्त्वज्ञानातही)
महत्त्वाची भूमिका बजावते. एखाद्या संचाचे सर्व घटक हेही शुद्ध संच
असतील, तर तो संच शुद्ध असतो. म्हणून रिक्त संच शुद्ध मानला जातो; पण एखाद्या
संचात संच नसलेली कोणतीही वस्तू घटक म्हणून असेल, तर तो संच शुद्ध
नसतो. म्हणजे शुद्ध संच हे फक्त रिक्त संचापासून, कोणतेही ``मूलघटक'' न
वापरता बनवलेले संच होत; मूलघटक म्हणजे स्वतः संच नसलेल्या वस्तू.

पुढील वाक्ये शुद्ध संचांच्या उपपत्तीसाठी स्वयंसिद्धक-प्रणाली मानता येतील:
\begin{align*}
& \lexists[x][\lnot \lexists[y][y \in x]]\\
& \lforall[x][\lforall[y][(\lforall[z](z \in x \liff z \in y) \lif
\eq[x][y])]]\\
& \lforall[x][\lforall[y][\lexists[z][\lforall[u][(u \in z \liff
(\eq[u][x] \lor \eq[u][y]))]]]]\\
& \lforall[x][\lexists[y][\lforall[z][(z \in y \liff \lexists[u][(z \in
        u \land u \in x)])]]]\\
\intertext{आणि पुढील रूपाची सर्व वाक्ये} &
\lexists[x][\lforall[y][(y \in x \liff \varphi(y))]]
\end{align*}
पहिले स्वयंसिद्धक सांगते की कोणतेही घटक नसलेला एक संच आहे (म्हणजे
$\emptyset$ अस्तित्वात आहे); दुसरे सांगते की संच विस्तारात्मक आहेत; तिसरे
सांगते की कोणत्याही संच~$X$ आणि~$Y$ साठी संच~$\{X, Y\}$ अस्तित्वात आहे;
आणि चौथे सांगते की कोणत्याही संच~$X$ साठी संच~$\cup X$ अस्तित्वात आहे,
जिथे~$\cup X$ हा~$X$ च्या सर्व घटकांचा संयोग आहे.

सर्वांत शेवटी उल्लेखलेल्या वाक्यांना एकत्रितपणे
\emph{अनौपचारिक गुणधर्माधारित संचरचना-रूपबंध} म्हणतात. तो मूलतः असे सांगतो
की प्रत्येक~$\varphi(x)$ साठी संच~$\Setabs{x}{\varphi(x)}$ अस्तित्वात आहे---म्हणून
पहिल्या नजरेत ते सत्य, उपयुक्त आणि कदाचित आवश्यकही वाटणारे स्वयंसिद्धक आहे.
त्याला ``अनौपचारिक'' म्हणतात कारण, प्रत्यक्षात, त्यामुळे ही उपपत्ती
अपूर्ततायोग्य होते: $\varphi(y)$ म्हणून $\lnot y \in y$ घेतल्यास पुढील
वाक्य मिळते:
\[
\lexists[x][\lforall[y][(y \in x \liff \lnot y \in y)]]
\]
आणि या वाक्याची कोणत्याही संरचनांमध्ये पूर्ति होत नाही.
\end{ex}

\begin{ex}
\emph{अवयवमीमांसेमध्ये} \emph{अवयवत्वाचा} संबंध हा एक मूलभूत संबंध आहे.
संचांच्या उपपत्तींप्रमाणेच अवयवत्वाच्याही अशा उपपत्ती आहेत, ज्या या
संबंधाविषयीच्या विविध (आणि कधीकधी परस्परविरोधी) संकल्पना
स्वयंसिद्धकांद्वारे निरूपित करतात.

अवयवमीमांसेच्या भाषेत एकच द्विस्थानी विधेयचिन्ह~$\Obj P$ असते आणि
$\Atom{\Obj P}{x, y}$ चा ``अर्थ'' $x$ हा~$y$ चा अवयव आहे असा असतो. हे
अर्थनिर्धारण मनात ठेवले असता, या भाषेतील संरचना ला
\emph{अवयवत्वाची रचना} म्हणतात. अर्थात, एकच द्विस्थानी विधेय असलेली प्रत्येक
रचना या नावास खरोखर पात्र असेलच असे नाही. ``अवयवत्व'' पकडण्याची शक्यता
असण्यासाठी $\Assign{\Obj P}{M}$ ने काही अटींची पूर्ति केली पाहिजे; त्या अटी
आपण अवयवत्वाच्या उपपत्तीची स्वयंसिद्धके म्हणून मांडू शकतो. उदाहरणार्थ,
वस्तूंवरील अवयवत्व हा अंशतः क्रम आहे: प्रत्येक वस्तू स्वतःचाच अवयव असते
(जरी तो \emph{अनुचित} अवयव असला तरी); कोणत्याही दोन भिन्न वस्तू एकमेकींचे
अवयव असू शकत नाहीत; आणि एखाद्या वस्तूच्या अवयवाचा अवयव हा स्वतः त्या
वस्तूचा अवयव असतो. लक्षात घ्या की या अर्थाने ``चा अवयव आहे'' हे ``चा
उपसंच आहे'' यासारखे आहे; पण ते ``चा घटक आहे'' यासारखे नाही, कारण घटक-संबंध
परावर्तीही नाही आणि संक्रमकही नाही.
\begin{align*}
& \lforall[x][\Atom{\Obj P}{x,x}] \\
& \lforall[x][\lforall[y][((\Part{x}{y} \land \Part{y}{x})
      \lif \eq[x][y])]] \\
& \lforall[x][\lforall[y][\lforall[z][((\Part{x}{y} \land
        \Part{y}{z}) \lif \Part{x}{z})]]]\\
\intertext{याखेरीज, कोणत्याही दोन वस्तूंचा अवयवी संयोग असतो (ज्याचे त्या दोन
  वस्तू अवयव असतात आणि जो या बाबतीत न्यूनतम असतो).}  &
\lforall[x][\lforall[y][\lexists[z][\lforall[u][(\Part{z}{u} \liff
        (\Part{x}{u} \land \Part{y}{u}))]]]]
\end{align*}
तत्त्वमीमांसकांनी विचारात घेतलेल्या अवयवत्वाच्या मूलभूत तत्त्वांपैकी ही
फक्त काही तत्त्वे आहेत. मात्र, आधी काही परिभाषित संबंध न मांडता पुढील
तत्त्वांची सूत्रे बनवणे किंवा ती लिहून काढणे लवकरच कठीण होते. उदाहरणार्थ,
अवयवमीमांसेत रस असलेले बहुतेक तत्त्वमीमांसक पुढील तत्त्वही वैध मानतात:
जेव्हा वस्तू~$x$ ला उचित अवयव~$y$ असतो, तेव्हा तिला असा अवयव~$z$ देखील
असतो की त्या अवयवाचा~$y$ शी कोणताही सामाईक अवयव नसतो आणि त्यामुळे~$y$ व
$z$ यांचा अवयवी संयोग~$x$ असतो.
\end{ex}












\section{संरचना मधील संबंध व्यक्त करणे}\label{fol:mat:exr:sec}

\begin{explain}
सूत्रांचा एक प्रमुख उपयोग म्हणजे संरचना~$\Struct M$ मधील
गुणधर्म आणि संबंध हे~$\Struct M$ ची भाषा~$\Lang L$ हिच्या आदिम संकल्पनांच्या
साहाय्याने व्यक्त करणे. याचा अर्थ पुढीलप्रमाणे आहे: $\Struct M$ चा
प्रांत हा वस्तूंचा संच असतो. स्थिरांक, फलने आणि
विधेये यांचे~$\Struct M$ मध्ये अनुक्रमे~$\Domain M$ मधील काही वस्तू,
$\Domain M$ वरील फलने आणि~$\Domain M$ वरील संबंध यांनी अर्थनिर्धारण केलेले
असते. उदाहरणार्थ, $\Obj A^2_0$ हे~$\Lang L$ मध्ये असेल, तर~$\Struct M$
त्याला संबंध~$R = \Assign{{\Obj A^2_0}}{M}$ नेमते. मग सूत्र
$\Atom{\Obj A^2_0}{\Obj v_1, \Obj v_2}$ पुढील अर्थाने नेमका तो संबंध
\emph{व्यक्त करते}: चर-मूल्यांकन~$s$ हे~$\Obj v_1$ ला
$a \in \Domain{M}$ आणि~$\Obj v_2$ ला $b \in \Domain M$ शी संगत करत असेल,
तर
\[
Rab \text{\quad iff\quad} \Sat{M}{\Atom{\Obj A^2_0}{\Obj v_1, \Obj v_2}}[s].
\]
येथे चर-मूल्यांकनांचा समावेश करणे आवश्यक आहे हे लक्षात घ्या: आपण फक्त
``$Rab$ iff $\Sat{M}{\Atom{\Obj A^2_0}{a, b}}$'' असे म्हणू शकत नाही, कारण
$a$ आणि~$b$ ही आपल्या भाषेतील चिन्हे नाहीत; ती~$\Domain{M}$ ची
घटक आहेत.

आपल्याकडे केवळ अणुसूत्र सूत्रे नसून, तार्किक संयोजके आणि संख्यापक
वापरून ती एकत्र करता येतात; त्यामुळे अधिक गुंतागुंतीची सूत्रे अशी इतर
संबंध परिभाषित करू शकतात जी~$\Struct M$ मध्ये थेट अंतर्भूत नसतात. ते कसे
करायचे आणि विशेषतः संरचनांमध्ये कोणते संबंध परिभाषित करता येतात,
यात आपल्याला रस आहे.
\end{explain}

\begin{defn}
$\varphi(\Obj v_1,\dots, \Obj v_n)$ हे~$\Lang L$ चे असे सूत्र असू द्या
की ज्यात फक्त~$\Obj v_1$,\dots, $\Obj v_n$ हीच चरे मुक्त येतात; आणि
$\Struct M$ ही~$\Lang L$ साठीची संरचना असू द्या. कोणत्याही अशा
चर-मूल्यांकन~$s$ साठी की $s(\Obj v_i) = a_i$ ($i = 1, \dots, n$),
\[
Ra_1\dots a_n \text{\quad iff\quad} \Sat{M}{\Atom{\varphi}{\Obj
    v_1,\dots, \Obj v_n}}[s]
\]
असे असेल, तर आणि केवळ तेव्हाच $\varphi(\Obj v_1,\dots, \Obj v_n)$ हे
$R \subseteq \Domain M^n$ हा \emph{संबंध व्यक्त करते}.
\end{defn}

\begin{ex}
अंकगणिताच्या मानक प्रतिमान~$\Struct N$ मध्ये सूत्र $\Obj
v_1 < \Obj v_2 \lor \eq[\Obj v_1][\Obj v_2]$ हे~$\Nat$ वरील~$\le$
संबंध व्यक्त करते. सूत्र $\eq[\Obj v_2][\Obj v_1']$ हे उत्तरवर्ती
संबंध व्यक्त करते; म्हणजे, $m$ हा~$n$ चा उत्तरवर्ती असेल तेव्हा
$Rnm$ असणारा संबंध~$R \subseteq \Nat^2$. सूत्र
$\eq[\Obj v_1][\Obj v_2']$ हे पूर्ववर्ती संबंध व्यक्त करते. सूत्रे
$\lexists[\Obj v_3][(\eq/[\Obj v_3][\Obj 0] \land
  \eq[\Obj v_2][(\Obj v_1 + \Obj v_3)])]$ आणि $\lexists[\Obj
  v_3][\eq[(\Obj v_1 + {\Obj v_3}')][\Obj v_2]]$ ही दोन्ही~$\Obj <$
संबंध व्यक्त करतात.

याचा अर्थ अंकगणिताच्या भाषेत विधेयचिन्ह~$<$ प्रत्यक्षात अनावश्यक आहे; ते
परिभाषित करता येते.
\end{ex}

\begin{explain}
ही कल्पना फक्त विशिष्ट संरचनांमध्येच रंजक नाही; एखादे अभिप्रेत
प्रतिमान किंवा प्रतिमाने वर्णन करण्यासाठी आपण भाषा वापरतो तेव्हा, म्हणजे
उपपत्तींचा विचार करताना, ती सर्वसाधारणपणेही रंजक ठरते. अशा उपपत्तींमध्ये
मूलभूत चिन्हे म्हणून अनेकदा फक्त काही विधेये असतात, पण ती ज्या
प्रांताचे वर्णन करतात त्यात इतर अनेक संबंध महत्त्वाची भूमिका बजावतात. या
इतर संबंधांना भाषेतील मूलभूत विधेयांचे अर्थनिर्धारण करणाऱ्या
संबंधांनी सुव्यवस्थित रीतीने व्यक्त करता येत असेल, तर आपण त्यांना त्या
भाषेत \emph{परिभाषित} करू शकतो असे म्हणतो.
\end{explain}

\begin{prob}
पुढील संबंध परिभाषित करणारी~$\Lang L_A$ मधील सूत्रे शोधा:
\begin{enumerate}
\item $n$ हे~$i$ आणि~$j$ यांच्या दरम्यान आहे;
\item $n$ ने~$m$ ला निःशेष भाग जातो (म्हणजे $m$ ही~$n$ ची पटीत संख्या आहे);
\item $n$ ही अविभाज्य संख्या आहे (म्हणजे $1$ आणि~$n$ यांखेरीज कोणत्याही
  संख्येने~$n$ ला निःशेष भाग जात नाही).
\end{enumerate}
\end{prob}

\begin{prob}
सूत्र $\varphi(\Obj v_1, \Obj v_2)$ हे संरचना~$\Struct M$ मध्ये संबंध
$R \subseteq \Domain M^2$ व्यक्त करते असे समजा. पुढील संबंध व्यक्त करणारी
सूत्रे शोधा:
\begin{enumerate}
\item $R$ चा व्युत्क्रम~$R^{-1}$;
\item सापेक्ष गुणाकार~$R \mid R$;
\end{enumerate}
$R$ चे संक्रमक संवरण~$R^+$ व्यक्त करण्याचा मार्ग तुम्हाला शोधता येतो का?
\end{prob}

\begin{prob}
फक्त एक द्विस्थानी विधेयचिन्ह~$<$ असलेली भाषा~$\Lang{L}$ असू द्या (इतर
कोणतीही स्थिरांक, फलने किंवा विधेये नाहीत---अर्थात
~$\eq$ सोडून). $\Domain{N} = \Nat$ आणि
$\Assign{<}{N} = \Setabs{\tuple{n,m}}{n < m}$ असणारी रचना~$\Struct{N}$
असू द्या. पुढील विधाने सिद्ध करा:
\begin{enumerate}
\item $\{ 0 \}$ हा~$\Struct{N}$ मध्ये परिभाष्य आहे;
\item $\{ 1 \}$ हा~$\Struct{N}$ मध्ये परिभाष्य आहे;
\item $\{ 2 \}$ हा~$\Struct{N}$ मध्ये परिभाष्य आहे;
\item प्रत्येक $n \in \Nat$ साठी संच~$\{ n \}$ हा~$\Struct{N}$ मध्ये
  परिभाष्य आहे;
\item $\Domain{N}$ चा प्रत्येक सांत उपसंच~$\Struct{N}$ मध्ये परिभाष्य आहे;
\item $\Domain{N}$ चा प्रत्येक सांत-पूरक उपसंच~$\Struct{N}$ मध्ये परिभाष्य
  आहे (जिथे $X \subseteq \Nat$ हा सांत-पूरक असतो तेव्हा आणि केवळ तेव्हाच,
  जेव्हा $\Nat \setminus X$ हा सांत असतो).
\end{enumerate}
\end{prob}












\section{संचांची उपपत्ती}\label{fol:mat:set:sec}

जवळजवळ संपूर्ण गणित संचांच्या उपपत्तीमध्ये विकसित करता येते. या
उपपत्तीमध्ये गणित विकसित करताना अनेक गोष्टी कराव्या लागतात. प्रथम,
संबंध~$\in$ साठी स्वयंसिद्धकांचा संच लागतो. वेगवेगळ्या अनेक
स्वयंसिद्धक-प्रणाली विकसित झाल्या असून, त्यांतील~$\in$ चे गुणधर्म कधीकधी
परस्परविरोधी असतात. निवडीच्या स्वयंसिद्धकासह झर्मेलो--फ्रेंकेल संच
उपपत्ती म्हणून ओळखली जाणारी स्वयंसिद्धक-प्रणाली~$\Log{ZFC}$ विशेष उठून
दिसते: ती आतापर्यंत सर्वाधिक वापरली आणि अभ्यासली गेलेली प्रणाली आहे, कारण
गणितज्ञांना ज्या जवळजवळ सर्व गोष्टी सिद्ध करता याव्यात अशी अपेक्षा असते,
त्या सिद्ध करण्यास तिची स्वयंसिद्धके पुरेशी ठरतात. पण हे स्थापित करण्यापूर्वी
गणितज्ञांना व्यक्त करायच्या सर्व गोष्टी आपण \emph{व्यक्त} तरी कशा करू शकतो,
हे प्रथम स्पष्ट करणे आवश्यक आहे. सुरुवातीलाच, भाषेत कोणतीही स्थिरांक
किंवा फलने नसतात; त्यामुळे विशिष्ट संचांविषयी (उदाहरणार्थ,
$\emptyset$ किंवा~$\Nat$), संचांवरील क्रियांविषयी (उदाहरणार्थ,
$X \cup Y$ आणि~$\Pow{X}$), इतकेच नव्हे तर संबंध आणि फलने यांसारख्या
संचांखेरीज इतर गोष्टींचा समावेश असलेल्या रचनांविषयी आपण बोलू शकतो हेही
पहिल्या नजरेत अस्पष्ट वाटते.

सुरुवातीला, ``चा घटक आहे'' हाच काही आपल्याला महत्त्वाचा असलेला एकमेव संबंध
नाही; ``चा उपसंच आहे'' हा जवळजवळ तितकाच महत्त्वाचा वाटतो. पण ``चा घटक
आहे'' याच्या आधारे ``चा उपसंच आहे'' हे आपण \emph{परिभाषित} करू शकतो. त्यासाठी
संच उपपत्तीच्या भाषेत असे सूत्र~$\varphi(x, y)$ शोधावे लागते की संचांची
जोडी~$\tuple{X, Y}$ त्याची पूर्ति तेव्हा आणि केवळ तेव्हाच करते, जेव्हा
$X \subseteq Y$. पण~$X$ हा~$Y$ चा उपसंच तेव्हा आणि केवळ तेव्हाच असतो,
जेव्हा~$X$ चे सर्व घटक हे~$Y$ चेही घटक असतात. म्हणून
पुढील सूत्राने~$\subseteq$ ची व्याख्या करता येते:
\[
\lforall[z][(z \in x \lif z \in y)]
\]
आता, एखाद्या सूत्रात~$\subseteq$ संबंध वापरायचा असेल तेव्हा त्याऐवजी हे
सूत्र वापरता येईल ($x$ आणि~$y$ यांच्या जागी योग्य पदे घालून आणि आवश्यक
असल्यास बद्ध चर~$z$ चे नाव बदलून). उदाहरणार्थ, संचांची विस्तारात्मकता
म्हणजे कोणतेही संच~$x$ आणि~$y$ हे परस्परांचे उपसंच असतील, तर~$x$ आणि~$y$
एकच संच असले पाहिजेत. हे
$\lforall[x][\lforall[y][((x \subseteq y \land y
    \subseteq x) \lif x = y)]]$ ने व्यक्त करता येते; किंवा वरच्या व्याख्येने
$\subseteq$ बदलल्यास पुढील सूत्राने:
\[
\lforall[x][\lforall[y][((\lforall[z][(z \in x \lif z \in y)] \land
    \lforall[z][(z \in y \lif z \in x)]) \lif x = y)]].
\]
हे खरे तर~$\Log{ZFC}$ च्या स्वयंसिद्धकांपैकी एक, ``विस्तारात्मकतेचे
स्वयंसिद्धक,'' आहे.

$\emptyset$ साठी कोणतेही स्थिरांक नाही, पण ``$x$ रिक्त आहे'' हे
$\lnot \lexists[y][y \in x]$ ने व्यक्त करता येते. मग ``$\emptyset$
अस्तित्वात आहे'' हे वाक्य~$\lexists[x][\lnot \lexists[y][y \in
    x]]$ बनते. हे~$\Log{ZFC}$ चे आणखी एक स्वयंसिद्धक आहे. (विस्तारात्मकतेच्या
स्वयंसिद्धकानुसार फक्त एकच रिक्त संच असतो, हे लक्षात घ्या.) संच उपपत्तीच्या
भाषेत~$\emptyset$ विषयी बोलायचे असेल तेव्हा आपण ते ``रिक्त असलेला एक संच
अस्तित्वात आहे आणि \dots'' असे लिहू. उदाहरणार्थ, $\emptyset$ हा प्रत्येक संचाचा
उपसंच आहे ही गोष्ट पुढीलप्रमाणे व्यक्त करता येईल:
\[
\lexists[x][(\lnot\lexists[y][y \in x] \land \lforall[z][x \subseteq
    z])]
\]
येथे अर्थातच $x \subseteq z$ हेसुद्धा त्याच्या व्याख्येने बदलावे लागेल.

$X \cup Y$ आणि~$\Pow{X}$ यांसारख्या संचांवरील क्रियांविषयी बोलण्यासाठी
अशीच युक्ती वापरावी लागते. संच उपपत्तीच्या भाषेत फलनचिन्हे नसतात, पण
$X \cup Y = Z$ आणि~$\Pow{X} = Y$ हे फलनीय संबंध पुढील सूत्रांनी व्यक्त
करता येतात:
\begin{align*}
& \lforall[u][((u \in x \lor u \in y) \liff u \in z)]\\
& \lforall[u][(u \subseteq x \liff u \in y )]
\end{align*}
कारण $X \cup Y$ चे घटक म्हणजे नेमके~$X$ चे घटक किंवा
~$Y$ चे घटक असलेले संच, आणि~$\Pow{X}$ चे घटक म्हणजे नेमके~$X$
चे उपसंच होत. पण त्यामुळे आपण $x \cup y$ किंवा~$\Pow{x}$ ही पदे असल्यासारखी
वापरू शकत नाही; $X \cup Y = Z$ आणि~$\Pow{X} = Y$ हे संबंध परिभाषित करणारी
संपूर्ण सूत्रे एवढीच वापरता येतात. प्रत्यक्षात, या संबंधांची कधी तरी
पूर्ति होते, म्हणजे संयोग आणि घातसंच नेहमी अस्तित्वात असतात, हे आपल्याला
अजून माहीत नाही. उदाहरणार्थ, वाक्य
$\lforall[x][\lexists[y][\Pow{x} = y]]$ हे~$\Log{ZFC}$ चे आणखी एक
स्वयंसिद्धक (घातसंचाचे स्वयंसिद्धक) आहे.

आता क्रमित जोड्या किंवा फलने यांविषयी कसे बोलायचे? येथे क्रमित जोड्या आणि
फलने हे विशिष्ट प्रकारचे संच म्हणून कसे मानता येतात, हे स्पष्ट करावे लागते.
क्रमित जोडी~$\tuple{x, y}$ ची व्याख्या संच~$\{\{x\}, \{x, y\}\}$ म्हणून
करणे हा एक मार्ग आहे. पण आधीप्रमाणेच या संचाला नाव देणारे फलन
आपण समाविष्ट करू शकत नाही; फक्त $\tuple{x, y} = z$, म्हणजे
$\{\{x\}, \{x, y\}\} = z$, हा संबंध परिभाषित करता येतो:
\[
\lforall[u][(u \in z \liff (\lforall[v][(v \in u \liff v = x)] \lor
  \lforall[v][(v \in u \liff (v = x \lor v = y))]))]
\]
याचा अर्थ~$z$ चे घटक~$u$ म्हणजे नेमके असे संच की ज्यांचा एकमेव
घटक हा~$x$ असतो किंवा ज्यांचे एकमेव घटक हे~$x$ आणि~$y$
असतात (दुसऱ्या शब्दांत, जे संच~$\{x\}$ शी किंवा~$\{x, y\}$ शी एकरूप
असतात). हे मिळाल्यानंतर आणखी गोष्टीही सांगता येतात; उदाहरणार्थ,
$X \times Y = Z$:
\[
\lforall[z][(z \in Z \liff \lexists[x][\lexists[y][(x \in
        X \land y \in Y \land \tuple{x, y} = z)]])]
\]

फलन~$f \colon X \to Y$ याचा संबंध~$f(x) = y$, म्हणजे जोड्यांचा संच
$\Setabs{\tuple{x,y}}{f(x) = y}$, म्हणून विचार करता येतो. मग संच~$f$ हे
$X$ पासून~$Y$ पर्यंतचे फलन आहे असे आपण तेव्हा म्हणू शकतो, जेव्हा (a) तो
$\subseteq X \times Y$ असा संबंध आहे, (b) तो सर्वत्र परिभाषित आहे,
म्हणजे प्रत्येक $x \in X$ साठी असा काही $y \in Y$ आहे की
$\tuple{x, y} \in f$, आणि (c) तो एकमूल्यी आहे, म्हणजे
$\tuple{x, y}, \tuple{x, y'} \in f$ असताना $y = y'$ (कारण फलनांची मूल्ये
एकमेव असली पाहिजेत). म्हणून ``$f$ हे~$X$ पासून~$Y$ पर्यंतचे फलन आहे'' हे
पुढीलप्रमाणे लिहिता येते:
\begin{align*}
 \lforall[u][(u \in f \lif {}] & \lexists[x][\lexists[y][(x \in X \land y \in
      Y \land \tuple{x, y} = u)]]) \land {}\\
 \lforall[x][(x \in X \lif {}] &
      (\lexists[y][(y \in Y \land \mathrm{maps}(f, x, y))] \land {}\\
& (\lforall[y][\lforall[y'][((\mathrm{maps}(f, x, y) \land
    \mathrm{maps}(f, x, y')) \lif y = y')]]))
\end{align*}
येथे $\mathrm{maps}(f, x, y)$ हे $\lexists[v][(v \in f \land
  \tuple{x, y} = v)]$ चे संक्षिप्त रूप आहे (हे सूत्र ``$f(x) = y$''
व्यक्त करते).

आता, उदाहरणार्थ, $f\colon X \to Y$ हे एकास-एक आहे हे व्यक्त करणेही
कठीण नाही:
\begin{multline*}
 f \colon X \to Y \land \lforall[x][\lforall[x'][((x \in X \land x' \in
     X \land {}]] \\
 \lexists[y][(\mathrm{maps}(f, x, y) \land \mathrm{maps}(f,
        x', y))]) \lif x = x')
\end{multline*}
फलन~$f\colon X \to Y$ हे एकास-एक तेव्हा आणि केवळ तेव्हाच असते, जेव्हा
$f$ हे $x, x' \in X$ या दोन्हींना एकाच~$y$ शी संगत करत असल्यास $x = x'$.
हे सूत्र~$\mathrm{inj}(f, X, Y)$ असे संक्षिप्तपणे लिहिल्यास, आता आपण संच
उपपत्तीच्या भाषेत कँटरच्या प्रमेयासारखी साधी नसलेली गोष्टही मांडू शकतो:
$\Pow{X}$ पासून~$X$ मध्ये जाणारे कोणतेही एकास-एक फलन नाही:
\[
\lforall[X][\lforall[Y][(\Pow{X} = Y \lif
    \lnot\lexists[f][\mathrm{inj}(f, Y, X)])]]
\]

प्रत्येक परिभाषक गुणधर्मासाठी संचाचे अस्तित्व हमीने देणारे आणखी एक
स्वयंसिद्धक संच उपपत्तीला आवश्यक आहे असे वाटू शकते. $\varphi(x)$ हे मुक्त
चर~$x$ असलेले संच उपपत्तीचे सूत्र असेल, तर पुढील वाक्य विचारात घेता
येते:
\[
\lexists[y][\lforall[x][(x \in y \liff \varphi(x))]].
\]
हे वाक्य सांगते की असा संच~$y$ आहे ज्याचे घटक म्हणजे नेमके
ते सर्व~$x$ होत जे~$\varphi(x)$ ची पूर्ति करतात. या रूपबंधाला ``गुणधर्माधारित
संचरचनेचे तत्त्व'' म्हणतात. ते फार उपयुक्त वाटते; पण दुर्दैवाने ते विसंगत
आहे. $\varphi(x) \ident \lnot x \in x$ घ्या; मग गुणधर्माधारित संचरचनेचे तत्त्व
पुढील विधान सांगते:
\[
\lexists[y][\lforall[x][(x \in y \liff x \notin x)]],
\]
म्हणजे स्वतःचे घटक नसलेल्या सर्व संचांच्या संचाचे अस्तित्व ते
सांगते. असा संच अस्तित्वात असू शकत नाही---ही रसेलची विरोधापत्ती आहे.
प्रत्यक्षात, $\Log{ZFC}$ मध्ये या तत्त्वाची मर्यादित---आणि सुसंगत---आवृत्ती,
म्हणजे विभक्तीकरण तत्त्व, समाविष्ट आहे:
\[
\lforall[z][\lexists[y][\lforall[x][(x \in y \liff (x \in z \land
      \varphi(x))]]].
\]

\begin{prob}
पुढील निष्पन्नता दाखवणारी निष्पत्ती देऊन गुणधर्माधारित संचरचनेचे
तत्त्व विसंगत आहे, हे दाखवा:
\[
\lexists[y][\lforall[x][(x \in y \liff x \notin x)]] \Proves \lfalse.
\]
प्रथम $(A \lif \lnot A) \land (\lnot A \lif A)
\Proves \lfalse$ दाखवल्यास मदत होईल.
\end{prob}











\section{संरचनेचे आकारमान व्यक्त करणे}\label{fol:mat:siz:sec}

\begin{explain}
भाषेतील अतार्किक चिन्हे न वापरताही रचनांचे काही गुणधर्म व्यक्त करता येतात.
उदाहरणार्थ, अशी वाक्ये आहेत की ती संरचनांमध्ये तेव्हा आणि केवळ
तेव्हाच सत्य असतात, जेव्हा त्या संरचनेच्या प्रांत मध्ये
किमान, जास्तीत जास्त किंवा नेमके दिलेल्या संख्येइतके~$n$ घटक असतात.
\end{explain}

\begin{prop}
वाक्य
\begin{multline*}
  \varphi_{\ge n} \ident \lexists[x_1][\lexists[x_2][\dots\lexists[x_n][{}]]]\\
  \begin{aligned}
  (\eq/[x_1][x_2] \land {}
  \eq/[x_1][x_3] \land \eq/[x_1][x_4] \land \dots \land \eq/[x_1][x_n] \land {}\\
\eq/[x_2][x_3] \land \eq/[x_2][x_4] \land \dots \land {} \eq/[x_2][x_n] \land {} \\
\vdots\\
\eq/[x_{n-1}][x_n])
\end{aligned}
\end{multline*}
हे संरचना~$\Struct M$ मध्ये तेव्हा आणि केवळ तेव्हाच सत्य असते,
जेव्हा $\Domain M$ मध्ये किमान $n$ घटक असतात. परिणामी,
$\Sat{M}{\lnot \varphi_{\ge n+1}}$ तेव्हा आणि केवळ तेव्हाच लागू होते, जेव्हा
$\Domain M$ मध्ये जास्तीत जास्त~$n$ घटक असतात.

\end{prop}

\begin{prop}
वाक्य
\begin{multline*}
\varphi_{= n} \ident \lexists[x_1][\lexists[x_2][\dots\lexists[x_n][{}]]] \\
\begin{aligned}
  (\eq/[x_1][x_2] \land {}
  \eq/[x_1][x_3] \land \eq/[x_1][x_4] \land \dots \land \eq/[x_1][x_n] \land {}\\
\eq/[x_2][x_3] \land \eq/[x_2][x_4] \land \dots \land {} \eq/[x_2][x_n] \land {} \\
\vdots\\
\eq/[x_{n-1}][x_n] \land {} \\
\lforall[y][(\eq[y][x_1] \lor \dots \lor \eq[y][x_n]]))
\end{aligned}
\end{multline*}
हे संरचना~$\Struct M$ मध्ये तेव्हा आणि केवळ तेव्हाच सत्य असते,
जेव्हा $\Domain M$ मध्ये नेमके $n$ घटक असतात.
\end{prop}

\begin{prop}
एखादी संरचना अनंत तेव्हा आणि केवळ तेव्हाच असते, जेव्हा ती पुढील
संचाचे प्रतिमान असते:
\[
\{\varphi_{\ge 1}, \varphi_{\ge 2}, \varphi_{\ge 3}, \dots \}.
\]
\end{prop}

असे एकही पूर्णतः तार्किक वाक्य नाही की जे~$\Struct M$ मध्ये तेव्हा आणि केवळ
तेव्हाच सत्य असेल, जेव्हा $\Domain M$ अनंत असेल. पण अशी अतार्किक
विधेये असलेली वाक्ये देता येतात ज्यांची फक्त अनंत प्रतिमाने
असतात (जरी प्रत्येक अनंत संरचना त्यांचे प्रतिमान नसते). सांत रचना
असण्याचा गुणधर्म आणि अगणनीय रचना असण्याचा गुणधर्म
वाक्यांच्या अनंत संचानेही व्यक्त करता येत नाही. ही तथ्ये संहतता आणि
लोव्हेनहाइम--स्कोलेम प्रमेयांवरून निष्पन्न होतात.



\chapter{प्रथम-क्रम तर्कशास्त्रापलीकडे}\label{fol:byd::chap}
\begin{editorial}
  जेरेमी अविगॅड यांच्या तर्कशास्त्रावरील टिपणांवरून रूपांतरित केलेला हा
  अध्याय इतर कोणकोणत्या तार्किक प्रणाली अस्तित्वात आहेत यांची अत्यंत संक्षिप्त
  झलक देतो. ज्या अभ्यासक्रमात या प्रणाली शिकवल्या जात नाहीत, त्यात पुढील
  अध्ययनासाठी विषय सुचवणारा अध्याय म्हणून तो अभिप्रेत आहे. येथे उल्लेखिलेल्या
  प्रत्येक विषयाला पुढे---अशी आशा आहे---ओपन लॉजिक प्रोजेक्टमध्ये स्वतंत्र
  भागाच्या पातळीवरील विवेचन मिळेल.
\end{editorial}









\section{आढावा}\label{fol:byd:int:sec}

प्रथम-क्रम तर्कशास्त्र हीच काही महत्त्वाची एकमेव तार्किक प्रणाली नाही:
प्रथम-क्रम तर्कशास्त्राचे अनेक विस्तार आणि अनेक भेद आहेत. सामान्यतः एखाद्या
तर्कशास्त्रात भाषेचे आकारिक विनिर्देशन, बहुतेकदा---पण नेहमीच नाही---एक
निगामी प्रणाली आणि बहुतेकदा---पण नेहमीच नाही---एक अभिप्रेत चिन्हार्थमीमांसा
यांचा समावेश असतो. पण या संज्ञेच्या तांत्रिक वापरामुळे एक उघड प्रश्न निर्माण
होतो: सहजप्रज्ञ किंवा तात्त्विक अर्थाने वापरल्या जाणाऱ्या ``तर्कशास्त्र'' या
शब्दाशी प्रथम-क्रम तर्कशास्त्र नसलेल्या तर्कशास्त्रांचा संबंध काय? खाली
वर्णिलेल्या सर्व प्रणाली कोणत्या ना कोणत्या प्रकारच्या युक्तिविचाराचे प्रतिमान
देण्यासाठी रचल्या आहेत; पण त्या तार्किक कशामुळे ठरतात, हे आपण सांगू शकतो का?

याची सोपी उत्तरे मिळत नाहीत. ``तर्कशास्त्र'' हा शब्द वेगवेगळ्या अर्थांनी आणि
वेगवेगळ्या संदर्भांत वापरला जातो; आणि ``सत्य'' या संकल्पनेप्रमाणेच या
संकल्पनेचेही अनेक तात्त्विक भूमिकांतून विश्लेषण झाले आहे. उदाहरणार्थ, कोणती
विधाने अनिवार्यतः सत्य, अनुभवपूर्व सत्य, अतार्किक पदांच्या अर्थनिर्धारणापासून
स्वतंत्रपणे सत्य, त्यांच्या आकारामुळे सत्य किंवा भाषिक संकेतामुळे सत्य आहेत
हे ठरवणे हे तार्किक युक्तिविचाराचे ध्येय आहे, असे कोणी मानू शकेल; आणि यांपैकी
प्रत्येक संकल्पना बरीच स्पष्ट करावी लागते. केवळ गणितात वापरल्या जाणाऱ्या
तर्कशास्त्रापुरते लक्ष मर्यादित केले, तरी त्याची व्याप्ती किती याविषयी फारसे
एकमत नाही. उदाहरणार्थ, फ्रेगे यांनी सुरू केलेल्या {\em तर्कमीमांसावादी}
परंपरेत रसेल आणि व्हाइटहेड यांनी \textit{प्रिन्सिपिया मॅथेमॅटिका} मध्ये
तर्कशास्त्राच्या पायावर गणित विकसित करण्याचा प्रयत्न केला. त्यांची तार्किक
प्रणाली खाली वर्णिलेल्या प्रणालीसारख्या उच्च-प्रकार तर्कशास्त्राचे एक रूप
होती. शेवटी त्यांना बहुतेक निकषांनुसार पूर्णतः तार्किक न वाटणारी स्वयंसिद्धके
समाविष्ट करावी लागली (विशेषतः अनंतत्वाचे स्वयंसिद्धक आणि न्यूनीकरणक्षमतेचे
स्वयंसिद्धक); तरीही उच्च-क्रम युक्तिविचाराची काही रूपे तार्किक म्हणून
स्वीकारली पाहिजेत, असे कोणी मानू शकेल. याउलट, ज्यांच्या सत्ताशास्त्रात
``विधानांना'' संवादाच्या वैध वस्तू म्हणून स्थान नाही असे क्वाइन म्हणतात की,
द्वितीय-क्रम आणि उच्च-क्रम तर्कशास्त्र ही मेंढीचे कातडे पांघरलेल्या संच
उपपत्तीचीच रूपे आहेत; म्हणजेच, विधेयांवर संख्यकीकरण करणाऱ्या प्रणाली पूर्णतः
तार्किक नाहीत.

आत्तासाठी असे तात्त्विक प्रश्न पुढे कधीतरी विचारात घेणेच योग्य ठरेल; आणि
खालील प्रणाली तार्किक असोत वा नसोत, विविध प्रकारच्या युक्तिविचारांची आकारिक
आदर्शीकरणे म्हणून त्यांच्याकडे पाहूया.












\section{बहुवर्गीय तर्कशास्त्र}\label{fol:byd:msl:sec}

प्रथम-क्रम तर्कशास्त्रात चरे आणि संख्यापक एकाच प्रांत वर लागू होतात.
पण अनेक (परस्परविभक्त) प्रांत असणे बऱ्याचदा उपयुक्त ठरते: उदाहरणार्थ,
संख्यांचा प्रांत, भूमितीय वस्तूंचा प्रांत, संख्यांपासून संख्यांकडे
जाणाऱ्या फलनांचा प्रांत, आबेली गटांचा प्रांत, इत्यादी ठेवावेसे वाटू
शकतात.

बहुवर्गीय तर्कशास्त्र अशी चौकट पुरवते. सुरुवात ``वर्गांच्या'' यादीने होते---
वस्तूचा ``वर्ग'' ती कोणत्या ``प्रांत'' मध्ये असणे अपेक्षित आहे ते दर्शवतो.
मग प्रत्येक वर्गासाठी स्वतंत्र चले आणि संख्यापक, तसेच (सामान्यतः)
स्वतंत्र एकरूपता असते. फलने आणि संबंध कोणत्या वर्गांतील वस्तू फलसाधक
म्हणून घेऊ शकतात यानुसार तेही ``प्रकारबद्ध'' केलेले असतात. याखेरीज
प्रथम-क्रम तर्कशास्त्राचे नेहमीचे नियम तसेच ठेवले जातात; संख्यापकांचे नियम
प्रत्येक वर्गासाठी स्वतंत्रपणे पुनरावृत्त होतात.

उदाहरणार्थ, आंतरराष्ट्रीय संबंधांचा अभ्यास करण्यासाठी आपण फ्रेंच नागरिक आणि
जर्मन नागरिक असे वस्तूंचे दोन वर्ग असलेली भाषा निवडू शकतो. पहिल्या वर्गातील
वस्तूंसाठी ``द्राक्षारस पितो'' हा एकस्थानी संबंध; दुसऱ्या वर्गातील वस्तूंसाठी
``वुर्स्ट खातो'' हा दुसरा एकस्थानी संबंध; आणि दोन फलसाधक घेणारा ``बहुराष्ट्रीय
विवाहित जोडपे बनवतात'' हा द्विस्थानी संबंध असू शकतो, ज्याचा पहिला फलसाधक
पहिल्या वर्गातील आणि दुसरा फलसाधक दुसऱ्या वर्गातील असतो. फ्रेंच नागरिकांसाठी
$a$, $b$, $c$ ही चरे आणि जर्मन नागरिकांसाठी $x$, $y$, $z$ ही चरे वापरली, तर
\[
\lforall[a][\lforall x][(\Atom{\Obj{MarriedTo}}{a,x} \lif
(\Atom{\Obj{DrinksWine}}{a} \lor \lnot \Atom{\Obj{EatsWurst}}{x})]]
\]
हे सूत्र असे सांगते की कोणत्याही फ्रेंच व्यक्तीचा विवाह एखाद्या जर्मन
व्यक्तीशी झाला असेल, तर एकतर ती फ्रेंच व्यक्ती द्राक्षारस पिते किंवा ती जर्मन
व्यक्ती वुर्स्ट खात नाही.

बहुवर्गीय तर्कशास्त्र प्रथम-क्रम तर्कशास्त्रात स्वाभाविक रीतीने अंतःस्थापित
करता येते: बहुवर्गीय प्रांत मधील सर्व वस्तू एकत्र करून प्रथम-क्रमाचा एक
प्रांत घ्यायचा, वर्गांची नोंद ठेवण्यासाठी एकस्थानी विधेये
वापरायची आणि संख्यापकांना योग्य वर्गापुरते सापेक्ष करायचे. उदाहरणार्थ, वरच्या
उदाहरणाशी संबंधित प्रथम-क्रम भाषेत वर्णिलेल्या इतर संबंधांबरोबर
``$\Obj{German}$'' आणि ``$\Obj{French}$'' ही एकस्थानी विधेये असतील;
मात्र वर्गाविषयीच्या अटी पुसलेल्या असतील. जर्मन वर्गातील चल~$x$
असलेला वर्गबद्ध संख्यापक $\lforall[x][\varphi]$ पुढीलप्रमाणे रूपांतरित होतो:
\[
\lforall[x][(\Atom{\Obj{German}}{x} \lif \varphi)].
\]
वर्ग परस्परांपासून वेगळे आहेत याची खात्री देणारी स्वयंसिद्धकेही जोडावी
लागतात---उदाहरणार्थ,
$\lforall[x][\lnot (\Atom{\Obj{German}}{x} \land
  \Atom{\Obj{French}}{x})]$---तसेच ``द्राक्षारस पितो'' हा संबंध फक्त
$\Atom{\Obj{French}}{x}$ या विधेयाची पूर्ति करणाऱ्या वस्तूंनाच लागू होतो
याची हमी देणारी स्वयंसिद्धकेही जोडावी लागतात. या संकेतांमुळे आणि
स्वयंसिद्धकांमुळे बहुवर्गीय वाक्ये प्रथम-क्रम वाक्ये मध्ये आणि
बहुवर्गीय निष्पत्त्या प्रथम-क्रम निष्पत्त्यांमध्ये रूपांतरित होतात,
हे दाखवणे अवघड नाही. तसेच बहुवर्गीय संरचना संबंधित प्रथम-क्रम
संरचनांमध्ये आणि उलट दिशेनेही ``रूपांतरित'' होतात; म्हणून बहुवर्गीय
तर्कशास्त्रासाठीही आपल्याला संपूर्णता प्रमेय मिळते.












\section{द्वितीय-क्रम तर्कशास्त्र}\label{fol:byd:sol:sec}

द्वितीय-क्रम तर्कशास्त्राची भाषा व्यक्तींच्या प्रांत वरच नव्हे, तर त्या
प्रांत वरील संबंधांवरही संख्यापन करण्याची मुभा देते. प्रथम-क्रम
भाषा~$\Lang{L}$ दिली असता, प्रत्येक $k$ साठी $k$-स्थानी संबंधांवर मान घेणारी
$R$ ही चले जोडली जातात आणि त्या चरांवर संख्यापन करण्याची मुभा दिली
जाते. $R$ हे एखाद्या $k$-स्थानी संबंधाचे चल असेल आणि $t_1$,
\dots,~$t_k$ ही नेहमीची (प्रथम-क्रम) पदे असतील, तर
$\Atom{R}{t_1,\dots,t_k}$ हे अण्वीय सूत्र आहे. याखेरीज सूत्रांचा
संच प्रथम-क्रम तर्कशास्त्राप्रमाणेच, फक्त द्वितीय-क्रम संख्यापनासाठी अतिरिक्त
उपनियम घालून, परिभाषित केला जातो. लक्षात घ्या की आपल्याकडे एकरूपता फक्त
प्रथम-क्रम पदांसाठी आहे: $R$ आणि $S$ ही समान स्थानसंख्या~$k$ असलेल्या
संबंधांची चले असतील, तर $\eq[R][S]$ हे पुढील सूत्राचे संक्षिप्त रूप
म्हणून परिभाषित करता येते:
\[
\lforall[x_1 \dots][\lforall[x_k][(\Atom{R}{x_1, \dots, x_k} \liff
  \Atom{S}{x_1, \dots, x_k})]].
\]

द्वितीय-क्रम तर्कशास्त्राचे नियम संख्यापकांचे नियम नव्या द्वितीय-क्रम चरांपर्यंत
विस्तारित करतात. मात्र ही चरे $\Lang{L}$ मधील विधेये शी आणि अधिक
सामान्यपणे $\Lang{L}$ मधील सूत्रे शी कशी परस्परक्रिया करतात, हे येथे
थोडे काळजीपूर्वक स्पष्ट करावे लागते. किमान चौकटीत संबंधचरे पदे म्हणून गणली
जातात; म्हणून पुढील प्रकारच्या अनुमानांना मुभा असते:
\[
\varphi(R) \Proves \lexists[R][\varphi(R)]
\]
पण $\Lang{L}$ ही अंकगणिताची भाषा असून $<$ हे स्थिर संबंधचिन्ह असेल, तर पुढील
अनुमानही ग्राह्य असावे अशी अपेक्षा असते:
\[
x < y \Proves \lexists[R][\Atom{R}{x,y}]
\]
किंवा दिलेल्या सूत्र~$\varphi$ साठी,
\[
\varphi(x_1, \dots, x_k) \Proves \lexists[R][\Atom{R}{x_1,\dots,x_k}]
\]
आणखी सामान्यपणे, पुढील प्रकारच्या अनुमानांना मुभा द्यावीशी वाटू शकते:
\[
\Subst{\varphi}{\lambd[\vec x][\psi(\vec x)]}{R} \Proves
\lexists[R][\varphi]
\]
येथे $\Subst{\varphi}{\lambd[\vec x][\psi(\vec x)]}{R}$ म्हणजे $\varphi$ मधील
$\Obj{R}{t_1,\dots,t_k}$ या रूपाचे प्रत्येक अण्वीय सूत्र
$\psi(t_1, \dots, t_k)$ ने बदलल्यावर मिळणारे सूत्र. हा शेवटचा नियम
{\em गुणधर्माधारित संबंधरचना-रूपबंध} असण्याशी समतुल्य आहे; म्हणजे
\[
\lexists[R][\lforall[x_1, \dots, x_k][(\varphi(x_1, \dots, x_k) \liff
\Atom{R}{x_1, \dots, x_k})]],
\]
या रूपाचे, द्वितीय-क्रम भाषेतील प्रत्येक सूत्र~$\varphi$ साठी एक,
स्वयंसिद्धक असते; त्यात $R$ हे मुक्त चर नसते. (सराव: $R$ ला $\varphi$ मध्ये येऊ
दिल्यास हा रूपबंध विसंगत होतो, हे दाखवा!)

तर्कशास्त्रज्ञ “द्वितीय-क्रम तर्कशास्त्राची स्वयंसिद्धके” म्हणतात तेव्हा
सहसा प्रथम-क्रम तर्कशास्त्राचा द्वितीय-क्रम संख्यापक-नियमांनी केलेला किमान
विस्तार आणि गुणधर्माधारित संबंधरचना-रूपबंध त्यांना अभिप्रेत असतो. पण या
स्वयंसिद्धकांच्या आणि नियमांच्या दुर्बल उपप्रणालींचा अभ्यास करणेही अनेकदा
उपयुक्त असते. उदाहरणार्थ, गुणधर्माधारित संबंधरचनेचा स्वयंसिद्धक-रूपबंध पूर्ण
सामान्यतेत \emph{स्वसमावेशक} आहे: द्वितीय-क्रम संख्यापक असलेल्या
सूत्र ने “परिभाषित” होणारा $\Atom{R}{x_1, \dots, x_k}$ हा संबंध
अस्तित्वात आहे असे विधान करण्याची तो मुभा देतो; आणि हे संख्यापक अशा सर्व
संबंधांच्या संचावर मान घेतात---ज्या संचात $R$ स्वतःही येतो!{} विसाव्या
शतकाच्या सुमारास रसेलच्या विरोधापत्तीला दिलेली एक सामान्य प्रतिक्रिया म्हणजे
अशा परिभाषांनाच तिच्यासाठी दोषी धरणे आणि गणिताचा पाया विकसित करताना त्या
टाळणे. सूत्र~$\varphi$ मध्ये द्वितीय-क्रम संख्यापक वापरण्यास मनाई केली, तर
गुणधर्माधारित संबंधरचनेचे काहीसे दुर्बल असे {\em स्वसमावेशरहित} रूप मिळते.

चिन्हार्थाच्या दृष्टीने द्वितीय-क्रम संरचनांमध्ये त्या भाषेसाठीची एक
प्रथम-क्रम संरचना आणि ज्या संबंधांवर द्वितीय-क्रम संख्यापक मान घेतात
असा त्या प्रांत वरील संबंधांचा संच असतो (अधिक नेमके सांगायचे, तर प्रत्येक
$k$ साठी स्थानसंख्या~$k$ असलेल्या संबंधांचा एक संच असतो). अर्थात,
निष्पत्ती प्रणालीत गुणधर्माधारित संबंधरचना समाविष्ट असेल, तर
“द्वितीय-क्रम भागात” ती स्वयंसिद्धके पूर्ण करण्याइतके संबंध असावेत ही
अतिरिक्त अट येते---नाहीतर निष्पत्ती प्रणाली निर्दोष राहत नाही!{} पुरेसे
संबंध उपलब्ध असल्याची खात्री करण्याचा एक सोपा मार्ग म्हणजे द्वितीय-क्रम भागात
प्रथम-क्रम भागावरील \emph{सर्व} संबंध घेणे. अशा संरचना ला
\emph{पूर्ण} म्हणतात आणि एका अर्थाने तीच भाषेची “अभिप्रेत संरचना”
असते. आपण फक्त पूर्ण संरचनेचा विचार केला, तर त्याला
\emph{पूर्ण} द्वितीय-क्रम चिन्हार्थमीमांसा म्हणतात. अशा वेळी संरचना
निर्दिष्ट करणे म्हणजे फक्त तिचा प्रथम-क्रम भाग निर्दिष्ट करणे; कारण
द्वितीय-क्रम भागातील घटक त्यावरून अप्रत्यक्षपणे ठरतात.

सारांशतः, द्वितीय-क्रम तर्कशास्त्राविषयी बोलताना काही संदिग्धता असते.
निष्पत्ती प्रणालीच्या दृष्टीने पुढीलपैकी कोणतीही गोष्ट अभिप्रेत असू शकते:
\begin{enumerate}
\item काही गुणधर्माधारित संबंधरचना-स्वयंसिद्धकांसह “किमान” द्वितीय-क्रम
  निष्पत्ती प्रणाली.
\item पूर्ण गुणधर्माधारित संबंधरचना असलेली “मानक” द्वितीय-क्रम
  निष्पत्ती प्रणाली.
\end{enumerate}
चिन्हार्थाच्या दृष्टीने पुढीलपैकी कशातही स्वारस्य असू शकते:
\begin{enumerate}
\item “दुर्बल” चिन्हार्थमीमांसा, जिच्यात संरचनांमध्ये प्रथम-क्रम
  भाग आणि गुणधर्माधारित संबंधरचना-स्वयंसिद्धके पूर्ण करण्याइतका मोठा
  द्वितीय-क्रम भाग असतो.
\item “मानक” द्वितीय-क्रम चिन्हार्थमीमांसा, जिच्यात फक्त पूर्ण
  संरचनेचा विचार केला जातो.
\end{enumerate}
तर्कशास्त्रज्ञांनी कोणती निष्पत्ती प्रणाली किंवा कोणती चिन्हार्थमीमांसा
अभिप्रेत आहे हे सांगितले नाही, तर सहसा प्रत्येक यादीतील दुसरी बाब त्यांना
अभिप्रेत असते. पुढे दिसेल त्याप्रमाणे, ही चिन्हार्थमीमांसा वापरण्याचा फायदा असा
की तिच्यात अनेक स्वाभाविक गणितीय रचनांची समरूपतेपर्यंत एकमेव वर्णने मिळतात;
त्याच वेळी निष्पत्ती प्रणाली बरीच समर्थ असून या चिन्हार्थमीमांसेसाठी
निर्दोष असते. मात्र निष्पत्ती प्रणाली या चिन्हार्थमीमांसेसाठी
\emph{संपूर्ण नाही}; इतकेच नव्हे, तर प्रभावी रीतीने दिलेली \emph{कोणतीही}
निष्पत्ती प्रणाली पूर्ण द्वितीय-क्रम चिन्हार्थमीमांसेसाठी संपूर्ण नसते.
दुसरीकडे, दुर्बल चिन्हार्थमीमांसेसाठी निष्पत्ती प्रणाली \emph{संपूर्ण
आहे}, हे आपण पाहू; म्हणून एखादे वाक्य सिद्ध करता येत नसेल, तर ते असत्य असणारी
\emph{एखादी तरी} संरचना असते, जरी ती पूर्ण रचना असेलच असे नाही.

द्वितीय-क्रम तर्कशास्त्राची भाषा बरीच अभिव्यक्तिसमर्थ आहे. एकस्थानी संबंधांची
प्रांताच्या उपसंचांशी ओळख करता येते; म्हणून विशेषतः या संचांवर संख्यापनही
करता येते. उदाहरणार्थ, स्वाभाविक संख्यांसाठीचे आगमन एकाच स्वयंसिद्धकाने
व्यक्त करता येते:
\[
\lforall[R][((\Atom{R}{\Obj{0}} \land \lforall[x][(\Atom{R}{x} \lif
    \Atom{R}{x'})]) \lif \lforall[x][\Atom{R}{x}])].
\]
अंकगणिताच्या भाषेत $\Obj 0, \Obj \prime, +, \times$ आणि $<$ ही चिन्हे घेतली,
तर त्यांचे वर्तन वर्णिण्यासाठी पुढील स्वयंसिद्धके जोडता येतात:
\begin{enumerate}
\item $\lforall[x][\lnot x' = \Obj 0]$
\item $\lforall[x][\lforall[y][(s(x) = s(y) \lif x = y)]]$
\item $\lforall[x][(x + \Obj 0) = x]$
\item $\lforall[x][\lforall[y][(x + y') = (x + y)']]$
\item $\lforall[x][(x \times \Obj 0) = \Obj 0]$
\item $\lforall[x][\lforall[y][(x \times y') = ((x \times y) + x)]]$
\item $\lforall[x][\lforall[y][(x < y \liff \lexists[z][y = (x + z')])]]$
\end{enumerate}
आपण पूर्ण द्वितीय-क्रम चिन्हार्थमीमांसा वापरत असू, तर ही स्वयंसिद्धके आणि
वरील आगमनाचे स्वयंसिद्धक मिळून अंकगणिताचे मानक प्रतिमान असलेल्या
संरचना~$\Struct{N}$ चे समरूपतेपर्यंत एकमेव वर्णन देतात, हे दाखवणे
अवघड नाही. ही स्वयंसिद्धके सत्य असलेली कोणतीही संरचना~$\Struct{M}$
दिली असता, $\Nat$ वरील नेहमीचे पुनरावर्तन वापरून $\Nat$ पासून
$\Struct{M}$ च्या प्रांत कडे जाणारे $f$ हे फलन असे परिभाषित करा की
$f(0) = \Assign{\Obj 0}{M}$ आणि $f(x+1) = \Assign{\prime}{M}(f(x))$.
$\Nat$ वरील नेहमीचे आगमन आणि $\Struct M$ मध्ये स्वयंसिद्धके (1) आणि~(2)
सत्य असण्याचा उपयोग केल्यावर $f$ हे एकास-एक आहे असे दिसते. $f$ हे
आच्छादक आहे हे पाहण्यासाठी, $f$ च्या व्याप्तीत असलेल्या
$\Domain{M}$ मधील घटकांचा संच $P$ घ्या. $\Struct M$ पूर्ण असल्यामुळे $P$
हा द्वितीय-क्रम प्रांत मध्ये असतो. $f$ च्या रचनेवरून
$\Assign{\Obj 0}{M}$ हा $P$ मध्ये असतो आणि $P$ हा
$\Assign{\prime}{M}$ खाली बंद असतो. $\Struct M$ मध्ये आगमनाचे स्वयंसिद्धक
सत्य असणे (विशेषतः $P$ साठी) $P$ हा $\Struct M$ च्या संपूर्ण प्रथम-क्रम
प्रांत इतकाच असल्याची हमी देते. यावरून $f$ हे एकास-एक आच्छादन आहे.
$\Nat$ वर नेहमीचे आगमन वारंवार वापरून $f$ हे रचनारक्षक प्रतिचित्रण आहे हे
दाखवणे यापेक्षा अवघड नाही.

संचसैद्धान्तिक भाषेत फलन हा संबंधाचाच एक विशेष प्रकार असतो; उदाहरणार्थ,
$\lforall[x][\lexists![y][R(x,y)]]$ पूर्ण करणाऱ्या $R$ या द्विस्थानी
संबंधाशी एकस्थानी फलन~$f$ ची ओळख करता येते. त्यामुळे फलनांवरही संख्यापन करता
येते. मग पूर्ण चिन्हार्थमीमांसा वापरून अनंत संरचनेचा वर्ग असा
परिभाषित करता येतो: त्या संरचना $\Struct M$ असतात ज्यांच्यासाठी
$\Struct M$ च्या प्रांत पासून त्याच्याच उचित उपसंचाकडे जाणारे एक
एकास-एक फलन असते:
\[
\lexists[f][(\lforall[x][\lforall[y][(\eq[f(x)][f(y)] \lif \eq[x][y])]]
  \land \lexists[y][\lforall[x][\eq/[f(x)][y]]])].
\]
या वाक्याचा निषेध मग सांत संरचनेचा वर्ग परिभाषित करतो.

याखेरीज, रेषीय क्रमाच्या परिभाषेत पुढील अट जोडून सुक्रमांचा वर्ग परिभाषित
करता येतो:
\[
\lforall[P][(\lexists[x][\Atom{P}{x}] \lif \lexists[x][(\Atom{P}{x}
    \land \lforall[y][(y < x \lif \lnot \Atom{P}{y})])])].
\]
“संच” आणि “एकस्थानी संबंध” यांची ओळख केल्यास, प्रत्येक रिक्तेतर संचाला
लघुतम घटक असतो असे हे सांगते. आणखी एका उदाहरणात, शिखरांची परस्पर न जोडलेल्या
भागांत कोणतीही अ-तुच्छ विभागणी नसते असे सांगून आलेखांची संयोजकता व्यक्त करता
येते:
\[
\lnot \lexists[A][(\lexists[x][A(x)] \land \lexists[y][\lnot A(y)]
  \land \lforall[w][\lforall[z][((\Atom{A}{w} \land \lnot \Atom{A}{z})
      \lif \lnot \Atom{R}{w,z})]])].
\]
दुसऱ्या उदाहरणासाठी, ज्या सांत संरचनेच्या प्रांताचा आकार सम
आहे त्यांचा वर्ग परिभाषित करण्याचा सराव तुम्ही करू शकता. आणखी लक्षवेधक
म्हणजे, परिमेय संख्या अंतर्भूत करणारे संपूर्ण क्रमित क्षेत्र म्हणून वास्तव
संख्यांचे समरूपतेपर्यंत एकमेव वर्णन देता येते.

थोडक्यात, द्वितीय-क्रम तर्कशास्त्र प्रथम-क्रम तर्कशास्त्रापेक्षा खूपच अधिक
अभिव्यक्तिसमर्थ आहे. ही चांगली बाजू; आता अडचण पाहू. पूर्ण द्वितीय-क्रम
चिन्हार्थमीमांसेसाठी संपूर्ण असणारी प्रभावी निष्पत्ती प्रणाली नाही, हे
आपण आधीच नमूद केले आहे. प्रथम-क्रम तर्कशास्त्राचे अनेक गुणधर्म येथे उपलब्ध
नसतात; त्यांत संहतता आणि लोव्हेनहाइम--स्कोलेम प्रमेयांचा समावेश होतो.

दुसरीकडे, पूर्ण द्वितीय-क्रम चिन्हार्थमीमांसा सोडून दुर्बल चिन्हार्थमीमांसा
स्वीकारली, तर किमान द्वितीय-क्रम निष्पत्ती प्रणाली या
चिन्हार्थमीमांसेसाठी संपूर्ण असते. दुसऱ्या शब्दांत, $\Proves$ चा अर्थ
“किमान प्रणालीत सिद्ध होते” आणि $\Entails$ चा अर्थ “दुर्बल
चिन्हार्थमीमांसेत तार्किकरीत्या निष्पन्न होते” असा घेतल्यास, $\Gamma \Entails
\varphi$ असेल तेव्हा $\Gamma \Proves \varphi$ असते, हे दाखवता येते.
निष्पत्ती प्रणालीत ठरावीक गुणधर्माधारित संबंधरचना-स्वयंसिद्धके समाविष्ट
करायची असतील, तर ही स्वयंसिद्धके पूर्ण करणाऱ्या द्वितीय-क्रम रचनांपुरती
चिन्हार्थमीमांसा मर्यादित करावी लागते: उदाहरणार्थ, $\Delta$ हा
गुणधर्माधारित संबंधरचना-स्वयंसिद्धकांचा संच (शक्यतो त्यांपैकी सर्वांचा संच)
असेल, तर $\Gamma \cup \Delta \Entails \varphi$ असल्यास $\Gamma \cup \Delta
\Proves \varphi$ असते. विशेषतः, आपण विचारात घेत असलेली गुणधर्माधारित
संबंधरचना-स्वयंसिद्धके वापरून $\varphi$ सिद्ध करता येत नसेल, तर ती स्वयंसिद्धके
सत्य असूनही $\lnot \varphi$ सत्य असणारे एक प्रतिमान असते.

दुर्बल चिन्हार्थमीमांसेसाठी संपूर्णता प्रमेय का लागू होते हे पाहण्याचा सर्वांत
सोपा मार्ग म्हणजे द्वितीय-क्रम तर्कशास्त्राचा पुढीलप्रमाणे बहुवर्गीय
तर्कशास्त्र म्हणून विचार करणे. एका वर्गाचा अर्थ नेहमीचा “प्रथम-क्रम”
प्रांत असा लावला जातो; मग प्रत्येक $k$ साठी “स्थानसंख्या $k$ असलेल्या
संबंधांचा” प्रांत असतो. भाषेत
“$\Atom{\Obj{true}_k}{R,x_1,\dots,x_k}$” ही अंगभूत संबंधचिन्हे घेतो. येथे
$R$ हे “$k$-स्थानी संबंध” या वर्गाचे चर आणि $x_1$, \dots,~$x_k$ या
प्रथम-क्रम वर्गातील वस्तू असताना, $R$ हा संबंध $x_1$, \dots,~$x_k$ यांना
लागू होतो असे हे चिन्ह सांगते.

ही ओळख स्वीकारल्यावर दुर्बल द्वितीय-क्रम चिन्हार्थमीमांसा मूलतः बहुवर्गीय
तर्कशास्त्राच्या नेहमीच्या चिन्हार्थमीमांसेसारखीच असते; आणि बहुवर्गीय
तर्कशास्त्राचे प्रथम-क्रम तर्कशास्त्रात अंतःस्थापन करता येते, हे आपण आधीच
पाहिले आहे. त्यामुळे दोन्ही दिशांतील रूपांतरणांचा विचार करता, द्वितीय-क्रम
तर्कशास्त्राची दुर्बल संकल्पना हे प्रत्यक्षात प्रथम-क्रम तर्कशास्त्राचेच
छद्मरूप आहे; त्यात प्रांत मध्ये योग्य स्वयंसिद्धकांच्या अधीन असलेल्या
“वस्तू” आणि “संबंध” या दोन्हींचा समावेश असतो.












\section{उच्च-क्रम तर्कशास्त्र}\label{fol:byd:hol:sec}

प्रथम-क्रम तर्कशास्त्रातून द्वितीय-क्रम तर्कशास्त्राकडे गेल्यामुळे औपचारिक
भाषेच्या चौकटीत प्रथम-क्रम प्रांत मधील वस्तूंच्या संचांविषयी बोलणे शक्य
झाले. तेथेच का थांबावे? उदाहरणार्थ, तृतीय-क्रम तर्कशास्त्रामुळे वस्तूंच्या
संचांचे संच, किंवा वस्तू आणि वस्तूंचे संच हे दोन्ही अंतर्भूत करणारे संचही
हाताळता यावेत. चतुर्थ-क्रम तर्कशास्त्रामुळे अशा प्रकारच्या वस्तूंच्या
संचांविषयी बोलता येईल. तुमच्या लक्षात आले असेलच की ही कल्पना कितीही वेळा
पुनरावृत्त करता येते.

व्यवहारात उच्च-क्रम तर्कशास्त्र संबंधांऐवजी फलनांच्या मदतीने
सूत्रted केले जाते. (स्वाभाविक ओळखी गृहीत धरल्यास हा फरक महत्त्वाचा
नाही.) काही मूलभूत “वर्ग” $A$, $B$, $C$,~\dots (ज्यांना आता आपण “प्रकार”
म्हणू) दिले असता, पुढील अट घालून नवे प्रकार तयार करता येतात:
\begin{quote}
$\sigma$ आणि $\tau$ हे सांत प्रकार असतील, तर $\sigma \to \tau$ हाही सांत
प्रकार असतो.
\end{quote}
प्रकारांना विन्यासात्मक “लेबले” समजा; आपल्या प्रांत मध्ये हव्या असलेल्या
वस्तूंचे ती वर्गीकरण करतात. $\sigma \to \tau$ हा प्रकार, प्रकार~$\sigma$ च्या
वस्तू घेऊन प्रकार~$\tau$ च्या वस्तू देणारी फलने असलेल्या वस्तूंचे वर्णन करतो.
उदाहरणार्थ, “सत्य” आणि “असत्य” या सत्यतामूल्यांचा $\Omega$ हा प्रकार आणि
नैसर्गिक संख्यांचा $\Nat$ हा प्रकार घ्यावासा वाटू शकतो. अशा वेळी $\Nat \to
\Omega$ प्रकारच्या वस्तू एकस्थानी संबंध किंवा $\Nat$ चे उपसंच मानता येतात;
$\Nat \to \Nat$ प्रकारच्या वस्तू ही नैसर्गिक संख्यांपासून नैसर्गिक संख्यांकडे
जाणारी फलने असतात; आणि $(\Nat \to \Nat) \to \Nat$ प्रकारच्या वस्तू
“फलनके” असतात, म्हणजे फलने घेऊन संख्या देणारी उच्च-प्रकार फलने.

द्वितीय-क्रम तर्कशास्त्राप्रमाणेच, उच्च-क्रम तर्कशास्त्राकडे बहुवर्गीय
तर्कशास्त्राचा एक प्रकार म्हणून पाहता येते; त्यात विचारात घ्यायच्या प्रत्येक
वस्तुप्रकारासाठी एक वर्ग असतो. पण उच्च-प्रकार तर्कशास्त्राचा विन्यास
मुळापासूनच परिभाषित करणे सहसा अधिक स्पष्ट ठरते. उदाहरणार्थ, सांत प्रकारांचा
संच पुढीलप्रमाणे विगमनाने परिभाषित करता येतो:
\begin{enumerate}
\item $\Nat$ हा सांत प्रकार आहे.
\item $\sigma$ आणि $\tau$ हे सांत प्रकार असतील, तर $\sigma
  \to \tau$ हाही सांत प्रकार आहे.
\item $\sigma$ आणि $\tau$ हे सांत प्रकार असतील, तर $\sigma \times
  \tau$ हाही सांत प्रकार आहे.
\end{enumerate}
सहज अर्थाने, $\Nat$ नैसर्गिक संख्यांचा प्रकार दर्शवतो, $\sigma \to \tau$
हा $\sigma$ पासून $\tau$ कडे जाणाऱ्या फलनांचा प्रकार दर्शवतो आणि
$\sigma \times \tau$ हा वस्तूंच्या जोड्यांचा प्रकार दर्शवतो, ज्यातील एक वस्तू
$\sigma$ मधील आणि दुसरी $\tau$ मधील असते. मग पदांचा संच पुढीलप्रमाणे
विगमनाने परिभाषित करता येतो:
\begin{enumerate}
\item प्रत्येक प्रकार $\sigma$ साठी $x$, $y$, $z$, \dots ही प्रकार~$\sigma$
  ची अनेक चरे असतात.
\item $\Obj 0$ हे $\Nat$ प्रकारचे पद आहे.
\item $\Obj S$ (उत्तरवर्ती) हे $\Nat \to \Nat$ प्रकारचे पद आहे.
\item $s$ हे $\sigma$ प्रकारचे पद आणि $t$ हे
  $\Nat \to (\sigma \to \sigma)$ प्रकारचे पद असेल, तर $\Obj{R}_{st}$ हे
  $\Nat \to \sigma$ प्रकारचे पद आहे.
\item $s$ हे $\tau \to \sigma$ प्रकारचे पद आणि $t$ हे प्रकार~$\tau$ चे पद
  असेल, तर $s(t)$ हे $\sigma$ प्रकारचे पद आहे.
\item $s$ हे प्रकार~$\sigma$ चे पद आणि $x$ हे प्रकार~$\tau$ चे चर असेल,
  तर $\lambd[x][s]$ हे $\tau \to \sigma$ प्रकारचे पद आहे.
\item $s$ हे प्रकार~$\sigma$ चे पद आणि $t$ हे प्रकार~$\tau$ चे पद असेल,
  तर $\tuple{s, t}$ हे $\sigma \times \tau$ प्रकारचे पद आहे.
\item $s$ हे प्रकार~$\sigma \times \tau$ चे पद असेल, तर $p_1(s)$ हे
  प्रकार~$\sigma$ चे पद आणि $p_2(s)$ हे प्रकार~$\tau$ चे पद आहे.
\end{enumerate}
सहज अर्थाने, $\Obj{R}_{st}$ हे पुढीलप्रमाणे पुनरावर्ती रीतीने परिभाषित केलेले
फलन दर्शवते:
\begin{align*}
\Obj{R}_{st}(0) & = s \\
\Obj{R}_{st}(x+1) & = t(x, R_{st}(x)),
\end{align*}
$\tuple{s, t}$ ही ज्याचा पहिला घटक~$s$ आणि दुसरा घटक~$t$ आहे अशी जोडी
दर्शवते, तर $p_1(s)$ आणि~$p_2(s)$ हे $s$ चे पहिले व दुसरे घटक
(“प्रक्षेपणे”) दर्शवतात. शेवटी, $\lambd[x][s]$ हे
\[
f(x) = s
\]
या अटीने प्रकार~$\tau$ च्या कोणत्याही~$x$ साठी परिभाषित झालेले फलन~$f$
दर्शवते.

म्हणून बाब (6) आपल्याला गुणधर्माधारित फलनरचनेचे एक रूप देते आणि
पदांच्या मदतीने फलने परिभाषित करणे शक्य करते. समान प्रकारच्या पदांमधील
एकरूपता विधाने $\eq[s][t]$, नेहमीचे विधानीय संयोजक आणि उच्च-प्रकार
संख्यापन यांपासून सूत्रे घडवली जातात. वर परिभाषित केलेल्या पदांचे
नियमन करणारी मूलभूत समीकरणे, तसेच संख्यापक आणि एकरूपता असलेले
तर्कशास्त्राचे नेहमीचे नियम, ही मग प्रणालीची स्वयंसिद्धके मानता येतात.

सांत प्रकारांच्या प्रणालीत सत्यतामूल्यांचा $\Omega$ हा प्रकार जोडला, तर
त्याचा वापर नियंत्रित करणारी स्वयंसिद्धकेही समाविष्ट करावी लागतात. खरे तर
योग्य युक्ती वापरून गुंतागुंतीची सूत्रे पूर्णपणे काढून टाकता येतात आणि
त्यांच्या जागी $\Omega$ प्रकारची पदे ठेवता येतात!{} त्यानुसार सिद्धता-पद्धतीतही
बदल करता येतो. त्यातून मूलतः १९३० च्या दशकात ॲलोन्झो चर्च यांनी मांडलेली
\emph{प्रकारांची साधी उपपत्ती} मिळते.

द्वितीय-क्रम तर्कशास्त्राप्रमाणेच, उच्च-प्रकार चिन्हार्थमीमांसेच्याही विविध
आवृत्त्या वापरता येतात. पूर्ण आवृत्तीत $\sigma \to \tau$ प्रकारची चरे,
प्रकार~$\sigma$ च्या वस्तूंपासून प्रकार~$\tau$ च्या वस्तूंकडे जाणाऱ्या
\emph{सर्व} फलनांच्या संचावर मान घेतात. अपेक्षेप्रमाणे ही चिन्हार्थमीमांसा
इतकी समर्थ आहे की तिच्यासाठी संपूर्ण आणि प्रभावी निष्पत्ती प्रणाली
असू शकत नाही. पण दुर्बल चिन्हार्थमीमांसेचाही विचार करता येतो; त्यात
संरचनांमध्ये प्रत्येक प्रकार $\tau$ साठी $T_\tau$ हे घटकांचे संच
आणि उपयोजन, प्रक्षेपण इत्यादींसाठी योग्य कृत्ये असतात. तपशील योग्य रीतीने
मांडल्यास वर वर्णिलेल्या प्रकारच्या निष्पत्ती प्रणालींसाठी संपूर्णता
प्रमेये मिळवता येतात.

उच्च-प्रकार तर्कशास्त्र आकर्षक आहे, कारण गणिताचा मोठा भाग स्वाभाविक रीतीने
अंतःस्थापित करता येईल अशी चौकट ते देते: $\Nat$ पासून सुरुवात करून वास्तव
संख्या, संतत फलने इत्यादी परिभाषित करता येतात. अंतःप्रज्ञावादी तर्कशास्त्राच्या
संदर्भातही ते विशेष आकर्षक आहे, कारण प्रकारांना स्पष्ट “रचनाशील” अर्थनिर्धारणे
आहेत. खरे तर, उच्च-प्रकार चिन्हार्थमीमांसेच्या रचनाशील आवृत्त्या
(अभिजात तर्कशास्त्राऐवजी अंतःप्रज्ञावादी तर्कशास्त्रावर आधारलेल्या) विकसित
करता येतात. त्या ही रचनाशील अर्थनिर्धारणे अतिशय स्पष्ट करतात आणि अनेक
दृष्टींनी त्यांच्या अभिजात समकक्षांपेक्षा अधिक रोचक असतात.












\section{अंतःप्रज्ञावादी तर्कशास्त्र}\label{fol:byd:il:sec}

द्वितीय-क्रम आणि उच्च-क्रम तर्कशास्त्राच्या उलट, अंतःप्रज्ञावादी प्रथम-क्रम
तर्कशास्त्र हे अभिजात आवृत्तीचे एक निर्बंधित रूप आहे; अधिक “रचनाशील” प्रकारच्या
तर्कविचाराचे प्रतिमान देणे हा त्याचा उद्देश आहे. त्यामागील काही प्रेरणा पुढील
उदाहरणांनी स्पष्ट होतील.

समजा, एखाद्या दिवशी कोणीतरी तुमच्याकडे येऊन सांगितले की त्यांनी अशी नैसर्गिक
संख्या~$x$ ठरवली आहे की $x$ अविभाज्य असेल तर रीमान गृहीतक सत्य असते आणि $x$
संयुक्त असेल तर रीमान गृहीतक असत्य असते. फारच चांगली बातमी!{} रीमान गृहीतक
सत्य आहे की नाही हा गणितातील मोठ्या अनिर्णित प्रश्नांपैकी एक आहे; आणि येथे तर
ही समस्या केवळ गणनेच्या समस्येत, म्हणजे एखादी ठरावीक संख्या अविभाज्य आहे की
नाही हे ठरवण्यात, रूपांतरित झाल्यासारखी दिसते.

$x$ चे ते अद्भुत मूल्य कोणते? ते त्याचे पुढीलप्रमाणे वर्णन करतात: रीमान गृहीतक
सत्य असेल तर $x$ ही $7$ इतकी आणि अन्यथा $9$ इतकी नैसर्गिक संख्या आहे.

तुम्ही चिडून आपले पैसे परत मागता. अभिजात दृष्टिकोनातून वरील वर्णन $x$ चे
एकमेव मूल्य खरोखर ठरवते; पण तुम्हाला प्रत्यक्षात हवे आहे ते
\emph{स्पष्टपणे} दिलेले $x$ चे मूल्य.

आणखी एक, कदाचित कमी कृत्रिम, उदाहरण पाहू. अपरिमेय संख्येचा परिमेय घात घेऊन
परिमेय उत्तर मिळू शकते, हे आपल्याला माहीत आहे. उदाहरणार्थ,
$\sqrt{2}^2 = 2$. अपरिमेय संख्येचा \emph{अपरिमेय} घात घेऊन परिमेय उत्तर
मिळू शकते की नाही, हे तितके स्पष्ट नाही. पुढील प्रमेय याचे होकारार्थी उत्तर
देते:

\begin{thm}
$a^b$ परिमेय असेल अशा $a$ आणि $b$ या अपरिमेय संख्या आहेत.
\end{thm}

\begin{proof}
$\sqrt{2}^{\sqrt{2}}$ चा विचार करा. ही संख्या परिमेय असेल, तर आपले काम झाले:
$a = b = \sqrt{2}$ घेता येईल. अन्यथा ती अपरिमेय आहे. मग
\[
(\sqrt{2}^{\sqrt{2}})^{\sqrt{2}} = \sqrt{2}^{\sqrt{2} \cdot
  \sqrt{2}} = \sqrt{2}^2 = 2,
\]
आणि ही संख्या निश्चितच परिमेय आहे. म्हणून या प्रकरणात $a$ म्हणून
$\sqrt{2}^{\sqrt{2}}$ आणि $b$ म्हणून~$\sqrt 2$ घ्या.
\end{proof}

याला वैध सिद्धता म्हणता येईल का? बहुतेक गणितज्ञांच्या मते होय. पण यात पुन्हा
काहीसे असमाधानकारक असे काही आहे: विशिष्ट गुणधर्म असलेल्या वास्तव संख्यांच्या
जोडीचे अस्तित्व आपण सिद्ध केले, पण ती संख्यांची \emph{कोणती} जोडी आहे हे
सांगता आलेले नाही. हाच निष्कर्ष अशा रीतीने सिद्ध करता येतो की सिद्धतेत
$a$, $b$ ही जोडी \emph{प्रत्यक्ष दिलेली} असते: $a = \sqrt{3}$ आणि
$b = \log_3 4$ घ्या. मग
\[
a^b = \sqrt{3}^{\log_3 4} = 3^{1/2 \cdot \log_3 4} = (3^{\log_3
  4})^{1/2} = 4^{1/2}= 2,
\]
कारण $3^{\log_3 x} = x$.

पहिल्या सिद्धतेतील पाऊल अमान्य असणाऱ्या तर्कविचाराचे प्रतिमान देण्यासाठी
अंतःप्रज्ञावादी तर्कशास्त्र रचले आहे. $\varphi(x)$ ची पूर्ति करणारा $x$ अस्तित्वात
आहे हे सिद्ध करणे म्हणजे दुसऱ्या सिद्धतेप्रमाणे एखादा ठरावीक~$x$ आणि तो
$\varphi$ ची पूर्ति करतो याची सिद्धता द्यावी लागते. $\varphi$ किंवा $\psi$ सत्य आहे हे
सिद्ध करण्यासाठी त्यांपैकी एकाची सिद्धता देता आली पाहिजे.

आकारिक भाषेत सांगायचे, तर प्रथम-क्रम तर्कशास्त्राची निष्पत्ती प्रणाली
ठरावीक रीतीने निर्बंधित केल्यावर अंतःप्रज्ञावादी प्रथम-क्रम तर्कशास्त्र मिळते.
त्याचप्रमाणे द्वितीय-क्रम किंवा उच्च-क्रम तर्कशास्त्राच्याही अंतःप्रज्ञावादी
आवृत्त्या असतात. गणितीय दृष्टिकोनातून या केवळ आकारिक निगमन-पद्धती आहेत; पण
आधी नमूद केल्याप्रमाणे, एका प्रकारच्या गणितीय तर्कविचाराचे प्रतिमान देणे हा
त्यांचा उद्देश आहे. हा तर्कविचार गणिताविषयीच्या एखाद्या तात्त्विक दृष्टिकोनाने
(उदाहरणार्थ, ब्रौवर यांच्या अंतःप्रज्ञावादाने) समर्थित मानता येतो; किंवा
बिशप यांच्या रचनाशीलतावादाच्या धर्तीवर अधिक “मूर्त” आणि समाधानकारक असा
गणितीय तर्कविचार मानता येतो; तसेच आकारिक वर्णन अनौपचारिक प्रेरणा खरोखर
पकडते की नाही यावर वादही घालता येतो. पण आपली तात्त्विक भूमिका कोणतीही असो,
अंतःप्रज्ञावादी तर्कशास्त्राचा आकारिक रीतीने मांडलेले तर्कशास्त्र म्हणून
अभ्यास करता येतो; आणि कारणे कोणतीही असोत, अनेक गणितीय तर्कशास्त्रज्ञांना तो
अभ्यास रंजक वाटतो.

अंतःप्रज्ञावादी संयोजकांचे एक अनौपचारिक रचनाशील अर्थनिर्धारण आहे. ते सहसा
BHK अर्थनिर्धारण म्हणून ओळखले जाते (ब्रौवर, हीटिंग आणि कोल्मोगोरोव यांच्या
नावांवरून). ते असे आहे: $\varphi \land \psi$ ची सिद्धता म्हणजे $\varphi$ ची सिद्धता
आणि $\psi$ ची सिद्धता यांची जोडी; $\varphi \lor \psi$ ची सिद्धता म्हणजे $\varphi$ ची
सिद्धता किंवा $\psi$ ची सिद्धता, तसेच यांपैकी नेमके कोणते प्रकरण आहे याची स्पष्ट
माहिती; $\varphi \lif \psi$ ची सिद्धता म्हणजे $\varphi$ च्या सिद्धतेचे $\psi$ च्या
सिद्धतेत रूपांतर करणारी प्रक्रिया; $\lforall[x][\varphi(x)]$ ची सिद्धता म्हणजे
$x$ च्या कोणत्याही मूल्यासाठी $\varphi(x)$ ची सिद्धता देणारी प्रक्रिया; आणि
$\lexists[x][\varphi(x)]$ ची सिद्धता म्हणजे $x$ चे एखादे मूल्य आणि ते मूल्य $\varphi$
ची पूर्ति करते याची सिद्धता. या अर्थनिर्धारणाचे संगणकीय वर्णन
“करी--हॉवर्ड समरूपण” किंवा “सूत्रे-हे-प्रकार प्रतिमान” या नावांनी करता
येते: सूत्र एखादा विशिष्ट दत्तांश-प्रकार ठरवते असे समजा आणि
सिद्धतांना त्या दत्तांश-प्रकारांच्या संगणकीय वस्तू समजा; या वस्तूंमुळे संबंधित
सूत्र सत्य आहे हे आपल्याला दिसते.

अंतःप्रज्ञावादी तर्कशास्त्राकडे अनेकदा “विमध्य सिद्धांत वजा केल्यावर उरणारे”
अभिजात तर्कशास्त्र म्हणून पाहिले जाते. पुढील प्रमेय हे अधिक नेमके करते.
\begin{thm}
अंतःप्रज्ञावादी तर्कशास्त्रात पुढील स्वयंसिद्धक-रूपबंध समतुल्य आहेत:
\begin{enumerate}
\item $(\lnot \varphi \lif \lfalse) \lif \varphi$.
\item $\varphi \lor \lnot \varphi$
\item $\lnot \lnot \varphi \lif \varphi$
\end{enumerate}
\end{thm}
कोणत्याही एका रूपबंधापासून इतर दोन्हींचे दृष्टांत मिळवणे हा अंतःप्रज्ञावादी
तर्कशास्त्रातील चांगला सराव आहे.

ब्रौवर यांच्या अंतःप्रज्ञावादाचे आकारीकरण म्हणून मांडलेल्या अंतःप्रज्ञावादी
विधानीय तर्कशास्त्राच्या पहिल्या निगमन-पद्धती कोल्मोगोरोव, ग्लिवेन्को आणि
हीटिंग यांनी स्वतंत्रपणे दिल्या. अंतःप्रज्ञावादी प्रथम-क्रम तर्कशास्त्राचे
(आणि अंतःप्रज्ञावादी गणिताच्या काही भागांचे) पहिले आकारीकरण हीटिंग यांनी
केले. अभिजात तर्कशास्त्रात वैध असलेले अनेक रूपबंध अंतःप्रज्ञावादी
तर्कशास्त्रात वैध नसले, तरी त्यांपैकी अनेक वैधही असतात.

\emph{द्विनकरण रूपांतरण} अभिजात आणि अंतःप्रज्ञावादी तर्कशास्त्रांतील एक
महत्त्वाचा संबंध वर्णिते. ते पुढीलप्रमाणे विगमनाने परिभाषित केले जाते
($\varphi^N$ हे अभिजात सूत्र~$\varphi$ चे “अंतःप्रज्ञावादी” रूपांतर समजा):
\begin{align*}
\varphi^N & \ident \lnot\lnot \varphi \quad \text{अण्वीय सूत्रे $\varphi$ साठी} \\
(\varphi \land \psi)^N & \ident (\varphi^N \land \psi^N) \\
(\varphi \lor \psi)^N & \ident  \lnot\lnot (\varphi^N \lor \psi^N) \\
(\varphi \lif \psi)^N & \ident (\varphi^N \lif \psi^N) \\
(\lforall[x][\varphi])^N & \ident \lforall[x][\varphi^N] \\
(\lexists[x][\varphi])^N & \ident \lnot\lnot\lexists[x][\varphi^N]
\end{align*}
कोल्मोगोरोव आणि ग्लिवेन्को यांनी विधानीय तर्कशास्त्रासाठी या रूपांतरणाच्या
आवृत्त्या दिल्या होत्या; विधेय तर्कशास्त्रासाठी ते गोडेल आणि गेंटझेन यांनी
स्वतंत्रपणे दिले. आपल्याला पुढील निष्कर्ष मिळतो.

\begin{thm}
\begin{enumerate}
\item $\varphi \liff \varphi^N$ हे अभिजात तर्कशास्त्रात सिद्धतायोग्य आहे.
\item $\varphi$ हे अभिजात तर्कशास्त्रात सिद्धतायोग्य असेल, तर $\varphi^N$ हे
  अंतःप्रज्ञावादी तर्कशास्त्रात सिद्धतायोग्य आहे.
\end{enumerate}
\end{thm}

आता पुढील संवादाची कल्पना करता येते. अभिजात गणितज्ञ: “मी $\varphi$ सिद्ध केले
आहे!” अंतःप्रज्ञावादी गणितज्ञ: “तुमची सिद्धता वैध नाही. तुम्ही प्रत्यक्षात
$\varphi^N$ सिद्ध केले आहे.” अभिजात गणितज्ञ: “मला तेही चालेल!” अभिजात
गणितज्ञाच्या दृष्टीने अंतःप्रज्ञावादी गणितज्ञ क्षुल्लक भेद करत आहे, कारण ती
दोन्ही समतुल्य आहेत. पण अंतःप्रज्ञावादी गणितज्ञाच्या मते त्यांत फरक आहे.

लक्षात घ्या की वरील रूपांतरण फक्त शुद्ध तर्कशास्त्राशी संबंधित आहे; अभिजात
आणि अंतःप्रज्ञावादी गणितासाठी योग्य \emph{अतार्किक} स्वयंसिद्धके कोणती किंवा
त्यांच्यात कोणता संबंध आहे, या प्रश्नाला ते हात घालत नाही. पण वरील प्रमेयाचा
पुढील किंचित विस्तार काही उपयुक्त माहिती देतो:

\begin{thm}
$\Gamma$ ने $\varphi$ अभिजात रीतीने सिद्ध होत असेल, तर $\Gamma^N$ ने $\varphi^N$
अंतःप्रज्ञावादी रीतीने सिद्ध होते.
\end{thm}

दुसऱ्या शब्दांत, काही गृहीतकांपासून $\varphi$ अभिजात रीतीने सिद्धतायोग्य असेल,
तर त्यांच्या द्विनकरण रूपांतरांपासून $\varphi^N$ अंतःप्रज्ञावादी रीतीने
सिद्धतायोग्य असते.

एखादे वाक्य किंवा विधानीय सूत्र अंतःप्रज्ञावादी रीतीने वैध आहे हे
दाखवण्यासाठी फक्त त्याची सिद्धता द्यावी लागते. पण ते वैध नाही हे कसे
दाखवणार? त्यासाठी आपल्याला निर्दोष, आणि शक्यतो संपूर्ण, चिन्हार्थमीमांसा हवी.
क्रिप्के यांनी दिलेली चिन्हार्थमीमांसा हे काम उत्तम रीतीने करते.

अभिजात तर्कशास्त्रासाठी केलेला खेळच येथेही खेळता येतो: चिन्हार्थमीमांसा
परिभाषित करायची आणि निर्दोषता व संपूर्णता सिद्ध करायची. मात्र पुढील भेद लक्षात
घेण्यासारखा आहे. अभिजात तर्कशास्त्रात चिन्हार्थमीमांसा एका अर्थाने
संयोजकांच्या अर्थातच अंतर्भूत असलेली “उघड” चिन्हार्थमीमांसा होती. क्रिप्के
चिन्हार्थमीमांसेसाठी काही अंतःप्रज्ञ स्पष्टीकरण देता येत असले, तरी तिच्यात
तशी अपरिहार्यता जाणवत नाही. शिवाय, अभिजात संरचना ही स्वाभाविक गणितीय
संकल्पना आहे; म्हणून संरचना ही संकल्पना अभिजात प्रथम-क्रम
तर्कशास्त्राच्या अभ्यासाचे साधन मानता येते, किंवा अभिजात प्रथम-क्रम
तर्कशास्त्र हे गणितीय संरचनेच्या अभ्यासाचे साधन मानता येते.
याउलट, क्रिप्के संरचना कडे फक्त तार्किक रचना म्हणूनच पाहता येते;
त्यांना स्वतंत्र गणितीय महत्त्व आहे असे दिसत नाही.

विधानीय भाषेसाठीची क्रिप्के संरचना~$\mModel M = \tuple{W, R, V}$
मध्ये $W$ हा संच, $W$ वरील लघुतम घटक असलेला $R$ हा अंशतः क्रम आणि
विधानीय चलांचे $W$ च्या घटक ना केलेले “एकस्वनिक” मूल्यांकन असते.
येथे $W$ चे घटक “जगे” किंवा “ज्ञानावस्था” दर्शवतात; $v \geq u$ हा
घटक $u$ ची “संभाव्य भावी अवस्था” दर्शवतो; आणि $u$ ला दिलेली विधानीय चरे ही
$u$ अवस्थेत सत्य असल्याचे माहीत असलेली विधाने असतात. बलन-संबंध
$\mSat{M}{\varphi}[w]$ मग हा संबंध भाषेतील कोणत्याही सूत्रे पर्यंत
विस्तारित करतो; $\mSat{M}{\varphi}[w]$ चा अर्थ “$w$ अवस्थेत $\varphi$ सत्य आहे” असा
वाचा. हा संबंध पुढीलप्रमाणे विगमनाने परिभाषित केला जातो:
\begin{enumerate}
\item $\mSat{M}{\Obj p_i}[w]$ तेव्हा आणि केवळ तेव्हाच, जेव्हा $\Obj p_i$
  हे $w$ ला दिलेल्या विधानीय चरांपैकी एक असते.
\item $\mSat/{M}{\lfalse}[w]$.
\item $\mSat{M}{(\varphi \land \psi)}[w]$ तेव्हा आणि केवळ तेव्हाच, जेव्हा
  $\mSat{M}{\varphi}[w]$ आणि $\mSat{M}{\psi}[w]$.
\item $\mSat{M}{(\varphi \lor \psi)}[w]$ तेव्हा आणि केवळ तेव्हाच, जेव्हा
  $\mSat{M}{\varphi}[w]$ किंवा $\mSat{M}{\psi}[w]$.
\item $\mSat{M}{(\varphi \lif \psi)}[w]$ तेव्हा आणि केवळ तेव्हाच, जेव्हा
  $w' \geq w$ आणि $\mSat{M}{\varphi}[w']$ असताना नेहमी $\mSat{M}{\psi}[w']$.
\end{enumerate}
प्रतिउदाहरण देणारी क्रिप्के संरचना रचून $\lnot (p \land q) \lif
(\lnot p \lor \lnot q)$ हे अंतःप्रज्ञावादी रीतीने वैध नाही हे दाखवण्याचा
प्रयत्न करणे हा चांगला सराव आहे.












\section{मोडल तर्कशास्त्रे}\label{fol:byd:mod:sec}

पुढील सशर्त वाक्याचे उदाहरण पाहा:
\begin{quote}
  जेरेमी त्या खोलीत एकटा असेल, तर त्याने मद्यपान केलेले आहे, तो विवस्त्र आहे
  आणि खुर्च्यांवर नाचत आहे.
\end{quote}
हे सशर्त विधान वास्तविक अभिव्यंजनाच्या अर्थाने सत्य असू शकते, पण तरीही
दिशाभूल करणारे असू शकते; कारण पूर्वांग असत्य आहे किंवा उत्तरांग सत्य आहे
एवढ्यापेक्षा पूर्वांग
आणि निष्कर्ष यांच्यात अधिक दृढ संबंध आहे असे ते सुचवताना दिसते. म्हणजे, हा
दावा फक्त या विशिष्ट जगातच सत्य नाही (जेरेमी खोलीत एकटा नसल्यामुळे येथे तो
क्षुल्लकपणे सत्य असू शकतो), तर पूर्वांग सत्य \emph{असते}, तर निष्कर्षही सत्य
\emph{असता}, असे त्याची शब्दरचना सुचवते. दुसऱ्या शब्दांत, हे विधान फक्त या
जगात नव्हे, तर कोणत्याही “संभाव्य” जगात सत्य आहे; किंवा ते या विशिष्ट जगात
केवळ सत्य असण्याऐवजी \emph{अनिवार्यपणे} सत्य आहे, असा त्याचा अर्थ घेता येतो.

अशा प्रकारच्या अनिवार्यतेचा अर्थ लावण्यासाठी मोडल तर्कशास्त्र रचले गेले.
नेहमीच्या विधानीय तर्कशास्त्रात चौकट-कारक जोडल्यावर मोडल विधानीय
तर्कशास्त्र मिळते; म्हणजे $\varphi$ हे सूत्र असेल, तर $\Box \varphi$ हेही
सूत्र असते. सहज अर्थाने, $\Box \varphi$ हे $\varphi$ \emph{अनिवार्यपणे} सत्य
आहे, किंवा कोणत्याही संभाव्य जगात सत्य आहे, असे सांगते. $\Diamond \varphi$ हे
सहसा $\lnot \Box \lnot \varphi$ चे संक्षिप्त रूप मानले जाते आणि $\varphi$
\emph{संभाव्यतः} सत्य आहे असे सांगणारे म्हणून वाचता येते. अर्थात, विधेय
तर्कशास्त्रातही मोडलता जोडता येते.

मोडल तर्कशास्त्राची चिन्हार्थमीमांसा देण्यासाठी क्रिप्के संरचना वापरता
येतात; खरे तर क्रिप्के यांनी ही चिन्हार्थमीमांसा प्रथम मोडल तर्कशास्त्र
डोळ्यांसमोर ठेवूनच रचली. अंशतः क्रमांपुरते मर्यादित न राहता, अधिक सामान्यपणे
“संभाव्य जगांचा” $P$ हा संच आणि जगांमधील $\Atom{R}{x,y}$ हा द्विपदी
“प्राप्यता” संबंध घेतला जातो. सहज अर्थाने, $\Atom{R}{p,q}$ हे जग~$q$ हे
जग~$p$ शी सुसंगत आहे असे सांगते; म्हणजे आपण जग~$p$ “मध्ये” असू, तर जग
$q$ सारखे असू शकले असते ही शक्यता विचारात घ्यावी लागते.

मोडल तर्कशास्त्राला कधी कधी “विस्तारात्मक” तर्कशास्त्राच्या विरुद्ध
“आशयात्मक” तर्कशास्त्र म्हणतात. अभिजात तर्कशास्त्रासारख्या विस्तारात्मक
तर्कशास्त्राची अभिप्रेत चिन्हार्थमीमांसा केवळ “प्रत्यक्ष” अशा एका जगाचा
निर्देश करते; पण “आशयात्मक” तर्कशास्त्राची चिन्हार्थमीमांसा अधिक विस्तृत
सत्ताशास्त्रावर अवलंबून असते. अनिवार्यतेची संरचना रचण्याखेरीज,
उपयोजनानुसार $\Box$ आणि $\Diamond$ यांचा नवा अर्थ लावून इतर भाषिक रचनांची
संरचना रचण्यासाठीही मोडलता वापरता येते. उदाहरणार्थ:
\begin{enumerate}
\item सिद्धतायोग्यता तर्कशास्त्रात $\Box \varphi$ चा अर्थ “$\varphi$ सिद्धतायोग्य
  आहे” आणि $\Diamond \varphi$ चा अर्थ “$\varphi$ सुसंगत आहे” असा घेतला जातो.
\item ज्ञानविषयक तर्कशास्त्रात $\Box \varphi$ चा अर्थ “मला $\varphi$ माहीत आहे” किंवा
  “माझा $\varphi$ वर विश्वास आहे” असा घेता येतो.
\item कालिक तर्कशास्त्रात $\Box \varphi$ चा अर्थ “$\varphi$ नेहमी सत्य आहे” आणि
  $\Diamond \varphi$ चा अर्थ “$\varphi$ कधीतरी सत्य असते” असा घेता येतो.
\end{enumerate}

मोडलता हाताळणारे नियम आणि स्वयंसिद्धके तर्कशास्त्रात जोडावीशी वाटतात.
उदाहरणार्थ, $\Log{S4}$ प्रणालीत विधानीय तर्कशास्त्राची नेहमीची स्वयंसिद्धके
आणि नियम, तसेच पुढील स्वयंसिद्धके असतात:
\begin{align*}
& \Box (\varphi \lif \psi) \lif (\Box \varphi \lif \Box \psi)\\
& \Box \varphi \lif \varphi \\
& \Box \varphi \lif \Box \Box \varphi
\intertext{आणि “$\varphi$ पासून $\Box \varphi$ हा निष्कर्ष काढा” हा नियमही असतो.
  $\Log{S5}$ मध्ये पुढील स्वयंसिद्धक जोडले जाते:}
& \Diamond \varphi \lif \Box \Diamond \varphi
\end{align*}
या स्वयंसिद्धकांच्या विविध आवृत्त्या वेगवेगळ्या उपयोजनांसाठी योग्य असू शकतात;
उदाहरणार्थ, S5 हे सहसा तार्किक अनिवार्यतेची संकल्पना लक्षणित करते असे मानतात.
चांगली गोष्ट अशी की क्रिप्के संरचना मधील प्राप्यता संबंध स्वाभाविक
रीतीने निर्बंधित करून, संबंधित निष्पत्ती प्रणाली निर्दोष आणि संपूर्ण
असेल अशी चिन्हार्थमीमांसा सहसा सापडते. उदाहरणार्थ, $\Log{S4}$ अशा क्रिप्के
संरचनेच्या वर्गाशी संबंधित आहे ज्यांत प्राप्यता संबंध परावर्ती आणि
संक्रमक असतो. $\Log{S5}$ अशा क्रिप्के संरचनेच्या वर्गाशी संबंधित
आहे ज्यांत प्राप्यता संबंध {\em वैश्विक} असतो, म्हणजे प्रत्येक जग इतर प्रत्येक
जगापासून प्राप्य असते; म्हणून $\Box \varphi$ तेव्हा आणि केवळ तेव्हाच लागू होते,
जेव्हा $\varphi$ प्रत्येक जगात लागू होते.












\section{इतर तर्कशास्त्रे}\label{fol:byd:oth:sec}

आतापर्यंत तुमच्या लक्षात आले असेलच की नवे तर्कशास्त्र रचणे अवघड नाही. तुम्हीही
स्वतःचा विन्यास तयार करू शकता, निगमन-पद्धती बनवू शकता आणि तिला साजेशी
चिन्हार्थमीमांसा घडवू शकता. निष्पत्ती प्रणाली चिन्हार्थमीमांसेसाठी
संपूर्ण हवी असेल, तर थोडे चातुर्य वापरावे लागू शकते; तसेच तुमचे तर्कशास्त्र
खरोखर रंजक आहे हे जगाला पटवून देण्यासाठी काही प्रयत्न करावे लागू शकतात. पण
त्याबदल्यात त्याच्या गणितीय आणि संगणकीय गुणधर्मांचा शोध घेताना अनेक तास
निर्भेळ आनंद मिळवता येईल.

अलीकडच्या दशकांत आकारिक तर्कशास्त्रांचा खरोखरच स्फोट झाला आहे. संदिग्ध
गुणधर्मांविषयीच्या तर्कविचाराचे प्रतिमान देण्यासाठी अस्पष्ट तर्कशास्त्र रचले
आहे. अनिश्चिततेविषयीच्या तर्कविचाराचे प्रतिमान देण्यासाठी संभाव्यताधारित
तर्कशास्त्र रचले आहे. खंडनीय प्रकारच्या तर्कविचाराचे प्रतिमान देण्यासाठी
पूर्वमान्य तर्कशास्त्रे आणि अ-एकस्वनिक तर्कशास्त्रे रचली आहेत; म्हणजे नवी
माहिती मिळाल्यावर नंतर उलटवता येतील असे “वाजवी” निष्कर्ष. ज्ञानाविषयीच्या
तर्कविचाराचे प्रतिमान देणारी ज्ञानविषयक तर्कशास्त्रे आहेत; कार्यकारण
संबंधांविषयीच्या तर्कविचाराचे प्रतिमान देणारी कार्यकारणविषयक तर्कशास्त्रे
आहेत; आणि नैतिक व नीतिविषयक कर्तव्यांविषयीच्या तर्कविचाराचे प्रतिमान देणारी
“कर्तव्यविषयक” तर्कशास्त्रेही आहेत. या प्रणाली मांडण्यामागील मुख्य प्रेरणा
तात्त्विक, गणितीय किंवा संगणकीय आहे यानुसार, अशा गोष्टींचा अभ्यास गणितीय
तर्कशास्त्र, तात्त्विक तर्कशास्त्र, कृत्रिम बुद्धिमत्ता, संज्ञानविज्ञान किंवा
इतर क्षेत्रांच्या नावाखाली झालेला आढळेल.

ही यादी सतत वाढत जाते आणि शक्यता अंतहीन दिसतात. सर्व मानवी तर्कबुद्धीचे गणनेत
रूपांतर करण्याचे लायप्निट्स यांचे स्वप्न आपण कदाचित कधीच साध्य करणार नाही---
पण त्यामुळे प्रयत्न करणे थांबवण्याचे कारण नाही.



\chapter{प्रतिमान उपपत्तीची मूलतत्त्वे}\label{mod:bas::chap}








\section{न्यूनीकृत आणि विस्तारित रचना}\label{mod:bas:red:sec}

ज्या भाषांत काही चिन्हे समान आहेत अशा भाषांची, तसेच त्या भाषांसाठीच्या
संरचनेची, तुलना करणे अनेकदा उपयुक्त किंवा आवश्यक असते. सर्वांत
नेहमीचे प्रकरण असे असते की भाषा~$\Lang{L}$ मधील सर्व चिन्हे
भाषा~$\Lang{L'}$ मध्येही असतात, म्हणजे $\Lang{L} \subseteq
\Lang{L'}$. मग अतिरिक्त चिन्हांची अर्थनिर्धारणे जोडून आणि समान चिन्हांची
अर्थनिर्धारणे तशीच ठेवून $\Lang{L}$-संरचना~$\Struct{M}$ चा विस्तार
नेहमीच $\Lang{L'}$-संरचनांमध्ये करता येतो. दुसरीकडे,
$\Lang{L}$ मध्ये नसलेल्या चिन्हांची अर्थनिर्धारणे केवळ ``विसरून''
$\Lang{L'}$-रचना~$\Struct{M'}$ पासून $\Lang{L}$-रचना मिळवता येते.

\begin{defn}
\label{mod:bas:red:defn:reduct}
समजा $\Lang L \subseteq \Lang L'$, $\Struct M$ ही
$\Lang L$-संरचना आणि $\Struct M'$ ही $\Lang L'$-संरचना आहे.
$\Struct M$ ही $\Struct M'$ ची $\Lang L$ पर्यंतची \emph{न्यूनीकृत रचना}
आणि $\Struct M'$ ही $\Struct M$ ची $\Lang L'$ पर्यंतची
\emph{विस्तारित रचना} असेल तेव्हा आणि केवळ तेव्हाच पुढील अटी पूर्ण होतात:
\begin{enumerate}
\item $\Domain{M} = \Domain{M'}$
\item प्रत्येक स्थिरांक~$c \in \Lang L$ साठी $\Assign{c}{M} =
  \Assign{c}{M'}$.
\item प्रत्येक फलन~$f \in \Lang L$ साठी $\Assign{f}{M} =
  \Assign{f}{M'}$.
\item प्रत्येक विधेय~$P \in \Lang L$ साठी $\Assign{P}{M} =
  \Assign{P}{M'}$.
\end{enumerate}
\end{defn}

\begin{prop}
\label{mod:bas:red:prop:reduct}
एखादी $\Lang{L}$-संरचना~$\Struct{M}$ ही एखाद्या
$\Lang{L'}$-संरचना~$\Struct{M'}$ ची न्यूनीकृत रचना असेल, तर सर्व
$\Lang{L}$-वाक्ये~$\varphi$ साठी,
\[
\Sat{M}{\varphi} \text{ iff } \Sat{M'}{\varphi}.
\]
\end{prop}

\begin{proof}
  सराव.
\end{proof}

\begin{prob}
\ref{mod:bas:red:prop:reduct} सिद्ध करा.
\end{prob}

\begin{defn}
आपल्याकडे $\Lang{L}$-रचना $\Struct{M}$ असेल, एका $n$-स्थानी
विधेय~$P$ ची भर घालून मिळालेला $\Lang{L}$ चा विस्तार
$\Lang{L'} = \Lang{L} \cup \{P\}$ असेल, आणि $R \subseteq \Domain{M}^n$
हा $n$-स्थानी संबंध असेल, तर $\Assign{P}{M'} = R$ असलेली
$\Struct{M}$ ची विस्तारित रचना~$\Struct{M'}$ दर्शवण्यासाठी आपण
$\Expan{M}{R}$ असे लिहितो.
\end{defn}











\section{उपसंरचना}\label{mod:bas:sub:sec}

एखाद्या संरचना~$\Struct{M}$ चा प्रांत हा दुसऱ्या
$\Struct{M'}$ चा उपसंच असू शकतो. पण $\Struct{M}$ हा $\Struct{M'}$ चा
``भाग'' आहे असे आपण अर्थातच फक्त तेव्हाच मानावे, जेव्हा $\Domain{M}
\subseteq \Domain{M'}$ असण्याबरोबरच भाषेतील चिन्हांचे अर्थनिर्धारण करताना
$\Struct{M}$ आणि $\Struct{M'}$ किमान सामाईक भाग~$\Domain{M}$ वर
``सहमत'' असतात.

\begin{defn}
\label{mod:bas:sub:defn:substructure}
त्याच भाषेसाठी~$\Lang L$ संरचना $\Struct M$ आणि
$\Struct M'$ दिल्या असता, $\Struct M$ ही $\Struct M'$ ची
\emph{उपसंरचना} आणि $\Struct M'$ ही $\Struct M$ ची
\emph{वर्धित रचना} आहे, असे आपण म्हणतो; $\Struct M \substruct \Struct M'$
असे लिहितो; आणि तेव्हा आणि केवळ तेव्हाच पुढील अटी पूर्ण होतात:
\begin{enumerate}
\item $\Domain{M} \subseteq \Domain{M'}$,
\item प्रत्येक स्थिर $c \in \Lang L$ साठी $\Assign{c}{M} =
    \Assign{c}{M'}$;
\item प्रत्येक $n$-स्थानी फलन $f \in \Lang L$ साठी
  सर्व $a_1$, \dots, $a_n \in \Domain{M}$ यांकरिता
  $\Assign{f}{M}(a_1, \dots, a_n) = \Assign{f}{M'}(a_1, \dots, a_n)$.
\item प्रत्येक $n$-स्थानी विधेय $R \in \Lang L$ साठी, सर्व
  $a_1$, \dots, $a_n \in \Domain{M}$ यांकरिता $\langle
  a_1, \dots, a_n\rangle \in \Assign{R}{M}$ तेव्हा आणि केवळ तेव्हाच,
  जेव्हा $\langle a_1, \dots, a_n\rangle \in \Assign{R}{M'}$.
\end{enumerate}
\end{defn}

\begin{rem}
\label{mod:bas:sub:rem:substructure}
भाषेत एकही स्थिर किंवा फलने नसतील, तर कोणताही $N
\subseteq \Domain{M}$ हा $\Assign{R}{N} = \Assign{R}{M} \cap N^n$
असे ठेवून $\Struct M$ ची, प्रांत~$\Domain{N} = N$ असलेली,
उपसंरचना~$\Struct{N}$ निश्चित करतो.
\end{rem}













\section{सांत मर्यादेपलीकडील प्रसरण}\label{mod:bas:ove:sec}

\begin{thm}
\label{mod:bas:ove:overspill} वाक्यांचा एखादा संच $\Gamma$ याला हवे तितके मोठे
सांत प्रतिमान असतील, तर त्याला एक अनंत प्रतिमान असते.
\end{thm}

\begin{proof}
$c_0$, $c_1$, \dots ही गणनीय इतकी नवीन स्थिरे जोडून $\Gamma$ च्या भाषेचा
विस्तार करा आणि $\Gamma \cup \{c_i \neq c_j : i \neq j\}$ हा संच विचारात
घ्या. $\Gamma$ ला हवे तितके मोठे सांत प्रतिमान आहेत, याचा अर्थ प्रत्येक
$m >0$ साठी असा $n\ge m$ असतो की $\Gamma$ ला आकारमान~$n$ चे प्रतिमान
असते. यावरून $\Gamma \cup \{c_i \neq c_j : i \neq j\}$ सांततः
पूर्ततायोग्य आहे. संहततेनुसार $\Gamma \cup \{c_i \neq c_j : i \neq j\}$
याला प्रतिमान $\Struct M$ असते; त्याचा प्रांत अनंत असलाच पाहिजे, कारण ते
$c_i \neq c_j$ या सर्व असमानतांची पूर्ति करते.
\end{proof}

\begin{prop}
\label{mod:bas:ove:inf-not-fo}
कोणत्याही प्रथम-क्रम भाषेत असे कोणतेही वाक्य $\varphi$ नाही की ते
संरचना~$\Struct M$ मध्ये तेव्हा आणि केवळ तेव्हाच सत्य असते, जेव्हा
त्या संरचनेचा प्रांत $\Domain{M}$ अनंत असतो.
\end{prop}

\begin{proof}
असे एखादे $\varphi$ असते, तर त्याचे नकरण $\lnot \varphi$ सर्व आणि फक्त सांत
संरचनांमध्ये सत्य असते. म्हणून त्याला हवे तितके मोठे सांत प्रतिमान
असतात, पण अनंत प्रतिमान नसते; हे \ref{mod:bas:ove:overspill} शी विसंगत आहे.
\end{proof}











\section{समरूपी रचना}\label{mod:bas:iso:sec}

प्रथम-क्रमाच्या संरचना दोन प्रकारे एकसारख्या असू शकतात. त्यांच्या
सारखेपणाचा एक प्रकार असा की त्या तीच वाक्ये सत्य ठरवतात. अशा
संरचना ना आपण \emph{प्राथमिक सममूल्य} म्हणतो. परंतु रचना अतिशय
भिन्न असूनही तीच वाक्ये सत्य ठरवू शकतात---उदाहरणार्थ, एक रचना
गणनीय असू शकते आणि दुसरी नसू शकते. याचे कारण असे की
संरचनेची पुष्कळ वैशिष्ट्ये प्रथम-क्रम भाषांमध्ये व्यक्त करता येत
नाहीत: कधी भाषा पुरेशी समृद्ध नसते, तर कधी लोव्हेनहाइम--स्कोलेम
प्रमेयासारख्या प्रथम-क्रम तर्कशास्त्राच्या मूलभूत मर्यादा आड येतात. म्हणून
संरचनेच्या सारखेपणाचा दुसरा, अधिक काटेकोर प्रकार असा की त्या
मूलतः एकच असतात: त्यांना घडवणाऱ्या वस्तू वेगळ्या असतात, पण त्यांची
रचनात्मक वैशिष्ट्ये वेगळी नसतात. \emph{समरूपते} च्या संकल्पनेने हे नेमके
मांडता येते.

\begin{defn}
\label{mod:bas:iso:defn:elem-equiv}
एकाच भाषा~$\Lang{L}$ साठीच्या दोन संरचना $\Struct{M}$ आणि
$\Struct M'$ दिलेल्या असताना, $\Lang{L}$ मधील प्रत्येक वाक्य~$\varphi$
साठी $\Sat{M}{\varphi}$ तेव्हा आणि केवळ तेव्हाच जेव्हा $\Sat{M'}{\varphi}$ असेल,
तर $\Struct{M}$ ही $\Struct M'$ शी \emph{प्राथमिक सममूल्य} आहे असे म्हणतो
आणि $\Struct{M} \equiv \Struct M'$ असे लिहितो.
\end{defn}

\begin{defn}
\label{mod:bas:iso:defn:isomorphism} एकाच भाषा~$\Lang L$ साठीच्या दोन
संरचना $\Struct{M}$ आणि $\Struct M'$ दिलेल्या असताना, पुढील अटी
पूर्ण करणारे $h \colon \Domain{M} \to \Domain{M'}$ फलन असेल, तर
$\Struct{M}$ ही~$\Struct M'$ शी \emph{समरूपी} आहे असे म्हणतो आणि
$\Struct{M} \simeq \Struct M'$ असे लिहितो:
\begin{enumerate}
\item $h$ हे एकास-एक आहे: $h(x) =
  h(y)$ असेल, तर $x = y$;
\item $h$ हे आच्छादक आहे: प्रत्येक $y \in \Domain{M'}$ साठी
  $h(x) = y$ असा $x \in \Domain{M}$ असतो;
\item \label{mod:bas:iso:defn:iso-const}प्रत्येक स्थिरांक $c$ साठी:
  $h(\Assign{c}{M}) = \Assign{c}{M'}$;
\item \label{mod:bas:iso:defn:iso-pred}प्रत्येक $n$-स्थानी विधेय~$P$ साठी:
  \[
  \tuple{a_1, \dots, a_n}\in \Assign{P}{M} \quad\text{तेव्हा आणि केवळ तेव्हाच जेव्हा}\quad
  \tuple{h(a_1), \dots, h(a_n)} \in \Assign{P}{M'};
  \]
\item \label{mod:bas:iso:defn:iso-func}प्रत्येक $n$-स्थानी फलन $f$ साठी:
  \[
  h(\Assign{f}{M}(a_1, \dots, a_n)) =
  \Assign{f}{M'}(h(a_1), \dots, h(a_n)).
  \]
\end{enumerate}
\end{defn}

\begin{thm}
\label{mod:bas:iso:thm:isom}
$\Struct{M} \iso \Struct M'$ असेल, तर $\Struct{M} \elemequiv
\Struct{M'}$.
\end{thm}

\begin{proof}
$h$ हे $\Struct{M}$ पासून $\Struct M'$ वरचे समरूपण असू द्या. कोणत्याही
मूल्यांकन~$s$ साठी $h \circ s$ हे $h$ आणि $s$ यांचे संयोजन असते, म्हणजे
ते $\Struct{M'}$ मधील असे मूल्यांकन आहे की $(h \circ s)(x) = h(s(x))$.
$t$ आणि $\varphi$ यांवरील विगमनाने पुढील अधिक सबळ दावे सिद्ध करता येतात:
\begin{enumerate}
  \item[a.] $h(\Value{t}{M}[s]) = \Value{t}{M'}[h\circ s]$.
  \item[b.] $\Sat{M}{\varphi}[s]$ तेव्हा आणि केवळ तेव्हाच जेव्हा $\Sat{M'}{\varphi}[h \circ s]$.
\end{enumerate}
पहिला दावा~$t$ च्या गुंतागुंतीवरील विगमनाने सिद्ध होतो.
\begin{enumerate}
\item $t \ident c$ असेल, तर $\Value{c}{M}[s] = \Assign{c}{M}$ आणि
  $\Value{c}{M'}[h \circ s] = \Assign{c}{M'}$. म्हणून
  $h(\Value{t}{M}[s]) = h(\Assign{c}{M}) = \Assign{c}{M'}$
  (\ref{mod:bas:iso:defn:isomorphism} मधील \ref{mod:bas:iso:defn:iso-const} नुसार) $=
  \Value{t}{M'}[h \circ s]$.
\item $t \ident x$ असेल, तर $\Value{x}{M}[s] = s(x)$ आणि
  $\Value{x}{M'}[h \circ s] = h(s(x))$. म्हणून $h(\Value{x}{M}[s]) =
  h(s(x)) = \Value{x}{M'}[h \circ s]$.
\item $t \ident f(t_1, \dots, t_n)$ असेल, तर
  \begin{align*}
    \Value{t}{M}[s] & = \Assign{f}{M}(\Value{t_1}{M}[s], \dots, \Value{t_n}{M}[s]) \quad\text{आणि}\\
  \Value{t}{M'}[h \circ s] & = \Assign{f}{M'}(\Value{t_1}{M'}[h \circ
    s], \dots, \Value{t_n}{M'}[h \circ s]).
  \end{align*}

  प्रत्येक $i$ साठी $h(\Value{t_i}{M}[s])
  = \Value{t_i}{M'}[h\circ s]$ हे विगमन गृहीतक आहे. म्हणून
  \begin{align}
    h(\Value{t}{M}[s])
    & = h(\Assign{f}{M}(\Value{t_1}{M}[s], \dots, \Value{t_n}{M}[s]) \notag\\
    & = \Assign{f}{M'}(h(\Value{t_1}{M}[s]), \dots,
    h(\Value{t_n}{M}[s])) \label{mod:bas:iso:iso-1}\\
    & = \Assign{f}{M'}(\Value{t_1}{M'}[h \circ s], \dots,
    \Value{t_n}{M'}[h \circ s]) \label{mod:bas:iso:iso-2}\\
    & = \Value{t}{M'}[h\circ s] \notag
  \end{align}
  येथे \ref{mod:bas:iso:iso-1} हे \ref{mod:bas:iso:defn:isomorphism} मधील
  \ref{mod:bas:iso:defn:iso-func} वरून, तर \ref{mod:bas:iso:iso-2} हे विगमन गृहीतकावरून मिळते.
\end{enumerate}
(b) भागाचा पुरावा सरावासाठी सोडला आहे.

$\varphi$ हे वाक्य असेल, तर मूल्यांकने~$s$ आणि $h \circ s$ महत्त्वाची राहत
नाहीत; म्हणून $\Sat{M}{\varphi}$ तेव्हा आणि केवळ तेव्हाच जेव्हा
$\Sat{M'}{\varphi}$.
\end{proof}

\begin{prob}
\ref{mod:bas:iso:thm:isom} मधील (b) चा पुरावा सविस्तर पूर्ण करा.
\ref{mod:bas:iso:defn:isomorphism} मध्ये समरूपणाचे वैशिष्ट्य ठरवणाऱ्या
पाचही गुणधर्मांपैकी प्रत्येक गुणधर्म कुठे वापरला आहे, ते नोंदवण्याची
खात्री करा.
\end{prob}

\begin{defn}
रचना $\Struct{M}$ ची \emph{स्वयंरूपता} म्हणजे $\Struct{M}$ पासून
त्याच रचनेवर असलेले समरूपण होय.
\end{defn}

\begin{prob}
कोणतीही रचना $\Struct{M}$ असू द्या. जर $X$ हा $\Struct{M}$ चा परिभाष्य
उपसंच असेल आणि $h$ ही $\Struct{M}$ ची स्वयंरूपता असेल, तर $X
= \Setabs{h(x)}{x \in X}$ हे दाखवा (म्हणजे $h$ अंतर्गत $X$ अपरिवर्तित
राहतो).
\end{prob}













\section{संरचनेची उपपत्ती}\label{mod:bas:thm:sec}

प्रत्येक संरचना~$\Struct{M}$ काही वाक्ये सत्य ठरवते आणि
काही असत्य ठरवते. ती सत्य ठरवत असलेल्या सर्व वाक्यांच्या संचाला
तिची \emph{उपपत्ती} म्हणतात. तो संच खरोखरच एक उपपत्ती असतो, कारण
त्याच्यापासून तार्किकरीत्या निष्पन्न होणारे काहीही त्याच्या सर्व
प्रतिमानांमध्ये, आणि त्यांत~$\Struct{M}$ चाही समावेश होतो, सत्य असले
पाहिजे.

\begin{defn}
  संरचना~$\Struct M$ दिलेली असताना, $\Struct{M}$ ची
  \emph{उपपत्ती} म्हणजे $\Struct{M}$ मध्ये सत्य असलेल्या वाक्ये
  चा संच $\Theory{M}$ होय; म्हणजे, $\Theory{M} =
  \Setabs{\varphi}{\Sat{M}{\varphi}}$.
\end{defn}

अभिप्रेत अर्थनिर्धारण असलेल्या वाक्यांच्या संचांसाठीही आपण
``उपपत्ती'' हा शब्द अनौपचारिकपणे वापरतो, मग ते संच निगामीदृष्ट्या बंद
असोत किंवा नसोत.

\begin{prop}
कोणत्याही $\Struct{M}$ साठी $\Theory{M}$ संपूर्ण असते.
\end{prop}

\begin{proof}
कोणत्याही वाक्य~$\varphi$ साठी एकतर $\Sat{M}{\varphi}$ किंवा
$\Sat{M}{\lnot \varphi}$; म्हणून एकतर $\varphi \in \Theory{M}$ किंवा
$\lnot \varphi \in \Theory{M}$.
\end{proof}

\begin{prop}\label{mod:bas:thm:prop:equiv}
  प्रत्येक $\varphi \in \Theory{M}$ साठी $\Struct{N} \models \varphi$ असेल,
  तर $\Struct{M} \elemequiv \Struct{N}$.
\end{prop}

\begin{proof}
सर्व $\varphi \in \Theory{M}$ साठी $\Sat{N}{\varphi}$ असल्यामुळे
$\Theory{M} \subseteq \Theory{N}$. जर $\Sat{N}{\varphi}$, तर
$\Sat/{N}{\lnot \varphi}$; म्हणून $\lnot \varphi \notin \Theory{M}$.
$\Theory{M}$ संपूर्ण असल्यामुळे $\varphi \in \Theory{M}$. म्हणून
$\Theory{N} \subseteq \Theory{M}$ आणि आपल्याला
$\Struct{M} \elemequiv \Struct{N}$ मिळते.
\end{proof}

\begin{rem}\label{mod:bas:thm:remark:R}
  $\Struct{R} = \langle\Real, <\rangle$ विचारात घ्या. या
  संरचनेचा प्रांत म्हणजे वास्तव संख्यांचा संच~$\Real$ असून,
  ती फक्त एक 2-स्थानी विधेय असलेल्या भाषा साठीची रचना
  आहे; त्या विधेयाचे अर्थनिर्धारण वास्तव संख्यांवरील $<$ संबंध असे केले
  आहे. स्पष्टपणे $\Struct{R}$ ही अगणनीय आहे; पण
  $\Theory{R}$ उघडपणे सुसंगत असल्यामुळे लोव्हेनहाइम--स्कोलेम प्रमेयाने
  तिला एक गणनीय प्रतिमान असते, त्याला $\Struct{S}$ म्हणू.
  \ref{mod:bas:thm:prop:equiv} नुसार $\Struct{R} \equiv \Struct{S}$. शिवाय,
  $\Struct{R}$ आणि $\Struct{S}$ या समरूपी नसल्यामुळे यावरून
  \ref{mod:bas:iso:thm:isom} चा उलट दावा सर्वसाधारणपणे खोटा ठरतो.
\end{rem}











\section{आंशिक समरूपणे}\label{mod:bas:pis:sec}

\begin{defn}
  दोन संरचना $\Struct{M}$ आणि $\Struct{N}$ दिलेल्या असताना,
  $\Struct{M}$ पासून $\Struct{N}$ कडेचे \emph{आंशिक समरूपण} म्हणजे
  $\Domain M$ मधील घटक आदान म्हणून घेऊन $\Domain N$ मध्ये मूल्ये देणारे
  असे सांत अंशतः फलन $p$, जे आपल्या प्रांतावर
  \ref{mod:bas:iso:defn:isomorphism} मधील समरूपणाच्या अटी पूर्ण करते:
  \begin{enumerate}
  \item $p$ हे एकास-एक आहे;
  \item प्रत्येक स्थिरांक~$c$ साठी: $p(\Assign{c}{M})$ परिभाषित
    असेल, तर $p(\Assign{c}{M}) = \Assign{c}{N}$;
  \item प्रत्येक $n$-स्थानी विधेय $P$ साठी: $a_1$, \dots, $a_n$
    हे $p$ च्या प्रांतात असतील, तर $\langle a_1, \dots, a_n\rangle \in
    \Assign P M$ तेव्हा आणि केवळ तेव्हाच जेव्हा $\langle p(a_1), \dots, p(a_n) \rangle
    \in \Assign P N$;
  \item प्रत्येक $n$-स्थानी फलन $f$ साठी: $a_1$, \dots, $a_n$
    हे $p$ च्या प्रांतात असतील, तर
    \begin{multline*}
    p(\Assign f M (a_1, \dots,a_n))\\
    = \Assign f N (p(a_1), \dots, p(a_n)).
    \end{multline*}
  \end{enumerate}
  $p$ सांत असण्याचा अर्थ $\dom{p}$ सांत असतो.
\end{defn}

लक्षात घ्या की रिक्त फलन~$\emptyset$ हे कोणत्याही दोन संरचना
मधील आंशिक समरूपण नेहमीच असते.

\begin{defn}\label{mod:bas:pis:defn:partialisom}
  दोन संरचना $\Struct{M}$ आणि $\Struct{N}$ यांना
  \emph{आंशिकरीत्या समरूपी} म्हणतात आणि $\Struct{M} \iso[p]
  \Struct{N}$ असे लिहितात, तेव्हा आणि केवळ तेव्हाच जेव्हा
  $\Struct{M}$ आणि $\Struct{N}$ यांमधील आंशिक समरूपणांचा अरिक्त संच $I$
  पुढील \emph{पुढे-मागे} गुणधर्म पूर्ण करतो:
  \begin{enumerate}
  \item (\emph{पुढे}) प्रत्येक $p \in I$ आणि $a \in \Domain M$ साठी असा
    $q \in I$ असतो की $p \subseteq q$ आणि $a$ हा $q$ च्या प्रांतात
    असतो;
  \item (\emph{मागे}) प्रत्येक $p \in I$ आणि $b \in \Domain N$ साठी असा
    $q \in I$ असतो की $p \subseteq q$ आणि $b$ हा $q$ च्या व्याप्तीत
    असतो.
  \end{enumerate}
\end{defn}

\begin{thm}\label{mod:bas:pis:thm:p-isom1}
  $\Struct{M} \iso[p] \Struct{N}$ असून $\Struct{M}$ आणि
  $\Struct{N}$ या गणनीय असतील, तर $\Struct{M} \iso
  \Struct{N}$.
\end{thm}

\begin{proof}
  $\Struct{M}$ आणि $\Struct{N}$ या गणनीय असल्यामुळे
  $\Domain{M} = \{a_0, a_1, \ldots \}$ आणि $\Domain{N} = \{b_0, b_1,
  \ldots \}$ असे मानू. मनमाना $p_0 \in I$ घेऊन सुरुवात करून, आंशिक
  समरूपणांची वाढती क्रमिका $p_0 \subseteq p_1 \subseteq p_2
  \subseteq \cdots$ पुढीलप्रमाणे परिभाषित करतो:
  \begin{enumerate}
  \item $n+1$ विषम असेल, म्हणजे $n = 2r$ असेल, तर पुढे गुणधर्म वापरून
    असा $p_{n+1} \in I$ शोधा की $p_n \subseteq p_{n+1}$
    आणि $a_r$ हा $p_{n+1}$ च्या प्रांतात असेल;
  \item $n+1$ सम असेल, म्हणजे $n+1 =2r$ असेल, तर मागे गुणधर्म वापरून
    असा $p_{n+1} \in I$ शोधा की $p_n \subseteq p_{n+1}$
    आणि $b_r$ हा $p_{n+1}$ च्या व्याप्तीत असेल.
  \end{enumerate}
आता
\[
p = \bigcup_{n\ge 0} p_n,
\]
असे ठेवले, तर $p$ हे $\Struct{M}$ आणि $\Struct{N}$ यांमधील समरूपण
असते.
\end{proof}

\begin{prob}
  \ref{mod:bas:pis:thm:p-isom1} मध्ये परिभाषित केलेले $p$ खरोखरच
  समरूपण आहे, हे सविस्तर दाखवा.
\end{prob}

\begin{thm}\label{mod:bas:pis:thm:p-isom2}
  $\Struct{M}$ आणि $\Struct{N}$ या निव्वळ संबंधात्मक भाषा
  साठीच्या संरचना आहेत असे समजा (म्हणजे फक्त विधेये
  असलेली आणि फलने किंवा स्थिरे नसलेली भाषा). मग
  $\Struct{M} \iso[p] \Struct{N}$ असेल, तर
  $\Struct{M} \elemequiv \Struct{N}$ सुद्धा असते.
\end{thm}

\begin{proof}
  सूत्रे वरील विगमनाने पुढील गोष्ट दाखवता येते. $a_1$, \dots,
  $a_n$ आणि $b_1$, \dots, $b_n$ यांसाठी प्रत्येक $a_i$ ला $b_i$ कडे
  नेणारे आंशिक समरूपण $p$ असेल, तसेच $s_1(x_i) =a_i$ आणि $s_2(x_i) =b_i$
  ($i =1$, \dots,~$n$ साठी) असेल, तर $\Sat{M}{\varphi}[s_1]$ तेव्हा आणि
  केवळ तेव्हाच जेव्हा $\Sat{N}{\varphi}[s_2]$. $n=0$ चे प्रकरण
  $\Struct{M} \elemequiv \Struct{N}$ देते.
\end{proof}

\begin{rem}
फलने उपस्थित असतील, तरी मागील निष्कर्ष सत्य असतो; पण $a_1$,
\dots, $a_n$ यांनी निर्मित $\Struct{M}$ च्या subसंरचना आणि
$b_1$, \dots, $b_n$ यांनी निर्मित $\Struct{N}$ च्या
subसंरचना यांमध्ये $p$ ने प्रेरित केलेले समरूपण विचारात घ्यावे
लागते.
\end{rem}

मागील निष्कर्षाचे टप्प्यांत ``विघटन'' करता येते: सूत्रामध्ये
एकमेकांत अंतर्भूत संख्यापकांची संख्या आणि संबंधित आंशिक समरूपणांचा
किती वेळा विस्तार करता येतो यांचा संबंध प्रस्थापित करावा लागतो.

\begin{defn}
  कोणतेही सूत्र~$\varphi$ दिले असता, $\varphi$ चा
  \emph{संख्यापक-प्रतांक}, $\QuantRank{\varphi} \in \Nat$, म्हणजे $\varphi$ मध्ये एकमेकांत अंतर्भूत
  संख्यापकांची कमाल संख्या; तो पुनरावर्ती रीतीने परिभाषित केला जातो.
  दोन संरचना $\Struct{M}$ आणि $\Struct{N}$ यांना
  \emph{$n$-सममूल्य} म्हणतात आणि $\Struct{M} \elemequiv[n] \Struct{N}$
  असे लिहितात, जर संख्यापक-प्रतांक जास्तीत जास्त~$n$ असलेल्या सर्व
  वाक्यांबाबत त्या एकमत असतील.
\end{defn}

\begin{prop}\label{mod:bas:pis:prop:qr-finite}
  $\Lang{L}$ ही सांत निव्वळ संबंधात्मक भाषा असू द्या, म्हणजे
  सांत इतकी विधेये आणि स्थिरांक असलेली, पण फलने
  नसलेली भाषा. मग प्रत्येक $n \in \Nat$ साठी, तार्किक
  सममूल्यतेपर्यंत, भाषा $\Lang{L}$ मध्ये संख्यापक-प्रतांक
  जास्तीत जास्त $n$ असलेली प्रथम-क्रम वाक्ये सांत इतकीच असतात.
\end{prop}

\begin{proof}
  $n$ वरील विगमनाने.
\end{proof}

\begin{defn}
  संरचना~$\Struct{M}$ दिलेली असताना, $\Domain M^{<\omega}$ हा
  $\Domain{M}$ वरील सर्व सांत क्रमिकांचा संच असू द्या. घटकांच्या सांत
  क्रमिकांवर $\mathbf{a}, \mathbf{b}, \mathbf{c}, \ldots$ ही चले वापरतो.
  जर $\mathbf{a} \in \Domain{M}^{<\omega}$ आणि $a \in \Domain{M}$,
  तर $\mathbf{a}a$ हे $\mathbf{a}$ आणि $a$ यांची \emph{जोडणी} दर्शवते.
\end{defn}

\begin{defn}
  संरचना $\Struct{M}$ आणि $\Struct{N}$ दिलेल्या असताना, समान
  लांबीच्या क्रमिकांमधील संबंध $I_n \subseteq \Domain M^{<\omega}
  \times \Domain N^{<\omega}$ हे $n$ वरील पुनरावर्तनाने पुढीलप्रमाणे
  परिभाषित करतो:
   \begin{enumerate}
   \item $I_0(\mathbf{a},\mathbf{b})$ तेव्हा आणि केवळ तेव्हाच जेव्हा
     $\mathbf{a}$ आणि $\mathbf{b}$ या अनुक्रमे $\Struct{M}$ आणि
     $\Struct{N}$ मध्ये तीच आण्विक सूत्रे पूर्ण करतात; म्हणजे,
     त्या क्रमिकांची समान लांबी k असेल, $s_1(x_i) = a_i$ आणि $s_2(x_i) =
     b_i$ असेल आणि $\varphi$ हे आण्विक असून त्यातील सर्व चले
     $x_1$, \dots,~$x_k$ यांपैकी असतील, तर $\Sat{M}{\varphi}[s_1]$ तेव्हा आणि
     केवळ तेव्हाच जेव्हा~$\Sat{N}{\varphi}[s_2]$.

   \item $I_{n+1} (\mathbf{a},\mathbf{b})$ तेव्हा आणि केवळ तेव्हाच
     जेव्हा प्रत्येक $a\in \Domain M$ साठी असा $b\in \Domain N$ असतो की
     $I_n (\mathbf{a}a,\mathbf{b}b)$, आणि उलटही.
   \end{enumerate}
\end{defn}


\begin{defn}
  $\Struct{M}$ आणि $\Struct{N}$ यांच्यासाठी
  $I_n(\emptyseq,\emptyseq)$ खरे असेल, तर $\Struct{M} \approx_n
  \Struct{N}$ असे लिहितो (येथे $\emptyseq$ ही रिक्त क्रमिका आहे).
\end{defn}

\begin{thm}\label{mod:bas:pis:thm:b-n-f}
  $\Lang{L}$ ही निव्वळ संबंधात्मक भाषा असू द्या. मग
  $I_n (\mathbf{a},\mathbf{b})$ असल्यास, $\QuantRank{\varphi} \le n$ असलेल्या
  प्रत्येक $\varphi$ साठी $\Sat{M}{\varphi}[\mathbf{a}]$ तेव्हा आणि केवळ तेव्हाच
  जेव्हा $\Sat{N}{\varphi}[\mathbf{b}]$ (येथे पुन्हा, $s(x_i) = a_i$ असणारे
  कोणतेही $s$ जर $\varphi$ पूर्ण करत असेल, तर $\mathbf{a}$ हे $\varphi$ पूर्ण
  करते). शिवाय, $\Lang{L}$ सांत असेल, तर उलट दावाही खरा असतो.
\end{thm}

\begin{proof}
  $I_n(\mathbf{a},\mathbf{b})$ यावरून $\mathbf{a}$ आणि $\mathbf{b}$
  संख्यापक-प्रतांक जास्तीत जास्त $n$ असलेली तीच सूत्रे पूर्ण करतात,
  हा पुरावा $\varphi$ वरील सोप्या विगमनाने मिळतो. उलट दाव्यासाठी आपण $n$
  वरील विगमन करतो आणि \ref{mod:bas:pis:prop:qr-finite} वापरतो; त्यावरून प्रत्येक
  $n$ साठी त्या संख्यापक-प्रतांकाची परस्पर तार्किकदृष्ट्या असममूल्य
  सूत्रे जास्तीत जास्त सांत इतकी असतात.

  $n=0$ साठी, $\mathbf{a}$ आणि $\mathbf{b}$ ही तीच संख्यापकमुक्त
  सूत्रे पूर्ण करतात या गृहीतकावरून ती तीच आण्विक सूत्रे पूर्ण करतात;
  म्हणून $I_0(\mathbf{a},\mathbf{b})$.

  $n+1$ च्या प्रकरणासाठी, $\mathbf{a}$ आणि $\mathbf{b}$ ही
  संख्यापक-प्रतांक जास्तीत जास्त $n+1$ असलेली तीच सूत्रे पूर्ण
  करतात असे समजा. $I_{n+1}(\mathbf{a},\mathbf{b})$ दाखवण्यासाठी प्रत्येक
  $a \in \Domain M$ साठी असा $b \in \Domain N$ असतो की
  $I_n(\mathbf{a}a,\mathbf{b}b)$, एवढे दाखवणे पुरेसे आहे; आणि विगमन
  गृहीतकाने प्रत्येक $a \in \Domain M$ साठी असा $b \in \Domain N$ असतो
  की $\mathbf{a}a$ आणि $\mathbf{b}b$ ही संख्यापक-प्रतांक जास्तीत जास्त
  $n$ असलेली तीच सूत्रे पूर्ण करतात, एवढे दाखवणे पुन्हा पुरेसे आहे.

  $a \in \Domain M$ दिलेला असताना, $\Struct{M}$ मध्ये $\mathbf{a}a$ ने
  पूर्ण केलेल्या, प्रतांक जास्तीत जास्त $n$ असलेल्या
  $\psi(x,\mathbf{y})$ या सूत्रांचा संच $!T^a_n$ असू द्या.
  $!T^a_n$ सांत आहे, म्हणून त्यातील सूत्रांचा संयोग असलेले एकच प्रथम-क्रम
  सूत्र म्हणून त्याचा विचार करू शकतो. यावरून $\mathbf{a}$ हे
  $\lexists[x][!T^a_n(x,\mathbf{y})]$ पूर्ण करते; या सूत्राचा
  संख्यापक-प्रतांक जास्तीत जास्त $n+1$ आहे. गृहीतकानुसार $\mathbf{b}$ हे
  $\Struct{N}$ मध्ये तेच सूत्र पूर्ण करते; म्हणून असा $b \in
  \Domain N$ असतो की $\mathbf{b}b$ हे $!T^a_n$ पूर्ण करते. विशेषतः,
  $\mathbf{b}b$ ही $\mathbf{a}a$ प्रमाणेच संख्यापक-प्रतांक जास्तीत जास्त
  $n$ असलेली सर्व सूत्रे पूर्ण करते. त्याचप्रमाणे, प्रत्येक $b \in
  \Domain N$ साठी असा $a\in \Domain M$ असतो की $\mathbf{a}a$ आणि
  $\mathbf{b}b$ ही संख्यापक-प्रतांक जास्तीत जास्त $n$ असलेली तीच
  सूत्रे पूर्ण करतात, हे दाखवता येते; त्यामुळे पुरावा पूर्ण होतो.
\end{proof}

\begin{cor}\label{mod:bas:pis:cor:b-n-f}
  $\Struct{M}$ आणि $\Struct{N}$ या सांत भाषा मधील निव्वळ
  संबंधात्मक संरचना असतील, तर $\Struct{M} \approx_n\Struct{N}$
  तेव्हा आणि केवळ तेव्हाच जेव्हा $\Struct{M} \elemequiv[n] \Struct{N}$.
  विशेषतः, $\Struct{M} \elemequiv \Struct{N}$ तेव्हा आणि केवळ तेव्हाच
  जेव्हा प्रत्येक $n$ साठी $\Struct{M} \approx_n \Struct{N}$.
\end{cor}











\section{सघन रेषीय क्रम}\label{mod:bas:dlo:sec}

\begin{defn}
  \emph{अंतबिंदुविरहित सघन रेषीय क्रम} म्हणजे एकच 2-स्थानी
  विधेय~$<$ असलेल्या भाषा साठीची आणि पुढील वाक्ये पूर्ण
  करणारी संरचना~$\Struct{M}$ होय:
  \begin{enumerate}
  \item $\lforall[x][\lnot x < x]$;
  \item $\lforall[x][\lforall[y][\lforall[z][(x < y \lif (y < z \lif x
    <z ))]]]$;
  \item $\lforall[x][\lforall[y][(x< y \lor \eq[x][y] \lor y < x)]]$;
  \item $\lforall[x][\lexists[y][x < y]]$;
  \item $\lforall[x][\lexists[y][y < x]]$;
  \item $\lforall[x][\lforall[y][(x < y \lif \lexists[z][(x < z \land
        z < y))]]]$.
 \end{enumerate}
\end{defn}

\begin{thm}\label{mod:bas:dlo:thm:cantorQ}
  कोणतेही दोन गणनीय अंतबिंदुविरहित सघन रेषीय क्रम समरूपी
  असतात.
\end{thm}

\begin{proof}
  $\Struct{M_1}$ आणि $\Struct{M_2}$ हे गणनीय अंतबिंदुविरहित
  सघन रेषीय क्रम असू द्या, जेथे ${<_1} = \Assign{<}{M_1}$ आणि ${<_2} =
  \Assign{<}{M_2}$; तसेच $\PIso{I}$ हा त्यांच्यामधील सर्व आंशिक
  समरूपणांचा संच असू द्या. किमान $\emptyset \in \PIso{I}$ असल्यामुळे
  $\PIso{I}$ रिक्त नाही. $\PIso{I}$ हा पुढे-मागे गुणधर्म पूर्ण करतो हे
  आपण दाखवतो. मग $\Struct{M_1} \iso[p] \Struct{M_2}$; आणि
  \ref{mod:bas:pis:thm:p-isom1} नुसार प्रमेय सिद्ध होते.

  $\PIso{I}$ हा पुढे गुणधर्म पूर्ण करतो हे दाखवण्यासाठी $p \in
  \PIso{I}$ असू द्या, $i = 1$, \dots,~$n$ साठी $p(a_i) = b_i$ असू द्या
  आणि सर्वसाधारणतेला बाधा न आणता $a_1 <_1 a_2 <_1 \cdots <_1 a_n$
  असे समजा. $a \in \Domain{M_1}$ दिलेला असताना, तो आधीच त्या नकाशाच्या
  प्रांतात असेल, तर तोच नकाशा अपेक्षित विस्तार आहे. नकाशा रिक्त असेल,
  तर दुसऱ्या रचनेच्या अरिक्त प्रांतातील कोणताही घटक प्रतिमा म्हणून
  निवडून त्याची दिलेल्या घटकाबरोबरची जोडी त्या नकाशात घालता येते.
  उरलेल्या प्रकरणांत
  पुढीलप्रमाणे $b \in \Domain{M_2}$ शोधा:

  \begin{enumerate}
  \item जर $a <_1 a_1$, तर $b <_2 b_1$ असेल असा $b \in
    \Domain{M_2}$ घ्या;
  \item जर $a_n <_1 a$, तर $b_n <_2 b$ असेल असा $b \in \Domain{M_2}$ घ्या;
 \item एखाद्या $i$ साठी $a_i <_1 a <_1 a_{i+1}$ असेल, तर $b_i <_2 b
   <_2 b_{i+1}$ असेल असा $b \in \Domain{M_2}$ घ्या.
  \end{enumerate}
  $\Struct{M_2}$ हा अंतबिंदुविरहित सघन रेषीय क्रम असल्यामुळे अपेक्षित
  गुणधर्म असलेला $b$ नेहमी शोधता येतो. $q = p \cup \{ \langle a, b
  \rangle \}$ असे परिभाषित करा; मग $q \in \PIso{I}$ हा $p$ चा अपेक्षित
  विस्तार आहे. यावरून पुढे गुणधर्म प्रस्थापित होतो. मागे गुणधर्मही
  याचप्रमाणे सिद्ध होतो. म्हणून $\Struct{M_1} \iso[p]
  \Struct{M_2}$; \ref{mod:bas:pis:thm:p-isom1} नुसार $\Struct{M_1} \iso
  \Struct{M_2}$.
\end{proof}

\begin{prob}
  $\PIso{I}$ हा मागे गुणधर्म पूर्ण करतो, हे तपासून
  \ref{mod:bas:dlo:thm:cantorQ} चा पुरावा पूर्ण करा.
\end{prob}

\begin{rem}
  $\Struct{S}$ हा कोणताही गणनीय अंतबिंदुविरहित सघन रेषीय क्रम
  असू द्या. मग (\ref{mod:bas:dlo:thm:cantorQ} नुसार) $\Struct{S} \iso
  \Struct{Q}$, जेथे $\Struct{Q} = (\Rat, <)$ हा परिमेय संख्यांचा संच
  $\Rat$ प्रांत म्हणून असलेला गणनीय सघन रेषीय क्रम आहे. आता
  \ref{mod:bas:thm:remark:R} मधील संरचना~$\Struct{R} = (\Real, <)$
  पुन्हा विचारात घ्या. असा गणनीय संरचना~$\Struct{S}$
  आहे की $\Struct{R} \elemequiv \Struct{S}$, हे आपण पाहिले. पण
  $\Struct{S}$ हा गणनीय अंतबिंदुविरहित सघन रेषीय क्रम आहे;
  म्हणून तो संरचना~$\Struct{Q}$ शी समरूपी (आणि म्हणून प्राथमिक
  सममूल्य) आहे. प्राथमिक सममूल्यतेच्या संक्रमणकतेने $\Struct{R}
  \elemequiv \Struct{Q}$. (हाच पुढे-मागे युक्तिवाद वापरून
  $\Struct{R} \iso[p] \Struct{Q}$ प्रस्थापित करून आपण हे थेटही दाखवू
  शकलो असतो.)
\end{rem}


\chapter{अंकगणिताची प्रतिमाने}\label{mod:mar::chap}








\section{परिचय}\label{mod:mar:int:sec}

अंकगणिताचे \emph{मानक प्रतिमान} म्हणजे अशी संरचना~$\Struct{N}$ की
$\Domain{N} = \Nat$ आणि तिच्यात $\Obj{0}$, $\prime$, $+$, $\times$ व
$<$ यांचे अपेक्षित नेहमीच्या अर्थाने अर्थनिर्धारण केलेले असते. म्हणजे,
$\Obj{0}$ हे $0$ असते, $\prime$ हे उत्तरवर्ती फलन असते, $+$ चा अर्थ
बेरीज आणि $\times$ चा अर्थ~$\Nat$ मधील संख्यांचा गुणाकार असा असतो.
अधिक स्पष्टपणे,
\begin{align*}
  \Assign{\Obj{0}}{N} & = 0\\
  \Assign{\prime}{N}(n) & = n + 1\\
  \Assign{+}{N}(n, m) & = n + m\\
  \Assign{\times}{N}(n, m) & = nm
\end{align*}
अर्थात, $\Lang{L_A}$ साठी अशाही रचना असतात की त्यांचे प्रांत~$\Nat$
नसतात. उदाहरणार्थ, $\Domain{M} = \{a\}^*$ (एकाच $a$ या चिन्हाच्या
सांत क्रमिका, म्हणजे $\emptyset$, $a$, $aa$, $aaa$, \dots) असा प्रांत
असलेली $\Struct{M}$ घेता येते आणि पुढील अर्थनिर्धारण करता येते:
\begin{align*}
  \Assign{\Obj{0}}{M} & = \emptyset\\
  \Assign{\prime}{M}(s) & = s \concat a\\
  \Assign{+}{M}(n, m) & = a^{n + m}\\
  \Assign{\times}{M}(n, m) & = a^{nm}
\end{align*}
या दोन रचना ``मूलतः एकच'' आहेत: फरक फक्त प्रांताचे
घटक कोणते आहेत यात आहे; अर्थनिर्धारण फलनांमुळे प्रांताचे
घटक परस्पर कसे संबंधित आहेत यात नाही. अशा दोन संरचना ना
\emph{समरूप} म्हणतात.

संहतता प्रमेयावरून सहज निष्पन्न होते की~$\Struct{N}$ मध्ये सत्य असलेल्या
कोणत्याही उपपत्तीला~$\Struct{N}$ शी समरूप नसलेली प्रतिमानेही असतात. अशा
रचनांना \emph{अमानक} म्हणतात. त्यांच्याबद्दलची रंजक गोष्ट अशी की मानक
प्रतिमानातील घटक (म्हणजे $\Struct{N}$ मधील, तसेच तिच्याशी समरूप
असलेल्या सर्व संरचना मधील घटक) मानक संख्यांक~$\num{n}$ यांच्या
मूल्यांनी पूर्णपणे व्यापलेले असतात, म्हणजे
\[
\Domain{N} = \Setabs{\Value{\num{n}}{N}}{n \in \Nat}
\]
पण अमानक प्रतिमानांत तसे नसते: $\Struct{M}$ अमानक असेल, तर किमान एक
$x \in \Domain{M}$ असा असतो की प्रत्येक~$n$ साठी
$x \neq \Value{\num{n}}{M}$.

हे अमानक घटक फार रंजक आहेत: ते ``अनंत नैसर्गिक संख्या'' आहेत. पण त्यांच्या
अस्तित्वामुळे अपूर्णता-घटनांचे एका अर्थी स्पष्टीकरणही मिळते. उदाहरणार्थ,
पेआनो अंकगणितासाठीचे सुसंगतता-विधान $\OCon[\Th{PA}]$, म्हणजे $\lnot
\lexists[x][\OPrf[\Th{PA}](x, \gn{\lfalse})]$, विचारात घ्या.
$\Th{PA}$ कडून $\OCon[\Th{PA}]$ किंवा $\lnot \OCon[\Th{PA}]$ यांपैकी
काहीही सिद्ध होत नसल्यामुळे यांपैकी कोणतेही विधान~$\Th{PA}$ मध्ये सुसंगतपणे
जोडता येते. $\Th{PA}$ सुसंगत असल्यामुळे
$\Sat{N}{\OCon[\Th{PA}]}$, आणि म्हणून $\Sat/{N}{\lnot
  \OCon[\Th{PA}]}$. त्यामुळे $\Struct{N}$ हे $\Th{PA}
\cup \{\lnot \OCon[\Th{PA}]\}$ चे \emph{प्रतिमान नाही}, आणि त्या उपपत्तीची
सर्व प्रतिमाने अमानकच असली पाहिजेत. $\Th{PA} \cup \{\lnot
\OCon[\Th{PA}]\}$ च्या प्रतिमानात असा एखादा घटक असलाच पाहिजे की
तो $\lexists[x][\OPrf[\Th{PA}](x, \gn{\lfalse})]$ सत्य करणारा साक्षी
ठरेल, म्हणजे~$\Th{PA}$ पासून व्याघाताच्या निष्पत्तीची गोडेल संख्या.

असा घटक मानक असू शकत नाही---कारण प्रत्येक~$n$ साठी
$\Th{PA} \Proves \lnot \OPrf[\Th{PA}](\num{n}, \gn{\lfalse})$.











\section{अंकगणिताची मानक प्रतिमाने}\label{mod:mar:stm:sec}

अंकगणिताची भाषा~$\Lang{L_A}$ ही अर्थातच संख्यांविषयी, विशेषतः नैसर्गिक
संख्यांविषयी, अभिप्रेत आहे. म्हणून ``ते'' मानक प्रतिमान~$\Struct{N}$ विशेष
आहे: आपल्याला ज्याविषयी बोलायचे आहे तेच हे प्रतिमान. पण तर्कशास्त्रात
आपल्याला अनेकदा फक्त रचनात्मक गुणधर्मांत रस असतो, आणि समरूप असलेल्या कोणत्याही
दोन संरचनांमध्ये ते समान असतात. त्यामुळे व्याप्ती थोडी वाढवून
$\Struct{N}$ शी समरूप असलेली कोणतीही संरचना ``मानक'' मानता येते.

\begin{defn}
$\Lang{L_A}$ साठीची एखादी संरचना जर~$\Struct{N}$ शी समरूप असेल,
तर ती \emph{मानक} आहे.
\end{defn}

\begin{prop}
\label{mod:mar:stm:prop:standard-domain} एखादी संरचना~$\Struct{M}$ मानक असेल,
तर तिचा प्रांत म्हणजे मानक संख्यांकांच्या मूल्यांचा संच असतो, म्हणजे
\[
\Domain{M} = \Setabs{\Value{\num{n}}{M}}{n \in \Nat}
\]
\end{prop}

\begin{proof}
प्रत्येक $\Value{\num{n}}{M} \in \Domain{M}$ हे स्पष्ट आहे. प्रत्येक
$x \in \Domain{M}$ हा कोणत्यातरी~$n$ साठी $\Value{\num{n}}{M}$ बरोबर आहे,
एवढेच दाखवायचे आहे. $\Struct{M}$ मानक असल्यामुळे ते~$\Struct{N}$ शी समरूप
आहे. $g\colon \Nat \to \Domain{M}$ हे समरूपण आहे असे समजा. मग
$g(n) = g(\Value{\num{n}}{N}) = \Value{\num{n}}{M}$. पण प्रत्येक
$x \in \Domain{M}$ साठी $g(n) = x$ असे एक~$n \in \Nat$ असते, कारण
$g$ हे आच्छादक आहे.
\end{proof}

\begin{explain}
$\Lang{L_A}$ साठीची एखादी रचना~$\Struct{M}$ मानक असेल, तर तिच्या
प्रांत मधील सर्व घटकांना $\num{0}$, $\num{1}$, $\num{2}$, \dots हे
मानक संख्यांक, म्हणजेच $\Obj{0}$, $\Obj{0}'$, $\Obj{0}''$ इत्यादी पदे,
नावे देऊ शकतात. अर्थात, याचा अर्थ $\Domain{M}$ मधील घटक
\emph{संख्याच आहेत} असा नाही; $\Domain{N}$ मधील संख्या जशा निवडून दाखवता
येतात, त्याच रीतीने हे घटक निवडून दाखवता येतात एवढाच त्याचा अर्थ आहे.
\end{explain}

\begin{prob}
\ref{mod:mar:stm:prop:standard-domain} चा व्यत्यास असत्य आहे हे दाखवा;
म्हणजे, $\Domain{M} = \Setabs{\Value{\num{n}}{M}}{n \in \Nat}$ असूनही
$\Struct{N}$ शी समरूप नसलेल्या एखाद्या संरचना~$\Struct{M}$ चे
उदाहरण द्या.
\end{prob}

\begin{prop}
\label{mod:mar:stm:prop:thq-standard}
जर $\Sat{M}{\Th{Q}}$ आणि $\Domain{M} =
\Setabs{\Value{\num{n}}{M}}{n \in \Nat}$, तर $\Struct{M}$ मानक आहे.
\end{prop}

\begin{proof}
$\Struct{M}$ हे~$\Struct{N}$ शी समरूप आहे हे दाखवायचे आहे.
$g(n) = \Value{\num{n}}{M}$ अशी व्याख्या असलेले
$g\colon \Nat \to \Domain{M}$ फलन विचारात घ्या. गृहीतकानुसार $g$ हे
आच्छादक आहे. ते एकास-एक सुद्धा आहे: $n \neq m$ असेल तेव्हा
$\Th{Q} \Proves \eq/[\num{n}][\num{m}]$. म्हणून
$\Sat{M}{\Th{Q}}$ असल्यामुळे, $n \neq m$ असेल तेव्हा
$\Sat{M}{\eq/[\num{n}][\num{m}]}$. त्यामुळे $n \neq m$ असल्यास
$\Value{\num{n}}{M} \neq \Value{\num{m}}{M}$, म्हणजे $g(n) \neq g(m)$.

$g$ हे समरूपण आहे हेही पडताळायचे आहे.
\begin{enumerate}
\item $\Assign{\Obj{0}}{N} = 0$ असल्यामुळे
  $g(\Assign{\Obj{0}}{N}) = g(0)$. $g$ च्या व्याख्येवरून
  $g(0) = \Value{\num{0}}{M}$. पण $\num{0}$ म्हणजे फक्त $\Obj{0}$;
  आणि एखादे पद स्थिरांक असेल, तर त्याचे मूल्य त्या संरचना ने
  त्या स्थिरांक ला नेमून दिलेले मूल्य असते, म्हणजे
  $\Value{\Obj{0}}{M} = \Assign{\Obj{0}}{M}$. म्हणून अपेक्षेप्रमाणे
  $g(\Assign{\Obj{0}}{N}) = \Assign{\Obj{0}}{M}$.
\item $\Struct{N}$ मध्ये $\prime$ हे~$\Nat$ वरील उत्तरवर्ती फलन असल्यामुळे
  $g(\Assign{\prime}{N}(n)) = g(n+1)$. पुढे, $g$ च्या व्याख्येवरून
  $g(n+1) = \Value{\num{n+1}}{M}$. पण $\num{n+1}$ आणि $\num{n}'$ हे
  एकच पद आहे, म्हणून $\Value{\num{n+1}}{M} = \Value{\num{n}'}{M}$.
  मूल्य-फलनाच्या व्याख्येवरून हे
  $= \Assign{\prime}{M}(\Value{\num{n}}{M})$. तसेच
  $\Value{\num{n}}{M} = g(n)$ असल्यामुळे
  $g(\Assign{\prime}{N}(n)) = \Assign{\prime}{M}(g(n))$ मिळते.
\item $\Struct{N}$ मध्ये $+$ हे~$\Nat$ वरील बेरीज-फलन असल्यामुळे
  $g(\Assign{+}{N}(n,m)) = g(n+m)$. पुढे, $g$ च्या व्याख्येवरून
  $g(n+m) = \Value{\num{n+m}}{M}$. पण $\Th{Q} \Proves
  \num{n+m} = (\num{n} + \num{m})$, म्हणून $\Value{\num{n+m}}{M} =
  \Value{\num{n}+\num{m}}{M}$. मूल्य-फलनाच्या व्याख्येवरून हे $=
  \Assign{+}{M}(\Value{\num{n}}{M},\Value{\num{m}}{M})$. तसेच
  $\Value{\num{n}}{M} = g(n)$ आणि $\Value{\num{m}}{M} = g(m)$ असल्यामुळे
  $g(\Assign{+}{N}(n, m)) = \Assign{+}{M}(g(n), g(m))$ मिळते.
\item $g(\Assign{\times}{N}(n, m)) = \Assign{\times}{M}(g(n), g(m))$:
  सराव.
\item $\tuple{n,m} \in \Assign{<}{N}$ हे तेव्हाच, जेव्हा $n < m$. जर
  $n < m$, तर $\Th{Q} \Proves \num{n} < \num{m}$ आणि
  $\Sat{M}{\num{n} < \num{m}}$ सुद्धा. त्यामुळे
  $\tuple{\Value{\num{n}}{M}, \Value{\num{m}}{M}} \in \Assign{<}{M}$,
  म्हणजे $\tuple{g(n), g(m)} \in \Assign{<}{M}$. जर $n \not< m$, तर
  $\Th{Q} \Proves \lnot \num{n} < \num{m}$ आणि परिणामी
  $\Sat/{M}{\num{n} < \num{m}}$. त्यामुळे पूर्वीप्रमाणे
  $\tuple{g(n), g(m)} \notin \Assign{<}{M}$. दोन्ही एकत्र केल्यास:
  $\tuple{n,m} \in \Assign{<}{N}$ हे तेव्हाच, जेव्हा
  $\tuple{g(n), g(m)} \in \Assign{<}{M}$.
\end{enumerate}
\end{proof}

\begin{explain}
$\Lang{L_A}$ साठीच्या इतर कोणत्याही संरचना~$\Struct{M}$ च्या
प्रांतात~$\Nat$ कडून संगती ठरवण्याचा सर्वांत साहजिक मार्ग म्हणजे फलन~$g$;
कारण अशा प्रत्येक $\Struct{M}$ मध्ये $\num{0}$, $\num{1}$, $\num{2}$
इत्यादींनी नाव दिलेले घटक असतात. त्यामुळे, $\Struct{M}$ ने
$\num{n}$ विषयीची किमान काही मूलभूत विधाने~$\Struct{N}$ प्रमाणेच सत्य
केली आणि $g$ हे एकास-एक व आच्छादकही असेल, तर $g$ हे समरूपण ठरेल यात नवल
नाही. किंबहुना, $\Domain{M}$ मध्ये $\num{n}$ यांनी नाव दिलेल्यांव्यतिरिक्त
दुसरे घटक नसतील, तर असे समरूपण एकमेव असते.
\end{explain}

\begin{prop}
\label{mod:mar:stm:prop:thq-unique-iso} जर $\Struct{M}$ मानक असेल, तर
\ref{mod:mar:stm:prop:thq-standard} च्या पुराव्यातील $g$ हे~$\Struct{N}$ कडून
$\Struct{M}$ कडे जाणारे एकमेव समरूपण आहे.
\end{prop}

\begin{proof}
$h\colon \Nat \to \Domain{M}$ हे $\Struct{N}$ आणि~$\Struct{M}$ यांमधील
समरूपण आहे असे समजा. $n$ वरील विगमनाने $g = h$ हे दाखवू. $n = 0$ असेल,
तर $g$ च्या व्याख्येवरून $g(0) = \Assign{\Obj{0}}{M}$. पण $h$ हे समरूपण
असल्यामुळे $h(0) = h(\Assign{\Obj{0}}{N}) =\Assign{\Obj{0}}{M}$;
म्हणून $g(0) = h(0)$.

आता $n+1$ चे प्रकरण विचारात घ्या. आपल्याला
\begin{align*}
  g(n+1) & = \Value{\num{n+1}}{M} \text{ हे~$g$ च्या व्याख्येवरून}\\
  & = \Value{\num{n}'}{M} \text{ कारण $\num{n+1}\ident \num{n}'$}\\
  & = \Assign{\prime}{M}(\Value{\num{n}}{M})
    \text{ हे $\Value{t'}{M}$ च्या व्याख्येवरून}\\
  & = \Assign{\prime}{M}(g(n))  \text{ हे~$g$ च्या व्याख्येवरून}\\
  & = \Assign{\prime}{M}(h(n)) \text{ हे विगमन गृहीतकावरून}\\
  & = h(\Assign{\prime}{N}(n)) \text{ कारण $h$ हे समरूपण आहे}\\
  & = h(n+1)
\end{align*}
मिळते.
\end{proof}

\begin{explain}
कोणत्याही गणनीय अनंत संच~$M$ साठी $\Nat$ आणि $M$ यांमध्ये
एकास-एक आच्छादन असते; म्हणून असा प्रत्येक संच~$M$ हा एखाद्या मानक
प्रतिमान~$\Struct{M}$ चा संभाव्य प्रांत आहे. किंबहुना, एकदा $z \in M$
ही वस्तू आणि योग्य फलन $s$ ही अनुक्रमे $\Assign{\Obj{0}}{M}$ आणि
$\Assign{\prime}{M}$ म्हणून निवडली, की $+$, $\times$ आणि $<$ यांचे
अर्थनिर्धारण आधीच निश्चित होते. $s\colon M \to M \setminus \{z\}$ अशी
एकास-एक आणि आच्छादक दोन्ही असलेली फलनेच मानक प्रतिमानात
$\Assign{\prime}{M}$ म्हणून योग्य असतात. $s$ च्या व्याप्तीत~$z$ असू शकत
नाही; अन्यथा $\lforall[x][\eq/[\Obj 0][x']]$ असत्य होईल. हे वाक्य
$\Struct{N}$ मध्ये सत्य आहे, म्हणून $\Struct{M}$ नेही ते सत्य केले पाहिजे.
फलन~$s$ हे एकास-एक असले पाहिजे, कारण $\Struct{N}$ मधील उत्तरवर्ती
फलन~$\Assign{\prime}{N}$ तसे आहे; आणि $\Assign{\prime}{N}$ हे
एकास-एक आहे हे~$\Struct{N}$ मध्ये सत्य असलेल्या वाक्य ने
व्यक्त होते. ते आच्छादक सुद्धा असले पाहिजे; अन्यथा असा एखादा
$x \in M \setminus \{z\}$ असेल की तो~$s$ च्या व्याप्तीत नसेल, म्हणजे
वाक्य $\lforall[x][(\eq[x][\Obj 0] \lor
\lexists[y][\eq[y'][x]])]$ हे~$\Struct{M}$ मध्ये असत्य होईल---पण ते
$\Struct{N}$ मध्ये सत्य आहे.

\end{explain}












\section{अमानक प्रतिमाने}\label{mod:mar:nst:sec}

\begin{explain}
$\Lang{L_A}$ साठीची एखादी संरचना~$\Struct{N}$ शी समरूप असेल,
तर तिला मानक म्हणतो. एखादी संरचना~$\Struct{N}$ शी समरूप नसेल,
तर तिला अमानक म्हणतो.
\end{explain}

\begin{defn}
$\Lang{L_A}$ साठीची संरचना~$\Struct{M}$ जर~$\Struct{N}$ शी समरूप
नसेल, तर ती \emph{अमानक} आहे. प्रांतातील जे घटक
$x \in \Domain{M}$ कोणत्यातरी $n \in \Nat$ साठी
$\Value{\num{n}}{M}$ बरोबर आहेत, त्यांना ($\Struct{M}$ मधील)
\emph{मानक संख्या} म्हणतात; आणि जे तसे नाहीत त्यांना \emph{अमानक संख्या}
म्हणतात.
\end{defn}

\begin{explain}
\ref{mod:mar:stm:prop:standard-domain} नुसार, $\Lang{L_A}$ साठीच्या कोणत्याही
मानक संरचनांमध्ये फक्त मानक घटक असतात. त्यामुळे अमानक
संरचनांमध्ये किमान एक अमानक घटक असलाच पाहिजे. किंबहुना, अमानक
घटक अस्तित्वात असणे हे ती संरचना अमानक असल्याची हमी देते.
\end{explain}

\begin{prop}
$\Lang{L_A}$ साठीच्या संरचना~$\Struct{M}$ मध्ये अमानक संख्या
असेल, तर $\Struct{M}$ अमानक आहे.
\end{prop}

\begin{proof}
तसे नाही असे समजा; म्हणजे $\Struct{M}$ मानक असूनही तिच्यात अमानक संख्या
$x$ आहे असे समजा. $g\colon \Nat \to \Domain{M}$ हे समरूपण असू द्या.
$n$ वरील विगमनाने $g(\Value{\num{n}}{N}) = \Value{\num{n}}{M}$ हे सहज
दिसते. दुसऱ्या शब्दांत, $g$ हे~$\Struct{N}$ मधील मानक संख्यांना
$\Struct{M}$ मधील मानक संख्यांवर पाठवते. $\Struct{M}$ मध्ये अमानक संख्या
असेल, तर गृहीतकाच्या विरुद्ध $g$ हे आच्छादक असू शकत नाही.
\end{proof}

\begin{prob}
$\Th{Q}$ मध्ये पुढील स्वयंसिद्धके आहेत हे आठवा:
\begin{align*}
& \lforall[x][\lforall[y][(\eq[x'][y'] \lif \eq[x][y])]] \tag{$!Q_1$}\\
& \lforall[x][\eq/[\Obj 0][x']] \tag{$!Q_2$}\\
& \lforall[x][(\eq[x][\Obj 0] \lor \lexists[y][\eq[x][y']])] \tag{$!Q_3$}
\end{align*}
पुढील अटी पूर्ण करणाऱ्या संरचना~$\Struct{M_1}$, $\Struct{M_2}$,
$\Struct{M_3}$ द्या:
\begin{enumerate}
\item $\Sat{M_1}{!Q_1}$, $\Sat{M_1}{!Q_2}$, $\Sat/{M_1}{!Q_3}$;
\item $\Sat{M_2}{!Q_1}$, $\Sat/{M_2}{!Q_2}$, $\Sat{M_2}{!Q_3}$; आणि
\item $\Sat/{M_3}{!Q_1}$, $\Sat{M_3}{!Q_2}$, $\Sat{M_3}{!Q_3}$.
\end{enumerate}
अर्थात, प्रत्येकासाठी फक्त~$\Assign{\Obj{0}}{M_i}$ आणि
$\Assign{\prime}{M_i}$ निश्चित करायची आहेत.
\end{prob}

\begin{explain}
$\Lang{L_A}$ साठी अमानक संरचना निश्चित करणे पुरेसे सोपे आहे.
उदाहरणार्थ, प्रांत~$\Int$ असलेली रचना घ्या आणि सर्व अतार्किक चिन्हांचे
नेहमीप्रमाणे अर्थनिर्धारण करा. ऋण संख्या कोणत्याही~$n$ साठी $\num{n}$ ची
मूल्ये नसल्यामुळे ही रचना अमानक आहे. अर्थात, $\Struct{N}$ प्रमाणेच ती
वाक्ये सत्य करते या अर्थाने ती अंकगणिताचे \emph{प्रतिमान} नसेल. उदाहरणार्थ,
$\lforall[x][\eq/[x'][\Obj{0}]]$ हे असत्य आहे. तरीसुद्धा, संहतता प्रमेय
वापरून अंकगणिताची अमानक प्रतिमाने अस्तित्वात आहेत हे सहज सिद्ध करता येते.
\end{explain}

\begin{prop}
$\Th{TA} = \Setabs{\varphi}{\Sat{N}{\varphi}}$ ही~$\Struct{N}$ ची उपपत्ती असू द्या.
$\Th{TA}$ ला गणनीय अमानक प्रतिमान आहे.
\end{prop}

\begin{proof}
$\Lang{L_A}$ मध्ये नवीन स्थिरांक~$c$ घालून तिचा विस्तार करा आणि
पुढील वाक्यांचा संच विचारात घ्या:
\[
\Gamma = \Th{TA} \cup \{\eq/[c][\num{0}], \eq/[c][\num{1}],
\eq/[c][\num{2}], \dots\}
\]
$\Gamma$ च्या कोणत्याही प्रतिमान~$\Struct{M^c}$ मध्ये
$x = \Assign{c}{M^c}$ असा अमानक घटक असेल, कारण प्रत्येक
$n \in \Nat$ साठी $x \neq \Value{\num{n}}{M}$. तसेच
$\Th{TA} \subseteq \Gamma$ असल्यामुळे अर्थातच $\Sat{M^c}{\Th{TA}}$.
फक्त~$c$ विसरून $\Struct{M^c}$ पासून $\Lang{L_A}$ साठीची
संरचना~$\Struct{M}$ बनवली, तरी तिच्या प्रांतात तो अमानक~$x$ राहतो
आणि $\Sat{M}{\Th{TA}}$ सुद्धा. $\Th{TA}$ मध्ये $c$ येत नसल्यामुळे या
दुसऱ्या दाव्याची हमी मिळते. म्हणून $\Gamma$ ला प्रतिमान आहे एवढे दाखवणे
पुरेसे आहे.


$\Gamma$ ला प्रतिमान आहे हे दाखवण्यासाठी संहतता प्रमेय वापरू. $\Gamma$ चा
प्रत्येक सांत उपसंच पूर्ततायोग्य असेल, तर $\Gamma$ सुद्धा पूर्ततायोग्य आहे.
कोणताही सांत उपसंच $\Gamma_0 \subseteq \Gamma$ विचारात घ्या.
$\Gamma_0$ मध्ये~$\Th{TA}$ मधील काही वाक्ये आणि
$\eq/[c][\num{n}]$ या रूपातील सांत वाक्ये येतात. या दुसऱ्या रूपातील एकही
वाक्य नसेल, तर c ला शून्य मूल्य देणारा N चा विस्तार त्या उपसंचाला थेट
पूर्ण करतो; म्हणून पुढे किमान एक असे वाक्य आहे असे समजा.

$\eq/[c][\num{k}] \in \Gamma_0$ असा $k$ हा सर्वांत मोठा अंक समजा.
$\Assign{c}{N_k} = k+1$ हे अर्थनिर्धारण समाविष्ट करण्यासाठी~$\Struct{N}$
चा विस्तार करून $\Struct{N_k}$ ची व्याख्या करा. $\Sat{N_k}{\Gamma_0}$:
जर $\varphi \in \Th{TA}$, तर $\Struct{N_k}$ ही~$c$ वगळता प्रत्येक बाबतीत
$\Struct{N}$ सारखीच असल्यामुळे आणि~$\varphi$ मध्ये $c$ येत नसल्यामुळे
$\Sat{N_k}{\varphi}$. तसेच $n \le k$ आणि $\Value{c}{N_k} = k+1$ असल्यामुळे
$\Sat{N_k}{\eq/[c][\num{n}]}$. म्हणून $\Gamma$ चा प्रत्येक सांत उपसंच
पूर्ततायोग्य आहे आणि संहततेनुसार संपूर्ण संचाला प्रतिमान आहे. अधोगामी
लोव्हेनहाइम--स्कोलेम प्रमेयानुसार हे प्रतिमान गणनीय घेता येते.

\end{proof}











\section{$\Th{Q}$ ची प्रतिमाने}\label{mod:mar:mdq:sec}

\begin{explain}
$\Struct{N}$ जी वाक्ये सत्य करते तीच सत्य करणाऱ्या अमानक
संरचना आहेत, म्हणजे त्या~$\Th{TA}$ ची प्रतिमाने आहेत, हे आपल्याला
माहीत आहे. $\Sat{N}{\Th{Q}}$ असल्यामुळे $\Th{TA}$ चे कोणतेही प्रतिमान
$\Th{Q}$ चेही प्रतिमान असते. $\Th{Q}$ ही~$\Th{TA}$ पेक्षा फारच दुर्बल
आहे; उदाहरणार्थ, $\Th{Q} \Proves/ \lforall[x][\lforall[y][\eq[(x +
      y)][(y+x)]]]$. दुर्बलतर उपपत्ती पूर्ण करणे सोपे असते: त्यांना अधिक
प्रतिमाने असतात. उदाहरणार्थ, $\Th{Q}$ च्या काही प्रतिमानांत
$\lforall[x][\lforall[y][\eq[(x + y)][(y+x)]]]$ असत्य असते; पण ती
$\Th{TA}$ ची किंवा त्या दृष्टीने $\Th{PA}$ चीसुद्धा प्रतिमाने असू शकत
नाहीत. $\Th{Q}$ ची प्रतिमाने तुलनेने सोपीही असतात: ती स्पष्टपणे निश्चित
करता येतात.
\end{explain}

\begin{ex}
\label{mod:mar:mdq:ex:model-K-of-Q}
$\Domain{K} = \Nat \cup \{a\}$ असा प्रांत आणि पुढील अर्थनिर्धारणे असलेली
संरचना~$\Struct{K}$ विचारात घ्या:
\begin{align*}
  \Assign{\Obj{0}}{K} & = 0\\
  \Assign{\prime}{K}(x) & =
  \begin{cases}
    x+1 & \text{जर $x\in \Nat$}\\
    a & \text{जर $x = a$}
  \end{cases}\\
  \Assign{+}{K}(x, y) & =
  \begin{cases}
    x+y & \text{जर $x$, $y \in\Nat$}\\
    a & \text{इतर प्रसंगी}
  \end{cases}\\
  \Assign{\times}{K}(x, y) & =
  \begin{cases}
    xy & \text{जर $x$, $y \in\Nat$}\\
    0 & \text{जर $x = 0$ किंवा $y = 0$}\\
    a & \text{इतर प्रसंगी}\\
  \end{cases}\\
  \Assign{<}{K} & =
  \Setabs{\tuple{x,y}}{x, y \in \Nat \text{ आणि } x<y} \cup
  \Setabs{\tuple{x,a}}{x \in \Domain{K}}
\end{align*}
$\Sat{K}{\Th{Q}}$ हे दाखवण्यासाठी $\Th{Q}$ ची सर्व स्वयंसिद्धके
$\Struct{K}$ मध्ये सत्य आहेत हे पडताळावे लागेल. सोयीसाठी,
$\Assign{\prime}{K}(x)$ करिता $x^\nssucc$ ($\Struct{K}$ मधील $x$ चा
``उत्तरवर्ती''), $\Assign{+}{K}(x, y)$ करिता $x \nsplus y$
($\Struct{K}$ मधील $x$ आणि $y$ यांची ``बेरीज''),
$\Assign{\times}{K}(x, y)$ करिता $x \nstimes y$ ($\Struct{K}$ मधील
$x$ आणि~$y$ यांचा ``गुणाकार''), आणि $\tuple{x,y} \in \Assign{<}{K}$
करिता $x \nsless y$ लिहू. या संक्षेपांनी $\Struct{K}$ मधील क्रिया अधिक
स्पष्टपणे अशा देता येतात:
\[
\begin{array}{c|c}
  x & x^\nssucc \\
  \hline
  n & n+1 \\
  a & a
\end{array}
\qquad
\begin{array}{c|ccc}
  x \nsplus y & 0 & m & a \\
  \hline
  0 & 0 & m & a \\
  n & n & n+m & a \\
  a & a & a & a \\
\end{array}
\qquad
\begin{array}{c|ccc}
  x \nstimes y & 0 & m & a \\
  \hline
  0 & 0 & 0 & 0 \\
  n & 0 & nm & a \\
  a & 0 & a & a \\
\end{array}
\]
$n$, $m \in \Nat$ यांच्यासाठी $n \nsless m$ हे तेव्हाच, जेव्हा $n<m$;
आणि प्रत्येक~$x \in \Domain{K}$ साठी $x \nsless a$.

$\nssucc$ हे एकास-एक असल्यामुळे
$\Sat{K}{\lforall[x][\lforall[y][(\eq[x'][y'] \lif \eq[x][y])]]}$.
$\Struct{K}$ मध्ये $0$ हा $\nssucc$-उत्तरवर्ती नसल्यामुळे
$\Sat{K}{\lforall[x][\eq/[\Obj 0][x']]}$. प्रत्येक $n>0$ साठी
$n = (n-1)^\nssucc$ आणि $a = a^\nssucc$ असल्यामुळे
$\Sat{K}{\lforall[x][(\eq[x][\Obj 0] \lor \lexists[y][\eq[x][y']])]}$.

$n \nsplus 0 = n+0 = n$ आणि~$\nsplus$ च्या व्याख्येवरून
$a\nsplus 0 = a$ असल्यामुळे $\Sat{K}{\lforall[x][\eq[(x + \Obj 0)][x]]}$.
$\Sat{K}{\lforall[x][\lforall[y][\eq[(x + y')][(x + y)']]]}$ हे थोडे
अधिक गुंतागुंतीचे आहे. $n$, $m$ या दोन्ही मानक असतील, तर:
\begin{align*}
(n \nsplus m^\nssucc) & = (n+(m+1)) = (n+m)+1 = (n \nsplus m)^\nssucc
\intertext{कारण मानक संख्यांवर $\nsplus$ आणि $^\nssucc$ ही अनुक्रमे $+$
  आणि $\prime$ यांच्याशी जुळतात. आता $x \in \Domain{K}$ असे समजा. मग}
(x \nsplus a^\nssucc) & = (x \nsplus a) = a = a^\nssucc = (x \nsplus a)^\nssucc
\intertext{उरलेले प्रकरण $y \in \Domain{K}$ पण $x =
  a$ असण्याचे आहे. येथे $y = n$ मानक आहे की $y = a$, यानुसारही प्रकरणे
  वेगळी करावी लागतात:}
(a \nsplus n^\nssucc) & = (a \nsplus (n+1)) = a = a^\nssucc = (a \nsplus n)^\nssucc\\
(a \nsplus a^\nssucc) & = (a \nsplus a) = a = a^\nssucc = (a \nsplus a)^\nssucc
\end{align*}

अर्थात, हा तपशील आवश्यकतेपेक्षा थोडा अधिक आहे. उदाहरणार्थ,
$z$ काहीही असले तरी $a \nsplus z = a$ असल्यामुळे
$a \nsplus a^\nssucc = a$ हे लगेच निष्पन्न होते. उरलेली स्वयंसिद्धकेही
याच रीतीने पडताळता येतात.

त्यामुळे $\Struct{K}$ हे~$\Th{Q}$ चे प्रतिमान आहे. तिची ``बेरीज''~$\nsplus$
ही क्रमनिरपेक्षसुद्धा आहे. पण $\Struct{N}$ मध्ये सत्य आणि $\Struct{K}$
मध्ये असत्य अशी इतर वाक्ये आहेत, तसेच याच्या उलटही आहे. उदाहरणार्थ,
$a \nsless a$, म्हणून $\Sat{K}{\lexists[x][x < x]}$ आणि
$\Sat/{K}{\lforall[x][\lnot x<x]}$. यावरून $\Th{Q} \Proves/
\lforall[x][\lnot x < x]$ हे दिसते.
\end{ex}

\begin{prob}
\ref{mod:mar:mdq:ex:model-K-of-Q} मधील $\Struct{K}$ ही~$\Th{Q}$ ची
उरलेली स्वयंसिद्धके पूर्ण करते हे सिद्ध करा:
\begin{align*}
  & \lforall[x][\eq[(x \times \Obj 0)][\Obj 0]] \tag{$!Q_6$}\\
  & \lforall[x][\lforall[y][\eq[(x \times y')][((x \times y) + x)]]] \tag{$!Q_7$}\\
  & \lforall[x][\lforall[y][(x < y \liff \lexists[z][\eq[(z' + x)][y]])]] \tag{$!Q_8$}
\end{align*}
फक्त~$\prime$ असलेले, $\Struct{N}$ मध्ये सत्य पण $\Struct{K}$ मध्ये असत्य
असे वाक्य शोधा.
\end{prob}

\begin{ex}
\label{mod:mar:mdq:ex:model-L-of-Q} $\Domain{L} = \Nat \cup \{a, b\}$ असा प्रांत
आणि पुढीलप्रमाणे दिलेली $\Assign{\prime}{L} = \nssucc$ व
$\Assign{+}{L} = \nsplus$ ही अर्थनिर्धारणे असलेली संरचना~$\Struct{L}$
विचारात घ्या:
\[
\begin{array}{c|c}
  x & x^\nssucc \\
  \hline
  n & n+1 \\
  a & a\\
  b & b
\end{array}
\qquad
\begin{array}{c|ccc}
  x \nsplus y & m & a & b\\
  \hline
  n & n+m & b & a\\
  a & a & b & a\\
  b & b & b & a
\end{array}
\]
$\nssucc$ हे एकास-एक आहे, $0$ त्याच्या व्याप्तीत नाही, आणि
$0$ व्यतिरिक्त प्रत्येक~$x \in \Domain{L}$ त्याच्या व्याप्तीत आहे;
म्हणून $!Q_1$--$!Q_3$ ही स्वयंसिद्धके~$\Struct{L}$ मध्ये सत्य आहेत.
प्रत्येक $x$ साठी $x \nsplus 0 = x$, म्हणून $!Q_4$ सुद्धा सत्य आहे.
$!Q_5$ साठी $x \nsplus y^\nssucc$ आणि $(x \nsplus y)^\nssucc$ विचारात
घ्या. $x$ आणि $y$ दोन्ही मानक असतील, तर ते समान आहेत; कारण तेव्हा
$\nssucc$ आणि $\nsplus$ ही $\prime$ आणि $+$ यांच्याशी जुळतात. $x$ अमानक
आणि $y$ मानक असेल, तर $x \nsplus y^\nssucc = x = x^\nssucc =
(x \nsplus y)^\nssucc$. $x$ आणि $y$ दोन्ही अमानक असतील, तर चार प्रकरणे
आहेत:
\begin{align*}
&  a \nsplus a^\nssucc  = b = b^\nssucc = (a \nsplus a)^\nssucc\\
&  b \nsplus b^\nssucc  = a = a^\nssucc = (b \nsplus b)^\nssucc\\
&  b \nsplus a^\nssucc  = b = b^\nssucc = (b \nsplus a)^\nssucc\\
&  a \nsplus b^\nssucc  = a = a^\nssucc = (a \nsplus b)^\nssucc\\
\intertext{जर $x$ मानक पण $y$ अमानक असेल, तर}
&  n \nsplus a^\nssucc  = n \nsplus a = b = b^\nssucc = (n \nsplus a)^\nssucc\\
&  n \nsplus b^\nssucc  = n \nsplus b = a = a^\nssucc = (n \nsplus b)^\nssucc
\end{align*}

त्यामुळे $\Sat{L}{!Q_5}$. पण $a \nsplus 0 \neq 0 \nsplus a$, म्हणून
$\Sat/{L}{\lforall[x][\lforall[y][\eq[(x+y)][(y+x)]]]}$.
\end{ex}

\begin{prob}
\ref{mod:mar:mdq:ex:model-L-of-Q} मधील $\Struct{L}$ मध्ये
$\times$ आणि $<$ यांचे अर्थनिर्धारण करणारी $\nstimes$ आणि $\nsless$
समाविष्ट करून तिचा विस्तार करा. तुमची रचना~$\Th{Q}$ ची उरलेली
स्वयंसिद्धके पूर्ण करते हे दाखवा:
\begin{align*}
& \lforall[x][\eq[(x \times \Obj 0)][\Obj 0]] \tag{$!Q_6$}\\ &
  \lforall[x][\lforall[y][\eq[(x \times y')][((x \times y) + x)]]]
  \tag{$!Q_7$}\\ & \lforall[x][\lforall[y][(x < y \liff \lexists[z][\eq[(z'+x)][y]])]] \tag{$!Q_8$}
\end{align*}
\end{prob}

\begin{prob}
\ref{mod:mar:mdq:ex:model-L-of-Q} मधील $\Struct{L}$ मध्ये $a^\nssucc
= a$ आणि $b^\nssucc = b$. ज्या~$\Th{Q}$ च्या प्रतिमानात
$a^\nssucc = b$ आणि $b^\nssucc = a$ असे आहे, असे प्रतिमान अस्तित्वात आहे का?
\end{prob}

\begin{explain}
अमानक घटक हे मानक घटकांच्या ``पलीकडे'' असलेली~$\Th{Q}$ ची
प्रतिमाने आपण स्पष्टपणे रचली आहेत. किंबहुना, स्वयंसिद्धकांनुसार तितके असणे
आवश्यक आहे. अमानक घटक~$x$ हा ${} \nsless 0$ असू शकत नाही, कारण
$\Th{Q} \Proves \lforall[x][\lnot x<0]$
(\readerexternalref{inc:req:min:lem:less-zero} पाहा). तसेच, प्रत्येक $n$ साठी
$\Th{Q} \Proves \lforall[x][(x < \num{n}' \lif
(\eq[x][\num{0}] \lor \eq[x][\num{1}] \lor \dots \lor
\eq[x][\num{n}]))]$ (\readerexternalref{inc:req:min:lem:less-nsucc}); म्हणून
कोणत्याही~$n>0$ साठी $a \nsless n$ असू शकत नाही.
\end{explain}











\section{$\Th{PA}$ ची प्रतिमाने}\label{mod:mar:mpa:sec}

\begin{explain}
$\Th{TA}$ चे कोणतेही अमानक प्रतिमान हे~$\Th{PA}$ चेही अमानक प्रतिमान
असते. $\Th{TA}$ ची आणि म्हणून~$\Th{PA}$ ची अमानक प्रतिमाने अस्तित्वात
आहेत हे आपल्याला माहीत आहे. अशा अमानक प्रतिमानांत अमानक ``संख्या'', म्हणजे
सर्व मानक ``संख्यांच्या'' ``पलीकडे'' असलेले प्रांताचे घटक, असतात
हेही आपल्याला माहीत आहे. पण त्यांची मांडणी कशी असते? त्या किती असतात?
दुर्बलतर उपपत्ती~$\Th{Q}$ च्या प्रतिमानात फक्त एक अमानक संख्या असू शकते
हे आपण पाहिले. पण या साध्या संरचना~$\Th{PA}$ किंवा $\Th{TA}$ ची
प्रतिमाने नाहीत.

$\Th{PA}$ किंवा $\Th{TA}$ च्या प्रतिमानांची रचना समजून घेण्याची गुरुकिल्ली
म्हणजे या उपपत्तींमध्ये कोणती तथ्ये निष्पन्न करता येण्याजोगे आहेत हे पाहणे. उदाहरणार्थ,
$\Th{PA}$ कडून आधीच $\lforall[x][\eq/[x][x']]$ आणि
$\lforall[x][\lforall[y][\eq[(x+y)][(y+x)]]]$ सिद्ध होतात; म्हणून ही
वाक्ये असत्य असलेल्या साध्या रचना~$\Th{PA}$ ची प्रतिमाने असू शकत
नाहीत.

$\Struct{M}$ हे~$\Th{PA}$ चे प्रतिमान आहे असे समजा. मग
$\Th{PA} \Proves \varphi$ असल्यास $\Sat{M}{\varphi}$. पुन्हा
$\Assign{\Obj{0}}{M}$ करिता $\nszero$, $\Assign{\prime}{M}$ करिता
$\nssucc$, $\Assign{+}{M}$ करिता $\nsplus$,
$\Assign{\times}{M}$ करिता $\nstimes$ आणि $\Assign{<}{M}$ करिता
$\nsless$ वापरू. मग कोणतेही वाक्य~$\varphi$ हे $\nszero$, $\nssucc$,
$\nsplus$, $\nstimes$ आणि $\nsless$ यांविषयी एखादी अट सांगते; आणि
$\Sat{M}{\varphi}$ असेल, तर ती अट पूर्ण झाली पाहिजे. उदाहरणार्थ,
$\Sat{M}{!Q_1}$, म्हणजे
$\Sat{M}{\lforall[x][\lforall[y][(\eq[x'][y'] \lif \eq[x][y])]]}$,
असेल तर $\nssucc$ हे एकास-एक असले पाहिजे.
\end{explain}

\begin{prop}
$\Struct{M}$ मध्ये $\nsless$ हा काटेकोर रेषीय क्रम आहे, म्हणजे तो पुढील
अटी पूर्ण करतो:
\begin{enumerate}
\item कोणत्याही~$x \in \Domain{M}$ साठी $x \nsless x$ नाही.
\item $x \nsless y$ आणि $y \nsless z$ असतील, तर $x \nsless z$.
\item कोणत्याही $x \neq y$ साठी $x \nsless y$ किंवा $y \nsless x$.
\end{enumerate}
\end{prop}

\begin{proof}
$\Th{PA}$ कडून पुढील सूत्रे सिद्ध होतात:
\begin{enumerate}
\item $\lforall[x][\lnot x < x]$
\item $\lforall[x][\lforall[y][\lforall[z][((x < y \land y < z) \lif x < z)]]]$
\item $\lforall[x][\lforall[y][((x < y \lor y < x) \lor \eq[x][y]))]]$
\end{enumerate}
\end{proof}

\begin{prop}
\label{mod:mar:mpa:prop:M-discrete} $\nsless$-क्रमात $\nszero$ हा~$\Domain{M}$ मधील
लघुतम घटक आहे. प्रत्येक $x$ साठी $x \nsless x^\nssucc$, आणि हा
गुणधर्म असलेला $\nsless$-लघुतम घटक म्हणजे $x^\nssucc$. शून्याहून
वेगळ्या कोणत्याही~$x$ साठी $y^\nssucc = x$ असा एकमेव $y$ असतो. ($y$ ला
$\Struct{M}$ मधील~$x$ चा ``पूर्ववर्ती'' म्हणतो आणि तो~$^\nssucc x$ असा
दर्शवतो.)

\end{prop}

\begin{proof}
सराव.
\end{proof}

\begin{prob}
\ref{mod:mar:mpa:prop:M-discrete} मधील $\nszero$, $\nssucc$ आणि
$\nsless$ यांचे गुणधर्म सुनिश्चित करणारी, $\Th{PA}$ मध्ये निष्पन्न करता येण्याजोगे
(आणि म्हणून~$\Struct{N}$ मध्ये सत्य) असलेली~$\Lang{L_A}$ मधील
वाक्ये शोधा.
\end{prob}

\begin{prop}
$\Struct{M}$ चे सर्व मानक घटक हे सर्व अमानक घटक पेक्षा
($\nsless$ नुसार) लहान असतात.
\end{prop}

\begin{proof}
$\Struct{M}$ चा मानक घटक~$\Value{\num{n}}{M}$ याचा संक्षेप म्हणून
$n$ वापरू. कोणत्याही~$n \in \Nat$ साठी
$\lforall[x][(x < \num{n}' \lif (\eq[x][\num{0}] \lor
  \eq[x][\num{1}] \lor \dots \lor \eq[x][\num{n}]))]$ हे $\Th{Q}$
कडून आधीच सिद्ध होते. $\nsless \nszero$ असलेले घटक
नाहीत. म्हणून $n$ मानक आणि $x$ अमानक असेल, तर $x \nsless n$ असू शकत
नाही. व्याख्येनुसार, अमानक घटक म्हणजे कोणत्याही~$n \in \Nat$ साठी
$\Value{\num{n}}{M}$ नसलेला घटक; म्हणून $x \neq n$ सुद्धा. $\nsless$
हा रेषीय क्रम असल्यामुळे $n \nsless x$ असलेच पाहिजे.
\end{proof}

\begin{prop}
$\Domain{M}$ मधील प्रत्येक अमानक घटक~$x$ हा पुढील उपसंचाचा घटक
असतो:
\[
\dots ^{\nssucc\nssucc\nssucc}x \nsless ^{\nssucc\nssucc}x \nsless
^{\nssucc}x \nsless x \nsless x^{\nssucc} \nsless x^{\nssucc\nssucc}
\nsless x^{\nssucc\nssucc\nssucc} \nsless \dots
\]
या उपसंचाला \emph{$x$ चा खंड} म्हणतो आणि तो $[x]$ असा लिहितो. त्याला
लघुतम किंवा महत्तम घटक नाही. एखाद्या मानक~$n$ साठी
$x \nsplus n = y$ किंवा $y \nsplus n = x$ असे असणाऱ्या
$y \in \Domain{M}$ चा संच, असे त्याचे लक्षण सांगता येते.
\end{prop}

\begin{proof}
असा संच~$[x]$ नेहमी अस्तित्वात असतो हे स्पष्ट आहे; कारण
$\Domain{M}$ मधील प्रत्येक अशून्य घटक~$y$ ला एकमेव
उत्तरवर्ती~$y^\nssucc$ आणि एकमेव पूर्ववर्ती~$^\nssucc y$ असतो. लागोपाठच्या
घटक $y$, $y^\nssucc$ साठी $y \nsless y^\nssucc$, आणि
$y$ ज्याच्यापेक्षा $\nsless$ आहे असा $\Domain{M}$ मधील $\nsless$-लघुतम
घटक म्हणजे $y^\nssucc$. नेहमी $^\nssucc y \nsless y$ आणि
$y \nsless y^\nssucc$ असल्यामुळे $[x]$ ला लघुतम किंवा महत्तम घटक
नाही. $y \in [x]$ असेल, तर $x \in [y]$; कारण मग एकतर
$y^{\nssucc\dots\nssucc} = x$ किंवा $x^{\nssucc\dots\nssucc} = y$.
$y^{\nssucc\dots\nssucc} = x$ ($n$ वेळा $\nssucc$) असेल, तर
$y \nsplus n = x$, आणि याचा व्यत्यासही खरा; कारण $n$ ही $\prime$ ची
संख्या असेल, तर $\Th{PA} \Proves
\lforall[x][\eq[x^{\prime\dots\prime}][(x + \num{n})]]$.
\end{proof}

\begin{prop}
$[x] \neq [y]$ आणि $x \nsless y$ असतील, तर कोणत्याही $u \in [x]$ आणि
कोणत्याही $v \in [y]$ साठी $u \nsless v$.
\end{prop}

\begin{proof}
$\Th{PA} \Proves \lforall[x][\lforall[y][(x < y \lif (x' < y
    \lor x' = y))]]$ हे लक्षात घ्या. म्हणून $u \nsless v$ आणि
$[u] \neq [v]$ असतील, तर कोणत्याही~$n$ साठी
$u \nsplus n^\nssucc \nsless v$ सुद्धा.

प्रत्येक $u \in [x]$ हा~$\nsless y$ आहे: गृहीतकानुसार
$x \nsless y$. $u \nsless x$ असेल, तर संक्रमणशीलतेने $u \nsless y$.
आणि $x \nsless u$ पण $u \in [x]$ असेल, तर कोणत्यातरी~$n$ साठी
$u = x\nsplus n^\nssucc$; म्हणून नुकत्याच सिद्ध केलेल्या तथ्यावरून
$u \nsless y$.

आता $v \in [y]$ हा~$\nsless y$ आहे, म्हणजे एखाद्या मानक~$m$
साठी $v \nsplus m^\nssucc = y$, असे समजा. यामुळे $v \nsless x$ असण्याची
शक्यता बाद होते; अन्यथा $y = v \nsplus m^\nssucc \nsless x$ झाले असते.
तसेच स्पष्टपणे $x \neq v$; अन्यथा
$x \nsplus m^\nssucc = v \nsplus m^\nssucc = y$ आणि $[x] = [y]$ झाले
असते. म्हणून $x \nsless v$. पण मग कोणत्याही~$n$ साठी
$x \nsplus n^\nssucc \nsless v$ सुद्धा. त्यामुळे $x \nsless u$ आणि
$u \in [x]$ असतील, तर $u \nsless v$. $u \nsless x$ असेल, तर
संक्रमणशीलतेने $u \nsless v$.

शेवटी, $y \nsless v$ असेल, तर आपण दाखवल्याप्रमाणे $u \nsless y$ आणि
$y \nsless v$ असल्यामुळे $u \nsless v$.
\end{proof}

\begin{cor}
$[x] \neq [y]$ असेल, तर $[x] \cap [y] = \emptyset$.
\end{cor}

\begin{proof}
$z \in [x]$ आणि $x \nsless y$ असे समजा. मग प्रत्येक $u \in [y]$ साठी
$z \nsless u$. जर $z \in [y]$ असेल, तर $z \nsless z$ मिळेल.
$y \nsless x$ असेल तेव्हाही याचप्रमाणे.
\end{proof}

\begin{explain}
याचा अर्थ खंडांनाच $\nsless$ चा आदर राखणारा क्रम देता येतो:
$[x] \nsless [y]$ हे तेव्हाच, जेव्हा $x \nsless y$; किंवा सममूल्य रीतीने,
कोणत्याही $u \in [x]$ आणि $v \in [y]$ साठी $u \nsless v$. मानक खंड
$[0]$ हा लघुतम खंड आहे हे स्पष्ट आहे. त्याचा कोणत्याही अमानक खंडाशी छेद
नाही, आणि कोणत्याही दोन अमानक खंडांचाही परस्परांशी छेद नाही. विशेषतः,
उत्तरवर्ती किंवा पूर्ववर्ती वारंवार घेऊन वेगळ्या खंडात ``पोहोचता'' येत नाही.
\end{explain}

\begin{prop}
$x$ आणि $y$ अमानक असतील, तर $x \nsless x \nsplus y$ आणि
$x \nsplus y \notin [x]$.
\end{prop}

\begin{proof}
$y$ अमानक असेल, तर $y \neq \nszero$. $\Th{PA} \Proves
\lforall[x][(y \neq \Obj{0} \lif x < (x+y))]$. आता
$x \nsplus y \in [x]$ असे समजा. $x \nsless x \nsplus y$ असल्यामुळे
$x \nsplus n^\nssucc = x \nsplus y$ असे झाले असते. पण $\Th{PA} \Proves
\lforall[x][\lforall[y][\lforall[z][(\eq[(x+y)][(x+z)] \lif y = z)]]]$
(बेरीजेसाठीचा निरसन नियम). याचा अर्थ एखाद्या मानक~$n$ साठी
$y = n^\nssucc$ असा झाला असता; पण $y$ अमानक आहे असे गृहीत आहे.
\end{proof}

\begin{prop}
लघुतम अमानक खंड नाही.
\end{prop}

\begin{proof}
$\Th{PA} \Proves \lforall[x][\lexists[y][(\eq[(y+y)][x] \lor
      \eq[(y+y)'][x])]]$; म्हणजे प्रत्येक $x$ ला~$2$ ने भाग जातो
(शक्यतो बाकी~$1$ ठेवून). $x$ अमानक असेल, तर $y$ सुद्धा अमानक असतो.
आधीच्या विधानावरून $y \nsless y \nsplus y$ आणि
$y \nsplus y \notin [y]$. मग $y \nsless (y \nsplus y)^\nssucc$ आणि
$(y \nsplus y)^\nssucc \notin [y]$ सुद्धा. पण $x = y \nsplus y$ किंवा
$x = (y \nsplus y)^\nssucc$, म्हणून $y \nsless x$ आणि $y \notin [x]$.
\end{proof}

\begin{prop}
महत्तम खंड नाही.
\end{prop}

\begin{proof}
सराव.
\end{proof}

\begin{prob}
$\Th{PA}$ च्या अमानक प्रतिमानात महत्तम खंड नसतो हे दाखवा.
\end{prob}

\begin{prop}
\label{mod:mar:mpa:prop:blocks-dense}
खंडांचा क्रम सघन आहे. म्हणजे, $x \nsless y$ आणि $[x] \neq [y]$ असतील,
तर त्या दोहोंपेक्षा वेगळा आणि त्यांच्यामध्ये असलेला एक खंड~$[z]$ असतो.
\end{prop}

\begin{proof}
$x \nsless y$ असे समजा. पूर्वीप्रमाणे, $x \nsplus y$ ला दोनने भाग जातो
(शक्यतो बाकी ठेवून): एखादा $z \in \Domain{M}$ असा असतो की एकतर
$x \nsplus y = z \nsplus z$ किंवा
$x \nsplus y = (z \nsplus z)^\nssucc$. घटक~$z$ हा $x$ आणि~$y$ यांची
``सरासरी'' आहे, आणि $x \nsless z$ व $z \nsless y$.

\end{proof}

\begin{prob}
\ref{mod:mar:mpa:prop:blocks-dense} चा तपशीलवार पुरावा लिहा.
$z$ चे अस्तित्व सुनिश्चित करण्यासाठी $\Th{PA}$ ने कोणते वाक्य
निष्पन्न केले पाहिजे? $x \nsless z$ आणि $z \nsless y$ का, तसेच
$[x] \neq [z]$ आणि $[z] \neq [y]$ का?
\end{prob}

\begin{explain}
प्रतिमान गणनीय असेल, तर अमानक खंडांचा क्रम परिमेय संख्यांच्या क्रमासारखा
असतो: ते अंतबिंदुविरहित गणनीय अनंत सघन रेषीय क्रम बनवतात. असे
कोणतेही दोन गणनीय अनंत क्रम समरूपी असतात हे दाखवता येते. त्यावरून
सत्य अंकगणिताच्या कोणत्याही दोन गणनीय अमानक प्रतिमान
$\Struct{M}_1$ आणि $\Struct{M_2}$ यांची फक्त $<$ आणि $=$ असलेल्या भाषेतील
न्यूनीकरणे समरूपी आहेत. खरेच, पुढीलप्रमाणे समरूपण~$h$ ची व्याख्या करता
येते: $\Struct{M_1}$ आणि $\Struct{M_2}$ यांचे मानक भाग मानक
प्रतिमान~$\Struct{N}$ शी आणि म्हणून परस्परांशी समरूपी आहेत. अमानक भाग
बनवणाऱ्या खंडांचा क्रम स्वतः परिमेय संख्यांसारखा असल्यामुळे ते समरूपी आहेत;
प्रत्येक खंडातील स्वेच्छ घटक परस्पर जुळवून आणि नंतर $\Struct{M_1}$ मधील
$x$ च्या उत्तरवर्तीची प्रतिमा ही $\Struct{M_2}$ मधील $x$ च्या प्रतिमेचा
उत्तरवर्ती घेऊन, खंडांचे समरूपण खंडांच्या \emph{आत} विस्तारित करता येते.
$\mathfrak{M}_1$ आणि $\mathfrak{M}_2$ या अंकगणिताच्या पूर्ण भाषेत समरूपी
आहेत असे यावरून \emph{निष्पन्न होत नाही} हे लक्षात घ्या (खरे तर समरूपण
नेहमी भाषा सापेक्ष असते), कारण $\Domain{M_1}$ आणि
$\Domain{M_2}$ वर बेरीज आणि गुणाकार परिभाषित करण्याचे परस्पर-असमरूपी मार्ग
आहेत. (पूर्ण उपपत्तीच्या गणनीय प्रतिमानांची संख्या~2 असू शकत नाही, या
वॉट यांच्या प्रसिद्ध प्रमेयावरूनही हे निष्पन्न होते.)

\end{explain}










\section{अंकगणिताची संगणनक्षम प्रतिमाने}\label{mod:mar:cmp:sec}

\begin{explain}
मानक प्रतिमान~$\Struct{N}$ ला दोन उपयुक्त वैशिष्ट्ये आहेत. त्याचा प्रांत
नैसर्गिक संख्यांचा संच~$\Nat$ आहे, म्हणजे त्याचे घटक हे नेमके तेच प्रकारचे
विषय आहेत ज्यांच्याविषयी अंकगणिताची भाषा वापरून आपल्याला बोलायचे असते, आणि
मानक संख्यांक~$\num{n}$ हा खरोखर~$n$ लाच निर्देशित करतो. दुसरे उपयुक्त
वैशिष्ट्य म्हणजे~$\Lang{L_A}$ मधील अतार्किक चिन्हांची सर्व अर्थनिर्धारणे
\emph{संगणनक्षम} आहेत. $\Assign{\prime}{N}$, $\Assign{+}{N}$ आणि
$\Assign{\times}{N}$ म्हणून काम करणारी उत्तरवर्ती, बेरीज आणि गुणाकार ही
फलने संख्यांवरील संगणनक्षम फलने आहेत. (म्हणजे ट्यूरिंग यंत्रांनी संगणन करता
येणारी किंवा, उदाहरणार्थ, आदिम पुनरावर्तनाने परिभाषित करता येणारी.) तसेच
$\Struct{N}$ वरील लघुतर संबंध, म्हणजे~$\Assign{<}{N}$, निर्णेय आहे.

$\Th{Q}$ आणि $\Th{PA}$ यांसारख्या अंकगणितीय उपपत्तींच्या अमानक प्रतिमानांत
अमानक घटक असलेच पाहिजेत. त्यामुळे त्यांच्या प्रांतांत सामान्यतः~$\Nat$
व्यतिरिक्तही घटक असतात. तरी कोणतीही गणनीय संरचना कोणत्याही
गणनीय अनंत संचावर, त्यात~$\Nat$ वरही, उभारता येते. म्हणून प्रांत
$\Nat$ असलेली अमानक प्रतिमानेही आहेत. अशा~$\Struct{M}$ प्रतिमानांत अर्थातच
किमान काही संख्या आपली नेहमीची भूमिका बजावू शकत नाहीत, कारण असा काही~$k$
असतो की प्रत्येक~$n \in \Nat$ साठी तो~$\Value{\num{n}}{M}$ पेक्षा
वेगळा असतो.
\end{explain}

\begin{defn}
$\Lang{L_A}$ साठीची संरचना~$\Struct{M}$ ही \emph{संगणनक्षम}
तेव्हाच, जेव्हा $\Domain{M} = \Nat$ असते, $\Assign{\prime}{M}$,
$\Assign{+}{M}$ आणि $\Assign{\times}{M}$ ही संगणनक्षम फलने असतात, आणि
$\Assign{<}{M}$ हा निर्णेय संबंध असतो.
\end{defn}

\begin{ex}\label{mod:mar:cmp:ex:comp-model-q}
\ref{mod:mar:mdq:ex:model-K-of-Q} मधील रचना~$\Struct{K}$ आठवा. तिचा प्रांत
$\Domain{K} = \Nat
\cup \{a\}$ होता आणि अर्थनिर्धारणे पुढीलप्रमाणे होती:
\begin{align*}
  \Assign{\Obj{0}}{K} & = 0\\
  \Assign{\prime}{K}(x) & =
  \begin{cases}
    x+1 & \text{जर $x\in \Nat$}\\
    a & \text{जर $x = a$}
  \end{cases}\\
  \Assign{+}{K}(x, y) & =
  \begin{cases}
    x+y & \text{जर $x$, $y \in\Nat$}\\
    a & \text{इतर प्रसंगी}
  \end{cases}\\
  \Assign{\times}{K}(x, y) & =
  \begin{cases}
    xy & \text{जर $x$, $y \in\Nat$}\\
    0 & \text{जर $x=0$ किंवा $y=0$}\\
    a & \text{इतर प्रसंगी}\\
  \end{cases}\\
  \Assign{<}{K} & =
  \Setabs{\tuple{x,y}}{x, y \in \Nat \text{ आणि } x<y} \cup
  \Setabs{\tuple{x,a}}{x \in \Domain{K}}
\end{align*}

पण $\Domain{K}$ हा गणनीय अनंत आहे आणि म्हणून तो~$\Nat$ शी तुल्यबल
आहे. उदाहरणार्थ, $g\colon \Nat \to \Domain{K}$ असे घेऊ, जिथे $g(0) =
a$ आणि $n>0$ साठी $g(n) = n-1$; हे एक एकास-एक आच्छादन आहे.

त्याचा वापर करून नव्या~$\Th{Q}$-प्रतिमान~$\Struct{K'}$ आणि
$\Struct{K}$ यांमध्ये समरूपता मिळवता येते. $\Struct{K'}$ मध्ये
$\Lang{L_A}$ च्या चिन्हांना वेगळी फलने आणि संबंध नेमावे लागतात, कारण
$\Nat$ मधील वेगवेगळे घटक मानक आणि अमानक संख्यांच्या भूमिका
बजावतात.

विशेषतः, आता~$0$ हा लघुतम मानक संख्या नसून~$a$ ची भूमिका बजावतो. आता~$1$
ही लघुतम मानक संख्या आहे. म्हणून $\Assign{\Obj{0}}{K'} = 1$ असे नेमू.
आता उत्तरवर्ती फलनही वेगळे आहे: एखादी मानक संख्या, म्हणजे $n > 0$, दिली
असता ते अजूनही $n+1$ परत करते. पण आता~$0$ हा~$a$ ची भूमिका बजावतो आणि
तो स्वतःचाच उत्तरवर्ती आहे. म्हणून $\Assign{\prime}{K'}(0) = 0$.
बेरीज आणि गुणाकारासाठी याचप्रमाणे
\begin{align*}
\Assign{+}{K'}(x, y) & =
  \begin{cases}
    x+y-1 & \text{जर $x$, $y >0$}\\
    0 & \text{इतर प्रसंगी}
  \end{cases}\\
  \Assign{\times}{K'}(x, y) & =
  \begin{cases}
    1 & \text{जर $x = 1$ किंवा $y = 1$}\\
    xy - x - y + 2 & \text{जर $x$, $y > 1$}\\
    0 & \text{इतर प्रसंगी}\\
  \end{cases}
\end{align*}
आणि $\tuple{x, y} \in \Assign{<}{K'}$ हे तेव्हाच, जेव्हा $x < y$, $x > 0$
आणि $y > 0$ असतात, किंवा $y = 0$ असतो.

ही सर्व फलने नैसर्गिक संख्यांची संगणनक्षम फलने आहेत आणि
$\Assign{<}{K'}$ हा~$\Nat$ वरील निर्णेय संबंध आहे---पण ती~$\Nat$ वरील
उत्तरवर्ती, बेरीज आणि गुणाकार हीच फलने नाहीत, तसेच $\Assign{<}{K'}$ हाही
$\Nat$ वरील~$<$ हाच संबंध नाही.
\end{ex}

\begin{prob}
\ref{mod:mar:mdq:ex:model-L-of-Q} मधील~$\Struct{L}$ शी समरूपी,
$\Domain{L'} = \Nat$ असलेली संरचना~$\Struct{L'}$ द्या.
\end{prob}

\begin{explain}
\ref{mod:mar:cmp:ex:comp-model-q} यावरून~$\Th{Q}$ ला प्रांत~$\Nat$ असलेली
संगणनक्षम अमानक प्रतिमाने आहेत हे दिसते. पण पुढील निष्कर्ष दाखवतो की
$\Th{PA}$ च्या प्रतिमानांसाठी (आणि म्हणून~$\Th{TA}$ च्या प्रतिमानांसाठीही)
हे खरे नाही.
\end{explain}

\begin{thm}[टेनेनबाउमचे प्रमेय]
$\Struct{N}$ हे~$\Th{PA}$ चे एकमेव संगणनक्षम प्रतिमान आहे.
\end{thm}



\chapter*{स्रोतावरील संपादकीय नोंदी}
\addcontentsline{toc}{chapter}{स्रोतावरील संपादकीय नोंदी}
या नोंदी अनुवादकाने वेगळ्या जोडल्या आहेत; त्या मूळ मजकुराचा भाग नाहीत.
स्रोताशी जुळवलेल्या TeX फाइलांतील मूळ चिन्हे बदललेली नाहीत.
\begin{enumerate}
\item विभाग \ref{sfr:rel:set:sec} मधील $K=L\cup I$ आणि $H=G\cup I$ येथे
$I$ चे स्वतंत्र नामकरण राहिले आहे. आधीच्या कर्णाच्या चर्चेनुसार त्याचा अर्थ
$\Id{\Nat}$ हा एकरूपता संबंध असा घ्यावा.
\item विभाग \ref{sfr:rel:tre:sec} मधील शाखेच्या व्याख्येत $z\in X\setminus B$
असे छापले आहे. वृक्षाचा आधारसंच $A$ आहे; येथे $X$ ऐवजी $A$ अभिप्रेत दिसतो.
\item त्याच विभागातील आरंभीच्या खंडांनुसार बंद असणाऱ्या उपवृक्षाच्या उदाहरणात
$A$ अरिक्त असणे आवश्यक आहे: आधीच्या व्याख्येनुसार वृक्षाला मूळ असते,
म्हणून रिक्त संच त्या व्याख्येत बसत नाही.
\item फलन प्रकरणातील OLFUN-001 ते OLFUN-005 या पाच
गोठवलेल्या-स्रोत निष्कर्षांची दुरुस्ती संबंधित परिच्छेदांलगत स्वतंत्र
``स्रोतदुरुस्ती'' नोंदींमध्ये दिली आहे. मूळ इंग्रजी बाइट्स बदललेले नाहीत.
\item संचांच्या आकारमानाच्या प्रकरणातील OLSIZ-001 ते OLSIZ-010
या दहा सामायिक गोठवलेल्या-स्रोत निष्कर्षांची आणि MRSIZ-001 ते MRSIZ-004
या चार स्थानिक निरीक्षणांची नोंद संबंधित परिच्छेदांलगत दिली आहे. मूळ
इंग्रजी बाइट्स बदललेले नाहीत.
\item सामायिक OLSIZ-011 सूचना अचूक बाइट-पुनर्तपासणीनंतर चुकीची ठरून
मागे घेण्यात आली; तिच्यावर आधारित कोणतीही मराठी दुरुस्ती केलेली नाही.
\item अंकगणितीकरण प्रकरणातील MRARITH-001 ते MRARITH-012 या बारा
गोठवलेल्या-स्रोत निरीक्षणांतील दुरुस्त्या किंवा स्पष्ट खुलासे संबंधित
परिच्छेदांलगत स्वतंत्र ``स्रोतदुरुस्ती'' नोंदींमध्ये दिले आहेत. मूळ
इंग्रजी बाइट्स बदललेले नाहीत.
\item अनंत संच प्रकरणातील MRINF-001 ते MRINF-003 या तीन स्रोत-निरीक्षणांची
नोंद ठेवली आहे. MRINF-001 मधील व्याकरणदोषाचा अभिप्रेत अर्थ थेट दिला आहे;
MRINF-002 मधील अप्रकट आधारसंचाचे गृहीतक तज्ज्ञ-पुनरावलोकनासाठी राखले आहे;
MRINF-003 मधील स्रोतदोषाचा दावा पुनर्तपासणीनंतर चुकीचा ठरून मागे घेतला आहे.
अंतःस्थ \texttt{cardeq} रचना वैध आहे; मराठी निष्कर्ष सममूल्य केंद्रित
मांडणी म्हणून जतन केला आहे. मूळ इंग्रजी बाइट्स बदललेले नाहीत.
\item विधानीय तर्कशास्त्र प्रकरणातील MRPL-001 आणि MRPL-002 या दोन
गोठवलेल्या-स्रोत चिन्हदोषांच्या दुरुस्त्या संबंधित परिच्छेदांलगत स्वतंत्र
``स्रोतदुरुस्ती'' नोंदींमध्ये दिल्या आहेत. मूळ इंग्रजी बाइट्स बदललेले नाहीत.
\item \textbf{MRPRF-001}: गोठवलेल्या इंग्रजी टॅब्लो उदाहरणातील
शेवटच्या दोन पायऱ्या $\sFmla{\True}{\varphi \land \psi}$ चे घटक उलगडतात,
पण त्यांची नियम-खूण चुकून $\TRule{\True}{\lif}$ अशी आहे. या वाचकात ती
अर्थानुसार $\TRule{\True}{\land}$ केली आहे; संरेखित इंग्रजी आणि मराठी
स्रोत-बाइट्स बदललेले नाहीत.

\item \textbf{OLSEQ-001}: OLP-0075 मधील चार पायऱ्यांत पूर्वपक्षातील
सूत्रांची अदलाबदल असूनही नियम-खूण उजवी अदलाबदल अशी आहे. या वाचकात त्या
चार खुणा डावी अदलाबदल अशा केल्या आहेत; संरेखित स्रोतांत मूळ रूप जतन आहे.
\item \textbf{OLSEQ-002}: त्याच उदाहरणातील दोन कथनसूत्रांत दुसऱ्या
सूत्रापूर्वीचे नकारचिन्ह गाळले आहे. लगतचे ध्येय, सममितीचे स्पष्टीकरण आणि
सिद्धता-वृक्ष यांनुसार या वाचकात ती दोन नकारचिन्हे पुनःस्थापित केली आहेत.
\item \textbf{OLSEQ-003}: OLP-0072 च्या शेवटच्या परिच्छेदातील ``पदावर
निर्बंध नाही'' याचा अर्थ पद बंद असण्याव्यतिरिक्त आणखी निर्बंध नाही असा आहे;
मराठी मजकूर ही आधी स्पष्ट केलेली अट सांगतो.
\item \textbf{OLND-001}: OLP-0087 मध्येही संख्यापक-नियमांसाठी पद बंद
असण्याची उभी अट शेवटच्या स्पष्टीकरणात स्पष्ट केली आहे.
\item \textbf{OLND-002}: OLP-0087 मधील आयगेन चलाचा एकत्रित सारांश सर्व
आधारविधानांतून चल वगळल्यासारखा वाचतो; मराठी सारांश प्रत्येक नियमाची योग्य
आधारविधान, निष्कर्ष आणि अमुक्त गृहीतकांची अट वेगवेगळी राखतो.
\item \textbf{OLND-003}: OLP-0089 मध्ये आकृतीतील उजवी शाखा चुकून डावी
म्हटली आहे; वाचकात आकृती आणि पुढील परिच्छेदाशी सुसंगत ``उजवीकडील'' आहे.
\item \textbf{OLND-004}: OLP-0089 मधील एकच अनुमान कधी नकार-विलोपन तर
कधी असत्यता-प्रवेशन असे खूणांकित आहे. अनुमान अर्थतः समान असल्याने मूळची
पर्यायी खूण जतन केली आहे.
\item \textbf{OLND-005}: OLP-0090 मधील आयगेन-चल तपासणीत मुख्य
आधारविधानातील नकार गाळला आहे आणि नियमाची व्याप्ती अपुरी सांगितली आहे;
मराठी मजकूर योग्य नकार व तात्पुरत्या गृहीतकाचा अपवाद स्पष्ट करतो.
\item \textbf{OLND-006}: OLP-0090 मधील सार्वत्रिक विलोपनासाठी ``कोणतेही
पद'' म्हणजे उभ्या नियमानुसार कोणतेही बंद पद; मराठी स्पष्टीकरण ती अट राखते.
\item \textbf{OLND-007}: OLP-0095 च्या एका सरावातील नियम-खूण
\verb|\Elim{\forall}| ही प्रकरणात इतरत्र असलेल्या
\verb|\Elim{\lforall}| शी दृश्यतः समतुल्य आहे. संरेखित रूप जतन केले आहे.

\item \textbf{OLTAB-001--OLTAB-004}: टॅब्लो प्रकरणातील आधीच्या
उदाहरणांमध्ये विषय, बंद पदांची अट आणि एका ओळ-संदर्भाशी संबंधित स्रोतदोष
दुरुस्त किंवा स्पष्ट केले आहेत; संरेखित इंग्रजी आणि मराठी स्रोत-बाइट्स जतन आहेत.
\item \textbf{OLTAB-005--OLTAB-009}: टॅब्लो सिद्धतायोग्यता आणि सुसंगतता
विभागातील संच-लेखन, नियम-पूर्वपद, ओळ-संदर्भ आणि शाखा-वर्णनातील निश्चित
स्रोतदोषांचे मराठी रूपात मर्यादित निराकरण केले आहे; मूळ इंग्रजी बाइट्स जतन आहेत.
\item \textbf{OLTAB-010}: विधानीय परिणामांच्या आठ टॅब्लो नोड्समधील
चिन्हांकित-सूत्र मॅक्रोची गहाळ युक्ती-सिमा पुनर्स्थापित केली आहे; संरेखित
स्रोतांत गोठवलेले रूप जतन आहे.
\item \textbf{OLTAB-011}: टॅब्लो निर्दोषता सिद्धतेतील सार्वत्रिक नियमांच्या
पुनरावृत्त B/A रूपबंध-विसंगतीत निष्कर्ष आणि अर्थविषयक पायऱ्यांशी सुसंगत
A-रूप वापरले आहे; गोठवलेले रूप संरेखित स्रोतांत जतन आहे.
\item \textbf{OLTAB-012}: एकरूपता-टॅब्लोच्या सममिती आणि संक्रमकता
स्पष्टीकरणांत लगतच्या वृक्षांशी विसंगत असलेली दोन कथनसूत्रे दुरुस्त केली आहेत;
औपचारिक वृक्ष अपरिवर्तित आहेत.

\item \textbf{OLAXD-001--OLAXD-009}: स्वयंसिद्धकीय निष्पत्ती प्रकरणातील
कंस, अंतिम निष्कर्ष, स्वयंसिद्धक-संदर्भ, संख्यापक-नियम आणि जुनी
फाइल-माहिती यांतील नऊ निश्चित स्रोतदोष मराठी संरेखित स्रोतांत मर्यादितपणे
दुरुस्त केले आहेत; गोठवलेले इंग्रजी बाइट्स जतन आहेत.
\item \textbf{OLCOM-001--OLCOM-006}: संपूर्णता प्रकरणातील सूत्र-युक्तिवाद,
चल-सुसंगती, बंद-पद अट, विरामचिन्हे, प्रतिनिधी आणि टॅग-संदर्भ यांतील सहा
निश्चित स्रोतदोषांची दुरुस्ती किंवा स्पष्टता मराठी स्रोतांत नोंदवली आहे.
 \item \textbf{OLFOL-001--OLFOL-040}: प्रथम-क्रम तर्कशास्त्र आणि प्रतिमान
उपपत्तीच्या पुढील प्रकरणांतील एकोणचाळीस सूत्रात्मक, निर्देशांकविषयक,
अर्थविषयक आणि सिद्धताविषयक स्रोतदोष किंवा अपूर्ण पायऱ्या मर्यादितपणे
दुरुस्त किंवा स्पष्ट केल्या आहेत. OLFOL-009 मधील अतिरिक्त आकुंचित कंसाची
सूचना पुनर्तपासणीनंतर चुकीची ठरून मागे घेतली आहे; गोठवलेली रचना आधीपासून
संतुलित आहे. प्रत्येक नोंदीचा अचूक स्रोत-स्थान, पर्याय आणि
पुनरावलोकन-प्रश्न सोबतच्या तज्ज्ञ-पुनरावलोकन नोंदींत आहे.
\end{enumerate}
\chapter*{संदर्भ}
\addcontentsline{toc}{chapter}{संदर्भ}
\phantomsection\label{bib:Benacerraf1965}
Paul Benacerraf (1965). ``What numbers could not be.''
\textit{The Philosophical Review}, 74(1), 47--73.

\par\medskip
Gottlob Frege (1884). \textit{Die Grundlagen der Arithmetik}.
Wilhelm Koebner.

Georg Cantor (1892). ``Über eine elementare Frage der
Mannigfaltigkeitslehre.'' \textit{Jahresbericht der deutschen
Mathematiker-Vereinigung}, 1, 75--78.

Michael Potter (2004). \textit{Set Theory and its Philosophy}.
Oxford University Press.

John Conway (2006). ``The Power of Mathematics.'' In Alan Blackwell
and David MacKay (eds.), \textit{Power}, Darwin College Lectures.
Cambridge University Press.

John J. O'Connor and Edmund F. Robertson (2005). ``The real numbers:
Stevin to Hilbert.''
\url{http://www-history.mcs.st-and.ac.uk/HistTopics/Real_numbers_2.html}.

Karin Usadi Katz and Mikhail G. Katz (2012). ``Stevin Numbers and
Reality.'' \textit{Foundations of Science}, 17(2), 109--123.

Richard Dedekind (1888). \textit{Was sind und was sollen die Zahlen?}
Vieweg, Braunschweig.

David Hilbert (2013). \textit{David Hilbert's Lectures on the Foundations
of Arithmetic and Logic 1917--1933}. William Bragg Ewald and Wilfried Sieg
(eds.). Springer, Heidelberg.


\par\medskip
P. D. Magnus, Tim Button, J. Robert Loftis, Aaron Thomas-Bolduc आणि
Richard Zach (2021). \textit{forall x: Calgary}. Fall 2021 आवृत्ती.

Raymond M. Smullyan (1968). \textit{First-Order Logic}. Springer.

M. M. Zuckerman (1973). ``On the Unique Readability of Formal Systems.''
\textit{Notre Dame Journal of Formal Logic}, 14(4), 546--548.
\end{document}
